%=====================================================================
% ĐỀ THI THỬ TỐT NGHIỆP THPT 2026 -- MÔN VẬT LÍ
% Nguồn: Trường THPT Yên Hòa, Hà Nội -- Lần 1, năm học 2025 - 2026 (PDF, không đáp án). Gõ lại theo template Eduzy.
%=====================================================================
\def\CheDo{solcolor}% loigiai | dethi
\def\tenchuong{THPT Yên Hòa, Hà Nội -- Lần 1}
\def\thuonghieu{VẬT LÍ 12 -- Eduzy.com.vn}
\def\TachCauDS{1}
\def\TachCauLG{1}
%============================================================
%  PREAMBLE EDUZY v2 - dùng chung cho đề thi, bài học, chuyên đề (gói ex_test)
%  Khác bản preamble-eduzy.tex ở đầu trang, chân trang: dải nền nâu đồng tràn ngang,
%  logo + Eduzy.com.vn, ô chữ vàng kem cạnh xiên, chân trang "Trang (số)";
%  logo Eduzy chìm giữa mọi trang (tắt: \logochimtat, bật lại: \logochimbat).
%  Cách dùng ở file chính (xem mau-chuan.tex, mau-chuyen-de.tex):
%      \def\CheDo{loigiai}          % loigiai | dethi | solcolor | color | book
%      \def\tenchuong{<dòng chữ trên đầu trang>}
%      \def\thuonghieu{<chữ ở chân trang lẻ>}   % không khai báo thì dùng mặc định
%      \def\tacgia{<chữ ở chân trang chẵn>}      % không khai báo thì dùng mặc định
%      \def\khauhieu{<chữ cạnh logo ở đầu trang>}  % tùy chọn, mặc định Eduzy.com.vn
%      %============================================================
%  PREAMBLE EDUZY v2 - dùng chung cho đề thi, bài học, chuyên đề (gói ex_test)
%  Khác bản preamble-eduzy.tex ở đầu trang, chân trang: dải nền nâu đồng tràn ngang,
%  logo + Eduzy.com.vn, ô chữ vàng kem cạnh xiên, chân trang "Trang (số)";
%  logo Eduzy chìm giữa mọi trang (tắt: \logochimtat, bật lại: \logochimbat).
%  Cách dùng ở file chính (xem mau-chuan.tex, mau-chuyen-de.tex):
%      \def\CheDo{loigiai}          % loigiai | dethi | solcolor | color | book
%      \def\tenchuong{<dòng chữ trên đầu trang>}
%      \def\thuonghieu{<chữ ở chân trang lẻ>}   % không khai báo thì dùng mặc định
%      \def\tacgia{<chữ ở chân trang chẵn>}      % không khai báo thì dùng mặc định
%      \def\khauhieu{<chữ cạnh logo ở đầu trang>}  % tùy chọn, mặc định Eduzy.com.vn
%      %============================================================
%  PREAMBLE EDUZY v2 - dùng chung cho đề thi, bài học, chuyên đề (gói ex_test)
%  Khác bản preamble-eduzy.tex ở đầu trang, chân trang: dải nền nâu đồng tràn ngang,
%  logo + Eduzy.com.vn, ô chữ vàng kem cạnh xiên, chân trang "Trang (số)";
%  logo Eduzy chìm giữa mọi trang (tắt: \logochimtat, bật lại: \logochimbat).
%  Cách dùng ở file chính (xem mau-chuan.tex, mau-chuyen-de.tex):
%      \def\CheDo{loigiai}          % loigiai | dethi | solcolor | color | book
%      \def\tenchuong{<dòng chữ trên đầu trang>}
%      \def\thuonghieu{<chữ ở chân trang lẻ>}   % không khai báo thì dùng mặc định
%      \def\tacgia{<chữ ở chân trang chẵn>}      % không khai báo thì dùng mặc định
%      \def\khauhieu{<chữ cạnh logo ở đầu trang>}  % tùy chọn, mặc định Eduzy.com.vn
%      \input{preamble-eduzy-v2}
%      \TLdapso                      % tài liệu: \TLfull | \TLdapso | \TLan
%      \begin{document} ... \end{document}
%  Logo Eduzy vẽ bằng TikZ ngay trong file (\eduzylogo), không cần file ảnh.
%  Biên dịch: pdfLaTeX (máy có gói vietnam) hoặc XeLaTeX.
%============================================================
\documentclass[12pt,a4paper,twoside]{article}%<-- twoside: bat dau trang/chan trang chan-le
\usepackage{iftex}
\usepackage{amsmath,amssymb,tikz,mathrsfs,tkz-tab}
\usepackage{graphicx}
\usepackage{pgfplots}% đồ thị hàm số
\pgfplotsset{compat=1.17}
%\usepackage{mathptmx}
\usepackage{mathpazo}
\ifXeTeX% XeLaTeX: font chữ có đủ dấu tiếng Việt, giữ nguyên font toán mathpazo
	\usepackage[no-math]{fontspec}
	\setmainfont{TeX Gyre Pagella}
	\newfontfamily\EZqag{TeX Gyre Adventor}[NFSSFamily=qag,Ligatures=TeX]
\fi
\usepackage[many]{tcolorbox}
\usetikzlibrary{math,through,calc,intersections,angles,quotes,shapes,shapes.geometric,arrows,patterns,matrix,chains,arrows.meta,decorations.shapes,decorations.fractals,decorations.markings,shadows,decorations.pathreplacing}
\usepackage[top=1.6cm, bottom=1.8cm, left=1.8cm, right=1.8cm]{geometry}% v2: top 1.6cm (bản cũ 1.5cm) để chữ không sát dải đầu trang
\providecommand{\CheDo}{loigiai}% loigiai | dethi | solcolor | color | book
\expandafter\usepackage\expandafter[\CheDo]{ex_test}
% \tenchuong: khai báo ở file chính (trước \input) hoặc \renewcommand sau \input
%============================================================
%  [GỘP TỪ eduzy-de.tex] EDUZY - MÀU NÂU ĐỒNG / VÀNG KIM
%  + ĐẦU TRANG, CHÂN TRANG KIỂU DẢI NGANG (v2), ĐỐI XỨNG TRANG CHẴN / TRANG LẺ.
%  (không cần file eduzy-de.tex riêng nữa)
%  Khối này phải đặt SAU \usepackage[...]{ex_test} và SAU \def\tenchuong{...}
%============================================================
\usepackage{lastpage}
\usepackage{needspace}
%--- DÀN TRANG ĐỀ THI: tách câu (THÊM 05/11/2026; chỉ thêm lệnh mới, KHÔNG đổi hành vi mặc định) ---
%  Hai công tắc, mặc định tắt (= preamble chạy y như cũ). Bật bằng cách đặt TRƯỚC \input{preamble-eduzy-v2}
%  trong file đề (tools/dan_trang/dan_trang.py tự ghi khi đặt \NeedspaceCau):
%    \def\TachCauDS{1}   câu Đúng/Sai được tách: câu dẫn (kèm hình) ở cuối trang, bảng Phát biểu sang trang sau
%                        (câu dẫn và hình không tách nhau, bảng không tách giữa hai trang)
%    \def\TachCauLG{1}   bản lời giải (\CheDo khác dethi): phần đề của câu cao hơn 1/3 trang được phép tách
%  Lệnh dùng trong file đề (giá trị do tools/dan_trang/dan_trang.py đo ở bản dethi):
%    \NeedspaceCau[<cao câu dẫn + hình>]{<cao cả khối câu>}   ngay trên \begin{ex} (câu đầu phần: trên tiêu đề phần)
%        công tắc tắt: đúng bằng \Needspace{<cao cả khối câu>} như cũ;
%        \TachCauDS=1 và có [..] (câu Đúng/Sai được chọn tách): chỉ giữ câu dẫn + hình;
%        \TachCauLG=1 ở bản lời giải: đòi tối đa 1/3 \textheight.
%    \NeedspaceBang{<cao bảng Phát biểu>}%   ngay trước \choiceTFt: \TachCauDS=1 thì giữ trọn bảng trên một trang;
%        công tắc tắt: không làm gì.
\providecommand{\TachCauDS}{0}
\providecommand{\TachCauLG}{0}
\newlength{\EZnsCau}
\newcommand{\NeedspaceCau}[2][]{%
	\setlength{\EZnsCau}{#2}%
	\ifnum\TachCauDS=1 \if\relax\detokenize{#1}\relax\else\setlength{\EZnsCau}{#1}\fi\fi
	\ifnum\TachCauLG=1 \ifdefstring{\kindOPTN}{dethi}{}{\ifdim\EZnsCau>0.3333\textheight\setlength{\EZnsCau}{0.3333\textheight}\fi}\fi
	\Needspace{\EZnsCau}}
\newcommand{\NeedspaceBang}[1]{\ifnum\TachCauDS=1 \Needspace{#1}\fi}

%--- Thang màu: đậm -> nhạt
\definecolor{gold1}{RGB}{126,96,0}    % nâu đồng đậm  #7E6000
\definecolor{gold2}{RGB}{192,144,3}   % vàng kim      #C09003
\definecolor{gold3}{RGB}{255,215,99}  % vàng nhạt     #FFD763
\definecolor{gold4}{RGB}{255,231,153} % vàng rất nhạt #FFE799
\definecolor{gold5}{RGB}{255,243,205} % vàng kem      #FFF3CD
%--- Ghi đè tông xanh/đỏ cũ
\definecolor{myblue}{RGB}{126,96,0}
\colorlet{blue}{gold1}
\colorlet{cyan}{gold2}
\colorlet{teal}{gold2}

%--- Tên chương mặc định (file đề nên tự khai báo lại)
\providecommand{\tenchuong}{}
%--- Chân trang
\providecommand{\thuonghieu}{TOÁN 12 -- Eduzy.com.vn}
\providecommand{\tacgia}{Tô Gia Bảo -- 0353703783}

%--- Chữ cạnh logo ở đầu trang
\providecommand{\khauhieu}{Eduzy.com.vn}

%====================== THÔNG SỐ ĐẦU TRANG, CHÂN TRANG (v2) ======================
%  Dải nền đậm tràn ngang trang; ô chữ vàng kem có cạnh xiên, kèm 2 vạch xiên chuyển màu.
%  Đầu trang lẻ : logo + \khauhieu (trái)   | ô chữ \tenchuong (phải)
%  Chân trang lẻ: ô chữ \thuonghieu (trái)  | Trang (số) (phải)
%  Trang chẵn soi gương trang lẻ; chân trang chẵn ghi \tacgia.
%--- Màu: chỉ dùng thang vàng kim Eduzy
\colorlet{ezNen}{gold1}% dải nền đậm
\colorlet{ezVach1}{gold3}% vạch xiên sát dải nền (đậm hơn)
\colorlet{ezVach2}{gold4}% vạch xiên sát ô chữ (nhạt hơn)
\colorlet{ezO}{gold5}% ô chữ vàng kem
%--- Kích thước
\newlength{\EZdai}\setlength{\EZdai}{1.1cm}% chiều cao dải nền
\newlength{\EZo}\setlength{\EZo}{0.68cm}% chiều cao ô chữ
\newlength{\EZxien}\setlength{\EZxien}{0.55cm}% độ lệch ngang của cạnh xiên
\newlength{\EZvach}\setlength{\EZvach}{0.75cm}% bề rộng mỗi vạch xiên
\newlength{\EZkhe}\setlength{\EZkhe}{0.12cm}% khe trắng giữa vạch và ô chữ
\newlength{\EZdem}\setlength{\EZdem}{0.4cm}% khoảng đệm từ chữ tới cạnh xiên
\newlength{\EZle}\setlength{\EZle}{1.8cm}% lề trái/phải của chữ (bằng lề thân bài)
\newlength{\EZrongmin}\setlength{\EZrongmin}{4.5cm}% ô chữ rộng tối thiểu
\newlength{\EZrong}% bề rộng chữ trong ô (tính tự động)
%--- Font
\newcommand{\EZfontO}{\fontsize{10.5pt}{12pt}\selectfont\bfseries}% chữ trong ô
\newcommand{\EZfontKH}{\fontfamily{qag}\fontsize{11pt}{12pt}\selectfont\bfseries}% \khauhieu cạnh logo
\newcommand{\EZfontTrang}{\fontsize{11.5pt}{12pt}\selectfont\bfseries}% chữ "Trang" và số trang
%--- Đo bề rộng chữ trong ô
\newcommand{\EZdorong}[1]{%
	\settowidth{\EZrong}{\EZfontO #1}%
	\ifdim\EZrong<\EZrongmin\setlength{\EZrong}{\EZrongmin}\fi}
%--- Logo + chữ trắng \khauhieu (dựng sẵn một lần khi vào \begin{document})
\newsavebox{\EZlogobox}
\AtBeginDocument{\sbox{\EZlogobox}{%
		\raisebox{-0.5\height}{\eduzylogo[1cm]}\hspace{0.15cm}%
		\raisebox{-0.5\height}{\color{white}\EZfontKH\khauhieu}}}
%--- Hai vạch xiên + khe trắng. #1: điểm gốc ở mép dưới, #2: hướng (1 sang phải, -1 sang trái),
%    #3: hướng nghiêng (1 kiểu "/", -1 kiểu "\")
\newcommand{\EZvachxien}[3]{%
	\fill[white] (#1) --++({(#3)*\EZxien},\EZo) --++({(#2)*\EZkhe},0) --++({-(#3)*\EZxien},-\EZo) --cycle;
	\fill[ezVach2] ([xshift=(#2)*\EZkhe]#1) --++({(#3)*\EZxien},\EZo) --++({(#2)*\EZvach},0) --++({-(#3)*\EZxien},-\EZo) --cycle;
	\fill[ezVach1] ([xshift=(#2)*(\EZkhe+\EZvach)]#1) --++({(#3)*\EZxien},\EZo) --++({(#2)*\EZvach},0) --++({-(#3)*\EZxien},-\EZo) --cycle;
	\draw[white,line width=0.6pt] ([xshift=(#2)*(\EZkhe+\EZvach)]#1) --++({(#3)*\EZxien},\EZo);
}
%--- Chữ trong ô (căn giữa theo chiều cao ô, không lệch vì dấu tiếng Việt)
\tikzset{ezchu/.style={inner sep=0pt,text=ezNen,text height=0.72em,text depth=0.22em}}

%====================== LOGO CHÌM GIỮA TRANG ======================
%  Vẽ cùng đầu trang nên nằm dưới chữ của thân bài, có trên mọi trang.
%  Tắt: \logochimtat (từ trang đó trở đi), bật lại: \logochimbat.
%  Đổi cỡ, độ đậm: khai báo trước \input, ví dụ \def\EZchimcao{9cm} \def\EZchimdo{0.1}
\providecommand{\EZchimcao}{10cm}% chiều cao logo chìm
\providecommand{\EZchimdo}{0.12}% độ đậm (0 = trong suốt, 1 = rõ hẳn)
\newif\ifEZchim \EZchimtrue
\newcommand{\logochimtat}{\EZchimfalse}
\newcommand{\logochimbat}{\EZchimtrue}
%  Dựng sẵn logo mờ một lần khi vào \begin{document}; transparency group làm mờ cả logo
%  một lượt nên không lộ vệt ở chỗ các dải màu chồng nhau.
\newsavebox{\EZchimtho}
\newsavebox{\EZchimbox}
\AtBeginDocument{%
	\sbox{\EZchimtho}{\eduzylogo[\EZchimcao]}%
	\sbox{\EZchimbox}{\begin{tikzpicture}
			\begin{scope}[transparency group,opacity=\EZchimdo]
				\node[inner sep=0pt] {\usebox{\EZchimtho}};
			\end{scope}
	\end{tikzpicture}}}
\newcommand{\EZvechim}{%
	\ifEZchim
		\node[inner sep=0pt] at (current page.center) {\usebox{\EZchimbox}};
	\fi}

%====================== ĐẦU TRANG - TRANG LẺ ======================
%  (logo bên trái, ô chữ \tenchuong bên phải)
\newcommand{\myfancyheadOdd}{%
	\EZdorong{\tenchuong}%
	\begin{tikzpicture}[remember picture,overlay]
		\coordinate (T) at (current page.north west);
		\coordinate (P) at (current page.north east);
		%--- logo chìm giữa trang
		\EZvechim
		%--- dải nền
		\fill[ezNen] (T) rectangle ([yshift=-\EZdai]P);
		%--- ô chữ: X là góc dưới bên trái, cạnh trái xiên kiểu "/"
		\coordinate (X) at ([xshift=-\EZle-\EZrong-\EZdem-\EZxien,yshift=-\EZdai]P);
		\fill[ezO] (X) --++(\EZxien,\EZo) -- ([yshift=-\EZdai+\EZo]P) -- ([yshift=-\EZdai]P) --cycle;
		\EZvachxien{X}{-1}{1}
		\node[ezchu,anchor=east] at ([xshift=-\EZle,yshift=-\EZdai+0.5*\EZo]P) {\EZfontO\tenchuong};
		%--- logo + tên
		\node[inner sep=0pt,anchor=west] at ([xshift=\EZle,yshift=-0.5*\EZdai]T) {\usebox{\EZlogobox}};
	\end{tikzpicture}%
}

%====================== ĐẦU TRANG - TRANG CHẴN =====================
%  (soi gương: ô chữ \tenchuong bên trái, logo bên phải)
\newcommand{\myfancyheadEven}{%
	\EZdorong{\tenchuong}%
	\begin{tikzpicture}[remember picture,overlay]
		\coordinate (T) at (current page.north west);
		\coordinate (P) at (current page.north east);
		%--- logo chìm giữa trang
		\EZvechim
		%--- dải nền
		\fill[ezNen] (T) rectangle ([yshift=-\EZdai]P);
		%--- ô chữ: X là góc dưới bên phải, cạnh phải xiên kiểu "\"
		\coordinate (X) at ([xshift=\EZle+\EZrong+\EZdem+\EZxien,yshift=-\EZdai]T);
		\fill[ezO] (X) --++(-\EZxien,\EZo) -- ([yshift=-\EZdai+\EZo]T) -- ([yshift=-\EZdai]T) --cycle;
		\EZvachxien{X}{1}{-1}
		\node[ezchu,anchor=west] at ([xshift=\EZle,yshift=-\EZdai+0.5*\EZo]T) {\EZfontO\tenchuong};
		%--- logo + tên
		\node[inner sep=0pt,anchor=east] at ([xshift=-\EZle,yshift=-0.5*\EZdai]P) {\usebox{\EZlogobox}};
	\end{tikzpicture}%
}

%====================== CHÂN TRANG - TRANG LẺ ======================
%  (ô chữ \thuonghieu bên trái, "Trang (số)" bên phải)
\newcommand{\myfancyfootOdd}{%
	\EZdorong{\thuonghieu}%
	\begin{tikzpicture}[remember picture,overlay]
		\coordinate (T) at (current page.south west);
		\coordinate (P) at (current page.south east);
		%--- dải nền
		\fill[ezNen] (T) rectangle ([yshift=\EZdai]P);
		%--- ô chữ sát mép trên dải: X là góc dưới bên phải, cạnh phải xiên kiểu "/"
		\coordinate (X) at ([xshift=\EZle+\EZrong+\EZdem,yshift=\EZdai-\EZo]T);
		\fill[ezO] ([yshift=\EZdai-\EZo]T) -- (X) --++(\EZxien,\EZo) -- ([yshift=\EZdai]T) --cycle;
		\EZvachxien{X}{1}{1}
		\node[ezchu,anchor=west] at ([xshift=\EZle,yshift=\EZdai-0.5*\EZo]T) {\EZfontO\thuonghieu};
		%--- số trang
		\node[circle,draw=white,line width=0.9pt,text=white,inner sep=0pt,minimum size=0.62cm,anchor=east]
		(so) at ([xshift=-\EZle,yshift=0.5*\EZdai]P) {\EZfontTrang\thepage};
		\node[inner sep=0pt,text=white,anchor=base east] (tr) at ([xshift=-0.15cm]so.base west) {\EZfontTrang Trang};
	\end{tikzpicture}%
}

%====================== CHÂN TRANG - TRANG CHẴN ====================
%  (soi gương: "Trang (số)" bên trái, ô chữ \tacgia bên phải)
\newcommand{\myfancyfootEven}{%
	\EZdorong{\tacgia}%
	\begin{tikzpicture}[remember picture,overlay]
		\coordinate (T) at (current page.south west);
		\coordinate (P) at (current page.south east);
		%--- dải nền
		\fill[ezNen] (T) rectangle ([yshift=\EZdai]P);
		%--- ô chữ sát mép trên dải: X là góc dưới bên trái, cạnh trái xiên kiểu "\"
		\coordinate (X) at ([xshift=-\EZle-\EZrong-\EZdem,yshift=\EZdai-\EZo]P);
		\fill[ezO] ([yshift=\EZdai-\EZo]P) -- (X) --++(-\EZxien,\EZo) -- ([yshift=\EZdai]P) --cycle;
		\EZvachxien{X}{-1}{-1}
		\node[ezchu,anchor=east] at ([xshift=-\EZle,yshift=\EZdai-0.5*\EZo]P) {\EZfontO\tacgia};
		%--- số trang
		%    (vòng tròn ẩn so0 nằm giữa dải để lấy đường chân chữ, giống trang lẻ)
		\node[circle,inner sep=0pt,minimum size=0.62cm] (so0) at ([yshift=0.5*\EZdai]T) {\EZfontTrang\phantom{\thepage}};
		\node[inner sep=0pt,text=white,anchor=base west] (tr) at ($(so0.base -| T)+(\EZle,0)$) {\EZfontTrang Trang};
		\node[circle,draw=white,line width=0.9pt,text=white,inner sep=0pt,minimum size=0.62cm,anchor=base west]
		(so) at ([xshift=0.15cm]tr.base east) {\EZfontTrang\thepage};
	\end{tikzpicture}%
}

%--- Giữ tương thích với các file cũ đang gọi \myfancyhead, \myfancyfoot
\newcommand{\myfancyhead}{\myfancyheadOdd}
\newcommand{\myfancyfoot}{\myfancyfootOdd}

%======================= GẮN VÀO TRANG =============================
\pagestyle{fancy}
\fancyhf{}
\fancyhead[LO]{\myfancyheadOdd}
\fancyhead[RE]{\myfancyheadEven}
\fancyfoot[LO]{\myfancyfootOdd}
\fancyfoot[RE]{\myfancyfootEven}
\renewcommand{\headrulewidth}{0pt}
\renewcommand{\footrulewidth}{0pt}
\fancypagestyle{plain}{%
	\fancyhf{}
	\fancyhead[LO]{\myfancyheadOdd}
	\fancyhead[RE]{\myfancyheadEven}
	\fancyfoot[LO]{\myfancyfootOdd}
	\fancyfoot[RE]{\myfancyfootEven}
	\renewcommand{\headrulewidth}{0pt}
	\renewcommand{\footrulewidth}{0pt}
}
%=================== HẾT PHẦN GỘP TỪ eduzy-de.tex ===================
\newcommand{\indam}[1]{
	{\fontsize{12pt}{9pt}\sffamily \selectfont\bfseries\color{gold1} #1}
}
\newcommand{\ind}[1]{
	{\fontsize{14pt}{9pt}\fontfamily{qag} \selectfont\bfseries\color{gold1} #1}
}
\renewcommand{\FalseEX}{\stepcounter{dapan}\textbf{\selectfont\color{gold1}\Alph{dapan}.}\ignorespaces}
\renewtheorem{ex}{\color{gold1}Câu}
\renewcommand{\baselinestretch}{1.1}
\newcommand{\hoac}[1]{
	\left[\begin{aligned}#1\end{aligned}\right.}
\newcommand{\heva}[1]{
	\left\{\begin{aligned}#1\end{aligned}\right.}
\usepackage{esvect}
\def\vec{\vv}
\def\vt{\vv}
\def\overrightarrow{\vv}
%====================== LOGO EDUZY VẼ BẰNG TIKZ ======================
%  Đường viền dò từ logo gốc; dải màu vàng kim đo trực tiếp trên ảnh gốc
%  (nghiêng 21 độ, tô bằng 80 dải hẹp nên hiển thị giống nhau trên mọi trình xem PDF).
%  Dùng: \eduzylogo          -> cao 2.4cm (mặc định)
%        \eduzylogo[1.5cm]   -> cao 1.5cm
\definecolor{ezText}{RGB}{245,200,46}
\newcommand{\eduzylogo}[1][2.4cm]{%
	\resizebox{!}{#1}{%
		\begin{tikzpicture}[x=0.01cm,y=0.01cm]
			\path[use as bounding box] (0,0) rectangle (224,-224);
			% mũ tốt nghiệp + vương miện
			\begin{scope}
				\clip
				(119.4,-176.1)..controls(118.9,-176.1)and(116.8,-176)..(114.6,-176)
				..controls(101.4,-176)and(85.8,-175.9)..(83.7,-175.7)
				..controls(81.8,-175.6)and(81,-175.6)..(80.2,-175.7)
				..controls(78.8,-175.9)and(77.7,-175.9)..(76.2,-175.7)
				..controls(74.9,-175.5)and(70.8,-175.5)..(64.3,-175.7)
				..controls(59.1,-175.9)and(58.3,-175.7)..(58.6,-173.7)
				..controls(58.6,-172.8)and(59.6,-169.6)..(60.3,-168)
				..controls(60.8,-166.9)and(61.4,-165.8)..(62.3,-165)
				..controls(63.1,-164.2)and(64.3,-163.4)..(65,-163.2)
				..controls(65.5,-163.1)and(72.3,-163)..(76.6,-163.1)
				..controls(77.9,-163.1)and(80.1,-163.1)..(81.4,-163)
				..controls(82.8,-162.9)and(86.3,-162.9)..(89.2,-162.9)
				..controls(93,-162.9)and(94.9,-162.8)..(96,-162.7)
				..controls(97.2,-162.5)and(101,-162.5)..(114,-162.4)
				..controls(123,-162.3)and(134.9,-162.2)..(140.2,-162.1)
				..controls(150.1,-162)and(152.2,-162.1)..(153.6,-162.4)
				..controls(154.4,-162.6)and(157.1,-163.8)..(158.1,-164.4)
				..controls(158.8,-164.8)and(159.7,-165.7)..(159.9,-166.3)
				..controls(160,-166.5)and(160.2,-167.1)..(160.2,-167.6)
				..controls(160.3,-168)and(160.6,-169.4)..(160.9,-170.6)
				..controls(161.2,-171.7)and(161.5,-173.1)..(161.6,-173.6)
				..controls(161.8,-174.4)and(161.8,-174.5)..(161.5,-174.9)
				..controls(161.2,-175.6)and(160.8,-175.9)..(159.9,-176)
				..controls(159.2,-176.1)and(154.1,-176.2)..(151.5,-176.2)
				..controls(148.1,-176.1)and(122.6,-176.1)..(121.6,-176.2)
				..controls(120.8,-176.2)and(119.8,-176.2)..(119.4,-176.1)--cycle
				(130.4,-160)..controls(129,-160)and(119.8,-160)..(110.1,-159.9)
				..controls(100.3,-159.8)and(88.5,-159.7)..(83.9,-159.6)
				..controls(77.8,-159.5)and(75.3,-159.5)..(74.5,-159.6)
				..controls(73.9,-159.7)and(72.5,-159.8)..(71.3,-159.8)
				..controls(69.6,-159.8)and(69,-159.8)..(68.5,-159.6)
				..controls(67.6,-159.3)and(66.8,-158.6)..(66.2,-157.6)
				..controls(66,-157.3)and(65.6,-156.5)..(65.3,-156.1)
				..controls(64.6,-155.1)and(64.5,-154.9)..(63.7,-152.9)
				..controls(63.2,-151.6)and(62.9,-150.9)..(61.8,-149.6)
				..controls(61.5,-149.2)and(61.1,-148.5)..(60.9,-147.9)
				..controls(60.7,-147.4)and(60.3,-146.6)..(60.1,-146.1)
				..controls(59.9,-145.5)and(59.5,-144.9)..(59,-144.3)
				..controls(58.5,-143.7)and(58.2,-143.2)..(57.9,-142.5)
				..controls(57.7,-141.9)and(57.3,-141.1)..(57,-140.7)
				..controls(56.6,-140.3)and(56.1,-139.3)..(55.8,-138.5)
				..controls(55.5,-137.8)and(55.1,-136.8)..(54.9,-136.5)
				..controls(54.7,-136.1)and(54.3,-135.2)..(53.9,-134.5)
				..controls(53.6,-133.8)and(53,-132.7)..(52.6,-132)
				..controls(51.6,-130.5)and(51.4,-130.2)..(51.1,-129)
				..controls(50.6,-127.5)and(50.6,-127.4)..(49.7,-126.3)
				..controls(49.2,-125.7)and(48.6,-124.8)..(48.3,-124.2)
				..controls(48,-123.6)and(47.5,-122.7)..(47.3,-122.2)
				..controls(47,-121.8)and(46.5,-120.7)..(46.1,-119.8)
				..controls(45.8,-118.9)and(45.2,-117.8)..(44.9,-117.2)
				..controls(44.2,-116.1)and(43.3,-114.4)..(42.3,-112.2)
				..controls(41.8,-111.3)and(41.3,-110.4)..(41.1,-110.1)
				..controls(40.5,-109.4)and(40.3,-109)..(40,-107.8)
				..controls(39.8,-107)and(39.6,-106.5)..(39.2,-106)
				..controls(38.9,-105.6)and(38.4,-104.7)..(38,-104)
				..controls(37.7,-103.4)and(37.2,-102.5)..(37,-102.1)
				..controls(36.7,-101.7)and(36.3,-100.8)..(36,-100.1)
				..controls(35.6,-98.9)and(35.3,-98.4)..(34.6,-97.6)
				..controls(34.2,-97.1)and(34.2,-96.3)..(34.6,-96)
				..controls(35,-95.7)and(35.5,-95.9)..(36.2,-96.7)
				..controls(36.5,-97)and(37.2,-97.5)..(37.6,-97.7)..controls(38,-98)and(38.6,-98.4)..(39.1,-98.9)
				..controls(39.6,-99.4)and(40.3,-100)..(40.6,-100.3)
				..controls(41.3,-100.8)and(42.4,-101.8)..(43.7,-102.9)
				..controls(44,-103.2)and(44.6,-103.7)..(45,-103.9)
				..controls(45.4,-104.2)and(45.9,-104.6)..(46.1,-104.8)
				..controls(47.1,-105.9)and(48.8,-107.1)..(50.5,-107.9)
				..controls(51.8,-108.6)and(52.3,-109)..(53.4,-110)
				..controls(54.8,-111.2)and(55.6,-111.6)..(57,-112.1)
				..controls(57.5,-112.3)and(58.5,-112.8)..(59.4,-113.3)
				..controls(60.2,-113.7)and(61.3,-114.3)..(61.9,-114.5)
				..controls(62.4,-114.7)and(63.3,-115.1)..(63.8,-115.3)
				..controls(64.4,-115.6)and(65.4,-115.9)..(66.2,-116.1)
				..controls(67,-116.3)and(68.4,-116.6)..(69.2,-116.9)
				..controls(72.3,-117.9)and(72.8,-118)..(78,-118.1)
				..controls(84,-118.3)and(86.1,-118.1)..(88.1,-117.1)
				..controls(89.4,-116.5)and(93,-115)..(93.7,-114.9)
				..controls(94.8,-114.6)and(95.4,-114.2)..(95.9,-113.4)
				..controls(96.2,-113)and(96.7,-112.5)..(97.1,-112.2)
				..controls(98,-111.6)and(101.3,-108.4)..(102.1,-107.2)
				..controls(102.4,-106.8)and(102.9,-106.2)..(103.2,-105.8)
				..controls(103.5,-105.5)and(104,-104.8)..(104.2,-104.2)
				..controls(104.5,-103.6)and(105,-102.8)..(105.3,-102.4)
				..controls(105.6,-102)and(106.1,-101.2)..(106.5,-100.6)
				..controls(107.1,-99.5)and(108.6,-97.5)..(109.5,-96.5)
				..controls(109.9,-96.1)and(110.1,-96)..(110.4,-96)
				..controls(111.1,-96)and(111.6,-96.5)..(111.8,-97.6)
				..controls(112,-98.6)and(112.3,-99.3)..(112.8,-99.9)
				..controls(113.7,-100.9)and(115.2,-103.1)..(115.7,-104.2)
				..controls(116.1,-105)and(116.5,-105.5)..(116.9,-106)
				..controls(117.3,-106.4)and(117.8,-106.9)..(118,-107.3)
				..controls(118.2,-107.6)and(118.6,-108)..(119,-108.3)
				..controls(119.6,-108.8)and(121.1,-110.3)..(121.8,-111.2)
				..controls(122.1,-111.7)and(122.5,-112)..(123.4,-112.4)
				..controls(124.8,-113.2)and(125.8,-113.8)..(127,-114.6)
				..controls(127.5,-114.9)and(128.1,-115.3)..(128.2,-115.4)
				..controls(128.4,-115.5)and(129.3,-115.8)..(130.2,-116)
				..controls(131.1,-116.2)and(132.5,-116.6)..(133.3,-116.9)
				..controls(135.7,-117.9)and(136.1,-117.9)..(142.5,-117.8)
				..controls(145.8,-117.8)and(148.4,-117.7)..(148.8,-117.6)
				..controls(149.2,-117.6)and(150.2,-117.2)..(151,-116.8)
				..controls(152,-116.4)and(152.8,-116)..(153.6,-115.9)
				..controls(154.5,-115.7)and(155.2,-115.4)..(156.3,-114.9)
				..controls(157.1,-114.5)and(158.2,-114)..(158.7,-113.8)
				..controls(159.1,-113.7)and(159.8,-113.3)..(160.2,-113)
				..controls(160.6,-112.7)and(161.3,-112.4)..(161.7,-112.2)
				..controls(162.7,-111.9)and(163.2,-111.7)..(164.1,-111.1)
				..controls(164.5,-110.8)and(165.3,-110.4)..(165.9,-110)
				..controls(166.5,-109.7)and(167.3,-109.2)..(167.6,-108.9)
				..controls(167.9,-108.6)and(168.6,-108.2)..(169,-108)
				..controls(169.4,-107.8)and(170.2,-107.3)..(170.7,-106.9)
				..controls(171.2,-106.5)and(172.1,-105.8)..(172.7,-105.4)
				..controls(173.2,-105)and(174.1,-104.4)..(174.5,-104)
				..controls(174.9,-103.7)and(175.7,-103.1)..(176.2,-102.6)
				..controls(176.7,-102.2)and(177.4,-101.5)..(177.9,-101.2)
				..controls(178.3,-100.8)and(178.8,-100.4)..(179,-100.2)
				..controls(179.2,-100)and(179.6,-99.6)..(180,-99.3)
				..controls(180.4,-99.1)and(180.9,-98.6)..(181.2,-98.4)
				..controls(181.5,-98.1)and(182.1,-97.7)..(182.5,-97.5)
				..controls(183.2,-97.2)and(183.5,-96.9)..(184.3,-95.9)
				..controls(184.6,-95.6)and(184.6,-95.6)..(184.8,-95.7)
				..controls(185,-95.8)and(185.1,-96)..(185.1,-96.2)
				..controls(185.1,-96.5)and(185,-96.7)..(184.6,-97.2)
				..controls(184.1,-97.8)and(184,-97.9)..(183.9,-98.8)
				..controls(183.7,-100.3)and(183.3,-101.3)..(182.9,-101.8)
				..controls(182.6,-102.1)and(182.3,-102.6)..(182.2,-102.9)
				..controls(182,-103.3)and(181.6,-104)..(181.2,-104.6)
				..controls(180.9,-105.2)and(180.5,-106.1)..(180.4,-106.6)
				..controls(179.9,-108)and(179.7,-108.5)..(179,-109.4)
				..controls(178.6,-109.8)and(178.2,-110.5)..(177.9,-111)
				..controls(177.7,-111.5)and(177.3,-112.3)..(177,-112.7)
				..controls(176.6,-113.2)and(176.4,-113.7)..(176.3,-114.2)
				..controls(176.2,-114.7)and(175.9,-115.2)..(175.4,-116)
				..controls(175,-116.6)and(174.4,-117.7)..(174,-118.5)
				..controls(173.7,-119.2)and(173.3,-120.1)..(173,-120.5)
				..controls(172.8,-120.9)and(172.4,-121.8)..(172.1,-122.5)
				..controls(171.8,-123.2)and(171.3,-124.1)..(171,-124.6)
				..controls(170.8,-125.1)and(170.3,-126)..(169.9,-126.6)
				..controls(169.6,-127.3)and(169.2,-128.1)..(169,-128.3)
				..controls(168.8,-128.6)and(168.5,-129.3)..(168.3,-129.8)
				..controls(168.1,-130.3)and(167.8,-131.1)..(167.5,-131.6)
				..controls(166.9,-132.5)and(164.2,-137.9)..(164.2,-138.1)
				..controls(164.2,-138.1)and(163.9,-138.8)..(163.4,-139.7)
				..controls(162.9,-140.5)and(162.4,-141.5)..(162.3,-141.9)
				..controls(161.9,-143)and(161.6,-143.4)..(160.9,-144.1)
				..controls(160.4,-144.6)and(160.3,-144.8)..(160.1,-145.4)
				..controls(160,-145.8)and(159.6,-146.8)..(159.2,-147.8)
				..controls(158.7,-148.8)and(158.1,-150)..(157.9,-150.6)
				..controls(157.6,-151.1)and(157.2,-151.9)..(156.9,-152.4)
				..controls(156.5,-152.8)and(156.1,-153.7)..(155.8,-154.4)
				..controls(155.6,-155)and(155.2,-155.9)..(154.9,-156.3)
				..controls(154.7,-156.7)and(154.4,-157.4)..(154.2,-157.9)
				..controls(154.1,-158.4)and(153.8,-159)..(153.5,-159.3)
				..controls(152.8,-160.1)and(153,-160.1)..(142.4,-160.1)
				..controls(137.2,-160.1)and(131.8,-160.1)..(130.4,-160)--cycle
				(141.5,-103.7)..controls(140.8,-103)and(139.8,-102.1)..(139.2,-101.6)
				..controls(138.5,-101.2)and(137.9,-100.6)..(137.8,-100.4)
				..controls(137,-99.5)and(134,-96.5)..(133.2,-96)
				..controls(132.7,-95.7)and(131.8,-94.9)..(131.1,-94.3)
				..controls(129.4,-92.7)and(127.9,-91.6)..(126.6,-91.1)
				..controls(125.2,-90.5)and(124.1,-90)..(122.9,-89.1)
				..controls(122.1,-88.6)and(121.3,-88.2)..(119.9,-87.7)
				..controls(118.8,-87.3)and(117.5,-86.8)..(116.9,-86.5)
				..controls(114.3,-85.5)and(113.7,-85.4)..(109.7,-85.5)
				..controls(105.3,-85.7)and(103.9,-85.9)..(101.6,-87)
				..controls(100.7,-87.4)and(99.4,-87.8)..(97.9,-88)
				..controls(97,-88.2)and(96.6,-88.3)..(95.8,-88.9)
				..controls(95.4,-89.2)and(94.7,-89.6)..(94.2,-89.8)
				..controls(92.9,-90.3)and(92.3,-90.7)..(90.9,-92.1)
				..controls(90.2,-92.7)and(89.4,-93.3)..(89,-93.5)..controls(87,-94.5)and(85.2,-96)..(83.6,-97.9)
				..controls(83.3,-98.2)and(82.5,-98.9)..(81.8,-99.4)
				..controls(80.7,-100.1)and(80.5,-100.3)..(80,-101.1)
				..controls(78.7,-103.1)and(77.2,-103.9)..(76.4,-103)
				..controls(75.7,-102.2)and(75.7,-103.2)..(75.7,-83.6)
				..controls(75.6,-73.6)and(75.6,-64.8)..(75.5,-64)
				..controls(75.5,-63.2)and(75.5,-61.4)..(75.6,-59.8)--(75.7,-57.1)--(76,-56.8)
				..controls(76.8,-56.2)and(78,-56.3)..(79.8,-57.2)
				..controls(80.4,-57.5)and(81.3,-57.8)..(81.8,-57.9)
				..controls(82.3,-58.1)and(83,-58.4)..(83.6,-58.8)..controls(84.3,-59.2)and(85,-59.5)..(86,-59.8)
				..controls(86.8,-60.1)and(88,-60.6)..(88.7,-60.9)
				..controls(89.3,-61.2)and(90.6,-61.7)..(91.4,-62)
				..controls(92.2,-62.3)and(93.7,-62.9)..(94.6,-63.3)
				..controls(95.6,-63.7)and(96.8,-64.1)..(97.4,-64.3)
				..controls(98,-64.5)and(99,-64.9)..(99.7,-65.2)
				..controls(100.4,-65.6)and(101.5,-66)..(102.3,-66.2)
				..controls(104.2,-66.8)and(105,-67.1)..(106.2,-67.8)
				..controls(107.9,-68.7)and(110.2,-69.4)..(111.3,-69.2)
				..controls(112.5,-69)and(115,-68.2)..(115.6,-67.8)
				..controls(117.3,-66.7)and(117.6,-66.6)..(119.5,-66.2)
				..controls(121.9,-65.8)and(122.1,-65.7)..(123,-65.1)
				..controls(124.1,-64.4)and(124.4,-64.3)..(126,-63.9)
				..controls(127,-63.7)and(127.7,-63.4)..(128.4,-63.1)
				..controls(129,-62.8)and(129.7,-62.5)..(130.6,-62.4)
				..controls(132.1,-62.1)and(134,-61.5)..(135.2,-61)
				..controls(135.7,-60.8)and(136.7,-60.5)..(137.3,-60.3)
				..controls(138.4,-60)and(139.3,-59.7)..(141.1,-59)
				..controls(141.6,-58.7)and(142,-58.6)..(142.5,-58.6)
				..controls(143.2,-58.6)and(143.3,-58.7)..(143.6,-59)
				..controls(144,-59.5)and(144.1,-59.9)..(144,-60.8)
				..controls(144,-61.2)and(144,-71.1)..(144,-82.8)--(143.9,-104.1)--(143.7,-104.5)
				..controls(143.5,-104.7)and(143.3,-104.8)..(143,-104.9)
				..controls(142.6,-104.9)and(142.5,-104.8)..(141.5,-103.7)--cycle
				(27.8,-91.4)..controls(27.1,-90.9)and(26.3,-90)..(26.1,-89.5)
				..controls(26,-89.2)and(25.8,-88.6)..(25.7,-87.9)
				..controls(25.4,-86.2)and(25.6,-84.8)..(26.5,-83.5)
				..controls(27.4,-82.3)and(28.3,-82)..(29.4,-82.3)
				..controls(31.1,-82.8)and(32.1,-84.9)..(32,-87.4)
				..controls(31.9,-89.3)and(31.2,-90.6)..(29.9,-91.4)
				..controls(29.1,-91.9)and(28.6,-91.9)..(27.8,-91.4)--cycle
				(190.1,-91.3)..controls(189.7,-91.2)and(189.3,-90.9)..(189.1,-90.7)
				..controls(188.3,-89.8)and(187.6,-87.7)..(187.6,-86.2)
				..controls(187.6,-85.5)and(187.7,-85.2)..(188.2,-84.4)
				..controls(189.1,-82.6)and(190.1,-81.8)..(191.3,-82.1)
				..controls(192,-82.3)and(193.5,-83.7)..(193.8,-84.4)
				..controls(194.1,-85.1)and(194.2,-87.1)..(194,-88.3)
				..controls(193.8,-89.8)and(193.3,-90.5)..(192.3,-91.3)
				..controls(191.7,-91.7)and(190.9,-91.7)..(190.1,-91.3)--cycle
				(173.7,-83.4)..controls(172.2,-82.5)and(172,-82.5)..(170.6,-82.4)
				..controls(169.7,-82.4)and(168.8,-82.4)..(168.3,-82.5)
				..controls(167.2,-82.7)and(166.9,-82.7)..(166.6,-82.2)
				..controls(166.3,-81.8)and(166.3,-81)..(166.8,-80)
				..controls(167.2,-78.9)and(167.4,-77.4)..(167.7,-74.1)
				..controls(167.8,-72.8)and(167.9,-71.8)..(167.9,-71.7)
				..controls(168,-71.7)and(168,-70.2)..(168,-68.4)..controls(168,-65.4)and(168,-64.9)..(167.7,-63)
				..controls(167.3,-60.5)and(167.1,-59.9)..(166.4,-58.4)
				..controls(166.1,-57.7)and(165.8,-56.9)..(165.7,-56.6)
				..controls(165.7,-56.2)and(165.4,-55.6)..(165.2,-55.2)
				..controls(164.6,-53.8)and(161.9,-50.7)..(161,-50.3)
				..controls(159.8,-49.7)and(157.9,-50)..(155.2,-51.2)
				..controls(154.7,-51.4)and(153.4,-51.8)..(152.2,-52.1)
				..controls(150.8,-52.4)and(149.4,-52.9)..(148.6,-53.2)
				..controls(147.8,-53.6)and(146.7,-54)..(146,-54.2)
				..controls(145.4,-54.4)and(144.5,-54.7)..(144,-54.9)
				..controls(143.6,-55.2)and(142.7,-55.5)..(141.8,-55.7)
				..controls(141,-55.9)and(139.6,-56.3)..(138.7,-56.6)
				..controls(137.8,-57)and(136.3,-57.4)..(135.3,-57.6)
				..controls(133.8,-57.9)and(133.2,-58.1)..(132.1,-58.7)
				..controls(131.1,-59.2)and(130.3,-59.4)..(129.4,-59.6)
				..controls(128.7,-59.8)and(127.8,-60)..(127.4,-60.1)
				..controls(127,-60.3)and(126,-60.5)..(125.1,-60.7)
				..controls(123.2,-61.1)and(121.7,-61.5)..(120.2,-62.1)
				..controls(119.7,-62.3)and(118.1,-62.8)..(116.7,-63.2)
				..controls(114.6,-63.9)and(114,-64.1)..(113,-64.7)
				..controls(111.9,-65.4)and(111.8,-65.4)..(111,-65.5)
				..controls(109.9,-65.5)and(109.3,-65.4)..(107.7,-64.7)
				..controls(107,-64.4)and(106,-64)..(105.3,-63.9)
				..controls(103.5,-63.5)and(102.2,-63)..(101.2,-62.5)
				..controls(100.1,-61.8)and(99.3,-61.5)..(97.8,-61.1)
				..controls(97.2,-61)and(96.2,-60.7)..(95.7,-60.4)
				..controls(95.1,-60.1)and(94.2,-59.8)..(93.5,-59.6)
				..controls(92.9,-59.4)and(91.7,-59)..(91,-58.7)..controls(90.2,-58.4)and(88.9,-58)..(88,-57.7)
				..controls(87.2,-57.5)and(86.1,-57.1)..(85.5,-56.8)
				..controls(85,-56.6)and(84.1,-56.2)..(83.5,-56)..controls(82.9,-55.7)and(81.8,-55.2)..(81,-54.9)
				..controls(80.1,-54.5)and(79.1,-54.2)..(78.7,-54.1)
				..controls(77.9,-53.8)and(77.3,-53.6)..(76.2,-52.9)
				..controls(75.7,-52.6)and(75.2,-52.4)..(74.1,-52.2)
				..controls(72.5,-51.9)and(72.4,-51.8)..(71.3,-51.1)
				..controls(69.9,-50.1)and(69.4,-49.9)..(67.9,-49.8)
				..controls(67,-49.7)and(66.6,-49.6)..(66.3,-49.4)
				..controls(66.1,-49.3)and(65.7,-49)..(65.4,-48.8)
				..controls(65.1,-48.7)and(64.2,-48.4)..(63.6,-48.2)
				..controls(62.9,-48)and(61.7,-47.6)..(60.9,-47.2)
				..controls(60.1,-46.8)and(58.9,-46.4)..(58.2,-46.2)
				..controls(57.6,-46)and(56.8,-45.7)..(56.5,-45.5)
				..controls(54.9,-44.6)and(54.6,-44.5)..(53.2,-44.2)
				..controls(51.7,-43.8)and(51,-43.6)..(49.9,-42.9)
				..controls(49.5,-42.6)and(48.9,-42.4)..(47.5,-42)
				..controls(45.7,-41.6)and(45.6,-41.5)..(44.3,-40.6)
				..controls(43.5,-40.1)and(42.8,-39.6)..(42.7,-39.4)
				..controls(42.4,-38.7)and(43,-38.1)..(44.5,-37.6)..controls(44.9,-37.5)and(46,-37)..(46.9,-36.6)
				..controls(47.9,-36.1)and(48.8,-35.8)..(49.6,-35.6)
				..controls(50.5,-35.4)and(51.3,-35.2)..(52.4,-34.6)
				..controls(53.3,-34.1)and(54.5,-33.7)..(54.9,-33.5)
				..controls(55.4,-33.4)and(56.4,-33)..(57.2,-32.6)
				..controls(58.6,-32)and(60.1,-31.5)..(61.6,-31.2)
				..controls(62.3,-31.1)and(62.7,-31)..(63.2,-30.7)
				..controls(64.1,-30)and(64.3,-29.9)..(65.3,-29.7)
				..controls(65.7,-29.6)and(66.5,-29.2)..(67.1,-28.9)
				..controls(67.6,-28.6)and(68.8,-28.2)..(69.6,-27.9)
				..controls(70.5,-27.7)and(71.6,-27.3)..(72.1,-27)
				..controls(72.6,-26.7)and(73.8,-26.2)..(74.8,-25.9)
				..controls(75.7,-25.5)and(76.7,-25.1)..(77.1,-25)
				..controls(77.4,-24.8)and(78.5,-24.4)..(79.6,-24.1)
				..controls(80.6,-23.8)and(82,-23.2)..(82.8,-22.9)
				..controls(83.5,-22.5)and(84.8,-22)..(85.8,-21.8)
				..controls(86.7,-21.5)and(87.8,-21.1)..(88.1,-20.9)
				..controls(88.4,-20.7)and(89.3,-20.3)..(90.1,-20.1)
				..controls(90.8,-19.8)and(91.9,-19.4)..(92.4,-19.1)
				..controls(93.5,-18.5)and(94.4,-18.2)..(95.4,-18.1)
				..controls(95.7,-18)and(96.4,-17.8)..(96.9,-17.5)
				..controls(97.5,-17.3)and(98.4,-16.9)..(99,-16.6)
				..controls(99.7,-16.3)and(100.9,-15.9)..(101.9,-15.6)
				..controls(102.9,-15.3)and(104.4,-14.7)..(105.3,-14.3)
				..controls(108.1,-13.2)and(108.5,-13.1)..(110.1,-13.1)
				..controls(111.4,-13.1)and(111.5,-13.1)..(112.4,-13.5)
				..controls(113.2,-14)and(115,-14.8)..(116.6,-15.4)
				..controls(116.9,-15.5)and(117.9,-15.8)..(118.8,-16.1)
				..controls(120.9,-16.6)and(121.5,-16.9)..(122.6,-17.9)
				..controls(123,-18.3)and(123.7,-18.8)..(123.9,-18.9)
				..controls(124.6,-19.4)and(126,-19.8)..(127.9,-20.2)
				..controls(129.6,-20.5)and(129.8,-20.6)..(131.1,-21.4)
				..controls(131.5,-21.7)and(132.4,-22.1)..(133.1,-22.3)
				..controls(136.1,-23.2)and(138.3,-24.1)..(140.2,-25)
				..controls(141,-25.5)and(141.5,-25.6)..(142.5,-25.8)
				..controls(145.1,-26.3)and(146,-26.6)..(149.4,-28.1)
				..controls(151.3,-28.9)and(153.5,-29.8)..(154.3,-30.1)
				..controls(155.1,-30.4)and(156.2,-30.8)..(156.6,-31)
				..controls(157.1,-31.2)and(158.2,-31.6)..(159.1,-31.8)
				..controls(160.6,-32.2)and(160.9,-32.3)..(161.7,-32.9)
				..controls(162.4,-33.3)and(162.9,-33.6)..(163.7,-33.8)
				..controls(164.3,-34)and(164.8,-34.1)..(164.9,-34.1)
				..controls(165,-34.1)and(165.6,-34.4)..(166.4,-34.8)
				..controls(167.3,-35.2)and(168.1,-35.4)..(169,-35.6)
				..controls(170,-35.8)and(170.5,-36)..(171.7,-36.6)
				..controls(172.4,-37)and(173.3,-37.4)..(173.6,-37.4)
				..controls(175.6,-37.7)and(176.2,-38)..(176.7,-39.1)
				..controls(176.9,-39.6)and(176.9,-39.7)..(176.7,-40.4)
				..controls(176.4,-41.6)and(175.8,-42.7)..(175.2,-43.4)
				..controls(174.5,-44.1)and(172.6,-45.5)..(171.7,-45.8)
				..controls(170.6,-46.2)and(169.1,-46)..(168.2,-45.4)
				..controls(167.9,-45.2)and(167.3,-44.8)..(166.9,-44.4)
				..controls(166.4,-44.1)and(165.8,-43.6)..(165.4,-43.2)
				..controls(165.1,-42.9)and(164.4,-42.3)..(163.9,-41.9)
				..controls(163.3,-41.5)and(162.7,-40.9)..(162.4,-40.5)
				..controls(161.4,-39.2)and(159.9,-38)..(158.8,-37.8)
				..controls(158.4,-37.8)and(158.3,-37.7)..(158.2,-37.2)
				..controls(158,-36.8)and(157.4,-36.1)..(157.1,-36.1)
				..controls(157,-36.1)and(156.7,-36.1)..(156.5,-36)
				..controls(156.1,-35.9)and(155.9,-35.7)..(155.7,-35.2)
				..controls(155.3,-34.6)and(155.1,-34.4)..(154,-34.1)
				..controls(153.6,-34)and(153.2,-33.7)..(152.9,-33.5)
				..controls(152.4,-32.9)and(151.6,-32.4)..(151.3,-32.4)
				..controls(150.9,-32.4)and(150.4,-32.6)..(150.4,-32.9)
				..controls(150.4,-33.2)and(150.8,-33.4)..(151.4,-33.6)
				..controls(152.1,-33.7)and(152.8,-34.1)..(153.5,-34.9)
				..controls(154.1,-35.5)and(154.3,-35.6)..(155,-35.8)
				..controls(155.7,-36)and(155.7,-36)..(156,-36.6)
				..controls(156.3,-37.3)and(156.7,-37.7)..(157.5,-37.8)
				..controls(158.1,-37.9)and(158.2,-38)..(158.4,-38.8)
				..controls(158.5,-39.4)and(158.6,-39.6)..(159.6,-40.4)
				..controls(160.2,-41)and(160.9,-41.6)..(161.3,-41.8)
				..controls(161.7,-42)and(161.9,-42.3)..(162.2,-42.8)
				..controls(162.4,-43.2)and(162.8,-43.8)..(163.1,-44.1)
				..controls(163.9,-44.8)and(165.1,-46.7)..(165.7,-48)
				..controls(166.2,-49.3)and(166.4,-49.7)..(167.2,-50.5)
				..controls(167.5,-50.8)and(167.9,-51.4)..(168.1,-51.8)
				..controls(168.3,-52.2)and(168.8,-53)..(169.1,-53.6)
				..controls(169.6,-54.3)and(169.8,-54.9)..(169.9,-55.6)
				..controls(170.4,-57.8)and(170.5,-58.1)..(171.2,-59.5)
				..controls(171.6,-60.3)and(172,-61.5)..(172.2,-62.1)
				..controls(172.3,-62.8)and(172.7,-63.7)..(172.9,-64.1)
				..controls(173.2,-64.6)and(173.4,-65.3)..(173.4,-65.6)
				..controls(173.8,-68.8)and(174.2,-70.6)..(174.9,-72.6)
				..controls(175.5,-74.5)and(175.9,-77.1)..(176.1,-80.2)
				..controls(176.3,-83.4)and(175.5,-84.4)..(173.7,-83.4)--cycle;
				\definecolor{ezs}{RGB}{247,179,26}\fill[ezs](-5.7,-156.6)--(-3.4,-157.5)--(61.9,12.7)--(59.7,13.6)--cycle;
				\definecolor{ezs}{RGB}{247,179,22}\fill[ezs](-3.8,-157.3)--(-1.5,-158.2)--(63.9,12)--(61.6,12.9)--cycle;
				\definecolor{ezs}{RGB}{248,179,19}\fill[ezs](-1.8,-158.1)--(0.4,-158.9)--(65.8,11.2)--(63.5,12.1)--cycle;
				\definecolor{ezs}{RGB}{248,178,16}\fill[ezs](0.1,-158.8)--(2.4,-159.7)--(67.7,10.5)--(65.4,11.4)--cycle;
				\definecolor{ezs}{RGB}{249,178,15}\fill[ezs](2,-159.5)--(4.3,-160.4)--(69.6,9.8)--(67.3,10.7)--cycle;
				\definecolor{ezs}{RGB}{251,177,17}\fill[ezs](3.9,-160.3)--(6.2,-161.1)--(71.5,9)--(69.2,9.9)--cycle;
				\definecolor{ezs}{RGB}{252,179,18}\fill[ezs](5.8,-161)--(8.1,-161.9)--(73.4,8.3)--(71.1,9.2)--cycle;
				\definecolor{ezs}{RGB}{252,182,17}\fill[ezs](7.7,-161.7)--(10,-162.6)--(75.3,7.6)--(73,8.5)--cycle;
				\definecolor{ezs}{RGB}{252,185,17}\fill[ezs](9.6,-162.5)--(11.9,-163.4)--(77.2,6.8)--(75,7.7)--cycle;
				\definecolor{ezs}{RGB}{251,189,18}\fill[ezs](11.5,-163.2)--(13.8,-164.1)--(79.2,6.1)--(76.9,7)--cycle;
				\definecolor{ezs}{RGB}{251,191,22}\fill[ezs](13.4,-163.9)--(15.7,-164.8)--(81.1,5.4)--(78.8,6.2)--cycle;
				\definecolor{ezs}{RGB}{252,193,25}\fill[ezs](15.4,-164.7)--(17.6,-165.6)--(83,4.6)--(80.7,5.5)--cycle;
				\definecolor{ezs}{RGB}{253,195,27}\fill[ezs](17.3,-165.4)--(19.6,-166.3)--(84.9,3.9)--(82.6,4.8)--cycle;
				\definecolor{ezs}{RGB}{253,196,31}\fill[ezs](19.2,-166.1)--(21.5,-167)--(86.8,3.2)--(84.5,4)--cycle;
				\definecolor{ezs}{RGB}{254,199,35}\fill[ezs](21.1,-166.9)--(23.4,-167.8)--(88.7,2.4)--(86.4,3.3)--cycle;
				\definecolor{ezs}{RGB}{253,201,38}\fill[ezs](23,-167.6)--(25.3,-168.5)--(90.6,1.7)--(88.3,2.6)--cycle;
				\definecolor{ezs}{RGB}{253,203,40}\fill[ezs](24.9,-168.3)--(27.2,-169.2)--(92.5,1)--(90.3,1.8)--cycle;
				\definecolor{ezs}{RGB}{253,205,42}\fill[ezs](26.8,-169.1)--(29.1,-170)--(94.4,0.2)--(92.2,1.1)--cycle;
				\definecolor{ezs}{RGB}{253,207,45}\fill[ezs](28.7,-169.8)--(31,-170.7)--(96.4,-0.5)--(94.1,0.4)--cycle;
				\definecolor{ezs}{RGB}{253,210,49}\fill[ezs](30.7,-170.5)--(32.9,-171.4)--(98.3,-1.2)--(96,-0.4)--cycle;
				\definecolor{ezs}{RGB}{253,212,54}\fill[ezs](32.6,-171.3)--(34.9,-172.2)--(100.2,-2)--(97.9,-1.1)--cycle;
				\definecolor{ezs}{RGB}{253,214,59}\fill[ezs](34.5,-172)--(36.8,-172.9)--(102.1,-2.7)--(99.8,-1.8)--cycle;
				\definecolor{ezs}{RGB}{253,216,63}\fill[ezs](36.4,-172.7)--(38.7,-173.6)--(104,-3.4)--(101.7,-2.6)--cycle;
				\definecolor{ezs}{RGB}{253,218,65}\fill[ezs](38.3,-173.5)--(40.6,-174.4)--(105.9,-4.2)--(103.6,-3.3)--cycle;
				\definecolor{ezs}{RGB}{254,220,66}\fill[ezs](40.2,-174.2)--(42.5,-175.1)--(107.8,-4.9)--(105.5,-4)--cycle;
				\definecolor{ezs}{RGB}{254,222,68}\fill[ezs](42.1,-175)--(44.4,-175.8)--(109.7,-5.6)--(107.5,-4.8)--cycle;
				\definecolor{ezs}{RGB}{254,224,71}\fill[ezs](44,-175.7)--(46.3,-176.6)--(111.7,-6.4)--(109.4,-5.5)--cycle;
				\definecolor{ezs}{RGB}{254,226,75}\fill[ezs](46,-176.4)--(48.2,-177.3)--(113.6,-7.1)--(111.3,-6.2)--cycle;
				\definecolor{ezs}{RGB}{254,228,79}\fill[ezs](47.9,-177.2)--(50.1,-178)--(115.5,-7.8)--(113.2,-7)--cycle;
				\definecolor{ezs}{RGB}{254,230,82}\fill[ezs](49.8,-177.9)--(52.1,-178.8)--(117.4,-8.6)--(115.1,-7.7)--cycle;
				\definecolor{ezs}{RGB}{254,232,85}\fill[ezs](51.7,-178.6)--(54,-179.5)--(119.3,-9.3)--(117,-8.4)--cycle;
				\definecolor{ezs}{RGB}{254,234,88}\fill[ezs](53.6,-179.4)--(55.9,-180.2)--(121.2,-10)--(118.9,-9.2)--cycle;
				\definecolor{ezs}{RGB}{254,237,91}\fill[ezs](55.5,-180.1)--(57.8,-181)--(123.1,-10.8)--(120.8,-9.9)--cycle;
				\definecolor{ezs}{RGB}{254,239,94}\fill[ezs](57.4,-180.8)--(59.7,-181.7)--(125,-11.5)--(122.8,-10.6)--cycle;
				\definecolor{ezs}{RGB}{254,241,97}\fill[ezs](59.3,-181.6)--(61.6,-182.4)--(127,-12.2)--(124.7,-11.4)--cycle;
				\definecolor{ezs}{RGB}{254,243,101}\fill[ezs](61.2,-182.3)--(63.5,-183.2)--(128.9,-13)--(126.6,-12.1)--cycle;
				\definecolor{ezs}{RGB}{254,245,103}\fill[ezs](63.2,-183)--(65.4,-183.9)--(130.8,-13.7)--(128.5,-12.8)--cycle;
				\definecolor{ezs}{RGB}{254,245,105}\fill[ezs](65.1,-183.8)--(67.4,-184.6)--(132.7,-14.4)--(130.4,-13.6)--cycle;
				\definecolor{ezs}{RGB}{254,245,105}\fill[ezs](67,-184.5)--(69.3,-185.4)--(134.6,-15.2)--(132.3,-14.3)--cycle;
				\definecolor{ezs}{RGB}{254,244,103}\fill[ezs](68.9,-185.2)--(71.2,-186.1)--(136.5,-15.9)--(134.2,-15)--cycle;
				\definecolor{ezs}{RGB}{253,243,101}\fill[ezs](70.8,-186)--(73.1,-186.8)--(138.4,-16.6)--(136.1,-15.8)--cycle;
				\definecolor{ezs}{RGB}{253,241,98}\fill[ezs](72.7,-186.7)--(75,-187.6)--(140.3,-17.4)--(138,-16.5)--cycle;
				\definecolor{ezs}{RGB}{253,239,95}\fill[ezs](74.6,-187.4)--(76.9,-188.3)--(142.2,-18.1)--(140,-17.2)--cycle;
				\definecolor{ezs}{RGB}{253,236,92}\fill[ezs](76.5,-188.2)--(78.8,-189)--(144.2,-18.8)--(141.9,-18)--cycle;
				\definecolor{ezs}{RGB}{253,234,88}\fill[ezs](78.5,-188.9)--(80.7,-189.8)--(146.1,-19.6)--(143.8,-18.7)--cycle;
				\definecolor{ezs}{RGB}{253,232,85}\fill[ezs](80.4,-189.6)--(82.7,-190.5)--(148,-20.3)--(145.7,-19.4)--cycle;
				\definecolor{ezs}{RGB}{253,230,82}\fill[ezs](82.3,-190.4)--(84.6,-191.2)--(149.9,-21)--(147.6,-20.2)--cycle;
				\definecolor{ezs}{RGB}{253,228,78}\fill[ezs](84.2,-191.1)--(86.5,-192)--(151.8,-21.8)--(149.5,-20.9)--cycle;
				\definecolor{ezs}{RGB}{253,226,74}\fill[ezs](86.1,-191.8)--(88.4,-192.7)--(153.7,-22.5)--(151.4,-21.6)--cycle;
				\definecolor{ezs}{RGB}{253,224,70}\fill[ezs](88,-192.6)--(90.3,-193.4)--(155.6,-23.3)--(153.3,-22.4)--cycle;
				\definecolor{ezs}{RGB}{253,222,67}\fill[ezs](89.9,-193.3)--(92.2,-194.2)--(157.5,-24)--(155.3,-23.1)--cycle;
				\definecolor{ezs}{RGB}{253,220,65}\fill[ezs](91.8,-194)--(94.1,-194.9)--(159.5,-24.7)--(157.2,-23.8)--cycle;
				\definecolor{ezs}{RGB}{253,217,63}\fill[ezs](93.7,-194.8)--(96,-195.6)--(161.4,-25.5)--(159.1,-24.6)--cycle;
				\definecolor{ezs}{RGB}{253,215,60}\fill[ezs](95.7,-195.5)--(97.9,-196.4)--(163.3,-26.2)--(161,-25.3)--cycle;
				\definecolor{ezs}{RGB}{253,213,56}\fill[ezs](97.6,-196.2)--(99.9,-197.1)--(165.2,-26.9)--(162.9,-26)--cycle;
				\definecolor{ezs}{RGB}{252,211,53}\fill[ezs](99.5,-197)--(101.8,-197.8)--(167.1,-27.7)--(164.8,-26.8)--cycle;
				\definecolor{ezs}{RGB}{252,209,49}\fill[ezs](101.4,-197.7)--(103.7,-198.6)--(169,-28.4)--(166.7,-27.5)--cycle;
				\definecolor{ezs}{RGB}{252,207,45}\fill[ezs](103.3,-198.4)--(105.6,-199.3)--(170.9,-29.1)--(168.6,-28.2)--cycle;
				\definecolor{ezs}{RGB}{253,205,42}\fill[ezs](105.2,-199.2)--(107.5,-200)--(172.8,-29.9)--(170.6,-29)--cycle;
				\definecolor{ezs}{RGB}{253,202,40}\fill[ezs](107.1,-199.9)--(109.4,-200.8)--(174.7,-30.6)--(172.5,-29.7)--cycle;
				\definecolor{ezs}{RGB}{253,200,38}\fill[ezs](109,-200.6)--(111.3,-201.5)--(176.7,-31.3)--(174.4,-30.4)--cycle;
				\definecolor{ezs}{RGB}{253,198,35}\fill[ezs](111,-201.4)--(113.2,-202.2)--(178.6,-32.1)--(176.3,-31.2)--cycle;
				\definecolor{ezs}{RGB}{253,196,31}\fill[ezs](112.9,-202.1)--(115.2,-203)--(180.5,-32.8)--(178.2,-31.9)--cycle;
				\definecolor{ezs}{RGB}{253,194,28}\fill[ezs](114.8,-202.8)--(117.1,-203.7)--(182.4,-33.5)--(180.1,-32.6)--cycle;
				\definecolor{ezs}{RGB}{253,192,26}\fill[ezs](116.7,-203.6)--(119,-204.5)--(184.3,-34.3)--(182,-33.4)--cycle;
				\definecolor{ezs}{RGB}{253,189,24}\fill[ezs](118.6,-204.3)--(120.9,-205.2)--(186.2,-35)--(183.9,-34.1)--cycle;
				\definecolor{ezs}{RGB}{253,186,20}\fill[ezs](120.5,-205)--(122.8,-205.9)--(188.1,-35.7)--(185.8,-34.8)--cycle;
				\definecolor{ezs}{RGB}{253,184,17}\fill[ezs](122.4,-205.8)--(124.7,-206.7)--(190,-36.5)--(187.8,-35.6)--cycle;
				\definecolor{ezs}{RGB}{252,183,15}\fill[ezs](124.3,-206.5)--(126.6,-207.4)--(192,-37.2)--(189.7,-36.3)--cycle;
				\definecolor{ezs}{RGB}{252,181,14}\fill[ezs](126.3,-207.2)--(128.5,-208.1)--(193.9,-37.9)--(191.6,-37.1)--cycle;
				\definecolor{ezs}{RGB}{252,179,13}\fill[ezs](128.2,-208)--(130.4,-208.9)--(195.8,-38.7)--(193.5,-37.8)--cycle;
				\definecolor{ezs}{RGB}{252,178,12}\fill[ezs](130.1,-208.7)--(132.4,-209.6)--(197.7,-39.4)--(195.4,-38.5)--cycle;
				\definecolor{ezs}{RGB}{251,178,10}\fill[ezs](132,-209.4)--(134.3,-210.3)--(199.6,-40.1)--(197.3,-39.3)--cycle;
				\definecolor{ezs}{RGB}{251,178,10}\fill[ezs](133.9,-210.2)--(136.2,-211.1)--(201.5,-40.9)--(199.2,-40)--cycle;
				\definecolor{ezs}{RGB}{250,177,11}\fill[ezs](135.8,-210.9)--(138.1,-211.8)--(203.4,-41.6)--(201.1,-40.7)--cycle;
				\definecolor{ezs}{RGB}{250,178,11}\fill[ezs](137.7,-211.6)--(140,-212.5)--(205.3,-42.3)--(203.1,-41.5)--cycle;
				\definecolor{ezs}{RGB}{248,180,10}\fill[ezs](139.6,-212.4)--(141.9,-213.3)--(207.3,-43.1)--(205,-42.2)--cycle;
				\definecolor{ezs}{RGB}{247,181,9}\fill[ezs](141.5,-213.1)--(143.8,-214)--(209.2,-43.8)--(206.9,-42.9)--cycle;
				\definecolor{ezs}{RGB}{246,182,9}\fill[ezs](143.5,-213.8)--(145.7,-214.7)--(211.1,-44.5)--(208.8,-43.7)--cycle;
				\definecolor{ezs}{RGB}{246,182,10}\fill[ezs](145.4,-214.6)--(147.7,-215.5)--(213,-45.3)--(210.7,-44.4)--cycle;
			\end{scope}
			% chữ EDUZY (Times đậm, làm dày nét bằng 5 lớp lệch nhẹ; căn đúng khung chữ logo gốc)
			\foreach \dx/\dy in {-0.7/0,0.7/0,0/-0.3,0/0.3,0/0}
			\node[anchor=center,inner sep=0pt,outer sep=0pt,text=ezText,align=center] at (109.6+\dx,-193.6+\dy)
			{\resizebox{1.01cm}{0.22cm}{\fontencoding{T1}\fontfamily{ptm}\fontseries{b}\selectfont EDUZY}};
		\end{tikzpicture}%
	}%
}
%====================== THIẾT LẬP THÂN ĐỀ ======================
%--- Môi trường tự luận {bt}: ex_test 3.3.4 chưa khai báo sẵn nên tạo ở đây
%    (ex_test đã tự móc \fixshowans và \listenumerate cho bt)
\makeatletter
\@ifundefined{bt}{\newtheorem{bt}{\color{gold1}\namebt}}{}
\makeatother
\renewcommand{\sectionmark}[1]{\markright{\thesection\ #1}}
%--- Khung tiêu đề (logo + tối đa 4 dòng; dòng 3, 4 để trống {} thì bỏ qua)
%  Đề thi:    \khungtieude{THI THỬ THPT QUỐC GIA NĂM 2026}{Môn Toán -- Đề số 21}
%                         {(Đề thi gồm: ...)}{Thời gian làm bài: 90 phút (không kể thời gian phát đề).}
%  Chuyên đề: \khungtieude{CHUYÊN ĐỀ 3 -- DÃY SỐ VÀ CÔNG THỨC TRUY HỒI}
%                         {Bồi dưỡng học sinh giỏi Toán THPT}{...}{}
\newcommand{\khungtieude}[4]{%
	\begin{tcolorbox}[colframe=gold1,colback=gold5!45,boxrule=0.5mm]
		\begin{minipage}[c]{3cm}
			\centerline{\eduzylogo[2.4cm]}
		\end{minipage}\hspace{0.5cm}
		\begin{minipage}[c]{12.5cm}
			\vskip 0.15 cm
			\centerline{\ind{#1}}
			\vskip 0.15 cm
			\centerline{\textbf{#2}}
			\if\relax\detokenize{#3}\relax\else
				\vskip 0.15 cm
				\centerline{\textit{#3}}
			\fi
			\if\relax\detokenize{#4}\relax\else
				\vskip 0.15 cm
				\centerline{\textit{#4}}
			\fi
			\vskip 0.15 cm
		\end{minipage}
	\end{tcolorbox}
}
%====================== BỔ SUNG CHO BÀI HỌC, CHUYÊN ĐỀ ======================
\usepackage[shortlabels]{enumitem}
\usepackage{tabularx}
%--- Danh sách gọn: \begin{enumerate}[gon] ... \end{enumerate}
\SetEnumitemKey{gon}{leftmargin=*,itemsep=1pt,parsep=0pt,topsep=3pt}
%--- Khung lý thuyết: \begin{khung}{} ... \end{khung}; tiêu đề chỉ dùng khi thật ngắn ("Mẹo", "Lỗi thường gặp")
\newtcolorbox{khung}[1]{breakable,enhanced,colframe=gold1,colback=gold5!45,
	boxrule=0.4mm,fonttitle=\bfseries\sffamily,coltitle=white,colbacktitle=gold1,
	attach boxed title to top left={xshift=6mm,yshift=-3mm},
	boxed title style={arc=1mm,boxrule=0.3mm},title={#1}}
%--- Tiêu đề mục nhỏ trong một dạng: \muc{Dấu hiệu}, \muc{Phương pháp}, \muc{Bài tự luyện}
\makeatletter
\newcommand{\muc}[1]{\par\medskip\noindent\indam{#1}\par\nobreak\smallskip\@afterindentfalse\@afterheading}
\makeatother
%--- Ô công thức cần nhớ (căn giữa, viền vàng, vừa khít nội dung)
%     \begin{center}\congthuc{$Q(n)=(n-1)!$}\end{center}
\newtcbox{\congthuc}{on line,colframe=gold1,colback=gold5!45,boxrule=0.5pt,arc=2pt,
	left=8pt,right=8pt,top=4pt,bottom=4pt}
%--- Công tắc lời giải bài tự luyện (ví dụ mẫu luôn dùng \loigiai)
%     \loigiaiTL[{đáp số}]{lời giải đầy đủ}
%     \TLfull -> in lời giải đầy đủ;  \TLdapso -> chỉ in một dòng "Đáp số";  \TLan -> ẩn hoàn toàn
%     (\CheDo là dethi thì mọi lời giải và đáp số đều ẩn)
\makeatletter
\newif\ifTLhien   \TLhientrue
\newif\ifTLdapso  \TLdapsofalse
\newcommand{\TLfull}{\TLhientrue\TLdapsofalse}
\newcommand{\TLdapso}{\TLhientrue\TLdapsotrue}
\newcommand{\TLan}{\TLhienfalse}
\newcommand{\loigiaiTL}[2][]{%
	\ifTLhien
		\ifTLdapso
			\def\@tlarg{#1}%
			\ifx\@tlarg\@empty\else
				\ifdefstring{\kindOPTN}{dethi}{}{\par\noindent{\bfseries\color{gold1}Đáp số:} #1\par}%
			\fi
		\else
			\loigiai{#2}%
		\fi
	\fi}
\makeatother
\TLdapso% mặc định: chỉ in đáp số bài tự luyện
%====================== BẢNG ĐÚNG/SAI (\choiceTFt) ======================
%  Hàng tiêu đề nền nâu đồng chữ trắng (Phát biểu | Đ | S), các hàng tô xen kẽ kem / trắng,
%  đường kẻ vàng kim, nhãn a) b) c) d) màu nâu đồng. Bản có đáp án: cột được chọn có vòng tròn
%  nâu đồng với dấu tích trắng (câu đúng tích ở cột Đúng, câu sai tích ở cột Sai).
%  Tiêu đề ghi đủ "Đúng | Sai": \OPTN{viettat=0} sau \input (cột Đúng/Sai rộng ra, phát biểu hẹp lại).
%  Dùng bảng gốc của ex_test: \def\BangDSgoc{1} trước \input{preamble-eduzy-v2}.
\providecommand{\BangDSgoc}{0}
\ifnum\BangDSgoc=0
\usepackage{colortbl,pifont}
\colorlet{ezDSke}{gold2!60!gold1}% đường kẻ bảng (nét 0,4pt như ex_test để không đổi chiều cao bảng)
\colorlet{ezDSle}{gold5!60}% hàng a), c)
\colorlet{ezDSchan}{white}% hàng b), d)
%--- dấu chọn ở cột Đúng/Sai
%    (bề rộng 0, căn giữa ô: không làm đổi bề rộng cột của ex_test, không báo tràn lề)
\newcommand{\EZdauchon}{\makebox[0pt]{\tikz[baseline=-1ex]\node[circle,fill=gold1,inner sep=0pt,
	minimum size=1.2em,text=white,font=\small]{\ding{51}};}}
\def\tickT{\EZdauchon}\def\tickF{\EZdauchon}
\AtBeginDocument{\sbox{\ccTFTab}{\EZdauchon}}
%--- nhãn a) b) c) d) màu nâu đồng (ở bản lời giải, nhãn câu đúng vẫn đóng khung như ex_test)
\renewcommand{\FalseTF}{\stepcounter{dapan}\textbf{\color{gold1}\DapAnTF\sepTF}\ignorespaces}
%--- trước khi dựng bảng: kẻ màu vàng kim (bề rộng cột giữ nguyên như ex_test); sau bảng trả lại màu đen
\newcount\EZhangTF
\newcommand{\EZbatdauTF}{%
	\global\EZhangTF=0\relax
	\arrayrulecolor{ezDSke}}
\newcommand{\EZketthucTF}{\arrayrulecolor{black}}
%--- hàng tiêu đề nền nâu đồng, chữ trắng
\makeatletter
\newcommand{\EZo@TF}[1]{\centering\bfseries\color{white}#1}
\newcommand{\EZtieudeTF}{%
	\ifcase\cotDS
		\ifnum\cotSTT=1\relax
			\def\rowHeader{\hline\rowcolor{gold1}\EZo@TF{TT}&\EZo@TF{\phatbieu}\tabularnewline}%
		\else
			\def\rowHeader{\hline\rowcolor{gold1}\EZo@TF{\phatbieu}\tabularnewline}%
		\fi
	\or
		\ifnum\cotSTT=1\relax
			\def\rowHeader{\hline\rowcolor{gold1}\EZo@TF{TT}&\EZo@TF{\phatbieu}&\EZo@TF{\Dung}\tabularnewline}%
		\else
			\def\rowHeader{\hline\rowcolor{gold1}\EZo@TF{\phatbieu}&\EZo@TF{\Dung}\tabularnewline}%
		\fi
	\or
		\ifnum\cotSTT=1\relax
			\def\rowHeader{\hline\rowcolor{gold1}\EZo@TF{TT}&\EZo@TF{\phatbieu}&\EZo@TF{\Dung}&\EZo@TF{\Sai}\tabularnewline}%
		\else
			\def\rowHeader{\hline\rowcolor{gold1}\EZo@TF{\phatbieu}&\EZo@TF{\Dung}&\EZo@TF{\Sai}\tabularnewline}%
		\fi
	\fi}
%--- tô xen kẽ từng hàng (gọi ngay sau \hline đầu mỗi hàng)
\newcommand{\EZmauhangTF}{%
	\noalign{\global\advance\EZhangTF\@ne}%
	\ifodd\EZhangTF\rowcolor{ezDSle}\else\rowcolor{ezDSchan}\fi}
%--- gắn vào \choiceTFt / \writeChoiceTFt của ex_test
\xpretocmd{\choiceTFt}{\EZbatdauTF}{}{\typeout{[preamble-eduzy-v2] Không gắn được bảng Đúng/Sai (đầu bảng)}}
\xpatchcmd{\choiceTFt}{\ifnum\dongPB=0\def\rowHeader{}\fi}{\EZtieudeTF\ifnum\dongPB=0\def\rowHeader{}\fi}
	{}{\typeout{[preamble-eduzy-v2] Không gắn được bảng Đúng/Sai (hàng tiêu đề)}}
\xpatchcmd{\writeChoiceTFt}{\hline\xdef\dotEX{\sdotEX}}{\hline\EZmauhangTF\xdef\dotEX{\sdotEX}}
	{}{\typeout{[preamble-eduzy-v2] Không gắn được bảng Đúng/Sai (tô hàng)}}
\xpatchcmd{\writeChoiceTFt}{{\end{longtable}}}{{\end{longtable}\EZketthucTF}}
	{}{\typeout{[preamble-eduzy-v2] Không gắn được bảng Đúng/Sai (cuối longtable)}}
\xpatchcmd{\writeChoiceTFt}{{\end{xtabular}\end{center}\medskip\noindent}}{{\end{xtabular}\end{center}\EZketthucTF\medskip\noindent}}
	{}{\typeout{[preamble-eduzy-v2] Không gắn được bảng Đúng/Sai (cuối xtabular)}}
\fi
%=====================================================================
\renewcommand{\parallel}{\mathrel{/\!\!/}}

%      \TLdapso                      % tài liệu: \TLfull | \TLdapso | \TLan
%      \begin{document} ... \end{document}
%  Logo Eduzy vẽ bằng TikZ ngay trong file (\eduzylogo), không cần file ảnh.
%  Biên dịch: pdfLaTeX (máy có gói vietnam) hoặc XeLaTeX.
%============================================================
\documentclass[12pt,a4paper,twoside]{article}%<-- twoside: bat dau trang/chan trang chan-le
\usepackage{iftex}
\usepackage{amsmath,amssymb,tikz,mathrsfs,tkz-tab}
\usepackage{graphicx}
\usepackage{pgfplots}% đồ thị hàm số
\pgfplotsset{compat=1.17}
%\usepackage{mathptmx}
\usepackage{mathpazo}
\ifXeTeX% XeLaTeX: font chữ có đủ dấu tiếng Việt, giữ nguyên font toán mathpazo
	\usepackage[no-math]{fontspec}
	\setmainfont{TeX Gyre Pagella}
	\newfontfamily\EZqag{TeX Gyre Adventor}[NFSSFamily=qag,Ligatures=TeX]
\fi
\usepackage[many]{tcolorbox}
\usetikzlibrary{math,through,calc,intersections,angles,quotes,shapes,shapes.geometric,arrows,patterns,matrix,chains,arrows.meta,decorations.shapes,decorations.fractals,decorations.markings,shadows,decorations.pathreplacing}
\usepackage[top=1.6cm, bottom=1.8cm, left=1.8cm, right=1.8cm]{geometry}% v2: top 1.6cm (bản cũ 1.5cm) để chữ không sát dải đầu trang
\providecommand{\CheDo}{loigiai}% loigiai | dethi | solcolor | color | book
\expandafter\usepackage\expandafter[\CheDo]{ex_test}
% \tenchuong: khai báo ở file chính (trước \input) hoặc \renewcommand sau \input
%============================================================
%  [GỘP TỪ eduzy-de.tex] EDUZY - MÀU NÂU ĐỒNG / VÀNG KIM
%  + ĐẦU TRANG, CHÂN TRANG KIỂU DẢI NGANG (v2), ĐỐI XỨNG TRANG CHẴN / TRANG LẺ.
%  (không cần file eduzy-de.tex riêng nữa)
%  Khối này phải đặt SAU \usepackage[...]{ex_test} và SAU \def\tenchuong{...}
%============================================================
\usepackage{lastpage}
\usepackage{needspace}
%--- DÀN TRANG ĐỀ THI: tách câu (THÊM 05/11/2026; chỉ thêm lệnh mới, KHÔNG đổi hành vi mặc định) ---
%  Hai công tắc, mặc định tắt (= preamble chạy y như cũ). Bật bằng cách đặt TRƯỚC %============================================================
%  PREAMBLE EDUZY v2 - dùng chung cho đề thi, bài học, chuyên đề (gói ex_test)
%  Khác bản preamble-eduzy.tex ở đầu trang, chân trang: dải nền nâu đồng tràn ngang,
%  logo + Eduzy.com.vn, ô chữ vàng kem cạnh xiên, chân trang "Trang (số)";
%  logo Eduzy chìm giữa mọi trang (tắt: \logochimtat, bật lại: \logochimbat).
%  Cách dùng ở file chính (xem mau-chuan.tex, mau-chuyen-de.tex):
%      \def\CheDo{loigiai}          % loigiai | dethi | solcolor | color | book
%      \def\tenchuong{<dòng chữ trên đầu trang>}
%      \def\thuonghieu{<chữ ở chân trang lẻ>}   % không khai báo thì dùng mặc định
%      \def\tacgia{<chữ ở chân trang chẵn>}      % không khai báo thì dùng mặc định
%      \def\khauhieu{<chữ cạnh logo ở đầu trang>}  % tùy chọn, mặc định Eduzy.com.vn
%      \input{preamble-eduzy-v2}
%      \TLdapso                      % tài liệu: \TLfull | \TLdapso | \TLan
%      \begin{document} ... \end{document}
%  Logo Eduzy vẽ bằng TikZ ngay trong file (\eduzylogo), không cần file ảnh.
%  Biên dịch: pdfLaTeX (máy có gói vietnam) hoặc XeLaTeX.
%============================================================
\documentclass[12pt,a4paper,twoside]{article}%<-- twoside: bat dau trang/chan trang chan-le
\usepackage{iftex}
\usepackage{amsmath,amssymb,tikz,mathrsfs,tkz-tab}
\usepackage{graphicx}
\usepackage{pgfplots}% đồ thị hàm số
\pgfplotsset{compat=1.17}
%\usepackage{mathptmx}
\usepackage{mathpazo}
\ifXeTeX% XeLaTeX: font chữ có đủ dấu tiếng Việt, giữ nguyên font toán mathpazo
	\usepackage[no-math]{fontspec}
	\setmainfont{TeX Gyre Pagella}
	\newfontfamily\EZqag{TeX Gyre Adventor}[NFSSFamily=qag,Ligatures=TeX]
\fi
\usepackage[many]{tcolorbox}
\usetikzlibrary{math,through,calc,intersections,angles,quotes,shapes,shapes.geometric,arrows,patterns,matrix,chains,arrows.meta,decorations.shapes,decorations.fractals,decorations.markings,shadows,decorations.pathreplacing}
\usepackage[top=1.6cm, bottom=1.8cm, left=1.8cm, right=1.8cm]{geometry}% v2: top 1.6cm (bản cũ 1.5cm) để chữ không sát dải đầu trang
\providecommand{\CheDo}{loigiai}% loigiai | dethi | solcolor | color | book
\expandafter\usepackage\expandafter[\CheDo]{ex_test}
% \tenchuong: khai báo ở file chính (trước \input) hoặc \renewcommand sau \input
%============================================================
%  [GỘP TỪ eduzy-de.tex] EDUZY - MÀU NÂU ĐỒNG / VÀNG KIM
%  + ĐẦU TRANG, CHÂN TRANG KIỂU DẢI NGANG (v2), ĐỐI XỨNG TRANG CHẴN / TRANG LẺ.
%  (không cần file eduzy-de.tex riêng nữa)
%  Khối này phải đặt SAU \usepackage[...]{ex_test} và SAU \def\tenchuong{...}
%============================================================
\usepackage{lastpage}
\usepackage{needspace}
%--- DÀN TRANG ĐỀ THI: tách câu (THÊM 05/11/2026; chỉ thêm lệnh mới, KHÔNG đổi hành vi mặc định) ---
%  Hai công tắc, mặc định tắt (= preamble chạy y như cũ). Bật bằng cách đặt TRƯỚC \input{preamble-eduzy-v2}
%  trong file đề (tools/dan_trang/dan_trang.py tự ghi khi đặt \NeedspaceCau):
%    \def\TachCauDS{1}   câu Đúng/Sai được tách: câu dẫn (kèm hình) ở cuối trang, bảng Phát biểu sang trang sau
%                        (câu dẫn và hình không tách nhau, bảng không tách giữa hai trang)
%    \def\TachCauLG{1}   bản lời giải (\CheDo khác dethi): phần đề của câu cao hơn 1/3 trang được phép tách
%  Lệnh dùng trong file đề (giá trị do tools/dan_trang/dan_trang.py đo ở bản dethi):
%    \NeedspaceCau[<cao câu dẫn + hình>]{<cao cả khối câu>}   ngay trên \begin{ex} (câu đầu phần: trên tiêu đề phần)
%        công tắc tắt: đúng bằng \Needspace{<cao cả khối câu>} như cũ;
%        \TachCauDS=1 và có [..] (câu Đúng/Sai được chọn tách): chỉ giữ câu dẫn + hình;
%        \TachCauLG=1 ở bản lời giải: đòi tối đa 1/3 \textheight.
%    \NeedspaceBang{<cao bảng Phát biểu>}%   ngay trước \choiceTFt: \TachCauDS=1 thì giữ trọn bảng trên một trang;
%        công tắc tắt: không làm gì.
\providecommand{\TachCauDS}{0}
\providecommand{\TachCauLG}{0}
\newlength{\EZnsCau}
\newcommand{\NeedspaceCau}[2][]{%
	\setlength{\EZnsCau}{#2}%
	\ifnum\TachCauDS=1 \if\relax\detokenize{#1}\relax\else\setlength{\EZnsCau}{#1}\fi\fi
	\ifnum\TachCauLG=1 \ifdefstring{\kindOPTN}{dethi}{}{\ifdim\EZnsCau>0.3333\textheight\setlength{\EZnsCau}{0.3333\textheight}\fi}\fi
	\Needspace{\EZnsCau}}
\newcommand{\NeedspaceBang}[1]{\ifnum\TachCauDS=1 \Needspace{#1}\fi}

%--- Thang màu: đậm -> nhạt
\definecolor{gold1}{RGB}{126,96,0}    % nâu đồng đậm  #7E6000
\definecolor{gold2}{RGB}{192,144,3}   % vàng kim      #C09003
\definecolor{gold3}{RGB}{255,215,99}  % vàng nhạt     #FFD763
\definecolor{gold4}{RGB}{255,231,153} % vàng rất nhạt #FFE799
\definecolor{gold5}{RGB}{255,243,205} % vàng kem      #FFF3CD
%--- Ghi đè tông xanh/đỏ cũ
\definecolor{myblue}{RGB}{126,96,0}
\colorlet{blue}{gold1}
\colorlet{cyan}{gold2}
\colorlet{teal}{gold2}

%--- Tên chương mặc định (file đề nên tự khai báo lại)
\providecommand{\tenchuong}{}
%--- Chân trang
\providecommand{\thuonghieu}{TOÁN 12 -- Eduzy.com.vn}
\providecommand{\tacgia}{Tô Gia Bảo -- 0353703783}

%--- Chữ cạnh logo ở đầu trang
\providecommand{\khauhieu}{Eduzy.com.vn}

%====================== THÔNG SỐ ĐẦU TRANG, CHÂN TRANG (v2) ======================
%  Dải nền đậm tràn ngang trang; ô chữ vàng kem có cạnh xiên, kèm 2 vạch xiên chuyển màu.
%  Đầu trang lẻ : logo + \khauhieu (trái)   | ô chữ \tenchuong (phải)
%  Chân trang lẻ: ô chữ \thuonghieu (trái)  | Trang (số) (phải)
%  Trang chẵn soi gương trang lẻ; chân trang chẵn ghi \tacgia.
%--- Màu: chỉ dùng thang vàng kim Eduzy
\colorlet{ezNen}{gold1}% dải nền đậm
\colorlet{ezVach1}{gold3}% vạch xiên sát dải nền (đậm hơn)
\colorlet{ezVach2}{gold4}% vạch xiên sát ô chữ (nhạt hơn)
\colorlet{ezO}{gold5}% ô chữ vàng kem
%--- Kích thước
\newlength{\EZdai}\setlength{\EZdai}{1.1cm}% chiều cao dải nền
\newlength{\EZo}\setlength{\EZo}{0.68cm}% chiều cao ô chữ
\newlength{\EZxien}\setlength{\EZxien}{0.55cm}% độ lệch ngang của cạnh xiên
\newlength{\EZvach}\setlength{\EZvach}{0.75cm}% bề rộng mỗi vạch xiên
\newlength{\EZkhe}\setlength{\EZkhe}{0.12cm}% khe trắng giữa vạch và ô chữ
\newlength{\EZdem}\setlength{\EZdem}{0.4cm}% khoảng đệm từ chữ tới cạnh xiên
\newlength{\EZle}\setlength{\EZle}{1.8cm}% lề trái/phải của chữ (bằng lề thân bài)
\newlength{\EZrongmin}\setlength{\EZrongmin}{4.5cm}% ô chữ rộng tối thiểu
\newlength{\EZrong}% bề rộng chữ trong ô (tính tự động)
%--- Font
\newcommand{\EZfontO}{\fontsize{10.5pt}{12pt}\selectfont\bfseries}% chữ trong ô
\newcommand{\EZfontKH}{\fontfamily{qag}\fontsize{11pt}{12pt}\selectfont\bfseries}% \khauhieu cạnh logo
\newcommand{\EZfontTrang}{\fontsize{11.5pt}{12pt}\selectfont\bfseries}% chữ "Trang" và số trang
%--- Đo bề rộng chữ trong ô
\newcommand{\EZdorong}[1]{%
	\settowidth{\EZrong}{\EZfontO #1}%
	\ifdim\EZrong<\EZrongmin\setlength{\EZrong}{\EZrongmin}\fi}
%--- Logo + chữ trắng \khauhieu (dựng sẵn một lần khi vào \begin{document})
\newsavebox{\EZlogobox}
\AtBeginDocument{\sbox{\EZlogobox}{%
		\raisebox{-0.5\height}{\eduzylogo[1cm]}\hspace{0.15cm}%
		\raisebox{-0.5\height}{\color{white}\EZfontKH\khauhieu}}}
%--- Hai vạch xiên + khe trắng. #1: điểm gốc ở mép dưới, #2: hướng (1 sang phải, -1 sang trái),
%    #3: hướng nghiêng (1 kiểu "/", -1 kiểu "\")
\newcommand{\EZvachxien}[3]{%
	\fill[white] (#1) --++({(#3)*\EZxien},\EZo) --++({(#2)*\EZkhe},0) --++({-(#3)*\EZxien},-\EZo) --cycle;
	\fill[ezVach2] ([xshift=(#2)*\EZkhe]#1) --++({(#3)*\EZxien},\EZo) --++({(#2)*\EZvach},0) --++({-(#3)*\EZxien},-\EZo) --cycle;
	\fill[ezVach1] ([xshift=(#2)*(\EZkhe+\EZvach)]#1) --++({(#3)*\EZxien},\EZo) --++({(#2)*\EZvach},0) --++({-(#3)*\EZxien},-\EZo) --cycle;
	\draw[white,line width=0.6pt] ([xshift=(#2)*(\EZkhe+\EZvach)]#1) --++({(#3)*\EZxien},\EZo);
}
%--- Chữ trong ô (căn giữa theo chiều cao ô, không lệch vì dấu tiếng Việt)
\tikzset{ezchu/.style={inner sep=0pt,text=ezNen,text height=0.72em,text depth=0.22em}}

%====================== LOGO CHÌM GIỮA TRANG ======================
%  Vẽ cùng đầu trang nên nằm dưới chữ của thân bài, có trên mọi trang.
%  Tắt: \logochimtat (từ trang đó trở đi), bật lại: \logochimbat.
%  Đổi cỡ, độ đậm: khai báo trước \input, ví dụ \def\EZchimcao{9cm} \def\EZchimdo{0.1}
\providecommand{\EZchimcao}{10cm}% chiều cao logo chìm
\providecommand{\EZchimdo}{0.12}% độ đậm (0 = trong suốt, 1 = rõ hẳn)
\newif\ifEZchim \EZchimtrue
\newcommand{\logochimtat}{\EZchimfalse}
\newcommand{\logochimbat}{\EZchimtrue}
%  Dựng sẵn logo mờ một lần khi vào \begin{document}; transparency group làm mờ cả logo
%  một lượt nên không lộ vệt ở chỗ các dải màu chồng nhau.
\newsavebox{\EZchimtho}
\newsavebox{\EZchimbox}
\AtBeginDocument{%
	\sbox{\EZchimtho}{\eduzylogo[\EZchimcao]}%
	\sbox{\EZchimbox}{\begin{tikzpicture}
			\begin{scope}[transparency group,opacity=\EZchimdo]
				\node[inner sep=0pt] {\usebox{\EZchimtho}};
			\end{scope}
	\end{tikzpicture}}}
\newcommand{\EZvechim}{%
	\ifEZchim
		\node[inner sep=0pt] at (current page.center) {\usebox{\EZchimbox}};
	\fi}

%====================== ĐẦU TRANG - TRANG LẺ ======================
%  (logo bên trái, ô chữ \tenchuong bên phải)
\newcommand{\myfancyheadOdd}{%
	\EZdorong{\tenchuong}%
	\begin{tikzpicture}[remember picture,overlay]
		\coordinate (T) at (current page.north west);
		\coordinate (P) at (current page.north east);
		%--- logo chìm giữa trang
		\EZvechim
		%--- dải nền
		\fill[ezNen] (T) rectangle ([yshift=-\EZdai]P);
		%--- ô chữ: X là góc dưới bên trái, cạnh trái xiên kiểu "/"
		\coordinate (X) at ([xshift=-\EZle-\EZrong-\EZdem-\EZxien,yshift=-\EZdai]P);
		\fill[ezO] (X) --++(\EZxien,\EZo) -- ([yshift=-\EZdai+\EZo]P) -- ([yshift=-\EZdai]P) --cycle;
		\EZvachxien{X}{-1}{1}
		\node[ezchu,anchor=east] at ([xshift=-\EZle,yshift=-\EZdai+0.5*\EZo]P) {\EZfontO\tenchuong};
		%--- logo + tên
		\node[inner sep=0pt,anchor=west] at ([xshift=\EZle,yshift=-0.5*\EZdai]T) {\usebox{\EZlogobox}};
	\end{tikzpicture}%
}

%====================== ĐẦU TRANG - TRANG CHẴN =====================
%  (soi gương: ô chữ \tenchuong bên trái, logo bên phải)
\newcommand{\myfancyheadEven}{%
	\EZdorong{\tenchuong}%
	\begin{tikzpicture}[remember picture,overlay]
		\coordinate (T) at (current page.north west);
		\coordinate (P) at (current page.north east);
		%--- logo chìm giữa trang
		\EZvechim
		%--- dải nền
		\fill[ezNen] (T) rectangle ([yshift=-\EZdai]P);
		%--- ô chữ: X là góc dưới bên phải, cạnh phải xiên kiểu "\"
		\coordinate (X) at ([xshift=\EZle+\EZrong+\EZdem+\EZxien,yshift=-\EZdai]T);
		\fill[ezO] (X) --++(-\EZxien,\EZo) -- ([yshift=-\EZdai+\EZo]T) -- ([yshift=-\EZdai]T) --cycle;
		\EZvachxien{X}{1}{-1}
		\node[ezchu,anchor=west] at ([xshift=\EZle,yshift=-\EZdai+0.5*\EZo]T) {\EZfontO\tenchuong};
		%--- logo + tên
		\node[inner sep=0pt,anchor=east] at ([xshift=-\EZle,yshift=-0.5*\EZdai]P) {\usebox{\EZlogobox}};
	\end{tikzpicture}%
}

%====================== CHÂN TRANG - TRANG LẺ ======================
%  (ô chữ \thuonghieu bên trái, "Trang (số)" bên phải)
\newcommand{\myfancyfootOdd}{%
	\EZdorong{\thuonghieu}%
	\begin{tikzpicture}[remember picture,overlay]
		\coordinate (T) at (current page.south west);
		\coordinate (P) at (current page.south east);
		%--- dải nền
		\fill[ezNen] (T) rectangle ([yshift=\EZdai]P);
		%--- ô chữ sát mép trên dải: X là góc dưới bên phải, cạnh phải xiên kiểu "/"
		\coordinate (X) at ([xshift=\EZle+\EZrong+\EZdem,yshift=\EZdai-\EZo]T);
		\fill[ezO] ([yshift=\EZdai-\EZo]T) -- (X) --++(\EZxien,\EZo) -- ([yshift=\EZdai]T) --cycle;
		\EZvachxien{X}{1}{1}
		\node[ezchu,anchor=west] at ([xshift=\EZle,yshift=\EZdai-0.5*\EZo]T) {\EZfontO\thuonghieu};
		%--- số trang
		\node[circle,draw=white,line width=0.9pt,text=white,inner sep=0pt,minimum size=0.62cm,anchor=east]
		(so) at ([xshift=-\EZle,yshift=0.5*\EZdai]P) {\EZfontTrang\thepage};
		\node[inner sep=0pt,text=white,anchor=base east] (tr) at ([xshift=-0.15cm]so.base west) {\EZfontTrang Trang};
	\end{tikzpicture}%
}

%====================== CHÂN TRANG - TRANG CHẴN ====================
%  (soi gương: "Trang (số)" bên trái, ô chữ \tacgia bên phải)
\newcommand{\myfancyfootEven}{%
	\EZdorong{\tacgia}%
	\begin{tikzpicture}[remember picture,overlay]
		\coordinate (T) at (current page.south west);
		\coordinate (P) at (current page.south east);
		%--- dải nền
		\fill[ezNen] (T) rectangle ([yshift=\EZdai]P);
		%--- ô chữ sát mép trên dải: X là góc dưới bên trái, cạnh trái xiên kiểu "\"
		\coordinate (X) at ([xshift=-\EZle-\EZrong-\EZdem,yshift=\EZdai-\EZo]P);
		\fill[ezO] ([yshift=\EZdai-\EZo]P) -- (X) --++(-\EZxien,\EZo) -- ([yshift=\EZdai]P) --cycle;
		\EZvachxien{X}{-1}{-1}
		\node[ezchu,anchor=east] at ([xshift=-\EZle,yshift=\EZdai-0.5*\EZo]P) {\EZfontO\tacgia};
		%--- số trang
		%    (vòng tròn ẩn so0 nằm giữa dải để lấy đường chân chữ, giống trang lẻ)
		\node[circle,inner sep=0pt,minimum size=0.62cm] (so0) at ([yshift=0.5*\EZdai]T) {\EZfontTrang\phantom{\thepage}};
		\node[inner sep=0pt,text=white,anchor=base west] (tr) at ($(so0.base -| T)+(\EZle,0)$) {\EZfontTrang Trang};
		\node[circle,draw=white,line width=0.9pt,text=white,inner sep=0pt,minimum size=0.62cm,anchor=base west]
		(so) at ([xshift=0.15cm]tr.base east) {\EZfontTrang\thepage};
	\end{tikzpicture}%
}

%--- Giữ tương thích với các file cũ đang gọi \myfancyhead, \myfancyfoot
\newcommand{\myfancyhead}{\myfancyheadOdd}
\newcommand{\myfancyfoot}{\myfancyfootOdd}

%======================= GẮN VÀO TRANG =============================
\pagestyle{fancy}
\fancyhf{}
\fancyhead[LO]{\myfancyheadOdd}
\fancyhead[RE]{\myfancyheadEven}
\fancyfoot[LO]{\myfancyfootOdd}
\fancyfoot[RE]{\myfancyfootEven}
\renewcommand{\headrulewidth}{0pt}
\renewcommand{\footrulewidth}{0pt}
\fancypagestyle{plain}{%
	\fancyhf{}
	\fancyhead[LO]{\myfancyheadOdd}
	\fancyhead[RE]{\myfancyheadEven}
	\fancyfoot[LO]{\myfancyfootOdd}
	\fancyfoot[RE]{\myfancyfootEven}
	\renewcommand{\headrulewidth}{0pt}
	\renewcommand{\footrulewidth}{0pt}
}
%=================== HẾT PHẦN GỘP TỪ eduzy-de.tex ===================
\newcommand{\indam}[1]{
	{\fontsize{12pt}{9pt}\sffamily \selectfont\bfseries\color{gold1} #1}
}
\newcommand{\ind}[1]{
	{\fontsize{14pt}{9pt}\fontfamily{qag} \selectfont\bfseries\color{gold1} #1}
}
\renewcommand{\FalseEX}{\stepcounter{dapan}\textbf{\selectfont\color{gold1}\Alph{dapan}.}\ignorespaces}
\renewtheorem{ex}{\color{gold1}Câu}
\renewcommand{\baselinestretch}{1.1}
\newcommand{\hoac}[1]{
	\left[\begin{aligned}#1\end{aligned}\right.}
\newcommand{\heva}[1]{
	\left\{\begin{aligned}#1\end{aligned}\right.}
\usepackage{esvect}
\def\vec{\vv}
\def\vt{\vv}
\def\overrightarrow{\vv}
%====================== LOGO EDUZY VẼ BẰNG TIKZ ======================
%  Đường viền dò từ logo gốc; dải màu vàng kim đo trực tiếp trên ảnh gốc
%  (nghiêng 21 độ, tô bằng 80 dải hẹp nên hiển thị giống nhau trên mọi trình xem PDF).
%  Dùng: \eduzylogo          -> cao 2.4cm (mặc định)
%        \eduzylogo[1.5cm]   -> cao 1.5cm
\definecolor{ezText}{RGB}{245,200,46}
\newcommand{\eduzylogo}[1][2.4cm]{%
	\resizebox{!}{#1}{%
		\begin{tikzpicture}[x=0.01cm,y=0.01cm]
			\path[use as bounding box] (0,0) rectangle (224,-224);
			% mũ tốt nghiệp + vương miện
			\begin{scope}
				\clip
				(119.4,-176.1)..controls(118.9,-176.1)and(116.8,-176)..(114.6,-176)
				..controls(101.4,-176)and(85.8,-175.9)..(83.7,-175.7)
				..controls(81.8,-175.6)and(81,-175.6)..(80.2,-175.7)
				..controls(78.8,-175.9)and(77.7,-175.9)..(76.2,-175.7)
				..controls(74.9,-175.5)and(70.8,-175.5)..(64.3,-175.7)
				..controls(59.1,-175.9)and(58.3,-175.7)..(58.6,-173.7)
				..controls(58.6,-172.8)and(59.6,-169.6)..(60.3,-168)
				..controls(60.8,-166.9)and(61.4,-165.8)..(62.3,-165)
				..controls(63.1,-164.2)and(64.3,-163.4)..(65,-163.2)
				..controls(65.5,-163.1)and(72.3,-163)..(76.6,-163.1)
				..controls(77.9,-163.1)and(80.1,-163.1)..(81.4,-163)
				..controls(82.8,-162.9)and(86.3,-162.9)..(89.2,-162.9)
				..controls(93,-162.9)and(94.9,-162.8)..(96,-162.7)
				..controls(97.2,-162.5)and(101,-162.5)..(114,-162.4)
				..controls(123,-162.3)and(134.9,-162.2)..(140.2,-162.1)
				..controls(150.1,-162)and(152.2,-162.1)..(153.6,-162.4)
				..controls(154.4,-162.6)and(157.1,-163.8)..(158.1,-164.4)
				..controls(158.8,-164.8)and(159.7,-165.7)..(159.9,-166.3)
				..controls(160,-166.5)and(160.2,-167.1)..(160.2,-167.6)
				..controls(160.3,-168)and(160.6,-169.4)..(160.9,-170.6)
				..controls(161.2,-171.7)and(161.5,-173.1)..(161.6,-173.6)
				..controls(161.8,-174.4)and(161.8,-174.5)..(161.5,-174.9)
				..controls(161.2,-175.6)and(160.8,-175.9)..(159.9,-176)
				..controls(159.2,-176.1)and(154.1,-176.2)..(151.5,-176.2)
				..controls(148.1,-176.1)and(122.6,-176.1)..(121.6,-176.2)
				..controls(120.8,-176.2)and(119.8,-176.2)..(119.4,-176.1)--cycle
				(130.4,-160)..controls(129,-160)and(119.8,-160)..(110.1,-159.9)
				..controls(100.3,-159.8)and(88.5,-159.7)..(83.9,-159.6)
				..controls(77.8,-159.5)and(75.3,-159.5)..(74.5,-159.6)
				..controls(73.9,-159.7)and(72.5,-159.8)..(71.3,-159.8)
				..controls(69.6,-159.8)and(69,-159.8)..(68.5,-159.6)
				..controls(67.6,-159.3)and(66.8,-158.6)..(66.2,-157.6)
				..controls(66,-157.3)and(65.6,-156.5)..(65.3,-156.1)
				..controls(64.6,-155.1)and(64.5,-154.9)..(63.7,-152.9)
				..controls(63.2,-151.6)and(62.9,-150.9)..(61.8,-149.6)
				..controls(61.5,-149.2)and(61.1,-148.5)..(60.9,-147.9)
				..controls(60.7,-147.4)and(60.3,-146.6)..(60.1,-146.1)
				..controls(59.9,-145.5)and(59.5,-144.9)..(59,-144.3)
				..controls(58.5,-143.7)and(58.2,-143.2)..(57.9,-142.5)
				..controls(57.7,-141.9)and(57.3,-141.1)..(57,-140.7)
				..controls(56.6,-140.3)and(56.1,-139.3)..(55.8,-138.5)
				..controls(55.5,-137.8)and(55.1,-136.8)..(54.9,-136.5)
				..controls(54.7,-136.1)and(54.3,-135.2)..(53.9,-134.5)
				..controls(53.6,-133.8)and(53,-132.7)..(52.6,-132)
				..controls(51.6,-130.5)and(51.4,-130.2)..(51.1,-129)
				..controls(50.6,-127.5)and(50.6,-127.4)..(49.7,-126.3)
				..controls(49.2,-125.7)and(48.6,-124.8)..(48.3,-124.2)
				..controls(48,-123.6)and(47.5,-122.7)..(47.3,-122.2)
				..controls(47,-121.8)and(46.5,-120.7)..(46.1,-119.8)
				..controls(45.8,-118.9)and(45.2,-117.8)..(44.9,-117.2)
				..controls(44.2,-116.1)and(43.3,-114.4)..(42.3,-112.2)
				..controls(41.8,-111.3)and(41.3,-110.4)..(41.1,-110.1)
				..controls(40.5,-109.4)and(40.3,-109)..(40,-107.8)
				..controls(39.8,-107)and(39.6,-106.5)..(39.2,-106)
				..controls(38.9,-105.6)and(38.4,-104.7)..(38,-104)
				..controls(37.7,-103.4)and(37.2,-102.5)..(37,-102.1)
				..controls(36.7,-101.7)and(36.3,-100.8)..(36,-100.1)
				..controls(35.6,-98.9)and(35.3,-98.4)..(34.6,-97.6)
				..controls(34.2,-97.1)and(34.2,-96.3)..(34.6,-96)
				..controls(35,-95.7)and(35.5,-95.9)..(36.2,-96.7)
				..controls(36.5,-97)and(37.2,-97.5)..(37.6,-97.7)..controls(38,-98)and(38.6,-98.4)..(39.1,-98.9)
				..controls(39.6,-99.4)and(40.3,-100)..(40.6,-100.3)
				..controls(41.3,-100.8)and(42.4,-101.8)..(43.7,-102.9)
				..controls(44,-103.2)and(44.6,-103.7)..(45,-103.9)
				..controls(45.4,-104.2)and(45.9,-104.6)..(46.1,-104.8)
				..controls(47.1,-105.9)and(48.8,-107.1)..(50.5,-107.9)
				..controls(51.8,-108.6)and(52.3,-109)..(53.4,-110)
				..controls(54.8,-111.2)and(55.6,-111.6)..(57,-112.1)
				..controls(57.5,-112.3)and(58.5,-112.8)..(59.4,-113.3)
				..controls(60.2,-113.7)and(61.3,-114.3)..(61.9,-114.5)
				..controls(62.4,-114.7)and(63.3,-115.1)..(63.8,-115.3)
				..controls(64.4,-115.6)and(65.4,-115.9)..(66.2,-116.1)
				..controls(67,-116.3)and(68.4,-116.6)..(69.2,-116.9)
				..controls(72.3,-117.9)and(72.8,-118)..(78,-118.1)
				..controls(84,-118.3)and(86.1,-118.1)..(88.1,-117.1)
				..controls(89.4,-116.5)and(93,-115)..(93.7,-114.9)
				..controls(94.8,-114.6)and(95.4,-114.2)..(95.9,-113.4)
				..controls(96.2,-113)and(96.7,-112.5)..(97.1,-112.2)
				..controls(98,-111.6)and(101.3,-108.4)..(102.1,-107.2)
				..controls(102.4,-106.8)and(102.9,-106.2)..(103.2,-105.8)
				..controls(103.5,-105.5)and(104,-104.8)..(104.2,-104.2)
				..controls(104.5,-103.6)and(105,-102.8)..(105.3,-102.4)
				..controls(105.6,-102)and(106.1,-101.2)..(106.5,-100.6)
				..controls(107.1,-99.5)and(108.6,-97.5)..(109.5,-96.5)
				..controls(109.9,-96.1)and(110.1,-96)..(110.4,-96)
				..controls(111.1,-96)and(111.6,-96.5)..(111.8,-97.6)
				..controls(112,-98.6)and(112.3,-99.3)..(112.8,-99.9)
				..controls(113.7,-100.9)and(115.2,-103.1)..(115.7,-104.2)
				..controls(116.1,-105)and(116.5,-105.5)..(116.9,-106)
				..controls(117.3,-106.4)and(117.8,-106.9)..(118,-107.3)
				..controls(118.2,-107.6)and(118.6,-108)..(119,-108.3)
				..controls(119.6,-108.8)and(121.1,-110.3)..(121.8,-111.2)
				..controls(122.1,-111.7)and(122.5,-112)..(123.4,-112.4)
				..controls(124.8,-113.2)and(125.8,-113.8)..(127,-114.6)
				..controls(127.5,-114.9)and(128.1,-115.3)..(128.2,-115.4)
				..controls(128.4,-115.5)and(129.3,-115.8)..(130.2,-116)
				..controls(131.1,-116.2)and(132.5,-116.6)..(133.3,-116.9)
				..controls(135.7,-117.9)and(136.1,-117.9)..(142.5,-117.8)
				..controls(145.8,-117.8)and(148.4,-117.7)..(148.8,-117.6)
				..controls(149.2,-117.6)and(150.2,-117.2)..(151,-116.8)
				..controls(152,-116.4)and(152.8,-116)..(153.6,-115.9)
				..controls(154.5,-115.7)and(155.2,-115.4)..(156.3,-114.9)
				..controls(157.1,-114.5)and(158.2,-114)..(158.7,-113.8)
				..controls(159.1,-113.7)and(159.8,-113.3)..(160.2,-113)
				..controls(160.6,-112.7)and(161.3,-112.4)..(161.7,-112.2)
				..controls(162.7,-111.9)and(163.2,-111.7)..(164.1,-111.1)
				..controls(164.5,-110.8)and(165.3,-110.4)..(165.9,-110)
				..controls(166.5,-109.7)and(167.3,-109.2)..(167.6,-108.9)
				..controls(167.9,-108.6)and(168.6,-108.2)..(169,-108)
				..controls(169.4,-107.8)and(170.2,-107.3)..(170.7,-106.9)
				..controls(171.2,-106.5)and(172.1,-105.8)..(172.7,-105.4)
				..controls(173.2,-105)and(174.1,-104.4)..(174.5,-104)
				..controls(174.9,-103.7)and(175.7,-103.1)..(176.2,-102.6)
				..controls(176.7,-102.2)and(177.4,-101.5)..(177.9,-101.2)
				..controls(178.3,-100.8)and(178.8,-100.4)..(179,-100.2)
				..controls(179.2,-100)and(179.6,-99.6)..(180,-99.3)
				..controls(180.4,-99.1)and(180.9,-98.6)..(181.2,-98.4)
				..controls(181.5,-98.1)and(182.1,-97.7)..(182.5,-97.5)
				..controls(183.2,-97.2)and(183.5,-96.9)..(184.3,-95.9)
				..controls(184.6,-95.6)and(184.6,-95.6)..(184.8,-95.7)
				..controls(185,-95.8)and(185.1,-96)..(185.1,-96.2)
				..controls(185.1,-96.5)and(185,-96.7)..(184.6,-97.2)
				..controls(184.1,-97.8)and(184,-97.9)..(183.9,-98.8)
				..controls(183.7,-100.3)and(183.3,-101.3)..(182.9,-101.8)
				..controls(182.6,-102.1)and(182.3,-102.6)..(182.2,-102.9)
				..controls(182,-103.3)and(181.6,-104)..(181.2,-104.6)
				..controls(180.9,-105.2)and(180.5,-106.1)..(180.4,-106.6)
				..controls(179.9,-108)and(179.7,-108.5)..(179,-109.4)
				..controls(178.6,-109.8)and(178.2,-110.5)..(177.9,-111)
				..controls(177.7,-111.5)and(177.3,-112.3)..(177,-112.7)
				..controls(176.6,-113.2)and(176.4,-113.7)..(176.3,-114.2)
				..controls(176.2,-114.7)and(175.9,-115.2)..(175.4,-116)
				..controls(175,-116.6)and(174.4,-117.7)..(174,-118.5)
				..controls(173.7,-119.2)and(173.3,-120.1)..(173,-120.5)
				..controls(172.8,-120.9)and(172.4,-121.8)..(172.1,-122.5)
				..controls(171.8,-123.2)and(171.3,-124.1)..(171,-124.6)
				..controls(170.8,-125.1)and(170.3,-126)..(169.9,-126.6)
				..controls(169.6,-127.3)and(169.2,-128.1)..(169,-128.3)
				..controls(168.8,-128.6)and(168.5,-129.3)..(168.3,-129.8)
				..controls(168.1,-130.3)and(167.8,-131.1)..(167.5,-131.6)
				..controls(166.9,-132.5)and(164.2,-137.9)..(164.2,-138.1)
				..controls(164.2,-138.1)and(163.9,-138.8)..(163.4,-139.7)
				..controls(162.9,-140.5)and(162.4,-141.5)..(162.3,-141.9)
				..controls(161.9,-143)and(161.6,-143.4)..(160.9,-144.1)
				..controls(160.4,-144.6)and(160.3,-144.8)..(160.1,-145.4)
				..controls(160,-145.8)and(159.6,-146.8)..(159.2,-147.8)
				..controls(158.7,-148.8)and(158.1,-150)..(157.9,-150.6)
				..controls(157.6,-151.1)and(157.2,-151.9)..(156.9,-152.4)
				..controls(156.5,-152.8)and(156.1,-153.7)..(155.8,-154.4)
				..controls(155.6,-155)and(155.2,-155.9)..(154.9,-156.3)
				..controls(154.7,-156.7)and(154.4,-157.4)..(154.2,-157.9)
				..controls(154.1,-158.4)and(153.8,-159)..(153.5,-159.3)
				..controls(152.8,-160.1)and(153,-160.1)..(142.4,-160.1)
				..controls(137.2,-160.1)and(131.8,-160.1)..(130.4,-160)--cycle
				(141.5,-103.7)..controls(140.8,-103)and(139.8,-102.1)..(139.2,-101.6)
				..controls(138.5,-101.2)and(137.9,-100.6)..(137.8,-100.4)
				..controls(137,-99.5)and(134,-96.5)..(133.2,-96)
				..controls(132.7,-95.7)and(131.8,-94.9)..(131.1,-94.3)
				..controls(129.4,-92.7)and(127.9,-91.6)..(126.6,-91.1)
				..controls(125.2,-90.5)and(124.1,-90)..(122.9,-89.1)
				..controls(122.1,-88.6)and(121.3,-88.2)..(119.9,-87.7)
				..controls(118.8,-87.3)and(117.5,-86.8)..(116.9,-86.5)
				..controls(114.3,-85.5)and(113.7,-85.4)..(109.7,-85.5)
				..controls(105.3,-85.7)and(103.9,-85.9)..(101.6,-87)
				..controls(100.7,-87.4)and(99.4,-87.8)..(97.9,-88)
				..controls(97,-88.2)and(96.6,-88.3)..(95.8,-88.9)
				..controls(95.4,-89.2)and(94.7,-89.6)..(94.2,-89.8)
				..controls(92.9,-90.3)and(92.3,-90.7)..(90.9,-92.1)
				..controls(90.2,-92.7)and(89.4,-93.3)..(89,-93.5)..controls(87,-94.5)and(85.2,-96)..(83.6,-97.9)
				..controls(83.3,-98.2)and(82.5,-98.9)..(81.8,-99.4)
				..controls(80.7,-100.1)and(80.5,-100.3)..(80,-101.1)
				..controls(78.7,-103.1)and(77.2,-103.9)..(76.4,-103)
				..controls(75.7,-102.2)and(75.7,-103.2)..(75.7,-83.6)
				..controls(75.6,-73.6)and(75.6,-64.8)..(75.5,-64)
				..controls(75.5,-63.2)and(75.5,-61.4)..(75.6,-59.8)--(75.7,-57.1)--(76,-56.8)
				..controls(76.8,-56.2)and(78,-56.3)..(79.8,-57.2)
				..controls(80.4,-57.5)and(81.3,-57.8)..(81.8,-57.9)
				..controls(82.3,-58.1)and(83,-58.4)..(83.6,-58.8)..controls(84.3,-59.2)and(85,-59.5)..(86,-59.8)
				..controls(86.8,-60.1)and(88,-60.6)..(88.7,-60.9)
				..controls(89.3,-61.2)and(90.6,-61.7)..(91.4,-62)
				..controls(92.2,-62.3)and(93.7,-62.9)..(94.6,-63.3)
				..controls(95.6,-63.7)and(96.8,-64.1)..(97.4,-64.3)
				..controls(98,-64.5)and(99,-64.9)..(99.7,-65.2)
				..controls(100.4,-65.6)and(101.5,-66)..(102.3,-66.2)
				..controls(104.2,-66.8)and(105,-67.1)..(106.2,-67.8)
				..controls(107.9,-68.7)and(110.2,-69.4)..(111.3,-69.2)
				..controls(112.5,-69)and(115,-68.2)..(115.6,-67.8)
				..controls(117.3,-66.7)and(117.6,-66.6)..(119.5,-66.2)
				..controls(121.9,-65.8)and(122.1,-65.7)..(123,-65.1)
				..controls(124.1,-64.4)and(124.4,-64.3)..(126,-63.9)
				..controls(127,-63.7)and(127.7,-63.4)..(128.4,-63.1)
				..controls(129,-62.8)and(129.7,-62.5)..(130.6,-62.4)
				..controls(132.1,-62.1)and(134,-61.5)..(135.2,-61)
				..controls(135.7,-60.8)and(136.7,-60.5)..(137.3,-60.3)
				..controls(138.4,-60)and(139.3,-59.7)..(141.1,-59)
				..controls(141.6,-58.7)and(142,-58.6)..(142.5,-58.6)
				..controls(143.2,-58.6)and(143.3,-58.7)..(143.6,-59)
				..controls(144,-59.5)and(144.1,-59.9)..(144,-60.8)
				..controls(144,-61.2)and(144,-71.1)..(144,-82.8)--(143.9,-104.1)--(143.7,-104.5)
				..controls(143.5,-104.7)and(143.3,-104.8)..(143,-104.9)
				..controls(142.6,-104.9)and(142.5,-104.8)..(141.5,-103.7)--cycle
				(27.8,-91.4)..controls(27.1,-90.9)and(26.3,-90)..(26.1,-89.5)
				..controls(26,-89.2)and(25.8,-88.6)..(25.7,-87.9)
				..controls(25.4,-86.2)and(25.6,-84.8)..(26.5,-83.5)
				..controls(27.4,-82.3)and(28.3,-82)..(29.4,-82.3)
				..controls(31.1,-82.8)and(32.1,-84.9)..(32,-87.4)
				..controls(31.9,-89.3)and(31.2,-90.6)..(29.9,-91.4)
				..controls(29.1,-91.9)and(28.6,-91.9)..(27.8,-91.4)--cycle
				(190.1,-91.3)..controls(189.7,-91.2)and(189.3,-90.9)..(189.1,-90.7)
				..controls(188.3,-89.8)and(187.6,-87.7)..(187.6,-86.2)
				..controls(187.6,-85.5)and(187.7,-85.2)..(188.2,-84.4)
				..controls(189.1,-82.6)and(190.1,-81.8)..(191.3,-82.1)
				..controls(192,-82.3)and(193.5,-83.7)..(193.8,-84.4)
				..controls(194.1,-85.1)and(194.2,-87.1)..(194,-88.3)
				..controls(193.8,-89.8)and(193.3,-90.5)..(192.3,-91.3)
				..controls(191.7,-91.7)and(190.9,-91.7)..(190.1,-91.3)--cycle
				(173.7,-83.4)..controls(172.2,-82.5)and(172,-82.5)..(170.6,-82.4)
				..controls(169.7,-82.4)and(168.8,-82.4)..(168.3,-82.5)
				..controls(167.2,-82.7)and(166.9,-82.7)..(166.6,-82.2)
				..controls(166.3,-81.8)and(166.3,-81)..(166.8,-80)
				..controls(167.2,-78.9)and(167.4,-77.4)..(167.7,-74.1)
				..controls(167.8,-72.8)and(167.9,-71.8)..(167.9,-71.7)
				..controls(168,-71.7)and(168,-70.2)..(168,-68.4)..controls(168,-65.4)and(168,-64.9)..(167.7,-63)
				..controls(167.3,-60.5)and(167.1,-59.9)..(166.4,-58.4)
				..controls(166.1,-57.7)and(165.8,-56.9)..(165.7,-56.6)
				..controls(165.7,-56.2)and(165.4,-55.6)..(165.2,-55.2)
				..controls(164.6,-53.8)and(161.9,-50.7)..(161,-50.3)
				..controls(159.8,-49.7)and(157.9,-50)..(155.2,-51.2)
				..controls(154.7,-51.4)and(153.4,-51.8)..(152.2,-52.1)
				..controls(150.8,-52.4)and(149.4,-52.9)..(148.6,-53.2)
				..controls(147.8,-53.6)and(146.7,-54)..(146,-54.2)
				..controls(145.4,-54.4)and(144.5,-54.7)..(144,-54.9)
				..controls(143.6,-55.2)and(142.7,-55.5)..(141.8,-55.7)
				..controls(141,-55.9)and(139.6,-56.3)..(138.7,-56.6)
				..controls(137.8,-57)and(136.3,-57.4)..(135.3,-57.6)
				..controls(133.8,-57.9)and(133.2,-58.1)..(132.1,-58.7)
				..controls(131.1,-59.2)and(130.3,-59.4)..(129.4,-59.6)
				..controls(128.7,-59.8)and(127.8,-60)..(127.4,-60.1)
				..controls(127,-60.3)and(126,-60.5)..(125.1,-60.7)
				..controls(123.2,-61.1)and(121.7,-61.5)..(120.2,-62.1)
				..controls(119.7,-62.3)and(118.1,-62.8)..(116.7,-63.2)
				..controls(114.6,-63.9)and(114,-64.1)..(113,-64.7)
				..controls(111.9,-65.4)and(111.8,-65.4)..(111,-65.5)
				..controls(109.9,-65.5)and(109.3,-65.4)..(107.7,-64.7)
				..controls(107,-64.4)and(106,-64)..(105.3,-63.9)
				..controls(103.5,-63.5)and(102.2,-63)..(101.2,-62.5)
				..controls(100.1,-61.8)and(99.3,-61.5)..(97.8,-61.1)
				..controls(97.2,-61)and(96.2,-60.7)..(95.7,-60.4)
				..controls(95.1,-60.1)and(94.2,-59.8)..(93.5,-59.6)
				..controls(92.9,-59.4)and(91.7,-59)..(91,-58.7)..controls(90.2,-58.4)and(88.9,-58)..(88,-57.7)
				..controls(87.2,-57.5)and(86.1,-57.1)..(85.5,-56.8)
				..controls(85,-56.6)and(84.1,-56.2)..(83.5,-56)..controls(82.9,-55.7)and(81.8,-55.2)..(81,-54.9)
				..controls(80.1,-54.5)and(79.1,-54.2)..(78.7,-54.1)
				..controls(77.9,-53.8)and(77.3,-53.6)..(76.2,-52.9)
				..controls(75.7,-52.6)and(75.2,-52.4)..(74.1,-52.2)
				..controls(72.5,-51.9)and(72.4,-51.8)..(71.3,-51.1)
				..controls(69.9,-50.1)and(69.4,-49.9)..(67.9,-49.8)
				..controls(67,-49.7)and(66.6,-49.6)..(66.3,-49.4)
				..controls(66.1,-49.3)and(65.7,-49)..(65.4,-48.8)
				..controls(65.1,-48.7)and(64.2,-48.4)..(63.6,-48.2)
				..controls(62.9,-48)and(61.7,-47.6)..(60.9,-47.2)
				..controls(60.1,-46.8)and(58.9,-46.4)..(58.2,-46.2)
				..controls(57.6,-46)and(56.8,-45.7)..(56.5,-45.5)
				..controls(54.9,-44.6)and(54.6,-44.5)..(53.2,-44.2)
				..controls(51.7,-43.8)and(51,-43.6)..(49.9,-42.9)
				..controls(49.5,-42.6)and(48.9,-42.4)..(47.5,-42)
				..controls(45.7,-41.6)and(45.6,-41.5)..(44.3,-40.6)
				..controls(43.5,-40.1)and(42.8,-39.6)..(42.7,-39.4)
				..controls(42.4,-38.7)and(43,-38.1)..(44.5,-37.6)..controls(44.9,-37.5)and(46,-37)..(46.9,-36.6)
				..controls(47.9,-36.1)and(48.8,-35.8)..(49.6,-35.6)
				..controls(50.5,-35.4)and(51.3,-35.2)..(52.4,-34.6)
				..controls(53.3,-34.1)and(54.5,-33.7)..(54.9,-33.5)
				..controls(55.4,-33.4)and(56.4,-33)..(57.2,-32.6)
				..controls(58.6,-32)and(60.1,-31.5)..(61.6,-31.2)
				..controls(62.3,-31.1)and(62.7,-31)..(63.2,-30.7)
				..controls(64.1,-30)and(64.3,-29.9)..(65.3,-29.7)
				..controls(65.7,-29.6)and(66.5,-29.2)..(67.1,-28.9)
				..controls(67.6,-28.6)and(68.8,-28.2)..(69.6,-27.9)
				..controls(70.5,-27.7)and(71.6,-27.3)..(72.1,-27)
				..controls(72.6,-26.7)and(73.8,-26.2)..(74.8,-25.9)
				..controls(75.7,-25.5)and(76.7,-25.1)..(77.1,-25)
				..controls(77.4,-24.8)and(78.5,-24.4)..(79.6,-24.1)
				..controls(80.6,-23.8)and(82,-23.2)..(82.8,-22.9)
				..controls(83.5,-22.5)and(84.8,-22)..(85.8,-21.8)
				..controls(86.7,-21.5)and(87.8,-21.1)..(88.1,-20.9)
				..controls(88.4,-20.7)and(89.3,-20.3)..(90.1,-20.1)
				..controls(90.8,-19.8)and(91.9,-19.4)..(92.4,-19.1)
				..controls(93.5,-18.5)and(94.4,-18.2)..(95.4,-18.1)
				..controls(95.7,-18)and(96.4,-17.8)..(96.9,-17.5)
				..controls(97.5,-17.3)and(98.4,-16.9)..(99,-16.6)
				..controls(99.7,-16.3)and(100.9,-15.9)..(101.9,-15.6)
				..controls(102.9,-15.3)and(104.4,-14.7)..(105.3,-14.3)
				..controls(108.1,-13.2)and(108.5,-13.1)..(110.1,-13.1)
				..controls(111.4,-13.1)and(111.5,-13.1)..(112.4,-13.5)
				..controls(113.2,-14)and(115,-14.8)..(116.6,-15.4)
				..controls(116.9,-15.5)and(117.9,-15.8)..(118.8,-16.1)
				..controls(120.9,-16.6)and(121.5,-16.9)..(122.6,-17.9)
				..controls(123,-18.3)and(123.7,-18.8)..(123.9,-18.9)
				..controls(124.6,-19.4)and(126,-19.8)..(127.9,-20.2)
				..controls(129.6,-20.5)and(129.8,-20.6)..(131.1,-21.4)
				..controls(131.5,-21.7)and(132.4,-22.1)..(133.1,-22.3)
				..controls(136.1,-23.2)and(138.3,-24.1)..(140.2,-25)
				..controls(141,-25.5)and(141.5,-25.6)..(142.5,-25.8)
				..controls(145.1,-26.3)and(146,-26.6)..(149.4,-28.1)
				..controls(151.3,-28.9)and(153.5,-29.8)..(154.3,-30.1)
				..controls(155.1,-30.4)and(156.2,-30.8)..(156.6,-31)
				..controls(157.1,-31.2)and(158.2,-31.6)..(159.1,-31.8)
				..controls(160.6,-32.2)and(160.9,-32.3)..(161.7,-32.9)
				..controls(162.4,-33.3)and(162.9,-33.6)..(163.7,-33.8)
				..controls(164.3,-34)and(164.8,-34.1)..(164.9,-34.1)
				..controls(165,-34.1)and(165.6,-34.4)..(166.4,-34.8)
				..controls(167.3,-35.2)and(168.1,-35.4)..(169,-35.6)
				..controls(170,-35.8)and(170.5,-36)..(171.7,-36.6)
				..controls(172.4,-37)and(173.3,-37.4)..(173.6,-37.4)
				..controls(175.6,-37.7)and(176.2,-38)..(176.7,-39.1)
				..controls(176.9,-39.6)and(176.9,-39.7)..(176.7,-40.4)
				..controls(176.4,-41.6)and(175.8,-42.7)..(175.2,-43.4)
				..controls(174.5,-44.1)and(172.6,-45.5)..(171.7,-45.8)
				..controls(170.6,-46.2)and(169.1,-46)..(168.2,-45.4)
				..controls(167.9,-45.2)and(167.3,-44.8)..(166.9,-44.4)
				..controls(166.4,-44.1)and(165.8,-43.6)..(165.4,-43.2)
				..controls(165.1,-42.9)and(164.4,-42.3)..(163.9,-41.9)
				..controls(163.3,-41.5)and(162.7,-40.9)..(162.4,-40.5)
				..controls(161.4,-39.2)and(159.9,-38)..(158.8,-37.8)
				..controls(158.4,-37.8)and(158.3,-37.7)..(158.2,-37.2)
				..controls(158,-36.8)and(157.4,-36.1)..(157.1,-36.1)
				..controls(157,-36.1)and(156.7,-36.1)..(156.5,-36)
				..controls(156.1,-35.9)and(155.9,-35.7)..(155.7,-35.2)
				..controls(155.3,-34.6)and(155.1,-34.4)..(154,-34.1)
				..controls(153.6,-34)and(153.2,-33.7)..(152.9,-33.5)
				..controls(152.4,-32.9)and(151.6,-32.4)..(151.3,-32.4)
				..controls(150.9,-32.4)and(150.4,-32.6)..(150.4,-32.9)
				..controls(150.4,-33.2)and(150.8,-33.4)..(151.4,-33.6)
				..controls(152.1,-33.7)and(152.8,-34.1)..(153.5,-34.9)
				..controls(154.1,-35.5)and(154.3,-35.6)..(155,-35.8)
				..controls(155.7,-36)and(155.7,-36)..(156,-36.6)
				..controls(156.3,-37.3)and(156.7,-37.7)..(157.5,-37.8)
				..controls(158.1,-37.9)and(158.2,-38)..(158.4,-38.8)
				..controls(158.5,-39.4)and(158.6,-39.6)..(159.6,-40.4)
				..controls(160.2,-41)and(160.9,-41.6)..(161.3,-41.8)
				..controls(161.7,-42)and(161.9,-42.3)..(162.2,-42.8)
				..controls(162.4,-43.2)and(162.8,-43.8)..(163.1,-44.1)
				..controls(163.9,-44.8)and(165.1,-46.7)..(165.7,-48)
				..controls(166.2,-49.3)and(166.4,-49.7)..(167.2,-50.5)
				..controls(167.5,-50.8)and(167.9,-51.4)..(168.1,-51.8)
				..controls(168.3,-52.2)and(168.8,-53)..(169.1,-53.6)
				..controls(169.6,-54.3)and(169.8,-54.9)..(169.9,-55.6)
				..controls(170.4,-57.8)and(170.5,-58.1)..(171.2,-59.5)
				..controls(171.6,-60.3)and(172,-61.5)..(172.2,-62.1)
				..controls(172.3,-62.8)and(172.7,-63.7)..(172.9,-64.1)
				..controls(173.2,-64.6)and(173.4,-65.3)..(173.4,-65.6)
				..controls(173.8,-68.8)and(174.2,-70.6)..(174.9,-72.6)
				..controls(175.5,-74.5)and(175.9,-77.1)..(176.1,-80.2)
				..controls(176.3,-83.4)and(175.5,-84.4)..(173.7,-83.4)--cycle;
				\definecolor{ezs}{RGB}{247,179,26}\fill[ezs](-5.7,-156.6)--(-3.4,-157.5)--(61.9,12.7)--(59.7,13.6)--cycle;
				\definecolor{ezs}{RGB}{247,179,22}\fill[ezs](-3.8,-157.3)--(-1.5,-158.2)--(63.9,12)--(61.6,12.9)--cycle;
				\definecolor{ezs}{RGB}{248,179,19}\fill[ezs](-1.8,-158.1)--(0.4,-158.9)--(65.8,11.2)--(63.5,12.1)--cycle;
				\definecolor{ezs}{RGB}{248,178,16}\fill[ezs](0.1,-158.8)--(2.4,-159.7)--(67.7,10.5)--(65.4,11.4)--cycle;
				\definecolor{ezs}{RGB}{249,178,15}\fill[ezs](2,-159.5)--(4.3,-160.4)--(69.6,9.8)--(67.3,10.7)--cycle;
				\definecolor{ezs}{RGB}{251,177,17}\fill[ezs](3.9,-160.3)--(6.2,-161.1)--(71.5,9)--(69.2,9.9)--cycle;
				\definecolor{ezs}{RGB}{252,179,18}\fill[ezs](5.8,-161)--(8.1,-161.9)--(73.4,8.3)--(71.1,9.2)--cycle;
				\definecolor{ezs}{RGB}{252,182,17}\fill[ezs](7.7,-161.7)--(10,-162.6)--(75.3,7.6)--(73,8.5)--cycle;
				\definecolor{ezs}{RGB}{252,185,17}\fill[ezs](9.6,-162.5)--(11.9,-163.4)--(77.2,6.8)--(75,7.7)--cycle;
				\definecolor{ezs}{RGB}{251,189,18}\fill[ezs](11.5,-163.2)--(13.8,-164.1)--(79.2,6.1)--(76.9,7)--cycle;
				\definecolor{ezs}{RGB}{251,191,22}\fill[ezs](13.4,-163.9)--(15.7,-164.8)--(81.1,5.4)--(78.8,6.2)--cycle;
				\definecolor{ezs}{RGB}{252,193,25}\fill[ezs](15.4,-164.7)--(17.6,-165.6)--(83,4.6)--(80.7,5.5)--cycle;
				\definecolor{ezs}{RGB}{253,195,27}\fill[ezs](17.3,-165.4)--(19.6,-166.3)--(84.9,3.9)--(82.6,4.8)--cycle;
				\definecolor{ezs}{RGB}{253,196,31}\fill[ezs](19.2,-166.1)--(21.5,-167)--(86.8,3.2)--(84.5,4)--cycle;
				\definecolor{ezs}{RGB}{254,199,35}\fill[ezs](21.1,-166.9)--(23.4,-167.8)--(88.7,2.4)--(86.4,3.3)--cycle;
				\definecolor{ezs}{RGB}{253,201,38}\fill[ezs](23,-167.6)--(25.3,-168.5)--(90.6,1.7)--(88.3,2.6)--cycle;
				\definecolor{ezs}{RGB}{253,203,40}\fill[ezs](24.9,-168.3)--(27.2,-169.2)--(92.5,1)--(90.3,1.8)--cycle;
				\definecolor{ezs}{RGB}{253,205,42}\fill[ezs](26.8,-169.1)--(29.1,-170)--(94.4,0.2)--(92.2,1.1)--cycle;
				\definecolor{ezs}{RGB}{253,207,45}\fill[ezs](28.7,-169.8)--(31,-170.7)--(96.4,-0.5)--(94.1,0.4)--cycle;
				\definecolor{ezs}{RGB}{253,210,49}\fill[ezs](30.7,-170.5)--(32.9,-171.4)--(98.3,-1.2)--(96,-0.4)--cycle;
				\definecolor{ezs}{RGB}{253,212,54}\fill[ezs](32.6,-171.3)--(34.9,-172.2)--(100.2,-2)--(97.9,-1.1)--cycle;
				\definecolor{ezs}{RGB}{253,214,59}\fill[ezs](34.5,-172)--(36.8,-172.9)--(102.1,-2.7)--(99.8,-1.8)--cycle;
				\definecolor{ezs}{RGB}{253,216,63}\fill[ezs](36.4,-172.7)--(38.7,-173.6)--(104,-3.4)--(101.7,-2.6)--cycle;
				\definecolor{ezs}{RGB}{253,218,65}\fill[ezs](38.3,-173.5)--(40.6,-174.4)--(105.9,-4.2)--(103.6,-3.3)--cycle;
				\definecolor{ezs}{RGB}{254,220,66}\fill[ezs](40.2,-174.2)--(42.5,-175.1)--(107.8,-4.9)--(105.5,-4)--cycle;
				\definecolor{ezs}{RGB}{254,222,68}\fill[ezs](42.1,-175)--(44.4,-175.8)--(109.7,-5.6)--(107.5,-4.8)--cycle;
				\definecolor{ezs}{RGB}{254,224,71}\fill[ezs](44,-175.7)--(46.3,-176.6)--(111.7,-6.4)--(109.4,-5.5)--cycle;
				\definecolor{ezs}{RGB}{254,226,75}\fill[ezs](46,-176.4)--(48.2,-177.3)--(113.6,-7.1)--(111.3,-6.2)--cycle;
				\definecolor{ezs}{RGB}{254,228,79}\fill[ezs](47.9,-177.2)--(50.1,-178)--(115.5,-7.8)--(113.2,-7)--cycle;
				\definecolor{ezs}{RGB}{254,230,82}\fill[ezs](49.8,-177.9)--(52.1,-178.8)--(117.4,-8.6)--(115.1,-7.7)--cycle;
				\definecolor{ezs}{RGB}{254,232,85}\fill[ezs](51.7,-178.6)--(54,-179.5)--(119.3,-9.3)--(117,-8.4)--cycle;
				\definecolor{ezs}{RGB}{254,234,88}\fill[ezs](53.6,-179.4)--(55.9,-180.2)--(121.2,-10)--(118.9,-9.2)--cycle;
				\definecolor{ezs}{RGB}{254,237,91}\fill[ezs](55.5,-180.1)--(57.8,-181)--(123.1,-10.8)--(120.8,-9.9)--cycle;
				\definecolor{ezs}{RGB}{254,239,94}\fill[ezs](57.4,-180.8)--(59.7,-181.7)--(125,-11.5)--(122.8,-10.6)--cycle;
				\definecolor{ezs}{RGB}{254,241,97}\fill[ezs](59.3,-181.6)--(61.6,-182.4)--(127,-12.2)--(124.7,-11.4)--cycle;
				\definecolor{ezs}{RGB}{254,243,101}\fill[ezs](61.2,-182.3)--(63.5,-183.2)--(128.9,-13)--(126.6,-12.1)--cycle;
				\definecolor{ezs}{RGB}{254,245,103}\fill[ezs](63.2,-183)--(65.4,-183.9)--(130.8,-13.7)--(128.5,-12.8)--cycle;
				\definecolor{ezs}{RGB}{254,245,105}\fill[ezs](65.1,-183.8)--(67.4,-184.6)--(132.7,-14.4)--(130.4,-13.6)--cycle;
				\definecolor{ezs}{RGB}{254,245,105}\fill[ezs](67,-184.5)--(69.3,-185.4)--(134.6,-15.2)--(132.3,-14.3)--cycle;
				\definecolor{ezs}{RGB}{254,244,103}\fill[ezs](68.9,-185.2)--(71.2,-186.1)--(136.5,-15.9)--(134.2,-15)--cycle;
				\definecolor{ezs}{RGB}{253,243,101}\fill[ezs](70.8,-186)--(73.1,-186.8)--(138.4,-16.6)--(136.1,-15.8)--cycle;
				\definecolor{ezs}{RGB}{253,241,98}\fill[ezs](72.7,-186.7)--(75,-187.6)--(140.3,-17.4)--(138,-16.5)--cycle;
				\definecolor{ezs}{RGB}{253,239,95}\fill[ezs](74.6,-187.4)--(76.9,-188.3)--(142.2,-18.1)--(140,-17.2)--cycle;
				\definecolor{ezs}{RGB}{253,236,92}\fill[ezs](76.5,-188.2)--(78.8,-189)--(144.2,-18.8)--(141.9,-18)--cycle;
				\definecolor{ezs}{RGB}{253,234,88}\fill[ezs](78.5,-188.9)--(80.7,-189.8)--(146.1,-19.6)--(143.8,-18.7)--cycle;
				\definecolor{ezs}{RGB}{253,232,85}\fill[ezs](80.4,-189.6)--(82.7,-190.5)--(148,-20.3)--(145.7,-19.4)--cycle;
				\definecolor{ezs}{RGB}{253,230,82}\fill[ezs](82.3,-190.4)--(84.6,-191.2)--(149.9,-21)--(147.6,-20.2)--cycle;
				\definecolor{ezs}{RGB}{253,228,78}\fill[ezs](84.2,-191.1)--(86.5,-192)--(151.8,-21.8)--(149.5,-20.9)--cycle;
				\definecolor{ezs}{RGB}{253,226,74}\fill[ezs](86.1,-191.8)--(88.4,-192.7)--(153.7,-22.5)--(151.4,-21.6)--cycle;
				\definecolor{ezs}{RGB}{253,224,70}\fill[ezs](88,-192.6)--(90.3,-193.4)--(155.6,-23.3)--(153.3,-22.4)--cycle;
				\definecolor{ezs}{RGB}{253,222,67}\fill[ezs](89.9,-193.3)--(92.2,-194.2)--(157.5,-24)--(155.3,-23.1)--cycle;
				\definecolor{ezs}{RGB}{253,220,65}\fill[ezs](91.8,-194)--(94.1,-194.9)--(159.5,-24.7)--(157.2,-23.8)--cycle;
				\definecolor{ezs}{RGB}{253,217,63}\fill[ezs](93.7,-194.8)--(96,-195.6)--(161.4,-25.5)--(159.1,-24.6)--cycle;
				\definecolor{ezs}{RGB}{253,215,60}\fill[ezs](95.7,-195.5)--(97.9,-196.4)--(163.3,-26.2)--(161,-25.3)--cycle;
				\definecolor{ezs}{RGB}{253,213,56}\fill[ezs](97.6,-196.2)--(99.9,-197.1)--(165.2,-26.9)--(162.9,-26)--cycle;
				\definecolor{ezs}{RGB}{252,211,53}\fill[ezs](99.5,-197)--(101.8,-197.8)--(167.1,-27.7)--(164.8,-26.8)--cycle;
				\definecolor{ezs}{RGB}{252,209,49}\fill[ezs](101.4,-197.7)--(103.7,-198.6)--(169,-28.4)--(166.7,-27.5)--cycle;
				\definecolor{ezs}{RGB}{252,207,45}\fill[ezs](103.3,-198.4)--(105.6,-199.3)--(170.9,-29.1)--(168.6,-28.2)--cycle;
				\definecolor{ezs}{RGB}{253,205,42}\fill[ezs](105.2,-199.2)--(107.5,-200)--(172.8,-29.9)--(170.6,-29)--cycle;
				\definecolor{ezs}{RGB}{253,202,40}\fill[ezs](107.1,-199.9)--(109.4,-200.8)--(174.7,-30.6)--(172.5,-29.7)--cycle;
				\definecolor{ezs}{RGB}{253,200,38}\fill[ezs](109,-200.6)--(111.3,-201.5)--(176.7,-31.3)--(174.4,-30.4)--cycle;
				\definecolor{ezs}{RGB}{253,198,35}\fill[ezs](111,-201.4)--(113.2,-202.2)--(178.6,-32.1)--(176.3,-31.2)--cycle;
				\definecolor{ezs}{RGB}{253,196,31}\fill[ezs](112.9,-202.1)--(115.2,-203)--(180.5,-32.8)--(178.2,-31.9)--cycle;
				\definecolor{ezs}{RGB}{253,194,28}\fill[ezs](114.8,-202.8)--(117.1,-203.7)--(182.4,-33.5)--(180.1,-32.6)--cycle;
				\definecolor{ezs}{RGB}{253,192,26}\fill[ezs](116.7,-203.6)--(119,-204.5)--(184.3,-34.3)--(182,-33.4)--cycle;
				\definecolor{ezs}{RGB}{253,189,24}\fill[ezs](118.6,-204.3)--(120.9,-205.2)--(186.2,-35)--(183.9,-34.1)--cycle;
				\definecolor{ezs}{RGB}{253,186,20}\fill[ezs](120.5,-205)--(122.8,-205.9)--(188.1,-35.7)--(185.8,-34.8)--cycle;
				\definecolor{ezs}{RGB}{253,184,17}\fill[ezs](122.4,-205.8)--(124.7,-206.7)--(190,-36.5)--(187.8,-35.6)--cycle;
				\definecolor{ezs}{RGB}{252,183,15}\fill[ezs](124.3,-206.5)--(126.6,-207.4)--(192,-37.2)--(189.7,-36.3)--cycle;
				\definecolor{ezs}{RGB}{252,181,14}\fill[ezs](126.3,-207.2)--(128.5,-208.1)--(193.9,-37.9)--(191.6,-37.1)--cycle;
				\definecolor{ezs}{RGB}{252,179,13}\fill[ezs](128.2,-208)--(130.4,-208.9)--(195.8,-38.7)--(193.5,-37.8)--cycle;
				\definecolor{ezs}{RGB}{252,178,12}\fill[ezs](130.1,-208.7)--(132.4,-209.6)--(197.7,-39.4)--(195.4,-38.5)--cycle;
				\definecolor{ezs}{RGB}{251,178,10}\fill[ezs](132,-209.4)--(134.3,-210.3)--(199.6,-40.1)--(197.3,-39.3)--cycle;
				\definecolor{ezs}{RGB}{251,178,10}\fill[ezs](133.9,-210.2)--(136.2,-211.1)--(201.5,-40.9)--(199.2,-40)--cycle;
				\definecolor{ezs}{RGB}{250,177,11}\fill[ezs](135.8,-210.9)--(138.1,-211.8)--(203.4,-41.6)--(201.1,-40.7)--cycle;
				\definecolor{ezs}{RGB}{250,178,11}\fill[ezs](137.7,-211.6)--(140,-212.5)--(205.3,-42.3)--(203.1,-41.5)--cycle;
				\definecolor{ezs}{RGB}{248,180,10}\fill[ezs](139.6,-212.4)--(141.9,-213.3)--(207.3,-43.1)--(205,-42.2)--cycle;
				\definecolor{ezs}{RGB}{247,181,9}\fill[ezs](141.5,-213.1)--(143.8,-214)--(209.2,-43.8)--(206.9,-42.9)--cycle;
				\definecolor{ezs}{RGB}{246,182,9}\fill[ezs](143.5,-213.8)--(145.7,-214.7)--(211.1,-44.5)--(208.8,-43.7)--cycle;
				\definecolor{ezs}{RGB}{246,182,10}\fill[ezs](145.4,-214.6)--(147.7,-215.5)--(213,-45.3)--(210.7,-44.4)--cycle;
			\end{scope}
			% chữ EDUZY (Times đậm, làm dày nét bằng 5 lớp lệch nhẹ; căn đúng khung chữ logo gốc)
			\foreach \dx/\dy in {-0.7/0,0.7/0,0/-0.3,0/0.3,0/0}
			\node[anchor=center,inner sep=0pt,outer sep=0pt,text=ezText,align=center] at (109.6+\dx,-193.6+\dy)
			{\resizebox{1.01cm}{0.22cm}{\fontencoding{T1}\fontfamily{ptm}\fontseries{b}\selectfont EDUZY}};
		\end{tikzpicture}%
	}%
}
%====================== THIẾT LẬP THÂN ĐỀ ======================
%--- Môi trường tự luận {bt}: ex_test 3.3.4 chưa khai báo sẵn nên tạo ở đây
%    (ex_test đã tự móc \fixshowans và \listenumerate cho bt)
\makeatletter
\@ifundefined{bt}{\newtheorem{bt}{\color{gold1}\namebt}}{}
\makeatother
\renewcommand{\sectionmark}[1]{\markright{\thesection\ #1}}
%--- Khung tiêu đề (logo + tối đa 4 dòng; dòng 3, 4 để trống {} thì bỏ qua)
%  Đề thi:    \khungtieude{THI THỬ THPT QUỐC GIA NĂM 2026}{Môn Toán -- Đề số 21}
%                         {(Đề thi gồm: ...)}{Thời gian làm bài: 90 phút (không kể thời gian phát đề).}
%  Chuyên đề: \khungtieude{CHUYÊN ĐỀ 3 -- DÃY SỐ VÀ CÔNG THỨC TRUY HỒI}
%                         {Bồi dưỡng học sinh giỏi Toán THPT}{...}{}
\newcommand{\khungtieude}[4]{%
	\begin{tcolorbox}[colframe=gold1,colback=gold5!45,boxrule=0.5mm]
		\begin{minipage}[c]{3cm}
			\centerline{\eduzylogo[2.4cm]}
		\end{minipage}\hspace{0.5cm}
		\begin{minipage}[c]{12.5cm}
			\vskip 0.15 cm
			\centerline{\ind{#1}}
			\vskip 0.15 cm
			\centerline{\textbf{#2}}
			\if\relax\detokenize{#3}\relax\else
				\vskip 0.15 cm
				\centerline{\textit{#3}}
			\fi
			\if\relax\detokenize{#4}\relax\else
				\vskip 0.15 cm
				\centerline{\textit{#4}}
			\fi
			\vskip 0.15 cm
		\end{minipage}
	\end{tcolorbox}
}
%====================== BỔ SUNG CHO BÀI HỌC, CHUYÊN ĐỀ ======================
\usepackage[shortlabels]{enumitem}
\usepackage{tabularx}
%--- Danh sách gọn: \begin{enumerate}[gon] ... \end{enumerate}
\SetEnumitemKey{gon}{leftmargin=*,itemsep=1pt,parsep=0pt,topsep=3pt}
%--- Khung lý thuyết: \begin{khung}{} ... \end{khung}; tiêu đề chỉ dùng khi thật ngắn ("Mẹo", "Lỗi thường gặp")
\newtcolorbox{khung}[1]{breakable,enhanced,colframe=gold1,colback=gold5!45,
	boxrule=0.4mm,fonttitle=\bfseries\sffamily,coltitle=white,colbacktitle=gold1,
	attach boxed title to top left={xshift=6mm,yshift=-3mm},
	boxed title style={arc=1mm,boxrule=0.3mm},title={#1}}
%--- Tiêu đề mục nhỏ trong một dạng: \muc{Dấu hiệu}, \muc{Phương pháp}, \muc{Bài tự luyện}
\makeatletter
\newcommand{\muc}[1]{\par\medskip\noindent\indam{#1}\par\nobreak\smallskip\@afterindentfalse\@afterheading}
\makeatother
%--- Ô công thức cần nhớ (căn giữa, viền vàng, vừa khít nội dung)
%     \begin{center}\congthuc{$Q(n)=(n-1)!$}\end{center}
\newtcbox{\congthuc}{on line,colframe=gold1,colback=gold5!45,boxrule=0.5pt,arc=2pt,
	left=8pt,right=8pt,top=4pt,bottom=4pt}
%--- Công tắc lời giải bài tự luyện (ví dụ mẫu luôn dùng \loigiai)
%     \loigiaiTL[{đáp số}]{lời giải đầy đủ}
%     \TLfull -> in lời giải đầy đủ;  \TLdapso -> chỉ in một dòng "Đáp số";  \TLan -> ẩn hoàn toàn
%     (\CheDo là dethi thì mọi lời giải và đáp số đều ẩn)
\makeatletter
\newif\ifTLhien   \TLhientrue
\newif\ifTLdapso  \TLdapsofalse
\newcommand{\TLfull}{\TLhientrue\TLdapsofalse}
\newcommand{\TLdapso}{\TLhientrue\TLdapsotrue}
\newcommand{\TLan}{\TLhienfalse}
\newcommand{\loigiaiTL}[2][]{%
	\ifTLhien
		\ifTLdapso
			\def\@tlarg{#1}%
			\ifx\@tlarg\@empty\else
				\ifdefstring{\kindOPTN}{dethi}{}{\par\noindent{\bfseries\color{gold1}Đáp số:} #1\par}%
			\fi
		\else
			\loigiai{#2}%
		\fi
	\fi}
\makeatother
\TLdapso% mặc định: chỉ in đáp số bài tự luyện
%====================== BẢNG ĐÚNG/SAI (\choiceTFt) ======================
%  Hàng tiêu đề nền nâu đồng chữ trắng (Phát biểu | Đ | S), các hàng tô xen kẽ kem / trắng,
%  đường kẻ vàng kim, nhãn a) b) c) d) màu nâu đồng. Bản có đáp án: cột được chọn có vòng tròn
%  nâu đồng với dấu tích trắng (câu đúng tích ở cột Đúng, câu sai tích ở cột Sai).
%  Tiêu đề ghi đủ "Đúng | Sai": \OPTN{viettat=0} sau \input (cột Đúng/Sai rộng ra, phát biểu hẹp lại).
%  Dùng bảng gốc của ex_test: \def\BangDSgoc{1} trước \input{preamble-eduzy-v2}.
\providecommand{\BangDSgoc}{0}
\ifnum\BangDSgoc=0
\usepackage{colortbl,pifont}
\colorlet{ezDSke}{gold2!60!gold1}% đường kẻ bảng (nét 0,4pt như ex_test để không đổi chiều cao bảng)
\colorlet{ezDSle}{gold5!60}% hàng a), c)
\colorlet{ezDSchan}{white}% hàng b), d)
%--- dấu chọn ở cột Đúng/Sai
%    (bề rộng 0, căn giữa ô: không làm đổi bề rộng cột của ex_test, không báo tràn lề)
\newcommand{\EZdauchon}{\makebox[0pt]{\tikz[baseline=-1ex]\node[circle,fill=gold1,inner sep=0pt,
	minimum size=1.2em,text=white,font=\small]{\ding{51}};}}
\def\tickT{\EZdauchon}\def\tickF{\EZdauchon}
\AtBeginDocument{\sbox{\ccTFTab}{\EZdauchon}}
%--- nhãn a) b) c) d) màu nâu đồng (ở bản lời giải, nhãn câu đúng vẫn đóng khung như ex_test)
\renewcommand{\FalseTF}{\stepcounter{dapan}\textbf{\color{gold1}\DapAnTF\sepTF}\ignorespaces}
%--- trước khi dựng bảng: kẻ màu vàng kim (bề rộng cột giữ nguyên như ex_test); sau bảng trả lại màu đen
\newcount\EZhangTF
\newcommand{\EZbatdauTF}{%
	\global\EZhangTF=0\relax
	\arrayrulecolor{ezDSke}}
\newcommand{\EZketthucTF}{\arrayrulecolor{black}}
%--- hàng tiêu đề nền nâu đồng, chữ trắng
\makeatletter
\newcommand{\EZo@TF}[1]{\centering\bfseries\color{white}#1}
\newcommand{\EZtieudeTF}{%
	\ifcase\cotDS
		\ifnum\cotSTT=1\relax
			\def\rowHeader{\hline\rowcolor{gold1}\EZo@TF{TT}&\EZo@TF{\phatbieu}\tabularnewline}%
		\else
			\def\rowHeader{\hline\rowcolor{gold1}\EZo@TF{\phatbieu}\tabularnewline}%
		\fi
	\or
		\ifnum\cotSTT=1\relax
			\def\rowHeader{\hline\rowcolor{gold1}\EZo@TF{TT}&\EZo@TF{\phatbieu}&\EZo@TF{\Dung}\tabularnewline}%
		\else
			\def\rowHeader{\hline\rowcolor{gold1}\EZo@TF{\phatbieu}&\EZo@TF{\Dung}\tabularnewline}%
		\fi
	\or
		\ifnum\cotSTT=1\relax
			\def\rowHeader{\hline\rowcolor{gold1}\EZo@TF{TT}&\EZo@TF{\phatbieu}&\EZo@TF{\Dung}&\EZo@TF{\Sai}\tabularnewline}%
		\else
			\def\rowHeader{\hline\rowcolor{gold1}\EZo@TF{\phatbieu}&\EZo@TF{\Dung}&\EZo@TF{\Sai}\tabularnewline}%
		\fi
	\fi}
%--- tô xen kẽ từng hàng (gọi ngay sau \hline đầu mỗi hàng)
\newcommand{\EZmauhangTF}{%
	\noalign{\global\advance\EZhangTF\@ne}%
	\ifodd\EZhangTF\rowcolor{ezDSle}\else\rowcolor{ezDSchan}\fi}
%--- gắn vào \choiceTFt / \writeChoiceTFt của ex_test
\xpretocmd{\choiceTFt}{\EZbatdauTF}{}{\typeout{[preamble-eduzy-v2] Không gắn được bảng Đúng/Sai (đầu bảng)}}
\xpatchcmd{\choiceTFt}{\ifnum\dongPB=0\def\rowHeader{}\fi}{\EZtieudeTF\ifnum\dongPB=0\def\rowHeader{}\fi}
	{}{\typeout{[preamble-eduzy-v2] Không gắn được bảng Đúng/Sai (hàng tiêu đề)}}
\xpatchcmd{\writeChoiceTFt}{\hline\xdef\dotEX{\sdotEX}}{\hline\EZmauhangTF\xdef\dotEX{\sdotEX}}
	{}{\typeout{[preamble-eduzy-v2] Không gắn được bảng Đúng/Sai (tô hàng)}}
\xpatchcmd{\writeChoiceTFt}{{\end{longtable}}}{{\end{longtable}\EZketthucTF}}
	{}{\typeout{[preamble-eduzy-v2] Không gắn được bảng Đúng/Sai (cuối longtable)}}
\xpatchcmd{\writeChoiceTFt}{{\end{xtabular}\end{center}\medskip\noindent}}{{\end{xtabular}\end{center}\EZketthucTF\medskip\noindent}}
	{}{\typeout{[preamble-eduzy-v2] Không gắn được bảng Đúng/Sai (cuối xtabular)}}
\fi
%=====================================================================
\renewcommand{\parallel}{\mathrel{/\!\!/}}

%  trong file đề (tools/dan_trang/dan_trang.py tự ghi khi đặt \NeedspaceCau):
%    \def\TachCauDS{1}   câu Đúng/Sai được tách: câu dẫn (kèm hình) ở cuối trang, bảng Phát biểu sang trang sau
%                        (câu dẫn và hình không tách nhau, bảng không tách giữa hai trang)
%    \def\TachCauLG{1}   bản lời giải (\CheDo khác dethi): phần đề của câu cao hơn 1/3 trang được phép tách
%  Lệnh dùng trong file đề (giá trị do tools/dan_trang/dan_trang.py đo ở bản dethi):
%    \NeedspaceCau[<cao câu dẫn + hình>]{<cao cả khối câu>}   ngay trên \begin{ex} (câu đầu phần: trên tiêu đề phần)
%        công tắc tắt: đúng bằng \Needspace{<cao cả khối câu>} như cũ;
%        \TachCauDS=1 và có [..] (câu Đúng/Sai được chọn tách): chỉ giữ câu dẫn + hình;
%        \TachCauLG=1 ở bản lời giải: đòi tối đa 1/3 \textheight.
%    \NeedspaceBang{<cao bảng Phát biểu>}%   ngay trước \choiceTFt: \TachCauDS=1 thì giữ trọn bảng trên một trang;
%        công tắc tắt: không làm gì.
\providecommand{\TachCauDS}{0}
\providecommand{\TachCauLG}{0}
\newlength{\EZnsCau}
\newcommand{\NeedspaceCau}[2][]{%
	\setlength{\EZnsCau}{#2}%
	\ifnum\TachCauDS=1 \if\relax\detokenize{#1}\relax\else\setlength{\EZnsCau}{#1}\fi\fi
	\ifnum\TachCauLG=1 \ifdefstring{\kindOPTN}{dethi}{}{\ifdim\EZnsCau>0.3333\textheight\setlength{\EZnsCau}{0.3333\textheight}\fi}\fi
	\Needspace{\EZnsCau}}
\newcommand{\NeedspaceBang}[1]{\ifnum\TachCauDS=1 \Needspace{#1}\fi}

%--- Thang màu: đậm -> nhạt
\definecolor{gold1}{RGB}{126,96,0}    % nâu đồng đậm  #7E6000
\definecolor{gold2}{RGB}{192,144,3}   % vàng kim      #C09003
\definecolor{gold3}{RGB}{255,215,99}  % vàng nhạt     #FFD763
\definecolor{gold4}{RGB}{255,231,153} % vàng rất nhạt #FFE799
\definecolor{gold5}{RGB}{255,243,205} % vàng kem      #FFF3CD
%--- Ghi đè tông xanh/đỏ cũ
\definecolor{myblue}{RGB}{126,96,0}
\colorlet{blue}{gold1}
\colorlet{cyan}{gold2}
\colorlet{teal}{gold2}

%--- Tên chương mặc định (file đề nên tự khai báo lại)
\providecommand{\tenchuong}{}
%--- Chân trang
\providecommand{\thuonghieu}{TOÁN 12 -- Eduzy.com.vn}
\providecommand{\tacgia}{Tô Gia Bảo -- 0353703783}

%--- Chữ cạnh logo ở đầu trang
\providecommand{\khauhieu}{Eduzy.com.vn}

%====================== THÔNG SỐ ĐẦU TRANG, CHÂN TRANG (v2) ======================
%  Dải nền đậm tràn ngang trang; ô chữ vàng kem có cạnh xiên, kèm 2 vạch xiên chuyển màu.
%  Đầu trang lẻ : logo + \khauhieu (trái)   | ô chữ \tenchuong (phải)
%  Chân trang lẻ: ô chữ \thuonghieu (trái)  | Trang (số) (phải)
%  Trang chẵn soi gương trang lẻ; chân trang chẵn ghi \tacgia.
%--- Màu: chỉ dùng thang vàng kim Eduzy
\colorlet{ezNen}{gold1}% dải nền đậm
\colorlet{ezVach1}{gold3}% vạch xiên sát dải nền (đậm hơn)
\colorlet{ezVach2}{gold4}% vạch xiên sát ô chữ (nhạt hơn)
\colorlet{ezO}{gold5}% ô chữ vàng kem
%--- Kích thước
\newlength{\EZdai}\setlength{\EZdai}{1.1cm}% chiều cao dải nền
\newlength{\EZo}\setlength{\EZo}{0.68cm}% chiều cao ô chữ
\newlength{\EZxien}\setlength{\EZxien}{0.55cm}% độ lệch ngang của cạnh xiên
\newlength{\EZvach}\setlength{\EZvach}{0.75cm}% bề rộng mỗi vạch xiên
\newlength{\EZkhe}\setlength{\EZkhe}{0.12cm}% khe trắng giữa vạch và ô chữ
\newlength{\EZdem}\setlength{\EZdem}{0.4cm}% khoảng đệm từ chữ tới cạnh xiên
\newlength{\EZle}\setlength{\EZle}{1.8cm}% lề trái/phải của chữ (bằng lề thân bài)
\newlength{\EZrongmin}\setlength{\EZrongmin}{4.5cm}% ô chữ rộng tối thiểu
\newlength{\EZrong}% bề rộng chữ trong ô (tính tự động)
%--- Font
\newcommand{\EZfontO}{\fontsize{10.5pt}{12pt}\selectfont\bfseries}% chữ trong ô
\newcommand{\EZfontKH}{\fontfamily{qag}\fontsize{11pt}{12pt}\selectfont\bfseries}% \khauhieu cạnh logo
\newcommand{\EZfontTrang}{\fontsize{11.5pt}{12pt}\selectfont\bfseries}% chữ "Trang" và số trang
%--- Đo bề rộng chữ trong ô
\newcommand{\EZdorong}[1]{%
	\settowidth{\EZrong}{\EZfontO #1}%
	\ifdim\EZrong<\EZrongmin\setlength{\EZrong}{\EZrongmin}\fi}
%--- Logo + chữ trắng \khauhieu (dựng sẵn một lần khi vào \begin{document})
\newsavebox{\EZlogobox}
\AtBeginDocument{\sbox{\EZlogobox}{%
		\raisebox{-0.5\height}{\eduzylogo[1cm]}\hspace{0.15cm}%
		\raisebox{-0.5\height}{\color{white}\EZfontKH\khauhieu}}}
%--- Hai vạch xiên + khe trắng. #1: điểm gốc ở mép dưới, #2: hướng (1 sang phải, -1 sang trái),
%    #3: hướng nghiêng (1 kiểu "/", -1 kiểu "\")
\newcommand{\EZvachxien}[3]{%
	\fill[white] (#1) --++({(#3)*\EZxien},\EZo) --++({(#2)*\EZkhe},0) --++({-(#3)*\EZxien},-\EZo) --cycle;
	\fill[ezVach2] ([xshift=(#2)*\EZkhe]#1) --++({(#3)*\EZxien},\EZo) --++({(#2)*\EZvach},0) --++({-(#3)*\EZxien},-\EZo) --cycle;
	\fill[ezVach1] ([xshift=(#2)*(\EZkhe+\EZvach)]#1) --++({(#3)*\EZxien},\EZo) --++({(#2)*\EZvach},0) --++({-(#3)*\EZxien},-\EZo) --cycle;
	\draw[white,line width=0.6pt] ([xshift=(#2)*(\EZkhe+\EZvach)]#1) --++({(#3)*\EZxien},\EZo);
}
%--- Chữ trong ô (căn giữa theo chiều cao ô, không lệch vì dấu tiếng Việt)
\tikzset{ezchu/.style={inner sep=0pt,text=ezNen,text height=0.72em,text depth=0.22em}}

%====================== LOGO CHÌM GIỮA TRANG ======================
%  Vẽ cùng đầu trang nên nằm dưới chữ của thân bài, có trên mọi trang.
%  Tắt: \logochimtat (từ trang đó trở đi), bật lại: \logochimbat.
%  Đổi cỡ, độ đậm: khai báo trước \input, ví dụ \def\EZchimcao{9cm} \def\EZchimdo{0.1}
\providecommand{\EZchimcao}{10cm}% chiều cao logo chìm
\providecommand{\EZchimdo}{0.12}% độ đậm (0 = trong suốt, 1 = rõ hẳn)
\newif\ifEZchim \EZchimtrue
\newcommand{\logochimtat}{\EZchimfalse}
\newcommand{\logochimbat}{\EZchimtrue}
%  Dựng sẵn logo mờ một lần khi vào \begin{document}; transparency group làm mờ cả logo
%  một lượt nên không lộ vệt ở chỗ các dải màu chồng nhau.
\newsavebox{\EZchimtho}
\newsavebox{\EZchimbox}
\AtBeginDocument{%
	\sbox{\EZchimtho}{\eduzylogo[\EZchimcao]}%
	\sbox{\EZchimbox}{\begin{tikzpicture}
			\begin{scope}[transparency group,opacity=\EZchimdo]
				\node[inner sep=0pt] {\usebox{\EZchimtho}};
			\end{scope}
	\end{tikzpicture}}}
\newcommand{\EZvechim}{%
	\ifEZchim
		\node[inner sep=0pt] at (current page.center) {\usebox{\EZchimbox}};
	\fi}

%====================== ĐẦU TRANG - TRANG LẺ ======================
%  (logo bên trái, ô chữ \tenchuong bên phải)
\newcommand{\myfancyheadOdd}{%
	\EZdorong{\tenchuong}%
	\begin{tikzpicture}[remember picture,overlay]
		\coordinate (T) at (current page.north west);
		\coordinate (P) at (current page.north east);
		%--- logo chìm giữa trang
		\EZvechim
		%--- dải nền
		\fill[ezNen] (T) rectangle ([yshift=-\EZdai]P);
		%--- ô chữ: X là góc dưới bên trái, cạnh trái xiên kiểu "/"
		\coordinate (X) at ([xshift=-\EZle-\EZrong-\EZdem-\EZxien,yshift=-\EZdai]P);
		\fill[ezO] (X) --++(\EZxien,\EZo) -- ([yshift=-\EZdai+\EZo]P) -- ([yshift=-\EZdai]P) --cycle;
		\EZvachxien{X}{-1}{1}
		\node[ezchu,anchor=east] at ([xshift=-\EZle,yshift=-\EZdai+0.5*\EZo]P) {\EZfontO\tenchuong};
		%--- logo + tên
		\node[inner sep=0pt,anchor=west] at ([xshift=\EZle,yshift=-0.5*\EZdai]T) {\usebox{\EZlogobox}};
	\end{tikzpicture}%
}

%====================== ĐẦU TRANG - TRANG CHẴN =====================
%  (soi gương: ô chữ \tenchuong bên trái, logo bên phải)
\newcommand{\myfancyheadEven}{%
	\EZdorong{\tenchuong}%
	\begin{tikzpicture}[remember picture,overlay]
		\coordinate (T) at (current page.north west);
		\coordinate (P) at (current page.north east);
		%--- logo chìm giữa trang
		\EZvechim
		%--- dải nền
		\fill[ezNen] (T) rectangle ([yshift=-\EZdai]P);
		%--- ô chữ: X là góc dưới bên phải, cạnh phải xiên kiểu "\"
		\coordinate (X) at ([xshift=\EZle+\EZrong+\EZdem+\EZxien,yshift=-\EZdai]T);
		\fill[ezO] (X) --++(-\EZxien,\EZo) -- ([yshift=-\EZdai+\EZo]T) -- ([yshift=-\EZdai]T) --cycle;
		\EZvachxien{X}{1}{-1}
		\node[ezchu,anchor=west] at ([xshift=\EZle,yshift=-\EZdai+0.5*\EZo]T) {\EZfontO\tenchuong};
		%--- logo + tên
		\node[inner sep=0pt,anchor=east] at ([xshift=-\EZle,yshift=-0.5*\EZdai]P) {\usebox{\EZlogobox}};
	\end{tikzpicture}%
}

%====================== CHÂN TRANG - TRANG LẺ ======================
%  (ô chữ \thuonghieu bên trái, "Trang (số)" bên phải)
\newcommand{\myfancyfootOdd}{%
	\EZdorong{\thuonghieu}%
	\begin{tikzpicture}[remember picture,overlay]
		\coordinate (T) at (current page.south west);
		\coordinate (P) at (current page.south east);
		%--- dải nền
		\fill[ezNen] (T) rectangle ([yshift=\EZdai]P);
		%--- ô chữ sát mép trên dải: X là góc dưới bên phải, cạnh phải xiên kiểu "/"
		\coordinate (X) at ([xshift=\EZle+\EZrong+\EZdem,yshift=\EZdai-\EZo]T);
		\fill[ezO] ([yshift=\EZdai-\EZo]T) -- (X) --++(\EZxien,\EZo) -- ([yshift=\EZdai]T) --cycle;
		\EZvachxien{X}{1}{1}
		\node[ezchu,anchor=west] at ([xshift=\EZle,yshift=\EZdai-0.5*\EZo]T) {\EZfontO\thuonghieu};
		%--- số trang
		\node[circle,draw=white,line width=0.9pt,text=white,inner sep=0pt,minimum size=0.62cm,anchor=east]
		(so) at ([xshift=-\EZle,yshift=0.5*\EZdai]P) {\EZfontTrang\thepage};
		\node[inner sep=0pt,text=white,anchor=base east] (tr) at ([xshift=-0.15cm]so.base west) {\EZfontTrang Trang};
	\end{tikzpicture}%
}

%====================== CHÂN TRANG - TRANG CHẴN ====================
%  (soi gương: "Trang (số)" bên trái, ô chữ \tacgia bên phải)
\newcommand{\myfancyfootEven}{%
	\EZdorong{\tacgia}%
	\begin{tikzpicture}[remember picture,overlay]
		\coordinate (T) at (current page.south west);
		\coordinate (P) at (current page.south east);
		%--- dải nền
		\fill[ezNen] (T) rectangle ([yshift=\EZdai]P);
		%--- ô chữ sát mép trên dải: X là góc dưới bên trái, cạnh trái xiên kiểu "\"
		\coordinate (X) at ([xshift=-\EZle-\EZrong-\EZdem,yshift=\EZdai-\EZo]P);
		\fill[ezO] ([yshift=\EZdai-\EZo]P) -- (X) --++(-\EZxien,\EZo) -- ([yshift=\EZdai]P) --cycle;
		\EZvachxien{X}{-1}{-1}
		\node[ezchu,anchor=east] at ([xshift=-\EZle,yshift=\EZdai-0.5*\EZo]P) {\EZfontO\tacgia};
		%--- số trang
		%    (vòng tròn ẩn so0 nằm giữa dải để lấy đường chân chữ, giống trang lẻ)
		\node[circle,inner sep=0pt,minimum size=0.62cm] (so0) at ([yshift=0.5*\EZdai]T) {\EZfontTrang\phantom{\thepage}};
		\node[inner sep=0pt,text=white,anchor=base west] (tr) at ($(so0.base -| T)+(\EZle,0)$) {\EZfontTrang Trang};
		\node[circle,draw=white,line width=0.9pt,text=white,inner sep=0pt,minimum size=0.62cm,anchor=base west]
		(so) at ([xshift=0.15cm]tr.base east) {\EZfontTrang\thepage};
	\end{tikzpicture}%
}

%--- Giữ tương thích với các file cũ đang gọi \myfancyhead, \myfancyfoot
\newcommand{\myfancyhead}{\myfancyheadOdd}
\newcommand{\myfancyfoot}{\myfancyfootOdd}

%======================= GẮN VÀO TRANG =============================
\pagestyle{fancy}
\fancyhf{}
\fancyhead[LO]{\myfancyheadOdd}
\fancyhead[RE]{\myfancyheadEven}
\fancyfoot[LO]{\myfancyfootOdd}
\fancyfoot[RE]{\myfancyfootEven}
\renewcommand{\headrulewidth}{0pt}
\renewcommand{\footrulewidth}{0pt}
\fancypagestyle{plain}{%
	\fancyhf{}
	\fancyhead[LO]{\myfancyheadOdd}
	\fancyhead[RE]{\myfancyheadEven}
	\fancyfoot[LO]{\myfancyfootOdd}
	\fancyfoot[RE]{\myfancyfootEven}
	\renewcommand{\headrulewidth}{0pt}
	\renewcommand{\footrulewidth}{0pt}
}
%=================== HẾT PHẦN GỘP TỪ eduzy-de.tex ===================
\newcommand{\indam}[1]{
	{\fontsize{12pt}{9pt}\sffamily \selectfont\bfseries\color{gold1} #1}
}
\newcommand{\ind}[1]{
	{\fontsize{14pt}{9pt}\fontfamily{qag} \selectfont\bfseries\color{gold1} #1}
}
\renewcommand{\FalseEX}{\stepcounter{dapan}\textbf{\selectfont\color{gold1}\Alph{dapan}.}\ignorespaces}
\renewtheorem{ex}{\color{gold1}Câu}
\renewcommand{\baselinestretch}{1.1}
\newcommand{\hoac}[1]{
	\left[\begin{aligned}#1\end{aligned}\right.}
\newcommand{\heva}[1]{
	\left\{\begin{aligned}#1\end{aligned}\right.}
\usepackage{esvect}
\def\vec{\vv}
\def\vt{\vv}
\def\overrightarrow{\vv}
%====================== LOGO EDUZY VẼ BẰNG TIKZ ======================
%  Đường viền dò từ logo gốc; dải màu vàng kim đo trực tiếp trên ảnh gốc
%  (nghiêng 21 độ, tô bằng 80 dải hẹp nên hiển thị giống nhau trên mọi trình xem PDF).
%  Dùng: \eduzylogo          -> cao 2.4cm (mặc định)
%        \eduzylogo[1.5cm]   -> cao 1.5cm
\definecolor{ezText}{RGB}{245,200,46}
\newcommand{\eduzylogo}[1][2.4cm]{%
	\resizebox{!}{#1}{%
		\begin{tikzpicture}[x=0.01cm,y=0.01cm]
			\path[use as bounding box] (0,0) rectangle (224,-224);
			% mũ tốt nghiệp + vương miện
			\begin{scope}
				\clip
				(119.4,-176.1)..controls(118.9,-176.1)and(116.8,-176)..(114.6,-176)
				..controls(101.4,-176)and(85.8,-175.9)..(83.7,-175.7)
				..controls(81.8,-175.6)and(81,-175.6)..(80.2,-175.7)
				..controls(78.8,-175.9)and(77.7,-175.9)..(76.2,-175.7)
				..controls(74.9,-175.5)and(70.8,-175.5)..(64.3,-175.7)
				..controls(59.1,-175.9)and(58.3,-175.7)..(58.6,-173.7)
				..controls(58.6,-172.8)and(59.6,-169.6)..(60.3,-168)
				..controls(60.8,-166.9)and(61.4,-165.8)..(62.3,-165)
				..controls(63.1,-164.2)and(64.3,-163.4)..(65,-163.2)
				..controls(65.5,-163.1)and(72.3,-163)..(76.6,-163.1)
				..controls(77.9,-163.1)and(80.1,-163.1)..(81.4,-163)
				..controls(82.8,-162.9)and(86.3,-162.9)..(89.2,-162.9)
				..controls(93,-162.9)and(94.9,-162.8)..(96,-162.7)
				..controls(97.2,-162.5)and(101,-162.5)..(114,-162.4)
				..controls(123,-162.3)and(134.9,-162.2)..(140.2,-162.1)
				..controls(150.1,-162)and(152.2,-162.1)..(153.6,-162.4)
				..controls(154.4,-162.6)and(157.1,-163.8)..(158.1,-164.4)
				..controls(158.8,-164.8)and(159.7,-165.7)..(159.9,-166.3)
				..controls(160,-166.5)and(160.2,-167.1)..(160.2,-167.6)
				..controls(160.3,-168)and(160.6,-169.4)..(160.9,-170.6)
				..controls(161.2,-171.7)and(161.5,-173.1)..(161.6,-173.6)
				..controls(161.8,-174.4)and(161.8,-174.5)..(161.5,-174.9)
				..controls(161.2,-175.6)and(160.8,-175.9)..(159.9,-176)
				..controls(159.2,-176.1)and(154.1,-176.2)..(151.5,-176.2)
				..controls(148.1,-176.1)and(122.6,-176.1)..(121.6,-176.2)
				..controls(120.8,-176.2)and(119.8,-176.2)..(119.4,-176.1)--cycle
				(130.4,-160)..controls(129,-160)and(119.8,-160)..(110.1,-159.9)
				..controls(100.3,-159.8)and(88.5,-159.7)..(83.9,-159.6)
				..controls(77.8,-159.5)and(75.3,-159.5)..(74.5,-159.6)
				..controls(73.9,-159.7)and(72.5,-159.8)..(71.3,-159.8)
				..controls(69.6,-159.8)and(69,-159.8)..(68.5,-159.6)
				..controls(67.6,-159.3)and(66.8,-158.6)..(66.2,-157.6)
				..controls(66,-157.3)and(65.6,-156.5)..(65.3,-156.1)
				..controls(64.6,-155.1)and(64.5,-154.9)..(63.7,-152.9)
				..controls(63.2,-151.6)and(62.9,-150.9)..(61.8,-149.6)
				..controls(61.5,-149.2)and(61.1,-148.5)..(60.9,-147.9)
				..controls(60.7,-147.4)and(60.3,-146.6)..(60.1,-146.1)
				..controls(59.9,-145.5)and(59.5,-144.9)..(59,-144.3)
				..controls(58.5,-143.7)and(58.2,-143.2)..(57.9,-142.5)
				..controls(57.7,-141.9)and(57.3,-141.1)..(57,-140.7)
				..controls(56.6,-140.3)and(56.1,-139.3)..(55.8,-138.5)
				..controls(55.5,-137.8)and(55.1,-136.8)..(54.9,-136.5)
				..controls(54.7,-136.1)and(54.3,-135.2)..(53.9,-134.5)
				..controls(53.6,-133.8)and(53,-132.7)..(52.6,-132)
				..controls(51.6,-130.5)and(51.4,-130.2)..(51.1,-129)
				..controls(50.6,-127.5)and(50.6,-127.4)..(49.7,-126.3)
				..controls(49.2,-125.7)and(48.6,-124.8)..(48.3,-124.2)
				..controls(48,-123.6)and(47.5,-122.7)..(47.3,-122.2)
				..controls(47,-121.8)and(46.5,-120.7)..(46.1,-119.8)
				..controls(45.8,-118.9)and(45.2,-117.8)..(44.9,-117.2)
				..controls(44.2,-116.1)and(43.3,-114.4)..(42.3,-112.2)
				..controls(41.8,-111.3)and(41.3,-110.4)..(41.1,-110.1)
				..controls(40.5,-109.4)and(40.3,-109)..(40,-107.8)
				..controls(39.8,-107)and(39.6,-106.5)..(39.2,-106)
				..controls(38.9,-105.6)and(38.4,-104.7)..(38,-104)
				..controls(37.7,-103.4)and(37.2,-102.5)..(37,-102.1)
				..controls(36.7,-101.7)and(36.3,-100.8)..(36,-100.1)
				..controls(35.6,-98.9)and(35.3,-98.4)..(34.6,-97.6)
				..controls(34.2,-97.1)and(34.2,-96.3)..(34.6,-96)
				..controls(35,-95.7)and(35.5,-95.9)..(36.2,-96.7)
				..controls(36.5,-97)and(37.2,-97.5)..(37.6,-97.7)..controls(38,-98)and(38.6,-98.4)..(39.1,-98.9)
				..controls(39.6,-99.4)and(40.3,-100)..(40.6,-100.3)
				..controls(41.3,-100.8)and(42.4,-101.8)..(43.7,-102.9)
				..controls(44,-103.2)and(44.6,-103.7)..(45,-103.9)
				..controls(45.4,-104.2)and(45.9,-104.6)..(46.1,-104.8)
				..controls(47.1,-105.9)and(48.8,-107.1)..(50.5,-107.9)
				..controls(51.8,-108.6)and(52.3,-109)..(53.4,-110)
				..controls(54.8,-111.2)and(55.6,-111.6)..(57,-112.1)
				..controls(57.5,-112.3)and(58.5,-112.8)..(59.4,-113.3)
				..controls(60.2,-113.7)and(61.3,-114.3)..(61.9,-114.5)
				..controls(62.4,-114.7)and(63.3,-115.1)..(63.8,-115.3)
				..controls(64.4,-115.6)and(65.4,-115.9)..(66.2,-116.1)
				..controls(67,-116.3)and(68.4,-116.6)..(69.2,-116.9)
				..controls(72.3,-117.9)and(72.8,-118)..(78,-118.1)
				..controls(84,-118.3)and(86.1,-118.1)..(88.1,-117.1)
				..controls(89.4,-116.5)and(93,-115)..(93.7,-114.9)
				..controls(94.8,-114.6)and(95.4,-114.2)..(95.9,-113.4)
				..controls(96.2,-113)and(96.7,-112.5)..(97.1,-112.2)
				..controls(98,-111.6)and(101.3,-108.4)..(102.1,-107.2)
				..controls(102.4,-106.8)and(102.9,-106.2)..(103.2,-105.8)
				..controls(103.5,-105.5)and(104,-104.8)..(104.2,-104.2)
				..controls(104.5,-103.6)and(105,-102.8)..(105.3,-102.4)
				..controls(105.6,-102)and(106.1,-101.2)..(106.5,-100.6)
				..controls(107.1,-99.5)and(108.6,-97.5)..(109.5,-96.5)
				..controls(109.9,-96.1)and(110.1,-96)..(110.4,-96)
				..controls(111.1,-96)and(111.6,-96.5)..(111.8,-97.6)
				..controls(112,-98.6)and(112.3,-99.3)..(112.8,-99.9)
				..controls(113.7,-100.9)and(115.2,-103.1)..(115.7,-104.2)
				..controls(116.1,-105)and(116.5,-105.5)..(116.9,-106)
				..controls(117.3,-106.4)and(117.8,-106.9)..(118,-107.3)
				..controls(118.2,-107.6)and(118.6,-108)..(119,-108.3)
				..controls(119.6,-108.8)and(121.1,-110.3)..(121.8,-111.2)
				..controls(122.1,-111.7)and(122.5,-112)..(123.4,-112.4)
				..controls(124.8,-113.2)and(125.8,-113.8)..(127,-114.6)
				..controls(127.5,-114.9)and(128.1,-115.3)..(128.2,-115.4)
				..controls(128.4,-115.5)and(129.3,-115.8)..(130.2,-116)
				..controls(131.1,-116.2)and(132.5,-116.6)..(133.3,-116.9)
				..controls(135.7,-117.9)and(136.1,-117.9)..(142.5,-117.8)
				..controls(145.8,-117.8)and(148.4,-117.7)..(148.8,-117.6)
				..controls(149.2,-117.6)and(150.2,-117.2)..(151,-116.8)
				..controls(152,-116.4)and(152.8,-116)..(153.6,-115.9)
				..controls(154.5,-115.7)and(155.2,-115.4)..(156.3,-114.9)
				..controls(157.1,-114.5)and(158.2,-114)..(158.7,-113.8)
				..controls(159.1,-113.7)and(159.8,-113.3)..(160.2,-113)
				..controls(160.6,-112.7)and(161.3,-112.4)..(161.7,-112.2)
				..controls(162.7,-111.9)and(163.2,-111.7)..(164.1,-111.1)
				..controls(164.5,-110.8)and(165.3,-110.4)..(165.9,-110)
				..controls(166.5,-109.7)and(167.3,-109.2)..(167.6,-108.9)
				..controls(167.9,-108.6)and(168.6,-108.2)..(169,-108)
				..controls(169.4,-107.8)and(170.2,-107.3)..(170.7,-106.9)
				..controls(171.2,-106.5)and(172.1,-105.8)..(172.7,-105.4)
				..controls(173.2,-105)and(174.1,-104.4)..(174.5,-104)
				..controls(174.9,-103.7)and(175.7,-103.1)..(176.2,-102.6)
				..controls(176.7,-102.2)and(177.4,-101.5)..(177.9,-101.2)
				..controls(178.3,-100.8)and(178.8,-100.4)..(179,-100.2)
				..controls(179.2,-100)and(179.6,-99.6)..(180,-99.3)
				..controls(180.4,-99.1)and(180.9,-98.6)..(181.2,-98.4)
				..controls(181.5,-98.1)and(182.1,-97.7)..(182.5,-97.5)
				..controls(183.2,-97.2)and(183.5,-96.9)..(184.3,-95.9)
				..controls(184.6,-95.6)and(184.6,-95.6)..(184.8,-95.7)
				..controls(185,-95.8)and(185.1,-96)..(185.1,-96.2)
				..controls(185.1,-96.5)and(185,-96.7)..(184.6,-97.2)
				..controls(184.1,-97.8)and(184,-97.9)..(183.9,-98.8)
				..controls(183.7,-100.3)and(183.3,-101.3)..(182.9,-101.8)
				..controls(182.6,-102.1)and(182.3,-102.6)..(182.2,-102.9)
				..controls(182,-103.3)and(181.6,-104)..(181.2,-104.6)
				..controls(180.9,-105.2)and(180.5,-106.1)..(180.4,-106.6)
				..controls(179.9,-108)and(179.7,-108.5)..(179,-109.4)
				..controls(178.6,-109.8)and(178.2,-110.5)..(177.9,-111)
				..controls(177.7,-111.5)and(177.3,-112.3)..(177,-112.7)
				..controls(176.6,-113.2)and(176.4,-113.7)..(176.3,-114.2)
				..controls(176.2,-114.7)and(175.9,-115.2)..(175.4,-116)
				..controls(175,-116.6)and(174.4,-117.7)..(174,-118.5)
				..controls(173.7,-119.2)and(173.3,-120.1)..(173,-120.5)
				..controls(172.8,-120.9)and(172.4,-121.8)..(172.1,-122.5)
				..controls(171.8,-123.2)and(171.3,-124.1)..(171,-124.6)
				..controls(170.8,-125.1)and(170.3,-126)..(169.9,-126.6)
				..controls(169.6,-127.3)and(169.2,-128.1)..(169,-128.3)
				..controls(168.8,-128.6)and(168.5,-129.3)..(168.3,-129.8)
				..controls(168.1,-130.3)and(167.8,-131.1)..(167.5,-131.6)
				..controls(166.9,-132.5)and(164.2,-137.9)..(164.2,-138.1)
				..controls(164.2,-138.1)and(163.9,-138.8)..(163.4,-139.7)
				..controls(162.9,-140.5)and(162.4,-141.5)..(162.3,-141.9)
				..controls(161.9,-143)and(161.6,-143.4)..(160.9,-144.1)
				..controls(160.4,-144.6)and(160.3,-144.8)..(160.1,-145.4)
				..controls(160,-145.8)and(159.6,-146.8)..(159.2,-147.8)
				..controls(158.7,-148.8)and(158.1,-150)..(157.9,-150.6)
				..controls(157.6,-151.1)and(157.2,-151.9)..(156.9,-152.4)
				..controls(156.5,-152.8)and(156.1,-153.7)..(155.8,-154.4)
				..controls(155.6,-155)and(155.2,-155.9)..(154.9,-156.3)
				..controls(154.7,-156.7)and(154.4,-157.4)..(154.2,-157.9)
				..controls(154.1,-158.4)and(153.8,-159)..(153.5,-159.3)
				..controls(152.8,-160.1)and(153,-160.1)..(142.4,-160.1)
				..controls(137.2,-160.1)and(131.8,-160.1)..(130.4,-160)--cycle
				(141.5,-103.7)..controls(140.8,-103)and(139.8,-102.1)..(139.2,-101.6)
				..controls(138.5,-101.2)and(137.9,-100.6)..(137.8,-100.4)
				..controls(137,-99.5)and(134,-96.5)..(133.2,-96)
				..controls(132.7,-95.7)and(131.8,-94.9)..(131.1,-94.3)
				..controls(129.4,-92.7)and(127.9,-91.6)..(126.6,-91.1)
				..controls(125.2,-90.5)and(124.1,-90)..(122.9,-89.1)
				..controls(122.1,-88.6)and(121.3,-88.2)..(119.9,-87.7)
				..controls(118.8,-87.3)and(117.5,-86.8)..(116.9,-86.5)
				..controls(114.3,-85.5)and(113.7,-85.4)..(109.7,-85.5)
				..controls(105.3,-85.7)and(103.9,-85.9)..(101.6,-87)
				..controls(100.7,-87.4)and(99.4,-87.8)..(97.9,-88)
				..controls(97,-88.2)and(96.6,-88.3)..(95.8,-88.9)
				..controls(95.4,-89.2)and(94.7,-89.6)..(94.2,-89.8)
				..controls(92.9,-90.3)and(92.3,-90.7)..(90.9,-92.1)
				..controls(90.2,-92.7)and(89.4,-93.3)..(89,-93.5)..controls(87,-94.5)and(85.2,-96)..(83.6,-97.9)
				..controls(83.3,-98.2)and(82.5,-98.9)..(81.8,-99.4)
				..controls(80.7,-100.1)and(80.5,-100.3)..(80,-101.1)
				..controls(78.7,-103.1)and(77.2,-103.9)..(76.4,-103)
				..controls(75.7,-102.2)and(75.7,-103.2)..(75.7,-83.6)
				..controls(75.6,-73.6)and(75.6,-64.8)..(75.5,-64)
				..controls(75.5,-63.2)and(75.5,-61.4)..(75.6,-59.8)--(75.7,-57.1)--(76,-56.8)
				..controls(76.8,-56.2)and(78,-56.3)..(79.8,-57.2)
				..controls(80.4,-57.5)and(81.3,-57.8)..(81.8,-57.9)
				..controls(82.3,-58.1)and(83,-58.4)..(83.6,-58.8)..controls(84.3,-59.2)and(85,-59.5)..(86,-59.8)
				..controls(86.8,-60.1)and(88,-60.6)..(88.7,-60.9)
				..controls(89.3,-61.2)and(90.6,-61.7)..(91.4,-62)
				..controls(92.2,-62.3)and(93.7,-62.9)..(94.6,-63.3)
				..controls(95.6,-63.7)and(96.8,-64.1)..(97.4,-64.3)
				..controls(98,-64.5)and(99,-64.9)..(99.7,-65.2)
				..controls(100.4,-65.6)and(101.5,-66)..(102.3,-66.2)
				..controls(104.2,-66.8)and(105,-67.1)..(106.2,-67.8)
				..controls(107.9,-68.7)and(110.2,-69.4)..(111.3,-69.2)
				..controls(112.5,-69)and(115,-68.2)..(115.6,-67.8)
				..controls(117.3,-66.7)and(117.6,-66.6)..(119.5,-66.2)
				..controls(121.9,-65.8)and(122.1,-65.7)..(123,-65.1)
				..controls(124.1,-64.4)and(124.4,-64.3)..(126,-63.9)
				..controls(127,-63.7)and(127.7,-63.4)..(128.4,-63.1)
				..controls(129,-62.8)and(129.7,-62.5)..(130.6,-62.4)
				..controls(132.1,-62.1)and(134,-61.5)..(135.2,-61)
				..controls(135.7,-60.8)and(136.7,-60.5)..(137.3,-60.3)
				..controls(138.4,-60)and(139.3,-59.7)..(141.1,-59)
				..controls(141.6,-58.7)and(142,-58.6)..(142.5,-58.6)
				..controls(143.2,-58.6)and(143.3,-58.7)..(143.6,-59)
				..controls(144,-59.5)and(144.1,-59.9)..(144,-60.8)
				..controls(144,-61.2)and(144,-71.1)..(144,-82.8)--(143.9,-104.1)--(143.7,-104.5)
				..controls(143.5,-104.7)and(143.3,-104.8)..(143,-104.9)
				..controls(142.6,-104.9)and(142.5,-104.8)..(141.5,-103.7)--cycle
				(27.8,-91.4)..controls(27.1,-90.9)and(26.3,-90)..(26.1,-89.5)
				..controls(26,-89.2)and(25.8,-88.6)..(25.7,-87.9)
				..controls(25.4,-86.2)and(25.6,-84.8)..(26.5,-83.5)
				..controls(27.4,-82.3)and(28.3,-82)..(29.4,-82.3)
				..controls(31.1,-82.8)and(32.1,-84.9)..(32,-87.4)
				..controls(31.9,-89.3)and(31.2,-90.6)..(29.9,-91.4)
				..controls(29.1,-91.9)and(28.6,-91.9)..(27.8,-91.4)--cycle
				(190.1,-91.3)..controls(189.7,-91.2)and(189.3,-90.9)..(189.1,-90.7)
				..controls(188.3,-89.8)and(187.6,-87.7)..(187.6,-86.2)
				..controls(187.6,-85.5)and(187.7,-85.2)..(188.2,-84.4)
				..controls(189.1,-82.6)and(190.1,-81.8)..(191.3,-82.1)
				..controls(192,-82.3)and(193.5,-83.7)..(193.8,-84.4)
				..controls(194.1,-85.1)and(194.2,-87.1)..(194,-88.3)
				..controls(193.8,-89.8)and(193.3,-90.5)..(192.3,-91.3)
				..controls(191.7,-91.7)and(190.9,-91.7)..(190.1,-91.3)--cycle
				(173.7,-83.4)..controls(172.2,-82.5)and(172,-82.5)..(170.6,-82.4)
				..controls(169.7,-82.4)and(168.8,-82.4)..(168.3,-82.5)
				..controls(167.2,-82.7)and(166.9,-82.7)..(166.6,-82.2)
				..controls(166.3,-81.8)and(166.3,-81)..(166.8,-80)
				..controls(167.2,-78.9)and(167.4,-77.4)..(167.7,-74.1)
				..controls(167.8,-72.8)and(167.9,-71.8)..(167.9,-71.7)
				..controls(168,-71.7)and(168,-70.2)..(168,-68.4)..controls(168,-65.4)and(168,-64.9)..(167.7,-63)
				..controls(167.3,-60.5)and(167.1,-59.9)..(166.4,-58.4)
				..controls(166.1,-57.7)and(165.8,-56.9)..(165.7,-56.6)
				..controls(165.7,-56.2)and(165.4,-55.6)..(165.2,-55.2)
				..controls(164.6,-53.8)and(161.9,-50.7)..(161,-50.3)
				..controls(159.8,-49.7)and(157.9,-50)..(155.2,-51.2)
				..controls(154.7,-51.4)and(153.4,-51.8)..(152.2,-52.1)
				..controls(150.8,-52.4)and(149.4,-52.9)..(148.6,-53.2)
				..controls(147.8,-53.6)and(146.7,-54)..(146,-54.2)
				..controls(145.4,-54.4)and(144.5,-54.7)..(144,-54.9)
				..controls(143.6,-55.2)and(142.7,-55.5)..(141.8,-55.7)
				..controls(141,-55.9)and(139.6,-56.3)..(138.7,-56.6)
				..controls(137.8,-57)and(136.3,-57.4)..(135.3,-57.6)
				..controls(133.8,-57.9)and(133.2,-58.1)..(132.1,-58.7)
				..controls(131.1,-59.2)and(130.3,-59.4)..(129.4,-59.6)
				..controls(128.7,-59.8)and(127.8,-60)..(127.4,-60.1)
				..controls(127,-60.3)and(126,-60.5)..(125.1,-60.7)
				..controls(123.2,-61.1)and(121.7,-61.5)..(120.2,-62.1)
				..controls(119.7,-62.3)and(118.1,-62.8)..(116.7,-63.2)
				..controls(114.6,-63.9)and(114,-64.1)..(113,-64.7)
				..controls(111.9,-65.4)and(111.8,-65.4)..(111,-65.5)
				..controls(109.9,-65.5)and(109.3,-65.4)..(107.7,-64.7)
				..controls(107,-64.4)and(106,-64)..(105.3,-63.9)
				..controls(103.5,-63.5)and(102.2,-63)..(101.2,-62.5)
				..controls(100.1,-61.8)and(99.3,-61.5)..(97.8,-61.1)
				..controls(97.2,-61)and(96.2,-60.7)..(95.7,-60.4)
				..controls(95.1,-60.1)and(94.2,-59.8)..(93.5,-59.6)
				..controls(92.9,-59.4)and(91.7,-59)..(91,-58.7)..controls(90.2,-58.4)and(88.9,-58)..(88,-57.7)
				..controls(87.2,-57.5)and(86.1,-57.1)..(85.5,-56.8)
				..controls(85,-56.6)and(84.1,-56.2)..(83.5,-56)..controls(82.9,-55.7)and(81.8,-55.2)..(81,-54.9)
				..controls(80.1,-54.5)and(79.1,-54.2)..(78.7,-54.1)
				..controls(77.9,-53.8)and(77.3,-53.6)..(76.2,-52.9)
				..controls(75.7,-52.6)and(75.2,-52.4)..(74.1,-52.2)
				..controls(72.5,-51.9)and(72.4,-51.8)..(71.3,-51.1)
				..controls(69.9,-50.1)and(69.4,-49.9)..(67.9,-49.8)
				..controls(67,-49.7)and(66.6,-49.6)..(66.3,-49.4)
				..controls(66.1,-49.3)and(65.7,-49)..(65.4,-48.8)
				..controls(65.1,-48.7)and(64.2,-48.4)..(63.6,-48.2)
				..controls(62.9,-48)and(61.7,-47.6)..(60.9,-47.2)
				..controls(60.1,-46.8)and(58.9,-46.4)..(58.2,-46.2)
				..controls(57.6,-46)and(56.8,-45.7)..(56.5,-45.5)
				..controls(54.9,-44.6)and(54.6,-44.5)..(53.2,-44.2)
				..controls(51.7,-43.8)and(51,-43.6)..(49.9,-42.9)
				..controls(49.5,-42.6)and(48.9,-42.4)..(47.5,-42)
				..controls(45.7,-41.6)and(45.6,-41.5)..(44.3,-40.6)
				..controls(43.5,-40.1)and(42.8,-39.6)..(42.7,-39.4)
				..controls(42.4,-38.7)and(43,-38.1)..(44.5,-37.6)..controls(44.9,-37.5)and(46,-37)..(46.9,-36.6)
				..controls(47.9,-36.1)and(48.8,-35.8)..(49.6,-35.6)
				..controls(50.5,-35.4)and(51.3,-35.2)..(52.4,-34.6)
				..controls(53.3,-34.1)and(54.5,-33.7)..(54.9,-33.5)
				..controls(55.4,-33.4)and(56.4,-33)..(57.2,-32.6)
				..controls(58.6,-32)and(60.1,-31.5)..(61.6,-31.2)
				..controls(62.3,-31.1)and(62.7,-31)..(63.2,-30.7)
				..controls(64.1,-30)and(64.3,-29.9)..(65.3,-29.7)
				..controls(65.7,-29.6)and(66.5,-29.2)..(67.1,-28.9)
				..controls(67.6,-28.6)and(68.8,-28.2)..(69.6,-27.9)
				..controls(70.5,-27.7)and(71.6,-27.3)..(72.1,-27)
				..controls(72.6,-26.7)and(73.8,-26.2)..(74.8,-25.9)
				..controls(75.7,-25.5)and(76.7,-25.1)..(77.1,-25)
				..controls(77.4,-24.8)and(78.5,-24.4)..(79.6,-24.1)
				..controls(80.6,-23.8)and(82,-23.2)..(82.8,-22.9)
				..controls(83.5,-22.5)and(84.8,-22)..(85.8,-21.8)
				..controls(86.7,-21.5)and(87.8,-21.1)..(88.1,-20.9)
				..controls(88.4,-20.7)and(89.3,-20.3)..(90.1,-20.1)
				..controls(90.8,-19.8)and(91.9,-19.4)..(92.4,-19.1)
				..controls(93.5,-18.5)and(94.4,-18.2)..(95.4,-18.1)
				..controls(95.7,-18)and(96.4,-17.8)..(96.9,-17.5)
				..controls(97.5,-17.3)and(98.4,-16.9)..(99,-16.6)
				..controls(99.7,-16.3)and(100.9,-15.9)..(101.9,-15.6)
				..controls(102.9,-15.3)and(104.4,-14.7)..(105.3,-14.3)
				..controls(108.1,-13.2)and(108.5,-13.1)..(110.1,-13.1)
				..controls(111.4,-13.1)and(111.5,-13.1)..(112.4,-13.5)
				..controls(113.2,-14)and(115,-14.8)..(116.6,-15.4)
				..controls(116.9,-15.5)and(117.9,-15.8)..(118.8,-16.1)
				..controls(120.9,-16.6)and(121.5,-16.9)..(122.6,-17.9)
				..controls(123,-18.3)and(123.7,-18.8)..(123.9,-18.9)
				..controls(124.6,-19.4)and(126,-19.8)..(127.9,-20.2)
				..controls(129.6,-20.5)and(129.8,-20.6)..(131.1,-21.4)
				..controls(131.5,-21.7)and(132.4,-22.1)..(133.1,-22.3)
				..controls(136.1,-23.2)and(138.3,-24.1)..(140.2,-25)
				..controls(141,-25.5)and(141.5,-25.6)..(142.5,-25.8)
				..controls(145.1,-26.3)and(146,-26.6)..(149.4,-28.1)
				..controls(151.3,-28.9)and(153.5,-29.8)..(154.3,-30.1)
				..controls(155.1,-30.4)and(156.2,-30.8)..(156.6,-31)
				..controls(157.1,-31.2)and(158.2,-31.6)..(159.1,-31.8)
				..controls(160.6,-32.2)and(160.9,-32.3)..(161.7,-32.9)
				..controls(162.4,-33.3)and(162.9,-33.6)..(163.7,-33.8)
				..controls(164.3,-34)and(164.8,-34.1)..(164.9,-34.1)
				..controls(165,-34.1)and(165.6,-34.4)..(166.4,-34.8)
				..controls(167.3,-35.2)and(168.1,-35.4)..(169,-35.6)
				..controls(170,-35.8)and(170.5,-36)..(171.7,-36.6)
				..controls(172.4,-37)and(173.3,-37.4)..(173.6,-37.4)
				..controls(175.6,-37.7)and(176.2,-38)..(176.7,-39.1)
				..controls(176.9,-39.6)and(176.9,-39.7)..(176.7,-40.4)
				..controls(176.4,-41.6)and(175.8,-42.7)..(175.2,-43.4)
				..controls(174.5,-44.1)and(172.6,-45.5)..(171.7,-45.8)
				..controls(170.6,-46.2)and(169.1,-46)..(168.2,-45.4)
				..controls(167.9,-45.2)and(167.3,-44.8)..(166.9,-44.4)
				..controls(166.4,-44.1)and(165.8,-43.6)..(165.4,-43.2)
				..controls(165.1,-42.9)and(164.4,-42.3)..(163.9,-41.9)
				..controls(163.3,-41.5)and(162.7,-40.9)..(162.4,-40.5)
				..controls(161.4,-39.2)and(159.9,-38)..(158.8,-37.8)
				..controls(158.4,-37.8)and(158.3,-37.7)..(158.2,-37.2)
				..controls(158,-36.8)and(157.4,-36.1)..(157.1,-36.1)
				..controls(157,-36.1)and(156.7,-36.1)..(156.5,-36)
				..controls(156.1,-35.9)and(155.9,-35.7)..(155.7,-35.2)
				..controls(155.3,-34.6)and(155.1,-34.4)..(154,-34.1)
				..controls(153.6,-34)and(153.2,-33.7)..(152.9,-33.5)
				..controls(152.4,-32.9)and(151.6,-32.4)..(151.3,-32.4)
				..controls(150.9,-32.4)and(150.4,-32.6)..(150.4,-32.9)
				..controls(150.4,-33.2)and(150.8,-33.4)..(151.4,-33.6)
				..controls(152.1,-33.7)and(152.8,-34.1)..(153.5,-34.9)
				..controls(154.1,-35.5)and(154.3,-35.6)..(155,-35.8)
				..controls(155.7,-36)and(155.7,-36)..(156,-36.6)
				..controls(156.3,-37.3)and(156.7,-37.7)..(157.5,-37.8)
				..controls(158.1,-37.9)and(158.2,-38)..(158.4,-38.8)
				..controls(158.5,-39.4)and(158.6,-39.6)..(159.6,-40.4)
				..controls(160.2,-41)and(160.9,-41.6)..(161.3,-41.8)
				..controls(161.7,-42)and(161.9,-42.3)..(162.2,-42.8)
				..controls(162.4,-43.2)and(162.8,-43.8)..(163.1,-44.1)
				..controls(163.9,-44.8)and(165.1,-46.7)..(165.7,-48)
				..controls(166.2,-49.3)and(166.4,-49.7)..(167.2,-50.5)
				..controls(167.5,-50.8)and(167.9,-51.4)..(168.1,-51.8)
				..controls(168.3,-52.2)and(168.8,-53)..(169.1,-53.6)
				..controls(169.6,-54.3)and(169.8,-54.9)..(169.9,-55.6)
				..controls(170.4,-57.8)and(170.5,-58.1)..(171.2,-59.5)
				..controls(171.6,-60.3)and(172,-61.5)..(172.2,-62.1)
				..controls(172.3,-62.8)and(172.7,-63.7)..(172.9,-64.1)
				..controls(173.2,-64.6)and(173.4,-65.3)..(173.4,-65.6)
				..controls(173.8,-68.8)and(174.2,-70.6)..(174.9,-72.6)
				..controls(175.5,-74.5)and(175.9,-77.1)..(176.1,-80.2)
				..controls(176.3,-83.4)and(175.5,-84.4)..(173.7,-83.4)--cycle;
				\definecolor{ezs}{RGB}{247,179,26}\fill[ezs](-5.7,-156.6)--(-3.4,-157.5)--(61.9,12.7)--(59.7,13.6)--cycle;
				\definecolor{ezs}{RGB}{247,179,22}\fill[ezs](-3.8,-157.3)--(-1.5,-158.2)--(63.9,12)--(61.6,12.9)--cycle;
				\definecolor{ezs}{RGB}{248,179,19}\fill[ezs](-1.8,-158.1)--(0.4,-158.9)--(65.8,11.2)--(63.5,12.1)--cycle;
				\definecolor{ezs}{RGB}{248,178,16}\fill[ezs](0.1,-158.8)--(2.4,-159.7)--(67.7,10.5)--(65.4,11.4)--cycle;
				\definecolor{ezs}{RGB}{249,178,15}\fill[ezs](2,-159.5)--(4.3,-160.4)--(69.6,9.8)--(67.3,10.7)--cycle;
				\definecolor{ezs}{RGB}{251,177,17}\fill[ezs](3.9,-160.3)--(6.2,-161.1)--(71.5,9)--(69.2,9.9)--cycle;
				\definecolor{ezs}{RGB}{252,179,18}\fill[ezs](5.8,-161)--(8.1,-161.9)--(73.4,8.3)--(71.1,9.2)--cycle;
				\definecolor{ezs}{RGB}{252,182,17}\fill[ezs](7.7,-161.7)--(10,-162.6)--(75.3,7.6)--(73,8.5)--cycle;
				\definecolor{ezs}{RGB}{252,185,17}\fill[ezs](9.6,-162.5)--(11.9,-163.4)--(77.2,6.8)--(75,7.7)--cycle;
				\definecolor{ezs}{RGB}{251,189,18}\fill[ezs](11.5,-163.2)--(13.8,-164.1)--(79.2,6.1)--(76.9,7)--cycle;
				\definecolor{ezs}{RGB}{251,191,22}\fill[ezs](13.4,-163.9)--(15.7,-164.8)--(81.1,5.4)--(78.8,6.2)--cycle;
				\definecolor{ezs}{RGB}{252,193,25}\fill[ezs](15.4,-164.7)--(17.6,-165.6)--(83,4.6)--(80.7,5.5)--cycle;
				\definecolor{ezs}{RGB}{253,195,27}\fill[ezs](17.3,-165.4)--(19.6,-166.3)--(84.9,3.9)--(82.6,4.8)--cycle;
				\definecolor{ezs}{RGB}{253,196,31}\fill[ezs](19.2,-166.1)--(21.5,-167)--(86.8,3.2)--(84.5,4)--cycle;
				\definecolor{ezs}{RGB}{254,199,35}\fill[ezs](21.1,-166.9)--(23.4,-167.8)--(88.7,2.4)--(86.4,3.3)--cycle;
				\definecolor{ezs}{RGB}{253,201,38}\fill[ezs](23,-167.6)--(25.3,-168.5)--(90.6,1.7)--(88.3,2.6)--cycle;
				\definecolor{ezs}{RGB}{253,203,40}\fill[ezs](24.9,-168.3)--(27.2,-169.2)--(92.5,1)--(90.3,1.8)--cycle;
				\definecolor{ezs}{RGB}{253,205,42}\fill[ezs](26.8,-169.1)--(29.1,-170)--(94.4,0.2)--(92.2,1.1)--cycle;
				\definecolor{ezs}{RGB}{253,207,45}\fill[ezs](28.7,-169.8)--(31,-170.7)--(96.4,-0.5)--(94.1,0.4)--cycle;
				\definecolor{ezs}{RGB}{253,210,49}\fill[ezs](30.7,-170.5)--(32.9,-171.4)--(98.3,-1.2)--(96,-0.4)--cycle;
				\definecolor{ezs}{RGB}{253,212,54}\fill[ezs](32.6,-171.3)--(34.9,-172.2)--(100.2,-2)--(97.9,-1.1)--cycle;
				\definecolor{ezs}{RGB}{253,214,59}\fill[ezs](34.5,-172)--(36.8,-172.9)--(102.1,-2.7)--(99.8,-1.8)--cycle;
				\definecolor{ezs}{RGB}{253,216,63}\fill[ezs](36.4,-172.7)--(38.7,-173.6)--(104,-3.4)--(101.7,-2.6)--cycle;
				\definecolor{ezs}{RGB}{253,218,65}\fill[ezs](38.3,-173.5)--(40.6,-174.4)--(105.9,-4.2)--(103.6,-3.3)--cycle;
				\definecolor{ezs}{RGB}{254,220,66}\fill[ezs](40.2,-174.2)--(42.5,-175.1)--(107.8,-4.9)--(105.5,-4)--cycle;
				\definecolor{ezs}{RGB}{254,222,68}\fill[ezs](42.1,-175)--(44.4,-175.8)--(109.7,-5.6)--(107.5,-4.8)--cycle;
				\definecolor{ezs}{RGB}{254,224,71}\fill[ezs](44,-175.7)--(46.3,-176.6)--(111.7,-6.4)--(109.4,-5.5)--cycle;
				\definecolor{ezs}{RGB}{254,226,75}\fill[ezs](46,-176.4)--(48.2,-177.3)--(113.6,-7.1)--(111.3,-6.2)--cycle;
				\definecolor{ezs}{RGB}{254,228,79}\fill[ezs](47.9,-177.2)--(50.1,-178)--(115.5,-7.8)--(113.2,-7)--cycle;
				\definecolor{ezs}{RGB}{254,230,82}\fill[ezs](49.8,-177.9)--(52.1,-178.8)--(117.4,-8.6)--(115.1,-7.7)--cycle;
				\definecolor{ezs}{RGB}{254,232,85}\fill[ezs](51.7,-178.6)--(54,-179.5)--(119.3,-9.3)--(117,-8.4)--cycle;
				\definecolor{ezs}{RGB}{254,234,88}\fill[ezs](53.6,-179.4)--(55.9,-180.2)--(121.2,-10)--(118.9,-9.2)--cycle;
				\definecolor{ezs}{RGB}{254,237,91}\fill[ezs](55.5,-180.1)--(57.8,-181)--(123.1,-10.8)--(120.8,-9.9)--cycle;
				\definecolor{ezs}{RGB}{254,239,94}\fill[ezs](57.4,-180.8)--(59.7,-181.7)--(125,-11.5)--(122.8,-10.6)--cycle;
				\definecolor{ezs}{RGB}{254,241,97}\fill[ezs](59.3,-181.6)--(61.6,-182.4)--(127,-12.2)--(124.7,-11.4)--cycle;
				\definecolor{ezs}{RGB}{254,243,101}\fill[ezs](61.2,-182.3)--(63.5,-183.2)--(128.9,-13)--(126.6,-12.1)--cycle;
				\definecolor{ezs}{RGB}{254,245,103}\fill[ezs](63.2,-183)--(65.4,-183.9)--(130.8,-13.7)--(128.5,-12.8)--cycle;
				\definecolor{ezs}{RGB}{254,245,105}\fill[ezs](65.1,-183.8)--(67.4,-184.6)--(132.7,-14.4)--(130.4,-13.6)--cycle;
				\definecolor{ezs}{RGB}{254,245,105}\fill[ezs](67,-184.5)--(69.3,-185.4)--(134.6,-15.2)--(132.3,-14.3)--cycle;
				\definecolor{ezs}{RGB}{254,244,103}\fill[ezs](68.9,-185.2)--(71.2,-186.1)--(136.5,-15.9)--(134.2,-15)--cycle;
				\definecolor{ezs}{RGB}{253,243,101}\fill[ezs](70.8,-186)--(73.1,-186.8)--(138.4,-16.6)--(136.1,-15.8)--cycle;
				\definecolor{ezs}{RGB}{253,241,98}\fill[ezs](72.7,-186.7)--(75,-187.6)--(140.3,-17.4)--(138,-16.5)--cycle;
				\definecolor{ezs}{RGB}{253,239,95}\fill[ezs](74.6,-187.4)--(76.9,-188.3)--(142.2,-18.1)--(140,-17.2)--cycle;
				\definecolor{ezs}{RGB}{253,236,92}\fill[ezs](76.5,-188.2)--(78.8,-189)--(144.2,-18.8)--(141.9,-18)--cycle;
				\definecolor{ezs}{RGB}{253,234,88}\fill[ezs](78.5,-188.9)--(80.7,-189.8)--(146.1,-19.6)--(143.8,-18.7)--cycle;
				\definecolor{ezs}{RGB}{253,232,85}\fill[ezs](80.4,-189.6)--(82.7,-190.5)--(148,-20.3)--(145.7,-19.4)--cycle;
				\definecolor{ezs}{RGB}{253,230,82}\fill[ezs](82.3,-190.4)--(84.6,-191.2)--(149.9,-21)--(147.6,-20.2)--cycle;
				\definecolor{ezs}{RGB}{253,228,78}\fill[ezs](84.2,-191.1)--(86.5,-192)--(151.8,-21.8)--(149.5,-20.9)--cycle;
				\definecolor{ezs}{RGB}{253,226,74}\fill[ezs](86.1,-191.8)--(88.4,-192.7)--(153.7,-22.5)--(151.4,-21.6)--cycle;
				\definecolor{ezs}{RGB}{253,224,70}\fill[ezs](88,-192.6)--(90.3,-193.4)--(155.6,-23.3)--(153.3,-22.4)--cycle;
				\definecolor{ezs}{RGB}{253,222,67}\fill[ezs](89.9,-193.3)--(92.2,-194.2)--(157.5,-24)--(155.3,-23.1)--cycle;
				\definecolor{ezs}{RGB}{253,220,65}\fill[ezs](91.8,-194)--(94.1,-194.9)--(159.5,-24.7)--(157.2,-23.8)--cycle;
				\definecolor{ezs}{RGB}{253,217,63}\fill[ezs](93.7,-194.8)--(96,-195.6)--(161.4,-25.5)--(159.1,-24.6)--cycle;
				\definecolor{ezs}{RGB}{253,215,60}\fill[ezs](95.7,-195.5)--(97.9,-196.4)--(163.3,-26.2)--(161,-25.3)--cycle;
				\definecolor{ezs}{RGB}{253,213,56}\fill[ezs](97.6,-196.2)--(99.9,-197.1)--(165.2,-26.9)--(162.9,-26)--cycle;
				\definecolor{ezs}{RGB}{252,211,53}\fill[ezs](99.5,-197)--(101.8,-197.8)--(167.1,-27.7)--(164.8,-26.8)--cycle;
				\definecolor{ezs}{RGB}{252,209,49}\fill[ezs](101.4,-197.7)--(103.7,-198.6)--(169,-28.4)--(166.7,-27.5)--cycle;
				\definecolor{ezs}{RGB}{252,207,45}\fill[ezs](103.3,-198.4)--(105.6,-199.3)--(170.9,-29.1)--(168.6,-28.2)--cycle;
				\definecolor{ezs}{RGB}{253,205,42}\fill[ezs](105.2,-199.2)--(107.5,-200)--(172.8,-29.9)--(170.6,-29)--cycle;
				\definecolor{ezs}{RGB}{253,202,40}\fill[ezs](107.1,-199.9)--(109.4,-200.8)--(174.7,-30.6)--(172.5,-29.7)--cycle;
				\definecolor{ezs}{RGB}{253,200,38}\fill[ezs](109,-200.6)--(111.3,-201.5)--(176.7,-31.3)--(174.4,-30.4)--cycle;
				\definecolor{ezs}{RGB}{253,198,35}\fill[ezs](111,-201.4)--(113.2,-202.2)--(178.6,-32.1)--(176.3,-31.2)--cycle;
				\definecolor{ezs}{RGB}{253,196,31}\fill[ezs](112.9,-202.1)--(115.2,-203)--(180.5,-32.8)--(178.2,-31.9)--cycle;
				\definecolor{ezs}{RGB}{253,194,28}\fill[ezs](114.8,-202.8)--(117.1,-203.7)--(182.4,-33.5)--(180.1,-32.6)--cycle;
				\definecolor{ezs}{RGB}{253,192,26}\fill[ezs](116.7,-203.6)--(119,-204.5)--(184.3,-34.3)--(182,-33.4)--cycle;
				\definecolor{ezs}{RGB}{253,189,24}\fill[ezs](118.6,-204.3)--(120.9,-205.2)--(186.2,-35)--(183.9,-34.1)--cycle;
				\definecolor{ezs}{RGB}{253,186,20}\fill[ezs](120.5,-205)--(122.8,-205.9)--(188.1,-35.7)--(185.8,-34.8)--cycle;
				\definecolor{ezs}{RGB}{253,184,17}\fill[ezs](122.4,-205.8)--(124.7,-206.7)--(190,-36.5)--(187.8,-35.6)--cycle;
				\definecolor{ezs}{RGB}{252,183,15}\fill[ezs](124.3,-206.5)--(126.6,-207.4)--(192,-37.2)--(189.7,-36.3)--cycle;
				\definecolor{ezs}{RGB}{252,181,14}\fill[ezs](126.3,-207.2)--(128.5,-208.1)--(193.9,-37.9)--(191.6,-37.1)--cycle;
				\definecolor{ezs}{RGB}{252,179,13}\fill[ezs](128.2,-208)--(130.4,-208.9)--(195.8,-38.7)--(193.5,-37.8)--cycle;
				\definecolor{ezs}{RGB}{252,178,12}\fill[ezs](130.1,-208.7)--(132.4,-209.6)--(197.7,-39.4)--(195.4,-38.5)--cycle;
				\definecolor{ezs}{RGB}{251,178,10}\fill[ezs](132,-209.4)--(134.3,-210.3)--(199.6,-40.1)--(197.3,-39.3)--cycle;
				\definecolor{ezs}{RGB}{251,178,10}\fill[ezs](133.9,-210.2)--(136.2,-211.1)--(201.5,-40.9)--(199.2,-40)--cycle;
				\definecolor{ezs}{RGB}{250,177,11}\fill[ezs](135.8,-210.9)--(138.1,-211.8)--(203.4,-41.6)--(201.1,-40.7)--cycle;
				\definecolor{ezs}{RGB}{250,178,11}\fill[ezs](137.7,-211.6)--(140,-212.5)--(205.3,-42.3)--(203.1,-41.5)--cycle;
				\definecolor{ezs}{RGB}{248,180,10}\fill[ezs](139.6,-212.4)--(141.9,-213.3)--(207.3,-43.1)--(205,-42.2)--cycle;
				\definecolor{ezs}{RGB}{247,181,9}\fill[ezs](141.5,-213.1)--(143.8,-214)--(209.2,-43.8)--(206.9,-42.9)--cycle;
				\definecolor{ezs}{RGB}{246,182,9}\fill[ezs](143.5,-213.8)--(145.7,-214.7)--(211.1,-44.5)--(208.8,-43.7)--cycle;
				\definecolor{ezs}{RGB}{246,182,10}\fill[ezs](145.4,-214.6)--(147.7,-215.5)--(213,-45.3)--(210.7,-44.4)--cycle;
			\end{scope}
			% chữ EDUZY (Times đậm, làm dày nét bằng 5 lớp lệch nhẹ; căn đúng khung chữ logo gốc)
			\foreach \dx/\dy in {-0.7/0,0.7/0,0/-0.3,0/0.3,0/0}
			\node[anchor=center,inner sep=0pt,outer sep=0pt,text=ezText,align=center] at (109.6+\dx,-193.6+\dy)
			{\resizebox{1.01cm}{0.22cm}{\fontencoding{T1}\fontfamily{ptm}\fontseries{b}\selectfont EDUZY}};
		\end{tikzpicture}%
	}%
}
%====================== THIẾT LẬP THÂN ĐỀ ======================
%--- Môi trường tự luận {bt}: ex_test 3.3.4 chưa khai báo sẵn nên tạo ở đây
%    (ex_test đã tự móc \fixshowans và \listenumerate cho bt)
\makeatletter
\@ifundefined{bt}{\newtheorem{bt}{\color{gold1}\namebt}}{}
\makeatother
\renewcommand{\sectionmark}[1]{\markright{\thesection\ #1}}
%--- Khung tiêu đề (logo + tối đa 4 dòng; dòng 3, 4 để trống {} thì bỏ qua)
%  Đề thi:    \khungtieude{THI THỬ THPT QUỐC GIA NĂM 2026}{Môn Toán -- Đề số 21}
%                         {(Đề thi gồm: ...)}{Thời gian làm bài: 90 phút (không kể thời gian phát đề).}
%  Chuyên đề: \khungtieude{CHUYÊN ĐỀ 3 -- DÃY SỐ VÀ CÔNG THỨC TRUY HỒI}
%                         {Bồi dưỡng học sinh giỏi Toán THPT}{...}{}
\newcommand{\khungtieude}[4]{%
	\begin{tcolorbox}[colframe=gold1,colback=gold5!45,boxrule=0.5mm]
		\begin{minipage}[c]{3cm}
			\centerline{\eduzylogo[2.4cm]}
		\end{minipage}\hspace{0.5cm}
		\begin{minipage}[c]{12.5cm}
			\vskip 0.15 cm
			\centerline{\ind{#1}}
			\vskip 0.15 cm
			\centerline{\textbf{#2}}
			\if\relax\detokenize{#3}\relax\else
				\vskip 0.15 cm
				\centerline{\textit{#3}}
			\fi
			\if\relax\detokenize{#4}\relax\else
				\vskip 0.15 cm
				\centerline{\textit{#4}}
			\fi
			\vskip 0.15 cm
		\end{minipage}
	\end{tcolorbox}
}
%====================== BỔ SUNG CHO BÀI HỌC, CHUYÊN ĐỀ ======================
\usepackage[shortlabels]{enumitem}
\usepackage{tabularx}
%--- Danh sách gọn: \begin{enumerate}[gon] ... \end{enumerate}
\SetEnumitemKey{gon}{leftmargin=*,itemsep=1pt,parsep=0pt,topsep=3pt}
%--- Khung lý thuyết: \begin{khung}{} ... \end{khung}; tiêu đề chỉ dùng khi thật ngắn ("Mẹo", "Lỗi thường gặp")
\newtcolorbox{khung}[1]{breakable,enhanced,colframe=gold1,colback=gold5!45,
	boxrule=0.4mm,fonttitle=\bfseries\sffamily,coltitle=white,colbacktitle=gold1,
	attach boxed title to top left={xshift=6mm,yshift=-3mm},
	boxed title style={arc=1mm,boxrule=0.3mm},title={#1}}
%--- Tiêu đề mục nhỏ trong một dạng: \muc{Dấu hiệu}, \muc{Phương pháp}, \muc{Bài tự luyện}
\makeatletter
\newcommand{\muc}[1]{\par\medskip\noindent\indam{#1}\par\nobreak\smallskip\@afterindentfalse\@afterheading}
\makeatother
%--- Ô công thức cần nhớ (căn giữa, viền vàng, vừa khít nội dung)
%     \begin{center}\congthuc{$Q(n)=(n-1)!$}\end{center}
\newtcbox{\congthuc}{on line,colframe=gold1,colback=gold5!45,boxrule=0.5pt,arc=2pt,
	left=8pt,right=8pt,top=4pt,bottom=4pt}
%--- Công tắc lời giải bài tự luyện (ví dụ mẫu luôn dùng \loigiai)
%     \loigiaiTL[{đáp số}]{lời giải đầy đủ}
%     \TLfull -> in lời giải đầy đủ;  \TLdapso -> chỉ in một dòng "Đáp số";  \TLan -> ẩn hoàn toàn
%     (\CheDo là dethi thì mọi lời giải và đáp số đều ẩn)
\makeatletter
\newif\ifTLhien   \TLhientrue
\newif\ifTLdapso  \TLdapsofalse
\newcommand{\TLfull}{\TLhientrue\TLdapsofalse}
\newcommand{\TLdapso}{\TLhientrue\TLdapsotrue}
\newcommand{\TLan}{\TLhienfalse}
\newcommand{\loigiaiTL}[2][]{%
	\ifTLhien
		\ifTLdapso
			\def\@tlarg{#1}%
			\ifx\@tlarg\@empty\else
				\ifdefstring{\kindOPTN}{dethi}{}{\par\noindent{\bfseries\color{gold1}Đáp số:} #1\par}%
			\fi
		\else
			\loigiai{#2}%
		\fi
	\fi}
\makeatother
\TLdapso% mặc định: chỉ in đáp số bài tự luyện
%====================== BẢNG ĐÚNG/SAI (\choiceTFt) ======================
%  Hàng tiêu đề nền nâu đồng chữ trắng (Phát biểu | Đ | S), các hàng tô xen kẽ kem / trắng,
%  đường kẻ vàng kim, nhãn a) b) c) d) màu nâu đồng. Bản có đáp án: cột được chọn có vòng tròn
%  nâu đồng với dấu tích trắng (câu đúng tích ở cột Đúng, câu sai tích ở cột Sai).
%  Tiêu đề ghi đủ "Đúng | Sai": \OPTN{viettat=0} sau \input (cột Đúng/Sai rộng ra, phát biểu hẹp lại).
%  Dùng bảng gốc của ex_test: \def\BangDSgoc{1} trước %============================================================
%  PREAMBLE EDUZY v2 - dùng chung cho đề thi, bài học, chuyên đề (gói ex_test)
%  Khác bản preamble-eduzy.tex ở đầu trang, chân trang: dải nền nâu đồng tràn ngang,
%  logo + Eduzy.com.vn, ô chữ vàng kem cạnh xiên, chân trang "Trang (số)";
%  logo Eduzy chìm giữa mọi trang (tắt: \logochimtat, bật lại: \logochimbat).
%  Cách dùng ở file chính (xem mau-chuan.tex, mau-chuyen-de.tex):
%      \def\CheDo{loigiai}          % loigiai | dethi | solcolor | color | book
%      \def\tenchuong{<dòng chữ trên đầu trang>}
%      \def\thuonghieu{<chữ ở chân trang lẻ>}   % không khai báo thì dùng mặc định
%      \def\tacgia{<chữ ở chân trang chẵn>}      % không khai báo thì dùng mặc định
%      \def\khauhieu{<chữ cạnh logo ở đầu trang>}  % tùy chọn, mặc định Eduzy.com.vn
%      \input{preamble-eduzy-v2}
%      \TLdapso                      % tài liệu: \TLfull | \TLdapso | \TLan
%      \begin{document} ... \end{document}
%  Logo Eduzy vẽ bằng TikZ ngay trong file (\eduzylogo), không cần file ảnh.
%  Biên dịch: pdfLaTeX (máy có gói vietnam) hoặc XeLaTeX.
%============================================================
\documentclass[12pt,a4paper,twoside]{article}%<-- twoside: bat dau trang/chan trang chan-le
\usepackage{iftex}
\usepackage{amsmath,amssymb,tikz,mathrsfs,tkz-tab}
\usepackage{graphicx}
\usepackage{pgfplots}% đồ thị hàm số
\pgfplotsset{compat=1.17}
%\usepackage{mathptmx}
\usepackage{mathpazo}
\ifXeTeX% XeLaTeX: font chữ có đủ dấu tiếng Việt, giữ nguyên font toán mathpazo
	\usepackage[no-math]{fontspec}
	\setmainfont{TeX Gyre Pagella}
	\newfontfamily\EZqag{TeX Gyre Adventor}[NFSSFamily=qag,Ligatures=TeX]
\fi
\usepackage[many]{tcolorbox}
\usetikzlibrary{math,through,calc,intersections,angles,quotes,shapes,shapes.geometric,arrows,patterns,matrix,chains,arrows.meta,decorations.shapes,decorations.fractals,decorations.markings,shadows,decorations.pathreplacing}
\usepackage[top=1.6cm, bottom=1.8cm, left=1.8cm, right=1.8cm]{geometry}% v2: top 1.6cm (bản cũ 1.5cm) để chữ không sát dải đầu trang
\providecommand{\CheDo}{loigiai}% loigiai | dethi | solcolor | color | book
\expandafter\usepackage\expandafter[\CheDo]{ex_test}
% \tenchuong: khai báo ở file chính (trước \input) hoặc \renewcommand sau \input
%============================================================
%  [GỘP TỪ eduzy-de.tex] EDUZY - MÀU NÂU ĐỒNG / VÀNG KIM
%  + ĐẦU TRANG, CHÂN TRANG KIỂU DẢI NGANG (v2), ĐỐI XỨNG TRANG CHẴN / TRANG LẺ.
%  (không cần file eduzy-de.tex riêng nữa)
%  Khối này phải đặt SAU \usepackage[...]{ex_test} và SAU \def\tenchuong{...}
%============================================================
\usepackage{lastpage}
\usepackage{needspace}
%--- DÀN TRANG ĐỀ THI: tách câu (THÊM 05/11/2026; chỉ thêm lệnh mới, KHÔNG đổi hành vi mặc định) ---
%  Hai công tắc, mặc định tắt (= preamble chạy y như cũ). Bật bằng cách đặt TRƯỚC \input{preamble-eduzy-v2}
%  trong file đề (tools/dan_trang/dan_trang.py tự ghi khi đặt \NeedspaceCau):
%    \def\TachCauDS{1}   câu Đúng/Sai được tách: câu dẫn (kèm hình) ở cuối trang, bảng Phát biểu sang trang sau
%                        (câu dẫn và hình không tách nhau, bảng không tách giữa hai trang)
%    \def\TachCauLG{1}   bản lời giải (\CheDo khác dethi): phần đề của câu cao hơn 1/3 trang được phép tách
%  Lệnh dùng trong file đề (giá trị do tools/dan_trang/dan_trang.py đo ở bản dethi):
%    \NeedspaceCau[<cao câu dẫn + hình>]{<cao cả khối câu>}   ngay trên \begin{ex} (câu đầu phần: trên tiêu đề phần)
%        công tắc tắt: đúng bằng \Needspace{<cao cả khối câu>} như cũ;
%        \TachCauDS=1 và có [..] (câu Đúng/Sai được chọn tách): chỉ giữ câu dẫn + hình;
%        \TachCauLG=1 ở bản lời giải: đòi tối đa 1/3 \textheight.
%    \NeedspaceBang{<cao bảng Phát biểu>}%   ngay trước \choiceTFt: \TachCauDS=1 thì giữ trọn bảng trên một trang;
%        công tắc tắt: không làm gì.
\providecommand{\TachCauDS}{0}
\providecommand{\TachCauLG}{0}
\newlength{\EZnsCau}
\newcommand{\NeedspaceCau}[2][]{%
	\setlength{\EZnsCau}{#2}%
	\ifnum\TachCauDS=1 \if\relax\detokenize{#1}\relax\else\setlength{\EZnsCau}{#1}\fi\fi
	\ifnum\TachCauLG=1 \ifdefstring{\kindOPTN}{dethi}{}{\ifdim\EZnsCau>0.3333\textheight\setlength{\EZnsCau}{0.3333\textheight}\fi}\fi
	\Needspace{\EZnsCau}}
\newcommand{\NeedspaceBang}[1]{\ifnum\TachCauDS=1 \Needspace{#1}\fi}

%--- Thang màu: đậm -> nhạt
\definecolor{gold1}{RGB}{126,96,0}    % nâu đồng đậm  #7E6000
\definecolor{gold2}{RGB}{192,144,3}   % vàng kim      #C09003
\definecolor{gold3}{RGB}{255,215,99}  % vàng nhạt     #FFD763
\definecolor{gold4}{RGB}{255,231,153} % vàng rất nhạt #FFE799
\definecolor{gold5}{RGB}{255,243,205} % vàng kem      #FFF3CD
%--- Ghi đè tông xanh/đỏ cũ
\definecolor{myblue}{RGB}{126,96,0}
\colorlet{blue}{gold1}
\colorlet{cyan}{gold2}
\colorlet{teal}{gold2}

%--- Tên chương mặc định (file đề nên tự khai báo lại)
\providecommand{\tenchuong}{}
%--- Chân trang
\providecommand{\thuonghieu}{TOÁN 12 -- Eduzy.com.vn}
\providecommand{\tacgia}{Tô Gia Bảo -- 0353703783}

%--- Chữ cạnh logo ở đầu trang
\providecommand{\khauhieu}{Eduzy.com.vn}

%====================== THÔNG SỐ ĐẦU TRANG, CHÂN TRANG (v2) ======================
%  Dải nền đậm tràn ngang trang; ô chữ vàng kem có cạnh xiên, kèm 2 vạch xiên chuyển màu.
%  Đầu trang lẻ : logo + \khauhieu (trái)   | ô chữ \tenchuong (phải)
%  Chân trang lẻ: ô chữ \thuonghieu (trái)  | Trang (số) (phải)
%  Trang chẵn soi gương trang lẻ; chân trang chẵn ghi \tacgia.
%--- Màu: chỉ dùng thang vàng kim Eduzy
\colorlet{ezNen}{gold1}% dải nền đậm
\colorlet{ezVach1}{gold3}% vạch xiên sát dải nền (đậm hơn)
\colorlet{ezVach2}{gold4}% vạch xiên sát ô chữ (nhạt hơn)
\colorlet{ezO}{gold5}% ô chữ vàng kem
%--- Kích thước
\newlength{\EZdai}\setlength{\EZdai}{1.1cm}% chiều cao dải nền
\newlength{\EZo}\setlength{\EZo}{0.68cm}% chiều cao ô chữ
\newlength{\EZxien}\setlength{\EZxien}{0.55cm}% độ lệch ngang của cạnh xiên
\newlength{\EZvach}\setlength{\EZvach}{0.75cm}% bề rộng mỗi vạch xiên
\newlength{\EZkhe}\setlength{\EZkhe}{0.12cm}% khe trắng giữa vạch và ô chữ
\newlength{\EZdem}\setlength{\EZdem}{0.4cm}% khoảng đệm từ chữ tới cạnh xiên
\newlength{\EZle}\setlength{\EZle}{1.8cm}% lề trái/phải của chữ (bằng lề thân bài)
\newlength{\EZrongmin}\setlength{\EZrongmin}{4.5cm}% ô chữ rộng tối thiểu
\newlength{\EZrong}% bề rộng chữ trong ô (tính tự động)
%--- Font
\newcommand{\EZfontO}{\fontsize{10.5pt}{12pt}\selectfont\bfseries}% chữ trong ô
\newcommand{\EZfontKH}{\fontfamily{qag}\fontsize{11pt}{12pt}\selectfont\bfseries}% \khauhieu cạnh logo
\newcommand{\EZfontTrang}{\fontsize{11.5pt}{12pt}\selectfont\bfseries}% chữ "Trang" và số trang
%--- Đo bề rộng chữ trong ô
\newcommand{\EZdorong}[1]{%
	\settowidth{\EZrong}{\EZfontO #1}%
	\ifdim\EZrong<\EZrongmin\setlength{\EZrong}{\EZrongmin}\fi}
%--- Logo + chữ trắng \khauhieu (dựng sẵn một lần khi vào \begin{document})
\newsavebox{\EZlogobox}
\AtBeginDocument{\sbox{\EZlogobox}{%
		\raisebox{-0.5\height}{\eduzylogo[1cm]}\hspace{0.15cm}%
		\raisebox{-0.5\height}{\color{white}\EZfontKH\khauhieu}}}
%--- Hai vạch xiên + khe trắng. #1: điểm gốc ở mép dưới, #2: hướng (1 sang phải, -1 sang trái),
%    #3: hướng nghiêng (1 kiểu "/", -1 kiểu "\")
\newcommand{\EZvachxien}[3]{%
	\fill[white] (#1) --++({(#3)*\EZxien},\EZo) --++({(#2)*\EZkhe},0) --++({-(#3)*\EZxien},-\EZo) --cycle;
	\fill[ezVach2] ([xshift=(#2)*\EZkhe]#1) --++({(#3)*\EZxien},\EZo) --++({(#2)*\EZvach},0) --++({-(#3)*\EZxien},-\EZo) --cycle;
	\fill[ezVach1] ([xshift=(#2)*(\EZkhe+\EZvach)]#1) --++({(#3)*\EZxien},\EZo) --++({(#2)*\EZvach},0) --++({-(#3)*\EZxien},-\EZo) --cycle;
	\draw[white,line width=0.6pt] ([xshift=(#2)*(\EZkhe+\EZvach)]#1) --++({(#3)*\EZxien},\EZo);
}
%--- Chữ trong ô (căn giữa theo chiều cao ô, không lệch vì dấu tiếng Việt)
\tikzset{ezchu/.style={inner sep=0pt,text=ezNen,text height=0.72em,text depth=0.22em}}

%====================== LOGO CHÌM GIỮA TRANG ======================
%  Vẽ cùng đầu trang nên nằm dưới chữ của thân bài, có trên mọi trang.
%  Tắt: \logochimtat (từ trang đó trở đi), bật lại: \logochimbat.
%  Đổi cỡ, độ đậm: khai báo trước \input, ví dụ \def\EZchimcao{9cm} \def\EZchimdo{0.1}
\providecommand{\EZchimcao}{10cm}% chiều cao logo chìm
\providecommand{\EZchimdo}{0.12}% độ đậm (0 = trong suốt, 1 = rõ hẳn)
\newif\ifEZchim \EZchimtrue
\newcommand{\logochimtat}{\EZchimfalse}
\newcommand{\logochimbat}{\EZchimtrue}
%  Dựng sẵn logo mờ một lần khi vào \begin{document}; transparency group làm mờ cả logo
%  một lượt nên không lộ vệt ở chỗ các dải màu chồng nhau.
\newsavebox{\EZchimtho}
\newsavebox{\EZchimbox}
\AtBeginDocument{%
	\sbox{\EZchimtho}{\eduzylogo[\EZchimcao]}%
	\sbox{\EZchimbox}{\begin{tikzpicture}
			\begin{scope}[transparency group,opacity=\EZchimdo]
				\node[inner sep=0pt] {\usebox{\EZchimtho}};
			\end{scope}
	\end{tikzpicture}}}
\newcommand{\EZvechim}{%
	\ifEZchim
		\node[inner sep=0pt] at (current page.center) {\usebox{\EZchimbox}};
	\fi}

%====================== ĐẦU TRANG - TRANG LẺ ======================
%  (logo bên trái, ô chữ \tenchuong bên phải)
\newcommand{\myfancyheadOdd}{%
	\EZdorong{\tenchuong}%
	\begin{tikzpicture}[remember picture,overlay]
		\coordinate (T) at (current page.north west);
		\coordinate (P) at (current page.north east);
		%--- logo chìm giữa trang
		\EZvechim
		%--- dải nền
		\fill[ezNen] (T) rectangle ([yshift=-\EZdai]P);
		%--- ô chữ: X là góc dưới bên trái, cạnh trái xiên kiểu "/"
		\coordinate (X) at ([xshift=-\EZle-\EZrong-\EZdem-\EZxien,yshift=-\EZdai]P);
		\fill[ezO] (X) --++(\EZxien,\EZo) -- ([yshift=-\EZdai+\EZo]P) -- ([yshift=-\EZdai]P) --cycle;
		\EZvachxien{X}{-1}{1}
		\node[ezchu,anchor=east] at ([xshift=-\EZle,yshift=-\EZdai+0.5*\EZo]P) {\EZfontO\tenchuong};
		%--- logo + tên
		\node[inner sep=0pt,anchor=west] at ([xshift=\EZle,yshift=-0.5*\EZdai]T) {\usebox{\EZlogobox}};
	\end{tikzpicture}%
}

%====================== ĐẦU TRANG - TRANG CHẴN =====================
%  (soi gương: ô chữ \tenchuong bên trái, logo bên phải)
\newcommand{\myfancyheadEven}{%
	\EZdorong{\tenchuong}%
	\begin{tikzpicture}[remember picture,overlay]
		\coordinate (T) at (current page.north west);
		\coordinate (P) at (current page.north east);
		%--- logo chìm giữa trang
		\EZvechim
		%--- dải nền
		\fill[ezNen] (T) rectangle ([yshift=-\EZdai]P);
		%--- ô chữ: X là góc dưới bên phải, cạnh phải xiên kiểu "\"
		\coordinate (X) at ([xshift=\EZle+\EZrong+\EZdem+\EZxien,yshift=-\EZdai]T);
		\fill[ezO] (X) --++(-\EZxien,\EZo) -- ([yshift=-\EZdai+\EZo]T) -- ([yshift=-\EZdai]T) --cycle;
		\EZvachxien{X}{1}{-1}
		\node[ezchu,anchor=west] at ([xshift=\EZle,yshift=-\EZdai+0.5*\EZo]T) {\EZfontO\tenchuong};
		%--- logo + tên
		\node[inner sep=0pt,anchor=east] at ([xshift=-\EZle,yshift=-0.5*\EZdai]P) {\usebox{\EZlogobox}};
	\end{tikzpicture}%
}

%====================== CHÂN TRANG - TRANG LẺ ======================
%  (ô chữ \thuonghieu bên trái, "Trang (số)" bên phải)
\newcommand{\myfancyfootOdd}{%
	\EZdorong{\thuonghieu}%
	\begin{tikzpicture}[remember picture,overlay]
		\coordinate (T) at (current page.south west);
		\coordinate (P) at (current page.south east);
		%--- dải nền
		\fill[ezNen] (T) rectangle ([yshift=\EZdai]P);
		%--- ô chữ sát mép trên dải: X là góc dưới bên phải, cạnh phải xiên kiểu "/"
		\coordinate (X) at ([xshift=\EZle+\EZrong+\EZdem,yshift=\EZdai-\EZo]T);
		\fill[ezO] ([yshift=\EZdai-\EZo]T) -- (X) --++(\EZxien,\EZo) -- ([yshift=\EZdai]T) --cycle;
		\EZvachxien{X}{1}{1}
		\node[ezchu,anchor=west] at ([xshift=\EZle,yshift=\EZdai-0.5*\EZo]T) {\EZfontO\thuonghieu};
		%--- số trang
		\node[circle,draw=white,line width=0.9pt,text=white,inner sep=0pt,minimum size=0.62cm,anchor=east]
		(so) at ([xshift=-\EZle,yshift=0.5*\EZdai]P) {\EZfontTrang\thepage};
		\node[inner sep=0pt,text=white,anchor=base east] (tr) at ([xshift=-0.15cm]so.base west) {\EZfontTrang Trang};
	\end{tikzpicture}%
}

%====================== CHÂN TRANG - TRANG CHẴN ====================
%  (soi gương: "Trang (số)" bên trái, ô chữ \tacgia bên phải)
\newcommand{\myfancyfootEven}{%
	\EZdorong{\tacgia}%
	\begin{tikzpicture}[remember picture,overlay]
		\coordinate (T) at (current page.south west);
		\coordinate (P) at (current page.south east);
		%--- dải nền
		\fill[ezNen] (T) rectangle ([yshift=\EZdai]P);
		%--- ô chữ sát mép trên dải: X là góc dưới bên trái, cạnh trái xiên kiểu "\"
		\coordinate (X) at ([xshift=-\EZle-\EZrong-\EZdem,yshift=\EZdai-\EZo]P);
		\fill[ezO] ([yshift=\EZdai-\EZo]P) -- (X) --++(-\EZxien,\EZo) -- ([yshift=\EZdai]P) --cycle;
		\EZvachxien{X}{-1}{-1}
		\node[ezchu,anchor=east] at ([xshift=-\EZle,yshift=\EZdai-0.5*\EZo]P) {\EZfontO\tacgia};
		%--- số trang
		%    (vòng tròn ẩn so0 nằm giữa dải để lấy đường chân chữ, giống trang lẻ)
		\node[circle,inner sep=0pt,minimum size=0.62cm] (so0) at ([yshift=0.5*\EZdai]T) {\EZfontTrang\phantom{\thepage}};
		\node[inner sep=0pt,text=white,anchor=base west] (tr) at ($(so0.base -| T)+(\EZle,0)$) {\EZfontTrang Trang};
		\node[circle,draw=white,line width=0.9pt,text=white,inner sep=0pt,minimum size=0.62cm,anchor=base west]
		(so) at ([xshift=0.15cm]tr.base east) {\EZfontTrang\thepage};
	\end{tikzpicture}%
}

%--- Giữ tương thích với các file cũ đang gọi \myfancyhead, \myfancyfoot
\newcommand{\myfancyhead}{\myfancyheadOdd}
\newcommand{\myfancyfoot}{\myfancyfootOdd}

%======================= GẮN VÀO TRANG =============================
\pagestyle{fancy}
\fancyhf{}
\fancyhead[LO]{\myfancyheadOdd}
\fancyhead[RE]{\myfancyheadEven}
\fancyfoot[LO]{\myfancyfootOdd}
\fancyfoot[RE]{\myfancyfootEven}
\renewcommand{\headrulewidth}{0pt}
\renewcommand{\footrulewidth}{0pt}
\fancypagestyle{plain}{%
	\fancyhf{}
	\fancyhead[LO]{\myfancyheadOdd}
	\fancyhead[RE]{\myfancyheadEven}
	\fancyfoot[LO]{\myfancyfootOdd}
	\fancyfoot[RE]{\myfancyfootEven}
	\renewcommand{\headrulewidth}{0pt}
	\renewcommand{\footrulewidth}{0pt}
}
%=================== HẾT PHẦN GỘP TỪ eduzy-de.tex ===================
\newcommand{\indam}[1]{
	{\fontsize{12pt}{9pt}\sffamily \selectfont\bfseries\color{gold1} #1}
}
\newcommand{\ind}[1]{
	{\fontsize{14pt}{9pt}\fontfamily{qag} \selectfont\bfseries\color{gold1} #1}
}
\renewcommand{\FalseEX}{\stepcounter{dapan}\textbf{\selectfont\color{gold1}\Alph{dapan}.}\ignorespaces}
\renewtheorem{ex}{\color{gold1}Câu}
\renewcommand{\baselinestretch}{1.1}
\newcommand{\hoac}[1]{
	\left[\begin{aligned}#1\end{aligned}\right.}
\newcommand{\heva}[1]{
	\left\{\begin{aligned}#1\end{aligned}\right.}
\usepackage{esvect}
\def\vec{\vv}
\def\vt{\vv}
\def\overrightarrow{\vv}
%====================== LOGO EDUZY VẼ BẰNG TIKZ ======================
%  Đường viền dò từ logo gốc; dải màu vàng kim đo trực tiếp trên ảnh gốc
%  (nghiêng 21 độ, tô bằng 80 dải hẹp nên hiển thị giống nhau trên mọi trình xem PDF).
%  Dùng: \eduzylogo          -> cao 2.4cm (mặc định)
%        \eduzylogo[1.5cm]   -> cao 1.5cm
\definecolor{ezText}{RGB}{245,200,46}
\newcommand{\eduzylogo}[1][2.4cm]{%
	\resizebox{!}{#1}{%
		\begin{tikzpicture}[x=0.01cm,y=0.01cm]
			\path[use as bounding box] (0,0) rectangle (224,-224);
			% mũ tốt nghiệp + vương miện
			\begin{scope}
				\clip
				(119.4,-176.1)..controls(118.9,-176.1)and(116.8,-176)..(114.6,-176)
				..controls(101.4,-176)and(85.8,-175.9)..(83.7,-175.7)
				..controls(81.8,-175.6)and(81,-175.6)..(80.2,-175.7)
				..controls(78.8,-175.9)and(77.7,-175.9)..(76.2,-175.7)
				..controls(74.9,-175.5)and(70.8,-175.5)..(64.3,-175.7)
				..controls(59.1,-175.9)and(58.3,-175.7)..(58.6,-173.7)
				..controls(58.6,-172.8)and(59.6,-169.6)..(60.3,-168)
				..controls(60.8,-166.9)and(61.4,-165.8)..(62.3,-165)
				..controls(63.1,-164.2)and(64.3,-163.4)..(65,-163.2)
				..controls(65.5,-163.1)and(72.3,-163)..(76.6,-163.1)
				..controls(77.9,-163.1)and(80.1,-163.1)..(81.4,-163)
				..controls(82.8,-162.9)and(86.3,-162.9)..(89.2,-162.9)
				..controls(93,-162.9)and(94.9,-162.8)..(96,-162.7)
				..controls(97.2,-162.5)and(101,-162.5)..(114,-162.4)
				..controls(123,-162.3)and(134.9,-162.2)..(140.2,-162.1)
				..controls(150.1,-162)and(152.2,-162.1)..(153.6,-162.4)
				..controls(154.4,-162.6)and(157.1,-163.8)..(158.1,-164.4)
				..controls(158.8,-164.8)and(159.7,-165.7)..(159.9,-166.3)
				..controls(160,-166.5)and(160.2,-167.1)..(160.2,-167.6)
				..controls(160.3,-168)and(160.6,-169.4)..(160.9,-170.6)
				..controls(161.2,-171.7)and(161.5,-173.1)..(161.6,-173.6)
				..controls(161.8,-174.4)and(161.8,-174.5)..(161.5,-174.9)
				..controls(161.2,-175.6)and(160.8,-175.9)..(159.9,-176)
				..controls(159.2,-176.1)and(154.1,-176.2)..(151.5,-176.2)
				..controls(148.1,-176.1)and(122.6,-176.1)..(121.6,-176.2)
				..controls(120.8,-176.2)and(119.8,-176.2)..(119.4,-176.1)--cycle
				(130.4,-160)..controls(129,-160)and(119.8,-160)..(110.1,-159.9)
				..controls(100.3,-159.8)and(88.5,-159.7)..(83.9,-159.6)
				..controls(77.8,-159.5)and(75.3,-159.5)..(74.5,-159.6)
				..controls(73.9,-159.7)and(72.5,-159.8)..(71.3,-159.8)
				..controls(69.6,-159.8)and(69,-159.8)..(68.5,-159.6)
				..controls(67.6,-159.3)and(66.8,-158.6)..(66.2,-157.6)
				..controls(66,-157.3)and(65.6,-156.5)..(65.3,-156.1)
				..controls(64.6,-155.1)and(64.5,-154.9)..(63.7,-152.9)
				..controls(63.2,-151.6)and(62.9,-150.9)..(61.8,-149.6)
				..controls(61.5,-149.2)and(61.1,-148.5)..(60.9,-147.9)
				..controls(60.7,-147.4)and(60.3,-146.6)..(60.1,-146.1)
				..controls(59.9,-145.5)and(59.5,-144.9)..(59,-144.3)
				..controls(58.5,-143.7)and(58.2,-143.2)..(57.9,-142.5)
				..controls(57.7,-141.9)and(57.3,-141.1)..(57,-140.7)
				..controls(56.6,-140.3)and(56.1,-139.3)..(55.8,-138.5)
				..controls(55.5,-137.8)and(55.1,-136.8)..(54.9,-136.5)
				..controls(54.7,-136.1)and(54.3,-135.2)..(53.9,-134.5)
				..controls(53.6,-133.8)and(53,-132.7)..(52.6,-132)
				..controls(51.6,-130.5)and(51.4,-130.2)..(51.1,-129)
				..controls(50.6,-127.5)and(50.6,-127.4)..(49.7,-126.3)
				..controls(49.2,-125.7)and(48.6,-124.8)..(48.3,-124.2)
				..controls(48,-123.6)and(47.5,-122.7)..(47.3,-122.2)
				..controls(47,-121.8)and(46.5,-120.7)..(46.1,-119.8)
				..controls(45.8,-118.9)and(45.2,-117.8)..(44.9,-117.2)
				..controls(44.2,-116.1)and(43.3,-114.4)..(42.3,-112.2)
				..controls(41.8,-111.3)and(41.3,-110.4)..(41.1,-110.1)
				..controls(40.5,-109.4)and(40.3,-109)..(40,-107.8)
				..controls(39.8,-107)and(39.6,-106.5)..(39.2,-106)
				..controls(38.9,-105.6)and(38.4,-104.7)..(38,-104)
				..controls(37.7,-103.4)and(37.2,-102.5)..(37,-102.1)
				..controls(36.7,-101.7)and(36.3,-100.8)..(36,-100.1)
				..controls(35.6,-98.9)and(35.3,-98.4)..(34.6,-97.6)
				..controls(34.2,-97.1)and(34.2,-96.3)..(34.6,-96)
				..controls(35,-95.7)and(35.5,-95.9)..(36.2,-96.7)
				..controls(36.5,-97)and(37.2,-97.5)..(37.6,-97.7)..controls(38,-98)and(38.6,-98.4)..(39.1,-98.9)
				..controls(39.6,-99.4)and(40.3,-100)..(40.6,-100.3)
				..controls(41.3,-100.8)and(42.4,-101.8)..(43.7,-102.9)
				..controls(44,-103.2)and(44.6,-103.7)..(45,-103.9)
				..controls(45.4,-104.2)and(45.9,-104.6)..(46.1,-104.8)
				..controls(47.1,-105.9)and(48.8,-107.1)..(50.5,-107.9)
				..controls(51.8,-108.6)and(52.3,-109)..(53.4,-110)
				..controls(54.8,-111.2)and(55.6,-111.6)..(57,-112.1)
				..controls(57.5,-112.3)and(58.5,-112.8)..(59.4,-113.3)
				..controls(60.2,-113.7)and(61.3,-114.3)..(61.9,-114.5)
				..controls(62.4,-114.7)and(63.3,-115.1)..(63.8,-115.3)
				..controls(64.4,-115.6)and(65.4,-115.9)..(66.2,-116.1)
				..controls(67,-116.3)and(68.4,-116.6)..(69.2,-116.9)
				..controls(72.3,-117.9)and(72.8,-118)..(78,-118.1)
				..controls(84,-118.3)and(86.1,-118.1)..(88.1,-117.1)
				..controls(89.4,-116.5)and(93,-115)..(93.7,-114.9)
				..controls(94.8,-114.6)and(95.4,-114.2)..(95.9,-113.4)
				..controls(96.2,-113)and(96.7,-112.5)..(97.1,-112.2)
				..controls(98,-111.6)and(101.3,-108.4)..(102.1,-107.2)
				..controls(102.4,-106.8)and(102.9,-106.2)..(103.2,-105.8)
				..controls(103.5,-105.5)and(104,-104.8)..(104.2,-104.2)
				..controls(104.5,-103.6)and(105,-102.8)..(105.3,-102.4)
				..controls(105.6,-102)and(106.1,-101.2)..(106.5,-100.6)
				..controls(107.1,-99.5)and(108.6,-97.5)..(109.5,-96.5)
				..controls(109.9,-96.1)and(110.1,-96)..(110.4,-96)
				..controls(111.1,-96)and(111.6,-96.5)..(111.8,-97.6)
				..controls(112,-98.6)and(112.3,-99.3)..(112.8,-99.9)
				..controls(113.7,-100.9)and(115.2,-103.1)..(115.7,-104.2)
				..controls(116.1,-105)and(116.5,-105.5)..(116.9,-106)
				..controls(117.3,-106.4)and(117.8,-106.9)..(118,-107.3)
				..controls(118.2,-107.6)and(118.6,-108)..(119,-108.3)
				..controls(119.6,-108.8)and(121.1,-110.3)..(121.8,-111.2)
				..controls(122.1,-111.7)and(122.5,-112)..(123.4,-112.4)
				..controls(124.8,-113.2)and(125.8,-113.8)..(127,-114.6)
				..controls(127.5,-114.9)and(128.1,-115.3)..(128.2,-115.4)
				..controls(128.4,-115.5)and(129.3,-115.8)..(130.2,-116)
				..controls(131.1,-116.2)and(132.5,-116.6)..(133.3,-116.9)
				..controls(135.7,-117.9)and(136.1,-117.9)..(142.5,-117.8)
				..controls(145.8,-117.8)and(148.4,-117.7)..(148.8,-117.6)
				..controls(149.2,-117.6)and(150.2,-117.2)..(151,-116.8)
				..controls(152,-116.4)and(152.8,-116)..(153.6,-115.9)
				..controls(154.5,-115.7)and(155.2,-115.4)..(156.3,-114.9)
				..controls(157.1,-114.5)and(158.2,-114)..(158.7,-113.8)
				..controls(159.1,-113.7)and(159.8,-113.3)..(160.2,-113)
				..controls(160.6,-112.7)and(161.3,-112.4)..(161.7,-112.2)
				..controls(162.7,-111.9)and(163.2,-111.7)..(164.1,-111.1)
				..controls(164.5,-110.8)and(165.3,-110.4)..(165.9,-110)
				..controls(166.5,-109.7)and(167.3,-109.2)..(167.6,-108.9)
				..controls(167.9,-108.6)and(168.6,-108.2)..(169,-108)
				..controls(169.4,-107.8)and(170.2,-107.3)..(170.7,-106.9)
				..controls(171.2,-106.5)and(172.1,-105.8)..(172.7,-105.4)
				..controls(173.2,-105)and(174.1,-104.4)..(174.5,-104)
				..controls(174.9,-103.7)and(175.7,-103.1)..(176.2,-102.6)
				..controls(176.7,-102.2)and(177.4,-101.5)..(177.9,-101.2)
				..controls(178.3,-100.8)and(178.8,-100.4)..(179,-100.2)
				..controls(179.2,-100)and(179.6,-99.6)..(180,-99.3)
				..controls(180.4,-99.1)and(180.9,-98.6)..(181.2,-98.4)
				..controls(181.5,-98.1)and(182.1,-97.7)..(182.5,-97.5)
				..controls(183.2,-97.2)and(183.5,-96.9)..(184.3,-95.9)
				..controls(184.6,-95.6)and(184.6,-95.6)..(184.8,-95.7)
				..controls(185,-95.8)and(185.1,-96)..(185.1,-96.2)
				..controls(185.1,-96.5)and(185,-96.7)..(184.6,-97.2)
				..controls(184.1,-97.8)and(184,-97.9)..(183.9,-98.8)
				..controls(183.7,-100.3)and(183.3,-101.3)..(182.9,-101.8)
				..controls(182.6,-102.1)and(182.3,-102.6)..(182.2,-102.9)
				..controls(182,-103.3)and(181.6,-104)..(181.2,-104.6)
				..controls(180.9,-105.2)and(180.5,-106.1)..(180.4,-106.6)
				..controls(179.9,-108)and(179.7,-108.5)..(179,-109.4)
				..controls(178.6,-109.8)and(178.2,-110.5)..(177.9,-111)
				..controls(177.7,-111.5)and(177.3,-112.3)..(177,-112.7)
				..controls(176.6,-113.2)and(176.4,-113.7)..(176.3,-114.2)
				..controls(176.2,-114.7)and(175.9,-115.2)..(175.4,-116)
				..controls(175,-116.6)and(174.4,-117.7)..(174,-118.5)
				..controls(173.7,-119.2)and(173.3,-120.1)..(173,-120.5)
				..controls(172.8,-120.9)and(172.4,-121.8)..(172.1,-122.5)
				..controls(171.8,-123.2)and(171.3,-124.1)..(171,-124.6)
				..controls(170.8,-125.1)and(170.3,-126)..(169.9,-126.6)
				..controls(169.6,-127.3)and(169.2,-128.1)..(169,-128.3)
				..controls(168.8,-128.6)and(168.5,-129.3)..(168.3,-129.8)
				..controls(168.1,-130.3)and(167.8,-131.1)..(167.5,-131.6)
				..controls(166.9,-132.5)and(164.2,-137.9)..(164.2,-138.1)
				..controls(164.2,-138.1)and(163.9,-138.8)..(163.4,-139.7)
				..controls(162.9,-140.5)and(162.4,-141.5)..(162.3,-141.9)
				..controls(161.9,-143)and(161.6,-143.4)..(160.9,-144.1)
				..controls(160.4,-144.6)and(160.3,-144.8)..(160.1,-145.4)
				..controls(160,-145.8)and(159.6,-146.8)..(159.2,-147.8)
				..controls(158.7,-148.8)and(158.1,-150)..(157.9,-150.6)
				..controls(157.6,-151.1)and(157.2,-151.9)..(156.9,-152.4)
				..controls(156.5,-152.8)and(156.1,-153.7)..(155.8,-154.4)
				..controls(155.6,-155)and(155.2,-155.9)..(154.9,-156.3)
				..controls(154.7,-156.7)and(154.4,-157.4)..(154.2,-157.9)
				..controls(154.1,-158.4)and(153.8,-159)..(153.5,-159.3)
				..controls(152.8,-160.1)and(153,-160.1)..(142.4,-160.1)
				..controls(137.2,-160.1)and(131.8,-160.1)..(130.4,-160)--cycle
				(141.5,-103.7)..controls(140.8,-103)and(139.8,-102.1)..(139.2,-101.6)
				..controls(138.5,-101.2)and(137.9,-100.6)..(137.8,-100.4)
				..controls(137,-99.5)and(134,-96.5)..(133.2,-96)
				..controls(132.7,-95.7)and(131.8,-94.9)..(131.1,-94.3)
				..controls(129.4,-92.7)and(127.9,-91.6)..(126.6,-91.1)
				..controls(125.2,-90.5)and(124.1,-90)..(122.9,-89.1)
				..controls(122.1,-88.6)and(121.3,-88.2)..(119.9,-87.7)
				..controls(118.8,-87.3)and(117.5,-86.8)..(116.9,-86.5)
				..controls(114.3,-85.5)and(113.7,-85.4)..(109.7,-85.5)
				..controls(105.3,-85.7)and(103.9,-85.9)..(101.6,-87)
				..controls(100.7,-87.4)and(99.4,-87.8)..(97.9,-88)
				..controls(97,-88.2)and(96.6,-88.3)..(95.8,-88.9)
				..controls(95.4,-89.2)and(94.7,-89.6)..(94.2,-89.8)
				..controls(92.9,-90.3)and(92.3,-90.7)..(90.9,-92.1)
				..controls(90.2,-92.7)and(89.4,-93.3)..(89,-93.5)..controls(87,-94.5)and(85.2,-96)..(83.6,-97.9)
				..controls(83.3,-98.2)and(82.5,-98.9)..(81.8,-99.4)
				..controls(80.7,-100.1)and(80.5,-100.3)..(80,-101.1)
				..controls(78.7,-103.1)and(77.2,-103.9)..(76.4,-103)
				..controls(75.7,-102.2)and(75.7,-103.2)..(75.7,-83.6)
				..controls(75.6,-73.6)and(75.6,-64.8)..(75.5,-64)
				..controls(75.5,-63.2)and(75.5,-61.4)..(75.6,-59.8)--(75.7,-57.1)--(76,-56.8)
				..controls(76.8,-56.2)and(78,-56.3)..(79.8,-57.2)
				..controls(80.4,-57.5)and(81.3,-57.8)..(81.8,-57.9)
				..controls(82.3,-58.1)and(83,-58.4)..(83.6,-58.8)..controls(84.3,-59.2)and(85,-59.5)..(86,-59.8)
				..controls(86.8,-60.1)and(88,-60.6)..(88.7,-60.9)
				..controls(89.3,-61.2)and(90.6,-61.7)..(91.4,-62)
				..controls(92.2,-62.3)and(93.7,-62.9)..(94.6,-63.3)
				..controls(95.6,-63.7)and(96.8,-64.1)..(97.4,-64.3)
				..controls(98,-64.5)and(99,-64.9)..(99.7,-65.2)
				..controls(100.4,-65.6)and(101.5,-66)..(102.3,-66.2)
				..controls(104.2,-66.8)and(105,-67.1)..(106.2,-67.8)
				..controls(107.9,-68.7)and(110.2,-69.4)..(111.3,-69.2)
				..controls(112.5,-69)and(115,-68.2)..(115.6,-67.8)
				..controls(117.3,-66.7)and(117.6,-66.6)..(119.5,-66.2)
				..controls(121.9,-65.8)and(122.1,-65.7)..(123,-65.1)
				..controls(124.1,-64.4)and(124.4,-64.3)..(126,-63.9)
				..controls(127,-63.7)and(127.7,-63.4)..(128.4,-63.1)
				..controls(129,-62.8)and(129.7,-62.5)..(130.6,-62.4)
				..controls(132.1,-62.1)and(134,-61.5)..(135.2,-61)
				..controls(135.7,-60.8)and(136.7,-60.5)..(137.3,-60.3)
				..controls(138.4,-60)and(139.3,-59.7)..(141.1,-59)
				..controls(141.6,-58.7)and(142,-58.6)..(142.5,-58.6)
				..controls(143.2,-58.6)and(143.3,-58.7)..(143.6,-59)
				..controls(144,-59.5)and(144.1,-59.9)..(144,-60.8)
				..controls(144,-61.2)and(144,-71.1)..(144,-82.8)--(143.9,-104.1)--(143.7,-104.5)
				..controls(143.5,-104.7)and(143.3,-104.8)..(143,-104.9)
				..controls(142.6,-104.9)and(142.5,-104.8)..(141.5,-103.7)--cycle
				(27.8,-91.4)..controls(27.1,-90.9)and(26.3,-90)..(26.1,-89.5)
				..controls(26,-89.2)and(25.8,-88.6)..(25.7,-87.9)
				..controls(25.4,-86.2)and(25.6,-84.8)..(26.5,-83.5)
				..controls(27.4,-82.3)and(28.3,-82)..(29.4,-82.3)
				..controls(31.1,-82.8)and(32.1,-84.9)..(32,-87.4)
				..controls(31.9,-89.3)and(31.2,-90.6)..(29.9,-91.4)
				..controls(29.1,-91.9)and(28.6,-91.9)..(27.8,-91.4)--cycle
				(190.1,-91.3)..controls(189.7,-91.2)and(189.3,-90.9)..(189.1,-90.7)
				..controls(188.3,-89.8)and(187.6,-87.7)..(187.6,-86.2)
				..controls(187.6,-85.5)and(187.7,-85.2)..(188.2,-84.4)
				..controls(189.1,-82.6)and(190.1,-81.8)..(191.3,-82.1)
				..controls(192,-82.3)and(193.5,-83.7)..(193.8,-84.4)
				..controls(194.1,-85.1)and(194.2,-87.1)..(194,-88.3)
				..controls(193.8,-89.8)and(193.3,-90.5)..(192.3,-91.3)
				..controls(191.7,-91.7)and(190.9,-91.7)..(190.1,-91.3)--cycle
				(173.7,-83.4)..controls(172.2,-82.5)and(172,-82.5)..(170.6,-82.4)
				..controls(169.7,-82.4)and(168.8,-82.4)..(168.3,-82.5)
				..controls(167.2,-82.7)and(166.9,-82.7)..(166.6,-82.2)
				..controls(166.3,-81.8)and(166.3,-81)..(166.8,-80)
				..controls(167.2,-78.9)and(167.4,-77.4)..(167.7,-74.1)
				..controls(167.8,-72.8)and(167.9,-71.8)..(167.9,-71.7)
				..controls(168,-71.7)and(168,-70.2)..(168,-68.4)..controls(168,-65.4)and(168,-64.9)..(167.7,-63)
				..controls(167.3,-60.5)and(167.1,-59.9)..(166.4,-58.4)
				..controls(166.1,-57.7)and(165.8,-56.9)..(165.7,-56.6)
				..controls(165.7,-56.2)and(165.4,-55.6)..(165.2,-55.2)
				..controls(164.6,-53.8)and(161.9,-50.7)..(161,-50.3)
				..controls(159.8,-49.7)and(157.9,-50)..(155.2,-51.2)
				..controls(154.7,-51.4)and(153.4,-51.8)..(152.2,-52.1)
				..controls(150.8,-52.4)and(149.4,-52.9)..(148.6,-53.2)
				..controls(147.8,-53.6)and(146.7,-54)..(146,-54.2)
				..controls(145.4,-54.4)and(144.5,-54.7)..(144,-54.9)
				..controls(143.6,-55.2)and(142.7,-55.5)..(141.8,-55.7)
				..controls(141,-55.9)and(139.6,-56.3)..(138.7,-56.6)
				..controls(137.8,-57)and(136.3,-57.4)..(135.3,-57.6)
				..controls(133.8,-57.9)and(133.2,-58.1)..(132.1,-58.7)
				..controls(131.1,-59.2)and(130.3,-59.4)..(129.4,-59.6)
				..controls(128.7,-59.8)and(127.8,-60)..(127.4,-60.1)
				..controls(127,-60.3)and(126,-60.5)..(125.1,-60.7)
				..controls(123.2,-61.1)and(121.7,-61.5)..(120.2,-62.1)
				..controls(119.7,-62.3)and(118.1,-62.8)..(116.7,-63.2)
				..controls(114.6,-63.9)and(114,-64.1)..(113,-64.7)
				..controls(111.9,-65.4)and(111.8,-65.4)..(111,-65.5)
				..controls(109.9,-65.5)and(109.3,-65.4)..(107.7,-64.7)
				..controls(107,-64.4)and(106,-64)..(105.3,-63.9)
				..controls(103.5,-63.5)and(102.2,-63)..(101.2,-62.5)
				..controls(100.1,-61.8)and(99.3,-61.5)..(97.8,-61.1)
				..controls(97.2,-61)and(96.2,-60.7)..(95.7,-60.4)
				..controls(95.1,-60.1)and(94.2,-59.8)..(93.5,-59.6)
				..controls(92.9,-59.4)and(91.7,-59)..(91,-58.7)..controls(90.2,-58.4)and(88.9,-58)..(88,-57.7)
				..controls(87.2,-57.5)and(86.1,-57.1)..(85.5,-56.8)
				..controls(85,-56.6)and(84.1,-56.2)..(83.5,-56)..controls(82.9,-55.7)and(81.8,-55.2)..(81,-54.9)
				..controls(80.1,-54.5)and(79.1,-54.2)..(78.7,-54.1)
				..controls(77.9,-53.8)and(77.3,-53.6)..(76.2,-52.9)
				..controls(75.7,-52.6)and(75.2,-52.4)..(74.1,-52.2)
				..controls(72.5,-51.9)and(72.4,-51.8)..(71.3,-51.1)
				..controls(69.9,-50.1)and(69.4,-49.9)..(67.9,-49.8)
				..controls(67,-49.7)and(66.6,-49.6)..(66.3,-49.4)
				..controls(66.1,-49.3)and(65.7,-49)..(65.4,-48.8)
				..controls(65.1,-48.7)and(64.2,-48.4)..(63.6,-48.2)
				..controls(62.9,-48)and(61.7,-47.6)..(60.9,-47.2)
				..controls(60.1,-46.8)and(58.9,-46.4)..(58.2,-46.2)
				..controls(57.6,-46)and(56.8,-45.7)..(56.5,-45.5)
				..controls(54.9,-44.6)and(54.6,-44.5)..(53.2,-44.2)
				..controls(51.7,-43.8)and(51,-43.6)..(49.9,-42.9)
				..controls(49.5,-42.6)and(48.9,-42.4)..(47.5,-42)
				..controls(45.7,-41.6)and(45.6,-41.5)..(44.3,-40.6)
				..controls(43.5,-40.1)and(42.8,-39.6)..(42.7,-39.4)
				..controls(42.4,-38.7)and(43,-38.1)..(44.5,-37.6)..controls(44.9,-37.5)and(46,-37)..(46.9,-36.6)
				..controls(47.9,-36.1)and(48.8,-35.8)..(49.6,-35.6)
				..controls(50.5,-35.4)and(51.3,-35.2)..(52.4,-34.6)
				..controls(53.3,-34.1)and(54.5,-33.7)..(54.9,-33.5)
				..controls(55.4,-33.4)and(56.4,-33)..(57.2,-32.6)
				..controls(58.6,-32)and(60.1,-31.5)..(61.6,-31.2)
				..controls(62.3,-31.1)and(62.7,-31)..(63.2,-30.7)
				..controls(64.1,-30)and(64.3,-29.9)..(65.3,-29.7)
				..controls(65.7,-29.6)and(66.5,-29.2)..(67.1,-28.9)
				..controls(67.6,-28.6)and(68.8,-28.2)..(69.6,-27.9)
				..controls(70.5,-27.7)and(71.6,-27.3)..(72.1,-27)
				..controls(72.6,-26.7)and(73.8,-26.2)..(74.8,-25.9)
				..controls(75.7,-25.5)and(76.7,-25.1)..(77.1,-25)
				..controls(77.4,-24.8)and(78.5,-24.4)..(79.6,-24.1)
				..controls(80.6,-23.8)and(82,-23.2)..(82.8,-22.9)
				..controls(83.5,-22.5)and(84.8,-22)..(85.8,-21.8)
				..controls(86.7,-21.5)and(87.8,-21.1)..(88.1,-20.9)
				..controls(88.4,-20.7)and(89.3,-20.3)..(90.1,-20.1)
				..controls(90.8,-19.8)and(91.9,-19.4)..(92.4,-19.1)
				..controls(93.5,-18.5)and(94.4,-18.2)..(95.4,-18.1)
				..controls(95.7,-18)and(96.4,-17.8)..(96.9,-17.5)
				..controls(97.5,-17.3)and(98.4,-16.9)..(99,-16.6)
				..controls(99.7,-16.3)and(100.9,-15.9)..(101.9,-15.6)
				..controls(102.9,-15.3)and(104.4,-14.7)..(105.3,-14.3)
				..controls(108.1,-13.2)and(108.5,-13.1)..(110.1,-13.1)
				..controls(111.4,-13.1)and(111.5,-13.1)..(112.4,-13.5)
				..controls(113.2,-14)and(115,-14.8)..(116.6,-15.4)
				..controls(116.9,-15.5)and(117.9,-15.8)..(118.8,-16.1)
				..controls(120.9,-16.6)and(121.5,-16.9)..(122.6,-17.9)
				..controls(123,-18.3)and(123.7,-18.8)..(123.9,-18.9)
				..controls(124.6,-19.4)and(126,-19.8)..(127.9,-20.2)
				..controls(129.6,-20.5)and(129.8,-20.6)..(131.1,-21.4)
				..controls(131.5,-21.7)and(132.4,-22.1)..(133.1,-22.3)
				..controls(136.1,-23.2)and(138.3,-24.1)..(140.2,-25)
				..controls(141,-25.5)and(141.5,-25.6)..(142.5,-25.8)
				..controls(145.1,-26.3)and(146,-26.6)..(149.4,-28.1)
				..controls(151.3,-28.9)and(153.5,-29.8)..(154.3,-30.1)
				..controls(155.1,-30.4)and(156.2,-30.8)..(156.6,-31)
				..controls(157.1,-31.2)and(158.2,-31.6)..(159.1,-31.8)
				..controls(160.6,-32.2)and(160.9,-32.3)..(161.7,-32.9)
				..controls(162.4,-33.3)and(162.9,-33.6)..(163.7,-33.8)
				..controls(164.3,-34)and(164.8,-34.1)..(164.9,-34.1)
				..controls(165,-34.1)and(165.6,-34.4)..(166.4,-34.8)
				..controls(167.3,-35.2)and(168.1,-35.4)..(169,-35.6)
				..controls(170,-35.8)and(170.5,-36)..(171.7,-36.6)
				..controls(172.4,-37)and(173.3,-37.4)..(173.6,-37.4)
				..controls(175.6,-37.7)and(176.2,-38)..(176.7,-39.1)
				..controls(176.9,-39.6)and(176.9,-39.7)..(176.7,-40.4)
				..controls(176.4,-41.6)and(175.8,-42.7)..(175.2,-43.4)
				..controls(174.5,-44.1)and(172.6,-45.5)..(171.7,-45.8)
				..controls(170.6,-46.2)and(169.1,-46)..(168.2,-45.4)
				..controls(167.9,-45.2)and(167.3,-44.8)..(166.9,-44.4)
				..controls(166.4,-44.1)and(165.8,-43.6)..(165.4,-43.2)
				..controls(165.1,-42.9)and(164.4,-42.3)..(163.9,-41.9)
				..controls(163.3,-41.5)and(162.7,-40.9)..(162.4,-40.5)
				..controls(161.4,-39.2)and(159.9,-38)..(158.8,-37.8)
				..controls(158.4,-37.8)and(158.3,-37.7)..(158.2,-37.2)
				..controls(158,-36.8)and(157.4,-36.1)..(157.1,-36.1)
				..controls(157,-36.1)and(156.7,-36.1)..(156.5,-36)
				..controls(156.1,-35.9)and(155.9,-35.7)..(155.7,-35.2)
				..controls(155.3,-34.6)and(155.1,-34.4)..(154,-34.1)
				..controls(153.6,-34)and(153.2,-33.7)..(152.9,-33.5)
				..controls(152.4,-32.9)and(151.6,-32.4)..(151.3,-32.4)
				..controls(150.9,-32.4)and(150.4,-32.6)..(150.4,-32.9)
				..controls(150.4,-33.2)and(150.8,-33.4)..(151.4,-33.6)
				..controls(152.1,-33.7)and(152.8,-34.1)..(153.5,-34.9)
				..controls(154.1,-35.5)and(154.3,-35.6)..(155,-35.8)
				..controls(155.7,-36)and(155.7,-36)..(156,-36.6)
				..controls(156.3,-37.3)and(156.7,-37.7)..(157.5,-37.8)
				..controls(158.1,-37.9)and(158.2,-38)..(158.4,-38.8)
				..controls(158.5,-39.4)and(158.6,-39.6)..(159.6,-40.4)
				..controls(160.2,-41)and(160.9,-41.6)..(161.3,-41.8)
				..controls(161.7,-42)and(161.9,-42.3)..(162.2,-42.8)
				..controls(162.4,-43.2)and(162.8,-43.8)..(163.1,-44.1)
				..controls(163.9,-44.8)and(165.1,-46.7)..(165.7,-48)
				..controls(166.2,-49.3)and(166.4,-49.7)..(167.2,-50.5)
				..controls(167.5,-50.8)and(167.9,-51.4)..(168.1,-51.8)
				..controls(168.3,-52.2)and(168.8,-53)..(169.1,-53.6)
				..controls(169.6,-54.3)and(169.8,-54.9)..(169.9,-55.6)
				..controls(170.4,-57.8)and(170.5,-58.1)..(171.2,-59.5)
				..controls(171.6,-60.3)and(172,-61.5)..(172.2,-62.1)
				..controls(172.3,-62.8)and(172.7,-63.7)..(172.9,-64.1)
				..controls(173.2,-64.6)and(173.4,-65.3)..(173.4,-65.6)
				..controls(173.8,-68.8)and(174.2,-70.6)..(174.9,-72.6)
				..controls(175.5,-74.5)and(175.9,-77.1)..(176.1,-80.2)
				..controls(176.3,-83.4)and(175.5,-84.4)..(173.7,-83.4)--cycle;
				\definecolor{ezs}{RGB}{247,179,26}\fill[ezs](-5.7,-156.6)--(-3.4,-157.5)--(61.9,12.7)--(59.7,13.6)--cycle;
				\definecolor{ezs}{RGB}{247,179,22}\fill[ezs](-3.8,-157.3)--(-1.5,-158.2)--(63.9,12)--(61.6,12.9)--cycle;
				\definecolor{ezs}{RGB}{248,179,19}\fill[ezs](-1.8,-158.1)--(0.4,-158.9)--(65.8,11.2)--(63.5,12.1)--cycle;
				\definecolor{ezs}{RGB}{248,178,16}\fill[ezs](0.1,-158.8)--(2.4,-159.7)--(67.7,10.5)--(65.4,11.4)--cycle;
				\definecolor{ezs}{RGB}{249,178,15}\fill[ezs](2,-159.5)--(4.3,-160.4)--(69.6,9.8)--(67.3,10.7)--cycle;
				\definecolor{ezs}{RGB}{251,177,17}\fill[ezs](3.9,-160.3)--(6.2,-161.1)--(71.5,9)--(69.2,9.9)--cycle;
				\definecolor{ezs}{RGB}{252,179,18}\fill[ezs](5.8,-161)--(8.1,-161.9)--(73.4,8.3)--(71.1,9.2)--cycle;
				\definecolor{ezs}{RGB}{252,182,17}\fill[ezs](7.7,-161.7)--(10,-162.6)--(75.3,7.6)--(73,8.5)--cycle;
				\definecolor{ezs}{RGB}{252,185,17}\fill[ezs](9.6,-162.5)--(11.9,-163.4)--(77.2,6.8)--(75,7.7)--cycle;
				\definecolor{ezs}{RGB}{251,189,18}\fill[ezs](11.5,-163.2)--(13.8,-164.1)--(79.2,6.1)--(76.9,7)--cycle;
				\definecolor{ezs}{RGB}{251,191,22}\fill[ezs](13.4,-163.9)--(15.7,-164.8)--(81.1,5.4)--(78.8,6.2)--cycle;
				\definecolor{ezs}{RGB}{252,193,25}\fill[ezs](15.4,-164.7)--(17.6,-165.6)--(83,4.6)--(80.7,5.5)--cycle;
				\definecolor{ezs}{RGB}{253,195,27}\fill[ezs](17.3,-165.4)--(19.6,-166.3)--(84.9,3.9)--(82.6,4.8)--cycle;
				\definecolor{ezs}{RGB}{253,196,31}\fill[ezs](19.2,-166.1)--(21.5,-167)--(86.8,3.2)--(84.5,4)--cycle;
				\definecolor{ezs}{RGB}{254,199,35}\fill[ezs](21.1,-166.9)--(23.4,-167.8)--(88.7,2.4)--(86.4,3.3)--cycle;
				\definecolor{ezs}{RGB}{253,201,38}\fill[ezs](23,-167.6)--(25.3,-168.5)--(90.6,1.7)--(88.3,2.6)--cycle;
				\definecolor{ezs}{RGB}{253,203,40}\fill[ezs](24.9,-168.3)--(27.2,-169.2)--(92.5,1)--(90.3,1.8)--cycle;
				\definecolor{ezs}{RGB}{253,205,42}\fill[ezs](26.8,-169.1)--(29.1,-170)--(94.4,0.2)--(92.2,1.1)--cycle;
				\definecolor{ezs}{RGB}{253,207,45}\fill[ezs](28.7,-169.8)--(31,-170.7)--(96.4,-0.5)--(94.1,0.4)--cycle;
				\definecolor{ezs}{RGB}{253,210,49}\fill[ezs](30.7,-170.5)--(32.9,-171.4)--(98.3,-1.2)--(96,-0.4)--cycle;
				\definecolor{ezs}{RGB}{253,212,54}\fill[ezs](32.6,-171.3)--(34.9,-172.2)--(100.2,-2)--(97.9,-1.1)--cycle;
				\definecolor{ezs}{RGB}{253,214,59}\fill[ezs](34.5,-172)--(36.8,-172.9)--(102.1,-2.7)--(99.8,-1.8)--cycle;
				\definecolor{ezs}{RGB}{253,216,63}\fill[ezs](36.4,-172.7)--(38.7,-173.6)--(104,-3.4)--(101.7,-2.6)--cycle;
				\definecolor{ezs}{RGB}{253,218,65}\fill[ezs](38.3,-173.5)--(40.6,-174.4)--(105.9,-4.2)--(103.6,-3.3)--cycle;
				\definecolor{ezs}{RGB}{254,220,66}\fill[ezs](40.2,-174.2)--(42.5,-175.1)--(107.8,-4.9)--(105.5,-4)--cycle;
				\definecolor{ezs}{RGB}{254,222,68}\fill[ezs](42.1,-175)--(44.4,-175.8)--(109.7,-5.6)--(107.5,-4.8)--cycle;
				\definecolor{ezs}{RGB}{254,224,71}\fill[ezs](44,-175.7)--(46.3,-176.6)--(111.7,-6.4)--(109.4,-5.5)--cycle;
				\definecolor{ezs}{RGB}{254,226,75}\fill[ezs](46,-176.4)--(48.2,-177.3)--(113.6,-7.1)--(111.3,-6.2)--cycle;
				\definecolor{ezs}{RGB}{254,228,79}\fill[ezs](47.9,-177.2)--(50.1,-178)--(115.5,-7.8)--(113.2,-7)--cycle;
				\definecolor{ezs}{RGB}{254,230,82}\fill[ezs](49.8,-177.9)--(52.1,-178.8)--(117.4,-8.6)--(115.1,-7.7)--cycle;
				\definecolor{ezs}{RGB}{254,232,85}\fill[ezs](51.7,-178.6)--(54,-179.5)--(119.3,-9.3)--(117,-8.4)--cycle;
				\definecolor{ezs}{RGB}{254,234,88}\fill[ezs](53.6,-179.4)--(55.9,-180.2)--(121.2,-10)--(118.9,-9.2)--cycle;
				\definecolor{ezs}{RGB}{254,237,91}\fill[ezs](55.5,-180.1)--(57.8,-181)--(123.1,-10.8)--(120.8,-9.9)--cycle;
				\definecolor{ezs}{RGB}{254,239,94}\fill[ezs](57.4,-180.8)--(59.7,-181.7)--(125,-11.5)--(122.8,-10.6)--cycle;
				\definecolor{ezs}{RGB}{254,241,97}\fill[ezs](59.3,-181.6)--(61.6,-182.4)--(127,-12.2)--(124.7,-11.4)--cycle;
				\definecolor{ezs}{RGB}{254,243,101}\fill[ezs](61.2,-182.3)--(63.5,-183.2)--(128.9,-13)--(126.6,-12.1)--cycle;
				\definecolor{ezs}{RGB}{254,245,103}\fill[ezs](63.2,-183)--(65.4,-183.9)--(130.8,-13.7)--(128.5,-12.8)--cycle;
				\definecolor{ezs}{RGB}{254,245,105}\fill[ezs](65.1,-183.8)--(67.4,-184.6)--(132.7,-14.4)--(130.4,-13.6)--cycle;
				\definecolor{ezs}{RGB}{254,245,105}\fill[ezs](67,-184.5)--(69.3,-185.4)--(134.6,-15.2)--(132.3,-14.3)--cycle;
				\definecolor{ezs}{RGB}{254,244,103}\fill[ezs](68.9,-185.2)--(71.2,-186.1)--(136.5,-15.9)--(134.2,-15)--cycle;
				\definecolor{ezs}{RGB}{253,243,101}\fill[ezs](70.8,-186)--(73.1,-186.8)--(138.4,-16.6)--(136.1,-15.8)--cycle;
				\definecolor{ezs}{RGB}{253,241,98}\fill[ezs](72.7,-186.7)--(75,-187.6)--(140.3,-17.4)--(138,-16.5)--cycle;
				\definecolor{ezs}{RGB}{253,239,95}\fill[ezs](74.6,-187.4)--(76.9,-188.3)--(142.2,-18.1)--(140,-17.2)--cycle;
				\definecolor{ezs}{RGB}{253,236,92}\fill[ezs](76.5,-188.2)--(78.8,-189)--(144.2,-18.8)--(141.9,-18)--cycle;
				\definecolor{ezs}{RGB}{253,234,88}\fill[ezs](78.5,-188.9)--(80.7,-189.8)--(146.1,-19.6)--(143.8,-18.7)--cycle;
				\definecolor{ezs}{RGB}{253,232,85}\fill[ezs](80.4,-189.6)--(82.7,-190.5)--(148,-20.3)--(145.7,-19.4)--cycle;
				\definecolor{ezs}{RGB}{253,230,82}\fill[ezs](82.3,-190.4)--(84.6,-191.2)--(149.9,-21)--(147.6,-20.2)--cycle;
				\definecolor{ezs}{RGB}{253,228,78}\fill[ezs](84.2,-191.1)--(86.5,-192)--(151.8,-21.8)--(149.5,-20.9)--cycle;
				\definecolor{ezs}{RGB}{253,226,74}\fill[ezs](86.1,-191.8)--(88.4,-192.7)--(153.7,-22.5)--(151.4,-21.6)--cycle;
				\definecolor{ezs}{RGB}{253,224,70}\fill[ezs](88,-192.6)--(90.3,-193.4)--(155.6,-23.3)--(153.3,-22.4)--cycle;
				\definecolor{ezs}{RGB}{253,222,67}\fill[ezs](89.9,-193.3)--(92.2,-194.2)--(157.5,-24)--(155.3,-23.1)--cycle;
				\definecolor{ezs}{RGB}{253,220,65}\fill[ezs](91.8,-194)--(94.1,-194.9)--(159.5,-24.7)--(157.2,-23.8)--cycle;
				\definecolor{ezs}{RGB}{253,217,63}\fill[ezs](93.7,-194.8)--(96,-195.6)--(161.4,-25.5)--(159.1,-24.6)--cycle;
				\definecolor{ezs}{RGB}{253,215,60}\fill[ezs](95.7,-195.5)--(97.9,-196.4)--(163.3,-26.2)--(161,-25.3)--cycle;
				\definecolor{ezs}{RGB}{253,213,56}\fill[ezs](97.6,-196.2)--(99.9,-197.1)--(165.2,-26.9)--(162.9,-26)--cycle;
				\definecolor{ezs}{RGB}{252,211,53}\fill[ezs](99.5,-197)--(101.8,-197.8)--(167.1,-27.7)--(164.8,-26.8)--cycle;
				\definecolor{ezs}{RGB}{252,209,49}\fill[ezs](101.4,-197.7)--(103.7,-198.6)--(169,-28.4)--(166.7,-27.5)--cycle;
				\definecolor{ezs}{RGB}{252,207,45}\fill[ezs](103.3,-198.4)--(105.6,-199.3)--(170.9,-29.1)--(168.6,-28.2)--cycle;
				\definecolor{ezs}{RGB}{253,205,42}\fill[ezs](105.2,-199.2)--(107.5,-200)--(172.8,-29.9)--(170.6,-29)--cycle;
				\definecolor{ezs}{RGB}{253,202,40}\fill[ezs](107.1,-199.9)--(109.4,-200.8)--(174.7,-30.6)--(172.5,-29.7)--cycle;
				\definecolor{ezs}{RGB}{253,200,38}\fill[ezs](109,-200.6)--(111.3,-201.5)--(176.7,-31.3)--(174.4,-30.4)--cycle;
				\definecolor{ezs}{RGB}{253,198,35}\fill[ezs](111,-201.4)--(113.2,-202.2)--(178.6,-32.1)--(176.3,-31.2)--cycle;
				\definecolor{ezs}{RGB}{253,196,31}\fill[ezs](112.9,-202.1)--(115.2,-203)--(180.5,-32.8)--(178.2,-31.9)--cycle;
				\definecolor{ezs}{RGB}{253,194,28}\fill[ezs](114.8,-202.8)--(117.1,-203.7)--(182.4,-33.5)--(180.1,-32.6)--cycle;
				\definecolor{ezs}{RGB}{253,192,26}\fill[ezs](116.7,-203.6)--(119,-204.5)--(184.3,-34.3)--(182,-33.4)--cycle;
				\definecolor{ezs}{RGB}{253,189,24}\fill[ezs](118.6,-204.3)--(120.9,-205.2)--(186.2,-35)--(183.9,-34.1)--cycle;
				\definecolor{ezs}{RGB}{253,186,20}\fill[ezs](120.5,-205)--(122.8,-205.9)--(188.1,-35.7)--(185.8,-34.8)--cycle;
				\definecolor{ezs}{RGB}{253,184,17}\fill[ezs](122.4,-205.8)--(124.7,-206.7)--(190,-36.5)--(187.8,-35.6)--cycle;
				\definecolor{ezs}{RGB}{252,183,15}\fill[ezs](124.3,-206.5)--(126.6,-207.4)--(192,-37.2)--(189.7,-36.3)--cycle;
				\definecolor{ezs}{RGB}{252,181,14}\fill[ezs](126.3,-207.2)--(128.5,-208.1)--(193.9,-37.9)--(191.6,-37.1)--cycle;
				\definecolor{ezs}{RGB}{252,179,13}\fill[ezs](128.2,-208)--(130.4,-208.9)--(195.8,-38.7)--(193.5,-37.8)--cycle;
				\definecolor{ezs}{RGB}{252,178,12}\fill[ezs](130.1,-208.7)--(132.4,-209.6)--(197.7,-39.4)--(195.4,-38.5)--cycle;
				\definecolor{ezs}{RGB}{251,178,10}\fill[ezs](132,-209.4)--(134.3,-210.3)--(199.6,-40.1)--(197.3,-39.3)--cycle;
				\definecolor{ezs}{RGB}{251,178,10}\fill[ezs](133.9,-210.2)--(136.2,-211.1)--(201.5,-40.9)--(199.2,-40)--cycle;
				\definecolor{ezs}{RGB}{250,177,11}\fill[ezs](135.8,-210.9)--(138.1,-211.8)--(203.4,-41.6)--(201.1,-40.7)--cycle;
				\definecolor{ezs}{RGB}{250,178,11}\fill[ezs](137.7,-211.6)--(140,-212.5)--(205.3,-42.3)--(203.1,-41.5)--cycle;
				\definecolor{ezs}{RGB}{248,180,10}\fill[ezs](139.6,-212.4)--(141.9,-213.3)--(207.3,-43.1)--(205,-42.2)--cycle;
				\definecolor{ezs}{RGB}{247,181,9}\fill[ezs](141.5,-213.1)--(143.8,-214)--(209.2,-43.8)--(206.9,-42.9)--cycle;
				\definecolor{ezs}{RGB}{246,182,9}\fill[ezs](143.5,-213.8)--(145.7,-214.7)--(211.1,-44.5)--(208.8,-43.7)--cycle;
				\definecolor{ezs}{RGB}{246,182,10}\fill[ezs](145.4,-214.6)--(147.7,-215.5)--(213,-45.3)--(210.7,-44.4)--cycle;
			\end{scope}
			% chữ EDUZY (Times đậm, làm dày nét bằng 5 lớp lệch nhẹ; căn đúng khung chữ logo gốc)
			\foreach \dx/\dy in {-0.7/0,0.7/0,0/-0.3,0/0.3,0/0}
			\node[anchor=center,inner sep=0pt,outer sep=0pt,text=ezText,align=center] at (109.6+\dx,-193.6+\dy)
			{\resizebox{1.01cm}{0.22cm}{\fontencoding{T1}\fontfamily{ptm}\fontseries{b}\selectfont EDUZY}};
		\end{tikzpicture}%
	}%
}
%====================== THIẾT LẬP THÂN ĐỀ ======================
%--- Môi trường tự luận {bt}: ex_test 3.3.4 chưa khai báo sẵn nên tạo ở đây
%    (ex_test đã tự móc \fixshowans và \listenumerate cho bt)
\makeatletter
\@ifundefined{bt}{\newtheorem{bt}{\color{gold1}\namebt}}{}
\makeatother
\renewcommand{\sectionmark}[1]{\markright{\thesection\ #1}}
%--- Khung tiêu đề (logo + tối đa 4 dòng; dòng 3, 4 để trống {} thì bỏ qua)
%  Đề thi:    \khungtieude{THI THỬ THPT QUỐC GIA NĂM 2026}{Môn Toán -- Đề số 21}
%                         {(Đề thi gồm: ...)}{Thời gian làm bài: 90 phút (không kể thời gian phát đề).}
%  Chuyên đề: \khungtieude{CHUYÊN ĐỀ 3 -- DÃY SỐ VÀ CÔNG THỨC TRUY HỒI}
%                         {Bồi dưỡng học sinh giỏi Toán THPT}{...}{}
\newcommand{\khungtieude}[4]{%
	\begin{tcolorbox}[colframe=gold1,colback=gold5!45,boxrule=0.5mm]
		\begin{minipage}[c]{3cm}
			\centerline{\eduzylogo[2.4cm]}
		\end{minipage}\hspace{0.5cm}
		\begin{minipage}[c]{12.5cm}
			\vskip 0.15 cm
			\centerline{\ind{#1}}
			\vskip 0.15 cm
			\centerline{\textbf{#2}}
			\if\relax\detokenize{#3}\relax\else
				\vskip 0.15 cm
				\centerline{\textit{#3}}
			\fi
			\if\relax\detokenize{#4}\relax\else
				\vskip 0.15 cm
				\centerline{\textit{#4}}
			\fi
			\vskip 0.15 cm
		\end{minipage}
	\end{tcolorbox}
}
%====================== BỔ SUNG CHO BÀI HỌC, CHUYÊN ĐỀ ======================
\usepackage[shortlabels]{enumitem}
\usepackage{tabularx}
%--- Danh sách gọn: \begin{enumerate}[gon] ... \end{enumerate}
\SetEnumitemKey{gon}{leftmargin=*,itemsep=1pt,parsep=0pt,topsep=3pt}
%--- Khung lý thuyết: \begin{khung}{} ... \end{khung}; tiêu đề chỉ dùng khi thật ngắn ("Mẹo", "Lỗi thường gặp")
\newtcolorbox{khung}[1]{breakable,enhanced,colframe=gold1,colback=gold5!45,
	boxrule=0.4mm,fonttitle=\bfseries\sffamily,coltitle=white,colbacktitle=gold1,
	attach boxed title to top left={xshift=6mm,yshift=-3mm},
	boxed title style={arc=1mm,boxrule=0.3mm},title={#1}}
%--- Tiêu đề mục nhỏ trong một dạng: \muc{Dấu hiệu}, \muc{Phương pháp}, \muc{Bài tự luyện}
\makeatletter
\newcommand{\muc}[1]{\par\medskip\noindent\indam{#1}\par\nobreak\smallskip\@afterindentfalse\@afterheading}
\makeatother
%--- Ô công thức cần nhớ (căn giữa, viền vàng, vừa khít nội dung)
%     \begin{center}\congthuc{$Q(n)=(n-1)!$}\end{center}
\newtcbox{\congthuc}{on line,colframe=gold1,colback=gold5!45,boxrule=0.5pt,arc=2pt,
	left=8pt,right=8pt,top=4pt,bottom=4pt}
%--- Công tắc lời giải bài tự luyện (ví dụ mẫu luôn dùng \loigiai)
%     \loigiaiTL[{đáp số}]{lời giải đầy đủ}
%     \TLfull -> in lời giải đầy đủ;  \TLdapso -> chỉ in một dòng "Đáp số";  \TLan -> ẩn hoàn toàn
%     (\CheDo là dethi thì mọi lời giải và đáp số đều ẩn)
\makeatletter
\newif\ifTLhien   \TLhientrue
\newif\ifTLdapso  \TLdapsofalse
\newcommand{\TLfull}{\TLhientrue\TLdapsofalse}
\newcommand{\TLdapso}{\TLhientrue\TLdapsotrue}
\newcommand{\TLan}{\TLhienfalse}
\newcommand{\loigiaiTL}[2][]{%
	\ifTLhien
		\ifTLdapso
			\def\@tlarg{#1}%
			\ifx\@tlarg\@empty\else
				\ifdefstring{\kindOPTN}{dethi}{}{\par\noindent{\bfseries\color{gold1}Đáp số:} #1\par}%
			\fi
		\else
			\loigiai{#2}%
		\fi
	\fi}
\makeatother
\TLdapso% mặc định: chỉ in đáp số bài tự luyện
%====================== BẢNG ĐÚNG/SAI (\choiceTFt) ======================
%  Hàng tiêu đề nền nâu đồng chữ trắng (Phát biểu | Đ | S), các hàng tô xen kẽ kem / trắng,
%  đường kẻ vàng kim, nhãn a) b) c) d) màu nâu đồng. Bản có đáp án: cột được chọn có vòng tròn
%  nâu đồng với dấu tích trắng (câu đúng tích ở cột Đúng, câu sai tích ở cột Sai).
%  Tiêu đề ghi đủ "Đúng | Sai": \OPTN{viettat=0} sau \input (cột Đúng/Sai rộng ra, phát biểu hẹp lại).
%  Dùng bảng gốc của ex_test: \def\BangDSgoc{1} trước \input{preamble-eduzy-v2}.
\providecommand{\BangDSgoc}{0}
\ifnum\BangDSgoc=0
\usepackage{colortbl,pifont}
\colorlet{ezDSke}{gold2!60!gold1}% đường kẻ bảng (nét 0,4pt như ex_test để không đổi chiều cao bảng)
\colorlet{ezDSle}{gold5!60}% hàng a), c)
\colorlet{ezDSchan}{white}% hàng b), d)
%--- dấu chọn ở cột Đúng/Sai
%    (bề rộng 0, căn giữa ô: không làm đổi bề rộng cột của ex_test, không báo tràn lề)
\newcommand{\EZdauchon}{\makebox[0pt]{\tikz[baseline=-1ex]\node[circle,fill=gold1,inner sep=0pt,
	minimum size=1.2em,text=white,font=\small]{\ding{51}};}}
\def\tickT{\EZdauchon}\def\tickF{\EZdauchon}
\AtBeginDocument{\sbox{\ccTFTab}{\EZdauchon}}
%--- nhãn a) b) c) d) màu nâu đồng (ở bản lời giải, nhãn câu đúng vẫn đóng khung như ex_test)
\renewcommand{\FalseTF}{\stepcounter{dapan}\textbf{\color{gold1}\DapAnTF\sepTF}\ignorespaces}
%--- trước khi dựng bảng: kẻ màu vàng kim (bề rộng cột giữ nguyên như ex_test); sau bảng trả lại màu đen
\newcount\EZhangTF
\newcommand{\EZbatdauTF}{%
	\global\EZhangTF=0\relax
	\arrayrulecolor{ezDSke}}
\newcommand{\EZketthucTF}{\arrayrulecolor{black}}
%--- hàng tiêu đề nền nâu đồng, chữ trắng
\makeatletter
\newcommand{\EZo@TF}[1]{\centering\bfseries\color{white}#1}
\newcommand{\EZtieudeTF}{%
	\ifcase\cotDS
		\ifnum\cotSTT=1\relax
			\def\rowHeader{\hline\rowcolor{gold1}\EZo@TF{TT}&\EZo@TF{\phatbieu}\tabularnewline}%
		\else
			\def\rowHeader{\hline\rowcolor{gold1}\EZo@TF{\phatbieu}\tabularnewline}%
		\fi
	\or
		\ifnum\cotSTT=1\relax
			\def\rowHeader{\hline\rowcolor{gold1}\EZo@TF{TT}&\EZo@TF{\phatbieu}&\EZo@TF{\Dung}\tabularnewline}%
		\else
			\def\rowHeader{\hline\rowcolor{gold1}\EZo@TF{\phatbieu}&\EZo@TF{\Dung}\tabularnewline}%
		\fi
	\or
		\ifnum\cotSTT=1\relax
			\def\rowHeader{\hline\rowcolor{gold1}\EZo@TF{TT}&\EZo@TF{\phatbieu}&\EZo@TF{\Dung}&\EZo@TF{\Sai}\tabularnewline}%
		\else
			\def\rowHeader{\hline\rowcolor{gold1}\EZo@TF{\phatbieu}&\EZo@TF{\Dung}&\EZo@TF{\Sai}\tabularnewline}%
		\fi
	\fi}
%--- tô xen kẽ từng hàng (gọi ngay sau \hline đầu mỗi hàng)
\newcommand{\EZmauhangTF}{%
	\noalign{\global\advance\EZhangTF\@ne}%
	\ifodd\EZhangTF\rowcolor{ezDSle}\else\rowcolor{ezDSchan}\fi}
%--- gắn vào \choiceTFt / \writeChoiceTFt của ex_test
\xpretocmd{\choiceTFt}{\EZbatdauTF}{}{\typeout{[preamble-eduzy-v2] Không gắn được bảng Đúng/Sai (đầu bảng)}}
\xpatchcmd{\choiceTFt}{\ifnum\dongPB=0\def\rowHeader{}\fi}{\EZtieudeTF\ifnum\dongPB=0\def\rowHeader{}\fi}
	{}{\typeout{[preamble-eduzy-v2] Không gắn được bảng Đúng/Sai (hàng tiêu đề)}}
\xpatchcmd{\writeChoiceTFt}{\hline\xdef\dotEX{\sdotEX}}{\hline\EZmauhangTF\xdef\dotEX{\sdotEX}}
	{}{\typeout{[preamble-eduzy-v2] Không gắn được bảng Đúng/Sai (tô hàng)}}
\xpatchcmd{\writeChoiceTFt}{{\end{longtable}}}{{\end{longtable}\EZketthucTF}}
	{}{\typeout{[preamble-eduzy-v2] Không gắn được bảng Đúng/Sai (cuối longtable)}}
\xpatchcmd{\writeChoiceTFt}{{\end{xtabular}\end{center}\medskip\noindent}}{{\end{xtabular}\end{center}\EZketthucTF\medskip\noindent}}
	{}{\typeout{[preamble-eduzy-v2] Không gắn được bảng Đúng/Sai (cuối xtabular)}}
\fi
%=====================================================================
\renewcommand{\parallel}{\mathrel{/\!\!/}}
.
\providecommand{\BangDSgoc}{0}
\ifnum\BangDSgoc=0
\usepackage{colortbl,pifont}
\colorlet{ezDSke}{gold2!60!gold1}% đường kẻ bảng (nét 0,4pt như ex_test để không đổi chiều cao bảng)
\colorlet{ezDSle}{gold5!60}% hàng a), c)
\colorlet{ezDSchan}{white}% hàng b), d)
%--- dấu chọn ở cột Đúng/Sai
%    (bề rộng 0, căn giữa ô: không làm đổi bề rộng cột của ex_test, không báo tràn lề)
\newcommand{\EZdauchon}{\makebox[0pt]{\tikz[baseline=-1ex]\node[circle,fill=gold1,inner sep=0pt,
	minimum size=1.2em,text=white,font=\small]{\ding{51}};}}
\def\tickT{\EZdauchon}\def\tickF{\EZdauchon}
\AtBeginDocument{\sbox{\ccTFTab}{\EZdauchon}}
%--- nhãn a) b) c) d) màu nâu đồng (ở bản lời giải, nhãn câu đúng vẫn đóng khung như ex_test)
\renewcommand{\FalseTF}{\stepcounter{dapan}\textbf{\color{gold1}\DapAnTF\sepTF}\ignorespaces}
%--- trước khi dựng bảng: kẻ màu vàng kim (bề rộng cột giữ nguyên như ex_test); sau bảng trả lại màu đen
\newcount\EZhangTF
\newcommand{\EZbatdauTF}{%
	\global\EZhangTF=0\relax
	\arrayrulecolor{ezDSke}}
\newcommand{\EZketthucTF}{\arrayrulecolor{black}}
%--- hàng tiêu đề nền nâu đồng, chữ trắng
\makeatletter
\newcommand{\EZo@TF}[1]{\centering\bfseries\color{white}#1}
\newcommand{\EZtieudeTF}{%
	\ifcase\cotDS
		\ifnum\cotSTT=1\relax
			\def\rowHeader{\hline\rowcolor{gold1}\EZo@TF{TT}&\EZo@TF{\phatbieu}\tabularnewline}%
		\else
			\def\rowHeader{\hline\rowcolor{gold1}\EZo@TF{\phatbieu}\tabularnewline}%
		\fi
	\or
		\ifnum\cotSTT=1\relax
			\def\rowHeader{\hline\rowcolor{gold1}\EZo@TF{TT}&\EZo@TF{\phatbieu}&\EZo@TF{\Dung}\tabularnewline}%
		\else
			\def\rowHeader{\hline\rowcolor{gold1}\EZo@TF{\phatbieu}&\EZo@TF{\Dung}\tabularnewline}%
		\fi
	\or
		\ifnum\cotSTT=1\relax
			\def\rowHeader{\hline\rowcolor{gold1}\EZo@TF{TT}&\EZo@TF{\phatbieu}&\EZo@TF{\Dung}&\EZo@TF{\Sai}\tabularnewline}%
		\else
			\def\rowHeader{\hline\rowcolor{gold1}\EZo@TF{\phatbieu}&\EZo@TF{\Dung}&\EZo@TF{\Sai}\tabularnewline}%
		\fi
	\fi}
%--- tô xen kẽ từng hàng (gọi ngay sau \hline đầu mỗi hàng)
\newcommand{\EZmauhangTF}{%
	\noalign{\global\advance\EZhangTF\@ne}%
	\ifodd\EZhangTF\rowcolor{ezDSle}\else\rowcolor{ezDSchan}\fi}
%--- gắn vào \choiceTFt / \writeChoiceTFt của ex_test
\xpretocmd{\choiceTFt}{\EZbatdauTF}{}{\typeout{[preamble-eduzy-v2] Không gắn được bảng Đúng/Sai (đầu bảng)}}
\xpatchcmd{\choiceTFt}{\ifnum\dongPB=0\def\rowHeader{}\fi}{\EZtieudeTF\ifnum\dongPB=0\def\rowHeader{}\fi}
	{}{\typeout{[preamble-eduzy-v2] Không gắn được bảng Đúng/Sai (hàng tiêu đề)}}
\xpatchcmd{\writeChoiceTFt}{\hline\xdef\dotEX{\sdotEX}}{\hline\EZmauhangTF\xdef\dotEX{\sdotEX}}
	{}{\typeout{[preamble-eduzy-v2] Không gắn được bảng Đúng/Sai (tô hàng)}}
\xpatchcmd{\writeChoiceTFt}{{\end{longtable}}}{{\end{longtable}\EZketthucTF}}
	{}{\typeout{[preamble-eduzy-v2] Không gắn được bảng Đúng/Sai (cuối longtable)}}
\xpatchcmd{\writeChoiceTFt}{{\end{xtabular}\end{center}\medskip\noindent}}{{\end{xtabular}\end{center}\EZketthucTF\medskip\noindent}}
	{}{\typeout{[preamble-eduzy-v2] Không gắn được bảng Đúng/Sai (cuối xtabular)}}
\fi
%=====================================================================
\renewcommand{\parallel}{\mathrel{/\!\!/}}

%      \TLdapso                      % tài liệu: \TLfull | \TLdapso | \TLan
%      \begin{document} ... \end{document}
%  Logo Eduzy vẽ bằng TikZ ngay trong file (\eduzylogo), không cần file ảnh.
%  Biên dịch: pdfLaTeX (máy có gói vietnam) hoặc XeLaTeX.
%============================================================
\documentclass[12pt,a4paper,twoside]{article}%<-- twoside: bat dau trang/chan trang chan-le
\usepackage{iftex}
\usepackage{amsmath,amssymb,tikz,mathrsfs,tkz-tab}
\usepackage{graphicx}
\usepackage{pgfplots}% đồ thị hàm số
\pgfplotsset{compat=1.17}
%\usepackage{mathptmx}
\usepackage{mathpazo}
\ifXeTeX% XeLaTeX: font chữ có đủ dấu tiếng Việt, giữ nguyên font toán mathpazo
	\usepackage[no-math]{fontspec}
	\setmainfont{TeX Gyre Pagella}
	\newfontfamily\EZqag{TeX Gyre Adventor}[NFSSFamily=qag,Ligatures=TeX]
\fi
\usepackage[many]{tcolorbox}
\usetikzlibrary{math,through,calc,intersections,angles,quotes,shapes,shapes.geometric,arrows,patterns,matrix,chains,arrows.meta,decorations.shapes,decorations.fractals,decorations.markings,shadows,decorations.pathreplacing}
\usepackage[top=1.6cm, bottom=1.8cm, left=1.8cm, right=1.8cm]{geometry}% v2: top 1.6cm (bản cũ 1.5cm) để chữ không sát dải đầu trang
\providecommand{\CheDo}{loigiai}% loigiai | dethi | solcolor | color | book
\expandafter\usepackage\expandafter[\CheDo]{ex_test}
% \tenchuong: khai báo ở file chính (trước \input) hoặc \renewcommand sau \input
%============================================================
%  [GỘP TỪ eduzy-de.tex] EDUZY - MÀU NÂU ĐỒNG / VÀNG KIM
%  + ĐẦU TRANG, CHÂN TRANG KIỂU DẢI NGANG (v2), ĐỐI XỨNG TRANG CHẴN / TRANG LẺ.
%  (không cần file eduzy-de.tex riêng nữa)
%  Khối này phải đặt SAU \usepackage[...]{ex_test} và SAU \def\tenchuong{...}
%============================================================
\usepackage{lastpage}
\usepackage{needspace}
%--- DÀN TRANG ĐỀ THI: tách câu (THÊM 05/11/2026; chỉ thêm lệnh mới, KHÔNG đổi hành vi mặc định) ---
%  Hai công tắc, mặc định tắt (= preamble chạy y như cũ). Bật bằng cách đặt TRƯỚC %============================================================
%  PREAMBLE EDUZY v2 - dùng chung cho đề thi, bài học, chuyên đề (gói ex_test)
%  Khác bản preamble-eduzy.tex ở đầu trang, chân trang: dải nền nâu đồng tràn ngang,
%  logo + Eduzy.com.vn, ô chữ vàng kem cạnh xiên, chân trang "Trang (số)";
%  logo Eduzy chìm giữa mọi trang (tắt: \logochimtat, bật lại: \logochimbat).
%  Cách dùng ở file chính (xem mau-chuan.tex, mau-chuyen-de.tex):
%      \def\CheDo{loigiai}          % loigiai | dethi | solcolor | color | book
%      \def\tenchuong{<dòng chữ trên đầu trang>}
%      \def\thuonghieu{<chữ ở chân trang lẻ>}   % không khai báo thì dùng mặc định
%      \def\tacgia{<chữ ở chân trang chẵn>}      % không khai báo thì dùng mặc định
%      \def\khauhieu{<chữ cạnh logo ở đầu trang>}  % tùy chọn, mặc định Eduzy.com.vn
%      %============================================================
%  PREAMBLE EDUZY v2 - dùng chung cho đề thi, bài học, chuyên đề (gói ex_test)
%  Khác bản preamble-eduzy.tex ở đầu trang, chân trang: dải nền nâu đồng tràn ngang,
%  logo + Eduzy.com.vn, ô chữ vàng kem cạnh xiên, chân trang "Trang (số)";
%  logo Eduzy chìm giữa mọi trang (tắt: \logochimtat, bật lại: \logochimbat).
%  Cách dùng ở file chính (xem mau-chuan.tex, mau-chuyen-de.tex):
%      \def\CheDo{loigiai}          % loigiai | dethi | solcolor | color | book
%      \def\tenchuong{<dòng chữ trên đầu trang>}
%      \def\thuonghieu{<chữ ở chân trang lẻ>}   % không khai báo thì dùng mặc định
%      \def\tacgia{<chữ ở chân trang chẵn>}      % không khai báo thì dùng mặc định
%      \def\khauhieu{<chữ cạnh logo ở đầu trang>}  % tùy chọn, mặc định Eduzy.com.vn
%      \input{preamble-eduzy-v2}
%      \TLdapso                      % tài liệu: \TLfull | \TLdapso | \TLan
%      \begin{document} ... \end{document}
%  Logo Eduzy vẽ bằng TikZ ngay trong file (\eduzylogo), không cần file ảnh.
%  Biên dịch: pdfLaTeX (máy có gói vietnam) hoặc XeLaTeX.
%============================================================
\documentclass[12pt,a4paper,twoside]{article}%<-- twoside: bat dau trang/chan trang chan-le
\usepackage{iftex}
\usepackage{amsmath,amssymb,tikz,mathrsfs,tkz-tab}
\usepackage{graphicx}
\usepackage{pgfplots}% đồ thị hàm số
\pgfplotsset{compat=1.17}
%\usepackage{mathptmx}
\usepackage{mathpazo}
\ifXeTeX% XeLaTeX: font chữ có đủ dấu tiếng Việt, giữ nguyên font toán mathpazo
	\usepackage[no-math]{fontspec}
	\setmainfont{TeX Gyre Pagella}
	\newfontfamily\EZqag{TeX Gyre Adventor}[NFSSFamily=qag,Ligatures=TeX]
\fi
\usepackage[many]{tcolorbox}
\usetikzlibrary{math,through,calc,intersections,angles,quotes,shapes,shapes.geometric,arrows,patterns,matrix,chains,arrows.meta,decorations.shapes,decorations.fractals,decorations.markings,shadows,decorations.pathreplacing}
\usepackage[top=1.6cm, bottom=1.8cm, left=1.8cm, right=1.8cm]{geometry}% v2: top 1.6cm (bản cũ 1.5cm) để chữ không sát dải đầu trang
\providecommand{\CheDo}{loigiai}% loigiai | dethi | solcolor | color | book
\expandafter\usepackage\expandafter[\CheDo]{ex_test}
% \tenchuong: khai báo ở file chính (trước \input) hoặc \renewcommand sau \input
%============================================================
%  [GỘP TỪ eduzy-de.tex] EDUZY - MÀU NÂU ĐỒNG / VÀNG KIM
%  + ĐẦU TRANG, CHÂN TRANG KIỂU DẢI NGANG (v2), ĐỐI XỨNG TRANG CHẴN / TRANG LẺ.
%  (không cần file eduzy-de.tex riêng nữa)
%  Khối này phải đặt SAU \usepackage[...]{ex_test} và SAU \def\tenchuong{...}
%============================================================
\usepackage{lastpage}
\usepackage{needspace}
%--- DÀN TRANG ĐỀ THI: tách câu (THÊM 05/11/2026; chỉ thêm lệnh mới, KHÔNG đổi hành vi mặc định) ---
%  Hai công tắc, mặc định tắt (= preamble chạy y như cũ). Bật bằng cách đặt TRƯỚC \input{preamble-eduzy-v2}
%  trong file đề (tools/dan_trang/dan_trang.py tự ghi khi đặt \NeedspaceCau):
%    \def\TachCauDS{1}   câu Đúng/Sai được tách: câu dẫn (kèm hình) ở cuối trang, bảng Phát biểu sang trang sau
%                        (câu dẫn và hình không tách nhau, bảng không tách giữa hai trang)
%    \def\TachCauLG{1}   bản lời giải (\CheDo khác dethi): phần đề của câu cao hơn 1/3 trang được phép tách
%  Lệnh dùng trong file đề (giá trị do tools/dan_trang/dan_trang.py đo ở bản dethi):
%    \NeedspaceCau[<cao câu dẫn + hình>]{<cao cả khối câu>}   ngay trên \begin{ex} (câu đầu phần: trên tiêu đề phần)
%        công tắc tắt: đúng bằng \Needspace{<cao cả khối câu>} như cũ;
%        \TachCauDS=1 và có [..] (câu Đúng/Sai được chọn tách): chỉ giữ câu dẫn + hình;
%        \TachCauLG=1 ở bản lời giải: đòi tối đa 1/3 \textheight.
%    \NeedspaceBang{<cao bảng Phát biểu>}%   ngay trước \choiceTFt: \TachCauDS=1 thì giữ trọn bảng trên một trang;
%        công tắc tắt: không làm gì.
\providecommand{\TachCauDS}{0}
\providecommand{\TachCauLG}{0}
\newlength{\EZnsCau}
\newcommand{\NeedspaceCau}[2][]{%
	\setlength{\EZnsCau}{#2}%
	\ifnum\TachCauDS=1 \if\relax\detokenize{#1}\relax\else\setlength{\EZnsCau}{#1}\fi\fi
	\ifnum\TachCauLG=1 \ifdefstring{\kindOPTN}{dethi}{}{\ifdim\EZnsCau>0.3333\textheight\setlength{\EZnsCau}{0.3333\textheight}\fi}\fi
	\Needspace{\EZnsCau}}
\newcommand{\NeedspaceBang}[1]{\ifnum\TachCauDS=1 \Needspace{#1}\fi}

%--- Thang màu: đậm -> nhạt
\definecolor{gold1}{RGB}{126,96,0}    % nâu đồng đậm  #7E6000
\definecolor{gold2}{RGB}{192,144,3}   % vàng kim      #C09003
\definecolor{gold3}{RGB}{255,215,99}  % vàng nhạt     #FFD763
\definecolor{gold4}{RGB}{255,231,153} % vàng rất nhạt #FFE799
\definecolor{gold5}{RGB}{255,243,205} % vàng kem      #FFF3CD
%--- Ghi đè tông xanh/đỏ cũ
\definecolor{myblue}{RGB}{126,96,0}
\colorlet{blue}{gold1}
\colorlet{cyan}{gold2}
\colorlet{teal}{gold2}

%--- Tên chương mặc định (file đề nên tự khai báo lại)
\providecommand{\tenchuong}{}
%--- Chân trang
\providecommand{\thuonghieu}{TOÁN 12 -- Eduzy.com.vn}
\providecommand{\tacgia}{Tô Gia Bảo -- 0353703783}

%--- Chữ cạnh logo ở đầu trang
\providecommand{\khauhieu}{Eduzy.com.vn}

%====================== THÔNG SỐ ĐẦU TRANG, CHÂN TRANG (v2) ======================
%  Dải nền đậm tràn ngang trang; ô chữ vàng kem có cạnh xiên, kèm 2 vạch xiên chuyển màu.
%  Đầu trang lẻ : logo + \khauhieu (trái)   | ô chữ \tenchuong (phải)
%  Chân trang lẻ: ô chữ \thuonghieu (trái)  | Trang (số) (phải)
%  Trang chẵn soi gương trang lẻ; chân trang chẵn ghi \tacgia.
%--- Màu: chỉ dùng thang vàng kim Eduzy
\colorlet{ezNen}{gold1}% dải nền đậm
\colorlet{ezVach1}{gold3}% vạch xiên sát dải nền (đậm hơn)
\colorlet{ezVach2}{gold4}% vạch xiên sát ô chữ (nhạt hơn)
\colorlet{ezO}{gold5}% ô chữ vàng kem
%--- Kích thước
\newlength{\EZdai}\setlength{\EZdai}{1.1cm}% chiều cao dải nền
\newlength{\EZo}\setlength{\EZo}{0.68cm}% chiều cao ô chữ
\newlength{\EZxien}\setlength{\EZxien}{0.55cm}% độ lệch ngang của cạnh xiên
\newlength{\EZvach}\setlength{\EZvach}{0.75cm}% bề rộng mỗi vạch xiên
\newlength{\EZkhe}\setlength{\EZkhe}{0.12cm}% khe trắng giữa vạch và ô chữ
\newlength{\EZdem}\setlength{\EZdem}{0.4cm}% khoảng đệm từ chữ tới cạnh xiên
\newlength{\EZle}\setlength{\EZle}{1.8cm}% lề trái/phải của chữ (bằng lề thân bài)
\newlength{\EZrongmin}\setlength{\EZrongmin}{4.5cm}% ô chữ rộng tối thiểu
\newlength{\EZrong}% bề rộng chữ trong ô (tính tự động)
%--- Font
\newcommand{\EZfontO}{\fontsize{10.5pt}{12pt}\selectfont\bfseries}% chữ trong ô
\newcommand{\EZfontKH}{\fontfamily{qag}\fontsize{11pt}{12pt}\selectfont\bfseries}% \khauhieu cạnh logo
\newcommand{\EZfontTrang}{\fontsize{11.5pt}{12pt}\selectfont\bfseries}% chữ "Trang" và số trang
%--- Đo bề rộng chữ trong ô
\newcommand{\EZdorong}[1]{%
	\settowidth{\EZrong}{\EZfontO #1}%
	\ifdim\EZrong<\EZrongmin\setlength{\EZrong}{\EZrongmin}\fi}
%--- Logo + chữ trắng \khauhieu (dựng sẵn một lần khi vào \begin{document})
\newsavebox{\EZlogobox}
\AtBeginDocument{\sbox{\EZlogobox}{%
		\raisebox{-0.5\height}{\eduzylogo[1cm]}\hspace{0.15cm}%
		\raisebox{-0.5\height}{\color{white}\EZfontKH\khauhieu}}}
%--- Hai vạch xiên + khe trắng. #1: điểm gốc ở mép dưới, #2: hướng (1 sang phải, -1 sang trái),
%    #3: hướng nghiêng (1 kiểu "/", -1 kiểu "\")
\newcommand{\EZvachxien}[3]{%
	\fill[white] (#1) --++({(#3)*\EZxien},\EZo) --++({(#2)*\EZkhe},0) --++({-(#3)*\EZxien},-\EZo) --cycle;
	\fill[ezVach2] ([xshift=(#2)*\EZkhe]#1) --++({(#3)*\EZxien},\EZo) --++({(#2)*\EZvach},0) --++({-(#3)*\EZxien},-\EZo) --cycle;
	\fill[ezVach1] ([xshift=(#2)*(\EZkhe+\EZvach)]#1) --++({(#3)*\EZxien},\EZo) --++({(#2)*\EZvach},0) --++({-(#3)*\EZxien},-\EZo) --cycle;
	\draw[white,line width=0.6pt] ([xshift=(#2)*(\EZkhe+\EZvach)]#1) --++({(#3)*\EZxien},\EZo);
}
%--- Chữ trong ô (căn giữa theo chiều cao ô, không lệch vì dấu tiếng Việt)
\tikzset{ezchu/.style={inner sep=0pt,text=ezNen,text height=0.72em,text depth=0.22em}}

%====================== LOGO CHÌM GIỮA TRANG ======================
%  Vẽ cùng đầu trang nên nằm dưới chữ của thân bài, có trên mọi trang.
%  Tắt: \logochimtat (từ trang đó trở đi), bật lại: \logochimbat.
%  Đổi cỡ, độ đậm: khai báo trước \input, ví dụ \def\EZchimcao{9cm} \def\EZchimdo{0.1}
\providecommand{\EZchimcao}{10cm}% chiều cao logo chìm
\providecommand{\EZchimdo}{0.12}% độ đậm (0 = trong suốt, 1 = rõ hẳn)
\newif\ifEZchim \EZchimtrue
\newcommand{\logochimtat}{\EZchimfalse}
\newcommand{\logochimbat}{\EZchimtrue}
%  Dựng sẵn logo mờ một lần khi vào \begin{document}; transparency group làm mờ cả logo
%  một lượt nên không lộ vệt ở chỗ các dải màu chồng nhau.
\newsavebox{\EZchimtho}
\newsavebox{\EZchimbox}
\AtBeginDocument{%
	\sbox{\EZchimtho}{\eduzylogo[\EZchimcao]}%
	\sbox{\EZchimbox}{\begin{tikzpicture}
			\begin{scope}[transparency group,opacity=\EZchimdo]
				\node[inner sep=0pt] {\usebox{\EZchimtho}};
			\end{scope}
	\end{tikzpicture}}}
\newcommand{\EZvechim}{%
	\ifEZchim
		\node[inner sep=0pt] at (current page.center) {\usebox{\EZchimbox}};
	\fi}

%====================== ĐẦU TRANG - TRANG LẺ ======================
%  (logo bên trái, ô chữ \tenchuong bên phải)
\newcommand{\myfancyheadOdd}{%
	\EZdorong{\tenchuong}%
	\begin{tikzpicture}[remember picture,overlay]
		\coordinate (T) at (current page.north west);
		\coordinate (P) at (current page.north east);
		%--- logo chìm giữa trang
		\EZvechim
		%--- dải nền
		\fill[ezNen] (T) rectangle ([yshift=-\EZdai]P);
		%--- ô chữ: X là góc dưới bên trái, cạnh trái xiên kiểu "/"
		\coordinate (X) at ([xshift=-\EZle-\EZrong-\EZdem-\EZxien,yshift=-\EZdai]P);
		\fill[ezO] (X) --++(\EZxien,\EZo) -- ([yshift=-\EZdai+\EZo]P) -- ([yshift=-\EZdai]P) --cycle;
		\EZvachxien{X}{-1}{1}
		\node[ezchu,anchor=east] at ([xshift=-\EZle,yshift=-\EZdai+0.5*\EZo]P) {\EZfontO\tenchuong};
		%--- logo + tên
		\node[inner sep=0pt,anchor=west] at ([xshift=\EZle,yshift=-0.5*\EZdai]T) {\usebox{\EZlogobox}};
	\end{tikzpicture}%
}

%====================== ĐẦU TRANG - TRANG CHẴN =====================
%  (soi gương: ô chữ \tenchuong bên trái, logo bên phải)
\newcommand{\myfancyheadEven}{%
	\EZdorong{\tenchuong}%
	\begin{tikzpicture}[remember picture,overlay]
		\coordinate (T) at (current page.north west);
		\coordinate (P) at (current page.north east);
		%--- logo chìm giữa trang
		\EZvechim
		%--- dải nền
		\fill[ezNen] (T) rectangle ([yshift=-\EZdai]P);
		%--- ô chữ: X là góc dưới bên phải, cạnh phải xiên kiểu "\"
		\coordinate (X) at ([xshift=\EZle+\EZrong+\EZdem+\EZxien,yshift=-\EZdai]T);
		\fill[ezO] (X) --++(-\EZxien,\EZo) -- ([yshift=-\EZdai+\EZo]T) -- ([yshift=-\EZdai]T) --cycle;
		\EZvachxien{X}{1}{-1}
		\node[ezchu,anchor=west] at ([xshift=\EZle,yshift=-\EZdai+0.5*\EZo]T) {\EZfontO\tenchuong};
		%--- logo + tên
		\node[inner sep=0pt,anchor=east] at ([xshift=-\EZle,yshift=-0.5*\EZdai]P) {\usebox{\EZlogobox}};
	\end{tikzpicture}%
}

%====================== CHÂN TRANG - TRANG LẺ ======================
%  (ô chữ \thuonghieu bên trái, "Trang (số)" bên phải)
\newcommand{\myfancyfootOdd}{%
	\EZdorong{\thuonghieu}%
	\begin{tikzpicture}[remember picture,overlay]
		\coordinate (T) at (current page.south west);
		\coordinate (P) at (current page.south east);
		%--- dải nền
		\fill[ezNen] (T) rectangle ([yshift=\EZdai]P);
		%--- ô chữ sát mép trên dải: X là góc dưới bên phải, cạnh phải xiên kiểu "/"
		\coordinate (X) at ([xshift=\EZle+\EZrong+\EZdem,yshift=\EZdai-\EZo]T);
		\fill[ezO] ([yshift=\EZdai-\EZo]T) -- (X) --++(\EZxien,\EZo) -- ([yshift=\EZdai]T) --cycle;
		\EZvachxien{X}{1}{1}
		\node[ezchu,anchor=west] at ([xshift=\EZle,yshift=\EZdai-0.5*\EZo]T) {\EZfontO\thuonghieu};
		%--- số trang
		\node[circle,draw=white,line width=0.9pt,text=white,inner sep=0pt,minimum size=0.62cm,anchor=east]
		(so) at ([xshift=-\EZle,yshift=0.5*\EZdai]P) {\EZfontTrang\thepage};
		\node[inner sep=0pt,text=white,anchor=base east] (tr) at ([xshift=-0.15cm]so.base west) {\EZfontTrang Trang};
	\end{tikzpicture}%
}

%====================== CHÂN TRANG - TRANG CHẴN ====================
%  (soi gương: "Trang (số)" bên trái, ô chữ \tacgia bên phải)
\newcommand{\myfancyfootEven}{%
	\EZdorong{\tacgia}%
	\begin{tikzpicture}[remember picture,overlay]
		\coordinate (T) at (current page.south west);
		\coordinate (P) at (current page.south east);
		%--- dải nền
		\fill[ezNen] (T) rectangle ([yshift=\EZdai]P);
		%--- ô chữ sát mép trên dải: X là góc dưới bên trái, cạnh trái xiên kiểu "\"
		\coordinate (X) at ([xshift=-\EZle-\EZrong-\EZdem,yshift=\EZdai-\EZo]P);
		\fill[ezO] ([yshift=\EZdai-\EZo]P) -- (X) --++(-\EZxien,\EZo) -- ([yshift=\EZdai]P) --cycle;
		\EZvachxien{X}{-1}{-1}
		\node[ezchu,anchor=east] at ([xshift=-\EZle,yshift=\EZdai-0.5*\EZo]P) {\EZfontO\tacgia};
		%--- số trang
		%    (vòng tròn ẩn so0 nằm giữa dải để lấy đường chân chữ, giống trang lẻ)
		\node[circle,inner sep=0pt,minimum size=0.62cm] (so0) at ([yshift=0.5*\EZdai]T) {\EZfontTrang\phantom{\thepage}};
		\node[inner sep=0pt,text=white,anchor=base west] (tr) at ($(so0.base -| T)+(\EZle,0)$) {\EZfontTrang Trang};
		\node[circle,draw=white,line width=0.9pt,text=white,inner sep=0pt,minimum size=0.62cm,anchor=base west]
		(so) at ([xshift=0.15cm]tr.base east) {\EZfontTrang\thepage};
	\end{tikzpicture}%
}

%--- Giữ tương thích với các file cũ đang gọi \myfancyhead, \myfancyfoot
\newcommand{\myfancyhead}{\myfancyheadOdd}
\newcommand{\myfancyfoot}{\myfancyfootOdd}

%======================= GẮN VÀO TRANG =============================
\pagestyle{fancy}
\fancyhf{}
\fancyhead[LO]{\myfancyheadOdd}
\fancyhead[RE]{\myfancyheadEven}
\fancyfoot[LO]{\myfancyfootOdd}
\fancyfoot[RE]{\myfancyfootEven}
\renewcommand{\headrulewidth}{0pt}
\renewcommand{\footrulewidth}{0pt}
\fancypagestyle{plain}{%
	\fancyhf{}
	\fancyhead[LO]{\myfancyheadOdd}
	\fancyhead[RE]{\myfancyheadEven}
	\fancyfoot[LO]{\myfancyfootOdd}
	\fancyfoot[RE]{\myfancyfootEven}
	\renewcommand{\headrulewidth}{0pt}
	\renewcommand{\footrulewidth}{0pt}
}
%=================== HẾT PHẦN GỘP TỪ eduzy-de.tex ===================
\newcommand{\indam}[1]{
	{\fontsize{12pt}{9pt}\sffamily \selectfont\bfseries\color{gold1} #1}
}
\newcommand{\ind}[1]{
	{\fontsize{14pt}{9pt}\fontfamily{qag} \selectfont\bfseries\color{gold1} #1}
}
\renewcommand{\FalseEX}{\stepcounter{dapan}\textbf{\selectfont\color{gold1}\Alph{dapan}.}\ignorespaces}
\renewtheorem{ex}{\color{gold1}Câu}
\renewcommand{\baselinestretch}{1.1}
\newcommand{\hoac}[1]{
	\left[\begin{aligned}#1\end{aligned}\right.}
\newcommand{\heva}[1]{
	\left\{\begin{aligned}#1\end{aligned}\right.}
\usepackage{esvect}
\def\vec{\vv}
\def\vt{\vv}
\def\overrightarrow{\vv}
%====================== LOGO EDUZY VẼ BẰNG TIKZ ======================
%  Đường viền dò từ logo gốc; dải màu vàng kim đo trực tiếp trên ảnh gốc
%  (nghiêng 21 độ, tô bằng 80 dải hẹp nên hiển thị giống nhau trên mọi trình xem PDF).
%  Dùng: \eduzylogo          -> cao 2.4cm (mặc định)
%        \eduzylogo[1.5cm]   -> cao 1.5cm
\definecolor{ezText}{RGB}{245,200,46}
\newcommand{\eduzylogo}[1][2.4cm]{%
	\resizebox{!}{#1}{%
		\begin{tikzpicture}[x=0.01cm,y=0.01cm]
			\path[use as bounding box] (0,0) rectangle (224,-224);
			% mũ tốt nghiệp + vương miện
			\begin{scope}
				\clip
				(119.4,-176.1)..controls(118.9,-176.1)and(116.8,-176)..(114.6,-176)
				..controls(101.4,-176)and(85.8,-175.9)..(83.7,-175.7)
				..controls(81.8,-175.6)and(81,-175.6)..(80.2,-175.7)
				..controls(78.8,-175.9)and(77.7,-175.9)..(76.2,-175.7)
				..controls(74.9,-175.5)and(70.8,-175.5)..(64.3,-175.7)
				..controls(59.1,-175.9)and(58.3,-175.7)..(58.6,-173.7)
				..controls(58.6,-172.8)and(59.6,-169.6)..(60.3,-168)
				..controls(60.8,-166.9)and(61.4,-165.8)..(62.3,-165)
				..controls(63.1,-164.2)and(64.3,-163.4)..(65,-163.2)
				..controls(65.5,-163.1)and(72.3,-163)..(76.6,-163.1)
				..controls(77.9,-163.1)and(80.1,-163.1)..(81.4,-163)
				..controls(82.8,-162.9)and(86.3,-162.9)..(89.2,-162.9)
				..controls(93,-162.9)and(94.9,-162.8)..(96,-162.7)
				..controls(97.2,-162.5)and(101,-162.5)..(114,-162.4)
				..controls(123,-162.3)and(134.9,-162.2)..(140.2,-162.1)
				..controls(150.1,-162)and(152.2,-162.1)..(153.6,-162.4)
				..controls(154.4,-162.6)and(157.1,-163.8)..(158.1,-164.4)
				..controls(158.8,-164.8)and(159.7,-165.7)..(159.9,-166.3)
				..controls(160,-166.5)and(160.2,-167.1)..(160.2,-167.6)
				..controls(160.3,-168)and(160.6,-169.4)..(160.9,-170.6)
				..controls(161.2,-171.7)and(161.5,-173.1)..(161.6,-173.6)
				..controls(161.8,-174.4)and(161.8,-174.5)..(161.5,-174.9)
				..controls(161.2,-175.6)and(160.8,-175.9)..(159.9,-176)
				..controls(159.2,-176.1)and(154.1,-176.2)..(151.5,-176.2)
				..controls(148.1,-176.1)and(122.6,-176.1)..(121.6,-176.2)
				..controls(120.8,-176.2)and(119.8,-176.2)..(119.4,-176.1)--cycle
				(130.4,-160)..controls(129,-160)and(119.8,-160)..(110.1,-159.9)
				..controls(100.3,-159.8)and(88.5,-159.7)..(83.9,-159.6)
				..controls(77.8,-159.5)and(75.3,-159.5)..(74.5,-159.6)
				..controls(73.9,-159.7)and(72.5,-159.8)..(71.3,-159.8)
				..controls(69.6,-159.8)and(69,-159.8)..(68.5,-159.6)
				..controls(67.6,-159.3)and(66.8,-158.6)..(66.2,-157.6)
				..controls(66,-157.3)and(65.6,-156.5)..(65.3,-156.1)
				..controls(64.6,-155.1)and(64.5,-154.9)..(63.7,-152.9)
				..controls(63.2,-151.6)and(62.9,-150.9)..(61.8,-149.6)
				..controls(61.5,-149.2)and(61.1,-148.5)..(60.9,-147.9)
				..controls(60.7,-147.4)and(60.3,-146.6)..(60.1,-146.1)
				..controls(59.9,-145.5)and(59.5,-144.9)..(59,-144.3)
				..controls(58.5,-143.7)and(58.2,-143.2)..(57.9,-142.5)
				..controls(57.7,-141.9)and(57.3,-141.1)..(57,-140.7)
				..controls(56.6,-140.3)and(56.1,-139.3)..(55.8,-138.5)
				..controls(55.5,-137.8)and(55.1,-136.8)..(54.9,-136.5)
				..controls(54.7,-136.1)and(54.3,-135.2)..(53.9,-134.5)
				..controls(53.6,-133.8)and(53,-132.7)..(52.6,-132)
				..controls(51.6,-130.5)and(51.4,-130.2)..(51.1,-129)
				..controls(50.6,-127.5)and(50.6,-127.4)..(49.7,-126.3)
				..controls(49.2,-125.7)and(48.6,-124.8)..(48.3,-124.2)
				..controls(48,-123.6)and(47.5,-122.7)..(47.3,-122.2)
				..controls(47,-121.8)and(46.5,-120.7)..(46.1,-119.8)
				..controls(45.8,-118.9)and(45.2,-117.8)..(44.9,-117.2)
				..controls(44.2,-116.1)and(43.3,-114.4)..(42.3,-112.2)
				..controls(41.8,-111.3)and(41.3,-110.4)..(41.1,-110.1)
				..controls(40.5,-109.4)and(40.3,-109)..(40,-107.8)
				..controls(39.8,-107)and(39.6,-106.5)..(39.2,-106)
				..controls(38.9,-105.6)and(38.4,-104.7)..(38,-104)
				..controls(37.7,-103.4)and(37.2,-102.5)..(37,-102.1)
				..controls(36.7,-101.7)and(36.3,-100.8)..(36,-100.1)
				..controls(35.6,-98.9)and(35.3,-98.4)..(34.6,-97.6)
				..controls(34.2,-97.1)and(34.2,-96.3)..(34.6,-96)
				..controls(35,-95.7)and(35.5,-95.9)..(36.2,-96.7)
				..controls(36.5,-97)and(37.2,-97.5)..(37.6,-97.7)..controls(38,-98)and(38.6,-98.4)..(39.1,-98.9)
				..controls(39.6,-99.4)and(40.3,-100)..(40.6,-100.3)
				..controls(41.3,-100.8)and(42.4,-101.8)..(43.7,-102.9)
				..controls(44,-103.2)and(44.6,-103.7)..(45,-103.9)
				..controls(45.4,-104.2)and(45.9,-104.6)..(46.1,-104.8)
				..controls(47.1,-105.9)and(48.8,-107.1)..(50.5,-107.9)
				..controls(51.8,-108.6)and(52.3,-109)..(53.4,-110)
				..controls(54.8,-111.2)and(55.6,-111.6)..(57,-112.1)
				..controls(57.5,-112.3)and(58.5,-112.8)..(59.4,-113.3)
				..controls(60.2,-113.7)and(61.3,-114.3)..(61.9,-114.5)
				..controls(62.4,-114.7)and(63.3,-115.1)..(63.8,-115.3)
				..controls(64.4,-115.6)and(65.4,-115.9)..(66.2,-116.1)
				..controls(67,-116.3)and(68.4,-116.6)..(69.2,-116.9)
				..controls(72.3,-117.9)and(72.8,-118)..(78,-118.1)
				..controls(84,-118.3)and(86.1,-118.1)..(88.1,-117.1)
				..controls(89.4,-116.5)and(93,-115)..(93.7,-114.9)
				..controls(94.8,-114.6)and(95.4,-114.2)..(95.9,-113.4)
				..controls(96.2,-113)and(96.7,-112.5)..(97.1,-112.2)
				..controls(98,-111.6)and(101.3,-108.4)..(102.1,-107.2)
				..controls(102.4,-106.8)and(102.9,-106.2)..(103.2,-105.8)
				..controls(103.5,-105.5)and(104,-104.8)..(104.2,-104.2)
				..controls(104.5,-103.6)and(105,-102.8)..(105.3,-102.4)
				..controls(105.6,-102)and(106.1,-101.2)..(106.5,-100.6)
				..controls(107.1,-99.5)and(108.6,-97.5)..(109.5,-96.5)
				..controls(109.9,-96.1)and(110.1,-96)..(110.4,-96)
				..controls(111.1,-96)and(111.6,-96.5)..(111.8,-97.6)
				..controls(112,-98.6)and(112.3,-99.3)..(112.8,-99.9)
				..controls(113.7,-100.9)and(115.2,-103.1)..(115.7,-104.2)
				..controls(116.1,-105)and(116.5,-105.5)..(116.9,-106)
				..controls(117.3,-106.4)and(117.8,-106.9)..(118,-107.3)
				..controls(118.2,-107.6)and(118.6,-108)..(119,-108.3)
				..controls(119.6,-108.8)and(121.1,-110.3)..(121.8,-111.2)
				..controls(122.1,-111.7)and(122.5,-112)..(123.4,-112.4)
				..controls(124.8,-113.2)and(125.8,-113.8)..(127,-114.6)
				..controls(127.5,-114.9)and(128.1,-115.3)..(128.2,-115.4)
				..controls(128.4,-115.5)and(129.3,-115.8)..(130.2,-116)
				..controls(131.1,-116.2)and(132.5,-116.6)..(133.3,-116.9)
				..controls(135.7,-117.9)and(136.1,-117.9)..(142.5,-117.8)
				..controls(145.8,-117.8)and(148.4,-117.7)..(148.8,-117.6)
				..controls(149.2,-117.6)and(150.2,-117.2)..(151,-116.8)
				..controls(152,-116.4)and(152.8,-116)..(153.6,-115.9)
				..controls(154.5,-115.7)and(155.2,-115.4)..(156.3,-114.9)
				..controls(157.1,-114.5)and(158.2,-114)..(158.7,-113.8)
				..controls(159.1,-113.7)and(159.8,-113.3)..(160.2,-113)
				..controls(160.6,-112.7)and(161.3,-112.4)..(161.7,-112.2)
				..controls(162.7,-111.9)and(163.2,-111.7)..(164.1,-111.1)
				..controls(164.5,-110.8)and(165.3,-110.4)..(165.9,-110)
				..controls(166.5,-109.7)and(167.3,-109.2)..(167.6,-108.9)
				..controls(167.9,-108.6)and(168.6,-108.2)..(169,-108)
				..controls(169.4,-107.8)and(170.2,-107.3)..(170.7,-106.9)
				..controls(171.2,-106.5)and(172.1,-105.8)..(172.7,-105.4)
				..controls(173.2,-105)and(174.1,-104.4)..(174.5,-104)
				..controls(174.9,-103.7)and(175.7,-103.1)..(176.2,-102.6)
				..controls(176.7,-102.2)and(177.4,-101.5)..(177.9,-101.2)
				..controls(178.3,-100.8)and(178.8,-100.4)..(179,-100.2)
				..controls(179.2,-100)and(179.6,-99.6)..(180,-99.3)
				..controls(180.4,-99.1)and(180.9,-98.6)..(181.2,-98.4)
				..controls(181.5,-98.1)and(182.1,-97.7)..(182.5,-97.5)
				..controls(183.2,-97.2)and(183.5,-96.9)..(184.3,-95.9)
				..controls(184.6,-95.6)and(184.6,-95.6)..(184.8,-95.7)
				..controls(185,-95.8)and(185.1,-96)..(185.1,-96.2)
				..controls(185.1,-96.5)and(185,-96.7)..(184.6,-97.2)
				..controls(184.1,-97.8)and(184,-97.9)..(183.9,-98.8)
				..controls(183.7,-100.3)and(183.3,-101.3)..(182.9,-101.8)
				..controls(182.6,-102.1)and(182.3,-102.6)..(182.2,-102.9)
				..controls(182,-103.3)and(181.6,-104)..(181.2,-104.6)
				..controls(180.9,-105.2)and(180.5,-106.1)..(180.4,-106.6)
				..controls(179.9,-108)and(179.7,-108.5)..(179,-109.4)
				..controls(178.6,-109.8)and(178.2,-110.5)..(177.9,-111)
				..controls(177.7,-111.5)and(177.3,-112.3)..(177,-112.7)
				..controls(176.6,-113.2)and(176.4,-113.7)..(176.3,-114.2)
				..controls(176.2,-114.7)and(175.9,-115.2)..(175.4,-116)
				..controls(175,-116.6)and(174.4,-117.7)..(174,-118.5)
				..controls(173.7,-119.2)and(173.3,-120.1)..(173,-120.5)
				..controls(172.8,-120.9)and(172.4,-121.8)..(172.1,-122.5)
				..controls(171.8,-123.2)and(171.3,-124.1)..(171,-124.6)
				..controls(170.8,-125.1)and(170.3,-126)..(169.9,-126.6)
				..controls(169.6,-127.3)and(169.2,-128.1)..(169,-128.3)
				..controls(168.8,-128.6)and(168.5,-129.3)..(168.3,-129.8)
				..controls(168.1,-130.3)and(167.8,-131.1)..(167.5,-131.6)
				..controls(166.9,-132.5)and(164.2,-137.9)..(164.2,-138.1)
				..controls(164.2,-138.1)and(163.9,-138.8)..(163.4,-139.7)
				..controls(162.9,-140.5)and(162.4,-141.5)..(162.3,-141.9)
				..controls(161.9,-143)and(161.6,-143.4)..(160.9,-144.1)
				..controls(160.4,-144.6)and(160.3,-144.8)..(160.1,-145.4)
				..controls(160,-145.8)and(159.6,-146.8)..(159.2,-147.8)
				..controls(158.7,-148.8)and(158.1,-150)..(157.9,-150.6)
				..controls(157.6,-151.1)and(157.2,-151.9)..(156.9,-152.4)
				..controls(156.5,-152.8)and(156.1,-153.7)..(155.8,-154.4)
				..controls(155.6,-155)and(155.2,-155.9)..(154.9,-156.3)
				..controls(154.7,-156.7)and(154.4,-157.4)..(154.2,-157.9)
				..controls(154.1,-158.4)and(153.8,-159)..(153.5,-159.3)
				..controls(152.8,-160.1)and(153,-160.1)..(142.4,-160.1)
				..controls(137.2,-160.1)and(131.8,-160.1)..(130.4,-160)--cycle
				(141.5,-103.7)..controls(140.8,-103)and(139.8,-102.1)..(139.2,-101.6)
				..controls(138.5,-101.2)and(137.9,-100.6)..(137.8,-100.4)
				..controls(137,-99.5)and(134,-96.5)..(133.2,-96)
				..controls(132.7,-95.7)and(131.8,-94.9)..(131.1,-94.3)
				..controls(129.4,-92.7)and(127.9,-91.6)..(126.6,-91.1)
				..controls(125.2,-90.5)and(124.1,-90)..(122.9,-89.1)
				..controls(122.1,-88.6)and(121.3,-88.2)..(119.9,-87.7)
				..controls(118.8,-87.3)and(117.5,-86.8)..(116.9,-86.5)
				..controls(114.3,-85.5)and(113.7,-85.4)..(109.7,-85.5)
				..controls(105.3,-85.7)and(103.9,-85.9)..(101.6,-87)
				..controls(100.7,-87.4)and(99.4,-87.8)..(97.9,-88)
				..controls(97,-88.2)and(96.6,-88.3)..(95.8,-88.9)
				..controls(95.4,-89.2)and(94.7,-89.6)..(94.2,-89.8)
				..controls(92.9,-90.3)and(92.3,-90.7)..(90.9,-92.1)
				..controls(90.2,-92.7)and(89.4,-93.3)..(89,-93.5)..controls(87,-94.5)and(85.2,-96)..(83.6,-97.9)
				..controls(83.3,-98.2)and(82.5,-98.9)..(81.8,-99.4)
				..controls(80.7,-100.1)and(80.5,-100.3)..(80,-101.1)
				..controls(78.7,-103.1)and(77.2,-103.9)..(76.4,-103)
				..controls(75.7,-102.2)and(75.7,-103.2)..(75.7,-83.6)
				..controls(75.6,-73.6)and(75.6,-64.8)..(75.5,-64)
				..controls(75.5,-63.2)and(75.5,-61.4)..(75.6,-59.8)--(75.7,-57.1)--(76,-56.8)
				..controls(76.8,-56.2)and(78,-56.3)..(79.8,-57.2)
				..controls(80.4,-57.5)and(81.3,-57.8)..(81.8,-57.9)
				..controls(82.3,-58.1)and(83,-58.4)..(83.6,-58.8)..controls(84.3,-59.2)and(85,-59.5)..(86,-59.8)
				..controls(86.8,-60.1)and(88,-60.6)..(88.7,-60.9)
				..controls(89.3,-61.2)and(90.6,-61.7)..(91.4,-62)
				..controls(92.2,-62.3)and(93.7,-62.9)..(94.6,-63.3)
				..controls(95.6,-63.7)and(96.8,-64.1)..(97.4,-64.3)
				..controls(98,-64.5)and(99,-64.9)..(99.7,-65.2)
				..controls(100.4,-65.6)and(101.5,-66)..(102.3,-66.2)
				..controls(104.2,-66.8)and(105,-67.1)..(106.2,-67.8)
				..controls(107.9,-68.7)and(110.2,-69.4)..(111.3,-69.2)
				..controls(112.5,-69)and(115,-68.2)..(115.6,-67.8)
				..controls(117.3,-66.7)and(117.6,-66.6)..(119.5,-66.2)
				..controls(121.9,-65.8)and(122.1,-65.7)..(123,-65.1)
				..controls(124.1,-64.4)and(124.4,-64.3)..(126,-63.9)
				..controls(127,-63.7)and(127.7,-63.4)..(128.4,-63.1)
				..controls(129,-62.8)and(129.7,-62.5)..(130.6,-62.4)
				..controls(132.1,-62.1)and(134,-61.5)..(135.2,-61)
				..controls(135.7,-60.8)and(136.7,-60.5)..(137.3,-60.3)
				..controls(138.4,-60)and(139.3,-59.7)..(141.1,-59)
				..controls(141.6,-58.7)and(142,-58.6)..(142.5,-58.6)
				..controls(143.2,-58.6)and(143.3,-58.7)..(143.6,-59)
				..controls(144,-59.5)and(144.1,-59.9)..(144,-60.8)
				..controls(144,-61.2)and(144,-71.1)..(144,-82.8)--(143.9,-104.1)--(143.7,-104.5)
				..controls(143.5,-104.7)and(143.3,-104.8)..(143,-104.9)
				..controls(142.6,-104.9)and(142.5,-104.8)..(141.5,-103.7)--cycle
				(27.8,-91.4)..controls(27.1,-90.9)and(26.3,-90)..(26.1,-89.5)
				..controls(26,-89.2)and(25.8,-88.6)..(25.7,-87.9)
				..controls(25.4,-86.2)and(25.6,-84.8)..(26.5,-83.5)
				..controls(27.4,-82.3)and(28.3,-82)..(29.4,-82.3)
				..controls(31.1,-82.8)and(32.1,-84.9)..(32,-87.4)
				..controls(31.9,-89.3)and(31.2,-90.6)..(29.9,-91.4)
				..controls(29.1,-91.9)and(28.6,-91.9)..(27.8,-91.4)--cycle
				(190.1,-91.3)..controls(189.7,-91.2)and(189.3,-90.9)..(189.1,-90.7)
				..controls(188.3,-89.8)and(187.6,-87.7)..(187.6,-86.2)
				..controls(187.6,-85.5)and(187.7,-85.2)..(188.2,-84.4)
				..controls(189.1,-82.6)and(190.1,-81.8)..(191.3,-82.1)
				..controls(192,-82.3)and(193.5,-83.7)..(193.8,-84.4)
				..controls(194.1,-85.1)and(194.2,-87.1)..(194,-88.3)
				..controls(193.8,-89.8)and(193.3,-90.5)..(192.3,-91.3)
				..controls(191.7,-91.7)and(190.9,-91.7)..(190.1,-91.3)--cycle
				(173.7,-83.4)..controls(172.2,-82.5)and(172,-82.5)..(170.6,-82.4)
				..controls(169.7,-82.4)and(168.8,-82.4)..(168.3,-82.5)
				..controls(167.2,-82.7)and(166.9,-82.7)..(166.6,-82.2)
				..controls(166.3,-81.8)and(166.3,-81)..(166.8,-80)
				..controls(167.2,-78.9)and(167.4,-77.4)..(167.7,-74.1)
				..controls(167.8,-72.8)and(167.9,-71.8)..(167.9,-71.7)
				..controls(168,-71.7)and(168,-70.2)..(168,-68.4)..controls(168,-65.4)and(168,-64.9)..(167.7,-63)
				..controls(167.3,-60.5)and(167.1,-59.9)..(166.4,-58.4)
				..controls(166.1,-57.7)and(165.8,-56.9)..(165.7,-56.6)
				..controls(165.7,-56.2)and(165.4,-55.6)..(165.2,-55.2)
				..controls(164.6,-53.8)and(161.9,-50.7)..(161,-50.3)
				..controls(159.8,-49.7)and(157.9,-50)..(155.2,-51.2)
				..controls(154.7,-51.4)and(153.4,-51.8)..(152.2,-52.1)
				..controls(150.8,-52.4)and(149.4,-52.9)..(148.6,-53.2)
				..controls(147.8,-53.6)and(146.7,-54)..(146,-54.2)
				..controls(145.4,-54.4)and(144.5,-54.7)..(144,-54.9)
				..controls(143.6,-55.2)and(142.7,-55.5)..(141.8,-55.7)
				..controls(141,-55.9)and(139.6,-56.3)..(138.7,-56.6)
				..controls(137.8,-57)and(136.3,-57.4)..(135.3,-57.6)
				..controls(133.8,-57.9)and(133.2,-58.1)..(132.1,-58.7)
				..controls(131.1,-59.2)and(130.3,-59.4)..(129.4,-59.6)
				..controls(128.7,-59.8)and(127.8,-60)..(127.4,-60.1)
				..controls(127,-60.3)and(126,-60.5)..(125.1,-60.7)
				..controls(123.2,-61.1)and(121.7,-61.5)..(120.2,-62.1)
				..controls(119.7,-62.3)and(118.1,-62.8)..(116.7,-63.2)
				..controls(114.6,-63.9)and(114,-64.1)..(113,-64.7)
				..controls(111.9,-65.4)and(111.8,-65.4)..(111,-65.5)
				..controls(109.9,-65.5)and(109.3,-65.4)..(107.7,-64.7)
				..controls(107,-64.4)and(106,-64)..(105.3,-63.9)
				..controls(103.5,-63.5)and(102.2,-63)..(101.2,-62.5)
				..controls(100.1,-61.8)and(99.3,-61.5)..(97.8,-61.1)
				..controls(97.2,-61)and(96.2,-60.7)..(95.7,-60.4)
				..controls(95.1,-60.1)and(94.2,-59.8)..(93.5,-59.6)
				..controls(92.9,-59.4)and(91.7,-59)..(91,-58.7)..controls(90.2,-58.4)and(88.9,-58)..(88,-57.7)
				..controls(87.2,-57.5)and(86.1,-57.1)..(85.5,-56.8)
				..controls(85,-56.6)and(84.1,-56.2)..(83.5,-56)..controls(82.9,-55.7)and(81.8,-55.2)..(81,-54.9)
				..controls(80.1,-54.5)and(79.1,-54.2)..(78.7,-54.1)
				..controls(77.9,-53.8)and(77.3,-53.6)..(76.2,-52.9)
				..controls(75.7,-52.6)and(75.2,-52.4)..(74.1,-52.2)
				..controls(72.5,-51.9)and(72.4,-51.8)..(71.3,-51.1)
				..controls(69.9,-50.1)and(69.4,-49.9)..(67.9,-49.8)
				..controls(67,-49.7)and(66.6,-49.6)..(66.3,-49.4)
				..controls(66.1,-49.3)and(65.7,-49)..(65.4,-48.8)
				..controls(65.1,-48.7)and(64.2,-48.4)..(63.6,-48.2)
				..controls(62.9,-48)and(61.7,-47.6)..(60.9,-47.2)
				..controls(60.1,-46.8)and(58.9,-46.4)..(58.2,-46.2)
				..controls(57.6,-46)and(56.8,-45.7)..(56.5,-45.5)
				..controls(54.9,-44.6)and(54.6,-44.5)..(53.2,-44.2)
				..controls(51.7,-43.8)and(51,-43.6)..(49.9,-42.9)
				..controls(49.5,-42.6)and(48.9,-42.4)..(47.5,-42)
				..controls(45.7,-41.6)and(45.6,-41.5)..(44.3,-40.6)
				..controls(43.5,-40.1)and(42.8,-39.6)..(42.7,-39.4)
				..controls(42.4,-38.7)and(43,-38.1)..(44.5,-37.6)..controls(44.9,-37.5)and(46,-37)..(46.9,-36.6)
				..controls(47.9,-36.1)and(48.8,-35.8)..(49.6,-35.6)
				..controls(50.5,-35.4)and(51.3,-35.2)..(52.4,-34.6)
				..controls(53.3,-34.1)and(54.5,-33.7)..(54.9,-33.5)
				..controls(55.4,-33.4)and(56.4,-33)..(57.2,-32.6)
				..controls(58.6,-32)and(60.1,-31.5)..(61.6,-31.2)
				..controls(62.3,-31.1)and(62.7,-31)..(63.2,-30.7)
				..controls(64.1,-30)and(64.3,-29.9)..(65.3,-29.7)
				..controls(65.7,-29.6)and(66.5,-29.2)..(67.1,-28.9)
				..controls(67.6,-28.6)and(68.8,-28.2)..(69.6,-27.9)
				..controls(70.5,-27.7)and(71.6,-27.3)..(72.1,-27)
				..controls(72.6,-26.7)and(73.8,-26.2)..(74.8,-25.9)
				..controls(75.7,-25.5)and(76.7,-25.1)..(77.1,-25)
				..controls(77.4,-24.8)and(78.5,-24.4)..(79.6,-24.1)
				..controls(80.6,-23.8)and(82,-23.2)..(82.8,-22.9)
				..controls(83.5,-22.5)and(84.8,-22)..(85.8,-21.8)
				..controls(86.7,-21.5)and(87.8,-21.1)..(88.1,-20.9)
				..controls(88.4,-20.7)and(89.3,-20.3)..(90.1,-20.1)
				..controls(90.8,-19.8)and(91.9,-19.4)..(92.4,-19.1)
				..controls(93.5,-18.5)and(94.4,-18.2)..(95.4,-18.1)
				..controls(95.7,-18)and(96.4,-17.8)..(96.9,-17.5)
				..controls(97.5,-17.3)and(98.4,-16.9)..(99,-16.6)
				..controls(99.7,-16.3)and(100.9,-15.9)..(101.9,-15.6)
				..controls(102.9,-15.3)and(104.4,-14.7)..(105.3,-14.3)
				..controls(108.1,-13.2)and(108.5,-13.1)..(110.1,-13.1)
				..controls(111.4,-13.1)and(111.5,-13.1)..(112.4,-13.5)
				..controls(113.2,-14)and(115,-14.8)..(116.6,-15.4)
				..controls(116.9,-15.5)and(117.9,-15.8)..(118.8,-16.1)
				..controls(120.9,-16.6)and(121.5,-16.9)..(122.6,-17.9)
				..controls(123,-18.3)and(123.7,-18.8)..(123.9,-18.9)
				..controls(124.6,-19.4)and(126,-19.8)..(127.9,-20.2)
				..controls(129.6,-20.5)and(129.8,-20.6)..(131.1,-21.4)
				..controls(131.5,-21.7)and(132.4,-22.1)..(133.1,-22.3)
				..controls(136.1,-23.2)and(138.3,-24.1)..(140.2,-25)
				..controls(141,-25.5)and(141.5,-25.6)..(142.5,-25.8)
				..controls(145.1,-26.3)and(146,-26.6)..(149.4,-28.1)
				..controls(151.3,-28.9)and(153.5,-29.8)..(154.3,-30.1)
				..controls(155.1,-30.4)and(156.2,-30.8)..(156.6,-31)
				..controls(157.1,-31.2)and(158.2,-31.6)..(159.1,-31.8)
				..controls(160.6,-32.2)and(160.9,-32.3)..(161.7,-32.9)
				..controls(162.4,-33.3)and(162.9,-33.6)..(163.7,-33.8)
				..controls(164.3,-34)and(164.8,-34.1)..(164.9,-34.1)
				..controls(165,-34.1)and(165.6,-34.4)..(166.4,-34.8)
				..controls(167.3,-35.2)and(168.1,-35.4)..(169,-35.6)
				..controls(170,-35.8)and(170.5,-36)..(171.7,-36.6)
				..controls(172.4,-37)and(173.3,-37.4)..(173.6,-37.4)
				..controls(175.6,-37.7)and(176.2,-38)..(176.7,-39.1)
				..controls(176.9,-39.6)and(176.9,-39.7)..(176.7,-40.4)
				..controls(176.4,-41.6)and(175.8,-42.7)..(175.2,-43.4)
				..controls(174.5,-44.1)and(172.6,-45.5)..(171.7,-45.8)
				..controls(170.6,-46.2)and(169.1,-46)..(168.2,-45.4)
				..controls(167.9,-45.2)and(167.3,-44.8)..(166.9,-44.4)
				..controls(166.4,-44.1)and(165.8,-43.6)..(165.4,-43.2)
				..controls(165.1,-42.9)and(164.4,-42.3)..(163.9,-41.9)
				..controls(163.3,-41.5)and(162.7,-40.9)..(162.4,-40.5)
				..controls(161.4,-39.2)and(159.9,-38)..(158.8,-37.8)
				..controls(158.4,-37.8)and(158.3,-37.7)..(158.2,-37.2)
				..controls(158,-36.8)and(157.4,-36.1)..(157.1,-36.1)
				..controls(157,-36.1)and(156.7,-36.1)..(156.5,-36)
				..controls(156.1,-35.9)and(155.9,-35.7)..(155.7,-35.2)
				..controls(155.3,-34.6)and(155.1,-34.4)..(154,-34.1)
				..controls(153.6,-34)and(153.2,-33.7)..(152.9,-33.5)
				..controls(152.4,-32.9)and(151.6,-32.4)..(151.3,-32.4)
				..controls(150.9,-32.4)and(150.4,-32.6)..(150.4,-32.9)
				..controls(150.4,-33.2)and(150.8,-33.4)..(151.4,-33.6)
				..controls(152.1,-33.7)and(152.8,-34.1)..(153.5,-34.9)
				..controls(154.1,-35.5)and(154.3,-35.6)..(155,-35.8)
				..controls(155.7,-36)and(155.7,-36)..(156,-36.6)
				..controls(156.3,-37.3)and(156.7,-37.7)..(157.5,-37.8)
				..controls(158.1,-37.9)and(158.2,-38)..(158.4,-38.8)
				..controls(158.5,-39.4)and(158.6,-39.6)..(159.6,-40.4)
				..controls(160.2,-41)and(160.9,-41.6)..(161.3,-41.8)
				..controls(161.7,-42)and(161.9,-42.3)..(162.2,-42.8)
				..controls(162.4,-43.2)and(162.8,-43.8)..(163.1,-44.1)
				..controls(163.9,-44.8)and(165.1,-46.7)..(165.7,-48)
				..controls(166.2,-49.3)and(166.4,-49.7)..(167.2,-50.5)
				..controls(167.5,-50.8)and(167.9,-51.4)..(168.1,-51.8)
				..controls(168.3,-52.2)and(168.8,-53)..(169.1,-53.6)
				..controls(169.6,-54.3)and(169.8,-54.9)..(169.9,-55.6)
				..controls(170.4,-57.8)and(170.5,-58.1)..(171.2,-59.5)
				..controls(171.6,-60.3)and(172,-61.5)..(172.2,-62.1)
				..controls(172.3,-62.8)and(172.7,-63.7)..(172.9,-64.1)
				..controls(173.2,-64.6)and(173.4,-65.3)..(173.4,-65.6)
				..controls(173.8,-68.8)and(174.2,-70.6)..(174.9,-72.6)
				..controls(175.5,-74.5)and(175.9,-77.1)..(176.1,-80.2)
				..controls(176.3,-83.4)and(175.5,-84.4)..(173.7,-83.4)--cycle;
				\definecolor{ezs}{RGB}{247,179,26}\fill[ezs](-5.7,-156.6)--(-3.4,-157.5)--(61.9,12.7)--(59.7,13.6)--cycle;
				\definecolor{ezs}{RGB}{247,179,22}\fill[ezs](-3.8,-157.3)--(-1.5,-158.2)--(63.9,12)--(61.6,12.9)--cycle;
				\definecolor{ezs}{RGB}{248,179,19}\fill[ezs](-1.8,-158.1)--(0.4,-158.9)--(65.8,11.2)--(63.5,12.1)--cycle;
				\definecolor{ezs}{RGB}{248,178,16}\fill[ezs](0.1,-158.8)--(2.4,-159.7)--(67.7,10.5)--(65.4,11.4)--cycle;
				\definecolor{ezs}{RGB}{249,178,15}\fill[ezs](2,-159.5)--(4.3,-160.4)--(69.6,9.8)--(67.3,10.7)--cycle;
				\definecolor{ezs}{RGB}{251,177,17}\fill[ezs](3.9,-160.3)--(6.2,-161.1)--(71.5,9)--(69.2,9.9)--cycle;
				\definecolor{ezs}{RGB}{252,179,18}\fill[ezs](5.8,-161)--(8.1,-161.9)--(73.4,8.3)--(71.1,9.2)--cycle;
				\definecolor{ezs}{RGB}{252,182,17}\fill[ezs](7.7,-161.7)--(10,-162.6)--(75.3,7.6)--(73,8.5)--cycle;
				\definecolor{ezs}{RGB}{252,185,17}\fill[ezs](9.6,-162.5)--(11.9,-163.4)--(77.2,6.8)--(75,7.7)--cycle;
				\definecolor{ezs}{RGB}{251,189,18}\fill[ezs](11.5,-163.2)--(13.8,-164.1)--(79.2,6.1)--(76.9,7)--cycle;
				\definecolor{ezs}{RGB}{251,191,22}\fill[ezs](13.4,-163.9)--(15.7,-164.8)--(81.1,5.4)--(78.8,6.2)--cycle;
				\definecolor{ezs}{RGB}{252,193,25}\fill[ezs](15.4,-164.7)--(17.6,-165.6)--(83,4.6)--(80.7,5.5)--cycle;
				\definecolor{ezs}{RGB}{253,195,27}\fill[ezs](17.3,-165.4)--(19.6,-166.3)--(84.9,3.9)--(82.6,4.8)--cycle;
				\definecolor{ezs}{RGB}{253,196,31}\fill[ezs](19.2,-166.1)--(21.5,-167)--(86.8,3.2)--(84.5,4)--cycle;
				\definecolor{ezs}{RGB}{254,199,35}\fill[ezs](21.1,-166.9)--(23.4,-167.8)--(88.7,2.4)--(86.4,3.3)--cycle;
				\definecolor{ezs}{RGB}{253,201,38}\fill[ezs](23,-167.6)--(25.3,-168.5)--(90.6,1.7)--(88.3,2.6)--cycle;
				\definecolor{ezs}{RGB}{253,203,40}\fill[ezs](24.9,-168.3)--(27.2,-169.2)--(92.5,1)--(90.3,1.8)--cycle;
				\definecolor{ezs}{RGB}{253,205,42}\fill[ezs](26.8,-169.1)--(29.1,-170)--(94.4,0.2)--(92.2,1.1)--cycle;
				\definecolor{ezs}{RGB}{253,207,45}\fill[ezs](28.7,-169.8)--(31,-170.7)--(96.4,-0.5)--(94.1,0.4)--cycle;
				\definecolor{ezs}{RGB}{253,210,49}\fill[ezs](30.7,-170.5)--(32.9,-171.4)--(98.3,-1.2)--(96,-0.4)--cycle;
				\definecolor{ezs}{RGB}{253,212,54}\fill[ezs](32.6,-171.3)--(34.9,-172.2)--(100.2,-2)--(97.9,-1.1)--cycle;
				\definecolor{ezs}{RGB}{253,214,59}\fill[ezs](34.5,-172)--(36.8,-172.9)--(102.1,-2.7)--(99.8,-1.8)--cycle;
				\definecolor{ezs}{RGB}{253,216,63}\fill[ezs](36.4,-172.7)--(38.7,-173.6)--(104,-3.4)--(101.7,-2.6)--cycle;
				\definecolor{ezs}{RGB}{253,218,65}\fill[ezs](38.3,-173.5)--(40.6,-174.4)--(105.9,-4.2)--(103.6,-3.3)--cycle;
				\definecolor{ezs}{RGB}{254,220,66}\fill[ezs](40.2,-174.2)--(42.5,-175.1)--(107.8,-4.9)--(105.5,-4)--cycle;
				\definecolor{ezs}{RGB}{254,222,68}\fill[ezs](42.1,-175)--(44.4,-175.8)--(109.7,-5.6)--(107.5,-4.8)--cycle;
				\definecolor{ezs}{RGB}{254,224,71}\fill[ezs](44,-175.7)--(46.3,-176.6)--(111.7,-6.4)--(109.4,-5.5)--cycle;
				\definecolor{ezs}{RGB}{254,226,75}\fill[ezs](46,-176.4)--(48.2,-177.3)--(113.6,-7.1)--(111.3,-6.2)--cycle;
				\definecolor{ezs}{RGB}{254,228,79}\fill[ezs](47.9,-177.2)--(50.1,-178)--(115.5,-7.8)--(113.2,-7)--cycle;
				\definecolor{ezs}{RGB}{254,230,82}\fill[ezs](49.8,-177.9)--(52.1,-178.8)--(117.4,-8.6)--(115.1,-7.7)--cycle;
				\definecolor{ezs}{RGB}{254,232,85}\fill[ezs](51.7,-178.6)--(54,-179.5)--(119.3,-9.3)--(117,-8.4)--cycle;
				\definecolor{ezs}{RGB}{254,234,88}\fill[ezs](53.6,-179.4)--(55.9,-180.2)--(121.2,-10)--(118.9,-9.2)--cycle;
				\definecolor{ezs}{RGB}{254,237,91}\fill[ezs](55.5,-180.1)--(57.8,-181)--(123.1,-10.8)--(120.8,-9.9)--cycle;
				\definecolor{ezs}{RGB}{254,239,94}\fill[ezs](57.4,-180.8)--(59.7,-181.7)--(125,-11.5)--(122.8,-10.6)--cycle;
				\definecolor{ezs}{RGB}{254,241,97}\fill[ezs](59.3,-181.6)--(61.6,-182.4)--(127,-12.2)--(124.7,-11.4)--cycle;
				\definecolor{ezs}{RGB}{254,243,101}\fill[ezs](61.2,-182.3)--(63.5,-183.2)--(128.9,-13)--(126.6,-12.1)--cycle;
				\definecolor{ezs}{RGB}{254,245,103}\fill[ezs](63.2,-183)--(65.4,-183.9)--(130.8,-13.7)--(128.5,-12.8)--cycle;
				\definecolor{ezs}{RGB}{254,245,105}\fill[ezs](65.1,-183.8)--(67.4,-184.6)--(132.7,-14.4)--(130.4,-13.6)--cycle;
				\definecolor{ezs}{RGB}{254,245,105}\fill[ezs](67,-184.5)--(69.3,-185.4)--(134.6,-15.2)--(132.3,-14.3)--cycle;
				\definecolor{ezs}{RGB}{254,244,103}\fill[ezs](68.9,-185.2)--(71.2,-186.1)--(136.5,-15.9)--(134.2,-15)--cycle;
				\definecolor{ezs}{RGB}{253,243,101}\fill[ezs](70.8,-186)--(73.1,-186.8)--(138.4,-16.6)--(136.1,-15.8)--cycle;
				\definecolor{ezs}{RGB}{253,241,98}\fill[ezs](72.7,-186.7)--(75,-187.6)--(140.3,-17.4)--(138,-16.5)--cycle;
				\definecolor{ezs}{RGB}{253,239,95}\fill[ezs](74.6,-187.4)--(76.9,-188.3)--(142.2,-18.1)--(140,-17.2)--cycle;
				\definecolor{ezs}{RGB}{253,236,92}\fill[ezs](76.5,-188.2)--(78.8,-189)--(144.2,-18.8)--(141.9,-18)--cycle;
				\definecolor{ezs}{RGB}{253,234,88}\fill[ezs](78.5,-188.9)--(80.7,-189.8)--(146.1,-19.6)--(143.8,-18.7)--cycle;
				\definecolor{ezs}{RGB}{253,232,85}\fill[ezs](80.4,-189.6)--(82.7,-190.5)--(148,-20.3)--(145.7,-19.4)--cycle;
				\definecolor{ezs}{RGB}{253,230,82}\fill[ezs](82.3,-190.4)--(84.6,-191.2)--(149.9,-21)--(147.6,-20.2)--cycle;
				\definecolor{ezs}{RGB}{253,228,78}\fill[ezs](84.2,-191.1)--(86.5,-192)--(151.8,-21.8)--(149.5,-20.9)--cycle;
				\definecolor{ezs}{RGB}{253,226,74}\fill[ezs](86.1,-191.8)--(88.4,-192.7)--(153.7,-22.5)--(151.4,-21.6)--cycle;
				\definecolor{ezs}{RGB}{253,224,70}\fill[ezs](88,-192.6)--(90.3,-193.4)--(155.6,-23.3)--(153.3,-22.4)--cycle;
				\definecolor{ezs}{RGB}{253,222,67}\fill[ezs](89.9,-193.3)--(92.2,-194.2)--(157.5,-24)--(155.3,-23.1)--cycle;
				\definecolor{ezs}{RGB}{253,220,65}\fill[ezs](91.8,-194)--(94.1,-194.9)--(159.5,-24.7)--(157.2,-23.8)--cycle;
				\definecolor{ezs}{RGB}{253,217,63}\fill[ezs](93.7,-194.8)--(96,-195.6)--(161.4,-25.5)--(159.1,-24.6)--cycle;
				\definecolor{ezs}{RGB}{253,215,60}\fill[ezs](95.7,-195.5)--(97.9,-196.4)--(163.3,-26.2)--(161,-25.3)--cycle;
				\definecolor{ezs}{RGB}{253,213,56}\fill[ezs](97.6,-196.2)--(99.9,-197.1)--(165.2,-26.9)--(162.9,-26)--cycle;
				\definecolor{ezs}{RGB}{252,211,53}\fill[ezs](99.5,-197)--(101.8,-197.8)--(167.1,-27.7)--(164.8,-26.8)--cycle;
				\definecolor{ezs}{RGB}{252,209,49}\fill[ezs](101.4,-197.7)--(103.7,-198.6)--(169,-28.4)--(166.7,-27.5)--cycle;
				\definecolor{ezs}{RGB}{252,207,45}\fill[ezs](103.3,-198.4)--(105.6,-199.3)--(170.9,-29.1)--(168.6,-28.2)--cycle;
				\definecolor{ezs}{RGB}{253,205,42}\fill[ezs](105.2,-199.2)--(107.5,-200)--(172.8,-29.9)--(170.6,-29)--cycle;
				\definecolor{ezs}{RGB}{253,202,40}\fill[ezs](107.1,-199.9)--(109.4,-200.8)--(174.7,-30.6)--(172.5,-29.7)--cycle;
				\definecolor{ezs}{RGB}{253,200,38}\fill[ezs](109,-200.6)--(111.3,-201.5)--(176.7,-31.3)--(174.4,-30.4)--cycle;
				\definecolor{ezs}{RGB}{253,198,35}\fill[ezs](111,-201.4)--(113.2,-202.2)--(178.6,-32.1)--(176.3,-31.2)--cycle;
				\definecolor{ezs}{RGB}{253,196,31}\fill[ezs](112.9,-202.1)--(115.2,-203)--(180.5,-32.8)--(178.2,-31.9)--cycle;
				\definecolor{ezs}{RGB}{253,194,28}\fill[ezs](114.8,-202.8)--(117.1,-203.7)--(182.4,-33.5)--(180.1,-32.6)--cycle;
				\definecolor{ezs}{RGB}{253,192,26}\fill[ezs](116.7,-203.6)--(119,-204.5)--(184.3,-34.3)--(182,-33.4)--cycle;
				\definecolor{ezs}{RGB}{253,189,24}\fill[ezs](118.6,-204.3)--(120.9,-205.2)--(186.2,-35)--(183.9,-34.1)--cycle;
				\definecolor{ezs}{RGB}{253,186,20}\fill[ezs](120.5,-205)--(122.8,-205.9)--(188.1,-35.7)--(185.8,-34.8)--cycle;
				\definecolor{ezs}{RGB}{253,184,17}\fill[ezs](122.4,-205.8)--(124.7,-206.7)--(190,-36.5)--(187.8,-35.6)--cycle;
				\definecolor{ezs}{RGB}{252,183,15}\fill[ezs](124.3,-206.5)--(126.6,-207.4)--(192,-37.2)--(189.7,-36.3)--cycle;
				\definecolor{ezs}{RGB}{252,181,14}\fill[ezs](126.3,-207.2)--(128.5,-208.1)--(193.9,-37.9)--(191.6,-37.1)--cycle;
				\definecolor{ezs}{RGB}{252,179,13}\fill[ezs](128.2,-208)--(130.4,-208.9)--(195.8,-38.7)--(193.5,-37.8)--cycle;
				\definecolor{ezs}{RGB}{252,178,12}\fill[ezs](130.1,-208.7)--(132.4,-209.6)--(197.7,-39.4)--(195.4,-38.5)--cycle;
				\definecolor{ezs}{RGB}{251,178,10}\fill[ezs](132,-209.4)--(134.3,-210.3)--(199.6,-40.1)--(197.3,-39.3)--cycle;
				\definecolor{ezs}{RGB}{251,178,10}\fill[ezs](133.9,-210.2)--(136.2,-211.1)--(201.5,-40.9)--(199.2,-40)--cycle;
				\definecolor{ezs}{RGB}{250,177,11}\fill[ezs](135.8,-210.9)--(138.1,-211.8)--(203.4,-41.6)--(201.1,-40.7)--cycle;
				\definecolor{ezs}{RGB}{250,178,11}\fill[ezs](137.7,-211.6)--(140,-212.5)--(205.3,-42.3)--(203.1,-41.5)--cycle;
				\definecolor{ezs}{RGB}{248,180,10}\fill[ezs](139.6,-212.4)--(141.9,-213.3)--(207.3,-43.1)--(205,-42.2)--cycle;
				\definecolor{ezs}{RGB}{247,181,9}\fill[ezs](141.5,-213.1)--(143.8,-214)--(209.2,-43.8)--(206.9,-42.9)--cycle;
				\definecolor{ezs}{RGB}{246,182,9}\fill[ezs](143.5,-213.8)--(145.7,-214.7)--(211.1,-44.5)--(208.8,-43.7)--cycle;
				\definecolor{ezs}{RGB}{246,182,10}\fill[ezs](145.4,-214.6)--(147.7,-215.5)--(213,-45.3)--(210.7,-44.4)--cycle;
			\end{scope}
			% chữ EDUZY (Times đậm, làm dày nét bằng 5 lớp lệch nhẹ; căn đúng khung chữ logo gốc)
			\foreach \dx/\dy in {-0.7/0,0.7/0,0/-0.3,0/0.3,0/0}
			\node[anchor=center,inner sep=0pt,outer sep=0pt,text=ezText,align=center] at (109.6+\dx,-193.6+\dy)
			{\resizebox{1.01cm}{0.22cm}{\fontencoding{T1}\fontfamily{ptm}\fontseries{b}\selectfont EDUZY}};
		\end{tikzpicture}%
	}%
}
%====================== THIẾT LẬP THÂN ĐỀ ======================
%--- Môi trường tự luận {bt}: ex_test 3.3.4 chưa khai báo sẵn nên tạo ở đây
%    (ex_test đã tự móc \fixshowans và \listenumerate cho bt)
\makeatletter
\@ifundefined{bt}{\newtheorem{bt}{\color{gold1}\namebt}}{}
\makeatother
\renewcommand{\sectionmark}[1]{\markright{\thesection\ #1}}
%--- Khung tiêu đề (logo + tối đa 4 dòng; dòng 3, 4 để trống {} thì bỏ qua)
%  Đề thi:    \khungtieude{THI THỬ THPT QUỐC GIA NĂM 2026}{Môn Toán -- Đề số 21}
%                         {(Đề thi gồm: ...)}{Thời gian làm bài: 90 phút (không kể thời gian phát đề).}
%  Chuyên đề: \khungtieude{CHUYÊN ĐỀ 3 -- DÃY SỐ VÀ CÔNG THỨC TRUY HỒI}
%                         {Bồi dưỡng học sinh giỏi Toán THPT}{...}{}
\newcommand{\khungtieude}[4]{%
	\begin{tcolorbox}[colframe=gold1,colback=gold5!45,boxrule=0.5mm]
		\begin{minipage}[c]{3cm}
			\centerline{\eduzylogo[2.4cm]}
		\end{minipage}\hspace{0.5cm}
		\begin{minipage}[c]{12.5cm}
			\vskip 0.15 cm
			\centerline{\ind{#1}}
			\vskip 0.15 cm
			\centerline{\textbf{#2}}
			\if\relax\detokenize{#3}\relax\else
				\vskip 0.15 cm
				\centerline{\textit{#3}}
			\fi
			\if\relax\detokenize{#4}\relax\else
				\vskip 0.15 cm
				\centerline{\textit{#4}}
			\fi
			\vskip 0.15 cm
		\end{minipage}
	\end{tcolorbox}
}
%====================== BỔ SUNG CHO BÀI HỌC, CHUYÊN ĐỀ ======================
\usepackage[shortlabels]{enumitem}
\usepackage{tabularx}
%--- Danh sách gọn: \begin{enumerate}[gon] ... \end{enumerate}
\SetEnumitemKey{gon}{leftmargin=*,itemsep=1pt,parsep=0pt,topsep=3pt}
%--- Khung lý thuyết: \begin{khung}{} ... \end{khung}; tiêu đề chỉ dùng khi thật ngắn ("Mẹo", "Lỗi thường gặp")
\newtcolorbox{khung}[1]{breakable,enhanced,colframe=gold1,colback=gold5!45,
	boxrule=0.4mm,fonttitle=\bfseries\sffamily,coltitle=white,colbacktitle=gold1,
	attach boxed title to top left={xshift=6mm,yshift=-3mm},
	boxed title style={arc=1mm,boxrule=0.3mm},title={#1}}
%--- Tiêu đề mục nhỏ trong một dạng: \muc{Dấu hiệu}, \muc{Phương pháp}, \muc{Bài tự luyện}
\makeatletter
\newcommand{\muc}[1]{\par\medskip\noindent\indam{#1}\par\nobreak\smallskip\@afterindentfalse\@afterheading}
\makeatother
%--- Ô công thức cần nhớ (căn giữa, viền vàng, vừa khít nội dung)
%     \begin{center}\congthuc{$Q(n)=(n-1)!$}\end{center}
\newtcbox{\congthuc}{on line,colframe=gold1,colback=gold5!45,boxrule=0.5pt,arc=2pt,
	left=8pt,right=8pt,top=4pt,bottom=4pt}
%--- Công tắc lời giải bài tự luyện (ví dụ mẫu luôn dùng \loigiai)
%     \loigiaiTL[{đáp số}]{lời giải đầy đủ}
%     \TLfull -> in lời giải đầy đủ;  \TLdapso -> chỉ in một dòng "Đáp số";  \TLan -> ẩn hoàn toàn
%     (\CheDo là dethi thì mọi lời giải và đáp số đều ẩn)
\makeatletter
\newif\ifTLhien   \TLhientrue
\newif\ifTLdapso  \TLdapsofalse
\newcommand{\TLfull}{\TLhientrue\TLdapsofalse}
\newcommand{\TLdapso}{\TLhientrue\TLdapsotrue}
\newcommand{\TLan}{\TLhienfalse}
\newcommand{\loigiaiTL}[2][]{%
	\ifTLhien
		\ifTLdapso
			\def\@tlarg{#1}%
			\ifx\@tlarg\@empty\else
				\ifdefstring{\kindOPTN}{dethi}{}{\par\noindent{\bfseries\color{gold1}Đáp số:} #1\par}%
			\fi
		\else
			\loigiai{#2}%
		\fi
	\fi}
\makeatother
\TLdapso% mặc định: chỉ in đáp số bài tự luyện
%====================== BẢNG ĐÚNG/SAI (\choiceTFt) ======================
%  Hàng tiêu đề nền nâu đồng chữ trắng (Phát biểu | Đ | S), các hàng tô xen kẽ kem / trắng,
%  đường kẻ vàng kim, nhãn a) b) c) d) màu nâu đồng. Bản có đáp án: cột được chọn có vòng tròn
%  nâu đồng với dấu tích trắng (câu đúng tích ở cột Đúng, câu sai tích ở cột Sai).
%  Tiêu đề ghi đủ "Đúng | Sai": \OPTN{viettat=0} sau \input (cột Đúng/Sai rộng ra, phát biểu hẹp lại).
%  Dùng bảng gốc của ex_test: \def\BangDSgoc{1} trước \input{preamble-eduzy-v2}.
\providecommand{\BangDSgoc}{0}
\ifnum\BangDSgoc=0
\usepackage{colortbl,pifont}
\colorlet{ezDSke}{gold2!60!gold1}% đường kẻ bảng (nét 0,4pt như ex_test để không đổi chiều cao bảng)
\colorlet{ezDSle}{gold5!60}% hàng a), c)
\colorlet{ezDSchan}{white}% hàng b), d)
%--- dấu chọn ở cột Đúng/Sai
%    (bề rộng 0, căn giữa ô: không làm đổi bề rộng cột của ex_test, không báo tràn lề)
\newcommand{\EZdauchon}{\makebox[0pt]{\tikz[baseline=-1ex]\node[circle,fill=gold1,inner sep=0pt,
	minimum size=1.2em,text=white,font=\small]{\ding{51}};}}
\def\tickT{\EZdauchon}\def\tickF{\EZdauchon}
\AtBeginDocument{\sbox{\ccTFTab}{\EZdauchon}}
%--- nhãn a) b) c) d) màu nâu đồng (ở bản lời giải, nhãn câu đúng vẫn đóng khung như ex_test)
\renewcommand{\FalseTF}{\stepcounter{dapan}\textbf{\color{gold1}\DapAnTF\sepTF}\ignorespaces}
%--- trước khi dựng bảng: kẻ màu vàng kim (bề rộng cột giữ nguyên như ex_test); sau bảng trả lại màu đen
\newcount\EZhangTF
\newcommand{\EZbatdauTF}{%
	\global\EZhangTF=0\relax
	\arrayrulecolor{ezDSke}}
\newcommand{\EZketthucTF}{\arrayrulecolor{black}}
%--- hàng tiêu đề nền nâu đồng, chữ trắng
\makeatletter
\newcommand{\EZo@TF}[1]{\centering\bfseries\color{white}#1}
\newcommand{\EZtieudeTF}{%
	\ifcase\cotDS
		\ifnum\cotSTT=1\relax
			\def\rowHeader{\hline\rowcolor{gold1}\EZo@TF{TT}&\EZo@TF{\phatbieu}\tabularnewline}%
		\else
			\def\rowHeader{\hline\rowcolor{gold1}\EZo@TF{\phatbieu}\tabularnewline}%
		\fi
	\or
		\ifnum\cotSTT=1\relax
			\def\rowHeader{\hline\rowcolor{gold1}\EZo@TF{TT}&\EZo@TF{\phatbieu}&\EZo@TF{\Dung}\tabularnewline}%
		\else
			\def\rowHeader{\hline\rowcolor{gold1}\EZo@TF{\phatbieu}&\EZo@TF{\Dung}\tabularnewline}%
		\fi
	\or
		\ifnum\cotSTT=1\relax
			\def\rowHeader{\hline\rowcolor{gold1}\EZo@TF{TT}&\EZo@TF{\phatbieu}&\EZo@TF{\Dung}&\EZo@TF{\Sai}\tabularnewline}%
		\else
			\def\rowHeader{\hline\rowcolor{gold1}\EZo@TF{\phatbieu}&\EZo@TF{\Dung}&\EZo@TF{\Sai}\tabularnewline}%
		\fi
	\fi}
%--- tô xen kẽ từng hàng (gọi ngay sau \hline đầu mỗi hàng)
\newcommand{\EZmauhangTF}{%
	\noalign{\global\advance\EZhangTF\@ne}%
	\ifodd\EZhangTF\rowcolor{ezDSle}\else\rowcolor{ezDSchan}\fi}
%--- gắn vào \choiceTFt / \writeChoiceTFt của ex_test
\xpretocmd{\choiceTFt}{\EZbatdauTF}{}{\typeout{[preamble-eduzy-v2] Không gắn được bảng Đúng/Sai (đầu bảng)}}
\xpatchcmd{\choiceTFt}{\ifnum\dongPB=0\def\rowHeader{}\fi}{\EZtieudeTF\ifnum\dongPB=0\def\rowHeader{}\fi}
	{}{\typeout{[preamble-eduzy-v2] Không gắn được bảng Đúng/Sai (hàng tiêu đề)}}
\xpatchcmd{\writeChoiceTFt}{\hline\xdef\dotEX{\sdotEX}}{\hline\EZmauhangTF\xdef\dotEX{\sdotEX}}
	{}{\typeout{[preamble-eduzy-v2] Không gắn được bảng Đúng/Sai (tô hàng)}}
\xpatchcmd{\writeChoiceTFt}{{\end{longtable}}}{{\end{longtable}\EZketthucTF}}
	{}{\typeout{[preamble-eduzy-v2] Không gắn được bảng Đúng/Sai (cuối longtable)}}
\xpatchcmd{\writeChoiceTFt}{{\end{xtabular}\end{center}\medskip\noindent}}{{\end{xtabular}\end{center}\EZketthucTF\medskip\noindent}}
	{}{\typeout{[preamble-eduzy-v2] Không gắn được bảng Đúng/Sai (cuối xtabular)}}
\fi
%=====================================================================
\renewcommand{\parallel}{\mathrel{/\!\!/}}

%      \TLdapso                      % tài liệu: \TLfull | \TLdapso | \TLan
%      \begin{document} ... \end{document}
%  Logo Eduzy vẽ bằng TikZ ngay trong file (\eduzylogo), không cần file ảnh.
%  Biên dịch: pdfLaTeX (máy có gói vietnam) hoặc XeLaTeX.
%============================================================
\documentclass[12pt,a4paper,twoside]{article}%<-- twoside: bat dau trang/chan trang chan-le
\usepackage{iftex}
\usepackage{amsmath,amssymb,tikz,mathrsfs,tkz-tab}
\usepackage{graphicx}
\usepackage{pgfplots}% đồ thị hàm số
\pgfplotsset{compat=1.17}
%\usepackage{mathptmx}
\usepackage{mathpazo}
\ifXeTeX% XeLaTeX: font chữ có đủ dấu tiếng Việt, giữ nguyên font toán mathpazo
	\usepackage[no-math]{fontspec}
	\setmainfont{TeX Gyre Pagella}
	\newfontfamily\EZqag{TeX Gyre Adventor}[NFSSFamily=qag,Ligatures=TeX]
\fi
\usepackage[many]{tcolorbox}
\usetikzlibrary{math,through,calc,intersections,angles,quotes,shapes,shapes.geometric,arrows,patterns,matrix,chains,arrows.meta,decorations.shapes,decorations.fractals,decorations.markings,shadows,decorations.pathreplacing}
\usepackage[top=1.6cm, bottom=1.8cm, left=1.8cm, right=1.8cm]{geometry}% v2: top 1.6cm (bản cũ 1.5cm) để chữ không sát dải đầu trang
\providecommand{\CheDo}{loigiai}% loigiai | dethi | solcolor | color | book
\expandafter\usepackage\expandafter[\CheDo]{ex_test}
% \tenchuong: khai báo ở file chính (trước \input) hoặc \renewcommand sau \input
%============================================================
%  [GỘP TỪ eduzy-de.tex] EDUZY - MÀU NÂU ĐỒNG / VÀNG KIM
%  + ĐẦU TRANG, CHÂN TRANG KIỂU DẢI NGANG (v2), ĐỐI XỨNG TRANG CHẴN / TRANG LẺ.
%  (không cần file eduzy-de.tex riêng nữa)
%  Khối này phải đặt SAU \usepackage[...]{ex_test} và SAU \def\tenchuong{...}
%============================================================
\usepackage{lastpage}
\usepackage{needspace}
%--- DÀN TRANG ĐỀ THI: tách câu (THÊM 05/11/2026; chỉ thêm lệnh mới, KHÔNG đổi hành vi mặc định) ---
%  Hai công tắc, mặc định tắt (= preamble chạy y như cũ). Bật bằng cách đặt TRƯỚC %============================================================
%  PREAMBLE EDUZY v2 - dùng chung cho đề thi, bài học, chuyên đề (gói ex_test)
%  Khác bản preamble-eduzy.tex ở đầu trang, chân trang: dải nền nâu đồng tràn ngang,
%  logo + Eduzy.com.vn, ô chữ vàng kem cạnh xiên, chân trang "Trang (số)";
%  logo Eduzy chìm giữa mọi trang (tắt: \logochimtat, bật lại: \logochimbat).
%  Cách dùng ở file chính (xem mau-chuan.tex, mau-chuyen-de.tex):
%      \def\CheDo{loigiai}          % loigiai | dethi | solcolor | color | book
%      \def\tenchuong{<dòng chữ trên đầu trang>}
%      \def\thuonghieu{<chữ ở chân trang lẻ>}   % không khai báo thì dùng mặc định
%      \def\tacgia{<chữ ở chân trang chẵn>}      % không khai báo thì dùng mặc định
%      \def\khauhieu{<chữ cạnh logo ở đầu trang>}  % tùy chọn, mặc định Eduzy.com.vn
%      \input{preamble-eduzy-v2}
%      \TLdapso                      % tài liệu: \TLfull | \TLdapso | \TLan
%      \begin{document} ... \end{document}
%  Logo Eduzy vẽ bằng TikZ ngay trong file (\eduzylogo), không cần file ảnh.
%  Biên dịch: pdfLaTeX (máy có gói vietnam) hoặc XeLaTeX.
%============================================================
\documentclass[12pt,a4paper,twoside]{article}%<-- twoside: bat dau trang/chan trang chan-le
\usepackage{iftex}
\usepackage{amsmath,amssymb,tikz,mathrsfs,tkz-tab}
\usepackage{graphicx}
\usepackage{pgfplots}% đồ thị hàm số
\pgfplotsset{compat=1.17}
%\usepackage{mathptmx}
\usepackage{mathpazo}
\ifXeTeX% XeLaTeX: font chữ có đủ dấu tiếng Việt, giữ nguyên font toán mathpazo
	\usepackage[no-math]{fontspec}
	\setmainfont{TeX Gyre Pagella}
	\newfontfamily\EZqag{TeX Gyre Adventor}[NFSSFamily=qag,Ligatures=TeX]
\fi
\usepackage[many]{tcolorbox}
\usetikzlibrary{math,through,calc,intersections,angles,quotes,shapes,shapes.geometric,arrows,patterns,matrix,chains,arrows.meta,decorations.shapes,decorations.fractals,decorations.markings,shadows,decorations.pathreplacing}
\usepackage[top=1.6cm, bottom=1.8cm, left=1.8cm, right=1.8cm]{geometry}% v2: top 1.6cm (bản cũ 1.5cm) để chữ không sát dải đầu trang
\providecommand{\CheDo}{loigiai}% loigiai | dethi | solcolor | color | book
\expandafter\usepackage\expandafter[\CheDo]{ex_test}
% \tenchuong: khai báo ở file chính (trước \input) hoặc \renewcommand sau \input
%============================================================
%  [GỘP TỪ eduzy-de.tex] EDUZY - MÀU NÂU ĐỒNG / VÀNG KIM
%  + ĐẦU TRANG, CHÂN TRANG KIỂU DẢI NGANG (v2), ĐỐI XỨNG TRANG CHẴN / TRANG LẺ.
%  (không cần file eduzy-de.tex riêng nữa)
%  Khối này phải đặt SAU \usepackage[...]{ex_test} và SAU \def\tenchuong{...}
%============================================================
\usepackage{lastpage}
\usepackage{needspace}
%--- DÀN TRANG ĐỀ THI: tách câu (THÊM 05/11/2026; chỉ thêm lệnh mới, KHÔNG đổi hành vi mặc định) ---
%  Hai công tắc, mặc định tắt (= preamble chạy y như cũ). Bật bằng cách đặt TRƯỚC \input{preamble-eduzy-v2}
%  trong file đề (tools/dan_trang/dan_trang.py tự ghi khi đặt \NeedspaceCau):
%    \def\TachCauDS{1}   câu Đúng/Sai được tách: câu dẫn (kèm hình) ở cuối trang, bảng Phát biểu sang trang sau
%                        (câu dẫn và hình không tách nhau, bảng không tách giữa hai trang)
%    \def\TachCauLG{1}   bản lời giải (\CheDo khác dethi): phần đề của câu cao hơn 1/3 trang được phép tách
%  Lệnh dùng trong file đề (giá trị do tools/dan_trang/dan_trang.py đo ở bản dethi):
%    \NeedspaceCau[<cao câu dẫn + hình>]{<cao cả khối câu>}   ngay trên \begin{ex} (câu đầu phần: trên tiêu đề phần)
%        công tắc tắt: đúng bằng \Needspace{<cao cả khối câu>} như cũ;
%        \TachCauDS=1 và có [..] (câu Đúng/Sai được chọn tách): chỉ giữ câu dẫn + hình;
%        \TachCauLG=1 ở bản lời giải: đòi tối đa 1/3 \textheight.
%    \NeedspaceBang{<cao bảng Phát biểu>}%   ngay trước \choiceTFt: \TachCauDS=1 thì giữ trọn bảng trên một trang;
%        công tắc tắt: không làm gì.
\providecommand{\TachCauDS}{0}
\providecommand{\TachCauLG}{0}
\newlength{\EZnsCau}
\newcommand{\NeedspaceCau}[2][]{%
	\setlength{\EZnsCau}{#2}%
	\ifnum\TachCauDS=1 \if\relax\detokenize{#1}\relax\else\setlength{\EZnsCau}{#1}\fi\fi
	\ifnum\TachCauLG=1 \ifdefstring{\kindOPTN}{dethi}{}{\ifdim\EZnsCau>0.3333\textheight\setlength{\EZnsCau}{0.3333\textheight}\fi}\fi
	\Needspace{\EZnsCau}}
\newcommand{\NeedspaceBang}[1]{\ifnum\TachCauDS=1 \Needspace{#1}\fi}

%--- Thang màu: đậm -> nhạt
\definecolor{gold1}{RGB}{126,96,0}    % nâu đồng đậm  #7E6000
\definecolor{gold2}{RGB}{192,144,3}   % vàng kim      #C09003
\definecolor{gold3}{RGB}{255,215,99}  % vàng nhạt     #FFD763
\definecolor{gold4}{RGB}{255,231,153} % vàng rất nhạt #FFE799
\definecolor{gold5}{RGB}{255,243,205} % vàng kem      #FFF3CD
%--- Ghi đè tông xanh/đỏ cũ
\definecolor{myblue}{RGB}{126,96,0}
\colorlet{blue}{gold1}
\colorlet{cyan}{gold2}
\colorlet{teal}{gold2}

%--- Tên chương mặc định (file đề nên tự khai báo lại)
\providecommand{\tenchuong}{}
%--- Chân trang
\providecommand{\thuonghieu}{TOÁN 12 -- Eduzy.com.vn}
\providecommand{\tacgia}{Tô Gia Bảo -- 0353703783}

%--- Chữ cạnh logo ở đầu trang
\providecommand{\khauhieu}{Eduzy.com.vn}

%====================== THÔNG SỐ ĐẦU TRANG, CHÂN TRANG (v2) ======================
%  Dải nền đậm tràn ngang trang; ô chữ vàng kem có cạnh xiên, kèm 2 vạch xiên chuyển màu.
%  Đầu trang lẻ : logo + \khauhieu (trái)   | ô chữ \tenchuong (phải)
%  Chân trang lẻ: ô chữ \thuonghieu (trái)  | Trang (số) (phải)
%  Trang chẵn soi gương trang lẻ; chân trang chẵn ghi \tacgia.
%--- Màu: chỉ dùng thang vàng kim Eduzy
\colorlet{ezNen}{gold1}% dải nền đậm
\colorlet{ezVach1}{gold3}% vạch xiên sát dải nền (đậm hơn)
\colorlet{ezVach2}{gold4}% vạch xiên sát ô chữ (nhạt hơn)
\colorlet{ezO}{gold5}% ô chữ vàng kem
%--- Kích thước
\newlength{\EZdai}\setlength{\EZdai}{1.1cm}% chiều cao dải nền
\newlength{\EZo}\setlength{\EZo}{0.68cm}% chiều cao ô chữ
\newlength{\EZxien}\setlength{\EZxien}{0.55cm}% độ lệch ngang của cạnh xiên
\newlength{\EZvach}\setlength{\EZvach}{0.75cm}% bề rộng mỗi vạch xiên
\newlength{\EZkhe}\setlength{\EZkhe}{0.12cm}% khe trắng giữa vạch và ô chữ
\newlength{\EZdem}\setlength{\EZdem}{0.4cm}% khoảng đệm từ chữ tới cạnh xiên
\newlength{\EZle}\setlength{\EZle}{1.8cm}% lề trái/phải của chữ (bằng lề thân bài)
\newlength{\EZrongmin}\setlength{\EZrongmin}{4.5cm}% ô chữ rộng tối thiểu
\newlength{\EZrong}% bề rộng chữ trong ô (tính tự động)
%--- Font
\newcommand{\EZfontO}{\fontsize{10.5pt}{12pt}\selectfont\bfseries}% chữ trong ô
\newcommand{\EZfontKH}{\fontfamily{qag}\fontsize{11pt}{12pt}\selectfont\bfseries}% \khauhieu cạnh logo
\newcommand{\EZfontTrang}{\fontsize{11.5pt}{12pt}\selectfont\bfseries}% chữ "Trang" và số trang
%--- Đo bề rộng chữ trong ô
\newcommand{\EZdorong}[1]{%
	\settowidth{\EZrong}{\EZfontO #1}%
	\ifdim\EZrong<\EZrongmin\setlength{\EZrong}{\EZrongmin}\fi}
%--- Logo + chữ trắng \khauhieu (dựng sẵn một lần khi vào \begin{document})
\newsavebox{\EZlogobox}
\AtBeginDocument{\sbox{\EZlogobox}{%
		\raisebox{-0.5\height}{\eduzylogo[1cm]}\hspace{0.15cm}%
		\raisebox{-0.5\height}{\color{white}\EZfontKH\khauhieu}}}
%--- Hai vạch xiên + khe trắng. #1: điểm gốc ở mép dưới, #2: hướng (1 sang phải, -1 sang trái),
%    #3: hướng nghiêng (1 kiểu "/", -1 kiểu "\")
\newcommand{\EZvachxien}[3]{%
	\fill[white] (#1) --++({(#3)*\EZxien},\EZo) --++({(#2)*\EZkhe},0) --++({-(#3)*\EZxien},-\EZo) --cycle;
	\fill[ezVach2] ([xshift=(#2)*\EZkhe]#1) --++({(#3)*\EZxien},\EZo) --++({(#2)*\EZvach},0) --++({-(#3)*\EZxien},-\EZo) --cycle;
	\fill[ezVach1] ([xshift=(#2)*(\EZkhe+\EZvach)]#1) --++({(#3)*\EZxien},\EZo) --++({(#2)*\EZvach},0) --++({-(#3)*\EZxien},-\EZo) --cycle;
	\draw[white,line width=0.6pt] ([xshift=(#2)*(\EZkhe+\EZvach)]#1) --++({(#3)*\EZxien},\EZo);
}
%--- Chữ trong ô (căn giữa theo chiều cao ô, không lệch vì dấu tiếng Việt)
\tikzset{ezchu/.style={inner sep=0pt,text=ezNen,text height=0.72em,text depth=0.22em}}

%====================== LOGO CHÌM GIỮA TRANG ======================
%  Vẽ cùng đầu trang nên nằm dưới chữ của thân bài, có trên mọi trang.
%  Tắt: \logochimtat (từ trang đó trở đi), bật lại: \logochimbat.
%  Đổi cỡ, độ đậm: khai báo trước \input, ví dụ \def\EZchimcao{9cm} \def\EZchimdo{0.1}
\providecommand{\EZchimcao}{10cm}% chiều cao logo chìm
\providecommand{\EZchimdo}{0.12}% độ đậm (0 = trong suốt, 1 = rõ hẳn)
\newif\ifEZchim \EZchimtrue
\newcommand{\logochimtat}{\EZchimfalse}
\newcommand{\logochimbat}{\EZchimtrue}
%  Dựng sẵn logo mờ một lần khi vào \begin{document}; transparency group làm mờ cả logo
%  một lượt nên không lộ vệt ở chỗ các dải màu chồng nhau.
\newsavebox{\EZchimtho}
\newsavebox{\EZchimbox}
\AtBeginDocument{%
	\sbox{\EZchimtho}{\eduzylogo[\EZchimcao]}%
	\sbox{\EZchimbox}{\begin{tikzpicture}
			\begin{scope}[transparency group,opacity=\EZchimdo]
				\node[inner sep=0pt] {\usebox{\EZchimtho}};
			\end{scope}
	\end{tikzpicture}}}
\newcommand{\EZvechim}{%
	\ifEZchim
		\node[inner sep=0pt] at (current page.center) {\usebox{\EZchimbox}};
	\fi}

%====================== ĐẦU TRANG - TRANG LẺ ======================
%  (logo bên trái, ô chữ \tenchuong bên phải)
\newcommand{\myfancyheadOdd}{%
	\EZdorong{\tenchuong}%
	\begin{tikzpicture}[remember picture,overlay]
		\coordinate (T) at (current page.north west);
		\coordinate (P) at (current page.north east);
		%--- logo chìm giữa trang
		\EZvechim
		%--- dải nền
		\fill[ezNen] (T) rectangle ([yshift=-\EZdai]P);
		%--- ô chữ: X là góc dưới bên trái, cạnh trái xiên kiểu "/"
		\coordinate (X) at ([xshift=-\EZle-\EZrong-\EZdem-\EZxien,yshift=-\EZdai]P);
		\fill[ezO] (X) --++(\EZxien,\EZo) -- ([yshift=-\EZdai+\EZo]P) -- ([yshift=-\EZdai]P) --cycle;
		\EZvachxien{X}{-1}{1}
		\node[ezchu,anchor=east] at ([xshift=-\EZle,yshift=-\EZdai+0.5*\EZo]P) {\EZfontO\tenchuong};
		%--- logo + tên
		\node[inner sep=0pt,anchor=west] at ([xshift=\EZle,yshift=-0.5*\EZdai]T) {\usebox{\EZlogobox}};
	\end{tikzpicture}%
}

%====================== ĐẦU TRANG - TRANG CHẴN =====================
%  (soi gương: ô chữ \tenchuong bên trái, logo bên phải)
\newcommand{\myfancyheadEven}{%
	\EZdorong{\tenchuong}%
	\begin{tikzpicture}[remember picture,overlay]
		\coordinate (T) at (current page.north west);
		\coordinate (P) at (current page.north east);
		%--- logo chìm giữa trang
		\EZvechim
		%--- dải nền
		\fill[ezNen] (T) rectangle ([yshift=-\EZdai]P);
		%--- ô chữ: X là góc dưới bên phải, cạnh phải xiên kiểu "\"
		\coordinate (X) at ([xshift=\EZle+\EZrong+\EZdem+\EZxien,yshift=-\EZdai]T);
		\fill[ezO] (X) --++(-\EZxien,\EZo) -- ([yshift=-\EZdai+\EZo]T) -- ([yshift=-\EZdai]T) --cycle;
		\EZvachxien{X}{1}{-1}
		\node[ezchu,anchor=west] at ([xshift=\EZle,yshift=-\EZdai+0.5*\EZo]T) {\EZfontO\tenchuong};
		%--- logo + tên
		\node[inner sep=0pt,anchor=east] at ([xshift=-\EZle,yshift=-0.5*\EZdai]P) {\usebox{\EZlogobox}};
	\end{tikzpicture}%
}

%====================== CHÂN TRANG - TRANG LẺ ======================
%  (ô chữ \thuonghieu bên trái, "Trang (số)" bên phải)
\newcommand{\myfancyfootOdd}{%
	\EZdorong{\thuonghieu}%
	\begin{tikzpicture}[remember picture,overlay]
		\coordinate (T) at (current page.south west);
		\coordinate (P) at (current page.south east);
		%--- dải nền
		\fill[ezNen] (T) rectangle ([yshift=\EZdai]P);
		%--- ô chữ sát mép trên dải: X là góc dưới bên phải, cạnh phải xiên kiểu "/"
		\coordinate (X) at ([xshift=\EZle+\EZrong+\EZdem,yshift=\EZdai-\EZo]T);
		\fill[ezO] ([yshift=\EZdai-\EZo]T) -- (X) --++(\EZxien,\EZo) -- ([yshift=\EZdai]T) --cycle;
		\EZvachxien{X}{1}{1}
		\node[ezchu,anchor=west] at ([xshift=\EZle,yshift=\EZdai-0.5*\EZo]T) {\EZfontO\thuonghieu};
		%--- số trang
		\node[circle,draw=white,line width=0.9pt,text=white,inner sep=0pt,minimum size=0.62cm,anchor=east]
		(so) at ([xshift=-\EZle,yshift=0.5*\EZdai]P) {\EZfontTrang\thepage};
		\node[inner sep=0pt,text=white,anchor=base east] (tr) at ([xshift=-0.15cm]so.base west) {\EZfontTrang Trang};
	\end{tikzpicture}%
}

%====================== CHÂN TRANG - TRANG CHẴN ====================
%  (soi gương: "Trang (số)" bên trái, ô chữ \tacgia bên phải)
\newcommand{\myfancyfootEven}{%
	\EZdorong{\tacgia}%
	\begin{tikzpicture}[remember picture,overlay]
		\coordinate (T) at (current page.south west);
		\coordinate (P) at (current page.south east);
		%--- dải nền
		\fill[ezNen] (T) rectangle ([yshift=\EZdai]P);
		%--- ô chữ sát mép trên dải: X là góc dưới bên trái, cạnh trái xiên kiểu "\"
		\coordinate (X) at ([xshift=-\EZle-\EZrong-\EZdem,yshift=\EZdai-\EZo]P);
		\fill[ezO] ([yshift=\EZdai-\EZo]P) -- (X) --++(-\EZxien,\EZo) -- ([yshift=\EZdai]P) --cycle;
		\EZvachxien{X}{-1}{-1}
		\node[ezchu,anchor=east] at ([xshift=-\EZle,yshift=\EZdai-0.5*\EZo]P) {\EZfontO\tacgia};
		%--- số trang
		%    (vòng tròn ẩn so0 nằm giữa dải để lấy đường chân chữ, giống trang lẻ)
		\node[circle,inner sep=0pt,minimum size=0.62cm] (so0) at ([yshift=0.5*\EZdai]T) {\EZfontTrang\phantom{\thepage}};
		\node[inner sep=0pt,text=white,anchor=base west] (tr) at ($(so0.base -| T)+(\EZle,0)$) {\EZfontTrang Trang};
		\node[circle,draw=white,line width=0.9pt,text=white,inner sep=0pt,minimum size=0.62cm,anchor=base west]
		(so) at ([xshift=0.15cm]tr.base east) {\EZfontTrang\thepage};
	\end{tikzpicture}%
}

%--- Giữ tương thích với các file cũ đang gọi \myfancyhead, \myfancyfoot
\newcommand{\myfancyhead}{\myfancyheadOdd}
\newcommand{\myfancyfoot}{\myfancyfootOdd}

%======================= GẮN VÀO TRANG =============================
\pagestyle{fancy}
\fancyhf{}
\fancyhead[LO]{\myfancyheadOdd}
\fancyhead[RE]{\myfancyheadEven}
\fancyfoot[LO]{\myfancyfootOdd}
\fancyfoot[RE]{\myfancyfootEven}
\renewcommand{\headrulewidth}{0pt}
\renewcommand{\footrulewidth}{0pt}
\fancypagestyle{plain}{%
	\fancyhf{}
	\fancyhead[LO]{\myfancyheadOdd}
	\fancyhead[RE]{\myfancyheadEven}
	\fancyfoot[LO]{\myfancyfootOdd}
	\fancyfoot[RE]{\myfancyfootEven}
	\renewcommand{\headrulewidth}{0pt}
	\renewcommand{\footrulewidth}{0pt}
}
%=================== HẾT PHẦN GỘP TỪ eduzy-de.tex ===================
\newcommand{\indam}[1]{
	{\fontsize{12pt}{9pt}\sffamily \selectfont\bfseries\color{gold1} #1}
}
\newcommand{\ind}[1]{
	{\fontsize{14pt}{9pt}\fontfamily{qag} \selectfont\bfseries\color{gold1} #1}
}
\renewcommand{\FalseEX}{\stepcounter{dapan}\textbf{\selectfont\color{gold1}\Alph{dapan}.}\ignorespaces}
\renewtheorem{ex}{\color{gold1}Câu}
\renewcommand{\baselinestretch}{1.1}
\newcommand{\hoac}[1]{
	\left[\begin{aligned}#1\end{aligned}\right.}
\newcommand{\heva}[1]{
	\left\{\begin{aligned}#1\end{aligned}\right.}
\usepackage{esvect}
\def\vec{\vv}
\def\vt{\vv}
\def\overrightarrow{\vv}
%====================== LOGO EDUZY VẼ BẰNG TIKZ ======================
%  Đường viền dò từ logo gốc; dải màu vàng kim đo trực tiếp trên ảnh gốc
%  (nghiêng 21 độ, tô bằng 80 dải hẹp nên hiển thị giống nhau trên mọi trình xem PDF).
%  Dùng: \eduzylogo          -> cao 2.4cm (mặc định)
%        \eduzylogo[1.5cm]   -> cao 1.5cm
\definecolor{ezText}{RGB}{245,200,46}
\newcommand{\eduzylogo}[1][2.4cm]{%
	\resizebox{!}{#1}{%
		\begin{tikzpicture}[x=0.01cm,y=0.01cm]
			\path[use as bounding box] (0,0) rectangle (224,-224);
			% mũ tốt nghiệp + vương miện
			\begin{scope}
				\clip
				(119.4,-176.1)..controls(118.9,-176.1)and(116.8,-176)..(114.6,-176)
				..controls(101.4,-176)and(85.8,-175.9)..(83.7,-175.7)
				..controls(81.8,-175.6)and(81,-175.6)..(80.2,-175.7)
				..controls(78.8,-175.9)and(77.7,-175.9)..(76.2,-175.7)
				..controls(74.9,-175.5)and(70.8,-175.5)..(64.3,-175.7)
				..controls(59.1,-175.9)and(58.3,-175.7)..(58.6,-173.7)
				..controls(58.6,-172.8)and(59.6,-169.6)..(60.3,-168)
				..controls(60.8,-166.9)and(61.4,-165.8)..(62.3,-165)
				..controls(63.1,-164.2)and(64.3,-163.4)..(65,-163.2)
				..controls(65.5,-163.1)and(72.3,-163)..(76.6,-163.1)
				..controls(77.9,-163.1)and(80.1,-163.1)..(81.4,-163)
				..controls(82.8,-162.9)and(86.3,-162.9)..(89.2,-162.9)
				..controls(93,-162.9)and(94.9,-162.8)..(96,-162.7)
				..controls(97.2,-162.5)and(101,-162.5)..(114,-162.4)
				..controls(123,-162.3)and(134.9,-162.2)..(140.2,-162.1)
				..controls(150.1,-162)and(152.2,-162.1)..(153.6,-162.4)
				..controls(154.4,-162.6)and(157.1,-163.8)..(158.1,-164.4)
				..controls(158.8,-164.8)and(159.7,-165.7)..(159.9,-166.3)
				..controls(160,-166.5)and(160.2,-167.1)..(160.2,-167.6)
				..controls(160.3,-168)and(160.6,-169.4)..(160.9,-170.6)
				..controls(161.2,-171.7)and(161.5,-173.1)..(161.6,-173.6)
				..controls(161.8,-174.4)and(161.8,-174.5)..(161.5,-174.9)
				..controls(161.2,-175.6)and(160.8,-175.9)..(159.9,-176)
				..controls(159.2,-176.1)and(154.1,-176.2)..(151.5,-176.2)
				..controls(148.1,-176.1)and(122.6,-176.1)..(121.6,-176.2)
				..controls(120.8,-176.2)and(119.8,-176.2)..(119.4,-176.1)--cycle
				(130.4,-160)..controls(129,-160)and(119.8,-160)..(110.1,-159.9)
				..controls(100.3,-159.8)and(88.5,-159.7)..(83.9,-159.6)
				..controls(77.8,-159.5)and(75.3,-159.5)..(74.5,-159.6)
				..controls(73.9,-159.7)and(72.5,-159.8)..(71.3,-159.8)
				..controls(69.6,-159.8)and(69,-159.8)..(68.5,-159.6)
				..controls(67.6,-159.3)and(66.8,-158.6)..(66.2,-157.6)
				..controls(66,-157.3)and(65.6,-156.5)..(65.3,-156.1)
				..controls(64.6,-155.1)and(64.5,-154.9)..(63.7,-152.9)
				..controls(63.2,-151.6)and(62.9,-150.9)..(61.8,-149.6)
				..controls(61.5,-149.2)and(61.1,-148.5)..(60.9,-147.9)
				..controls(60.7,-147.4)and(60.3,-146.6)..(60.1,-146.1)
				..controls(59.9,-145.5)and(59.5,-144.9)..(59,-144.3)
				..controls(58.5,-143.7)and(58.2,-143.2)..(57.9,-142.5)
				..controls(57.7,-141.9)and(57.3,-141.1)..(57,-140.7)
				..controls(56.6,-140.3)and(56.1,-139.3)..(55.8,-138.5)
				..controls(55.5,-137.8)and(55.1,-136.8)..(54.9,-136.5)
				..controls(54.7,-136.1)and(54.3,-135.2)..(53.9,-134.5)
				..controls(53.6,-133.8)and(53,-132.7)..(52.6,-132)
				..controls(51.6,-130.5)and(51.4,-130.2)..(51.1,-129)
				..controls(50.6,-127.5)and(50.6,-127.4)..(49.7,-126.3)
				..controls(49.2,-125.7)and(48.6,-124.8)..(48.3,-124.2)
				..controls(48,-123.6)and(47.5,-122.7)..(47.3,-122.2)
				..controls(47,-121.8)and(46.5,-120.7)..(46.1,-119.8)
				..controls(45.8,-118.9)and(45.2,-117.8)..(44.9,-117.2)
				..controls(44.2,-116.1)and(43.3,-114.4)..(42.3,-112.2)
				..controls(41.8,-111.3)and(41.3,-110.4)..(41.1,-110.1)
				..controls(40.5,-109.4)and(40.3,-109)..(40,-107.8)
				..controls(39.8,-107)and(39.6,-106.5)..(39.2,-106)
				..controls(38.9,-105.6)and(38.4,-104.7)..(38,-104)
				..controls(37.7,-103.4)and(37.2,-102.5)..(37,-102.1)
				..controls(36.7,-101.7)and(36.3,-100.8)..(36,-100.1)
				..controls(35.6,-98.9)and(35.3,-98.4)..(34.6,-97.6)
				..controls(34.2,-97.1)and(34.2,-96.3)..(34.6,-96)
				..controls(35,-95.7)and(35.5,-95.9)..(36.2,-96.7)
				..controls(36.5,-97)and(37.2,-97.5)..(37.6,-97.7)..controls(38,-98)and(38.6,-98.4)..(39.1,-98.9)
				..controls(39.6,-99.4)and(40.3,-100)..(40.6,-100.3)
				..controls(41.3,-100.8)and(42.4,-101.8)..(43.7,-102.9)
				..controls(44,-103.2)and(44.6,-103.7)..(45,-103.9)
				..controls(45.4,-104.2)and(45.9,-104.6)..(46.1,-104.8)
				..controls(47.1,-105.9)and(48.8,-107.1)..(50.5,-107.9)
				..controls(51.8,-108.6)and(52.3,-109)..(53.4,-110)
				..controls(54.8,-111.2)and(55.6,-111.6)..(57,-112.1)
				..controls(57.5,-112.3)and(58.5,-112.8)..(59.4,-113.3)
				..controls(60.2,-113.7)and(61.3,-114.3)..(61.9,-114.5)
				..controls(62.4,-114.7)and(63.3,-115.1)..(63.8,-115.3)
				..controls(64.4,-115.6)and(65.4,-115.9)..(66.2,-116.1)
				..controls(67,-116.3)and(68.4,-116.6)..(69.2,-116.9)
				..controls(72.3,-117.9)and(72.8,-118)..(78,-118.1)
				..controls(84,-118.3)and(86.1,-118.1)..(88.1,-117.1)
				..controls(89.4,-116.5)and(93,-115)..(93.7,-114.9)
				..controls(94.8,-114.6)and(95.4,-114.2)..(95.9,-113.4)
				..controls(96.2,-113)and(96.7,-112.5)..(97.1,-112.2)
				..controls(98,-111.6)and(101.3,-108.4)..(102.1,-107.2)
				..controls(102.4,-106.8)and(102.9,-106.2)..(103.2,-105.8)
				..controls(103.5,-105.5)and(104,-104.8)..(104.2,-104.2)
				..controls(104.5,-103.6)and(105,-102.8)..(105.3,-102.4)
				..controls(105.6,-102)and(106.1,-101.2)..(106.5,-100.6)
				..controls(107.1,-99.5)and(108.6,-97.5)..(109.5,-96.5)
				..controls(109.9,-96.1)and(110.1,-96)..(110.4,-96)
				..controls(111.1,-96)and(111.6,-96.5)..(111.8,-97.6)
				..controls(112,-98.6)and(112.3,-99.3)..(112.8,-99.9)
				..controls(113.7,-100.9)and(115.2,-103.1)..(115.7,-104.2)
				..controls(116.1,-105)and(116.5,-105.5)..(116.9,-106)
				..controls(117.3,-106.4)and(117.8,-106.9)..(118,-107.3)
				..controls(118.2,-107.6)and(118.6,-108)..(119,-108.3)
				..controls(119.6,-108.8)and(121.1,-110.3)..(121.8,-111.2)
				..controls(122.1,-111.7)and(122.5,-112)..(123.4,-112.4)
				..controls(124.8,-113.2)and(125.8,-113.8)..(127,-114.6)
				..controls(127.5,-114.9)and(128.1,-115.3)..(128.2,-115.4)
				..controls(128.4,-115.5)and(129.3,-115.8)..(130.2,-116)
				..controls(131.1,-116.2)and(132.5,-116.6)..(133.3,-116.9)
				..controls(135.7,-117.9)and(136.1,-117.9)..(142.5,-117.8)
				..controls(145.8,-117.8)and(148.4,-117.7)..(148.8,-117.6)
				..controls(149.2,-117.6)and(150.2,-117.2)..(151,-116.8)
				..controls(152,-116.4)and(152.8,-116)..(153.6,-115.9)
				..controls(154.5,-115.7)and(155.2,-115.4)..(156.3,-114.9)
				..controls(157.1,-114.5)and(158.2,-114)..(158.7,-113.8)
				..controls(159.1,-113.7)and(159.8,-113.3)..(160.2,-113)
				..controls(160.6,-112.7)and(161.3,-112.4)..(161.7,-112.2)
				..controls(162.7,-111.9)and(163.2,-111.7)..(164.1,-111.1)
				..controls(164.5,-110.8)and(165.3,-110.4)..(165.9,-110)
				..controls(166.5,-109.7)and(167.3,-109.2)..(167.6,-108.9)
				..controls(167.9,-108.6)and(168.6,-108.2)..(169,-108)
				..controls(169.4,-107.8)and(170.2,-107.3)..(170.7,-106.9)
				..controls(171.2,-106.5)and(172.1,-105.8)..(172.7,-105.4)
				..controls(173.2,-105)and(174.1,-104.4)..(174.5,-104)
				..controls(174.9,-103.7)and(175.7,-103.1)..(176.2,-102.6)
				..controls(176.7,-102.2)and(177.4,-101.5)..(177.9,-101.2)
				..controls(178.3,-100.8)and(178.8,-100.4)..(179,-100.2)
				..controls(179.2,-100)and(179.6,-99.6)..(180,-99.3)
				..controls(180.4,-99.1)and(180.9,-98.6)..(181.2,-98.4)
				..controls(181.5,-98.1)and(182.1,-97.7)..(182.5,-97.5)
				..controls(183.2,-97.2)and(183.5,-96.9)..(184.3,-95.9)
				..controls(184.6,-95.6)and(184.6,-95.6)..(184.8,-95.7)
				..controls(185,-95.8)and(185.1,-96)..(185.1,-96.2)
				..controls(185.1,-96.5)and(185,-96.7)..(184.6,-97.2)
				..controls(184.1,-97.8)and(184,-97.9)..(183.9,-98.8)
				..controls(183.7,-100.3)and(183.3,-101.3)..(182.9,-101.8)
				..controls(182.6,-102.1)and(182.3,-102.6)..(182.2,-102.9)
				..controls(182,-103.3)and(181.6,-104)..(181.2,-104.6)
				..controls(180.9,-105.2)and(180.5,-106.1)..(180.4,-106.6)
				..controls(179.9,-108)and(179.7,-108.5)..(179,-109.4)
				..controls(178.6,-109.8)and(178.2,-110.5)..(177.9,-111)
				..controls(177.7,-111.5)and(177.3,-112.3)..(177,-112.7)
				..controls(176.6,-113.2)and(176.4,-113.7)..(176.3,-114.2)
				..controls(176.2,-114.7)and(175.9,-115.2)..(175.4,-116)
				..controls(175,-116.6)and(174.4,-117.7)..(174,-118.5)
				..controls(173.7,-119.2)and(173.3,-120.1)..(173,-120.5)
				..controls(172.8,-120.9)and(172.4,-121.8)..(172.1,-122.5)
				..controls(171.8,-123.2)and(171.3,-124.1)..(171,-124.6)
				..controls(170.8,-125.1)and(170.3,-126)..(169.9,-126.6)
				..controls(169.6,-127.3)and(169.2,-128.1)..(169,-128.3)
				..controls(168.8,-128.6)and(168.5,-129.3)..(168.3,-129.8)
				..controls(168.1,-130.3)and(167.8,-131.1)..(167.5,-131.6)
				..controls(166.9,-132.5)and(164.2,-137.9)..(164.2,-138.1)
				..controls(164.2,-138.1)and(163.9,-138.8)..(163.4,-139.7)
				..controls(162.9,-140.5)and(162.4,-141.5)..(162.3,-141.9)
				..controls(161.9,-143)and(161.6,-143.4)..(160.9,-144.1)
				..controls(160.4,-144.6)and(160.3,-144.8)..(160.1,-145.4)
				..controls(160,-145.8)and(159.6,-146.8)..(159.2,-147.8)
				..controls(158.7,-148.8)and(158.1,-150)..(157.9,-150.6)
				..controls(157.6,-151.1)and(157.2,-151.9)..(156.9,-152.4)
				..controls(156.5,-152.8)and(156.1,-153.7)..(155.8,-154.4)
				..controls(155.6,-155)and(155.2,-155.9)..(154.9,-156.3)
				..controls(154.7,-156.7)and(154.4,-157.4)..(154.2,-157.9)
				..controls(154.1,-158.4)and(153.8,-159)..(153.5,-159.3)
				..controls(152.8,-160.1)and(153,-160.1)..(142.4,-160.1)
				..controls(137.2,-160.1)and(131.8,-160.1)..(130.4,-160)--cycle
				(141.5,-103.7)..controls(140.8,-103)and(139.8,-102.1)..(139.2,-101.6)
				..controls(138.5,-101.2)and(137.9,-100.6)..(137.8,-100.4)
				..controls(137,-99.5)and(134,-96.5)..(133.2,-96)
				..controls(132.7,-95.7)and(131.8,-94.9)..(131.1,-94.3)
				..controls(129.4,-92.7)and(127.9,-91.6)..(126.6,-91.1)
				..controls(125.2,-90.5)and(124.1,-90)..(122.9,-89.1)
				..controls(122.1,-88.6)and(121.3,-88.2)..(119.9,-87.7)
				..controls(118.8,-87.3)and(117.5,-86.8)..(116.9,-86.5)
				..controls(114.3,-85.5)and(113.7,-85.4)..(109.7,-85.5)
				..controls(105.3,-85.7)and(103.9,-85.9)..(101.6,-87)
				..controls(100.7,-87.4)and(99.4,-87.8)..(97.9,-88)
				..controls(97,-88.2)and(96.6,-88.3)..(95.8,-88.9)
				..controls(95.4,-89.2)and(94.7,-89.6)..(94.2,-89.8)
				..controls(92.9,-90.3)and(92.3,-90.7)..(90.9,-92.1)
				..controls(90.2,-92.7)and(89.4,-93.3)..(89,-93.5)..controls(87,-94.5)and(85.2,-96)..(83.6,-97.9)
				..controls(83.3,-98.2)and(82.5,-98.9)..(81.8,-99.4)
				..controls(80.7,-100.1)and(80.5,-100.3)..(80,-101.1)
				..controls(78.7,-103.1)and(77.2,-103.9)..(76.4,-103)
				..controls(75.7,-102.2)and(75.7,-103.2)..(75.7,-83.6)
				..controls(75.6,-73.6)and(75.6,-64.8)..(75.5,-64)
				..controls(75.5,-63.2)and(75.5,-61.4)..(75.6,-59.8)--(75.7,-57.1)--(76,-56.8)
				..controls(76.8,-56.2)and(78,-56.3)..(79.8,-57.2)
				..controls(80.4,-57.5)and(81.3,-57.8)..(81.8,-57.9)
				..controls(82.3,-58.1)and(83,-58.4)..(83.6,-58.8)..controls(84.3,-59.2)and(85,-59.5)..(86,-59.8)
				..controls(86.8,-60.1)and(88,-60.6)..(88.7,-60.9)
				..controls(89.3,-61.2)and(90.6,-61.7)..(91.4,-62)
				..controls(92.2,-62.3)and(93.7,-62.9)..(94.6,-63.3)
				..controls(95.6,-63.7)and(96.8,-64.1)..(97.4,-64.3)
				..controls(98,-64.5)and(99,-64.9)..(99.7,-65.2)
				..controls(100.4,-65.6)and(101.5,-66)..(102.3,-66.2)
				..controls(104.2,-66.8)and(105,-67.1)..(106.2,-67.8)
				..controls(107.9,-68.7)and(110.2,-69.4)..(111.3,-69.2)
				..controls(112.5,-69)and(115,-68.2)..(115.6,-67.8)
				..controls(117.3,-66.7)and(117.6,-66.6)..(119.5,-66.2)
				..controls(121.9,-65.8)and(122.1,-65.7)..(123,-65.1)
				..controls(124.1,-64.4)and(124.4,-64.3)..(126,-63.9)
				..controls(127,-63.7)and(127.7,-63.4)..(128.4,-63.1)
				..controls(129,-62.8)and(129.7,-62.5)..(130.6,-62.4)
				..controls(132.1,-62.1)and(134,-61.5)..(135.2,-61)
				..controls(135.7,-60.8)and(136.7,-60.5)..(137.3,-60.3)
				..controls(138.4,-60)and(139.3,-59.7)..(141.1,-59)
				..controls(141.6,-58.7)and(142,-58.6)..(142.5,-58.6)
				..controls(143.2,-58.6)and(143.3,-58.7)..(143.6,-59)
				..controls(144,-59.5)and(144.1,-59.9)..(144,-60.8)
				..controls(144,-61.2)and(144,-71.1)..(144,-82.8)--(143.9,-104.1)--(143.7,-104.5)
				..controls(143.5,-104.7)and(143.3,-104.8)..(143,-104.9)
				..controls(142.6,-104.9)and(142.5,-104.8)..(141.5,-103.7)--cycle
				(27.8,-91.4)..controls(27.1,-90.9)and(26.3,-90)..(26.1,-89.5)
				..controls(26,-89.2)and(25.8,-88.6)..(25.7,-87.9)
				..controls(25.4,-86.2)and(25.6,-84.8)..(26.5,-83.5)
				..controls(27.4,-82.3)and(28.3,-82)..(29.4,-82.3)
				..controls(31.1,-82.8)and(32.1,-84.9)..(32,-87.4)
				..controls(31.9,-89.3)and(31.2,-90.6)..(29.9,-91.4)
				..controls(29.1,-91.9)and(28.6,-91.9)..(27.8,-91.4)--cycle
				(190.1,-91.3)..controls(189.7,-91.2)and(189.3,-90.9)..(189.1,-90.7)
				..controls(188.3,-89.8)and(187.6,-87.7)..(187.6,-86.2)
				..controls(187.6,-85.5)and(187.7,-85.2)..(188.2,-84.4)
				..controls(189.1,-82.6)and(190.1,-81.8)..(191.3,-82.1)
				..controls(192,-82.3)and(193.5,-83.7)..(193.8,-84.4)
				..controls(194.1,-85.1)and(194.2,-87.1)..(194,-88.3)
				..controls(193.8,-89.8)and(193.3,-90.5)..(192.3,-91.3)
				..controls(191.7,-91.7)and(190.9,-91.7)..(190.1,-91.3)--cycle
				(173.7,-83.4)..controls(172.2,-82.5)and(172,-82.5)..(170.6,-82.4)
				..controls(169.7,-82.4)and(168.8,-82.4)..(168.3,-82.5)
				..controls(167.2,-82.7)and(166.9,-82.7)..(166.6,-82.2)
				..controls(166.3,-81.8)and(166.3,-81)..(166.8,-80)
				..controls(167.2,-78.9)and(167.4,-77.4)..(167.7,-74.1)
				..controls(167.8,-72.8)and(167.9,-71.8)..(167.9,-71.7)
				..controls(168,-71.7)and(168,-70.2)..(168,-68.4)..controls(168,-65.4)and(168,-64.9)..(167.7,-63)
				..controls(167.3,-60.5)and(167.1,-59.9)..(166.4,-58.4)
				..controls(166.1,-57.7)and(165.8,-56.9)..(165.7,-56.6)
				..controls(165.7,-56.2)and(165.4,-55.6)..(165.2,-55.2)
				..controls(164.6,-53.8)and(161.9,-50.7)..(161,-50.3)
				..controls(159.8,-49.7)and(157.9,-50)..(155.2,-51.2)
				..controls(154.7,-51.4)and(153.4,-51.8)..(152.2,-52.1)
				..controls(150.8,-52.4)and(149.4,-52.9)..(148.6,-53.2)
				..controls(147.8,-53.6)and(146.7,-54)..(146,-54.2)
				..controls(145.4,-54.4)and(144.5,-54.7)..(144,-54.9)
				..controls(143.6,-55.2)and(142.7,-55.5)..(141.8,-55.7)
				..controls(141,-55.9)and(139.6,-56.3)..(138.7,-56.6)
				..controls(137.8,-57)and(136.3,-57.4)..(135.3,-57.6)
				..controls(133.8,-57.9)and(133.2,-58.1)..(132.1,-58.7)
				..controls(131.1,-59.2)and(130.3,-59.4)..(129.4,-59.6)
				..controls(128.7,-59.8)and(127.8,-60)..(127.4,-60.1)
				..controls(127,-60.3)and(126,-60.5)..(125.1,-60.7)
				..controls(123.2,-61.1)and(121.7,-61.5)..(120.2,-62.1)
				..controls(119.7,-62.3)and(118.1,-62.8)..(116.7,-63.2)
				..controls(114.6,-63.9)and(114,-64.1)..(113,-64.7)
				..controls(111.9,-65.4)and(111.8,-65.4)..(111,-65.5)
				..controls(109.9,-65.5)and(109.3,-65.4)..(107.7,-64.7)
				..controls(107,-64.4)and(106,-64)..(105.3,-63.9)
				..controls(103.5,-63.5)and(102.2,-63)..(101.2,-62.5)
				..controls(100.1,-61.8)and(99.3,-61.5)..(97.8,-61.1)
				..controls(97.2,-61)and(96.2,-60.7)..(95.7,-60.4)
				..controls(95.1,-60.1)and(94.2,-59.8)..(93.5,-59.6)
				..controls(92.9,-59.4)and(91.7,-59)..(91,-58.7)..controls(90.2,-58.4)and(88.9,-58)..(88,-57.7)
				..controls(87.2,-57.5)and(86.1,-57.1)..(85.5,-56.8)
				..controls(85,-56.6)and(84.1,-56.2)..(83.5,-56)..controls(82.9,-55.7)and(81.8,-55.2)..(81,-54.9)
				..controls(80.1,-54.5)and(79.1,-54.2)..(78.7,-54.1)
				..controls(77.9,-53.8)and(77.3,-53.6)..(76.2,-52.9)
				..controls(75.7,-52.6)and(75.2,-52.4)..(74.1,-52.2)
				..controls(72.5,-51.9)and(72.4,-51.8)..(71.3,-51.1)
				..controls(69.9,-50.1)and(69.4,-49.9)..(67.9,-49.8)
				..controls(67,-49.7)and(66.6,-49.6)..(66.3,-49.4)
				..controls(66.1,-49.3)and(65.7,-49)..(65.4,-48.8)
				..controls(65.1,-48.7)and(64.2,-48.4)..(63.6,-48.2)
				..controls(62.9,-48)and(61.7,-47.6)..(60.9,-47.2)
				..controls(60.1,-46.8)and(58.9,-46.4)..(58.2,-46.2)
				..controls(57.6,-46)and(56.8,-45.7)..(56.5,-45.5)
				..controls(54.9,-44.6)and(54.6,-44.5)..(53.2,-44.2)
				..controls(51.7,-43.8)and(51,-43.6)..(49.9,-42.9)
				..controls(49.5,-42.6)and(48.9,-42.4)..(47.5,-42)
				..controls(45.7,-41.6)and(45.6,-41.5)..(44.3,-40.6)
				..controls(43.5,-40.1)and(42.8,-39.6)..(42.7,-39.4)
				..controls(42.4,-38.7)and(43,-38.1)..(44.5,-37.6)..controls(44.9,-37.5)and(46,-37)..(46.9,-36.6)
				..controls(47.9,-36.1)and(48.8,-35.8)..(49.6,-35.6)
				..controls(50.5,-35.4)and(51.3,-35.2)..(52.4,-34.6)
				..controls(53.3,-34.1)and(54.5,-33.7)..(54.9,-33.5)
				..controls(55.4,-33.4)and(56.4,-33)..(57.2,-32.6)
				..controls(58.6,-32)and(60.1,-31.5)..(61.6,-31.2)
				..controls(62.3,-31.1)and(62.7,-31)..(63.2,-30.7)
				..controls(64.1,-30)and(64.3,-29.9)..(65.3,-29.7)
				..controls(65.7,-29.6)and(66.5,-29.2)..(67.1,-28.9)
				..controls(67.6,-28.6)and(68.8,-28.2)..(69.6,-27.9)
				..controls(70.5,-27.7)and(71.6,-27.3)..(72.1,-27)
				..controls(72.6,-26.7)and(73.8,-26.2)..(74.8,-25.9)
				..controls(75.7,-25.5)and(76.7,-25.1)..(77.1,-25)
				..controls(77.4,-24.8)and(78.5,-24.4)..(79.6,-24.1)
				..controls(80.6,-23.8)and(82,-23.2)..(82.8,-22.9)
				..controls(83.5,-22.5)and(84.8,-22)..(85.8,-21.8)
				..controls(86.7,-21.5)and(87.8,-21.1)..(88.1,-20.9)
				..controls(88.4,-20.7)and(89.3,-20.3)..(90.1,-20.1)
				..controls(90.8,-19.8)and(91.9,-19.4)..(92.4,-19.1)
				..controls(93.5,-18.5)and(94.4,-18.2)..(95.4,-18.1)
				..controls(95.7,-18)and(96.4,-17.8)..(96.9,-17.5)
				..controls(97.5,-17.3)and(98.4,-16.9)..(99,-16.6)
				..controls(99.7,-16.3)and(100.9,-15.9)..(101.9,-15.6)
				..controls(102.9,-15.3)and(104.4,-14.7)..(105.3,-14.3)
				..controls(108.1,-13.2)and(108.5,-13.1)..(110.1,-13.1)
				..controls(111.4,-13.1)and(111.5,-13.1)..(112.4,-13.5)
				..controls(113.2,-14)and(115,-14.8)..(116.6,-15.4)
				..controls(116.9,-15.5)and(117.9,-15.8)..(118.8,-16.1)
				..controls(120.9,-16.6)and(121.5,-16.9)..(122.6,-17.9)
				..controls(123,-18.3)and(123.7,-18.8)..(123.9,-18.9)
				..controls(124.6,-19.4)and(126,-19.8)..(127.9,-20.2)
				..controls(129.6,-20.5)and(129.8,-20.6)..(131.1,-21.4)
				..controls(131.5,-21.7)and(132.4,-22.1)..(133.1,-22.3)
				..controls(136.1,-23.2)and(138.3,-24.1)..(140.2,-25)
				..controls(141,-25.5)and(141.5,-25.6)..(142.5,-25.8)
				..controls(145.1,-26.3)and(146,-26.6)..(149.4,-28.1)
				..controls(151.3,-28.9)and(153.5,-29.8)..(154.3,-30.1)
				..controls(155.1,-30.4)and(156.2,-30.8)..(156.6,-31)
				..controls(157.1,-31.2)and(158.2,-31.6)..(159.1,-31.8)
				..controls(160.6,-32.2)and(160.9,-32.3)..(161.7,-32.9)
				..controls(162.4,-33.3)and(162.9,-33.6)..(163.7,-33.8)
				..controls(164.3,-34)and(164.8,-34.1)..(164.9,-34.1)
				..controls(165,-34.1)and(165.6,-34.4)..(166.4,-34.8)
				..controls(167.3,-35.2)and(168.1,-35.4)..(169,-35.6)
				..controls(170,-35.8)and(170.5,-36)..(171.7,-36.6)
				..controls(172.4,-37)and(173.3,-37.4)..(173.6,-37.4)
				..controls(175.6,-37.7)and(176.2,-38)..(176.7,-39.1)
				..controls(176.9,-39.6)and(176.9,-39.7)..(176.7,-40.4)
				..controls(176.4,-41.6)and(175.8,-42.7)..(175.2,-43.4)
				..controls(174.5,-44.1)and(172.6,-45.5)..(171.7,-45.8)
				..controls(170.6,-46.2)and(169.1,-46)..(168.2,-45.4)
				..controls(167.9,-45.2)and(167.3,-44.8)..(166.9,-44.4)
				..controls(166.4,-44.1)and(165.8,-43.6)..(165.4,-43.2)
				..controls(165.1,-42.9)and(164.4,-42.3)..(163.9,-41.9)
				..controls(163.3,-41.5)and(162.7,-40.9)..(162.4,-40.5)
				..controls(161.4,-39.2)and(159.9,-38)..(158.8,-37.8)
				..controls(158.4,-37.8)and(158.3,-37.7)..(158.2,-37.2)
				..controls(158,-36.8)and(157.4,-36.1)..(157.1,-36.1)
				..controls(157,-36.1)and(156.7,-36.1)..(156.5,-36)
				..controls(156.1,-35.9)and(155.9,-35.7)..(155.7,-35.2)
				..controls(155.3,-34.6)and(155.1,-34.4)..(154,-34.1)
				..controls(153.6,-34)and(153.2,-33.7)..(152.9,-33.5)
				..controls(152.4,-32.9)and(151.6,-32.4)..(151.3,-32.4)
				..controls(150.9,-32.4)and(150.4,-32.6)..(150.4,-32.9)
				..controls(150.4,-33.2)and(150.8,-33.4)..(151.4,-33.6)
				..controls(152.1,-33.7)and(152.8,-34.1)..(153.5,-34.9)
				..controls(154.1,-35.5)and(154.3,-35.6)..(155,-35.8)
				..controls(155.7,-36)and(155.7,-36)..(156,-36.6)
				..controls(156.3,-37.3)and(156.7,-37.7)..(157.5,-37.8)
				..controls(158.1,-37.9)and(158.2,-38)..(158.4,-38.8)
				..controls(158.5,-39.4)and(158.6,-39.6)..(159.6,-40.4)
				..controls(160.2,-41)and(160.9,-41.6)..(161.3,-41.8)
				..controls(161.7,-42)and(161.9,-42.3)..(162.2,-42.8)
				..controls(162.4,-43.2)and(162.8,-43.8)..(163.1,-44.1)
				..controls(163.9,-44.8)and(165.1,-46.7)..(165.7,-48)
				..controls(166.2,-49.3)and(166.4,-49.7)..(167.2,-50.5)
				..controls(167.5,-50.8)and(167.9,-51.4)..(168.1,-51.8)
				..controls(168.3,-52.2)and(168.8,-53)..(169.1,-53.6)
				..controls(169.6,-54.3)and(169.8,-54.9)..(169.9,-55.6)
				..controls(170.4,-57.8)and(170.5,-58.1)..(171.2,-59.5)
				..controls(171.6,-60.3)and(172,-61.5)..(172.2,-62.1)
				..controls(172.3,-62.8)and(172.7,-63.7)..(172.9,-64.1)
				..controls(173.2,-64.6)and(173.4,-65.3)..(173.4,-65.6)
				..controls(173.8,-68.8)and(174.2,-70.6)..(174.9,-72.6)
				..controls(175.5,-74.5)and(175.9,-77.1)..(176.1,-80.2)
				..controls(176.3,-83.4)and(175.5,-84.4)..(173.7,-83.4)--cycle;
				\definecolor{ezs}{RGB}{247,179,26}\fill[ezs](-5.7,-156.6)--(-3.4,-157.5)--(61.9,12.7)--(59.7,13.6)--cycle;
				\definecolor{ezs}{RGB}{247,179,22}\fill[ezs](-3.8,-157.3)--(-1.5,-158.2)--(63.9,12)--(61.6,12.9)--cycle;
				\definecolor{ezs}{RGB}{248,179,19}\fill[ezs](-1.8,-158.1)--(0.4,-158.9)--(65.8,11.2)--(63.5,12.1)--cycle;
				\definecolor{ezs}{RGB}{248,178,16}\fill[ezs](0.1,-158.8)--(2.4,-159.7)--(67.7,10.5)--(65.4,11.4)--cycle;
				\definecolor{ezs}{RGB}{249,178,15}\fill[ezs](2,-159.5)--(4.3,-160.4)--(69.6,9.8)--(67.3,10.7)--cycle;
				\definecolor{ezs}{RGB}{251,177,17}\fill[ezs](3.9,-160.3)--(6.2,-161.1)--(71.5,9)--(69.2,9.9)--cycle;
				\definecolor{ezs}{RGB}{252,179,18}\fill[ezs](5.8,-161)--(8.1,-161.9)--(73.4,8.3)--(71.1,9.2)--cycle;
				\definecolor{ezs}{RGB}{252,182,17}\fill[ezs](7.7,-161.7)--(10,-162.6)--(75.3,7.6)--(73,8.5)--cycle;
				\definecolor{ezs}{RGB}{252,185,17}\fill[ezs](9.6,-162.5)--(11.9,-163.4)--(77.2,6.8)--(75,7.7)--cycle;
				\definecolor{ezs}{RGB}{251,189,18}\fill[ezs](11.5,-163.2)--(13.8,-164.1)--(79.2,6.1)--(76.9,7)--cycle;
				\definecolor{ezs}{RGB}{251,191,22}\fill[ezs](13.4,-163.9)--(15.7,-164.8)--(81.1,5.4)--(78.8,6.2)--cycle;
				\definecolor{ezs}{RGB}{252,193,25}\fill[ezs](15.4,-164.7)--(17.6,-165.6)--(83,4.6)--(80.7,5.5)--cycle;
				\definecolor{ezs}{RGB}{253,195,27}\fill[ezs](17.3,-165.4)--(19.6,-166.3)--(84.9,3.9)--(82.6,4.8)--cycle;
				\definecolor{ezs}{RGB}{253,196,31}\fill[ezs](19.2,-166.1)--(21.5,-167)--(86.8,3.2)--(84.5,4)--cycle;
				\definecolor{ezs}{RGB}{254,199,35}\fill[ezs](21.1,-166.9)--(23.4,-167.8)--(88.7,2.4)--(86.4,3.3)--cycle;
				\definecolor{ezs}{RGB}{253,201,38}\fill[ezs](23,-167.6)--(25.3,-168.5)--(90.6,1.7)--(88.3,2.6)--cycle;
				\definecolor{ezs}{RGB}{253,203,40}\fill[ezs](24.9,-168.3)--(27.2,-169.2)--(92.5,1)--(90.3,1.8)--cycle;
				\definecolor{ezs}{RGB}{253,205,42}\fill[ezs](26.8,-169.1)--(29.1,-170)--(94.4,0.2)--(92.2,1.1)--cycle;
				\definecolor{ezs}{RGB}{253,207,45}\fill[ezs](28.7,-169.8)--(31,-170.7)--(96.4,-0.5)--(94.1,0.4)--cycle;
				\definecolor{ezs}{RGB}{253,210,49}\fill[ezs](30.7,-170.5)--(32.9,-171.4)--(98.3,-1.2)--(96,-0.4)--cycle;
				\definecolor{ezs}{RGB}{253,212,54}\fill[ezs](32.6,-171.3)--(34.9,-172.2)--(100.2,-2)--(97.9,-1.1)--cycle;
				\definecolor{ezs}{RGB}{253,214,59}\fill[ezs](34.5,-172)--(36.8,-172.9)--(102.1,-2.7)--(99.8,-1.8)--cycle;
				\definecolor{ezs}{RGB}{253,216,63}\fill[ezs](36.4,-172.7)--(38.7,-173.6)--(104,-3.4)--(101.7,-2.6)--cycle;
				\definecolor{ezs}{RGB}{253,218,65}\fill[ezs](38.3,-173.5)--(40.6,-174.4)--(105.9,-4.2)--(103.6,-3.3)--cycle;
				\definecolor{ezs}{RGB}{254,220,66}\fill[ezs](40.2,-174.2)--(42.5,-175.1)--(107.8,-4.9)--(105.5,-4)--cycle;
				\definecolor{ezs}{RGB}{254,222,68}\fill[ezs](42.1,-175)--(44.4,-175.8)--(109.7,-5.6)--(107.5,-4.8)--cycle;
				\definecolor{ezs}{RGB}{254,224,71}\fill[ezs](44,-175.7)--(46.3,-176.6)--(111.7,-6.4)--(109.4,-5.5)--cycle;
				\definecolor{ezs}{RGB}{254,226,75}\fill[ezs](46,-176.4)--(48.2,-177.3)--(113.6,-7.1)--(111.3,-6.2)--cycle;
				\definecolor{ezs}{RGB}{254,228,79}\fill[ezs](47.9,-177.2)--(50.1,-178)--(115.5,-7.8)--(113.2,-7)--cycle;
				\definecolor{ezs}{RGB}{254,230,82}\fill[ezs](49.8,-177.9)--(52.1,-178.8)--(117.4,-8.6)--(115.1,-7.7)--cycle;
				\definecolor{ezs}{RGB}{254,232,85}\fill[ezs](51.7,-178.6)--(54,-179.5)--(119.3,-9.3)--(117,-8.4)--cycle;
				\definecolor{ezs}{RGB}{254,234,88}\fill[ezs](53.6,-179.4)--(55.9,-180.2)--(121.2,-10)--(118.9,-9.2)--cycle;
				\definecolor{ezs}{RGB}{254,237,91}\fill[ezs](55.5,-180.1)--(57.8,-181)--(123.1,-10.8)--(120.8,-9.9)--cycle;
				\definecolor{ezs}{RGB}{254,239,94}\fill[ezs](57.4,-180.8)--(59.7,-181.7)--(125,-11.5)--(122.8,-10.6)--cycle;
				\definecolor{ezs}{RGB}{254,241,97}\fill[ezs](59.3,-181.6)--(61.6,-182.4)--(127,-12.2)--(124.7,-11.4)--cycle;
				\definecolor{ezs}{RGB}{254,243,101}\fill[ezs](61.2,-182.3)--(63.5,-183.2)--(128.9,-13)--(126.6,-12.1)--cycle;
				\definecolor{ezs}{RGB}{254,245,103}\fill[ezs](63.2,-183)--(65.4,-183.9)--(130.8,-13.7)--(128.5,-12.8)--cycle;
				\definecolor{ezs}{RGB}{254,245,105}\fill[ezs](65.1,-183.8)--(67.4,-184.6)--(132.7,-14.4)--(130.4,-13.6)--cycle;
				\definecolor{ezs}{RGB}{254,245,105}\fill[ezs](67,-184.5)--(69.3,-185.4)--(134.6,-15.2)--(132.3,-14.3)--cycle;
				\definecolor{ezs}{RGB}{254,244,103}\fill[ezs](68.9,-185.2)--(71.2,-186.1)--(136.5,-15.9)--(134.2,-15)--cycle;
				\definecolor{ezs}{RGB}{253,243,101}\fill[ezs](70.8,-186)--(73.1,-186.8)--(138.4,-16.6)--(136.1,-15.8)--cycle;
				\definecolor{ezs}{RGB}{253,241,98}\fill[ezs](72.7,-186.7)--(75,-187.6)--(140.3,-17.4)--(138,-16.5)--cycle;
				\definecolor{ezs}{RGB}{253,239,95}\fill[ezs](74.6,-187.4)--(76.9,-188.3)--(142.2,-18.1)--(140,-17.2)--cycle;
				\definecolor{ezs}{RGB}{253,236,92}\fill[ezs](76.5,-188.2)--(78.8,-189)--(144.2,-18.8)--(141.9,-18)--cycle;
				\definecolor{ezs}{RGB}{253,234,88}\fill[ezs](78.5,-188.9)--(80.7,-189.8)--(146.1,-19.6)--(143.8,-18.7)--cycle;
				\definecolor{ezs}{RGB}{253,232,85}\fill[ezs](80.4,-189.6)--(82.7,-190.5)--(148,-20.3)--(145.7,-19.4)--cycle;
				\definecolor{ezs}{RGB}{253,230,82}\fill[ezs](82.3,-190.4)--(84.6,-191.2)--(149.9,-21)--(147.6,-20.2)--cycle;
				\definecolor{ezs}{RGB}{253,228,78}\fill[ezs](84.2,-191.1)--(86.5,-192)--(151.8,-21.8)--(149.5,-20.9)--cycle;
				\definecolor{ezs}{RGB}{253,226,74}\fill[ezs](86.1,-191.8)--(88.4,-192.7)--(153.7,-22.5)--(151.4,-21.6)--cycle;
				\definecolor{ezs}{RGB}{253,224,70}\fill[ezs](88,-192.6)--(90.3,-193.4)--(155.6,-23.3)--(153.3,-22.4)--cycle;
				\definecolor{ezs}{RGB}{253,222,67}\fill[ezs](89.9,-193.3)--(92.2,-194.2)--(157.5,-24)--(155.3,-23.1)--cycle;
				\definecolor{ezs}{RGB}{253,220,65}\fill[ezs](91.8,-194)--(94.1,-194.9)--(159.5,-24.7)--(157.2,-23.8)--cycle;
				\definecolor{ezs}{RGB}{253,217,63}\fill[ezs](93.7,-194.8)--(96,-195.6)--(161.4,-25.5)--(159.1,-24.6)--cycle;
				\definecolor{ezs}{RGB}{253,215,60}\fill[ezs](95.7,-195.5)--(97.9,-196.4)--(163.3,-26.2)--(161,-25.3)--cycle;
				\definecolor{ezs}{RGB}{253,213,56}\fill[ezs](97.6,-196.2)--(99.9,-197.1)--(165.2,-26.9)--(162.9,-26)--cycle;
				\definecolor{ezs}{RGB}{252,211,53}\fill[ezs](99.5,-197)--(101.8,-197.8)--(167.1,-27.7)--(164.8,-26.8)--cycle;
				\definecolor{ezs}{RGB}{252,209,49}\fill[ezs](101.4,-197.7)--(103.7,-198.6)--(169,-28.4)--(166.7,-27.5)--cycle;
				\definecolor{ezs}{RGB}{252,207,45}\fill[ezs](103.3,-198.4)--(105.6,-199.3)--(170.9,-29.1)--(168.6,-28.2)--cycle;
				\definecolor{ezs}{RGB}{253,205,42}\fill[ezs](105.2,-199.2)--(107.5,-200)--(172.8,-29.9)--(170.6,-29)--cycle;
				\definecolor{ezs}{RGB}{253,202,40}\fill[ezs](107.1,-199.9)--(109.4,-200.8)--(174.7,-30.6)--(172.5,-29.7)--cycle;
				\definecolor{ezs}{RGB}{253,200,38}\fill[ezs](109,-200.6)--(111.3,-201.5)--(176.7,-31.3)--(174.4,-30.4)--cycle;
				\definecolor{ezs}{RGB}{253,198,35}\fill[ezs](111,-201.4)--(113.2,-202.2)--(178.6,-32.1)--(176.3,-31.2)--cycle;
				\definecolor{ezs}{RGB}{253,196,31}\fill[ezs](112.9,-202.1)--(115.2,-203)--(180.5,-32.8)--(178.2,-31.9)--cycle;
				\definecolor{ezs}{RGB}{253,194,28}\fill[ezs](114.8,-202.8)--(117.1,-203.7)--(182.4,-33.5)--(180.1,-32.6)--cycle;
				\definecolor{ezs}{RGB}{253,192,26}\fill[ezs](116.7,-203.6)--(119,-204.5)--(184.3,-34.3)--(182,-33.4)--cycle;
				\definecolor{ezs}{RGB}{253,189,24}\fill[ezs](118.6,-204.3)--(120.9,-205.2)--(186.2,-35)--(183.9,-34.1)--cycle;
				\definecolor{ezs}{RGB}{253,186,20}\fill[ezs](120.5,-205)--(122.8,-205.9)--(188.1,-35.7)--(185.8,-34.8)--cycle;
				\definecolor{ezs}{RGB}{253,184,17}\fill[ezs](122.4,-205.8)--(124.7,-206.7)--(190,-36.5)--(187.8,-35.6)--cycle;
				\definecolor{ezs}{RGB}{252,183,15}\fill[ezs](124.3,-206.5)--(126.6,-207.4)--(192,-37.2)--(189.7,-36.3)--cycle;
				\definecolor{ezs}{RGB}{252,181,14}\fill[ezs](126.3,-207.2)--(128.5,-208.1)--(193.9,-37.9)--(191.6,-37.1)--cycle;
				\definecolor{ezs}{RGB}{252,179,13}\fill[ezs](128.2,-208)--(130.4,-208.9)--(195.8,-38.7)--(193.5,-37.8)--cycle;
				\definecolor{ezs}{RGB}{252,178,12}\fill[ezs](130.1,-208.7)--(132.4,-209.6)--(197.7,-39.4)--(195.4,-38.5)--cycle;
				\definecolor{ezs}{RGB}{251,178,10}\fill[ezs](132,-209.4)--(134.3,-210.3)--(199.6,-40.1)--(197.3,-39.3)--cycle;
				\definecolor{ezs}{RGB}{251,178,10}\fill[ezs](133.9,-210.2)--(136.2,-211.1)--(201.5,-40.9)--(199.2,-40)--cycle;
				\definecolor{ezs}{RGB}{250,177,11}\fill[ezs](135.8,-210.9)--(138.1,-211.8)--(203.4,-41.6)--(201.1,-40.7)--cycle;
				\definecolor{ezs}{RGB}{250,178,11}\fill[ezs](137.7,-211.6)--(140,-212.5)--(205.3,-42.3)--(203.1,-41.5)--cycle;
				\definecolor{ezs}{RGB}{248,180,10}\fill[ezs](139.6,-212.4)--(141.9,-213.3)--(207.3,-43.1)--(205,-42.2)--cycle;
				\definecolor{ezs}{RGB}{247,181,9}\fill[ezs](141.5,-213.1)--(143.8,-214)--(209.2,-43.8)--(206.9,-42.9)--cycle;
				\definecolor{ezs}{RGB}{246,182,9}\fill[ezs](143.5,-213.8)--(145.7,-214.7)--(211.1,-44.5)--(208.8,-43.7)--cycle;
				\definecolor{ezs}{RGB}{246,182,10}\fill[ezs](145.4,-214.6)--(147.7,-215.5)--(213,-45.3)--(210.7,-44.4)--cycle;
			\end{scope}
			% chữ EDUZY (Times đậm, làm dày nét bằng 5 lớp lệch nhẹ; căn đúng khung chữ logo gốc)
			\foreach \dx/\dy in {-0.7/0,0.7/0,0/-0.3,0/0.3,0/0}
			\node[anchor=center,inner sep=0pt,outer sep=0pt,text=ezText,align=center] at (109.6+\dx,-193.6+\dy)
			{\resizebox{1.01cm}{0.22cm}{\fontencoding{T1}\fontfamily{ptm}\fontseries{b}\selectfont EDUZY}};
		\end{tikzpicture}%
	}%
}
%====================== THIẾT LẬP THÂN ĐỀ ======================
%--- Môi trường tự luận {bt}: ex_test 3.3.4 chưa khai báo sẵn nên tạo ở đây
%    (ex_test đã tự móc \fixshowans và \listenumerate cho bt)
\makeatletter
\@ifundefined{bt}{\newtheorem{bt}{\color{gold1}\namebt}}{}
\makeatother
\renewcommand{\sectionmark}[1]{\markright{\thesection\ #1}}
%--- Khung tiêu đề (logo + tối đa 4 dòng; dòng 3, 4 để trống {} thì bỏ qua)
%  Đề thi:    \khungtieude{THI THỬ THPT QUỐC GIA NĂM 2026}{Môn Toán -- Đề số 21}
%                         {(Đề thi gồm: ...)}{Thời gian làm bài: 90 phút (không kể thời gian phát đề).}
%  Chuyên đề: \khungtieude{CHUYÊN ĐỀ 3 -- DÃY SỐ VÀ CÔNG THỨC TRUY HỒI}
%                         {Bồi dưỡng học sinh giỏi Toán THPT}{...}{}
\newcommand{\khungtieude}[4]{%
	\begin{tcolorbox}[colframe=gold1,colback=gold5!45,boxrule=0.5mm]
		\begin{minipage}[c]{3cm}
			\centerline{\eduzylogo[2.4cm]}
		\end{minipage}\hspace{0.5cm}
		\begin{minipage}[c]{12.5cm}
			\vskip 0.15 cm
			\centerline{\ind{#1}}
			\vskip 0.15 cm
			\centerline{\textbf{#2}}
			\if\relax\detokenize{#3}\relax\else
				\vskip 0.15 cm
				\centerline{\textit{#3}}
			\fi
			\if\relax\detokenize{#4}\relax\else
				\vskip 0.15 cm
				\centerline{\textit{#4}}
			\fi
			\vskip 0.15 cm
		\end{minipage}
	\end{tcolorbox}
}
%====================== BỔ SUNG CHO BÀI HỌC, CHUYÊN ĐỀ ======================
\usepackage[shortlabels]{enumitem}
\usepackage{tabularx}
%--- Danh sách gọn: \begin{enumerate}[gon] ... \end{enumerate}
\SetEnumitemKey{gon}{leftmargin=*,itemsep=1pt,parsep=0pt,topsep=3pt}
%--- Khung lý thuyết: \begin{khung}{} ... \end{khung}; tiêu đề chỉ dùng khi thật ngắn ("Mẹo", "Lỗi thường gặp")
\newtcolorbox{khung}[1]{breakable,enhanced,colframe=gold1,colback=gold5!45,
	boxrule=0.4mm,fonttitle=\bfseries\sffamily,coltitle=white,colbacktitle=gold1,
	attach boxed title to top left={xshift=6mm,yshift=-3mm},
	boxed title style={arc=1mm,boxrule=0.3mm},title={#1}}
%--- Tiêu đề mục nhỏ trong một dạng: \muc{Dấu hiệu}, \muc{Phương pháp}, \muc{Bài tự luyện}
\makeatletter
\newcommand{\muc}[1]{\par\medskip\noindent\indam{#1}\par\nobreak\smallskip\@afterindentfalse\@afterheading}
\makeatother
%--- Ô công thức cần nhớ (căn giữa, viền vàng, vừa khít nội dung)
%     \begin{center}\congthuc{$Q(n)=(n-1)!$}\end{center}
\newtcbox{\congthuc}{on line,colframe=gold1,colback=gold5!45,boxrule=0.5pt,arc=2pt,
	left=8pt,right=8pt,top=4pt,bottom=4pt}
%--- Công tắc lời giải bài tự luyện (ví dụ mẫu luôn dùng \loigiai)
%     \loigiaiTL[{đáp số}]{lời giải đầy đủ}
%     \TLfull -> in lời giải đầy đủ;  \TLdapso -> chỉ in một dòng "Đáp số";  \TLan -> ẩn hoàn toàn
%     (\CheDo là dethi thì mọi lời giải và đáp số đều ẩn)
\makeatletter
\newif\ifTLhien   \TLhientrue
\newif\ifTLdapso  \TLdapsofalse
\newcommand{\TLfull}{\TLhientrue\TLdapsofalse}
\newcommand{\TLdapso}{\TLhientrue\TLdapsotrue}
\newcommand{\TLan}{\TLhienfalse}
\newcommand{\loigiaiTL}[2][]{%
	\ifTLhien
		\ifTLdapso
			\def\@tlarg{#1}%
			\ifx\@tlarg\@empty\else
				\ifdefstring{\kindOPTN}{dethi}{}{\par\noindent{\bfseries\color{gold1}Đáp số:} #1\par}%
			\fi
		\else
			\loigiai{#2}%
		\fi
	\fi}
\makeatother
\TLdapso% mặc định: chỉ in đáp số bài tự luyện
%====================== BẢNG ĐÚNG/SAI (\choiceTFt) ======================
%  Hàng tiêu đề nền nâu đồng chữ trắng (Phát biểu | Đ | S), các hàng tô xen kẽ kem / trắng,
%  đường kẻ vàng kim, nhãn a) b) c) d) màu nâu đồng. Bản có đáp án: cột được chọn có vòng tròn
%  nâu đồng với dấu tích trắng (câu đúng tích ở cột Đúng, câu sai tích ở cột Sai).
%  Tiêu đề ghi đủ "Đúng | Sai": \OPTN{viettat=0} sau \input (cột Đúng/Sai rộng ra, phát biểu hẹp lại).
%  Dùng bảng gốc của ex_test: \def\BangDSgoc{1} trước \input{preamble-eduzy-v2}.
\providecommand{\BangDSgoc}{0}
\ifnum\BangDSgoc=0
\usepackage{colortbl,pifont}
\colorlet{ezDSke}{gold2!60!gold1}% đường kẻ bảng (nét 0,4pt như ex_test để không đổi chiều cao bảng)
\colorlet{ezDSle}{gold5!60}% hàng a), c)
\colorlet{ezDSchan}{white}% hàng b), d)
%--- dấu chọn ở cột Đúng/Sai
%    (bề rộng 0, căn giữa ô: không làm đổi bề rộng cột của ex_test, không báo tràn lề)
\newcommand{\EZdauchon}{\makebox[0pt]{\tikz[baseline=-1ex]\node[circle,fill=gold1,inner sep=0pt,
	minimum size=1.2em,text=white,font=\small]{\ding{51}};}}
\def\tickT{\EZdauchon}\def\tickF{\EZdauchon}
\AtBeginDocument{\sbox{\ccTFTab}{\EZdauchon}}
%--- nhãn a) b) c) d) màu nâu đồng (ở bản lời giải, nhãn câu đúng vẫn đóng khung như ex_test)
\renewcommand{\FalseTF}{\stepcounter{dapan}\textbf{\color{gold1}\DapAnTF\sepTF}\ignorespaces}
%--- trước khi dựng bảng: kẻ màu vàng kim (bề rộng cột giữ nguyên như ex_test); sau bảng trả lại màu đen
\newcount\EZhangTF
\newcommand{\EZbatdauTF}{%
	\global\EZhangTF=0\relax
	\arrayrulecolor{ezDSke}}
\newcommand{\EZketthucTF}{\arrayrulecolor{black}}
%--- hàng tiêu đề nền nâu đồng, chữ trắng
\makeatletter
\newcommand{\EZo@TF}[1]{\centering\bfseries\color{white}#1}
\newcommand{\EZtieudeTF}{%
	\ifcase\cotDS
		\ifnum\cotSTT=1\relax
			\def\rowHeader{\hline\rowcolor{gold1}\EZo@TF{TT}&\EZo@TF{\phatbieu}\tabularnewline}%
		\else
			\def\rowHeader{\hline\rowcolor{gold1}\EZo@TF{\phatbieu}\tabularnewline}%
		\fi
	\or
		\ifnum\cotSTT=1\relax
			\def\rowHeader{\hline\rowcolor{gold1}\EZo@TF{TT}&\EZo@TF{\phatbieu}&\EZo@TF{\Dung}\tabularnewline}%
		\else
			\def\rowHeader{\hline\rowcolor{gold1}\EZo@TF{\phatbieu}&\EZo@TF{\Dung}\tabularnewline}%
		\fi
	\or
		\ifnum\cotSTT=1\relax
			\def\rowHeader{\hline\rowcolor{gold1}\EZo@TF{TT}&\EZo@TF{\phatbieu}&\EZo@TF{\Dung}&\EZo@TF{\Sai}\tabularnewline}%
		\else
			\def\rowHeader{\hline\rowcolor{gold1}\EZo@TF{\phatbieu}&\EZo@TF{\Dung}&\EZo@TF{\Sai}\tabularnewline}%
		\fi
	\fi}
%--- tô xen kẽ từng hàng (gọi ngay sau \hline đầu mỗi hàng)
\newcommand{\EZmauhangTF}{%
	\noalign{\global\advance\EZhangTF\@ne}%
	\ifodd\EZhangTF\rowcolor{ezDSle}\else\rowcolor{ezDSchan}\fi}
%--- gắn vào \choiceTFt / \writeChoiceTFt của ex_test
\xpretocmd{\choiceTFt}{\EZbatdauTF}{}{\typeout{[preamble-eduzy-v2] Không gắn được bảng Đúng/Sai (đầu bảng)}}
\xpatchcmd{\choiceTFt}{\ifnum\dongPB=0\def\rowHeader{}\fi}{\EZtieudeTF\ifnum\dongPB=0\def\rowHeader{}\fi}
	{}{\typeout{[preamble-eduzy-v2] Không gắn được bảng Đúng/Sai (hàng tiêu đề)}}
\xpatchcmd{\writeChoiceTFt}{\hline\xdef\dotEX{\sdotEX}}{\hline\EZmauhangTF\xdef\dotEX{\sdotEX}}
	{}{\typeout{[preamble-eduzy-v2] Không gắn được bảng Đúng/Sai (tô hàng)}}
\xpatchcmd{\writeChoiceTFt}{{\end{longtable}}}{{\end{longtable}\EZketthucTF}}
	{}{\typeout{[preamble-eduzy-v2] Không gắn được bảng Đúng/Sai (cuối longtable)}}
\xpatchcmd{\writeChoiceTFt}{{\end{xtabular}\end{center}\medskip\noindent}}{{\end{xtabular}\end{center}\EZketthucTF\medskip\noindent}}
	{}{\typeout{[preamble-eduzy-v2] Không gắn được bảng Đúng/Sai (cuối xtabular)}}
\fi
%=====================================================================
\renewcommand{\parallel}{\mathrel{/\!\!/}}

%  trong file đề (tools/dan_trang/dan_trang.py tự ghi khi đặt \NeedspaceCau):
%    \def\TachCauDS{1}   câu Đúng/Sai được tách: câu dẫn (kèm hình) ở cuối trang, bảng Phát biểu sang trang sau
%                        (câu dẫn và hình không tách nhau, bảng không tách giữa hai trang)
%    \def\TachCauLG{1}   bản lời giải (\CheDo khác dethi): phần đề của câu cao hơn 1/3 trang được phép tách
%  Lệnh dùng trong file đề (giá trị do tools/dan_trang/dan_trang.py đo ở bản dethi):
%    \NeedspaceCau[<cao câu dẫn + hình>]{<cao cả khối câu>}   ngay trên \begin{ex} (câu đầu phần: trên tiêu đề phần)
%        công tắc tắt: đúng bằng \Needspace{<cao cả khối câu>} như cũ;
%        \TachCauDS=1 và có [..] (câu Đúng/Sai được chọn tách): chỉ giữ câu dẫn + hình;
%        \TachCauLG=1 ở bản lời giải: đòi tối đa 1/3 \textheight.
%    \NeedspaceBang{<cao bảng Phát biểu>}%   ngay trước \choiceTFt: \TachCauDS=1 thì giữ trọn bảng trên một trang;
%        công tắc tắt: không làm gì.
\providecommand{\TachCauDS}{0}
\providecommand{\TachCauLG}{0}
\newlength{\EZnsCau}
\newcommand{\NeedspaceCau}[2][]{%
	\setlength{\EZnsCau}{#2}%
	\ifnum\TachCauDS=1 \if\relax\detokenize{#1}\relax\else\setlength{\EZnsCau}{#1}\fi\fi
	\ifnum\TachCauLG=1 \ifdefstring{\kindOPTN}{dethi}{}{\ifdim\EZnsCau>0.3333\textheight\setlength{\EZnsCau}{0.3333\textheight}\fi}\fi
	\Needspace{\EZnsCau}}
\newcommand{\NeedspaceBang}[1]{\ifnum\TachCauDS=1 \Needspace{#1}\fi}

%--- Thang màu: đậm -> nhạt
\definecolor{gold1}{RGB}{126,96,0}    % nâu đồng đậm  #7E6000
\definecolor{gold2}{RGB}{192,144,3}   % vàng kim      #C09003
\definecolor{gold3}{RGB}{255,215,99}  % vàng nhạt     #FFD763
\definecolor{gold4}{RGB}{255,231,153} % vàng rất nhạt #FFE799
\definecolor{gold5}{RGB}{255,243,205} % vàng kem      #FFF3CD
%--- Ghi đè tông xanh/đỏ cũ
\definecolor{myblue}{RGB}{126,96,0}
\colorlet{blue}{gold1}
\colorlet{cyan}{gold2}
\colorlet{teal}{gold2}

%--- Tên chương mặc định (file đề nên tự khai báo lại)
\providecommand{\tenchuong}{}
%--- Chân trang
\providecommand{\thuonghieu}{TOÁN 12 -- Eduzy.com.vn}
\providecommand{\tacgia}{Tô Gia Bảo -- 0353703783}

%--- Chữ cạnh logo ở đầu trang
\providecommand{\khauhieu}{Eduzy.com.vn}

%====================== THÔNG SỐ ĐẦU TRANG, CHÂN TRANG (v2) ======================
%  Dải nền đậm tràn ngang trang; ô chữ vàng kem có cạnh xiên, kèm 2 vạch xiên chuyển màu.
%  Đầu trang lẻ : logo + \khauhieu (trái)   | ô chữ \tenchuong (phải)
%  Chân trang lẻ: ô chữ \thuonghieu (trái)  | Trang (số) (phải)
%  Trang chẵn soi gương trang lẻ; chân trang chẵn ghi \tacgia.
%--- Màu: chỉ dùng thang vàng kim Eduzy
\colorlet{ezNen}{gold1}% dải nền đậm
\colorlet{ezVach1}{gold3}% vạch xiên sát dải nền (đậm hơn)
\colorlet{ezVach2}{gold4}% vạch xiên sát ô chữ (nhạt hơn)
\colorlet{ezO}{gold5}% ô chữ vàng kem
%--- Kích thước
\newlength{\EZdai}\setlength{\EZdai}{1.1cm}% chiều cao dải nền
\newlength{\EZo}\setlength{\EZo}{0.68cm}% chiều cao ô chữ
\newlength{\EZxien}\setlength{\EZxien}{0.55cm}% độ lệch ngang của cạnh xiên
\newlength{\EZvach}\setlength{\EZvach}{0.75cm}% bề rộng mỗi vạch xiên
\newlength{\EZkhe}\setlength{\EZkhe}{0.12cm}% khe trắng giữa vạch và ô chữ
\newlength{\EZdem}\setlength{\EZdem}{0.4cm}% khoảng đệm từ chữ tới cạnh xiên
\newlength{\EZle}\setlength{\EZle}{1.8cm}% lề trái/phải của chữ (bằng lề thân bài)
\newlength{\EZrongmin}\setlength{\EZrongmin}{4.5cm}% ô chữ rộng tối thiểu
\newlength{\EZrong}% bề rộng chữ trong ô (tính tự động)
%--- Font
\newcommand{\EZfontO}{\fontsize{10.5pt}{12pt}\selectfont\bfseries}% chữ trong ô
\newcommand{\EZfontKH}{\fontfamily{qag}\fontsize{11pt}{12pt}\selectfont\bfseries}% \khauhieu cạnh logo
\newcommand{\EZfontTrang}{\fontsize{11.5pt}{12pt}\selectfont\bfseries}% chữ "Trang" và số trang
%--- Đo bề rộng chữ trong ô
\newcommand{\EZdorong}[1]{%
	\settowidth{\EZrong}{\EZfontO #1}%
	\ifdim\EZrong<\EZrongmin\setlength{\EZrong}{\EZrongmin}\fi}
%--- Logo + chữ trắng \khauhieu (dựng sẵn một lần khi vào \begin{document})
\newsavebox{\EZlogobox}
\AtBeginDocument{\sbox{\EZlogobox}{%
		\raisebox{-0.5\height}{\eduzylogo[1cm]}\hspace{0.15cm}%
		\raisebox{-0.5\height}{\color{white}\EZfontKH\khauhieu}}}
%--- Hai vạch xiên + khe trắng. #1: điểm gốc ở mép dưới, #2: hướng (1 sang phải, -1 sang trái),
%    #3: hướng nghiêng (1 kiểu "/", -1 kiểu "\")
\newcommand{\EZvachxien}[3]{%
	\fill[white] (#1) --++({(#3)*\EZxien},\EZo) --++({(#2)*\EZkhe},0) --++({-(#3)*\EZxien},-\EZo) --cycle;
	\fill[ezVach2] ([xshift=(#2)*\EZkhe]#1) --++({(#3)*\EZxien},\EZo) --++({(#2)*\EZvach},0) --++({-(#3)*\EZxien},-\EZo) --cycle;
	\fill[ezVach1] ([xshift=(#2)*(\EZkhe+\EZvach)]#1) --++({(#3)*\EZxien},\EZo) --++({(#2)*\EZvach},0) --++({-(#3)*\EZxien},-\EZo) --cycle;
	\draw[white,line width=0.6pt] ([xshift=(#2)*(\EZkhe+\EZvach)]#1) --++({(#3)*\EZxien},\EZo);
}
%--- Chữ trong ô (căn giữa theo chiều cao ô, không lệch vì dấu tiếng Việt)
\tikzset{ezchu/.style={inner sep=0pt,text=ezNen,text height=0.72em,text depth=0.22em}}

%====================== LOGO CHÌM GIỮA TRANG ======================
%  Vẽ cùng đầu trang nên nằm dưới chữ của thân bài, có trên mọi trang.
%  Tắt: \logochimtat (từ trang đó trở đi), bật lại: \logochimbat.
%  Đổi cỡ, độ đậm: khai báo trước \input, ví dụ \def\EZchimcao{9cm} \def\EZchimdo{0.1}
\providecommand{\EZchimcao}{10cm}% chiều cao logo chìm
\providecommand{\EZchimdo}{0.12}% độ đậm (0 = trong suốt, 1 = rõ hẳn)
\newif\ifEZchim \EZchimtrue
\newcommand{\logochimtat}{\EZchimfalse}
\newcommand{\logochimbat}{\EZchimtrue}
%  Dựng sẵn logo mờ một lần khi vào \begin{document}; transparency group làm mờ cả logo
%  một lượt nên không lộ vệt ở chỗ các dải màu chồng nhau.
\newsavebox{\EZchimtho}
\newsavebox{\EZchimbox}
\AtBeginDocument{%
	\sbox{\EZchimtho}{\eduzylogo[\EZchimcao]}%
	\sbox{\EZchimbox}{\begin{tikzpicture}
			\begin{scope}[transparency group,opacity=\EZchimdo]
				\node[inner sep=0pt] {\usebox{\EZchimtho}};
			\end{scope}
	\end{tikzpicture}}}
\newcommand{\EZvechim}{%
	\ifEZchim
		\node[inner sep=0pt] at (current page.center) {\usebox{\EZchimbox}};
	\fi}

%====================== ĐẦU TRANG - TRANG LẺ ======================
%  (logo bên trái, ô chữ \tenchuong bên phải)
\newcommand{\myfancyheadOdd}{%
	\EZdorong{\tenchuong}%
	\begin{tikzpicture}[remember picture,overlay]
		\coordinate (T) at (current page.north west);
		\coordinate (P) at (current page.north east);
		%--- logo chìm giữa trang
		\EZvechim
		%--- dải nền
		\fill[ezNen] (T) rectangle ([yshift=-\EZdai]P);
		%--- ô chữ: X là góc dưới bên trái, cạnh trái xiên kiểu "/"
		\coordinate (X) at ([xshift=-\EZle-\EZrong-\EZdem-\EZxien,yshift=-\EZdai]P);
		\fill[ezO] (X) --++(\EZxien,\EZo) -- ([yshift=-\EZdai+\EZo]P) -- ([yshift=-\EZdai]P) --cycle;
		\EZvachxien{X}{-1}{1}
		\node[ezchu,anchor=east] at ([xshift=-\EZle,yshift=-\EZdai+0.5*\EZo]P) {\EZfontO\tenchuong};
		%--- logo + tên
		\node[inner sep=0pt,anchor=west] at ([xshift=\EZle,yshift=-0.5*\EZdai]T) {\usebox{\EZlogobox}};
	\end{tikzpicture}%
}

%====================== ĐẦU TRANG - TRANG CHẴN =====================
%  (soi gương: ô chữ \tenchuong bên trái, logo bên phải)
\newcommand{\myfancyheadEven}{%
	\EZdorong{\tenchuong}%
	\begin{tikzpicture}[remember picture,overlay]
		\coordinate (T) at (current page.north west);
		\coordinate (P) at (current page.north east);
		%--- logo chìm giữa trang
		\EZvechim
		%--- dải nền
		\fill[ezNen] (T) rectangle ([yshift=-\EZdai]P);
		%--- ô chữ: X là góc dưới bên phải, cạnh phải xiên kiểu "\"
		\coordinate (X) at ([xshift=\EZle+\EZrong+\EZdem+\EZxien,yshift=-\EZdai]T);
		\fill[ezO] (X) --++(-\EZxien,\EZo) -- ([yshift=-\EZdai+\EZo]T) -- ([yshift=-\EZdai]T) --cycle;
		\EZvachxien{X}{1}{-1}
		\node[ezchu,anchor=west] at ([xshift=\EZle,yshift=-\EZdai+0.5*\EZo]T) {\EZfontO\tenchuong};
		%--- logo + tên
		\node[inner sep=0pt,anchor=east] at ([xshift=-\EZle,yshift=-0.5*\EZdai]P) {\usebox{\EZlogobox}};
	\end{tikzpicture}%
}

%====================== CHÂN TRANG - TRANG LẺ ======================
%  (ô chữ \thuonghieu bên trái, "Trang (số)" bên phải)
\newcommand{\myfancyfootOdd}{%
	\EZdorong{\thuonghieu}%
	\begin{tikzpicture}[remember picture,overlay]
		\coordinate (T) at (current page.south west);
		\coordinate (P) at (current page.south east);
		%--- dải nền
		\fill[ezNen] (T) rectangle ([yshift=\EZdai]P);
		%--- ô chữ sát mép trên dải: X là góc dưới bên phải, cạnh phải xiên kiểu "/"
		\coordinate (X) at ([xshift=\EZle+\EZrong+\EZdem,yshift=\EZdai-\EZo]T);
		\fill[ezO] ([yshift=\EZdai-\EZo]T) -- (X) --++(\EZxien,\EZo) -- ([yshift=\EZdai]T) --cycle;
		\EZvachxien{X}{1}{1}
		\node[ezchu,anchor=west] at ([xshift=\EZle,yshift=\EZdai-0.5*\EZo]T) {\EZfontO\thuonghieu};
		%--- số trang
		\node[circle,draw=white,line width=0.9pt,text=white,inner sep=0pt,minimum size=0.62cm,anchor=east]
		(so) at ([xshift=-\EZle,yshift=0.5*\EZdai]P) {\EZfontTrang\thepage};
		\node[inner sep=0pt,text=white,anchor=base east] (tr) at ([xshift=-0.15cm]so.base west) {\EZfontTrang Trang};
	\end{tikzpicture}%
}

%====================== CHÂN TRANG - TRANG CHẴN ====================
%  (soi gương: "Trang (số)" bên trái, ô chữ \tacgia bên phải)
\newcommand{\myfancyfootEven}{%
	\EZdorong{\tacgia}%
	\begin{tikzpicture}[remember picture,overlay]
		\coordinate (T) at (current page.south west);
		\coordinate (P) at (current page.south east);
		%--- dải nền
		\fill[ezNen] (T) rectangle ([yshift=\EZdai]P);
		%--- ô chữ sát mép trên dải: X là góc dưới bên trái, cạnh trái xiên kiểu "\"
		\coordinate (X) at ([xshift=-\EZle-\EZrong-\EZdem,yshift=\EZdai-\EZo]P);
		\fill[ezO] ([yshift=\EZdai-\EZo]P) -- (X) --++(-\EZxien,\EZo) -- ([yshift=\EZdai]P) --cycle;
		\EZvachxien{X}{-1}{-1}
		\node[ezchu,anchor=east] at ([xshift=-\EZle,yshift=\EZdai-0.5*\EZo]P) {\EZfontO\tacgia};
		%--- số trang
		%    (vòng tròn ẩn so0 nằm giữa dải để lấy đường chân chữ, giống trang lẻ)
		\node[circle,inner sep=0pt,minimum size=0.62cm] (so0) at ([yshift=0.5*\EZdai]T) {\EZfontTrang\phantom{\thepage}};
		\node[inner sep=0pt,text=white,anchor=base west] (tr) at ($(so0.base -| T)+(\EZle,0)$) {\EZfontTrang Trang};
		\node[circle,draw=white,line width=0.9pt,text=white,inner sep=0pt,minimum size=0.62cm,anchor=base west]
		(so) at ([xshift=0.15cm]tr.base east) {\EZfontTrang\thepage};
	\end{tikzpicture}%
}

%--- Giữ tương thích với các file cũ đang gọi \myfancyhead, \myfancyfoot
\newcommand{\myfancyhead}{\myfancyheadOdd}
\newcommand{\myfancyfoot}{\myfancyfootOdd}

%======================= GẮN VÀO TRANG =============================
\pagestyle{fancy}
\fancyhf{}
\fancyhead[LO]{\myfancyheadOdd}
\fancyhead[RE]{\myfancyheadEven}
\fancyfoot[LO]{\myfancyfootOdd}
\fancyfoot[RE]{\myfancyfootEven}
\renewcommand{\headrulewidth}{0pt}
\renewcommand{\footrulewidth}{0pt}
\fancypagestyle{plain}{%
	\fancyhf{}
	\fancyhead[LO]{\myfancyheadOdd}
	\fancyhead[RE]{\myfancyheadEven}
	\fancyfoot[LO]{\myfancyfootOdd}
	\fancyfoot[RE]{\myfancyfootEven}
	\renewcommand{\headrulewidth}{0pt}
	\renewcommand{\footrulewidth}{0pt}
}
%=================== HẾT PHẦN GỘP TỪ eduzy-de.tex ===================
\newcommand{\indam}[1]{
	{\fontsize{12pt}{9pt}\sffamily \selectfont\bfseries\color{gold1} #1}
}
\newcommand{\ind}[1]{
	{\fontsize{14pt}{9pt}\fontfamily{qag} \selectfont\bfseries\color{gold1} #1}
}
\renewcommand{\FalseEX}{\stepcounter{dapan}\textbf{\selectfont\color{gold1}\Alph{dapan}.}\ignorespaces}
\renewtheorem{ex}{\color{gold1}Câu}
\renewcommand{\baselinestretch}{1.1}
\newcommand{\hoac}[1]{
	\left[\begin{aligned}#1\end{aligned}\right.}
\newcommand{\heva}[1]{
	\left\{\begin{aligned}#1\end{aligned}\right.}
\usepackage{esvect}
\def\vec{\vv}
\def\vt{\vv}
\def\overrightarrow{\vv}
%====================== LOGO EDUZY VẼ BẰNG TIKZ ======================
%  Đường viền dò từ logo gốc; dải màu vàng kim đo trực tiếp trên ảnh gốc
%  (nghiêng 21 độ, tô bằng 80 dải hẹp nên hiển thị giống nhau trên mọi trình xem PDF).
%  Dùng: \eduzylogo          -> cao 2.4cm (mặc định)
%        \eduzylogo[1.5cm]   -> cao 1.5cm
\definecolor{ezText}{RGB}{245,200,46}
\newcommand{\eduzylogo}[1][2.4cm]{%
	\resizebox{!}{#1}{%
		\begin{tikzpicture}[x=0.01cm,y=0.01cm]
			\path[use as bounding box] (0,0) rectangle (224,-224);
			% mũ tốt nghiệp + vương miện
			\begin{scope}
				\clip
				(119.4,-176.1)..controls(118.9,-176.1)and(116.8,-176)..(114.6,-176)
				..controls(101.4,-176)and(85.8,-175.9)..(83.7,-175.7)
				..controls(81.8,-175.6)and(81,-175.6)..(80.2,-175.7)
				..controls(78.8,-175.9)and(77.7,-175.9)..(76.2,-175.7)
				..controls(74.9,-175.5)and(70.8,-175.5)..(64.3,-175.7)
				..controls(59.1,-175.9)and(58.3,-175.7)..(58.6,-173.7)
				..controls(58.6,-172.8)and(59.6,-169.6)..(60.3,-168)
				..controls(60.8,-166.9)and(61.4,-165.8)..(62.3,-165)
				..controls(63.1,-164.2)and(64.3,-163.4)..(65,-163.2)
				..controls(65.5,-163.1)and(72.3,-163)..(76.6,-163.1)
				..controls(77.9,-163.1)and(80.1,-163.1)..(81.4,-163)
				..controls(82.8,-162.9)and(86.3,-162.9)..(89.2,-162.9)
				..controls(93,-162.9)and(94.9,-162.8)..(96,-162.7)
				..controls(97.2,-162.5)and(101,-162.5)..(114,-162.4)
				..controls(123,-162.3)and(134.9,-162.2)..(140.2,-162.1)
				..controls(150.1,-162)and(152.2,-162.1)..(153.6,-162.4)
				..controls(154.4,-162.6)and(157.1,-163.8)..(158.1,-164.4)
				..controls(158.8,-164.8)and(159.7,-165.7)..(159.9,-166.3)
				..controls(160,-166.5)and(160.2,-167.1)..(160.2,-167.6)
				..controls(160.3,-168)and(160.6,-169.4)..(160.9,-170.6)
				..controls(161.2,-171.7)and(161.5,-173.1)..(161.6,-173.6)
				..controls(161.8,-174.4)and(161.8,-174.5)..(161.5,-174.9)
				..controls(161.2,-175.6)and(160.8,-175.9)..(159.9,-176)
				..controls(159.2,-176.1)and(154.1,-176.2)..(151.5,-176.2)
				..controls(148.1,-176.1)and(122.6,-176.1)..(121.6,-176.2)
				..controls(120.8,-176.2)and(119.8,-176.2)..(119.4,-176.1)--cycle
				(130.4,-160)..controls(129,-160)and(119.8,-160)..(110.1,-159.9)
				..controls(100.3,-159.8)and(88.5,-159.7)..(83.9,-159.6)
				..controls(77.8,-159.5)and(75.3,-159.5)..(74.5,-159.6)
				..controls(73.9,-159.7)and(72.5,-159.8)..(71.3,-159.8)
				..controls(69.6,-159.8)and(69,-159.8)..(68.5,-159.6)
				..controls(67.6,-159.3)and(66.8,-158.6)..(66.2,-157.6)
				..controls(66,-157.3)and(65.6,-156.5)..(65.3,-156.1)
				..controls(64.6,-155.1)and(64.5,-154.9)..(63.7,-152.9)
				..controls(63.2,-151.6)and(62.9,-150.9)..(61.8,-149.6)
				..controls(61.5,-149.2)and(61.1,-148.5)..(60.9,-147.9)
				..controls(60.7,-147.4)and(60.3,-146.6)..(60.1,-146.1)
				..controls(59.9,-145.5)and(59.5,-144.9)..(59,-144.3)
				..controls(58.5,-143.7)and(58.2,-143.2)..(57.9,-142.5)
				..controls(57.7,-141.9)and(57.3,-141.1)..(57,-140.7)
				..controls(56.6,-140.3)and(56.1,-139.3)..(55.8,-138.5)
				..controls(55.5,-137.8)and(55.1,-136.8)..(54.9,-136.5)
				..controls(54.7,-136.1)and(54.3,-135.2)..(53.9,-134.5)
				..controls(53.6,-133.8)and(53,-132.7)..(52.6,-132)
				..controls(51.6,-130.5)and(51.4,-130.2)..(51.1,-129)
				..controls(50.6,-127.5)and(50.6,-127.4)..(49.7,-126.3)
				..controls(49.2,-125.7)and(48.6,-124.8)..(48.3,-124.2)
				..controls(48,-123.6)and(47.5,-122.7)..(47.3,-122.2)
				..controls(47,-121.8)and(46.5,-120.7)..(46.1,-119.8)
				..controls(45.8,-118.9)and(45.2,-117.8)..(44.9,-117.2)
				..controls(44.2,-116.1)and(43.3,-114.4)..(42.3,-112.2)
				..controls(41.8,-111.3)and(41.3,-110.4)..(41.1,-110.1)
				..controls(40.5,-109.4)and(40.3,-109)..(40,-107.8)
				..controls(39.8,-107)and(39.6,-106.5)..(39.2,-106)
				..controls(38.9,-105.6)and(38.4,-104.7)..(38,-104)
				..controls(37.7,-103.4)and(37.2,-102.5)..(37,-102.1)
				..controls(36.7,-101.7)and(36.3,-100.8)..(36,-100.1)
				..controls(35.6,-98.9)and(35.3,-98.4)..(34.6,-97.6)
				..controls(34.2,-97.1)and(34.2,-96.3)..(34.6,-96)
				..controls(35,-95.7)and(35.5,-95.9)..(36.2,-96.7)
				..controls(36.5,-97)and(37.2,-97.5)..(37.6,-97.7)..controls(38,-98)and(38.6,-98.4)..(39.1,-98.9)
				..controls(39.6,-99.4)and(40.3,-100)..(40.6,-100.3)
				..controls(41.3,-100.8)and(42.4,-101.8)..(43.7,-102.9)
				..controls(44,-103.2)and(44.6,-103.7)..(45,-103.9)
				..controls(45.4,-104.2)and(45.9,-104.6)..(46.1,-104.8)
				..controls(47.1,-105.9)and(48.8,-107.1)..(50.5,-107.9)
				..controls(51.8,-108.6)and(52.3,-109)..(53.4,-110)
				..controls(54.8,-111.2)and(55.6,-111.6)..(57,-112.1)
				..controls(57.5,-112.3)and(58.5,-112.8)..(59.4,-113.3)
				..controls(60.2,-113.7)and(61.3,-114.3)..(61.9,-114.5)
				..controls(62.4,-114.7)and(63.3,-115.1)..(63.8,-115.3)
				..controls(64.4,-115.6)and(65.4,-115.9)..(66.2,-116.1)
				..controls(67,-116.3)and(68.4,-116.6)..(69.2,-116.9)
				..controls(72.3,-117.9)and(72.8,-118)..(78,-118.1)
				..controls(84,-118.3)and(86.1,-118.1)..(88.1,-117.1)
				..controls(89.4,-116.5)and(93,-115)..(93.7,-114.9)
				..controls(94.8,-114.6)and(95.4,-114.2)..(95.9,-113.4)
				..controls(96.2,-113)and(96.7,-112.5)..(97.1,-112.2)
				..controls(98,-111.6)and(101.3,-108.4)..(102.1,-107.2)
				..controls(102.4,-106.8)and(102.9,-106.2)..(103.2,-105.8)
				..controls(103.5,-105.5)and(104,-104.8)..(104.2,-104.2)
				..controls(104.5,-103.6)and(105,-102.8)..(105.3,-102.4)
				..controls(105.6,-102)and(106.1,-101.2)..(106.5,-100.6)
				..controls(107.1,-99.5)and(108.6,-97.5)..(109.5,-96.5)
				..controls(109.9,-96.1)and(110.1,-96)..(110.4,-96)
				..controls(111.1,-96)and(111.6,-96.5)..(111.8,-97.6)
				..controls(112,-98.6)and(112.3,-99.3)..(112.8,-99.9)
				..controls(113.7,-100.9)and(115.2,-103.1)..(115.7,-104.2)
				..controls(116.1,-105)and(116.5,-105.5)..(116.9,-106)
				..controls(117.3,-106.4)and(117.8,-106.9)..(118,-107.3)
				..controls(118.2,-107.6)and(118.6,-108)..(119,-108.3)
				..controls(119.6,-108.8)and(121.1,-110.3)..(121.8,-111.2)
				..controls(122.1,-111.7)and(122.5,-112)..(123.4,-112.4)
				..controls(124.8,-113.2)and(125.8,-113.8)..(127,-114.6)
				..controls(127.5,-114.9)and(128.1,-115.3)..(128.2,-115.4)
				..controls(128.4,-115.5)and(129.3,-115.8)..(130.2,-116)
				..controls(131.1,-116.2)and(132.5,-116.6)..(133.3,-116.9)
				..controls(135.7,-117.9)and(136.1,-117.9)..(142.5,-117.8)
				..controls(145.8,-117.8)and(148.4,-117.7)..(148.8,-117.6)
				..controls(149.2,-117.6)and(150.2,-117.2)..(151,-116.8)
				..controls(152,-116.4)and(152.8,-116)..(153.6,-115.9)
				..controls(154.5,-115.7)and(155.2,-115.4)..(156.3,-114.9)
				..controls(157.1,-114.5)and(158.2,-114)..(158.7,-113.8)
				..controls(159.1,-113.7)and(159.8,-113.3)..(160.2,-113)
				..controls(160.6,-112.7)and(161.3,-112.4)..(161.7,-112.2)
				..controls(162.7,-111.9)and(163.2,-111.7)..(164.1,-111.1)
				..controls(164.5,-110.8)and(165.3,-110.4)..(165.9,-110)
				..controls(166.5,-109.7)and(167.3,-109.2)..(167.6,-108.9)
				..controls(167.9,-108.6)and(168.6,-108.2)..(169,-108)
				..controls(169.4,-107.8)and(170.2,-107.3)..(170.7,-106.9)
				..controls(171.2,-106.5)and(172.1,-105.8)..(172.7,-105.4)
				..controls(173.2,-105)and(174.1,-104.4)..(174.5,-104)
				..controls(174.9,-103.7)and(175.7,-103.1)..(176.2,-102.6)
				..controls(176.7,-102.2)and(177.4,-101.5)..(177.9,-101.2)
				..controls(178.3,-100.8)and(178.8,-100.4)..(179,-100.2)
				..controls(179.2,-100)and(179.6,-99.6)..(180,-99.3)
				..controls(180.4,-99.1)and(180.9,-98.6)..(181.2,-98.4)
				..controls(181.5,-98.1)and(182.1,-97.7)..(182.5,-97.5)
				..controls(183.2,-97.2)and(183.5,-96.9)..(184.3,-95.9)
				..controls(184.6,-95.6)and(184.6,-95.6)..(184.8,-95.7)
				..controls(185,-95.8)and(185.1,-96)..(185.1,-96.2)
				..controls(185.1,-96.5)and(185,-96.7)..(184.6,-97.2)
				..controls(184.1,-97.8)and(184,-97.9)..(183.9,-98.8)
				..controls(183.7,-100.3)and(183.3,-101.3)..(182.9,-101.8)
				..controls(182.6,-102.1)and(182.3,-102.6)..(182.2,-102.9)
				..controls(182,-103.3)and(181.6,-104)..(181.2,-104.6)
				..controls(180.9,-105.2)and(180.5,-106.1)..(180.4,-106.6)
				..controls(179.9,-108)and(179.7,-108.5)..(179,-109.4)
				..controls(178.6,-109.8)and(178.2,-110.5)..(177.9,-111)
				..controls(177.7,-111.5)and(177.3,-112.3)..(177,-112.7)
				..controls(176.6,-113.2)and(176.4,-113.7)..(176.3,-114.2)
				..controls(176.2,-114.7)and(175.9,-115.2)..(175.4,-116)
				..controls(175,-116.6)and(174.4,-117.7)..(174,-118.5)
				..controls(173.7,-119.2)and(173.3,-120.1)..(173,-120.5)
				..controls(172.8,-120.9)and(172.4,-121.8)..(172.1,-122.5)
				..controls(171.8,-123.2)and(171.3,-124.1)..(171,-124.6)
				..controls(170.8,-125.1)and(170.3,-126)..(169.9,-126.6)
				..controls(169.6,-127.3)and(169.2,-128.1)..(169,-128.3)
				..controls(168.8,-128.6)and(168.5,-129.3)..(168.3,-129.8)
				..controls(168.1,-130.3)and(167.8,-131.1)..(167.5,-131.6)
				..controls(166.9,-132.5)and(164.2,-137.9)..(164.2,-138.1)
				..controls(164.2,-138.1)and(163.9,-138.8)..(163.4,-139.7)
				..controls(162.9,-140.5)and(162.4,-141.5)..(162.3,-141.9)
				..controls(161.9,-143)and(161.6,-143.4)..(160.9,-144.1)
				..controls(160.4,-144.6)and(160.3,-144.8)..(160.1,-145.4)
				..controls(160,-145.8)and(159.6,-146.8)..(159.2,-147.8)
				..controls(158.7,-148.8)and(158.1,-150)..(157.9,-150.6)
				..controls(157.6,-151.1)and(157.2,-151.9)..(156.9,-152.4)
				..controls(156.5,-152.8)and(156.1,-153.7)..(155.8,-154.4)
				..controls(155.6,-155)and(155.2,-155.9)..(154.9,-156.3)
				..controls(154.7,-156.7)and(154.4,-157.4)..(154.2,-157.9)
				..controls(154.1,-158.4)and(153.8,-159)..(153.5,-159.3)
				..controls(152.8,-160.1)and(153,-160.1)..(142.4,-160.1)
				..controls(137.2,-160.1)and(131.8,-160.1)..(130.4,-160)--cycle
				(141.5,-103.7)..controls(140.8,-103)and(139.8,-102.1)..(139.2,-101.6)
				..controls(138.5,-101.2)and(137.9,-100.6)..(137.8,-100.4)
				..controls(137,-99.5)and(134,-96.5)..(133.2,-96)
				..controls(132.7,-95.7)and(131.8,-94.9)..(131.1,-94.3)
				..controls(129.4,-92.7)and(127.9,-91.6)..(126.6,-91.1)
				..controls(125.2,-90.5)and(124.1,-90)..(122.9,-89.1)
				..controls(122.1,-88.6)and(121.3,-88.2)..(119.9,-87.7)
				..controls(118.8,-87.3)and(117.5,-86.8)..(116.9,-86.5)
				..controls(114.3,-85.5)and(113.7,-85.4)..(109.7,-85.5)
				..controls(105.3,-85.7)and(103.9,-85.9)..(101.6,-87)
				..controls(100.7,-87.4)and(99.4,-87.8)..(97.9,-88)
				..controls(97,-88.2)and(96.6,-88.3)..(95.8,-88.9)
				..controls(95.4,-89.2)and(94.7,-89.6)..(94.2,-89.8)
				..controls(92.9,-90.3)and(92.3,-90.7)..(90.9,-92.1)
				..controls(90.2,-92.7)and(89.4,-93.3)..(89,-93.5)..controls(87,-94.5)and(85.2,-96)..(83.6,-97.9)
				..controls(83.3,-98.2)and(82.5,-98.9)..(81.8,-99.4)
				..controls(80.7,-100.1)and(80.5,-100.3)..(80,-101.1)
				..controls(78.7,-103.1)and(77.2,-103.9)..(76.4,-103)
				..controls(75.7,-102.2)and(75.7,-103.2)..(75.7,-83.6)
				..controls(75.6,-73.6)and(75.6,-64.8)..(75.5,-64)
				..controls(75.5,-63.2)and(75.5,-61.4)..(75.6,-59.8)--(75.7,-57.1)--(76,-56.8)
				..controls(76.8,-56.2)and(78,-56.3)..(79.8,-57.2)
				..controls(80.4,-57.5)and(81.3,-57.8)..(81.8,-57.9)
				..controls(82.3,-58.1)and(83,-58.4)..(83.6,-58.8)..controls(84.3,-59.2)and(85,-59.5)..(86,-59.8)
				..controls(86.8,-60.1)and(88,-60.6)..(88.7,-60.9)
				..controls(89.3,-61.2)and(90.6,-61.7)..(91.4,-62)
				..controls(92.2,-62.3)and(93.7,-62.9)..(94.6,-63.3)
				..controls(95.6,-63.7)and(96.8,-64.1)..(97.4,-64.3)
				..controls(98,-64.5)and(99,-64.9)..(99.7,-65.2)
				..controls(100.4,-65.6)and(101.5,-66)..(102.3,-66.2)
				..controls(104.2,-66.8)and(105,-67.1)..(106.2,-67.8)
				..controls(107.9,-68.7)and(110.2,-69.4)..(111.3,-69.2)
				..controls(112.5,-69)and(115,-68.2)..(115.6,-67.8)
				..controls(117.3,-66.7)and(117.6,-66.6)..(119.5,-66.2)
				..controls(121.9,-65.8)and(122.1,-65.7)..(123,-65.1)
				..controls(124.1,-64.4)and(124.4,-64.3)..(126,-63.9)
				..controls(127,-63.7)and(127.7,-63.4)..(128.4,-63.1)
				..controls(129,-62.8)and(129.7,-62.5)..(130.6,-62.4)
				..controls(132.1,-62.1)and(134,-61.5)..(135.2,-61)
				..controls(135.7,-60.8)and(136.7,-60.5)..(137.3,-60.3)
				..controls(138.4,-60)and(139.3,-59.7)..(141.1,-59)
				..controls(141.6,-58.7)and(142,-58.6)..(142.5,-58.6)
				..controls(143.2,-58.6)and(143.3,-58.7)..(143.6,-59)
				..controls(144,-59.5)and(144.1,-59.9)..(144,-60.8)
				..controls(144,-61.2)and(144,-71.1)..(144,-82.8)--(143.9,-104.1)--(143.7,-104.5)
				..controls(143.5,-104.7)and(143.3,-104.8)..(143,-104.9)
				..controls(142.6,-104.9)and(142.5,-104.8)..(141.5,-103.7)--cycle
				(27.8,-91.4)..controls(27.1,-90.9)and(26.3,-90)..(26.1,-89.5)
				..controls(26,-89.2)and(25.8,-88.6)..(25.7,-87.9)
				..controls(25.4,-86.2)and(25.6,-84.8)..(26.5,-83.5)
				..controls(27.4,-82.3)and(28.3,-82)..(29.4,-82.3)
				..controls(31.1,-82.8)and(32.1,-84.9)..(32,-87.4)
				..controls(31.9,-89.3)and(31.2,-90.6)..(29.9,-91.4)
				..controls(29.1,-91.9)and(28.6,-91.9)..(27.8,-91.4)--cycle
				(190.1,-91.3)..controls(189.7,-91.2)and(189.3,-90.9)..(189.1,-90.7)
				..controls(188.3,-89.8)and(187.6,-87.7)..(187.6,-86.2)
				..controls(187.6,-85.5)and(187.7,-85.2)..(188.2,-84.4)
				..controls(189.1,-82.6)and(190.1,-81.8)..(191.3,-82.1)
				..controls(192,-82.3)and(193.5,-83.7)..(193.8,-84.4)
				..controls(194.1,-85.1)and(194.2,-87.1)..(194,-88.3)
				..controls(193.8,-89.8)and(193.3,-90.5)..(192.3,-91.3)
				..controls(191.7,-91.7)and(190.9,-91.7)..(190.1,-91.3)--cycle
				(173.7,-83.4)..controls(172.2,-82.5)and(172,-82.5)..(170.6,-82.4)
				..controls(169.7,-82.4)and(168.8,-82.4)..(168.3,-82.5)
				..controls(167.2,-82.7)and(166.9,-82.7)..(166.6,-82.2)
				..controls(166.3,-81.8)and(166.3,-81)..(166.8,-80)
				..controls(167.2,-78.9)and(167.4,-77.4)..(167.7,-74.1)
				..controls(167.8,-72.8)and(167.9,-71.8)..(167.9,-71.7)
				..controls(168,-71.7)and(168,-70.2)..(168,-68.4)..controls(168,-65.4)and(168,-64.9)..(167.7,-63)
				..controls(167.3,-60.5)and(167.1,-59.9)..(166.4,-58.4)
				..controls(166.1,-57.7)and(165.8,-56.9)..(165.7,-56.6)
				..controls(165.7,-56.2)and(165.4,-55.6)..(165.2,-55.2)
				..controls(164.6,-53.8)and(161.9,-50.7)..(161,-50.3)
				..controls(159.8,-49.7)and(157.9,-50)..(155.2,-51.2)
				..controls(154.7,-51.4)and(153.4,-51.8)..(152.2,-52.1)
				..controls(150.8,-52.4)and(149.4,-52.9)..(148.6,-53.2)
				..controls(147.8,-53.6)and(146.7,-54)..(146,-54.2)
				..controls(145.4,-54.4)and(144.5,-54.7)..(144,-54.9)
				..controls(143.6,-55.2)and(142.7,-55.5)..(141.8,-55.7)
				..controls(141,-55.9)and(139.6,-56.3)..(138.7,-56.6)
				..controls(137.8,-57)and(136.3,-57.4)..(135.3,-57.6)
				..controls(133.8,-57.9)and(133.2,-58.1)..(132.1,-58.7)
				..controls(131.1,-59.2)and(130.3,-59.4)..(129.4,-59.6)
				..controls(128.7,-59.8)and(127.8,-60)..(127.4,-60.1)
				..controls(127,-60.3)and(126,-60.5)..(125.1,-60.7)
				..controls(123.2,-61.1)and(121.7,-61.5)..(120.2,-62.1)
				..controls(119.7,-62.3)and(118.1,-62.8)..(116.7,-63.2)
				..controls(114.6,-63.9)and(114,-64.1)..(113,-64.7)
				..controls(111.9,-65.4)and(111.8,-65.4)..(111,-65.5)
				..controls(109.9,-65.5)and(109.3,-65.4)..(107.7,-64.7)
				..controls(107,-64.4)and(106,-64)..(105.3,-63.9)
				..controls(103.5,-63.5)and(102.2,-63)..(101.2,-62.5)
				..controls(100.1,-61.8)and(99.3,-61.5)..(97.8,-61.1)
				..controls(97.2,-61)and(96.2,-60.7)..(95.7,-60.4)
				..controls(95.1,-60.1)and(94.2,-59.8)..(93.5,-59.6)
				..controls(92.9,-59.4)and(91.7,-59)..(91,-58.7)..controls(90.2,-58.4)and(88.9,-58)..(88,-57.7)
				..controls(87.2,-57.5)and(86.1,-57.1)..(85.5,-56.8)
				..controls(85,-56.6)and(84.1,-56.2)..(83.5,-56)..controls(82.9,-55.7)and(81.8,-55.2)..(81,-54.9)
				..controls(80.1,-54.5)and(79.1,-54.2)..(78.7,-54.1)
				..controls(77.9,-53.8)and(77.3,-53.6)..(76.2,-52.9)
				..controls(75.7,-52.6)and(75.2,-52.4)..(74.1,-52.2)
				..controls(72.5,-51.9)and(72.4,-51.8)..(71.3,-51.1)
				..controls(69.9,-50.1)and(69.4,-49.9)..(67.9,-49.8)
				..controls(67,-49.7)and(66.6,-49.6)..(66.3,-49.4)
				..controls(66.1,-49.3)and(65.7,-49)..(65.4,-48.8)
				..controls(65.1,-48.7)and(64.2,-48.4)..(63.6,-48.2)
				..controls(62.9,-48)and(61.7,-47.6)..(60.9,-47.2)
				..controls(60.1,-46.8)and(58.9,-46.4)..(58.2,-46.2)
				..controls(57.6,-46)and(56.8,-45.7)..(56.5,-45.5)
				..controls(54.9,-44.6)and(54.6,-44.5)..(53.2,-44.2)
				..controls(51.7,-43.8)and(51,-43.6)..(49.9,-42.9)
				..controls(49.5,-42.6)and(48.9,-42.4)..(47.5,-42)
				..controls(45.7,-41.6)and(45.6,-41.5)..(44.3,-40.6)
				..controls(43.5,-40.1)and(42.8,-39.6)..(42.7,-39.4)
				..controls(42.4,-38.7)and(43,-38.1)..(44.5,-37.6)..controls(44.9,-37.5)and(46,-37)..(46.9,-36.6)
				..controls(47.9,-36.1)and(48.8,-35.8)..(49.6,-35.6)
				..controls(50.5,-35.4)and(51.3,-35.2)..(52.4,-34.6)
				..controls(53.3,-34.1)and(54.5,-33.7)..(54.9,-33.5)
				..controls(55.4,-33.4)and(56.4,-33)..(57.2,-32.6)
				..controls(58.6,-32)and(60.1,-31.5)..(61.6,-31.2)
				..controls(62.3,-31.1)and(62.7,-31)..(63.2,-30.7)
				..controls(64.1,-30)and(64.3,-29.9)..(65.3,-29.7)
				..controls(65.7,-29.6)and(66.5,-29.2)..(67.1,-28.9)
				..controls(67.6,-28.6)and(68.8,-28.2)..(69.6,-27.9)
				..controls(70.5,-27.7)and(71.6,-27.3)..(72.1,-27)
				..controls(72.6,-26.7)and(73.8,-26.2)..(74.8,-25.9)
				..controls(75.7,-25.5)and(76.7,-25.1)..(77.1,-25)
				..controls(77.4,-24.8)and(78.5,-24.4)..(79.6,-24.1)
				..controls(80.6,-23.8)and(82,-23.2)..(82.8,-22.9)
				..controls(83.5,-22.5)and(84.8,-22)..(85.8,-21.8)
				..controls(86.7,-21.5)and(87.8,-21.1)..(88.1,-20.9)
				..controls(88.4,-20.7)and(89.3,-20.3)..(90.1,-20.1)
				..controls(90.8,-19.8)and(91.9,-19.4)..(92.4,-19.1)
				..controls(93.5,-18.5)and(94.4,-18.2)..(95.4,-18.1)
				..controls(95.7,-18)and(96.4,-17.8)..(96.9,-17.5)
				..controls(97.5,-17.3)and(98.4,-16.9)..(99,-16.6)
				..controls(99.7,-16.3)and(100.9,-15.9)..(101.9,-15.6)
				..controls(102.9,-15.3)and(104.4,-14.7)..(105.3,-14.3)
				..controls(108.1,-13.2)and(108.5,-13.1)..(110.1,-13.1)
				..controls(111.4,-13.1)and(111.5,-13.1)..(112.4,-13.5)
				..controls(113.2,-14)and(115,-14.8)..(116.6,-15.4)
				..controls(116.9,-15.5)and(117.9,-15.8)..(118.8,-16.1)
				..controls(120.9,-16.6)and(121.5,-16.9)..(122.6,-17.9)
				..controls(123,-18.3)and(123.7,-18.8)..(123.9,-18.9)
				..controls(124.6,-19.4)and(126,-19.8)..(127.9,-20.2)
				..controls(129.6,-20.5)and(129.8,-20.6)..(131.1,-21.4)
				..controls(131.5,-21.7)and(132.4,-22.1)..(133.1,-22.3)
				..controls(136.1,-23.2)and(138.3,-24.1)..(140.2,-25)
				..controls(141,-25.5)and(141.5,-25.6)..(142.5,-25.8)
				..controls(145.1,-26.3)and(146,-26.6)..(149.4,-28.1)
				..controls(151.3,-28.9)and(153.5,-29.8)..(154.3,-30.1)
				..controls(155.1,-30.4)and(156.2,-30.8)..(156.6,-31)
				..controls(157.1,-31.2)and(158.2,-31.6)..(159.1,-31.8)
				..controls(160.6,-32.2)and(160.9,-32.3)..(161.7,-32.9)
				..controls(162.4,-33.3)and(162.9,-33.6)..(163.7,-33.8)
				..controls(164.3,-34)and(164.8,-34.1)..(164.9,-34.1)
				..controls(165,-34.1)and(165.6,-34.4)..(166.4,-34.8)
				..controls(167.3,-35.2)and(168.1,-35.4)..(169,-35.6)
				..controls(170,-35.8)and(170.5,-36)..(171.7,-36.6)
				..controls(172.4,-37)and(173.3,-37.4)..(173.6,-37.4)
				..controls(175.6,-37.7)and(176.2,-38)..(176.7,-39.1)
				..controls(176.9,-39.6)and(176.9,-39.7)..(176.7,-40.4)
				..controls(176.4,-41.6)and(175.8,-42.7)..(175.2,-43.4)
				..controls(174.5,-44.1)and(172.6,-45.5)..(171.7,-45.8)
				..controls(170.6,-46.2)and(169.1,-46)..(168.2,-45.4)
				..controls(167.9,-45.2)and(167.3,-44.8)..(166.9,-44.4)
				..controls(166.4,-44.1)and(165.8,-43.6)..(165.4,-43.2)
				..controls(165.1,-42.9)and(164.4,-42.3)..(163.9,-41.9)
				..controls(163.3,-41.5)and(162.7,-40.9)..(162.4,-40.5)
				..controls(161.4,-39.2)and(159.9,-38)..(158.8,-37.8)
				..controls(158.4,-37.8)and(158.3,-37.7)..(158.2,-37.2)
				..controls(158,-36.8)and(157.4,-36.1)..(157.1,-36.1)
				..controls(157,-36.1)and(156.7,-36.1)..(156.5,-36)
				..controls(156.1,-35.9)and(155.9,-35.7)..(155.7,-35.2)
				..controls(155.3,-34.6)and(155.1,-34.4)..(154,-34.1)
				..controls(153.6,-34)and(153.2,-33.7)..(152.9,-33.5)
				..controls(152.4,-32.9)and(151.6,-32.4)..(151.3,-32.4)
				..controls(150.9,-32.4)and(150.4,-32.6)..(150.4,-32.9)
				..controls(150.4,-33.2)and(150.8,-33.4)..(151.4,-33.6)
				..controls(152.1,-33.7)and(152.8,-34.1)..(153.5,-34.9)
				..controls(154.1,-35.5)and(154.3,-35.6)..(155,-35.8)
				..controls(155.7,-36)and(155.7,-36)..(156,-36.6)
				..controls(156.3,-37.3)and(156.7,-37.7)..(157.5,-37.8)
				..controls(158.1,-37.9)and(158.2,-38)..(158.4,-38.8)
				..controls(158.5,-39.4)and(158.6,-39.6)..(159.6,-40.4)
				..controls(160.2,-41)and(160.9,-41.6)..(161.3,-41.8)
				..controls(161.7,-42)and(161.9,-42.3)..(162.2,-42.8)
				..controls(162.4,-43.2)and(162.8,-43.8)..(163.1,-44.1)
				..controls(163.9,-44.8)and(165.1,-46.7)..(165.7,-48)
				..controls(166.2,-49.3)and(166.4,-49.7)..(167.2,-50.5)
				..controls(167.5,-50.8)and(167.9,-51.4)..(168.1,-51.8)
				..controls(168.3,-52.2)and(168.8,-53)..(169.1,-53.6)
				..controls(169.6,-54.3)and(169.8,-54.9)..(169.9,-55.6)
				..controls(170.4,-57.8)and(170.5,-58.1)..(171.2,-59.5)
				..controls(171.6,-60.3)and(172,-61.5)..(172.2,-62.1)
				..controls(172.3,-62.8)and(172.7,-63.7)..(172.9,-64.1)
				..controls(173.2,-64.6)and(173.4,-65.3)..(173.4,-65.6)
				..controls(173.8,-68.8)and(174.2,-70.6)..(174.9,-72.6)
				..controls(175.5,-74.5)and(175.9,-77.1)..(176.1,-80.2)
				..controls(176.3,-83.4)and(175.5,-84.4)..(173.7,-83.4)--cycle;
				\definecolor{ezs}{RGB}{247,179,26}\fill[ezs](-5.7,-156.6)--(-3.4,-157.5)--(61.9,12.7)--(59.7,13.6)--cycle;
				\definecolor{ezs}{RGB}{247,179,22}\fill[ezs](-3.8,-157.3)--(-1.5,-158.2)--(63.9,12)--(61.6,12.9)--cycle;
				\definecolor{ezs}{RGB}{248,179,19}\fill[ezs](-1.8,-158.1)--(0.4,-158.9)--(65.8,11.2)--(63.5,12.1)--cycle;
				\definecolor{ezs}{RGB}{248,178,16}\fill[ezs](0.1,-158.8)--(2.4,-159.7)--(67.7,10.5)--(65.4,11.4)--cycle;
				\definecolor{ezs}{RGB}{249,178,15}\fill[ezs](2,-159.5)--(4.3,-160.4)--(69.6,9.8)--(67.3,10.7)--cycle;
				\definecolor{ezs}{RGB}{251,177,17}\fill[ezs](3.9,-160.3)--(6.2,-161.1)--(71.5,9)--(69.2,9.9)--cycle;
				\definecolor{ezs}{RGB}{252,179,18}\fill[ezs](5.8,-161)--(8.1,-161.9)--(73.4,8.3)--(71.1,9.2)--cycle;
				\definecolor{ezs}{RGB}{252,182,17}\fill[ezs](7.7,-161.7)--(10,-162.6)--(75.3,7.6)--(73,8.5)--cycle;
				\definecolor{ezs}{RGB}{252,185,17}\fill[ezs](9.6,-162.5)--(11.9,-163.4)--(77.2,6.8)--(75,7.7)--cycle;
				\definecolor{ezs}{RGB}{251,189,18}\fill[ezs](11.5,-163.2)--(13.8,-164.1)--(79.2,6.1)--(76.9,7)--cycle;
				\definecolor{ezs}{RGB}{251,191,22}\fill[ezs](13.4,-163.9)--(15.7,-164.8)--(81.1,5.4)--(78.8,6.2)--cycle;
				\definecolor{ezs}{RGB}{252,193,25}\fill[ezs](15.4,-164.7)--(17.6,-165.6)--(83,4.6)--(80.7,5.5)--cycle;
				\definecolor{ezs}{RGB}{253,195,27}\fill[ezs](17.3,-165.4)--(19.6,-166.3)--(84.9,3.9)--(82.6,4.8)--cycle;
				\definecolor{ezs}{RGB}{253,196,31}\fill[ezs](19.2,-166.1)--(21.5,-167)--(86.8,3.2)--(84.5,4)--cycle;
				\definecolor{ezs}{RGB}{254,199,35}\fill[ezs](21.1,-166.9)--(23.4,-167.8)--(88.7,2.4)--(86.4,3.3)--cycle;
				\definecolor{ezs}{RGB}{253,201,38}\fill[ezs](23,-167.6)--(25.3,-168.5)--(90.6,1.7)--(88.3,2.6)--cycle;
				\definecolor{ezs}{RGB}{253,203,40}\fill[ezs](24.9,-168.3)--(27.2,-169.2)--(92.5,1)--(90.3,1.8)--cycle;
				\definecolor{ezs}{RGB}{253,205,42}\fill[ezs](26.8,-169.1)--(29.1,-170)--(94.4,0.2)--(92.2,1.1)--cycle;
				\definecolor{ezs}{RGB}{253,207,45}\fill[ezs](28.7,-169.8)--(31,-170.7)--(96.4,-0.5)--(94.1,0.4)--cycle;
				\definecolor{ezs}{RGB}{253,210,49}\fill[ezs](30.7,-170.5)--(32.9,-171.4)--(98.3,-1.2)--(96,-0.4)--cycle;
				\definecolor{ezs}{RGB}{253,212,54}\fill[ezs](32.6,-171.3)--(34.9,-172.2)--(100.2,-2)--(97.9,-1.1)--cycle;
				\definecolor{ezs}{RGB}{253,214,59}\fill[ezs](34.5,-172)--(36.8,-172.9)--(102.1,-2.7)--(99.8,-1.8)--cycle;
				\definecolor{ezs}{RGB}{253,216,63}\fill[ezs](36.4,-172.7)--(38.7,-173.6)--(104,-3.4)--(101.7,-2.6)--cycle;
				\definecolor{ezs}{RGB}{253,218,65}\fill[ezs](38.3,-173.5)--(40.6,-174.4)--(105.9,-4.2)--(103.6,-3.3)--cycle;
				\definecolor{ezs}{RGB}{254,220,66}\fill[ezs](40.2,-174.2)--(42.5,-175.1)--(107.8,-4.9)--(105.5,-4)--cycle;
				\definecolor{ezs}{RGB}{254,222,68}\fill[ezs](42.1,-175)--(44.4,-175.8)--(109.7,-5.6)--(107.5,-4.8)--cycle;
				\definecolor{ezs}{RGB}{254,224,71}\fill[ezs](44,-175.7)--(46.3,-176.6)--(111.7,-6.4)--(109.4,-5.5)--cycle;
				\definecolor{ezs}{RGB}{254,226,75}\fill[ezs](46,-176.4)--(48.2,-177.3)--(113.6,-7.1)--(111.3,-6.2)--cycle;
				\definecolor{ezs}{RGB}{254,228,79}\fill[ezs](47.9,-177.2)--(50.1,-178)--(115.5,-7.8)--(113.2,-7)--cycle;
				\definecolor{ezs}{RGB}{254,230,82}\fill[ezs](49.8,-177.9)--(52.1,-178.8)--(117.4,-8.6)--(115.1,-7.7)--cycle;
				\definecolor{ezs}{RGB}{254,232,85}\fill[ezs](51.7,-178.6)--(54,-179.5)--(119.3,-9.3)--(117,-8.4)--cycle;
				\definecolor{ezs}{RGB}{254,234,88}\fill[ezs](53.6,-179.4)--(55.9,-180.2)--(121.2,-10)--(118.9,-9.2)--cycle;
				\definecolor{ezs}{RGB}{254,237,91}\fill[ezs](55.5,-180.1)--(57.8,-181)--(123.1,-10.8)--(120.8,-9.9)--cycle;
				\definecolor{ezs}{RGB}{254,239,94}\fill[ezs](57.4,-180.8)--(59.7,-181.7)--(125,-11.5)--(122.8,-10.6)--cycle;
				\definecolor{ezs}{RGB}{254,241,97}\fill[ezs](59.3,-181.6)--(61.6,-182.4)--(127,-12.2)--(124.7,-11.4)--cycle;
				\definecolor{ezs}{RGB}{254,243,101}\fill[ezs](61.2,-182.3)--(63.5,-183.2)--(128.9,-13)--(126.6,-12.1)--cycle;
				\definecolor{ezs}{RGB}{254,245,103}\fill[ezs](63.2,-183)--(65.4,-183.9)--(130.8,-13.7)--(128.5,-12.8)--cycle;
				\definecolor{ezs}{RGB}{254,245,105}\fill[ezs](65.1,-183.8)--(67.4,-184.6)--(132.7,-14.4)--(130.4,-13.6)--cycle;
				\definecolor{ezs}{RGB}{254,245,105}\fill[ezs](67,-184.5)--(69.3,-185.4)--(134.6,-15.2)--(132.3,-14.3)--cycle;
				\definecolor{ezs}{RGB}{254,244,103}\fill[ezs](68.9,-185.2)--(71.2,-186.1)--(136.5,-15.9)--(134.2,-15)--cycle;
				\definecolor{ezs}{RGB}{253,243,101}\fill[ezs](70.8,-186)--(73.1,-186.8)--(138.4,-16.6)--(136.1,-15.8)--cycle;
				\definecolor{ezs}{RGB}{253,241,98}\fill[ezs](72.7,-186.7)--(75,-187.6)--(140.3,-17.4)--(138,-16.5)--cycle;
				\definecolor{ezs}{RGB}{253,239,95}\fill[ezs](74.6,-187.4)--(76.9,-188.3)--(142.2,-18.1)--(140,-17.2)--cycle;
				\definecolor{ezs}{RGB}{253,236,92}\fill[ezs](76.5,-188.2)--(78.8,-189)--(144.2,-18.8)--(141.9,-18)--cycle;
				\definecolor{ezs}{RGB}{253,234,88}\fill[ezs](78.5,-188.9)--(80.7,-189.8)--(146.1,-19.6)--(143.8,-18.7)--cycle;
				\definecolor{ezs}{RGB}{253,232,85}\fill[ezs](80.4,-189.6)--(82.7,-190.5)--(148,-20.3)--(145.7,-19.4)--cycle;
				\definecolor{ezs}{RGB}{253,230,82}\fill[ezs](82.3,-190.4)--(84.6,-191.2)--(149.9,-21)--(147.6,-20.2)--cycle;
				\definecolor{ezs}{RGB}{253,228,78}\fill[ezs](84.2,-191.1)--(86.5,-192)--(151.8,-21.8)--(149.5,-20.9)--cycle;
				\definecolor{ezs}{RGB}{253,226,74}\fill[ezs](86.1,-191.8)--(88.4,-192.7)--(153.7,-22.5)--(151.4,-21.6)--cycle;
				\definecolor{ezs}{RGB}{253,224,70}\fill[ezs](88,-192.6)--(90.3,-193.4)--(155.6,-23.3)--(153.3,-22.4)--cycle;
				\definecolor{ezs}{RGB}{253,222,67}\fill[ezs](89.9,-193.3)--(92.2,-194.2)--(157.5,-24)--(155.3,-23.1)--cycle;
				\definecolor{ezs}{RGB}{253,220,65}\fill[ezs](91.8,-194)--(94.1,-194.9)--(159.5,-24.7)--(157.2,-23.8)--cycle;
				\definecolor{ezs}{RGB}{253,217,63}\fill[ezs](93.7,-194.8)--(96,-195.6)--(161.4,-25.5)--(159.1,-24.6)--cycle;
				\definecolor{ezs}{RGB}{253,215,60}\fill[ezs](95.7,-195.5)--(97.9,-196.4)--(163.3,-26.2)--(161,-25.3)--cycle;
				\definecolor{ezs}{RGB}{253,213,56}\fill[ezs](97.6,-196.2)--(99.9,-197.1)--(165.2,-26.9)--(162.9,-26)--cycle;
				\definecolor{ezs}{RGB}{252,211,53}\fill[ezs](99.5,-197)--(101.8,-197.8)--(167.1,-27.7)--(164.8,-26.8)--cycle;
				\definecolor{ezs}{RGB}{252,209,49}\fill[ezs](101.4,-197.7)--(103.7,-198.6)--(169,-28.4)--(166.7,-27.5)--cycle;
				\definecolor{ezs}{RGB}{252,207,45}\fill[ezs](103.3,-198.4)--(105.6,-199.3)--(170.9,-29.1)--(168.6,-28.2)--cycle;
				\definecolor{ezs}{RGB}{253,205,42}\fill[ezs](105.2,-199.2)--(107.5,-200)--(172.8,-29.9)--(170.6,-29)--cycle;
				\definecolor{ezs}{RGB}{253,202,40}\fill[ezs](107.1,-199.9)--(109.4,-200.8)--(174.7,-30.6)--(172.5,-29.7)--cycle;
				\definecolor{ezs}{RGB}{253,200,38}\fill[ezs](109,-200.6)--(111.3,-201.5)--(176.7,-31.3)--(174.4,-30.4)--cycle;
				\definecolor{ezs}{RGB}{253,198,35}\fill[ezs](111,-201.4)--(113.2,-202.2)--(178.6,-32.1)--(176.3,-31.2)--cycle;
				\definecolor{ezs}{RGB}{253,196,31}\fill[ezs](112.9,-202.1)--(115.2,-203)--(180.5,-32.8)--(178.2,-31.9)--cycle;
				\definecolor{ezs}{RGB}{253,194,28}\fill[ezs](114.8,-202.8)--(117.1,-203.7)--(182.4,-33.5)--(180.1,-32.6)--cycle;
				\definecolor{ezs}{RGB}{253,192,26}\fill[ezs](116.7,-203.6)--(119,-204.5)--(184.3,-34.3)--(182,-33.4)--cycle;
				\definecolor{ezs}{RGB}{253,189,24}\fill[ezs](118.6,-204.3)--(120.9,-205.2)--(186.2,-35)--(183.9,-34.1)--cycle;
				\definecolor{ezs}{RGB}{253,186,20}\fill[ezs](120.5,-205)--(122.8,-205.9)--(188.1,-35.7)--(185.8,-34.8)--cycle;
				\definecolor{ezs}{RGB}{253,184,17}\fill[ezs](122.4,-205.8)--(124.7,-206.7)--(190,-36.5)--(187.8,-35.6)--cycle;
				\definecolor{ezs}{RGB}{252,183,15}\fill[ezs](124.3,-206.5)--(126.6,-207.4)--(192,-37.2)--(189.7,-36.3)--cycle;
				\definecolor{ezs}{RGB}{252,181,14}\fill[ezs](126.3,-207.2)--(128.5,-208.1)--(193.9,-37.9)--(191.6,-37.1)--cycle;
				\definecolor{ezs}{RGB}{252,179,13}\fill[ezs](128.2,-208)--(130.4,-208.9)--(195.8,-38.7)--(193.5,-37.8)--cycle;
				\definecolor{ezs}{RGB}{252,178,12}\fill[ezs](130.1,-208.7)--(132.4,-209.6)--(197.7,-39.4)--(195.4,-38.5)--cycle;
				\definecolor{ezs}{RGB}{251,178,10}\fill[ezs](132,-209.4)--(134.3,-210.3)--(199.6,-40.1)--(197.3,-39.3)--cycle;
				\definecolor{ezs}{RGB}{251,178,10}\fill[ezs](133.9,-210.2)--(136.2,-211.1)--(201.5,-40.9)--(199.2,-40)--cycle;
				\definecolor{ezs}{RGB}{250,177,11}\fill[ezs](135.8,-210.9)--(138.1,-211.8)--(203.4,-41.6)--(201.1,-40.7)--cycle;
				\definecolor{ezs}{RGB}{250,178,11}\fill[ezs](137.7,-211.6)--(140,-212.5)--(205.3,-42.3)--(203.1,-41.5)--cycle;
				\definecolor{ezs}{RGB}{248,180,10}\fill[ezs](139.6,-212.4)--(141.9,-213.3)--(207.3,-43.1)--(205,-42.2)--cycle;
				\definecolor{ezs}{RGB}{247,181,9}\fill[ezs](141.5,-213.1)--(143.8,-214)--(209.2,-43.8)--(206.9,-42.9)--cycle;
				\definecolor{ezs}{RGB}{246,182,9}\fill[ezs](143.5,-213.8)--(145.7,-214.7)--(211.1,-44.5)--(208.8,-43.7)--cycle;
				\definecolor{ezs}{RGB}{246,182,10}\fill[ezs](145.4,-214.6)--(147.7,-215.5)--(213,-45.3)--(210.7,-44.4)--cycle;
			\end{scope}
			% chữ EDUZY (Times đậm, làm dày nét bằng 5 lớp lệch nhẹ; căn đúng khung chữ logo gốc)
			\foreach \dx/\dy in {-0.7/0,0.7/0,0/-0.3,0/0.3,0/0}
			\node[anchor=center,inner sep=0pt,outer sep=0pt,text=ezText,align=center] at (109.6+\dx,-193.6+\dy)
			{\resizebox{1.01cm}{0.22cm}{\fontencoding{T1}\fontfamily{ptm}\fontseries{b}\selectfont EDUZY}};
		\end{tikzpicture}%
	}%
}
%====================== THIẾT LẬP THÂN ĐỀ ======================
%--- Môi trường tự luận {bt}: ex_test 3.3.4 chưa khai báo sẵn nên tạo ở đây
%    (ex_test đã tự móc \fixshowans và \listenumerate cho bt)
\makeatletter
\@ifundefined{bt}{\newtheorem{bt}{\color{gold1}\namebt}}{}
\makeatother
\renewcommand{\sectionmark}[1]{\markright{\thesection\ #1}}
%--- Khung tiêu đề (logo + tối đa 4 dòng; dòng 3, 4 để trống {} thì bỏ qua)
%  Đề thi:    \khungtieude{THI THỬ THPT QUỐC GIA NĂM 2026}{Môn Toán -- Đề số 21}
%                         {(Đề thi gồm: ...)}{Thời gian làm bài: 90 phút (không kể thời gian phát đề).}
%  Chuyên đề: \khungtieude{CHUYÊN ĐỀ 3 -- DÃY SỐ VÀ CÔNG THỨC TRUY HỒI}
%                         {Bồi dưỡng học sinh giỏi Toán THPT}{...}{}
\newcommand{\khungtieude}[4]{%
	\begin{tcolorbox}[colframe=gold1,colback=gold5!45,boxrule=0.5mm]
		\begin{minipage}[c]{3cm}
			\centerline{\eduzylogo[2.4cm]}
		\end{minipage}\hspace{0.5cm}
		\begin{minipage}[c]{12.5cm}
			\vskip 0.15 cm
			\centerline{\ind{#1}}
			\vskip 0.15 cm
			\centerline{\textbf{#2}}
			\if\relax\detokenize{#3}\relax\else
				\vskip 0.15 cm
				\centerline{\textit{#3}}
			\fi
			\if\relax\detokenize{#4}\relax\else
				\vskip 0.15 cm
				\centerline{\textit{#4}}
			\fi
			\vskip 0.15 cm
		\end{minipage}
	\end{tcolorbox}
}
%====================== BỔ SUNG CHO BÀI HỌC, CHUYÊN ĐỀ ======================
\usepackage[shortlabels]{enumitem}
\usepackage{tabularx}
%--- Danh sách gọn: \begin{enumerate}[gon] ... \end{enumerate}
\SetEnumitemKey{gon}{leftmargin=*,itemsep=1pt,parsep=0pt,topsep=3pt}
%--- Khung lý thuyết: \begin{khung}{} ... \end{khung}; tiêu đề chỉ dùng khi thật ngắn ("Mẹo", "Lỗi thường gặp")
\newtcolorbox{khung}[1]{breakable,enhanced,colframe=gold1,colback=gold5!45,
	boxrule=0.4mm,fonttitle=\bfseries\sffamily,coltitle=white,colbacktitle=gold1,
	attach boxed title to top left={xshift=6mm,yshift=-3mm},
	boxed title style={arc=1mm,boxrule=0.3mm},title={#1}}
%--- Tiêu đề mục nhỏ trong một dạng: \muc{Dấu hiệu}, \muc{Phương pháp}, \muc{Bài tự luyện}
\makeatletter
\newcommand{\muc}[1]{\par\medskip\noindent\indam{#1}\par\nobreak\smallskip\@afterindentfalse\@afterheading}
\makeatother
%--- Ô công thức cần nhớ (căn giữa, viền vàng, vừa khít nội dung)
%     \begin{center}\congthuc{$Q(n)=(n-1)!$}\end{center}
\newtcbox{\congthuc}{on line,colframe=gold1,colback=gold5!45,boxrule=0.5pt,arc=2pt,
	left=8pt,right=8pt,top=4pt,bottom=4pt}
%--- Công tắc lời giải bài tự luyện (ví dụ mẫu luôn dùng \loigiai)
%     \loigiaiTL[{đáp số}]{lời giải đầy đủ}
%     \TLfull -> in lời giải đầy đủ;  \TLdapso -> chỉ in một dòng "Đáp số";  \TLan -> ẩn hoàn toàn
%     (\CheDo là dethi thì mọi lời giải và đáp số đều ẩn)
\makeatletter
\newif\ifTLhien   \TLhientrue
\newif\ifTLdapso  \TLdapsofalse
\newcommand{\TLfull}{\TLhientrue\TLdapsofalse}
\newcommand{\TLdapso}{\TLhientrue\TLdapsotrue}
\newcommand{\TLan}{\TLhienfalse}
\newcommand{\loigiaiTL}[2][]{%
	\ifTLhien
		\ifTLdapso
			\def\@tlarg{#1}%
			\ifx\@tlarg\@empty\else
				\ifdefstring{\kindOPTN}{dethi}{}{\par\noindent{\bfseries\color{gold1}Đáp số:} #1\par}%
			\fi
		\else
			\loigiai{#2}%
		\fi
	\fi}
\makeatother
\TLdapso% mặc định: chỉ in đáp số bài tự luyện
%====================== BẢNG ĐÚNG/SAI (\choiceTFt) ======================
%  Hàng tiêu đề nền nâu đồng chữ trắng (Phát biểu | Đ | S), các hàng tô xen kẽ kem / trắng,
%  đường kẻ vàng kim, nhãn a) b) c) d) màu nâu đồng. Bản có đáp án: cột được chọn có vòng tròn
%  nâu đồng với dấu tích trắng (câu đúng tích ở cột Đúng, câu sai tích ở cột Sai).
%  Tiêu đề ghi đủ "Đúng | Sai": \OPTN{viettat=0} sau \input (cột Đúng/Sai rộng ra, phát biểu hẹp lại).
%  Dùng bảng gốc của ex_test: \def\BangDSgoc{1} trước %============================================================
%  PREAMBLE EDUZY v2 - dùng chung cho đề thi, bài học, chuyên đề (gói ex_test)
%  Khác bản preamble-eduzy.tex ở đầu trang, chân trang: dải nền nâu đồng tràn ngang,
%  logo + Eduzy.com.vn, ô chữ vàng kem cạnh xiên, chân trang "Trang (số)";
%  logo Eduzy chìm giữa mọi trang (tắt: \logochimtat, bật lại: \logochimbat).
%  Cách dùng ở file chính (xem mau-chuan.tex, mau-chuyen-de.tex):
%      \def\CheDo{loigiai}          % loigiai | dethi | solcolor | color | book
%      \def\tenchuong{<dòng chữ trên đầu trang>}
%      \def\thuonghieu{<chữ ở chân trang lẻ>}   % không khai báo thì dùng mặc định
%      \def\tacgia{<chữ ở chân trang chẵn>}      % không khai báo thì dùng mặc định
%      \def\khauhieu{<chữ cạnh logo ở đầu trang>}  % tùy chọn, mặc định Eduzy.com.vn
%      \input{preamble-eduzy-v2}
%      \TLdapso                      % tài liệu: \TLfull | \TLdapso | \TLan
%      \begin{document} ... \end{document}
%  Logo Eduzy vẽ bằng TikZ ngay trong file (\eduzylogo), không cần file ảnh.
%  Biên dịch: pdfLaTeX (máy có gói vietnam) hoặc XeLaTeX.
%============================================================
\documentclass[12pt,a4paper,twoside]{article}%<-- twoside: bat dau trang/chan trang chan-le
\usepackage{iftex}
\usepackage{amsmath,amssymb,tikz,mathrsfs,tkz-tab}
\usepackage{graphicx}
\usepackage{pgfplots}% đồ thị hàm số
\pgfplotsset{compat=1.17}
%\usepackage{mathptmx}
\usepackage{mathpazo}
\ifXeTeX% XeLaTeX: font chữ có đủ dấu tiếng Việt, giữ nguyên font toán mathpazo
	\usepackage[no-math]{fontspec}
	\setmainfont{TeX Gyre Pagella}
	\newfontfamily\EZqag{TeX Gyre Adventor}[NFSSFamily=qag,Ligatures=TeX]
\fi
\usepackage[many]{tcolorbox}
\usetikzlibrary{math,through,calc,intersections,angles,quotes,shapes,shapes.geometric,arrows,patterns,matrix,chains,arrows.meta,decorations.shapes,decorations.fractals,decorations.markings,shadows,decorations.pathreplacing}
\usepackage[top=1.6cm, bottom=1.8cm, left=1.8cm, right=1.8cm]{geometry}% v2: top 1.6cm (bản cũ 1.5cm) để chữ không sát dải đầu trang
\providecommand{\CheDo}{loigiai}% loigiai | dethi | solcolor | color | book
\expandafter\usepackage\expandafter[\CheDo]{ex_test}
% \tenchuong: khai báo ở file chính (trước \input) hoặc \renewcommand sau \input
%============================================================
%  [GỘP TỪ eduzy-de.tex] EDUZY - MÀU NÂU ĐỒNG / VÀNG KIM
%  + ĐẦU TRANG, CHÂN TRANG KIỂU DẢI NGANG (v2), ĐỐI XỨNG TRANG CHẴN / TRANG LẺ.
%  (không cần file eduzy-de.tex riêng nữa)
%  Khối này phải đặt SAU \usepackage[...]{ex_test} và SAU \def\tenchuong{...}
%============================================================
\usepackage{lastpage}
\usepackage{needspace}
%--- DÀN TRANG ĐỀ THI: tách câu (THÊM 05/11/2026; chỉ thêm lệnh mới, KHÔNG đổi hành vi mặc định) ---
%  Hai công tắc, mặc định tắt (= preamble chạy y như cũ). Bật bằng cách đặt TRƯỚC \input{preamble-eduzy-v2}
%  trong file đề (tools/dan_trang/dan_trang.py tự ghi khi đặt \NeedspaceCau):
%    \def\TachCauDS{1}   câu Đúng/Sai được tách: câu dẫn (kèm hình) ở cuối trang, bảng Phát biểu sang trang sau
%                        (câu dẫn và hình không tách nhau, bảng không tách giữa hai trang)
%    \def\TachCauLG{1}   bản lời giải (\CheDo khác dethi): phần đề của câu cao hơn 1/3 trang được phép tách
%  Lệnh dùng trong file đề (giá trị do tools/dan_trang/dan_trang.py đo ở bản dethi):
%    \NeedspaceCau[<cao câu dẫn + hình>]{<cao cả khối câu>}   ngay trên \begin{ex} (câu đầu phần: trên tiêu đề phần)
%        công tắc tắt: đúng bằng \Needspace{<cao cả khối câu>} như cũ;
%        \TachCauDS=1 và có [..] (câu Đúng/Sai được chọn tách): chỉ giữ câu dẫn + hình;
%        \TachCauLG=1 ở bản lời giải: đòi tối đa 1/3 \textheight.
%    \NeedspaceBang{<cao bảng Phát biểu>}%   ngay trước \choiceTFt: \TachCauDS=1 thì giữ trọn bảng trên một trang;
%        công tắc tắt: không làm gì.
\providecommand{\TachCauDS}{0}
\providecommand{\TachCauLG}{0}
\newlength{\EZnsCau}
\newcommand{\NeedspaceCau}[2][]{%
	\setlength{\EZnsCau}{#2}%
	\ifnum\TachCauDS=1 \if\relax\detokenize{#1}\relax\else\setlength{\EZnsCau}{#1}\fi\fi
	\ifnum\TachCauLG=1 \ifdefstring{\kindOPTN}{dethi}{}{\ifdim\EZnsCau>0.3333\textheight\setlength{\EZnsCau}{0.3333\textheight}\fi}\fi
	\Needspace{\EZnsCau}}
\newcommand{\NeedspaceBang}[1]{\ifnum\TachCauDS=1 \Needspace{#1}\fi}

%--- Thang màu: đậm -> nhạt
\definecolor{gold1}{RGB}{126,96,0}    % nâu đồng đậm  #7E6000
\definecolor{gold2}{RGB}{192,144,3}   % vàng kim      #C09003
\definecolor{gold3}{RGB}{255,215,99}  % vàng nhạt     #FFD763
\definecolor{gold4}{RGB}{255,231,153} % vàng rất nhạt #FFE799
\definecolor{gold5}{RGB}{255,243,205} % vàng kem      #FFF3CD
%--- Ghi đè tông xanh/đỏ cũ
\definecolor{myblue}{RGB}{126,96,0}
\colorlet{blue}{gold1}
\colorlet{cyan}{gold2}
\colorlet{teal}{gold2}

%--- Tên chương mặc định (file đề nên tự khai báo lại)
\providecommand{\tenchuong}{}
%--- Chân trang
\providecommand{\thuonghieu}{TOÁN 12 -- Eduzy.com.vn}
\providecommand{\tacgia}{Tô Gia Bảo -- 0353703783}

%--- Chữ cạnh logo ở đầu trang
\providecommand{\khauhieu}{Eduzy.com.vn}

%====================== THÔNG SỐ ĐẦU TRANG, CHÂN TRANG (v2) ======================
%  Dải nền đậm tràn ngang trang; ô chữ vàng kem có cạnh xiên, kèm 2 vạch xiên chuyển màu.
%  Đầu trang lẻ : logo + \khauhieu (trái)   | ô chữ \tenchuong (phải)
%  Chân trang lẻ: ô chữ \thuonghieu (trái)  | Trang (số) (phải)
%  Trang chẵn soi gương trang lẻ; chân trang chẵn ghi \tacgia.
%--- Màu: chỉ dùng thang vàng kim Eduzy
\colorlet{ezNen}{gold1}% dải nền đậm
\colorlet{ezVach1}{gold3}% vạch xiên sát dải nền (đậm hơn)
\colorlet{ezVach2}{gold4}% vạch xiên sát ô chữ (nhạt hơn)
\colorlet{ezO}{gold5}% ô chữ vàng kem
%--- Kích thước
\newlength{\EZdai}\setlength{\EZdai}{1.1cm}% chiều cao dải nền
\newlength{\EZo}\setlength{\EZo}{0.68cm}% chiều cao ô chữ
\newlength{\EZxien}\setlength{\EZxien}{0.55cm}% độ lệch ngang của cạnh xiên
\newlength{\EZvach}\setlength{\EZvach}{0.75cm}% bề rộng mỗi vạch xiên
\newlength{\EZkhe}\setlength{\EZkhe}{0.12cm}% khe trắng giữa vạch và ô chữ
\newlength{\EZdem}\setlength{\EZdem}{0.4cm}% khoảng đệm từ chữ tới cạnh xiên
\newlength{\EZle}\setlength{\EZle}{1.8cm}% lề trái/phải của chữ (bằng lề thân bài)
\newlength{\EZrongmin}\setlength{\EZrongmin}{4.5cm}% ô chữ rộng tối thiểu
\newlength{\EZrong}% bề rộng chữ trong ô (tính tự động)
%--- Font
\newcommand{\EZfontO}{\fontsize{10.5pt}{12pt}\selectfont\bfseries}% chữ trong ô
\newcommand{\EZfontKH}{\fontfamily{qag}\fontsize{11pt}{12pt}\selectfont\bfseries}% \khauhieu cạnh logo
\newcommand{\EZfontTrang}{\fontsize{11.5pt}{12pt}\selectfont\bfseries}% chữ "Trang" và số trang
%--- Đo bề rộng chữ trong ô
\newcommand{\EZdorong}[1]{%
	\settowidth{\EZrong}{\EZfontO #1}%
	\ifdim\EZrong<\EZrongmin\setlength{\EZrong}{\EZrongmin}\fi}
%--- Logo + chữ trắng \khauhieu (dựng sẵn một lần khi vào \begin{document})
\newsavebox{\EZlogobox}
\AtBeginDocument{\sbox{\EZlogobox}{%
		\raisebox{-0.5\height}{\eduzylogo[1cm]}\hspace{0.15cm}%
		\raisebox{-0.5\height}{\color{white}\EZfontKH\khauhieu}}}
%--- Hai vạch xiên + khe trắng. #1: điểm gốc ở mép dưới, #2: hướng (1 sang phải, -1 sang trái),
%    #3: hướng nghiêng (1 kiểu "/", -1 kiểu "\")
\newcommand{\EZvachxien}[3]{%
	\fill[white] (#1) --++({(#3)*\EZxien},\EZo) --++({(#2)*\EZkhe},0) --++({-(#3)*\EZxien},-\EZo) --cycle;
	\fill[ezVach2] ([xshift=(#2)*\EZkhe]#1) --++({(#3)*\EZxien},\EZo) --++({(#2)*\EZvach},0) --++({-(#3)*\EZxien},-\EZo) --cycle;
	\fill[ezVach1] ([xshift=(#2)*(\EZkhe+\EZvach)]#1) --++({(#3)*\EZxien},\EZo) --++({(#2)*\EZvach},0) --++({-(#3)*\EZxien},-\EZo) --cycle;
	\draw[white,line width=0.6pt] ([xshift=(#2)*(\EZkhe+\EZvach)]#1) --++({(#3)*\EZxien},\EZo);
}
%--- Chữ trong ô (căn giữa theo chiều cao ô, không lệch vì dấu tiếng Việt)
\tikzset{ezchu/.style={inner sep=0pt,text=ezNen,text height=0.72em,text depth=0.22em}}

%====================== LOGO CHÌM GIỮA TRANG ======================
%  Vẽ cùng đầu trang nên nằm dưới chữ của thân bài, có trên mọi trang.
%  Tắt: \logochimtat (từ trang đó trở đi), bật lại: \logochimbat.
%  Đổi cỡ, độ đậm: khai báo trước \input, ví dụ \def\EZchimcao{9cm} \def\EZchimdo{0.1}
\providecommand{\EZchimcao}{10cm}% chiều cao logo chìm
\providecommand{\EZchimdo}{0.12}% độ đậm (0 = trong suốt, 1 = rõ hẳn)
\newif\ifEZchim \EZchimtrue
\newcommand{\logochimtat}{\EZchimfalse}
\newcommand{\logochimbat}{\EZchimtrue}
%  Dựng sẵn logo mờ một lần khi vào \begin{document}; transparency group làm mờ cả logo
%  một lượt nên không lộ vệt ở chỗ các dải màu chồng nhau.
\newsavebox{\EZchimtho}
\newsavebox{\EZchimbox}
\AtBeginDocument{%
	\sbox{\EZchimtho}{\eduzylogo[\EZchimcao]}%
	\sbox{\EZchimbox}{\begin{tikzpicture}
			\begin{scope}[transparency group,opacity=\EZchimdo]
				\node[inner sep=0pt] {\usebox{\EZchimtho}};
			\end{scope}
	\end{tikzpicture}}}
\newcommand{\EZvechim}{%
	\ifEZchim
		\node[inner sep=0pt] at (current page.center) {\usebox{\EZchimbox}};
	\fi}

%====================== ĐẦU TRANG - TRANG LẺ ======================
%  (logo bên trái, ô chữ \tenchuong bên phải)
\newcommand{\myfancyheadOdd}{%
	\EZdorong{\tenchuong}%
	\begin{tikzpicture}[remember picture,overlay]
		\coordinate (T) at (current page.north west);
		\coordinate (P) at (current page.north east);
		%--- logo chìm giữa trang
		\EZvechim
		%--- dải nền
		\fill[ezNen] (T) rectangle ([yshift=-\EZdai]P);
		%--- ô chữ: X là góc dưới bên trái, cạnh trái xiên kiểu "/"
		\coordinate (X) at ([xshift=-\EZle-\EZrong-\EZdem-\EZxien,yshift=-\EZdai]P);
		\fill[ezO] (X) --++(\EZxien,\EZo) -- ([yshift=-\EZdai+\EZo]P) -- ([yshift=-\EZdai]P) --cycle;
		\EZvachxien{X}{-1}{1}
		\node[ezchu,anchor=east] at ([xshift=-\EZle,yshift=-\EZdai+0.5*\EZo]P) {\EZfontO\tenchuong};
		%--- logo + tên
		\node[inner sep=0pt,anchor=west] at ([xshift=\EZle,yshift=-0.5*\EZdai]T) {\usebox{\EZlogobox}};
	\end{tikzpicture}%
}

%====================== ĐẦU TRANG - TRANG CHẴN =====================
%  (soi gương: ô chữ \tenchuong bên trái, logo bên phải)
\newcommand{\myfancyheadEven}{%
	\EZdorong{\tenchuong}%
	\begin{tikzpicture}[remember picture,overlay]
		\coordinate (T) at (current page.north west);
		\coordinate (P) at (current page.north east);
		%--- logo chìm giữa trang
		\EZvechim
		%--- dải nền
		\fill[ezNen] (T) rectangle ([yshift=-\EZdai]P);
		%--- ô chữ: X là góc dưới bên phải, cạnh phải xiên kiểu "\"
		\coordinate (X) at ([xshift=\EZle+\EZrong+\EZdem+\EZxien,yshift=-\EZdai]T);
		\fill[ezO] (X) --++(-\EZxien,\EZo) -- ([yshift=-\EZdai+\EZo]T) -- ([yshift=-\EZdai]T) --cycle;
		\EZvachxien{X}{1}{-1}
		\node[ezchu,anchor=west] at ([xshift=\EZle,yshift=-\EZdai+0.5*\EZo]T) {\EZfontO\tenchuong};
		%--- logo + tên
		\node[inner sep=0pt,anchor=east] at ([xshift=-\EZle,yshift=-0.5*\EZdai]P) {\usebox{\EZlogobox}};
	\end{tikzpicture}%
}

%====================== CHÂN TRANG - TRANG LẺ ======================
%  (ô chữ \thuonghieu bên trái, "Trang (số)" bên phải)
\newcommand{\myfancyfootOdd}{%
	\EZdorong{\thuonghieu}%
	\begin{tikzpicture}[remember picture,overlay]
		\coordinate (T) at (current page.south west);
		\coordinate (P) at (current page.south east);
		%--- dải nền
		\fill[ezNen] (T) rectangle ([yshift=\EZdai]P);
		%--- ô chữ sát mép trên dải: X là góc dưới bên phải, cạnh phải xiên kiểu "/"
		\coordinate (X) at ([xshift=\EZle+\EZrong+\EZdem,yshift=\EZdai-\EZo]T);
		\fill[ezO] ([yshift=\EZdai-\EZo]T) -- (X) --++(\EZxien,\EZo) -- ([yshift=\EZdai]T) --cycle;
		\EZvachxien{X}{1}{1}
		\node[ezchu,anchor=west] at ([xshift=\EZle,yshift=\EZdai-0.5*\EZo]T) {\EZfontO\thuonghieu};
		%--- số trang
		\node[circle,draw=white,line width=0.9pt,text=white,inner sep=0pt,minimum size=0.62cm,anchor=east]
		(so) at ([xshift=-\EZle,yshift=0.5*\EZdai]P) {\EZfontTrang\thepage};
		\node[inner sep=0pt,text=white,anchor=base east] (tr) at ([xshift=-0.15cm]so.base west) {\EZfontTrang Trang};
	\end{tikzpicture}%
}

%====================== CHÂN TRANG - TRANG CHẴN ====================
%  (soi gương: "Trang (số)" bên trái, ô chữ \tacgia bên phải)
\newcommand{\myfancyfootEven}{%
	\EZdorong{\tacgia}%
	\begin{tikzpicture}[remember picture,overlay]
		\coordinate (T) at (current page.south west);
		\coordinate (P) at (current page.south east);
		%--- dải nền
		\fill[ezNen] (T) rectangle ([yshift=\EZdai]P);
		%--- ô chữ sát mép trên dải: X là góc dưới bên trái, cạnh trái xiên kiểu "\"
		\coordinate (X) at ([xshift=-\EZle-\EZrong-\EZdem,yshift=\EZdai-\EZo]P);
		\fill[ezO] ([yshift=\EZdai-\EZo]P) -- (X) --++(-\EZxien,\EZo) -- ([yshift=\EZdai]P) --cycle;
		\EZvachxien{X}{-1}{-1}
		\node[ezchu,anchor=east] at ([xshift=-\EZle,yshift=\EZdai-0.5*\EZo]P) {\EZfontO\tacgia};
		%--- số trang
		%    (vòng tròn ẩn so0 nằm giữa dải để lấy đường chân chữ, giống trang lẻ)
		\node[circle,inner sep=0pt,minimum size=0.62cm] (so0) at ([yshift=0.5*\EZdai]T) {\EZfontTrang\phantom{\thepage}};
		\node[inner sep=0pt,text=white,anchor=base west] (tr) at ($(so0.base -| T)+(\EZle,0)$) {\EZfontTrang Trang};
		\node[circle,draw=white,line width=0.9pt,text=white,inner sep=0pt,minimum size=0.62cm,anchor=base west]
		(so) at ([xshift=0.15cm]tr.base east) {\EZfontTrang\thepage};
	\end{tikzpicture}%
}

%--- Giữ tương thích với các file cũ đang gọi \myfancyhead, \myfancyfoot
\newcommand{\myfancyhead}{\myfancyheadOdd}
\newcommand{\myfancyfoot}{\myfancyfootOdd}

%======================= GẮN VÀO TRANG =============================
\pagestyle{fancy}
\fancyhf{}
\fancyhead[LO]{\myfancyheadOdd}
\fancyhead[RE]{\myfancyheadEven}
\fancyfoot[LO]{\myfancyfootOdd}
\fancyfoot[RE]{\myfancyfootEven}
\renewcommand{\headrulewidth}{0pt}
\renewcommand{\footrulewidth}{0pt}
\fancypagestyle{plain}{%
	\fancyhf{}
	\fancyhead[LO]{\myfancyheadOdd}
	\fancyhead[RE]{\myfancyheadEven}
	\fancyfoot[LO]{\myfancyfootOdd}
	\fancyfoot[RE]{\myfancyfootEven}
	\renewcommand{\headrulewidth}{0pt}
	\renewcommand{\footrulewidth}{0pt}
}
%=================== HẾT PHẦN GỘP TỪ eduzy-de.tex ===================
\newcommand{\indam}[1]{
	{\fontsize{12pt}{9pt}\sffamily \selectfont\bfseries\color{gold1} #1}
}
\newcommand{\ind}[1]{
	{\fontsize{14pt}{9pt}\fontfamily{qag} \selectfont\bfseries\color{gold1} #1}
}
\renewcommand{\FalseEX}{\stepcounter{dapan}\textbf{\selectfont\color{gold1}\Alph{dapan}.}\ignorespaces}
\renewtheorem{ex}{\color{gold1}Câu}
\renewcommand{\baselinestretch}{1.1}
\newcommand{\hoac}[1]{
	\left[\begin{aligned}#1\end{aligned}\right.}
\newcommand{\heva}[1]{
	\left\{\begin{aligned}#1\end{aligned}\right.}
\usepackage{esvect}
\def\vec{\vv}
\def\vt{\vv}
\def\overrightarrow{\vv}
%====================== LOGO EDUZY VẼ BẰNG TIKZ ======================
%  Đường viền dò từ logo gốc; dải màu vàng kim đo trực tiếp trên ảnh gốc
%  (nghiêng 21 độ, tô bằng 80 dải hẹp nên hiển thị giống nhau trên mọi trình xem PDF).
%  Dùng: \eduzylogo          -> cao 2.4cm (mặc định)
%        \eduzylogo[1.5cm]   -> cao 1.5cm
\definecolor{ezText}{RGB}{245,200,46}
\newcommand{\eduzylogo}[1][2.4cm]{%
	\resizebox{!}{#1}{%
		\begin{tikzpicture}[x=0.01cm,y=0.01cm]
			\path[use as bounding box] (0,0) rectangle (224,-224);
			% mũ tốt nghiệp + vương miện
			\begin{scope}
				\clip
				(119.4,-176.1)..controls(118.9,-176.1)and(116.8,-176)..(114.6,-176)
				..controls(101.4,-176)and(85.8,-175.9)..(83.7,-175.7)
				..controls(81.8,-175.6)and(81,-175.6)..(80.2,-175.7)
				..controls(78.8,-175.9)and(77.7,-175.9)..(76.2,-175.7)
				..controls(74.9,-175.5)and(70.8,-175.5)..(64.3,-175.7)
				..controls(59.1,-175.9)and(58.3,-175.7)..(58.6,-173.7)
				..controls(58.6,-172.8)and(59.6,-169.6)..(60.3,-168)
				..controls(60.8,-166.9)and(61.4,-165.8)..(62.3,-165)
				..controls(63.1,-164.2)and(64.3,-163.4)..(65,-163.2)
				..controls(65.5,-163.1)and(72.3,-163)..(76.6,-163.1)
				..controls(77.9,-163.1)and(80.1,-163.1)..(81.4,-163)
				..controls(82.8,-162.9)and(86.3,-162.9)..(89.2,-162.9)
				..controls(93,-162.9)and(94.9,-162.8)..(96,-162.7)
				..controls(97.2,-162.5)and(101,-162.5)..(114,-162.4)
				..controls(123,-162.3)and(134.9,-162.2)..(140.2,-162.1)
				..controls(150.1,-162)and(152.2,-162.1)..(153.6,-162.4)
				..controls(154.4,-162.6)and(157.1,-163.8)..(158.1,-164.4)
				..controls(158.8,-164.8)and(159.7,-165.7)..(159.9,-166.3)
				..controls(160,-166.5)and(160.2,-167.1)..(160.2,-167.6)
				..controls(160.3,-168)and(160.6,-169.4)..(160.9,-170.6)
				..controls(161.2,-171.7)and(161.5,-173.1)..(161.6,-173.6)
				..controls(161.8,-174.4)and(161.8,-174.5)..(161.5,-174.9)
				..controls(161.2,-175.6)and(160.8,-175.9)..(159.9,-176)
				..controls(159.2,-176.1)and(154.1,-176.2)..(151.5,-176.2)
				..controls(148.1,-176.1)and(122.6,-176.1)..(121.6,-176.2)
				..controls(120.8,-176.2)and(119.8,-176.2)..(119.4,-176.1)--cycle
				(130.4,-160)..controls(129,-160)and(119.8,-160)..(110.1,-159.9)
				..controls(100.3,-159.8)and(88.5,-159.7)..(83.9,-159.6)
				..controls(77.8,-159.5)and(75.3,-159.5)..(74.5,-159.6)
				..controls(73.9,-159.7)and(72.5,-159.8)..(71.3,-159.8)
				..controls(69.6,-159.8)and(69,-159.8)..(68.5,-159.6)
				..controls(67.6,-159.3)and(66.8,-158.6)..(66.2,-157.6)
				..controls(66,-157.3)and(65.6,-156.5)..(65.3,-156.1)
				..controls(64.6,-155.1)and(64.5,-154.9)..(63.7,-152.9)
				..controls(63.2,-151.6)and(62.9,-150.9)..(61.8,-149.6)
				..controls(61.5,-149.2)and(61.1,-148.5)..(60.9,-147.9)
				..controls(60.7,-147.4)and(60.3,-146.6)..(60.1,-146.1)
				..controls(59.9,-145.5)and(59.5,-144.9)..(59,-144.3)
				..controls(58.5,-143.7)and(58.2,-143.2)..(57.9,-142.5)
				..controls(57.7,-141.9)and(57.3,-141.1)..(57,-140.7)
				..controls(56.6,-140.3)and(56.1,-139.3)..(55.8,-138.5)
				..controls(55.5,-137.8)and(55.1,-136.8)..(54.9,-136.5)
				..controls(54.7,-136.1)and(54.3,-135.2)..(53.9,-134.5)
				..controls(53.6,-133.8)and(53,-132.7)..(52.6,-132)
				..controls(51.6,-130.5)and(51.4,-130.2)..(51.1,-129)
				..controls(50.6,-127.5)and(50.6,-127.4)..(49.7,-126.3)
				..controls(49.2,-125.7)and(48.6,-124.8)..(48.3,-124.2)
				..controls(48,-123.6)and(47.5,-122.7)..(47.3,-122.2)
				..controls(47,-121.8)and(46.5,-120.7)..(46.1,-119.8)
				..controls(45.8,-118.9)and(45.2,-117.8)..(44.9,-117.2)
				..controls(44.2,-116.1)and(43.3,-114.4)..(42.3,-112.2)
				..controls(41.8,-111.3)and(41.3,-110.4)..(41.1,-110.1)
				..controls(40.5,-109.4)and(40.3,-109)..(40,-107.8)
				..controls(39.8,-107)and(39.6,-106.5)..(39.2,-106)
				..controls(38.9,-105.6)and(38.4,-104.7)..(38,-104)
				..controls(37.7,-103.4)and(37.2,-102.5)..(37,-102.1)
				..controls(36.7,-101.7)and(36.3,-100.8)..(36,-100.1)
				..controls(35.6,-98.9)and(35.3,-98.4)..(34.6,-97.6)
				..controls(34.2,-97.1)and(34.2,-96.3)..(34.6,-96)
				..controls(35,-95.7)and(35.5,-95.9)..(36.2,-96.7)
				..controls(36.5,-97)and(37.2,-97.5)..(37.6,-97.7)..controls(38,-98)and(38.6,-98.4)..(39.1,-98.9)
				..controls(39.6,-99.4)and(40.3,-100)..(40.6,-100.3)
				..controls(41.3,-100.8)and(42.4,-101.8)..(43.7,-102.9)
				..controls(44,-103.2)and(44.6,-103.7)..(45,-103.9)
				..controls(45.4,-104.2)and(45.9,-104.6)..(46.1,-104.8)
				..controls(47.1,-105.9)and(48.8,-107.1)..(50.5,-107.9)
				..controls(51.8,-108.6)and(52.3,-109)..(53.4,-110)
				..controls(54.8,-111.2)and(55.6,-111.6)..(57,-112.1)
				..controls(57.5,-112.3)and(58.5,-112.8)..(59.4,-113.3)
				..controls(60.2,-113.7)and(61.3,-114.3)..(61.9,-114.5)
				..controls(62.4,-114.7)and(63.3,-115.1)..(63.8,-115.3)
				..controls(64.4,-115.6)and(65.4,-115.9)..(66.2,-116.1)
				..controls(67,-116.3)and(68.4,-116.6)..(69.2,-116.9)
				..controls(72.3,-117.9)and(72.8,-118)..(78,-118.1)
				..controls(84,-118.3)and(86.1,-118.1)..(88.1,-117.1)
				..controls(89.4,-116.5)and(93,-115)..(93.7,-114.9)
				..controls(94.8,-114.6)and(95.4,-114.2)..(95.9,-113.4)
				..controls(96.2,-113)and(96.7,-112.5)..(97.1,-112.2)
				..controls(98,-111.6)and(101.3,-108.4)..(102.1,-107.2)
				..controls(102.4,-106.8)and(102.9,-106.2)..(103.2,-105.8)
				..controls(103.5,-105.5)and(104,-104.8)..(104.2,-104.2)
				..controls(104.5,-103.6)and(105,-102.8)..(105.3,-102.4)
				..controls(105.6,-102)and(106.1,-101.2)..(106.5,-100.6)
				..controls(107.1,-99.5)and(108.6,-97.5)..(109.5,-96.5)
				..controls(109.9,-96.1)and(110.1,-96)..(110.4,-96)
				..controls(111.1,-96)and(111.6,-96.5)..(111.8,-97.6)
				..controls(112,-98.6)and(112.3,-99.3)..(112.8,-99.9)
				..controls(113.7,-100.9)and(115.2,-103.1)..(115.7,-104.2)
				..controls(116.1,-105)and(116.5,-105.5)..(116.9,-106)
				..controls(117.3,-106.4)and(117.8,-106.9)..(118,-107.3)
				..controls(118.2,-107.6)and(118.6,-108)..(119,-108.3)
				..controls(119.6,-108.8)and(121.1,-110.3)..(121.8,-111.2)
				..controls(122.1,-111.7)and(122.5,-112)..(123.4,-112.4)
				..controls(124.8,-113.2)and(125.8,-113.8)..(127,-114.6)
				..controls(127.5,-114.9)and(128.1,-115.3)..(128.2,-115.4)
				..controls(128.4,-115.5)and(129.3,-115.8)..(130.2,-116)
				..controls(131.1,-116.2)and(132.5,-116.6)..(133.3,-116.9)
				..controls(135.7,-117.9)and(136.1,-117.9)..(142.5,-117.8)
				..controls(145.8,-117.8)and(148.4,-117.7)..(148.8,-117.6)
				..controls(149.2,-117.6)and(150.2,-117.2)..(151,-116.8)
				..controls(152,-116.4)and(152.8,-116)..(153.6,-115.9)
				..controls(154.5,-115.7)and(155.2,-115.4)..(156.3,-114.9)
				..controls(157.1,-114.5)and(158.2,-114)..(158.7,-113.8)
				..controls(159.1,-113.7)and(159.8,-113.3)..(160.2,-113)
				..controls(160.6,-112.7)and(161.3,-112.4)..(161.7,-112.2)
				..controls(162.7,-111.9)and(163.2,-111.7)..(164.1,-111.1)
				..controls(164.5,-110.8)and(165.3,-110.4)..(165.9,-110)
				..controls(166.5,-109.7)and(167.3,-109.2)..(167.6,-108.9)
				..controls(167.9,-108.6)and(168.6,-108.2)..(169,-108)
				..controls(169.4,-107.8)and(170.2,-107.3)..(170.7,-106.9)
				..controls(171.2,-106.5)and(172.1,-105.8)..(172.7,-105.4)
				..controls(173.2,-105)and(174.1,-104.4)..(174.5,-104)
				..controls(174.9,-103.7)and(175.7,-103.1)..(176.2,-102.6)
				..controls(176.7,-102.2)and(177.4,-101.5)..(177.9,-101.2)
				..controls(178.3,-100.8)and(178.8,-100.4)..(179,-100.2)
				..controls(179.2,-100)and(179.6,-99.6)..(180,-99.3)
				..controls(180.4,-99.1)and(180.9,-98.6)..(181.2,-98.4)
				..controls(181.5,-98.1)and(182.1,-97.7)..(182.5,-97.5)
				..controls(183.2,-97.2)and(183.5,-96.9)..(184.3,-95.9)
				..controls(184.6,-95.6)and(184.6,-95.6)..(184.8,-95.7)
				..controls(185,-95.8)and(185.1,-96)..(185.1,-96.2)
				..controls(185.1,-96.5)and(185,-96.7)..(184.6,-97.2)
				..controls(184.1,-97.8)and(184,-97.9)..(183.9,-98.8)
				..controls(183.7,-100.3)and(183.3,-101.3)..(182.9,-101.8)
				..controls(182.6,-102.1)and(182.3,-102.6)..(182.2,-102.9)
				..controls(182,-103.3)and(181.6,-104)..(181.2,-104.6)
				..controls(180.9,-105.2)and(180.5,-106.1)..(180.4,-106.6)
				..controls(179.9,-108)and(179.7,-108.5)..(179,-109.4)
				..controls(178.6,-109.8)and(178.2,-110.5)..(177.9,-111)
				..controls(177.7,-111.5)and(177.3,-112.3)..(177,-112.7)
				..controls(176.6,-113.2)and(176.4,-113.7)..(176.3,-114.2)
				..controls(176.2,-114.7)and(175.9,-115.2)..(175.4,-116)
				..controls(175,-116.6)and(174.4,-117.7)..(174,-118.5)
				..controls(173.7,-119.2)and(173.3,-120.1)..(173,-120.5)
				..controls(172.8,-120.9)and(172.4,-121.8)..(172.1,-122.5)
				..controls(171.8,-123.2)and(171.3,-124.1)..(171,-124.6)
				..controls(170.8,-125.1)and(170.3,-126)..(169.9,-126.6)
				..controls(169.6,-127.3)and(169.2,-128.1)..(169,-128.3)
				..controls(168.8,-128.6)and(168.5,-129.3)..(168.3,-129.8)
				..controls(168.1,-130.3)and(167.8,-131.1)..(167.5,-131.6)
				..controls(166.9,-132.5)and(164.2,-137.9)..(164.2,-138.1)
				..controls(164.2,-138.1)and(163.9,-138.8)..(163.4,-139.7)
				..controls(162.9,-140.5)and(162.4,-141.5)..(162.3,-141.9)
				..controls(161.9,-143)and(161.6,-143.4)..(160.9,-144.1)
				..controls(160.4,-144.6)and(160.3,-144.8)..(160.1,-145.4)
				..controls(160,-145.8)and(159.6,-146.8)..(159.2,-147.8)
				..controls(158.7,-148.8)and(158.1,-150)..(157.9,-150.6)
				..controls(157.6,-151.1)and(157.2,-151.9)..(156.9,-152.4)
				..controls(156.5,-152.8)and(156.1,-153.7)..(155.8,-154.4)
				..controls(155.6,-155)and(155.2,-155.9)..(154.9,-156.3)
				..controls(154.7,-156.7)and(154.4,-157.4)..(154.2,-157.9)
				..controls(154.1,-158.4)and(153.8,-159)..(153.5,-159.3)
				..controls(152.8,-160.1)and(153,-160.1)..(142.4,-160.1)
				..controls(137.2,-160.1)and(131.8,-160.1)..(130.4,-160)--cycle
				(141.5,-103.7)..controls(140.8,-103)and(139.8,-102.1)..(139.2,-101.6)
				..controls(138.5,-101.2)and(137.9,-100.6)..(137.8,-100.4)
				..controls(137,-99.5)and(134,-96.5)..(133.2,-96)
				..controls(132.7,-95.7)and(131.8,-94.9)..(131.1,-94.3)
				..controls(129.4,-92.7)and(127.9,-91.6)..(126.6,-91.1)
				..controls(125.2,-90.5)and(124.1,-90)..(122.9,-89.1)
				..controls(122.1,-88.6)and(121.3,-88.2)..(119.9,-87.7)
				..controls(118.8,-87.3)and(117.5,-86.8)..(116.9,-86.5)
				..controls(114.3,-85.5)and(113.7,-85.4)..(109.7,-85.5)
				..controls(105.3,-85.7)and(103.9,-85.9)..(101.6,-87)
				..controls(100.7,-87.4)and(99.4,-87.8)..(97.9,-88)
				..controls(97,-88.2)and(96.6,-88.3)..(95.8,-88.9)
				..controls(95.4,-89.2)and(94.7,-89.6)..(94.2,-89.8)
				..controls(92.9,-90.3)and(92.3,-90.7)..(90.9,-92.1)
				..controls(90.2,-92.7)and(89.4,-93.3)..(89,-93.5)..controls(87,-94.5)and(85.2,-96)..(83.6,-97.9)
				..controls(83.3,-98.2)and(82.5,-98.9)..(81.8,-99.4)
				..controls(80.7,-100.1)and(80.5,-100.3)..(80,-101.1)
				..controls(78.7,-103.1)and(77.2,-103.9)..(76.4,-103)
				..controls(75.7,-102.2)and(75.7,-103.2)..(75.7,-83.6)
				..controls(75.6,-73.6)and(75.6,-64.8)..(75.5,-64)
				..controls(75.5,-63.2)and(75.5,-61.4)..(75.6,-59.8)--(75.7,-57.1)--(76,-56.8)
				..controls(76.8,-56.2)and(78,-56.3)..(79.8,-57.2)
				..controls(80.4,-57.5)and(81.3,-57.8)..(81.8,-57.9)
				..controls(82.3,-58.1)and(83,-58.4)..(83.6,-58.8)..controls(84.3,-59.2)and(85,-59.5)..(86,-59.8)
				..controls(86.8,-60.1)and(88,-60.6)..(88.7,-60.9)
				..controls(89.3,-61.2)and(90.6,-61.7)..(91.4,-62)
				..controls(92.2,-62.3)and(93.7,-62.9)..(94.6,-63.3)
				..controls(95.6,-63.7)and(96.8,-64.1)..(97.4,-64.3)
				..controls(98,-64.5)and(99,-64.9)..(99.7,-65.2)
				..controls(100.4,-65.6)and(101.5,-66)..(102.3,-66.2)
				..controls(104.2,-66.8)and(105,-67.1)..(106.2,-67.8)
				..controls(107.9,-68.7)and(110.2,-69.4)..(111.3,-69.2)
				..controls(112.5,-69)and(115,-68.2)..(115.6,-67.8)
				..controls(117.3,-66.7)and(117.6,-66.6)..(119.5,-66.2)
				..controls(121.9,-65.8)and(122.1,-65.7)..(123,-65.1)
				..controls(124.1,-64.4)and(124.4,-64.3)..(126,-63.9)
				..controls(127,-63.7)and(127.7,-63.4)..(128.4,-63.1)
				..controls(129,-62.8)and(129.7,-62.5)..(130.6,-62.4)
				..controls(132.1,-62.1)and(134,-61.5)..(135.2,-61)
				..controls(135.7,-60.8)and(136.7,-60.5)..(137.3,-60.3)
				..controls(138.4,-60)and(139.3,-59.7)..(141.1,-59)
				..controls(141.6,-58.7)and(142,-58.6)..(142.5,-58.6)
				..controls(143.2,-58.6)and(143.3,-58.7)..(143.6,-59)
				..controls(144,-59.5)and(144.1,-59.9)..(144,-60.8)
				..controls(144,-61.2)and(144,-71.1)..(144,-82.8)--(143.9,-104.1)--(143.7,-104.5)
				..controls(143.5,-104.7)and(143.3,-104.8)..(143,-104.9)
				..controls(142.6,-104.9)and(142.5,-104.8)..(141.5,-103.7)--cycle
				(27.8,-91.4)..controls(27.1,-90.9)and(26.3,-90)..(26.1,-89.5)
				..controls(26,-89.2)and(25.8,-88.6)..(25.7,-87.9)
				..controls(25.4,-86.2)and(25.6,-84.8)..(26.5,-83.5)
				..controls(27.4,-82.3)and(28.3,-82)..(29.4,-82.3)
				..controls(31.1,-82.8)and(32.1,-84.9)..(32,-87.4)
				..controls(31.9,-89.3)and(31.2,-90.6)..(29.9,-91.4)
				..controls(29.1,-91.9)and(28.6,-91.9)..(27.8,-91.4)--cycle
				(190.1,-91.3)..controls(189.7,-91.2)and(189.3,-90.9)..(189.1,-90.7)
				..controls(188.3,-89.8)and(187.6,-87.7)..(187.6,-86.2)
				..controls(187.6,-85.5)and(187.7,-85.2)..(188.2,-84.4)
				..controls(189.1,-82.6)and(190.1,-81.8)..(191.3,-82.1)
				..controls(192,-82.3)and(193.5,-83.7)..(193.8,-84.4)
				..controls(194.1,-85.1)and(194.2,-87.1)..(194,-88.3)
				..controls(193.8,-89.8)and(193.3,-90.5)..(192.3,-91.3)
				..controls(191.7,-91.7)and(190.9,-91.7)..(190.1,-91.3)--cycle
				(173.7,-83.4)..controls(172.2,-82.5)and(172,-82.5)..(170.6,-82.4)
				..controls(169.7,-82.4)and(168.8,-82.4)..(168.3,-82.5)
				..controls(167.2,-82.7)and(166.9,-82.7)..(166.6,-82.2)
				..controls(166.3,-81.8)and(166.3,-81)..(166.8,-80)
				..controls(167.2,-78.9)and(167.4,-77.4)..(167.7,-74.1)
				..controls(167.8,-72.8)and(167.9,-71.8)..(167.9,-71.7)
				..controls(168,-71.7)and(168,-70.2)..(168,-68.4)..controls(168,-65.4)and(168,-64.9)..(167.7,-63)
				..controls(167.3,-60.5)and(167.1,-59.9)..(166.4,-58.4)
				..controls(166.1,-57.7)and(165.8,-56.9)..(165.7,-56.6)
				..controls(165.7,-56.2)and(165.4,-55.6)..(165.2,-55.2)
				..controls(164.6,-53.8)and(161.9,-50.7)..(161,-50.3)
				..controls(159.8,-49.7)and(157.9,-50)..(155.2,-51.2)
				..controls(154.7,-51.4)and(153.4,-51.8)..(152.2,-52.1)
				..controls(150.8,-52.4)and(149.4,-52.9)..(148.6,-53.2)
				..controls(147.8,-53.6)and(146.7,-54)..(146,-54.2)
				..controls(145.4,-54.4)and(144.5,-54.7)..(144,-54.9)
				..controls(143.6,-55.2)and(142.7,-55.5)..(141.8,-55.7)
				..controls(141,-55.9)and(139.6,-56.3)..(138.7,-56.6)
				..controls(137.8,-57)and(136.3,-57.4)..(135.3,-57.6)
				..controls(133.8,-57.9)and(133.2,-58.1)..(132.1,-58.7)
				..controls(131.1,-59.2)and(130.3,-59.4)..(129.4,-59.6)
				..controls(128.7,-59.8)and(127.8,-60)..(127.4,-60.1)
				..controls(127,-60.3)and(126,-60.5)..(125.1,-60.7)
				..controls(123.2,-61.1)and(121.7,-61.5)..(120.2,-62.1)
				..controls(119.7,-62.3)and(118.1,-62.8)..(116.7,-63.2)
				..controls(114.6,-63.9)and(114,-64.1)..(113,-64.7)
				..controls(111.9,-65.4)and(111.8,-65.4)..(111,-65.5)
				..controls(109.9,-65.5)and(109.3,-65.4)..(107.7,-64.7)
				..controls(107,-64.4)and(106,-64)..(105.3,-63.9)
				..controls(103.5,-63.5)and(102.2,-63)..(101.2,-62.5)
				..controls(100.1,-61.8)and(99.3,-61.5)..(97.8,-61.1)
				..controls(97.2,-61)and(96.2,-60.7)..(95.7,-60.4)
				..controls(95.1,-60.1)and(94.2,-59.8)..(93.5,-59.6)
				..controls(92.9,-59.4)and(91.7,-59)..(91,-58.7)..controls(90.2,-58.4)and(88.9,-58)..(88,-57.7)
				..controls(87.2,-57.5)and(86.1,-57.1)..(85.5,-56.8)
				..controls(85,-56.6)and(84.1,-56.2)..(83.5,-56)..controls(82.9,-55.7)and(81.8,-55.2)..(81,-54.9)
				..controls(80.1,-54.5)and(79.1,-54.2)..(78.7,-54.1)
				..controls(77.9,-53.8)and(77.3,-53.6)..(76.2,-52.9)
				..controls(75.7,-52.6)and(75.2,-52.4)..(74.1,-52.2)
				..controls(72.5,-51.9)and(72.4,-51.8)..(71.3,-51.1)
				..controls(69.9,-50.1)and(69.4,-49.9)..(67.9,-49.8)
				..controls(67,-49.7)and(66.6,-49.6)..(66.3,-49.4)
				..controls(66.1,-49.3)and(65.7,-49)..(65.4,-48.8)
				..controls(65.1,-48.7)and(64.2,-48.4)..(63.6,-48.2)
				..controls(62.9,-48)and(61.7,-47.6)..(60.9,-47.2)
				..controls(60.1,-46.8)and(58.9,-46.4)..(58.2,-46.2)
				..controls(57.6,-46)and(56.8,-45.7)..(56.5,-45.5)
				..controls(54.9,-44.6)and(54.6,-44.5)..(53.2,-44.2)
				..controls(51.7,-43.8)and(51,-43.6)..(49.9,-42.9)
				..controls(49.5,-42.6)and(48.9,-42.4)..(47.5,-42)
				..controls(45.7,-41.6)and(45.6,-41.5)..(44.3,-40.6)
				..controls(43.5,-40.1)and(42.8,-39.6)..(42.7,-39.4)
				..controls(42.4,-38.7)and(43,-38.1)..(44.5,-37.6)..controls(44.9,-37.5)and(46,-37)..(46.9,-36.6)
				..controls(47.9,-36.1)and(48.8,-35.8)..(49.6,-35.6)
				..controls(50.5,-35.4)and(51.3,-35.2)..(52.4,-34.6)
				..controls(53.3,-34.1)and(54.5,-33.7)..(54.9,-33.5)
				..controls(55.4,-33.4)and(56.4,-33)..(57.2,-32.6)
				..controls(58.6,-32)and(60.1,-31.5)..(61.6,-31.2)
				..controls(62.3,-31.1)and(62.7,-31)..(63.2,-30.7)
				..controls(64.1,-30)and(64.3,-29.9)..(65.3,-29.7)
				..controls(65.7,-29.6)and(66.5,-29.2)..(67.1,-28.9)
				..controls(67.6,-28.6)and(68.8,-28.2)..(69.6,-27.9)
				..controls(70.5,-27.7)and(71.6,-27.3)..(72.1,-27)
				..controls(72.6,-26.7)and(73.8,-26.2)..(74.8,-25.9)
				..controls(75.7,-25.5)and(76.7,-25.1)..(77.1,-25)
				..controls(77.4,-24.8)and(78.5,-24.4)..(79.6,-24.1)
				..controls(80.6,-23.8)and(82,-23.2)..(82.8,-22.9)
				..controls(83.5,-22.5)and(84.8,-22)..(85.8,-21.8)
				..controls(86.7,-21.5)and(87.8,-21.1)..(88.1,-20.9)
				..controls(88.4,-20.7)and(89.3,-20.3)..(90.1,-20.1)
				..controls(90.8,-19.8)and(91.9,-19.4)..(92.4,-19.1)
				..controls(93.5,-18.5)and(94.4,-18.2)..(95.4,-18.1)
				..controls(95.7,-18)and(96.4,-17.8)..(96.9,-17.5)
				..controls(97.5,-17.3)and(98.4,-16.9)..(99,-16.6)
				..controls(99.7,-16.3)and(100.9,-15.9)..(101.9,-15.6)
				..controls(102.9,-15.3)and(104.4,-14.7)..(105.3,-14.3)
				..controls(108.1,-13.2)and(108.5,-13.1)..(110.1,-13.1)
				..controls(111.4,-13.1)and(111.5,-13.1)..(112.4,-13.5)
				..controls(113.2,-14)and(115,-14.8)..(116.6,-15.4)
				..controls(116.9,-15.5)and(117.9,-15.8)..(118.8,-16.1)
				..controls(120.9,-16.6)and(121.5,-16.9)..(122.6,-17.9)
				..controls(123,-18.3)and(123.7,-18.8)..(123.9,-18.9)
				..controls(124.6,-19.4)and(126,-19.8)..(127.9,-20.2)
				..controls(129.6,-20.5)and(129.8,-20.6)..(131.1,-21.4)
				..controls(131.5,-21.7)and(132.4,-22.1)..(133.1,-22.3)
				..controls(136.1,-23.2)and(138.3,-24.1)..(140.2,-25)
				..controls(141,-25.5)and(141.5,-25.6)..(142.5,-25.8)
				..controls(145.1,-26.3)and(146,-26.6)..(149.4,-28.1)
				..controls(151.3,-28.9)and(153.5,-29.8)..(154.3,-30.1)
				..controls(155.1,-30.4)and(156.2,-30.8)..(156.6,-31)
				..controls(157.1,-31.2)and(158.2,-31.6)..(159.1,-31.8)
				..controls(160.6,-32.2)and(160.9,-32.3)..(161.7,-32.9)
				..controls(162.4,-33.3)and(162.9,-33.6)..(163.7,-33.8)
				..controls(164.3,-34)and(164.8,-34.1)..(164.9,-34.1)
				..controls(165,-34.1)and(165.6,-34.4)..(166.4,-34.8)
				..controls(167.3,-35.2)and(168.1,-35.4)..(169,-35.6)
				..controls(170,-35.8)and(170.5,-36)..(171.7,-36.6)
				..controls(172.4,-37)and(173.3,-37.4)..(173.6,-37.4)
				..controls(175.6,-37.7)and(176.2,-38)..(176.7,-39.1)
				..controls(176.9,-39.6)and(176.9,-39.7)..(176.7,-40.4)
				..controls(176.4,-41.6)and(175.8,-42.7)..(175.2,-43.4)
				..controls(174.5,-44.1)and(172.6,-45.5)..(171.7,-45.8)
				..controls(170.6,-46.2)and(169.1,-46)..(168.2,-45.4)
				..controls(167.9,-45.2)and(167.3,-44.8)..(166.9,-44.4)
				..controls(166.4,-44.1)and(165.8,-43.6)..(165.4,-43.2)
				..controls(165.1,-42.9)and(164.4,-42.3)..(163.9,-41.9)
				..controls(163.3,-41.5)and(162.7,-40.9)..(162.4,-40.5)
				..controls(161.4,-39.2)and(159.9,-38)..(158.8,-37.8)
				..controls(158.4,-37.8)and(158.3,-37.7)..(158.2,-37.2)
				..controls(158,-36.8)and(157.4,-36.1)..(157.1,-36.1)
				..controls(157,-36.1)and(156.7,-36.1)..(156.5,-36)
				..controls(156.1,-35.9)and(155.9,-35.7)..(155.7,-35.2)
				..controls(155.3,-34.6)and(155.1,-34.4)..(154,-34.1)
				..controls(153.6,-34)and(153.2,-33.7)..(152.9,-33.5)
				..controls(152.4,-32.9)and(151.6,-32.4)..(151.3,-32.4)
				..controls(150.9,-32.4)and(150.4,-32.6)..(150.4,-32.9)
				..controls(150.4,-33.2)and(150.8,-33.4)..(151.4,-33.6)
				..controls(152.1,-33.7)and(152.8,-34.1)..(153.5,-34.9)
				..controls(154.1,-35.5)and(154.3,-35.6)..(155,-35.8)
				..controls(155.7,-36)and(155.7,-36)..(156,-36.6)
				..controls(156.3,-37.3)and(156.7,-37.7)..(157.5,-37.8)
				..controls(158.1,-37.9)and(158.2,-38)..(158.4,-38.8)
				..controls(158.5,-39.4)and(158.6,-39.6)..(159.6,-40.4)
				..controls(160.2,-41)and(160.9,-41.6)..(161.3,-41.8)
				..controls(161.7,-42)and(161.9,-42.3)..(162.2,-42.8)
				..controls(162.4,-43.2)and(162.8,-43.8)..(163.1,-44.1)
				..controls(163.9,-44.8)and(165.1,-46.7)..(165.7,-48)
				..controls(166.2,-49.3)and(166.4,-49.7)..(167.2,-50.5)
				..controls(167.5,-50.8)and(167.9,-51.4)..(168.1,-51.8)
				..controls(168.3,-52.2)and(168.8,-53)..(169.1,-53.6)
				..controls(169.6,-54.3)and(169.8,-54.9)..(169.9,-55.6)
				..controls(170.4,-57.8)and(170.5,-58.1)..(171.2,-59.5)
				..controls(171.6,-60.3)and(172,-61.5)..(172.2,-62.1)
				..controls(172.3,-62.8)and(172.7,-63.7)..(172.9,-64.1)
				..controls(173.2,-64.6)and(173.4,-65.3)..(173.4,-65.6)
				..controls(173.8,-68.8)and(174.2,-70.6)..(174.9,-72.6)
				..controls(175.5,-74.5)and(175.9,-77.1)..(176.1,-80.2)
				..controls(176.3,-83.4)and(175.5,-84.4)..(173.7,-83.4)--cycle;
				\definecolor{ezs}{RGB}{247,179,26}\fill[ezs](-5.7,-156.6)--(-3.4,-157.5)--(61.9,12.7)--(59.7,13.6)--cycle;
				\definecolor{ezs}{RGB}{247,179,22}\fill[ezs](-3.8,-157.3)--(-1.5,-158.2)--(63.9,12)--(61.6,12.9)--cycle;
				\definecolor{ezs}{RGB}{248,179,19}\fill[ezs](-1.8,-158.1)--(0.4,-158.9)--(65.8,11.2)--(63.5,12.1)--cycle;
				\definecolor{ezs}{RGB}{248,178,16}\fill[ezs](0.1,-158.8)--(2.4,-159.7)--(67.7,10.5)--(65.4,11.4)--cycle;
				\definecolor{ezs}{RGB}{249,178,15}\fill[ezs](2,-159.5)--(4.3,-160.4)--(69.6,9.8)--(67.3,10.7)--cycle;
				\definecolor{ezs}{RGB}{251,177,17}\fill[ezs](3.9,-160.3)--(6.2,-161.1)--(71.5,9)--(69.2,9.9)--cycle;
				\definecolor{ezs}{RGB}{252,179,18}\fill[ezs](5.8,-161)--(8.1,-161.9)--(73.4,8.3)--(71.1,9.2)--cycle;
				\definecolor{ezs}{RGB}{252,182,17}\fill[ezs](7.7,-161.7)--(10,-162.6)--(75.3,7.6)--(73,8.5)--cycle;
				\definecolor{ezs}{RGB}{252,185,17}\fill[ezs](9.6,-162.5)--(11.9,-163.4)--(77.2,6.8)--(75,7.7)--cycle;
				\definecolor{ezs}{RGB}{251,189,18}\fill[ezs](11.5,-163.2)--(13.8,-164.1)--(79.2,6.1)--(76.9,7)--cycle;
				\definecolor{ezs}{RGB}{251,191,22}\fill[ezs](13.4,-163.9)--(15.7,-164.8)--(81.1,5.4)--(78.8,6.2)--cycle;
				\definecolor{ezs}{RGB}{252,193,25}\fill[ezs](15.4,-164.7)--(17.6,-165.6)--(83,4.6)--(80.7,5.5)--cycle;
				\definecolor{ezs}{RGB}{253,195,27}\fill[ezs](17.3,-165.4)--(19.6,-166.3)--(84.9,3.9)--(82.6,4.8)--cycle;
				\definecolor{ezs}{RGB}{253,196,31}\fill[ezs](19.2,-166.1)--(21.5,-167)--(86.8,3.2)--(84.5,4)--cycle;
				\definecolor{ezs}{RGB}{254,199,35}\fill[ezs](21.1,-166.9)--(23.4,-167.8)--(88.7,2.4)--(86.4,3.3)--cycle;
				\definecolor{ezs}{RGB}{253,201,38}\fill[ezs](23,-167.6)--(25.3,-168.5)--(90.6,1.7)--(88.3,2.6)--cycle;
				\definecolor{ezs}{RGB}{253,203,40}\fill[ezs](24.9,-168.3)--(27.2,-169.2)--(92.5,1)--(90.3,1.8)--cycle;
				\definecolor{ezs}{RGB}{253,205,42}\fill[ezs](26.8,-169.1)--(29.1,-170)--(94.4,0.2)--(92.2,1.1)--cycle;
				\definecolor{ezs}{RGB}{253,207,45}\fill[ezs](28.7,-169.8)--(31,-170.7)--(96.4,-0.5)--(94.1,0.4)--cycle;
				\definecolor{ezs}{RGB}{253,210,49}\fill[ezs](30.7,-170.5)--(32.9,-171.4)--(98.3,-1.2)--(96,-0.4)--cycle;
				\definecolor{ezs}{RGB}{253,212,54}\fill[ezs](32.6,-171.3)--(34.9,-172.2)--(100.2,-2)--(97.9,-1.1)--cycle;
				\definecolor{ezs}{RGB}{253,214,59}\fill[ezs](34.5,-172)--(36.8,-172.9)--(102.1,-2.7)--(99.8,-1.8)--cycle;
				\definecolor{ezs}{RGB}{253,216,63}\fill[ezs](36.4,-172.7)--(38.7,-173.6)--(104,-3.4)--(101.7,-2.6)--cycle;
				\definecolor{ezs}{RGB}{253,218,65}\fill[ezs](38.3,-173.5)--(40.6,-174.4)--(105.9,-4.2)--(103.6,-3.3)--cycle;
				\definecolor{ezs}{RGB}{254,220,66}\fill[ezs](40.2,-174.2)--(42.5,-175.1)--(107.8,-4.9)--(105.5,-4)--cycle;
				\definecolor{ezs}{RGB}{254,222,68}\fill[ezs](42.1,-175)--(44.4,-175.8)--(109.7,-5.6)--(107.5,-4.8)--cycle;
				\definecolor{ezs}{RGB}{254,224,71}\fill[ezs](44,-175.7)--(46.3,-176.6)--(111.7,-6.4)--(109.4,-5.5)--cycle;
				\definecolor{ezs}{RGB}{254,226,75}\fill[ezs](46,-176.4)--(48.2,-177.3)--(113.6,-7.1)--(111.3,-6.2)--cycle;
				\definecolor{ezs}{RGB}{254,228,79}\fill[ezs](47.9,-177.2)--(50.1,-178)--(115.5,-7.8)--(113.2,-7)--cycle;
				\definecolor{ezs}{RGB}{254,230,82}\fill[ezs](49.8,-177.9)--(52.1,-178.8)--(117.4,-8.6)--(115.1,-7.7)--cycle;
				\definecolor{ezs}{RGB}{254,232,85}\fill[ezs](51.7,-178.6)--(54,-179.5)--(119.3,-9.3)--(117,-8.4)--cycle;
				\definecolor{ezs}{RGB}{254,234,88}\fill[ezs](53.6,-179.4)--(55.9,-180.2)--(121.2,-10)--(118.9,-9.2)--cycle;
				\definecolor{ezs}{RGB}{254,237,91}\fill[ezs](55.5,-180.1)--(57.8,-181)--(123.1,-10.8)--(120.8,-9.9)--cycle;
				\definecolor{ezs}{RGB}{254,239,94}\fill[ezs](57.4,-180.8)--(59.7,-181.7)--(125,-11.5)--(122.8,-10.6)--cycle;
				\definecolor{ezs}{RGB}{254,241,97}\fill[ezs](59.3,-181.6)--(61.6,-182.4)--(127,-12.2)--(124.7,-11.4)--cycle;
				\definecolor{ezs}{RGB}{254,243,101}\fill[ezs](61.2,-182.3)--(63.5,-183.2)--(128.9,-13)--(126.6,-12.1)--cycle;
				\definecolor{ezs}{RGB}{254,245,103}\fill[ezs](63.2,-183)--(65.4,-183.9)--(130.8,-13.7)--(128.5,-12.8)--cycle;
				\definecolor{ezs}{RGB}{254,245,105}\fill[ezs](65.1,-183.8)--(67.4,-184.6)--(132.7,-14.4)--(130.4,-13.6)--cycle;
				\definecolor{ezs}{RGB}{254,245,105}\fill[ezs](67,-184.5)--(69.3,-185.4)--(134.6,-15.2)--(132.3,-14.3)--cycle;
				\definecolor{ezs}{RGB}{254,244,103}\fill[ezs](68.9,-185.2)--(71.2,-186.1)--(136.5,-15.9)--(134.2,-15)--cycle;
				\definecolor{ezs}{RGB}{253,243,101}\fill[ezs](70.8,-186)--(73.1,-186.8)--(138.4,-16.6)--(136.1,-15.8)--cycle;
				\definecolor{ezs}{RGB}{253,241,98}\fill[ezs](72.7,-186.7)--(75,-187.6)--(140.3,-17.4)--(138,-16.5)--cycle;
				\definecolor{ezs}{RGB}{253,239,95}\fill[ezs](74.6,-187.4)--(76.9,-188.3)--(142.2,-18.1)--(140,-17.2)--cycle;
				\definecolor{ezs}{RGB}{253,236,92}\fill[ezs](76.5,-188.2)--(78.8,-189)--(144.2,-18.8)--(141.9,-18)--cycle;
				\definecolor{ezs}{RGB}{253,234,88}\fill[ezs](78.5,-188.9)--(80.7,-189.8)--(146.1,-19.6)--(143.8,-18.7)--cycle;
				\definecolor{ezs}{RGB}{253,232,85}\fill[ezs](80.4,-189.6)--(82.7,-190.5)--(148,-20.3)--(145.7,-19.4)--cycle;
				\definecolor{ezs}{RGB}{253,230,82}\fill[ezs](82.3,-190.4)--(84.6,-191.2)--(149.9,-21)--(147.6,-20.2)--cycle;
				\definecolor{ezs}{RGB}{253,228,78}\fill[ezs](84.2,-191.1)--(86.5,-192)--(151.8,-21.8)--(149.5,-20.9)--cycle;
				\definecolor{ezs}{RGB}{253,226,74}\fill[ezs](86.1,-191.8)--(88.4,-192.7)--(153.7,-22.5)--(151.4,-21.6)--cycle;
				\definecolor{ezs}{RGB}{253,224,70}\fill[ezs](88,-192.6)--(90.3,-193.4)--(155.6,-23.3)--(153.3,-22.4)--cycle;
				\definecolor{ezs}{RGB}{253,222,67}\fill[ezs](89.9,-193.3)--(92.2,-194.2)--(157.5,-24)--(155.3,-23.1)--cycle;
				\definecolor{ezs}{RGB}{253,220,65}\fill[ezs](91.8,-194)--(94.1,-194.9)--(159.5,-24.7)--(157.2,-23.8)--cycle;
				\definecolor{ezs}{RGB}{253,217,63}\fill[ezs](93.7,-194.8)--(96,-195.6)--(161.4,-25.5)--(159.1,-24.6)--cycle;
				\definecolor{ezs}{RGB}{253,215,60}\fill[ezs](95.7,-195.5)--(97.9,-196.4)--(163.3,-26.2)--(161,-25.3)--cycle;
				\definecolor{ezs}{RGB}{253,213,56}\fill[ezs](97.6,-196.2)--(99.9,-197.1)--(165.2,-26.9)--(162.9,-26)--cycle;
				\definecolor{ezs}{RGB}{252,211,53}\fill[ezs](99.5,-197)--(101.8,-197.8)--(167.1,-27.7)--(164.8,-26.8)--cycle;
				\definecolor{ezs}{RGB}{252,209,49}\fill[ezs](101.4,-197.7)--(103.7,-198.6)--(169,-28.4)--(166.7,-27.5)--cycle;
				\definecolor{ezs}{RGB}{252,207,45}\fill[ezs](103.3,-198.4)--(105.6,-199.3)--(170.9,-29.1)--(168.6,-28.2)--cycle;
				\definecolor{ezs}{RGB}{253,205,42}\fill[ezs](105.2,-199.2)--(107.5,-200)--(172.8,-29.9)--(170.6,-29)--cycle;
				\definecolor{ezs}{RGB}{253,202,40}\fill[ezs](107.1,-199.9)--(109.4,-200.8)--(174.7,-30.6)--(172.5,-29.7)--cycle;
				\definecolor{ezs}{RGB}{253,200,38}\fill[ezs](109,-200.6)--(111.3,-201.5)--(176.7,-31.3)--(174.4,-30.4)--cycle;
				\definecolor{ezs}{RGB}{253,198,35}\fill[ezs](111,-201.4)--(113.2,-202.2)--(178.6,-32.1)--(176.3,-31.2)--cycle;
				\definecolor{ezs}{RGB}{253,196,31}\fill[ezs](112.9,-202.1)--(115.2,-203)--(180.5,-32.8)--(178.2,-31.9)--cycle;
				\definecolor{ezs}{RGB}{253,194,28}\fill[ezs](114.8,-202.8)--(117.1,-203.7)--(182.4,-33.5)--(180.1,-32.6)--cycle;
				\definecolor{ezs}{RGB}{253,192,26}\fill[ezs](116.7,-203.6)--(119,-204.5)--(184.3,-34.3)--(182,-33.4)--cycle;
				\definecolor{ezs}{RGB}{253,189,24}\fill[ezs](118.6,-204.3)--(120.9,-205.2)--(186.2,-35)--(183.9,-34.1)--cycle;
				\definecolor{ezs}{RGB}{253,186,20}\fill[ezs](120.5,-205)--(122.8,-205.9)--(188.1,-35.7)--(185.8,-34.8)--cycle;
				\definecolor{ezs}{RGB}{253,184,17}\fill[ezs](122.4,-205.8)--(124.7,-206.7)--(190,-36.5)--(187.8,-35.6)--cycle;
				\definecolor{ezs}{RGB}{252,183,15}\fill[ezs](124.3,-206.5)--(126.6,-207.4)--(192,-37.2)--(189.7,-36.3)--cycle;
				\definecolor{ezs}{RGB}{252,181,14}\fill[ezs](126.3,-207.2)--(128.5,-208.1)--(193.9,-37.9)--(191.6,-37.1)--cycle;
				\definecolor{ezs}{RGB}{252,179,13}\fill[ezs](128.2,-208)--(130.4,-208.9)--(195.8,-38.7)--(193.5,-37.8)--cycle;
				\definecolor{ezs}{RGB}{252,178,12}\fill[ezs](130.1,-208.7)--(132.4,-209.6)--(197.7,-39.4)--(195.4,-38.5)--cycle;
				\definecolor{ezs}{RGB}{251,178,10}\fill[ezs](132,-209.4)--(134.3,-210.3)--(199.6,-40.1)--(197.3,-39.3)--cycle;
				\definecolor{ezs}{RGB}{251,178,10}\fill[ezs](133.9,-210.2)--(136.2,-211.1)--(201.5,-40.9)--(199.2,-40)--cycle;
				\definecolor{ezs}{RGB}{250,177,11}\fill[ezs](135.8,-210.9)--(138.1,-211.8)--(203.4,-41.6)--(201.1,-40.7)--cycle;
				\definecolor{ezs}{RGB}{250,178,11}\fill[ezs](137.7,-211.6)--(140,-212.5)--(205.3,-42.3)--(203.1,-41.5)--cycle;
				\definecolor{ezs}{RGB}{248,180,10}\fill[ezs](139.6,-212.4)--(141.9,-213.3)--(207.3,-43.1)--(205,-42.2)--cycle;
				\definecolor{ezs}{RGB}{247,181,9}\fill[ezs](141.5,-213.1)--(143.8,-214)--(209.2,-43.8)--(206.9,-42.9)--cycle;
				\definecolor{ezs}{RGB}{246,182,9}\fill[ezs](143.5,-213.8)--(145.7,-214.7)--(211.1,-44.5)--(208.8,-43.7)--cycle;
				\definecolor{ezs}{RGB}{246,182,10}\fill[ezs](145.4,-214.6)--(147.7,-215.5)--(213,-45.3)--(210.7,-44.4)--cycle;
			\end{scope}
			% chữ EDUZY (Times đậm, làm dày nét bằng 5 lớp lệch nhẹ; căn đúng khung chữ logo gốc)
			\foreach \dx/\dy in {-0.7/0,0.7/0,0/-0.3,0/0.3,0/0}
			\node[anchor=center,inner sep=0pt,outer sep=0pt,text=ezText,align=center] at (109.6+\dx,-193.6+\dy)
			{\resizebox{1.01cm}{0.22cm}{\fontencoding{T1}\fontfamily{ptm}\fontseries{b}\selectfont EDUZY}};
		\end{tikzpicture}%
	}%
}
%====================== THIẾT LẬP THÂN ĐỀ ======================
%--- Môi trường tự luận {bt}: ex_test 3.3.4 chưa khai báo sẵn nên tạo ở đây
%    (ex_test đã tự móc \fixshowans và \listenumerate cho bt)
\makeatletter
\@ifundefined{bt}{\newtheorem{bt}{\color{gold1}\namebt}}{}
\makeatother
\renewcommand{\sectionmark}[1]{\markright{\thesection\ #1}}
%--- Khung tiêu đề (logo + tối đa 4 dòng; dòng 3, 4 để trống {} thì bỏ qua)
%  Đề thi:    \khungtieude{THI THỬ THPT QUỐC GIA NĂM 2026}{Môn Toán -- Đề số 21}
%                         {(Đề thi gồm: ...)}{Thời gian làm bài: 90 phút (không kể thời gian phát đề).}
%  Chuyên đề: \khungtieude{CHUYÊN ĐỀ 3 -- DÃY SỐ VÀ CÔNG THỨC TRUY HỒI}
%                         {Bồi dưỡng học sinh giỏi Toán THPT}{...}{}
\newcommand{\khungtieude}[4]{%
	\begin{tcolorbox}[colframe=gold1,colback=gold5!45,boxrule=0.5mm]
		\begin{minipage}[c]{3cm}
			\centerline{\eduzylogo[2.4cm]}
		\end{minipage}\hspace{0.5cm}
		\begin{minipage}[c]{12.5cm}
			\vskip 0.15 cm
			\centerline{\ind{#1}}
			\vskip 0.15 cm
			\centerline{\textbf{#2}}
			\if\relax\detokenize{#3}\relax\else
				\vskip 0.15 cm
				\centerline{\textit{#3}}
			\fi
			\if\relax\detokenize{#4}\relax\else
				\vskip 0.15 cm
				\centerline{\textit{#4}}
			\fi
			\vskip 0.15 cm
		\end{minipage}
	\end{tcolorbox}
}
%====================== BỔ SUNG CHO BÀI HỌC, CHUYÊN ĐỀ ======================
\usepackage[shortlabels]{enumitem}
\usepackage{tabularx}
%--- Danh sách gọn: \begin{enumerate}[gon] ... \end{enumerate}
\SetEnumitemKey{gon}{leftmargin=*,itemsep=1pt,parsep=0pt,topsep=3pt}
%--- Khung lý thuyết: \begin{khung}{} ... \end{khung}; tiêu đề chỉ dùng khi thật ngắn ("Mẹo", "Lỗi thường gặp")
\newtcolorbox{khung}[1]{breakable,enhanced,colframe=gold1,colback=gold5!45,
	boxrule=0.4mm,fonttitle=\bfseries\sffamily,coltitle=white,colbacktitle=gold1,
	attach boxed title to top left={xshift=6mm,yshift=-3mm},
	boxed title style={arc=1mm,boxrule=0.3mm},title={#1}}
%--- Tiêu đề mục nhỏ trong một dạng: \muc{Dấu hiệu}, \muc{Phương pháp}, \muc{Bài tự luyện}
\makeatletter
\newcommand{\muc}[1]{\par\medskip\noindent\indam{#1}\par\nobreak\smallskip\@afterindentfalse\@afterheading}
\makeatother
%--- Ô công thức cần nhớ (căn giữa, viền vàng, vừa khít nội dung)
%     \begin{center}\congthuc{$Q(n)=(n-1)!$}\end{center}
\newtcbox{\congthuc}{on line,colframe=gold1,colback=gold5!45,boxrule=0.5pt,arc=2pt,
	left=8pt,right=8pt,top=4pt,bottom=4pt}
%--- Công tắc lời giải bài tự luyện (ví dụ mẫu luôn dùng \loigiai)
%     \loigiaiTL[{đáp số}]{lời giải đầy đủ}
%     \TLfull -> in lời giải đầy đủ;  \TLdapso -> chỉ in một dòng "Đáp số";  \TLan -> ẩn hoàn toàn
%     (\CheDo là dethi thì mọi lời giải và đáp số đều ẩn)
\makeatletter
\newif\ifTLhien   \TLhientrue
\newif\ifTLdapso  \TLdapsofalse
\newcommand{\TLfull}{\TLhientrue\TLdapsofalse}
\newcommand{\TLdapso}{\TLhientrue\TLdapsotrue}
\newcommand{\TLan}{\TLhienfalse}
\newcommand{\loigiaiTL}[2][]{%
	\ifTLhien
		\ifTLdapso
			\def\@tlarg{#1}%
			\ifx\@tlarg\@empty\else
				\ifdefstring{\kindOPTN}{dethi}{}{\par\noindent{\bfseries\color{gold1}Đáp số:} #1\par}%
			\fi
		\else
			\loigiai{#2}%
		\fi
	\fi}
\makeatother
\TLdapso% mặc định: chỉ in đáp số bài tự luyện
%====================== BẢNG ĐÚNG/SAI (\choiceTFt) ======================
%  Hàng tiêu đề nền nâu đồng chữ trắng (Phát biểu | Đ | S), các hàng tô xen kẽ kem / trắng,
%  đường kẻ vàng kim, nhãn a) b) c) d) màu nâu đồng. Bản có đáp án: cột được chọn có vòng tròn
%  nâu đồng với dấu tích trắng (câu đúng tích ở cột Đúng, câu sai tích ở cột Sai).
%  Tiêu đề ghi đủ "Đúng | Sai": \OPTN{viettat=0} sau \input (cột Đúng/Sai rộng ra, phát biểu hẹp lại).
%  Dùng bảng gốc của ex_test: \def\BangDSgoc{1} trước \input{preamble-eduzy-v2}.
\providecommand{\BangDSgoc}{0}
\ifnum\BangDSgoc=0
\usepackage{colortbl,pifont}
\colorlet{ezDSke}{gold2!60!gold1}% đường kẻ bảng (nét 0,4pt như ex_test để không đổi chiều cao bảng)
\colorlet{ezDSle}{gold5!60}% hàng a), c)
\colorlet{ezDSchan}{white}% hàng b), d)
%--- dấu chọn ở cột Đúng/Sai
%    (bề rộng 0, căn giữa ô: không làm đổi bề rộng cột của ex_test, không báo tràn lề)
\newcommand{\EZdauchon}{\makebox[0pt]{\tikz[baseline=-1ex]\node[circle,fill=gold1,inner sep=0pt,
	minimum size=1.2em,text=white,font=\small]{\ding{51}};}}
\def\tickT{\EZdauchon}\def\tickF{\EZdauchon}
\AtBeginDocument{\sbox{\ccTFTab}{\EZdauchon}}
%--- nhãn a) b) c) d) màu nâu đồng (ở bản lời giải, nhãn câu đúng vẫn đóng khung như ex_test)
\renewcommand{\FalseTF}{\stepcounter{dapan}\textbf{\color{gold1}\DapAnTF\sepTF}\ignorespaces}
%--- trước khi dựng bảng: kẻ màu vàng kim (bề rộng cột giữ nguyên như ex_test); sau bảng trả lại màu đen
\newcount\EZhangTF
\newcommand{\EZbatdauTF}{%
	\global\EZhangTF=0\relax
	\arrayrulecolor{ezDSke}}
\newcommand{\EZketthucTF}{\arrayrulecolor{black}}
%--- hàng tiêu đề nền nâu đồng, chữ trắng
\makeatletter
\newcommand{\EZo@TF}[1]{\centering\bfseries\color{white}#1}
\newcommand{\EZtieudeTF}{%
	\ifcase\cotDS
		\ifnum\cotSTT=1\relax
			\def\rowHeader{\hline\rowcolor{gold1}\EZo@TF{TT}&\EZo@TF{\phatbieu}\tabularnewline}%
		\else
			\def\rowHeader{\hline\rowcolor{gold1}\EZo@TF{\phatbieu}\tabularnewline}%
		\fi
	\or
		\ifnum\cotSTT=1\relax
			\def\rowHeader{\hline\rowcolor{gold1}\EZo@TF{TT}&\EZo@TF{\phatbieu}&\EZo@TF{\Dung}\tabularnewline}%
		\else
			\def\rowHeader{\hline\rowcolor{gold1}\EZo@TF{\phatbieu}&\EZo@TF{\Dung}\tabularnewline}%
		\fi
	\or
		\ifnum\cotSTT=1\relax
			\def\rowHeader{\hline\rowcolor{gold1}\EZo@TF{TT}&\EZo@TF{\phatbieu}&\EZo@TF{\Dung}&\EZo@TF{\Sai}\tabularnewline}%
		\else
			\def\rowHeader{\hline\rowcolor{gold1}\EZo@TF{\phatbieu}&\EZo@TF{\Dung}&\EZo@TF{\Sai}\tabularnewline}%
		\fi
	\fi}
%--- tô xen kẽ từng hàng (gọi ngay sau \hline đầu mỗi hàng)
\newcommand{\EZmauhangTF}{%
	\noalign{\global\advance\EZhangTF\@ne}%
	\ifodd\EZhangTF\rowcolor{ezDSle}\else\rowcolor{ezDSchan}\fi}
%--- gắn vào \choiceTFt / \writeChoiceTFt của ex_test
\xpretocmd{\choiceTFt}{\EZbatdauTF}{}{\typeout{[preamble-eduzy-v2] Không gắn được bảng Đúng/Sai (đầu bảng)}}
\xpatchcmd{\choiceTFt}{\ifnum\dongPB=0\def\rowHeader{}\fi}{\EZtieudeTF\ifnum\dongPB=0\def\rowHeader{}\fi}
	{}{\typeout{[preamble-eduzy-v2] Không gắn được bảng Đúng/Sai (hàng tiêu đề)}}
\xpatchcmd{\writeChoiceTFt}{\hline\xdef\dotEX{\sdotEX}}{\hline\EZmauhangTF\xdef\dotEX{\sdotEX}}
	{}{\typeout{[preamble-eduzy-v2] Không gắn được bảng Đúng/Sai (tô hàng)}}
\xpatchcmd{\writeChoiceTFt}{{\end{longtable}}}{{\end{longtable}\EZketthucTF}}
	{}{\typeout{[preamble-eduzy-v2] Không gắn được bảng Đúng/Sai (cuối longtable)}}
\xpatchcmd{\writeChoiceTFt}{{\end{xtabular}\end{center}\medskip\noindent}}{{\end{xtabular}\end{center}\EZketthucTF\medskip\noindent}}
	{}{\typeout{[preamble-eduzy-v2] Không gắn được bảng Đúng/Sai (cuối xtabular)}}
\fi
%=====================================================================
\renewcommand{\parallel}{\mathrel{/\!\!/}}
.
\providecommand{\BangDSgoc}{0}
\ifnum\BangDSgoc=0
\usepackage{colortbl,pifont}
\colorlet{ezDSke}{gold2!60!gold1}% đường kẻ bảng (nét 0,4pt như ex_test để không đổi chiều cao bảng)
\colorlet{ezDSle}{gold5!60}% hàng a), c)
\colorlet{ezDSchan}{white}% hàng b), d)
%--- dấu chọn ở cột Đúng/Sai
%    (bề rộng 0, căn giữa ô: không làm đổi bề rộng cột của ex_test, không báo tràn lề)
\newcommand{\EZdauchon}{\makebox[0pt]{\tikz[baseline=-1ex]\node[circle,fill=gold1,inner sep=0pt,
	minimum size=1.2em,text=white,font=\small]{\ding{51}};}}
\def\tickT{\EZdauchon}\def\tickF{\EZdauchon}
\AtBeginDocument{\sbox{\ccTFTab}{\EZdauchon}}
%--- nhãn a) b) c) d) màu nâu đồng (ở bản lời giải, nhãn câu đúng vẫn đóng khung như ex_test)
\renewcommand{\FalseTF}{\stepcounter{dapan}\textbf{\color{gold1}\DapAnTF\sepTF}\ignorespaces}
%--- trước khi dựng bảng: kẻ màu vàng kim (bề rộng cột giữ nguyên như ex_test); sau bảng trả lại màu đen
\newcount\EZhangTF
\newcommand{\EZbatdauTF}{%
	\global\EZhangTF=0\relax
	\arrayrulecolor{ezDSke}}
\newcommand{\EZketthucTF}{\arrayrulecolor{black}}
%--- hàng tiêu đề nền nâu đồng, chữ trắng
\makeatletter
\newcommand{\EZo@TF}[1]{\centering\bfseries\color{white}#1}
\newcommand{\EZtieudeTF}{%
	\ifcase\cotDS
		\ifnum\cotSTT=1\relax
			\def\rowHeader{\hline\rowcolor{gold1}\EZo@TF{TT}&\EZo@TF{\phatbieu}\tabularnewline}%
		\else
			\def\rowHeader{\hline\rowcolor{gold1}\EZo@TF{\phatbieu}\tabularnewline}%
		\fi
	\or
		\ifnum\cotSTT=1\relax
			\def\rowHeader{\hline\rowcolor{gold1}\EZo@TF{TT}&\EZo@TF{\phatbieu}&\EZo@TF{\Dung}\tabularnewline}%
		\else
			\def\rowHeader{\hline\rowcolor{gold1}\EZo@TF{\phatbieu}&\EZo@TF{\Dung}\tabularnewline}%
		\fi
	\or
		\ifnum\cotSTT=1\relax
			\def\rowHeader{\hline\rowcolor{gold1}\EZo@TF{TT}&\EZo@TF{\phatbieu}&\EZo@TF{\Dung}&\EZo@TF{\Sai}\tabularnewline}%
		\else
			\def\rowHeader{\hline\rowcolor{gold1}\EZo@TF{\phatbieu}&\EZo@TF{\Dung}&\EZo@TF{\Sai}\tabularnewline}%
		\fi
	\fi}
%--- tô xen kẽ từng hàng (gọi ngay sau \hline đầu mỗi hàng)
\newcommand{\EZmauhangTF}{%
	\noalign{\global\advance\EZhangTF\@ne}%
	\ifodd\EZhangTF\rowcolor{ezDSle}\else\rowcolor{ezDSchan}\fi}
%--- gắn vào \choiceTFt / \writeChoiceTFt của ex_test
\xpretocmd{\choiceTFt}{\EZbatdauTF}{}{\typeout{[preamble-eduzy-v2] Không gắn được bảng Đúng/Sai (đầu bảng)}}
\xpatchcmd{\choiceTFt}{\ifnum\dongPB=0\def\rowHeader{}\fi}{\EZtieudeTF\ifnum\dongPB=0\def\rowHeader{}\fi}
	{}{\typeout{[preamble-eduzy-v2] Không gắn được bảng Đúng/Sai (hàng tiêu đề)}}
\xpatchcmd{\writeChoiceTFt}{\hline\xdef\dotEX{\sdotEX}}{\hline\EZmauhangTF\xdef\dotEX{\sdotEX}}
	{}{\typeout{[preamble-eduzy-v2] Không gắn được bảng Đúng/Sai (tô hàng)}}
\xpatchcmd{\writeChoiceTFt}{{\end{longtable}}}{{\end{longtable}\EZketthucTF}}
	{}{\typeout{[preamble-eduzy-v2] Không gắn được bảng Đúng/Sai (cuối longtable)}}
\xpatchcmd{\writeChoiceTFt}{{\end{xtabular}\end{center}\medskip\noindent}}{{\end{xtabular}\end{center}\EZketthucTF\medskip\noindent}}
	{}{\typeout{[preamble-eduzy-v2] Không gắn được bảng Đúng/Sai (cuối xtabular)}}
\fi
%=====================================================================
\renewcommand{\parallel}{\mathrel{/\!\!/}}

%  trong file đề (tools/dan_trang/dan_trang.py tự ghi khi đặt \NeedspaceCau):
%    \def\TachCauDS{1}   câu Đúng/Sai được tách: câu dẫn (kèm hình) ở cuối trang, bảng Phát biểu sang trang sau
%                        (câu dẫn và hình không tách nhau, bảng không tách giữa hai trang)
%    \def\TachCauLG{1}   bản lời giải (\CheDo khác dethi): phần đề của câu cao hơn 1/3 trang được phép tách
%  Lệnh dùng trong file đề (giá trị do tools/dan_trang/dan_trang.py đo ở bản dethi):
%    \NeedspaceCau[<cao câu dẫn + hình>]{<cao cả khối câu>}   ngay trên \begin{ex} (câu đầu phần: trên tiêu đề phần)
%        công tắc tắt: đúng bằng \Needspace{<cao cả khối câu>} như cũ;
%        \TachCauDS=1 và có [..] (câu Đúng/Sai được chọn tách): chỉ giữ câu dẫn + hình;
%        \TachCauLG=1 ở bản lời giải: đòi tối đa 1/3 \textheight.
%    \NeedspaceBang{<cao bảng Phát biểu>}%   ngay trước \choiceTFt: \TachCauDS=1 thì giữ trọn bảng trên một trang;
%        công tắc tắt: không làm gì.
\providecommand{\TachCauDS}{0}
\providecommand{\TachCauLG}{0}
\newlength{\EZnsCau}
\newcommand{\NeedspaceCau}[2][]{%
	\setlength{\EZnsCau}{#2}%
	\ifnum\TachCauDS=1 \if\relax\detokenize{#1}\relax\else\setlength{\EZnsCau}{#1}\fi\fi
	\ifnum\TachCauLG=1 \ifdefstring{\kindOPTN}{dethi}{}{\ifdim\EZnsCau>0.3333\textheight\setlength{\EZnsCau}{0.3333\textheight}\fi}\fi
	\Needspace{\EZnsCau}}
\newcommand{\NeedspaceBang}[1]{\ifnum\TachCauDS=1 \Needspace{#1}\fi}

%--- Thang màu: đậm -> nhạt
\definecolor{gold1}{RGB}{126,96,0}    % nâu đồng đậm  #7E6000
\definecolor{gold2}{RGB}{192,144,3}   % vàng kim      #C09003
\definecolor{gold3}{RGB}{255,215,99}  % vàng nhạt     #FFD763
\definecolor{gold4}{RGB}{255,231,153} % vàng rất nhạt #FFE799
\definecolor{gold5}{RGB}{255,243,205} % vàng kem      #FFF3CD
%--- Ghi đè tông xanh/đỏ cũ
\definecolor{myblue}{RGB}{126,96,0}
\colorlet{blue}{gold1}
\colorlet{cyan}{gold2}
\colorlet{teal}{gold2}

%--- Tên chương mặc định (file đề nên tự khai báo lại)
\providecommand{\tenchuong}{}
%--- Chân trang
\providecommand{\thuonghieu}{TOÁN 12 -- Eduzy.com.vn}
\providecommand{\tacgia}{Tô Gia Bảo -- 0353703783}

%--- Chữ cạnh logo ở đầu trang
\providecommand{\khauhieu}{Eduzy.com.vn}

%====================== THÔNG SỐ ĐẦU TRANG, CHÂN TRANG (v2) ======================
%  Dải nền đậm tràn ngang trang; ô chữ vàng kem có cạnh xiên, kèm 2 vạch xiên chuyển màu.
%  Đầu trang lẻ : logo + \khauhieu (trái)   | ô chữ \tenchuong (phải)
%  Chân trang lẻ: ô chữ \thuonghieu (trái)  | Trang (số) (phải)
%  Trang chẵn soi gương trang lẻ; chân trang chẵn ghi \tacgia.
%--- Màu: chỉ dùng thang vàng kim Eduzy
\colorlet{ezNen}{gold1}% dải nền đậm
\colorlet{ezVach1}{gold3}% vạch xiên sát dải nền (đậm hơn)
\colorlet{ezVach2}{gold4}% vạch xiên sát ô chữ (nhạt hơn)
\colorlet{ezO}{gold5}% ô chữ vàng kem
%--- Kích thước
\newlength{\EZdai}\setlength{\EZdai}{1.1cm}% chiều cao dải nền
\newlength{\EZo}\setlength{\EZo}{0.68cm}% chiều cao ô chữ
\newlength{\EZxien}\setlength{\EZxien}{0.55cm}% độ lệch ngang của cạnh xiên
\newlength{\EZvach}\setlength{\EZvach}{0.75cm}% bề rộng mỗi vạch xiên
\newlength{\EZkhe}\setlength{\EZkhe}{0.12cm}% khe trắng giữa vạch và ô chữ
\newlength{\EZdem}\setlength{\EZdem}{0.4cm}% khoảng đệm từ chữ tới cạnh xiên
\newlength{\EZle}\setlength{\EZle}{1.8cm}% lề trái/phải của chữ (bằng lề thân bài)
\newlength{\EZrongmin}\setlength{\EZrongmin}{4.5cm}% ô chữ rộng tối thiểu
\newlength{\EZrong}% bề rộng chữ trong ô (tính tự động)
%--- Font
\newcommand{\EZfontO}{\fontsize{10.5pt}{12pt}\selectfont\bfseries}% chữ trong ô
\newcommand{\EZfontKH}{\fontfamily{qag}\fontsize{11pt}{12pt}\selectfont\bfseries}% \khauhieu cạnh logo
\newcommand{\EZfontTrang}{\fontsize{11.5pt}{12pt}\selectfont\bfseries}% chữ "Trang" và số trang
%--- Đo bề rộng chữ trong ô
\newcommand{\EZdorong}[1]{%
	\settowidth{\EZrong}{\EZfontO #1}%
	\ifdim\EZrong<\EZrongmin\setlength{\EZrong}{\EZrongmin}\fi}
%--- Logo + chữ trắng \khauhieu (dựng sẵn một lần khi vào \begin{document})
\newsavebox{\EZlogobox}
\AtBeginDocument{\sbox{\EZlogobox}{%
		\raisebox{-0.5\height}{\eduzylogo[1cm]}\hspace{0.15cm}%
		\raisebox{-0.5\height}{\color{white}\EZfontKH\khauhieu}}}
%--- Hai vạch xiên + khe trắng. #1: điểm gốc ở mép dưới, #2: hướng (1 sang phải, -1 sang trái),
%    #3: hướng nghiêng (1 kiểu "/", -1 kiểu "\")
\newcommand{\EZvachxien}[3]{%
	\fill[white] (#1) --++({(#3)*\EZxien},\EZo) --++({(#2)*\EZkhe},0) --++({-(#3)*\EZxien},-\EZo) --cycle;
	\fill[ezVach2] ([xshift=(#2)*\EZkhe]#1) --++({(#3)*\EZxien},\EZo) --++({(#2)*\EZvach},0) --++({-(#3)*\EZxien},-\EZo) --cycle;
	\fill[ezVach1] ([xshift=(#2)*(\EZkhe+\EZvach)]#1) --++({(#3)*\EZxien},\EZo) --++({(#2)*\EZvach},0) --++({-(#3)*\EZxien},-\EZo) --cycle;
	\draw[white,line width=0.6pt] ([xshift=(#2)*(\EZkhe+\EZvach)]#1) --++({(#3)*\EZxien},\EZo);
}
%--- Chữ trong ô (căn giữa theo chiều cao ô, không lệch vì dấu tiếng Việt)
\tikzset{ezchu/.style={inner sep=0pt,text=ezNen,text height=0.72em,text depth=0.22em}}

%====================== LOGO CHÌM GIỮA TRANG ======================
%  Vẽ cùng đầu trang nên nằm dưới chữ của thân bài, có trên mọi trang.
%  Tắt: \logochimtat (từ trang đó trở đi), bật lại: \logochimbat.
%  Đổi cỡ, độ đậm: khai báo trước \input, ví dụ \def\EZchimcao{9cm} \def\EZchimdo{0.1}
\providecommand{\EZchimcao}{10cm}% chiều cao logo chìm
\providecommand{\EZchimdo}{0.12}% độ đậm (0 = trong suốt, 1 = rõ hẳn)
\newif\ifEZchim \EZchimtrue
\newcommand{\logochimtat}{\EZchimfalse}
\newcommand{\logochimbat}{\EZchimtrue}
%  Dựng sẵn logo mờ một lần khi vào \begin{document}; transparency group làm mờ cả logo
%  một lượt nên không lộ vệt ở chỗ các dải màu chồng nhau.
\newsavebox{\EZchimtho}
\newsavebox{\EZchimbox}
\AtBeginDocument{%
	\sbox{\EZchimtho}{\eduzylogo[\EZchimcao]}%
	\sbox{\EZchimbox}{\begin{tikzpicture}
			\begin{scope}[transparency group,opacity=\EZchimdo]
				\node[inner sep=0pt] {\usebox{\EZchimtho}};
			\end{scope}
	\end{tikzpicture}}}
\newcommand{\EZvechim}{%
	\ifEZchim
		\node[inner sep=0pt] at (current page.center) {\usebox{\EZchimbox}};
	\fi}

%====================== ĐẦU TRANG - TRANG LẺ ======================
%  (logo bên trái, ô chữ \tenchuong bên phải)
\newcommand{\myfancyheadOdd}{%
	\EZdorong{\tenchuong}%
	\begin{tikzpicture}[remember picture,overlay]
		\coordinate (T) at (current page.north west);
		\coordinate (P) at (current page.north east);
		%--- logo chìm giữa trang
		\EZvechim
		%--- dải nền
		\fill[ezNen] (T) rectangle ([yshift=-\EZdai]P);
		%--- ô chữ: X là góc dưới bên trái, cạnh trái xiên kiểu "/"
		\coordinate (X) at ([xshift=-\EZle-\EZrong-\EZdem-\EZxien,yshift=-\EZdai]P);
		\fill[ezO] (X) --++(\EZxien,\EZo) -- ([yshift=-\EZdai+\EZo]P) -- ([yshift=-\EZdai]P) --cycle;
		\EZvachxien{X}{-1}{1}
		\node[ezchu,anchor=east] at ([xshift=-\EZle,yshift=-\EZdai+0.5*\EZo]P) {\EZfontO\tenchuong};
		%--- logo + tên
		\node[inner sep=0pt,anchor=west] at ([xshift=\EZle,yshift=-0.5*\EZdai]T) {\usebox{\EZlogobox}};
	\end{tikzpicture}%
}

%====================== ĐẦU TRANG - TRANG CHẴN =====================
%  (soi gương: ô chữ \tenchuong bên trái, logo bên phải)
\newcommand{\myfancyheadEven}{%
	\EZdorong{\tenchuong}%
	\begin{tikzpicture}[remember picture,overlay]
		\coordinate (T) at (current page.north west);
		\coordinate (P) at (current page.north east);
		%--- logo chìm giữa trang
		\EZvechim
		%--- dải nền
		\fill[ezNen] (T) rectangle ([yshift=-\EZdai]P);
		%--- ô chữ: X là góc dưới bên phải, cạnh phải xiên kiểu "\"
		\coordinate (X) at ([xshift=\EZle+\EZrong+\EZdem+\EZxien,yshift=-\EZdai]T);
		\fill[ezO] (X) --++(-\EZxien,\EZo) -- ([yshift=-\EZdai+\EZo]T) -- ([yshift=-\EZdai]T) --cycle;
		\EZvachxien{X}{1}{-1}
		\node[ezchu,anchor=west] at ([xshift=\EZle,yshift=-\EZdai+0.5*\EZo]T) {\EZfontO\tenchuong};
		%--- logo + tên
		\node[inner sep=0pt,anchor=east] at ([xshift=-\EZle,yshift=-0.5*\EZdai]P) {\usebox{\EZlogobox}};
	\end{tikzpicture}%
}

%====================== CHÂN TRANG - TRANG LẺ ======================
%  (ô chữ \thuonghieu bên trái, "Trang (số)" bên phải)
\newcommand{\myfancyfootOdd}{%
	\EZdorong{\thuonghieu}%
	\begin{tikzpicture}[remember picture,overlay]
		\coordinate (T) at (current page.south west);
		\coordinate (P) at (current page.south east);
		%--- dải nền
		\fill[ezNen] (T) rectangle ([yshift=\EZdai]P);
		%--- ô chữ sát mép trên dải: X là góc dưới bên phải, cạnh phải xiên kiểu "/"
		\coordinate (X) at ([xshift=\EZle+\EZrong+\EZdem,yshift=\EZdai-\EZo]T);
		\fill[ezO] ([yshift=\EZdai-\EZo]T) -- (X) --++(\EZxien,\EZo) -- ([yshift=\EZdai]T) --cycle;
		\EZvachxien{X}{1}{1}
		\node[ezchu,anchor=west] at ([xshift=\EZle,yshift=\EZdai-0.5*\EZo]T) {\EZfontO\thuonghieu};
		%--- số trang
		\node[circle,draw=white,line width=0.9pt,text=white,inner sep=0pt,minimum size=0.62cm,anchor=east]
		(so) at ([xshift=-\EZle,yshift=0.5*\EZdai]P) {\EZfontTrang\thepage};
		\node[inner sep=0pt,text=white,anchor=base east] (tr) at ([xshift=-0.15cm]so.base west) {\EZfontTrang Trang};
	\end{tikzpicture}%
}

%====================== CHÂN TRANG - TRANG CHẴN ====================
%  (soi gương: "Trang (số)" bên trái, ô chữ \tacgia bên phải)
\newcommand{\myfancyfootEven}{%
	\EZdorong{\tacgia}%
	\begin{tikzpicture}[remember picture,overlay]
		\coordinate (T) at (current page.south west);
		\coordinate (P) at (current page.south east);
		%--- dải nền
		\fill[ezNen] (T) rectangle ([yshift=\EZdai]P);
		%--- ô chữ sát mép trên dải: X là góc dưới bên trái, cạnh trái xiên kiểu "\"
		\coordinate (X) at ([xshift=-\EZle-\EZrong-\EZdem,yshift=\EZdai-\EZo]P);
		\fill[ezO] ([yshift=\EZdai-\EZo]P) -- (X) --++(-\EZxien,\EZo) -- ([yshift=\EZdai]P) --cycle;
		\EZvachxien{X}{-1}{-1}
		\node[ezchu,anchor=east] at ([xshift=-\EZle,yshift=\EZdai-0.5*\EZo]P) {\EZfontO\tacgia};
		%--- số trang
		%    (vòng tròn ẩn so0 nằm giữa dải để lấy đường chân chữ, giống trang lẻ)
		\node[circle,inner sep=0pt,minimum size=0.62cm] (so0) at ([yshift=0.5*\EZdai]T) {\EZfontTrang\phantom{\thepage}};
		\node[inner sep=0pt,text=white,anchor=base west] (tr) at ($(so0.base -| T)+(\EZle,0)$) {\EZfontTrang Trang};
		\node[circle,draw=white,line width=0.9pt,text=white,inner sep=0pt,minimum size=0.62cm,anchor=base west]
		(so) at ([xshift=0.15cm]tr.base east) {\EZfontTrang\thepage};
	\end{tikzpicture}%
}

%--- Giữ tương thích với các file cũ đang gọi \myfancyhead, \myfancyfoot
\newcommand{\myfancyhead}{\myfancyheadOdd}
\newcommand{\myfancyfoot}{\myfancyfootOdd}

%======================= GẮN VÀO TRANG =============================
\pagestyle{fancy}
\fancyhf{}
\fancyhead[LO]{\myfancyheadOdd}
\fancyhead[RE]{\myfancyheadEven}
\fancyfoot[LO]{\myfancyfootOdd}
\fancyfoot[RE]{\myfancyfootEven}
\renewcommand{\headrulewidth}{0pt}
\renewcommand{\footrulewidth}{0pt}
\fancypagestyle{plain}{%
	\fancyhf{}
	\fancyhead[LO]{\myfancyheadOdd}
	\fancyhead[RE]{\myfancyheadEven}
	\fancyfoot[LO]{\myfancyfootOdd}
	\fancyfoot[RE]{\myfancyfootEven}
	\renewcommand{\headrulewidth}{0pt}
	\renewcommand{\footrulewidth}{0pt}
}
%=================== HẾT PHẦN GỘP TỪ eduzy-de.tex ===================
\newcommand{\indam}[1]{
	{\fontsize{12pt}{9pt}\sffamily \selectfont\bfseries\color{gold1} #1}
}
\newcommand{\ind}[1]{
	{\fontsize{14pt}{9pt}\fontfamily{qag} \selectfont\bfseries\color{gold1} #1}
}
\renewcommand{\FalseEX}{\stepcounter{dapan}\textbf{\selectfont\color{gold1}\Alph{dapan}.}\ignorespaces}
\renewtheorem{ex}{\color{gold1}Câu}
\renewcommand{\baselinestretch}{1.1}
\newcommand{\hoac}[1]{
	\left[\begin{aligned}#1\end{aligned}\right.}
\newcommand{\heva}[1]{
	\left\{\begin{aligned}#1\end{aligned}\right.}
\usepackage{esvect}
\def\vec{\vv}
\def\vt{\vv}
\def\overrightarrow{\vv}
%====================== LOGO EDUZY VẼ BẰNG TIKZ ======================
%  Đường viền dò từ logo gốc; dải màu vàng kim đo trực tiếp trên ảnh gốc
%  (nghiêng 21 độ, tô bằng 80 dải hẹp nên hiển thị giống nhau trên mọi trình xem PDF).
%  Dùng: \eduzylogo          -> cao 2.4cm (mặc định)
%        \eduzylogo[1.5cm]   -> cao 1.5cm
\definecolor{ezText}{RGB}{245,200,46}
\newcommand{\eduzylogo}[1][2.4cm]{%
	\resizebox{!}{#1}{%
		\begin{tikzpicture}[x=0.01cm,y=0.01cm]
			\path[use as bounding box] (0,0) rectangle (224,-224);
			% mũ tốt nghiệp + vương miện
			\begin{scope}
				\clip
				(119.4,-176.1)..controls(118.9,-176.1)and(116.8,-176)..(114.6,-176)
				..controls(101.4,-176)and(85.8,-175.9)..(83.7,-175.7)
				..controls(81.8,-175.6)and(81,-175.6)..(80.2,-175.7)
				..controls(78.8,-175.9)and(77.7,-175.9)..(76.2,-175.7)
				..controls(74.9,-175.5)and(70.8,-175.5)..(64.3,-175.7)
				..controls(59.1,-175.9)and(58.3,-175.7)..(58.6,-173.7)
				..controls(58.6,-172.8)and(59.6,-169.6)..(60.3,-168)
				..controls(60.8,-166.9)and(61.4,-165.8)..(62.3,-165)
				..controls(63.1,-164.2)and(64.3,-163.4)..(65,-163.2)
				..controls(65.5,-163.1)and(72.3,-163)..(76.6,-163.1)
				..controls(77.9,-163.1)and(80.1,-163.1)..(81.4,-163)
				..controls(82.8,-162.9)and(86.3,-162.9)..(89.2,-162.9)
				..controls(93,-162.9)and(94.9,-162.8)..(96,-162.7)
				..controls(97.2,-162.5)and(101,-162.5)..(114,-162.4)
				..controls(123,-162.3)and(134.9,-162.2)..(140.2,-162.1)
				..controls(150.1,-162)and(152.2,-162.1)..(153.6,-162.4)
				..controls(154.4,-162.6)and(157.1,-163.8)..(158.1,-164.4)
				..controls(158.8,-164.8)and(159.7,-165.7)..(159.9,-166.3)
				..controls(160,-166.5)and(160.2,-167.1)..(160.2,-167.6)
				..controls(160.3,-168)and(160.6,-169.4)..(160.9,-170.6)
				..controls(161.2,-171.7)and(161.5,-173.1)..(161.6,-173.6)
				..controls(161.8,-174.4)and(161.8,-174.5)..(161.5,-174.9)
				..controls(161.2,-175.6)and(160.8,-175.9)..(159.9,-176)
				..controls(159.2,-176.1)and(154.1,-176.2)..(151.5,-176.2)
				..controls(148.1,-176.1)and(122.6,-176.1)..(121.6,-176.2)
				..controls(120.8,-176.2)and(119.8,-176.2)..(119.4,-176.1)--cycle
				(130.4,-160)..controls(129,-160)and(119.8,-160)..(110.1,-159.9)
				..controls(100.3,-159.8)and(88.5,-159.7)..(83.9,-159.6)
				..controls(77.8,-159.5)and(75.3,-159.5)..(74.5,-159.6)
				..controls(73.9,-159.7)and(72.5,-159.8)..(71.3,-159.8)
				..controls(69.6,-159.8)and(69,-159.8)..(68.5,-159.6)
				..controls(67.6,-159.3)and(66.8,-158.6)..(66.2,-157.6)
				..controls(66,-157.3)and(65.6,-156.5)..(65.3,-156.1)
				..controls(64.6,-155.1)and(64.5,-154.9)..(63.7,-152.9)
				..controls(63.2,-151.6)and(62.9,-150.9)..(61.8,-149.6)
				..controls(61.5,-149.2)and(61.1,-148.5)..(60.9,-147.9)
				..controls(60.7,-147.4)and(60.3,-146.6)..(60.1,-146.1)
				..controls(59.9,-145.5)and(59.5,-144.9)..(59,-144.3)
				..controls(58.5,-143.7)and(58.2,-143.2)..(57.9,-142.5)
				..controls(57.7,-141.9)and(57.3,-141.1)..(57,-140.7)
				..controls(56.6,-140.3)and(56.1,-139.3)..(55.8,-138.5)
				..controls(55.5,-137.8)and(55.1,-136.8)..(54.9,-136.5)
				..controls(54.7,-136.1)and(54.3,-135.2)..(53.9,-134.5)
				..controls(53.6,-133.8)and(53,-132.7)..(52.6,-132)
				..controls(51.6,-130.5)and(51.4,-130.2)..(51.1,-129)
				..controls(50.6,-127.5)and(50.6,-127.4)..(49.7,-126.3)
				..controls(49.2,-125.7)and(48.6,-124.8)..(48.3,-124.2)
				..controls(48,-123.6)and(47.5,-122.7)..(47.3,-122.2)
				..controls(47,-121.8)and(46.5,-120.7)..(46.1,-119.8)
				..controls(45.8,-118.9)and(45.2,-117.8)..(44.9,-117.2)
				..controls(44.2,-116.1)and(43.3,-114.4)..(42.3,-112.2)
				..controls(41.8,-111.3)and(41.3,-110.4)..(41.1,-110.1)
				..controls(40.5,-109.4)and(40.3,-109)..(40,-107.8)
				..controls(39.8,-107)and(39.6,-106.5)..(39.2,-106)
				..controls(38.9,-105.6)and(38.4,-104.7)..(38,-104)
				..controls(37.7,-103.4)and(37.2,-102.5)..(37,-102.1)
				..controls(36.7,-101.7)and(36.3,-100.8)..(36,-100.1)
				..controls(35.6,-98.9)and(35.3,-98.4)..(34.6,-97.6)
				..controls(34.2,-97.1)and(34.2,-96.3)..(34.6,-96)
				..controls(35,-95.7)and(35.5,-95.9)..(36.2,-96.7)
				..controls(36.5,-97)and(37.2,-97.5)..(37.6,-97.7)..controls(38,-98)and(38.6,-98.4)..(39.1,-98.9)
				..controls(39.6,-99.4)and(40.3,-100)..(40.6,-100.3)
				..controls(41.3,-100.8)and(42.4,-101.8)..(43.7,-102.9)
				..controls(44,-103.2)and(44.6,-103.7)..(45,-103.9)
				..controls(45.4,-104.2)and(45.9,-104.6)..(46.1,-104.8)
				..controls(47.1,-105.9)and(48.8,-107.1)..(50.5,-107.9)
				..controls(51.8,-108.6)and(52.3,-109)..(53.4,-110)
				..controls(54.8,-111.2)and(55.6,-111.6)..(57,-112.1)
				..controls(57.5,-112.3)and(58.5,-112.8)..(59.4,-113.3)
				..controls(60.2,-113.7)and(61.3,-114.3)..(61.9,-114.5)
				..controls(62.4,-114.7)and(63.3,-115.1)..(63.8,-115.3)
				..controls(64.4,-115.6)and(65.4,-115.9)..(66.2,-116.1)
				..controls(67,-116.3)and(68.4,-116.6)..(69.2,-116.9)
				..controls(72.3,-117.9)and(72.8,-118)..(78,-118.1)
				..controls(84,-118.3)and(86.1,-118.1)..(88.1,-117.1)
				..controls(89.4,-116.5)and(93,-115)..(93.7,-114.9)
				..controls(94.8,-114.6)and(95.4,-114.2)..(95.9,-113.4)
				..controls(96.2,-113)and(96.7,-112.5)..(97.1,-112.2)
				..controls(98,-111.6)and(101.3,-108.4)..(102.1,-107.2)
				..controls(102.4,-106.8)and(102.9,-106.2)..(103.2,-105.8)
				..controls(103.5,-105.5)and(104,-104.8)..(104.2,-104.2)
				..controls(104.5,-103.6)and(105,-102.8)..(105.3,-102.4)
				..controls(105.6,-102)and(106.1,-101.2)..(106.5,-100.6)
				..controls(107.1,-99.5)and(108.6,-97.5)..(109.5,-96.5)
				..controls(109.9,-96.1)and(110.1,-96)..(110.4,-96)
				..controls(111.1,-96)and(111.6,-96.5)..(111.8,-97.6)
				..controls(112,-98.6)and(112.3,-99.3)..(112.8,-99.9)
				..controls(113.7,-100.9)and(115.2,-103.1)..(115.7,-104.2)
				..controls(116.1,-105)and(116.5,-105.5)..(116.9,-106)
				..controls(117.3,-106.4)and(117.8,-106.9)..(118,-107.3)
				..controls(118.2,-107.6)and(118.6,-108)..(119,-108.3)
				..controls(119.6,-108.8)and(121.1,-110.3)..(121.8,-111.2)
				..controls(122.1,-111.7)and(122.5,-112)..(123.4,-112.4)
				..controls(124.8,-113.2)and(125.8,-113.8)..(127,-114.6)
				..controls(127.5,-114.9)and(128.1,-115.3)..(128.2,-115.4)
				..controls(128.4,-115.5)and(129.3,-115.8)..(130.2,-116)
				..controls(131.1,-116.2)and(132.5,-116.6)..(133.3,-116.9)
				..controls(135.7,-117.9)and(136.1,-117.9)..(142.5,-117.8)
				..controls(145.8,-117.8)and(148.4,-117.7)..(148.8,-117.6)
				..controls(149.2,-117.6)and(150.2,-117.2)..(151,-116.8)
				..controls(152,-116.4)and(152.8,-116)..(153.6,-115.9)
				..controls(154.5,-115.7)and(155.2,-115.4)..(156.3,-114.9)
				..controls(157.1,-114.5)and(158.2,-114)..(158.7,-113.8)
				..controls(159.1,-113.7)and(159.8,-113.3)..(160.2,-113)
				..controls(160.6,-112.7)and(161.3,-112.4)..(161.7,-112.2)
				..controls(162.7,-111.9)and(163.2,-111.7)..(164.1,-111.1)
				..controls(164.5,-110.8)and(165.3,-110.4)..(165.9,-110)
				..controls(166.5,-109.7)and(167.3,-109.2)..(167.6,-108.9)
				..controls(167.9,-108.6)and(168.6,-108.2)..(169,-108)
				..controls(169.4,-107.8)and(170.2,-107.3)..(170.7,-106.9)
				..controls(171.2,-106.5)and(172.1,-105.8)..(172.7,-105.4)
				..controls(173.2,-105)and(174.1,-104.4)..(174.5,-104)
				..controls(174.9,-103.7)and(175.7,-103.1)..(176.2,-102.6)
				..controls(176.7,-102.2)and(177.4,-101.5)..(177.9,-101.2)
				..controls(178.3,-100.8)and(178.8,-100.4)..(179,-100.2)
				..controls(179.2,-100)and(179.6,-99.6)..(180,-99.3)
				..controls(180.4,-99.1)and(180.9,-98.6)..(181.2,-98.4)
				..controls(181.5,-98.1)and(182.1,-97.7)..(182.5,-97.5)
				..controls(183.2,-97.2)and(183.5,-96.9)..(184.3,-95.9)
				..controls(184.6,-95.6)and(184.6,-95.6)..(184.8,-95.7)
				..controls(185,-95.8)and(185.1,-96)..(185.1,-96.2)
				..controls(185.1,-96.5)and(185,-96.7)..(184.6,-97.2)
				..controls(184.1,-97.8)and(184,-97.9)..(183.9,-98.8)
				..controls(183.7,-100.3)and(183.3,-101.3)..(182.9,-101.8)
				..controls(182.6,-102.1)and(182.3,-102.6)..(182.2,-102.9)
				..controls(182,-103.3)and(181.6,-104)..(181.2,-104.6)
				..controls(180.9,-105.2)and(180.5,-106.1)..(180.4,-106.6)
				..controls(179.9,-108)and(179.7,-108.5)..(179,-109.4)
				..controls(178.6,-109.8)and(178.2,-110.5)..(177.9,-111)
				..controls(177.7,-111.5)and(177.3,-112.3)..(177,-112.7)
				..controls(176.6,-113.2)and(176.4,-113.7)..(176.3,-114.2)
				..controls(176.2,-114.7)and(175.9,-115.2)..(175.4,-116)
				..controls(175,-116.6)and(174.4,-117.7)..(174,-118.5)
				..controls(173.7,-119.2)and(173.3,-120.1)..(173,-120.5)
				..controls(172.8,-120.9)and(172.4,-121.8)..(172.1,-122.5)
				..controls(171.8,-123.2)and(171.3,-124.1)..(171,-124.6)
				..controls(170.8,-125.1)and(170.3,-126)..(169.9,-126.6)
				..controls(169.6,-127.3)and(169.2,-128.1)..(169,-128.3)
				..controls(168.8,-128.6)and(168.5,-129.3)..(168.3,-129.8)
				..controls(168.1,-130.3)and(167.8,-131.1)..(167.5,-131.6)
				..controls(166.9,-132.5)and(164.2,-137.9)..(164.2,-138.1)
				..controls(164.2,-138.1)and(163.9,-138.8)..(163.4,-139.7)
				..controls(162.9,-140.5)and(162.4,-141.5)..(162.3,-141.9)
				..controls(161.9,-143)and(161.6,-143.4)..(160.9,-144.1)
				..controls(160.4,-144.6)and(160.3,-144.8)..(160.1,-145.4)
				..controls(160,-145.8)and(159.6,-146.8)..(159.2,-147.8)
				..controls(158.7,-148.8)and(158.1,-150)..(157.9,-150.6)
				..controls(157.6,-151.1)and(157.2,-151.9)..(156.9,-152.4)
				..controls(156.5,-152.8)and(156.1,-153.7)..(155.8,-154.4)
				..controls(155.6,-155)and(155.2,-155.9)..(154.9,-156.3)
				..controls(154.7,-156.7)and(154.4,-157.4)..(154.2,-157.9)
				..controls(154.1,-158.4)and(153.8,-159)..(153.5,-159.3)
				..controls(152.8,-160.1)and(153,-160.1)..(142.4,-160.1)
				..controls(137.2,-160.1)and(131.8,-160.1)..(130.4,-160)--cycle
				(141.5,-103.7)..controls(140.8,-103)and(139.8,-102.1)..(139.2,-101.6)
				..controls(138.5,-101.2)and(137.9,-100.6)..(137.8,-100.4)
				..controls(137,-99.5)and(134,-96.5)..(133.2,-96)
				..controls(132.7,-95.7)and(131.8,-94.9)..(131.1,-94.3)
				..controls(129.4,-92.7)and(127.9,-91.6)..(126.6,-91.1)
				..controls(125.2,-90.5)and(124.1,-90)..(122.9,-89.1)
				..controls(122.1,-88.6)and(121.3,-88.2)..(119.9,-87.7)
				..controls(118.8,-87.3)and(117.5,-86.8)..(116.9,-86.5)
				..controls(114.3,-85.5)and(113.7,-85.4)..(109.7,-85.5)
				..controls(105.3,-85.7)and(103.9,-85.9)..(101.6,-87)
				..controls(100.7,-87.4)and(99.4,-87.8)..(97.9,-88)
				..controls(97,-88.2)and(96.6,-88.3)..(95.8,-88.9)
				..controls(95.4,-89.2)and(94.7,-89.6)..(94.2,-89.8)
				..controls(92.9,-90.3)and(92.3,-90.7)..(90.9,-92.1)
				..controls(90.2,-92.7)and(89.4,-93.3)..(89,-93.5)..controls(87,-94.5)and(85.2,-96)..(83.6,-97.9)
				..controls(83.3,-98.2)and(82.5,-98.9)..(81.8,-99.4)
				..controls(80.7,-100.1)and(80.5,-100.3)..(80,-101.1)
				..controls(78.7,-103.1)and(77.2,-103.9)..(76.4,-103)
				..controls(75.7,-102.2)and(75.7,-103.2)..(75.7,-83.6)
				..controls(75.6,-73.6)and(75.6,-64.8)..(75.5,-64)
				..controls(75.5,-63.2)and(75.5,-61.4)..(75.6,-59.8)--(75.7,-57.1)--(76,-56.8)
				..controls(76.8,-56.2)and(78,-56.3)..(79.8,-57.2)
				..controls(80.4,-57.5)and(81.3,-57.8)..(81.8,-57.9)
				..controls(82.3,-58.1)and(83,-58.4)..(83.6,-58.8)..controls(84.3,-59.2)and(85,-59.5)..(86,-59.8)
				..controls(86.8,-60.1)and(88,-60.6)..(88.7,-60.9)
				..controls(89.3,-61.2)and(90.6,-61.7)..(91.4,-62)
				..controls(92.2,-62.3)and(93.7,-62.9)..(94.6,-63.3)
				..controls(95.6,-63.7)and(96.8,-64.1)..(97.4,-64.3)
				..controls(98,-64.5)and(99,-64.9)..(99.7,-65.2)
				..controls(100.4,-65.6)and(101.5,-66)..(102.3,-66.2)
				..controls(104.2,-66.8)and(105,-67.1)..(106.2,-67.8)
				..controls(107.9,-68.7)and(110.2,-69.4)..(111.3,-69.2)
				..controls(112.5,-69)and(115,-68.2)..(115.6,-67.8)
				..controls(117.3,-66.7)and(117.6,-66.6)..(119.5,-66.2)
				..controls(121.9,-65.8)and(122.1,-65.7)..(123,-65.1)
				..controls(124.1,-64.4)and(124.4,-64.3)..(126,-63.9)
				..controls(127,-63.7)and(127.7,-63.4)..(128.4,-63.1)
				..controls(129,-62.8)and(129.7,-62.5)..(130.6,-62.4)
				..controls(132.1,-62.1)and(134,-61.5)..(135.2,-61)
				..controls(135.7,-60.8)and(136.7,-60.5)..(137.3,-60.3)
				..controls(138.4,-60)and(139.3,-59.7)..(141.1,-59)
				..controls(141.6,-58.7)and(142,-58.6)..(142.5,-58.6)
				..controls(143.2,-58.6)and(143.3,-58.7)..(143.6,-59)
				..controls(144,-59.5)and(144.1,-59.9)..(144,-60.8)
				..controls(144,-61.2)and(144,-71.1)..(144,-82.8)--(143.9,-104.1)--(143.7,-104.5)
				..controls(143.5,-104.7)and(143.3,-104.8)..(143,-104.9)
				..controls(142.6,-104.9)and(142.5,-104.8)..(141.5,-103.7)--cycle
				(27.8,-91.4)..controls(27.1,-90.9)and(26.3,-90)..(26.1,-89.5)
				..controls(26,-89.2)and(25.8,-88.6)..(25.7,-87.9)
				..controls(25.4,-86.2)and(25.6,-84.8)..(26.5,-83.5)
				..controls(27.4,-82.3)and(28.3,-82)..(29.4,-82.3)
				..controls(31.1,-82.8)and(32.1,-84.9)..(32,-87.4)
				..controls(31.9,-89.3)and(31.2,-90.6)..(29.9,-91.4)
				..controls(29.1,-91.9)and(28.6,-91.9)..(27.8,-91.4)--cycle
				(190.1,-91.3)..controls(189.7,-91.2)and(189.3,-90.9)..(189.1,-90.7)
				..controls(188.3,-89.8)and(187.6,-87.7)..(187.6,-86.2)
				..controls(187.6,-85.5)and(187.7,-85.2)..(188.2,-84.4)
				..controls(189.1,-82.6)and(190.1,-81.8)..(191.3,-82.1)
				..controls(192,-82.3)and(193.5,-83.7)..(193.8,-84.4)
				..controls(194.1,-85.1)and(194.2,-87.1)..(194,-88.3)
				..controls(193.8,-89.8)and(193.3,-90.5)..(192.3,-91.3)
				..controls(191.7,-91.7)and(190.9,-91.7)..(190.1,-91.3)--cycle
				(173.7,-83.4)..controls(172.2,-82.5)and(172,-82.5)..(170.6,-82.4)
				..controls(169.7,-82.4)and(168.8,-82.4)..(168.3,-82.5)
				..controls(167.2,-82.7)and(166.9,-82.7)..(166.6,-82.2)
				..controls(166.3,-81.8)and(166.3,-81)..(166.8,-80)
				..controls(167.2,-78.9)and(167.4,-77.4)..(167.7,-74.1)
				..controls(167.8,-72.8)and(167.9,-71.8)..(167.9,-71.7)
				..controls(168,-71.7)and(168,-70.2)..(168,-68.4)..controls(168,-65.4)and(168,-64.9)..(167.7,-63)
				..controls(167.3,-60.5)and(167.1,-59.9)..(166.4,-58.4)
				..controls(166.1,-57.7)and(165.8,-56.9)..(165.7,-56.6)
				..controls(165.7,-56.2)and(165.4,-55.6)..(165.2,-55.2)
				..controls(164.6,-53.8)and(161.9,-50.7)..(161,-50.3)
				..controls(159.8,-49.7)and(157.9,-50)..(155.2,-51.2)
				..controls(154.7,-51.4)and(153.4,-51.8)..(152.2,-52.1)
				..controls(150.8,-52.4)and(149.4,-52.9)..(148.6,-53.2)
				..controls(147.8,-53.6)and(146.7,-54)..(146,-54.2)
				..controls(145.4,-54.4)and(144.5,-54.7)..(144,-54.9)
				..controls(143.6,-55.2)and(142.7,-55.5)..(141.8,-55.7)
				..controls(141,-55.9)and(139.6,-56.3)..(138.7,-56.6)
				..controls(137.8,-57)and(136.3,-57.4)..(135.3,-57.6)
				..controls(133.8,-57.9)and(133.2,-58.1)..(132.1,-58.7)
				..controls(131.1,-59.2)and(130.3,-59.4)..(129.4,-59.6)
				..controls(128.7,-59.8)and(127.8,-60)..(127.4,-60.1)
				..controls(127,-60.3)and(126,-60.5)..(125.1,-60.7)
				..controls(123.2,-61.1)and(121.7,-61.5)..(120.2,-62.1)
				..controls(119.7,-62.3)and(118.1,-62.8)..(116.7,-63.2)
				..controls(114.6,-63.9)and(114,-64.1)..(113,-64.7)
				..controls(111.9,-65.4)and(111.8,-65.4)..(111,-65.5)
				..controls(109.9,-65.5)and(109.3,-65.4)..(107.7,-64.7)
				..controls(107,-64.4)and(106,-64)..(105.3,-63.9)
				..controls(103.5,-63.5)and(102.2,-63)..(101.2,-62.5)
				..controls(100.1,-61.8)and(99.3,-61.5)..(97.8,-61.1)
				..controls(97.2,-61)and(96.2,-60.7)..(95.7,-60.4)
				..controls(95.1,-60.1)and(94.2,-59.8)..(93.5,-59.6)
				..controls(92.9,-59.4)and(91.7,-59)..(91,-58.7)..controls(90.2,-58.4)and(88.9,-58)..(88,-57.7)
				..controls(87.2,-57.5)and(86.1,-57.1)..(85.5,-56.8)
				..controls(85,-56.6)and(84.1,-56.2)..(83.5,-56)..controls(82.9,-55.7)and(81.8,-55.2)..(81,-54.9)
				..controls(80.1,-54.5)and(79.1,-54.2)..(78.7,-54.1)
				..controls(77.9,-53.8)and(77.3,-53.6)..(76.2,-52.9)
				..controls(75.7,-52.6)and(75.2,-52.4)..(74.1,-52.2)
				..controls(72.5,-51.9)and(72.4,-51.8)..(71.3,-51.1)
				..controls(69.9,-50.1)and(69.4,-49.9)..(67.9,-49.8)
				..controls(67,-49.7)and(66.6,-49.6)..(66.3,-49.4)
				..controls(66.1,-49.3)and(65.7,-49)..(65.4,-48.8)
				..controls(65.1,-48.7)and(64.2,-48.4)..(63.6,-48.2)
				..controls(62.9,-48)and(61.7,-47.6)..(60.9,-47.2)
				..controls(60.1,-46.8)and(58.9,-46.4)..(58.2,-46.2)
				..controls(57.6,-46)and(56.8,-45.7)..(56.5,-45.5)
				..controls(54.9,-44.6)and(54.6,-44.5)..(53.2,-44.2)
				..controls(51.7,-43.8)and(51,-43.6)..(49.9,-42.9)
				..controls(49.5,-42.6)and(48.9,-42.4)..(47.5,-42)
				..controls(45.7,-41.6)and(45.6,-41.5)..(44.3,-40.6)
				..controls(43.5,-40.1)and(42.8,-39.6)..(42.7,-39.4)
				..controls(42.4,-38.7)and(43,-38.1)..(44.5,-37.6)..controls(44.9,-37.5)and(46,-37)..(46.9,-36.6)
				..controls(47.9,-36.1)and(48.8,-35.8)..(49.6,-35.6)
				..controls(50.5,-35.4)and(51.3,-35.2)..(52.4,-34.6)
				..controls(53.3,-34.1)and(54.5,-33.7)..(54.9,-33.5)
				..controls(55.4,-33.4)and(56.4,-33)..(57.2,-32.6)
				..controls(58.6,-32)and(60.1,-31.5)..(61.6,-31.2)
				..controls(62.3,-31.1)and(62.7,-31)..(63.2,-30.7)
				..controls(64.1,-30)and(64.3,-29.9)..(65.3,-29.7)
				..controls(65.7,-29.6)and(66.5,-29.2)..(67.1,-28.9)
				..controls(67.6,-28.6)and(68.8,-28.2)..(69.6,-27.9)
				..controls(70.5,-27.7)and(71.6,-27.3)..(72.1,-27)
				..controls(72.6,-26.7)and(73.8,-26.2)..(74.8,-25.9)
				..controls(75.7,-25.5)and(76.7,-25.1)..(77.1,-25)
				..controls(77.4,-24.8)and(78.5,-24.4)..(79.6,-24.1)
				..controls(80.6,-23.8)and(82,-23.2)..(82.8,-22.9)
				..controls(83.5,-22.5)and(84.8,-22)..(85.8,-21.8)
				..controls(86.7,-21.5)and(87.8,-21.1)..(88.1,-20.9)
				..controls(88.4,-20.7)and(89.3,-20.3)..(90.1,-20.1)
				..controls(90.8,-19.8)and(91.9,-19.4)..(92.4,-19.1)
				..controls(93.5,-18.5)and(94.4,-18.2)..(95.4,-18.1)
				..controls(95.7,-18)and(96.4,-17.8)..(96.9,-17.5)
				..controls(97.5,-17.3)and(98.4,-16.9)..(99,-16.6)
				..controls(99.7,-16.3)and(100.9,-15.9)..(101.9,-15.6)
				..controls(102.9,-15.3)and(104.4,-14.7)..(105.3,-14.3)
				..controls(108.1,-13.2)and(108.5,-13.1)..(110.1,-13.1)
				..controls(111.4,-13.1)and(111.5,-13.1)..(112.4,-13.5)
				..controls(113.2,-14)and(115,-14.8)..(116.6,-15.4)
				..controls(116.9,-15.5)and(117.9,-15.8)..(118.8,-16.1)
				..controls(120.9,-16.6)and(121.5,-16.9)..(122.6,-17.9)
				..controls(123,-18.3)and(123.7,-18.8)..(123.9,-18.9)
				..controls(124.6,-19.4)and(126,-19.8)..(127.9,-20.2)
				..controls(129.6,-20.5)and(129.8,-20.6)..(131.1,-21.4)
				..controls(131.5,-21.7)and(132.4,-22.1)..(133.1,-22.3)
				..controls(136.1,-23.2)and(138.3,-24.1)..(140.2,-25)
				..controls(141,-25.5)and(141.5,-25.6)..(142.5,-25.8)
				..controls(145.1,-26.3)and(146,-26.6)..(149.4,-28.1)
				..controls(151.3,-28.9)and(153.5,-29.8)..(154.3,-30.1)
				..controls(155.1,-30.4)and(156.2,-30.8)..(156.6,-31)
				..controls(157.1,-31.2)and(158.2,-31.6)..(159.1,-31.8)
				..controls(160.6,-32.2)and(160.9,-32.3)..(161.7,-32.9)
				..controls(162.4,-33.3)and(162.9,-33.6)..(163.7,-33.8)
				..controls(164.3,-34)and(164.8,-34.1)..(164.9,-34.1)
				..controls(165,-34.1)and(165.6,-34.4)..(166.4,-34.8)
				..controls(167.3,-35.2)and(168.1,-35.4)..(169,-35.6)
				..controls(170,-35.8)and(170.5,-36)..(171.7,-36.6)
				..controls(172.4,-37)and(173.3,-37.4)..(173.6,-37.4)
				..controls(175.6,-37.7)and(176.2,-38)..(176.7,-39.1)
				..controls(176.9,-39.6)and(176.9,-39.7)..(176.7,-40.4)
				..controls(176.4,-41.6)and(175.8,-42.7)..(175.2,-43.4)
				..controls(174.5,-44.1)and(172.6,-45.5)..(171.7,-45.8)
				..controls(170.6,-46.2)and(169.1,-46)..(168.2,-45.4)
				..controls(167.9,-45.2)and(167.3,-44.8)..(166.9,-44.4)
				..controls(166.4,-44.1)and(165.8,-43.6)..(165.4,-43.2)
				..controls(165.1,-42.9)and(164.4,-42.3)..(163.9,-41.9)
				..controls(163.3,-41.5)and(162.7,-40.9)..(162.4,-40.5)
				..controls(161.4,-39.2)and(159.9,-38)..(158.8,-37.8)
				..controls(158.4,-37.8)and(158.3,-37.7)..(158.2,-37.2)
				..controls(158,-36.8)and(157.4,-36.1)..(157.1,-36.1)
				..controls(157,-36.1)and(156.7,-36.1)..(156.5,-36)
				..controls(156.1,-35.9)and(155.9,-35.7)..(155.7,-35.2)
				..controls(155.3,-34.6)and(155.1,-34.4)..(154,-34.1)
				..controls(153.6,-34)and(153.2,-33.7)..(152.9,-33.5)
				..controls(152.4,-32.9)and(151.6,-32.4)..(151.3,-32.4)
				..controls(150.9,-32.4)and(150.4,-32.6)..(150.4,-32.9)
				..controls(150.4,-33.2)and(150.8,-33.4)..(151.4,-33.6)
				..controls(152.1,-33.7)and(152.8,-34.1)..(153.5,-34.9)
				..controls(154.1,-35.5)and(154.3,-35.6)..(155,-35.8)
				..controls(155.7,-36)and(155.7,-36)..(156,-36.6)
				..controls(156.3,-37.3)and(156.7,-37.7)..(157.5,-37.8)
				..controls(158.1,-37.9)and(158.2,-38)..(158.4,-38.8)
				..controls(158.5,-39.4)and(158.6,-39.6)..(159.6,-40.4)
				..controls(160.2,-41)and(160.9,-41.6)..(161.3,-41.8)
				..controls(161.7,-42)and(161.9,-42.3)..(162.2,-42.8)
				..controls(162.4,-43.2)and(162.8,-43.8)..(163.1,-44.1)
				..controls(163.9,-44.8)and(165.1,-46.7)..(165.7,-48)
				..controls(166.2,-49.3)and(166.4,-49.7)..(167.2,-50.5)
				..controls(167.5,-50.8)and(167.9,-51.4)..(168.1,-51.8)
				..controls(168.3,-52.2)and(168.8,-53)..(169.1,-53.6)
				..controls(169.6,-54.3)and(169.8,-54.9)..(169.9,-55.6)
				..controls(170.4,-57.8)and(170.5,-58.1)..(171.2,-59.5)
				..controls(171.6,-60.3)and(172,-61.5)..(172.2,-62.1)
				..controls(172.3,-62.8)and(172.7,-63.7)..(172.9,-64.1)
				..controls(173.2,-64.6)and(173.4,-65.3)..(173.4,-65.6)
				..controls(173.8,-68.8)and(174.2,-70.6)..(174.9,-72.6)
				..controls(175.5,-74.5)and(175.9,-77.1)..(176.1,-80.2)
				..controls(176.3,-83.4)and(175.5,-84.4)..(173.7,-83.4)--cycle;
				\definecolor{ezs}{RGB}{247,179,26}\fill[ezs](-5.7,-156.6)--(-3.4,-157.5)--(61.9,12.7)--(59.7,13.6)--cycle;
				\definecolor{ezs}{RGB}{247,179,22}\fill[ezs](-3.8,-157.3)--(-1.5,-158.2)--(63.9,12)--(61.6,12.9)--cycle;
				\definecolor{ezs}{RGB}{248,179,19}\fill[ezs](-1.8,-158.1)--(0.4,-158.9)--(65.8,11.2)--(63.5,12.1)--cycle;
				\definecolor{ezs}{RGB}{248,178,16}\fill[ezs](0.1,-158.8)--(2.4,-159.7)--(67.7,10.5)--(65.4,11.4)--cycle;
				\definecolor{ezs}{RGB}{249,178,15}\fill[ezs](2,-159.5)--(4.3,-160.4)--(69.6,9.8)--(67.3,10.7)--cycle;
				\definecolor{ezs}{RGB}{251,177,17}\fill[ezs](3.9,-160.3)--(6.2,-161.1)--(71.5,9)--(69.2,9.9)--cycle;
				\definecolor{ezs}{RGB}{252,179,18}\fill[ezs](5.8,-161)--(8.1,-161.9)--(73.4,8.3)--(71.1,9.2)--cycle;
				\definecolor{ezs}{RGB}{252,182,17}\fill[ezs](7.7,-161.7)--(10,-162.6)--(75.3,7.6)--(73,8.5)--cycle;
				\definecolor{ezs}{RGB}{252,185,17}\fill[ezs](9.6,-162.5)--(11.9,-163.4)--(77.2,6.8)--(75,7.7)--cycle;
				\definecolor{ezs}{RGB}{251,189,18}\fill[ezs](11.5,-163.2)--(13.8,-164.1)--(79.2,6.1)--(76.9,7)--cycle;
				\definecolor{ezs}{RGB}{251,191,22}\fill[ezs](13.4,-163.9)--(15.7,-164.8)--(81.1,5.4)--(78.8,6.2)--cycle;
				\definecolor{ezs}{RGB}{252,193,25}\fill[ezs](15.4,-164.7)--(17.6,-165.6)--(83,4.6)--(80.7,5.5)--cycle;
				\definecolor{ezs}{RGB}{253,195,27}\fill[ezs](17.3,-165.4)--(19.6,-166.3)--(84.9,3.9)--(82.6,4.8)--cycle;
				\definecolor{ezs}{RGB}{253,196,31}\fill[ezs](19.2,-166.1)--(21.5,-167)--(86.8,3.2)--(84.5,4)--cycle;
				\definecolor{ezs}{RGB}{254,199,35}\fill[ezs](21.1,-166.9)--(23.4,-167.8)--(88.7,2.4)--(86.4,3.3)--cycle;
				\definecolor{ezs}{RGB}{253,201,38}\fill[ezs](23,-167.6)--(25.3,-168.5)--(90.6,1.7)--(88.3,2.6)--cycle;
				\definecolor{ezs}{RGB}{253,203,40}\fill[ezs](24.9,-168.3)--(27.2,-169.2)--(92.5,1)--(90.3,1.8)--cycle;
				\definecolor{ezs}{RGB}{253,205,42}\fill[ezs](26.8,-169.1)--(29.1,-170)--(94.4,0.2)--(92.2,1.1)--cycle;
				\definecolor{ezs}{RGB}{253,207,45}\fill[ezs](28.7,-169.8)--(31,-170.7)--(96.4,-0.5)--(94.1,0.4)--cycle;
				\definecolor{ezs}{RGB}{253,210,49}\fill[ezs](30.7,-170.5)--(32.9,-171.4)--(98.3,-1.2)--(96,-0.4)--cycle;
				\definecolor{ezs}{RGB}{253,212,54}\fill[ezs](32.6,-171.3)--(34.9,-172.2)--(100.2,-2)--(97.9,-1.1)--cycle;
				\definecolor{ezs}{RGB}{253,214,59}\fill[ezs](34.5,-172)--(36.8,-172.9)--(102.1,-2.7)--(99.8,-1.8)--cycle;
				\definecolor{ezs}{RGB}{253,216,63}\fill[ezs](36.4,-172.7)--(38.7,-173.6)--(104,-3.4)--(101.7,-2.6)--cycle;
				\definecolor{ezs}{RGB}{253,218,65}\fill[ezs](38.3,-173.5)--(40.6,-174.4)--(105.9,-4.2)--(103.6,-3.3)--cycle;
				\definecolor{ezs}{RGB}{254,220,66}\fill[ezs](40.2,-174.2)--(42.5,-175.1)--(107.8,-4.9)--(105.5,-4)--cycle;
				\definecolor{ezs}{RGB}{254,222,68}\fill[ezs](42.1,-175)--(44.4,-175.8)--(109.7,-5.6)--(107.5,-4.8)--cycle;
				\definecolor{ezs}{RGB}{254,224,71}\fill[ezs](44,-175.7)--(46.3,-176.6)--(111.7,-6.4)--(109.4,-5.5)--cycle;
				\definecolor{ezs}{RGB}{254,226,75}\fill[ezs](46,-176.4)--(48.2,-177.3)--(113.6,-7.1)--(111.3,-6.2)--cycle;
				\definecolor{ezs}{RGB}{254,228,79}\fill[ezs](47.9,-177.2)--(50.1,-178)--(115.5,-7.8)--(113.2,-7)--cycle;
				\definecolor{ezs}{RGB}{254,230,82}\fill[ezs](49.8,-177.9)--(52.1,-178.8)--(117.4,-8.6)--(115.1,-7.7)--cycle;
				\definecolor{ezs}{RGB}{254,232,85}\fill[ezs](51.7,-178.6)--(54,-179.5)--(119.3,-9.3)--(117,-8.4)--cycle;
				\definecolor{ezs}{RGB}{254,234,88}\fill[ezs](53.6,-179.4)--(55.9,-180.2)--(121.2,-10)--(118.9,-9.2)--cycle;
				\definecolor{ezs}{RGB}{254,237,91}\fill[ezs](55.5,-180.1)--(57.8,-181)--(123.1,-10.8)--(120.8,-9.9)--cycle;
				\definecolor{ezs}{RGB}{254,239,94}\fill[ezs](57.4,-180.8)--(59.7,-181.7)--(125,-11.5)--(122.8,-10.6)--cycle;
				\definecolor{ezs}{RGB}{254,241,97}\fill[ezs](59.3,-181.6)--(61.6,-182.4)--(127,-12.2)--(124.7,-11.4)--cycle;
				\definecolor{ezs}{RGB}{254,243,101}\fill[ezs](61.2,-182.3)--(63.5,-183.2)--(128.9,-13)--(126.6,-12.1)--cycle;
				\definecolor{ezs}{RGB}{254,245,103}\fill[ezs](63.2,-183)--(65.4,-183.9)--(130.8,-13.7)--(128.5,-12.8)--cycle;
				\definecolor{ezs}{RGB}{254,245,105}\fill[ezs](65.1,-183.8)--(67.4,-184.6)--(132.7,-14.4)--(130.4,-13.6)--cycle;
				\definecolor{ezs}{RGB}{254,245,105}\fill[ezs](67,-184.5)--(69.3,-185.4)--(134.6,-15.2)--(132.3,-14.3)--cycle;
				\definecolor{ezs}{RGB}{254,244,103}\fill[ezs](68.9,-185.2)--(71.2,-186.1)--(136.5,-15.9)--(134.2,-15)--cycle;
				\definecolor{ezs}{RGB}{253,243,101}\fill[ezs](70.8,-186)--(73.1,-186.8)--(138.4,-16.6)--(136.1,-15.8)--cycle;
				\definecolor{ezs}{RGB}{253,241,98}\fill[ezs](72.7,-186.7)--(75,-187.6)--(140.3,-17.4)--(138,-16.5)--cycle;
				\definecolor{ezs}{RGB}{253,239,95}\fill[ezs](74.6,-187.4)--(76.9,-188.3)--(142.2,-18.1)--(140,-17.2)--cycle;
				\definecolor{ezs}{RGB}{253,236,92}\fill[ezs](76.5,-188.2)--(78.8,-189)--(144.2,-18.8)--(141.9,-18)--cycle;
				\definecolor{ezs}{RGB}{253,234,88}\fill[ezs](78.5,-188.9)--(80.7,-189.8)--(146.1,-19.6)--(143.8,-18.7)--cycle;
				\definecolor{ezs}{RGB}{253,232,85}\fill[ezs](80.4,-189.6)--(82.7,-190.5)--(148,-20.3)--(145.7,-19.4)--cycle;
				\definecolor{ezs}{RGB}{253,230,82}\fill[ezs](82.3,-190.4)--(84.6,-191.2)--(149.9,-21)--(147.6,-20.2)--cycle;
				\definecolor{ezs}{RGB}{253,228,78}\fill[ezs](84.2,-191.1)--(86.5,-192)--(151.8,-21.8)--(149.5,-20.9)--cycle;
				\definecolor{ezs}{RGB}{253,226,74}\fill[ezs](86.1,-191.8)--(88.4,-192.7)--(153.7,-22.5)--(151.4,-21.6)--cycle;
				\definecolor{ezs}{RGB}{253,224,70}\fill[ezs](88,-192.6)--(90.3,-193.4)--(155.6,-23.3)--(153.3,-22.4)--cycle;
				\definecolor{ezs}{RGB}{253,222,67}\fill[ezs](89.9,-193.3)--(92.2,-194.2)--(157.5,-24)--(155.3,-23.1)--cycle;
				\definecolor{ezs}{RGB}{253,220,65}\fill[ezs](91.8,-194)--(94.1,-194.9)--(159.5,-24.7)--(157.2,-23.8)--cycle;
				\definecolor{ezs}{RGB}{253,217,63}\fill[ezs](93.7,-194.8)--(96,-195.6)--(161.4,-25.5)--(159.1,-24.6)--cycle;
				\definecolor{ezs}{RGB}{253,215,60}\fill[ezs](95.7,-195.5)--(97.9,-196.4)--(163.3,-26.2)--(161,-25.3)--cycle;
				\definecolor{ezs}{RGB}{253,213,56}\fill[ezs](97.6,-196.2)--(99.9,-197.1)--(165.2,-26.9)--(162.9,-26)--cycle;
				\definecolor{ezs}{RGB}{252,211,53}\fill[ezs](99.5,-197)--(101.8,-197.8)--(167.1,-27.7)--(164.8,-26.8)--cycle;
				\definecolor{ezs}{RGB}{252,209,49}\fill[ezs](101.4,-197.7)--(103.7,-198.6)--(169,-28.4)--(166.7,-27.5)--cycle;
				\definecolor{ezs}{RGB}{252,207,45}\fill[ezs](103.3,-198.4)--(105.6,-199.3)--(170.9,-29.1)--(168.6,-28.2)--cycle;
				\definecolor{ezs}{RGB}{253,205,42}\fill[ezs](105.2,-199.2)--(107.5,-200)--(172.8,-29.9)--(170.6,-29)--cycle;
				\definecolor{ezs}{RGB}{253,202,40}\fill[ezs](107.1,-199.9)--(109.4,-200.8)--(174.7,-30.6)--(172.5,-29.7)--cycle;
				\definecolor{ezs}{RGB}{253,200,38}\fill[ezs](109,-200.6)--(111.3,-201.5)--(176.7,-31.3)--(174.4,-30.4)--cycle;
				\definecolor{ezs}{RGB}{253,198,35}\fill[ezs](111,-201.4)--(113.2,-202.2)--(178.6,-32.1)--(176.3,-31.2)--cycle;
				\definecolor{ezs}{RGB}{253,196,31}\fill[ezs](112.9,-202.1)--(115.2,-203)--(180.5,-32.8)--(178.2,-31.9)--cycle;
				\definecolor{ezs}{RGB}{253,194,28}\fill[ezs](114.8,-202.8)--(117.1,-203.7)--(182.4,-33.5)--(180.1,-32.6)--cycle;
				\definecolor{ezs}{RGB}{253,192,26}\fill[ezs](116.7,-203.6)--(119,-204.5)--(184.3,-34.3)--(182,-33.4)--cycle;
				\definecolor{ezs}{RGB}{253,189,24}\fill[ezs](118.6,-204.3)--(120.9,-205.2)--(186.2,-35)--(183.9,-34.1)--cycle;
				\definecolor{ezs}{RGB}{253,186,20}\fill[ezs](120.5,-205)--(122.8,-205.9)--(188.1,-35.7)--(185.8,-34.8)--cycle;
				\definecolor{ezs}{RGB}{253,184,17}\fill[ezs](122.4,-205.8)--(124.7,-206.7)--(190,-36.5)--(187.8,-35.6)--cycle;
				\definecolor{ezs}{RGB}{252,183,15}\fill[ezs](124.3,-206.5)--(126.6,-207.4)--(192,-37.2)--(189.7,-36.3)--cycle;
				\definecolor{ezs}{RGB}{252,181,14}\fill[ezs](126.3,-207.2)--(128.5,-208.1)--(193.9,-37.9)--(191.6,-37.1)--cycle;
				\definecolor{ezs}{RGB}{252,179,13}\fill[ezs](128.2,-208)--(130.4,-208.9)--(195.8,-38.7)--(193.5,-37.8)--cycle;
				\definecolor{ezs}{RGB}{252,178,12}\fill[ezs](130.1,-208.7)--(132.4,-209.6)--(197.7,-39.4)--(195.4,-38.5)--cycle;
				\definecolor{ezs}{RGB}{251,178,10}\fill[ezs](132,-209.4)--(134.3,-210.3)--(199.6,-40.1)--(197.3,-39.3)--cycle;
				\definecolor{ezs}{RGB}{251,178,10}\fill[ezs](133.9,-210.2)--(136.2,-211.1)--(201.5,-40.9)--(199.2,-40)--cycle;
				\definecolor{ezs}{RGB}{250,177,11}\fill[ezs](135.8,-210.9)--(138.1,-211.8)--(203.4,-41.6)--(201.1,-40.7)--cycle;
				\definecolor{ezs}{RGB}{250,178,11}\fill[ezs](137.7,-211.6)--(140,-212.5)--(205.3,-42.3)--(203.1,-41.5)--cycle;
				\definecolor{ezs}{RGB}{248,180,10}\fill[ezs](139.6,-212.4)--(141.9,-213.3)--(207.3,-43.1)--(205,-42.2)--cycle;
				\definecolor{ezs}{RGB}{247,181,9}\fill[ezs](141.5,-213.1)--(143.8,-214)--(209.2,-43.8)--(206.9,-42.9)--cycle;
				\definecolor{ezs}{RGB}{246,182,9}\fill[ezs](143.5,-213.8)--(145.7,-214.7)--(211.1,-44.5)--(208.8,-43.7)--cycle;
				\definecolor{ezs}{RGB}{246,182,10}\fill[ezs](145.4,-214.6)--(147.7,-215.5)--(213,-45.3)--(210.7,-44.4)--cycle;
			\end{scope}
			% chữ EDUZY (Times đậm, làm dày nét bằng 5 lớp lệch nhẹ; căn đúng khung chữ logo gốc)
			\foreach \dx/\dy in {-0.7/0,0.7/0,0/-0.3,0/0.3,0/0}
			\node[anchor=center,inner sep=0pt,outer sep=0pt,text=ezText,align=center] at (109.6+\dx,-193.6+\dy)
			{\resizebox{1.01cm}{0.22cm}{\fontencoding{T1}\fontfamily{ptm}\fontseries{b}\selectfont EDUZY}};
		\end{tikzpicture}%
	}%
}
%====================== THIẾT LẬP THÂN ĐỀ ======================
%--- Môi trường tự luận {bt}: ex_test 3.3.4 chưa khai báo sẵn nên tạo ở đây
%    (ex_test đã tự móc \fixshowans và \listenumerate cho bt)
\makeatletter
\@ifundefined{bt}{\newtheorem{bt}{\color{gold1}\namebt}}{}
\makeatother
\renewcommand{\sectionmark}[1]{\markright{\thesection\ #1}}
%--- Khung tiêu đề (logo + tối đa 4 dòng; dòng 3, 4 để trống {} thì bỏ qua)
%  Đề thi:    \khungtieude{THI THỬ THPT QUỐC GIA NĂM 2026}{Môn Toán -- Đề số 21}
%                         {(Đề thi gồm: ...)}{Thời gian làm bài: 90 phút (không kể thời gian phát đề).}
%  Chuyên đề: \khungtieude{CHUYÊN ĐỀ 3 -- DÃY SỐ VÀ CÔNG THỨC TRUY HỒI}
%                         {Bồi dưỡng học sinh giỏi Toán THPT}{...}{}
\newcommand{\khungtieude}[4]{%
	\begin{tcolorbox}[colframe=gold1,colback=gold5!45,boxrule=0.5mm]
		\begin{minipage}[c]{3cm}
			\centerline{\eduzylogo[2.4cm]}
		\end{minipage}\hspace{0.5cm}
		\begin{minipage}[c]{12.5cm}
			\vskip 0.15 cm
			\centerline{\ind{#1}}
			\vskip 0.15 cm
			\centerline{\textbf{#2}}
			\if\relax\detokenize{#3}\relax\else
				\vskip 0.15 cm
				\centerline{\textit{#3}}
			\fi
			\if\relax\detokenize{#4}\relax\else
				\vskip 0.15 cm
				\centerline{\textit{#4}}
			\fi
			\vskip 0.15 cm
		\end{minipage}
	\end{tcolorbox}
}
%====================== BỔ SUNG CHO BÀI HỌC, CHUYÊN ĐỀ ======================
\usepackage[shortlabels]{enumitem}
\usepackage{tabularx}
%--- Danh sách gọn: \begin{enumerate}[gon] ... \end{enumerate}
\SetEnumitemKey{gon}{leftmargin=*,itemsep=1pt,parsep=0pt,topsep=3pt}
%--- Khung lý thuyết: \begin{khung}{} ... \end{khung}; tiêu đề chỉ dùng khi thật ngắn ("Mẹo", "Lỗi thường gặp")
\newtcolorbox{khung}[1]{breakable,enhanced,colframe=gold1,colback=gold5!45,
	boxrule=0.4mm,fonttitle=\bfseries\sffamily,coltitle=white,colbacktitle=gold1,
	attach boxed title to top left={xshift=6mm,yshift=-3mm},
	boxed title style={arc=1mm,boxrule=0.3mm},title={#1}}
%--- Tiêu đề mục nhỏ trong một dạng: \muc{Dấu hiệu}, \muc{Phương pháp}, \muc{Bài tự luyện}
\makeatletter
\newcommand{\muc}[1]{\par\medskip\noindent\indam{#1}\par\nobreak\smallskip\@afterindentfalse\@afterheading}
\makeatother
%--- Ô công thức cần nhớ (căn giữa, viền vàng, vừa khít nội dung)
%     \begin{center}\congthuc{$Q(n)=(n-1)!$}\end{center}
\newtcbox{\congthuc}{on line,colframe=gold1,colback=gold5!45,boxrule=0.5pt,arc=2pt,
	left=8pt,right=8pt,top=4pt,bottom=4pt}
%--- Công tắc lời giải bài tự luyện (ví dụ mẫu luôn dùng \loigiai)
%     \loigiaiTL[{đáp số}]{lời giải đầy đủ}
%     \TLfull -> in lời giải đầy đủ;  \TLdapso -> chỉ in một dòng "Đáp số";  \TLan -> ẩn hoàn toàn
%     (\CheDo là dethi thì mọi lời giải và đáp số đều ẩn)
\makeatletter
\newif\ifTLhien   \TLhientrue
\newif\ifTLdapso  \TLdapsofalse
\newcommand{\TLfull}{\TLhientrue\TLdapsofalse}
\newcommand{\TLdapso}{\TLhientrue\TLdapsotrue}
\newcommand{\TLan}{\TLhienfalse}
\newcommand{\loigiaiTL}[2][]{%
	\ifTLhien
		\ifTLdapso
			\def\@tlarg{#1}%
			\ifx\@tlarg\@empty\else
				\ifdefstring{\kindOPTN}{dethi}{}{\par\noindent{\bfseries\color{gold1}Đáp số:} #1\par}%
			\fi
		\else
			\loigiai{#2}%
		\fi
	\fi}
\makeatother
\TLdapso% mặc định: chỉ in đáp số bài tự luyện
%====================== BẢNG ĐÚNG/SAI (\choiceTFt) ======================
%  Hàng tiêu đề nền nâu đồng chữ trắng (Phát biểu | Đ | S), các hàng tô xen kẽ kem / trắng,
%  đường kẻ vàng kim, nhãn a) b) c) d) màu nâu đồng. Bản có đáp án: cột được chọn có vòng tròn
%  nâu đồng với dấu tích trắng (câu đúng tích ở cột Đúng, câu sai tích ở cột Sai).
%  Tiêu đề ghi đủ "Đúng | Sai": \OPTN{viettat=0} sau \input (cột Đúng/Sai rộng ra, phát biểu hẹp lại).
%  Dùng bảng gốc của ex_test: \def\BangDSgoc{1} trước %============================================================
%  PREAMBLE EDUZY v2 - dùng chung cho đề thi, bài học, chuyên đề (gói ex_test)
%  Khác bản preamble-eduzy.tex ở đầu trang, chân trang: dải nền nâu đồng tràn ngang,
%  logo + Eduzy.com.vn, ô chữ vàng kem cạnh xiên, chân trang "Trang (số)";
%  logo Eduzy chìm giữa mọi trang (tắt: \logochimtat, bật lại: \logochimbat).
%  Cách dùng ở file chính (xem mau-chuan.tex, mau-chuyen-de.tex):
%      \def\CheDo{loigiai}          % loigiai | dethi | solcolor | color | book
%      \def\tenchuong{<dòng chữ trên đầu trang>}
%      \def\thuonghieu{<chữ ở chân trang lẻ>}   % không khai báo thì dùng mặc định
%      \def\tacgia{<chữ ở chân trang chẵn>}      % không khai báo thì dùng mặc định
%      \def\khauhieu{<chữ cạnh logo ở đầu trang>}  % tùy chọn, mặc định Eduzy.com.vn
%      %============================================================
%  PREAMBLE EDUZY v2 - dùng chung cho đề thi, bài học, chuyên đề (gói ex_test)
%  Khác bản preamble-eduzy.tex ở đầu trang, chân trang: dải nền nâu đồng tràn ngang,
%  logo + Eduzy.com.vn, ô chữ vàng kem cạnh xiên, chân trang "Trang (số)";
%  logo Eduzy chìm giữa mọi trang (tắt: \logochimtat, bật lại: \logochimbat).
%  Cách dùng ở file chính (xem mau-chuan.tex, mau-chuyen-de.tex):
%      \def\CheDo{loigiai}          % loigiai | dethi | solcolor | color | book
%      \def\tenchuong{<dòng chữ trên đầu trang>}
%      \def\thuonghieu{<chữ ở chân trang lẻ>}   % không khai báo thì dùng mặc định
%      \def\tacgia{<chữ ở chân trang chẵn>}      % không khai báo thì dùng mặc định
%      \def\khauhieu{<chữ cạnh logo ở đầu trang>}  % tùy chọn, mặc định Eduzy.com.vn
%      \input{preamble-eduzy-v2}
%      \TLdapso                      % tài liệu: \TLfull | \TLdapso | \TLan
%      \begin{document} ... \end{document}
%  Logo Eduzy vẽ bằng TikZ ngay trong file (\eduzylogo), không cần file ảnh.
%  Biên dịch: pdfLaTeX (máy có gói vietnam) hoặc XeLaTeX.
%============================================================
\documentclass[12pt,a4paper,twoside]{article}%<-- twoside: bat dau trang/chan trang chan-le
\usepackage{iftex}
\usepackage{amsmath,amssymb,tikz,mathrsfs,tkz-tab}
\usepackage{graphicx}
\usepackage{pgfplots}% đồ thị hàm số
\pgfplotsset{compat=1.17}
%\usepackage{mathptmx}
\usepackage{mathpazo}
\ifXeTeX% XeLaTeX: font chữ có đủ dấu tiếng Việt, giữ nguyên font toán mathpazo
	\usepackage[no-math]{fontspec}
	\setmainfont{TeX Gyre Pagella}
	\newfontfamily\EZqag{TeX Gyre Adventor}[NFSSFamily=qag,Ligatures=TeX]
\fi
\usepackage[many]{tcolorbox}
\usetikzlibrary{math,through,calc,intersections,angles,quotes,shapes,shapes.geometric,arrows,patterns,matrix,chains,arrows.meta,decorations.shapes,decorations.fractals,decorations.markings,shadows,decorations.pathreplacing}
\usepackage[top=1.6cm, bottom=1.8cm, left=1.8cm, right=1.8cm]{geometry}% v2: top 1.6cm (bản cũ 1.5cm) để chữ không sát dải đầu trang
\providecommand{\CheDo}{loigiai}% loigiai | dethi | solcolor | color | book
\expandafter\usepackage\expandafter[\CheDo]{ex_test}
% \tenchuong: khai báo ở file chính (trước \input) hoặc \renewcommand sau \input
%============================================================
%  [GỘP TỪ eduzy-de.tex] EDUZY - MÀU NÂU ĐỒNG / VÀNG KIM
%  + ĐẦU TRANG, CHÂN TRANG KIỂU DẢI NGANG (v2), ĐỐI XỨNG TRANG CHẴN / TRANG LẺ.
%  (không cần file eduzy-de.tex riêng nữa)
%  Khối này phải đặt SAU \usepackage[...]{ex_test} và SAU \def\tenchuong{...}
%============================================================
\usepackage{lastpage}
\usepackage{needspace}
%--- DÀN TRANG ĐỀ THI: tách câu (THÊM 05/11/2026; chỉ thêm lệnh mới, KHÔNG đổi hành vi mặc định) ---
%  Hai công tắc, mặc định tắt (= preamble chạy y như cũ). Bật bằng cách đặt TRƯỚC \input{preamble-eduzy-v2}
%  trong file đề (tools/dan_trang/dan_trang.py tự ghi khi đặt \NeedspaceCau):
%    \def\TachCauDS{1}   câu Đúng/Sai được tách: câu dẫn (kèm hình) ở cuối trang, bảng Phát biểu sang trang sau
%                        (câu dẫn và hình không tách nhau, bảng không tách giữa hai trang)
%    \def\TachCauLG{1}   bản lời giải (\CheDo khác dethi): phần đề của câu cao hơn 1/3 trang được phép tách
%  Lệnh dùng trong file đề (giá trị do tools/dan_trang/dan_trang.py đo ở bản dethi):
%    \NeedspaceCau[<cao câu dẫn + hình>]{<cao cả khối câu>}   ngay trên \begin{ex} (câu đầu phần: trên tiêu đề phần)
%        công tắc tắt: đúng bằng \Needspace{<cao cả khối câu>} như cũ;
%        \TachCauDS=1 và có [..] (câu Đúng/Sai được chọn tách): chỉ giữ câu dẫn + hình;
%        \TachCauLG=1 ở bản lời giải: đòi tối đa 1/3 \textheight.
%    \NeedspaceBang{<cao bảng Phát biểu>}%   ngay trước \choiceTFt: \TachCauDS=1 thì giữ trọn bảng trên một trang;
%        công tắc tắt: không làm gì.
\providecommand{\TachCauDS}{0}
\providecommand{\TachCauLG}{0}
\newlength{\EZnsCau}
\newcommand{\NeedspaceCau}[2][]{%
	\setlength{\EZnsCau}{#2}%
	\ifnum\TachCauDS=1 \if\relax\detokenize{#1}\relax\else\setlength{\EZnsCau}{#1}\fi\fi
	\ifnum\TachCauLG=1 \ifdefstring{\kindOPTN}{dethi}{}{\ifdim\EZnsCau>0.3333\textheight\setlength{\EZnsCau}{0.3333\textheight}\fi}\fi
	\Needspace{\EZnsCau}}
\newcommand{\NeedspaceBang}[1]{\ifnum\TachCauDS=1 \Needspace{#1}\fi}

%--- Thang màu: đậm -> nhạt
\definecolor{gold1}{RGB}{126,96,0}    % nâu đồng đậm  #7E6000
\definecolor{gold2}{RGB}{192,144,3}   % vàng kim      #C09003
\definecolor{gold3}{RGB}{255,215,99}  % vàng nhạt     #FFD763
\definecolor{gold4}{RGB}{255,231,153} % vàng rất nhạt #FFE799
\definecolor{gold5}{RGB}{255,243,205} % vàng kem      #FFF3CD
%--- Ghi đè tông xanh/đỏ cũ
\definecolor{myblue}{RGB}{126,96,0}
\colorlet{blue}{gold1}
\colorlet{cyan}{gold2}
\colorlet{teal}{gold2}

%--- Tên chương mặc định (file đề nên tự khai báo lại)
\providecommand{\tenchuong}{}
%--- Chân trang
\providecommand{\thuonghieu}{TOÁN 12 -- Eduzy.com.vn}
\providecommand{\tacgia}{Tô Gia Bảo -- 0353703783}

%--- Chữ cạnh logo ở đầu trang
\providecommand{\khauhieu}{Eduzy.com.vn}

%====================== THÔNG SỐ ĐẦU TRANG, CHÂN TRANG (v2) ======================
%  Dải nền đậm tràn ngang trang; ô chữ vàng kem có cạnh xiên, kèm 2 vạch xiên chuyển màu.
%  Đầu trang lẻ : logo + \khauhieu (trái)   | ô chữ \tenchuong (phải)
%  Chân trang lẻ: ô chữ \thuonghieu (trái)  | Trang (số) (phải)
%  Trang chẵn soi gương trang lẻ; chân trang chẵn ghi \tacgia.
%--- Màu: chỉ dùng thang vàng kim Eduzy
\colorlet{ezNen}{gold1}% dải nền đậm
\colorlet{ezVach1}{gold3}% vạch xiên sát dải nền (đậm hơn)
\colorlet{ezVach2}{gold4}% vạch xiên sát ô chữ (nhạt hơn)
\colorlet{ezO}{gold5}% ô chữ vàng kem
%--- Kích thước
\newlength{\EZdai}\setlength{\EZdai}{1.1cm}% chiều cao dải nền
\newlength{\EZo}\setlength{\EZo}{0.68cm}% chiều cao ô chữ
\newlength{\EZxien}\setlength{\EZxien}{0.55cm}% độ lệch ngang của cạnh xiên
\newlength{\EZvach}\setlength{\EZvach}{0.75cm}% bề rộng mỗi vạch xiên
\newlength{\EZkhe}\setlength{\EZkhe}{0.12cm}% khe trắng giữa vạch và ô chữ
\newlength{\EZdem}\setlength{\EZdem}{0.4cm}% khoảng đệm từ chữ tới cạnh xiên
\newlength{\EZle}\setlength{\EZle}{1.8cm}% lề trái/phải của chữ (bằng lề thân bài)
\newlength{\EZrongmin}\setlength{\EZrongmin}{4.5cm}% ô chữ rộng tối thiểu
\newlength{\EZrong}% bề rộng chữ trong ô (tính tự động)
%--- Font
\newcommand{\EZfontO}{\fontsize{10.5pt}{12pt}\selectfont\bfseries}% chữ trong ô
\newcommand{\EZfontKH}{\fontfamily{qag}\fontsize{11pt}{12pt}\selectfont\bfseries}% \khauhieu cạnh logo
\newcommand{\EZfontTrang}{\fontsize{11.5pt}{12pt}\selectfont\bfseries}% chữ "Trang" và số trang
%--- Đo bề rộng chữ trong ô
\newcommand{\EZdorong}[1]{%
	\settowidth{\EZrong}{\EZfontO #1}%
	\ifdim\EZrong<\EZrongmin\setlength{\EZrong}{\EZrongmin}\fi}
%--- Logo + chữ trắng \khauhieu (dựng sẵn một lần khi vào \begin{document})
\newsavebox{\EZlogobox}
\AtBeginDocument{\sbox{\EZlogobox}{%
		\raisebox{-0.5\height}{\eduzylogo[1cm]}\hspace{0.15cm}%
		\raisebox{-0.5\height}{\color{white}\EZfontKH\khauhieu}}}
%--- Hai vạch xiên + khe trắng. #1: điểm gốc ở mép dưới, #2: hướng (1 sang phải, -1 sang trái),
%    #3: hướng nghiêng (1 kiểu "/", -1 kiểu "\")
\newcommand{\EZvachxien}[3]{%
	\fill[white] (#1) --++({(#3)*\EZxien},\EZo) --++({(#2)*\EZkhe},0) --++({-(#3)*\EZxien},-\EZo) --cycle;
	\fill[ezVach2] ([xshift=(#2)*\EZkhe]#1) --++({(#3)*\EZxien},\EZo) --++({(#2)*\EZvach},0) --++({-(#3)*\EZxien},-\EZo) --cycle;
	\fill[ezVach1] ([xshift=(#2)*(\EZkhe+\EZvach)]#1) --++({(#3)*\EZxien},\EZo) --++({(#2)*\EZvach},0) --++({-(#3)*\EZxien},-\EZo) --cycle;
	\draw[white,line width=0.6pt] ([xshift=(#2)*(\EZkhe+\EZvach)]#1) --++({(#3)*\EZxien},\EZo);
}
%--- Chữ trong ô (căn giữa theo chiều cao ô, không lệch vì dấu tiếng Việt)
\tikzset{ezchu/.style={inner sep=0pt,text=ezNen,text height=0.72em,text depth=0.22em}}

%====================== LOGO CHÌM GIỮA TRANG ======================
%  Vẽ cùng đầu trang nên nằm dưới chữ của thân bài, có trên mọi trang.
%  Tắt: \logochimtat (từ trang đó trở đi), bật lại: \logochimbat.
%  Đổi cỡ, độ đậm: khai báo trước \input, ví dụ \def\EZchimcao{9cm} \def\EZchimdo{0.1}
\providecommand{\EZchimcao}{10cm}% chiều cao logo chìm
\providecommand{\EZchimdo}{0.12}% độ đậm (0 = trong suốt, 1 = rõ hẳn)
\newif\ifEZchim \EZchimtrue
\newcommand{\logochimtat}{\EZchimfalse}
\newcommand{\logochimbat}{\EZchimtrue}
%  Dựng sẵn logo mờ một lần khi vào \begin{document}; transparency group làm mờ cả logo
%  một lượt nên không lộ vệt ở chỗ các dải màu chồng nhau.
\newsavebox{\EZchimtho}
\newsavebox{\EZchimbox}
\AtBeginDocument{%
	\sbox{\EZchimtho}{\eduzylogo[\EZchimcao]}%
	\sbox{\EZchimbox}{\begin{tikzpicture}
			\begin{scope}[transparency group,opacity=\EZchimdo]
				\node[inner sep=0pt] {\usebox{\EZchimtho}};
			\end{scope}
	\end{tikzpicture}}}
\newcommand{\EZvechim}{%
	\ifEZchim
		\node[inner sep=0pt] at (current page.center) {\usebox{\EZchimbox}};
	\fi}

%====================== ĐẦU TRANG - TRANG LẺ ======================
%  (logo bên trái, ô chữ \tenchuong bên phải)
\newcommand{\myfancyheadOdd}{%
	\EZdorong{\tenchuong}%
	\begin{tikzpicture}[remember picture,overlay]
		\coordinate (T) at (current page.north west);
		\coordinate (P) at (current page.north east);
		%--- logo chìm giữa trang
		\EZvechim
		%--- dải nền
		\fill[ezNen] (T) rectangle ([yshift=-\EZdai]P);
		%--- ô chữ: X là góc dưới bên trái, cạnh trái xiên kiểu "/"
		\coordinate (X) at ([xshift=-\EZle-\EZrong-\EZdem-\EZxien,yshift=-\EZdai]P);
		\fill[ezO] (X) --++(\EZxien,\EZo) -- ([yshift=-\EZdai+\EZo]P) -- ([yshift=-\EZdai]P) --cycle;
		\EZvachxien{X}{-1}{1}
		\node[ezchu,anchor=east] at ([xshift=-\EZle,yshift=-\EZdai+0.5*\EZo]P) {\EZfontO\tenchuong};
		%--- logo + tên
		\node[inner sep=0pt,anchor=west] at ([xshift=\EZle,yshift=-0.5*\EZdai]T) {\usebox{\EZlogobox}};
	\end{tikzpicture}%
}

%====================== ĐẦU TRANG - TRANG CHẴN =====================
%  (soi gương: ô chữ \tenchuong bên trái, logo bên phải)
\newcommand{\myfancyheadEven}{%
	\EZdorong{\tenchuong}%
	\begin{tikzpicture}[remember picture,overlay]
		\coordinate (T) at (current page.north west);
		\coordinate (P) at (current page.north east);
		%--- logo chìm giữa trang
		\EZvechim
		%--- dải nền
		\fill[ezNen] (T) rectangle ([yshift=-\EZdai]P);
		%--- ô chữ: X là góc dưới bên phải, cạnh phải xiên kiểu "\"
		\coordinate (X) at ([xshift=\EZle+\EZrong+\EZdem+\EZxien,yshift=-\EZdai]T);
		\fill[ezO] (X) --++(-\EZxien,\EZo) -- ([yshift=-\EZdai+\EZo]T) -- ([yshift=-\EZdai]T) --cycle;
		\EZvachxien{X}{1}{-1}
		\node[ezchu,anchor=west] at ([xshift=\EZle,yshift=-\EZdai+0.5*\EZo]T) {\EZfontO\tenchuong};
		%--- logo + tên
		\node[inner sep=0pt,anchor=east] at ([xshift=-\EZle,yshift=-0.5*\EZdai]P) {\usebox{\EZlogobox}};
	\end{tikzpicture}%
}

%====================== CHÂN TRANG - TRANG LẺ ======================
%  (ô chữ \thuonghieu bên trái, "Trang (số)" bên phải)
\newcommand{\myfancyfootOdd}{%
	\EZdorong{\thuonghieu}%
	\begin{tikzpicture}[remember picture,overlay]
		\coordinate (T) at (current page.south west);
		\coordinate (P) at (current page.south east);
		%--- dải nền
		\fill[ezNen] (T) rectangle ([yshift=\EZdai]P);
		%--- ô chữ sát mép trên dải: X là góc dưới bên phải, cạnh phải xiên kiểu "/"
		\coordinate (X) at ([xshift=\EZle+\EZrong+\EZdem,yshift=\EZdai-\EZo]T);
		\fill[ezO] ([yshift=\EZdai-\EZo]T) -- (X) --++(\EZxien,\EZo) -- ([yshift=\EZdai]T) --cycle;
		\EZvachxien{X}{1}{1}
		\node[ezchu,anchor=west] at ([xshift=\EZle,yshift=\EZdai-0.5*\EZo]T) {\EZfontO\thuonghieu};
		%--- số trang
		\node[circle,draw=white,line width=0.9pt,text=white,inner sep=0pt,minimum size=0.62cm,anchor=east]
		(so) at ([xshift=-\EZle,yshift=0.5*\EZdai]P) {\EZfontTrang\thepage};
		\node[inner sep=0pt,text=white,anchor=base east] (tr) at ([xshift=-0.15cm]so.base west) {\EZfontTrang Trang};
	\end{tikzpicture}%
}

%====================== CHÂN TRANG - TRANG CHẴN ====================
%  (soi gương: "Trang (số)" bên trái, ô chữ \tacgia bên phải)
\newcommand{\myfancyfootEven}{%
	\EZdorong{\tacgia}%
	\begin{tikzpicture}[remember picture,overlay]
		\coordinate (T) at (current page.south west);
		\coordinate (P) at (current page.south east);
		%--- dải nền
		\fill[ezNen] (T) rectangle ([yshift=\EZdai]P);
		%--- ô chữ sát mép trên dải: X là góc dưới bên trái, cạnh trái xiên kiểu "\"
		\coordinate (X) at ([xshift=-\EZle-\EZrong-\EZdem,yshift=\EZdai-\EZo]P);
		\fill[ezO] ([yshift=\EZdai-\EZo]P) -- (X) --++(-\EZxien,\EZo) -- ([yshift=\EZdai]P) --cycle;
		\EZvachxien{X}{-1}{-1}
		\node[ezchu,anchor=east] at ([xshift=-\EZle,yshift=\EZdai-0.5*\EZo]P) {\EZfontO\tacgia};
		%--- số trang
		%    (vòng tròn ẩn so0 nằm giữa dải để lấy đường chân chữ, giống trang lẻ)
		\node[circle,inner sep=0pt,minimum size=0.62cm] (so0) at ([yshift=0.5*\EZdai]T) {\EZfontTrang\phantom{\thepage}};
		\node[inner sep=0pt,text=white,anchor=base west] (tr) at ($(so0.base -| T)+(\EZle,0)$) {\EZfontTrang Trang};
		\node[circle,draw=white,line width=0.9pt,text=white,inner sep=0pt,minimum size=0.62cm,anchor=base west]
		(so) at ([xshift=0.15cm]tr.base east) {\EZfontTrang\thepage};
	\end{tikzpicture}%
}

%--- Giữ tương thích với các file cũ đang gọi \myfancyhead, \myfancyfoot
\newcommand{\myfancyhead}{\myfancyheadOdd}
\newcommand{\myfancyfoot}{\myfancyfootOdd}

%======================= GẮN VÀO TRANG =============================
\pagestyle{fancy}
\fancyhf{}
\fancyhead[LO]{\myfancyheadOdd}
\fancyhead[RE]{\myfancyheadEven}
\fancyfoot[LO]{\myfancyfootOdd}
\fancyfoot[RE]{\myfancyfootEven}
\renewcommand{\headrulewidth}{0pt}
\renewcommand{\footrulewidth}{0pt}
\fancypagestyle{plain}{%
	\fancyhf{}
	\fancyhead[LO]{\myfancyheadOdd}
	\fancyhead[RE]{\myfancyheadEven}
	\fancyfoot[LO]{\myfancyfootOdd}
	\fancyfoot[RE]{\myfancyfootEven}
	\renewcommand{\headrulewidth}{0pt}
	\renewcommand{\footrulewidth}{0pt}
}
%=================== HẾT PHẦN GỘP TỪ eduzy-de.tex ===================
\newcommand{\indam}[1]{
	{\fontsize{12pt}{9pt}\sffamily \selectfont\bfseries\color{gold1} #1}
}
\newcommand{\ind}[1]{
	{\fontsize{14pt}{9pt}\fontfamily{qag} \selectfont\bfseries\color{gold1} #1}
}
\renewcommand{\FalseEX}{\stepcounter{dapan}\textbf{\selectfont\color{gold1}\Alph{dapan}.}\ignorespaces}
\renewtheorem{ex}{\color{gold1}Câu}
\renewcommand{\baselinestretch}{1.1}
\newcommand{\hoac}[1]{
	\left[\begin{aligned}#1\end{aligned}\right.}
\newcommand{\heva}[1]{
	\left\{\begin{aligned}#1\end{aligned}\right.}
\usepackage{esvect}
\def\vec{\vv}
\def\vt{\vv}
\def\overrightarrow{\vv}
%====================== LOGO EDUZY VẼ BẰNG TIKZ ======================
%  Đường viền dò từ logo gốc; dải màu vàng kim đo trực tiếp trên ảnh gốc
%  (nghiêng 21 độ, tô bằng 80 dải hẹp nên hiển thị giống nhau trên mọi trình xem PDF).
%  Dùng: \eduzylogo          -> cao 2.4cm (mặc định)
%        \eduzylogo[1.5cm]   -> cao 1.5cm
\definecolor{ezText}{RGB}{245,200,46}
\newcommand{\eduzylogo}[1][2.4cm]{%
	\resizebox{!}{#1}{%
		\begin{tikzpicture}[x=0.01cm,y=0.01cm]
			\path[use as bounding box] (0,0) rectangle (224,-224);
			% mũ tốt nghiệp + vương miện
			\begin{scope}
				\clip
				(119.4,-176.1)..controls(118.9,-176.1)and(116.8,-176)..(114.6,-176)
				..controls(101.4,-176)and(85.8,-175.9)..(83.7,-175.7)
				..controls(81.8,-175.6)and(81,-175.6)..(80.2,-175.7)
				..controls(78.8,-175.9)and(77.7,-175.9)..(76.2,-175.7)
				..controls(74.9,-175.5)and(70.8,-175.5)..(64.3,-175.7)
				..controls(59.1,-175.9)and(58.3,-175.7)..(58.6,-173.7)
				..controls(58.6,-172.8)and(59.6,-169.6)..(60.3,-168)
				..controls(60.8,-166.9)and(61.4,-165.8)..(62.3,-165)
				..controls(63.1,-164.2)and(64.3,-163.4)..(65,-163.2)
				..controls(65.5,-163.1)and(72.3,-163)..(76.6,-163.1)
				..controls(77.9,-163.1)and(80.1,-163.1)..(81.4,-163)
				..controls(82.8,-162.9)and(86.3,-162.9)..(89.2,-162.9)
				..controls(93,-162.9)and(94.9,-162.8)..(96,-162.7)
				..controls(97.2,-162.5)and(101,-162.5)..(114,-162.4)
				..controls(123,-162.3)and(134.9,-162.2)..(140.2,-162.1)
				..controls(150.1,-162)and(152.2,-162.1)..(153.6,-162.4)
				..controls(154.4,-162.6)and(157.1,-163.8)..(158.1,-164.4)
				..controls(158.8,-164.8)and(159.7,-165.7)..(159.9,-166.3)
				..controls(160,-166.5)and(160.2,-167.1)..(160.2,-167.6)
				..controls(160.3,-168)and(160.6,-169.4)..(160.9,-170.6)
				..controls(161.2,-171.7)and(161.5,-173.1)..(161.6,-173.6)
				..controls(161.8,-174.4)and(161.8,-174.5)..(161.5,-174.9)
				..controls(161.2,-175.6)and(160.8,-175.9)..(159.9,-176)
				..controls(159.2,-176.1)and(154.1,-176.2)..(151.5,-176.2)
				..controls(148.1,-176.1)and(122.6,-176.1)..(121.6,-176.2)
				..controls(120.8,-176.2)and(119.8,-176.2)..(119.4,-176.1)--cycle
				(130.4,-160)..controls(129,-160)and(119.8,-160)..(110.1,-159.9)
				..controls(100.3,-159.8)and(88.5,-159.7)..(83.9,-159.6)
				..controls(77.8,-159.5)and(75.3,-159.5)..(74.5,-159.6)
				..controls(73.9,-159.7)and(72.5,-159.8)..(71.3,-159.8)
				..controls(69.6,-159.8)and(69,-159.8)..(68.5,-159.6)
				..controls(67.6,-159.3)and(66.8,-158.6)..(66.2,-157.6)
				..controls(66,-157.3)and(65.6,-156.5)..(65.3,-156.1)
				..controls(64.6,-155.1)and(64.5,-154.9)..(63.7,-152.9)
				..controls(63.2,-151.6)and(62.9,-150.9)..(61.8,-149.6)
				..controls(61.5,-149.2)and(61.1,-148.5)..(60.9,-147.9)
				..controls(60.7,-147.4)and(60.3,-146.6)..(60.1,-146.1)
				..controls(59.9,-145.5)and(59.5,-144.9)..(59,-144.3)
				..controls(58.5,-143.7)and(58.2,-143.2)..(57.9,-142.5)
				..controls(57.7,-141.9)and(57.3,-141.1)..(57,-140.7)
				..controls(56.6,-140.3)and(56.1,-139.3)..(55.8,-138.5)
				..controls(55.5,-137.8)and(55.1,-136.8)..(54.9,-136.5)
				..controls(54.7,-136.1)and(54.3,-135.2)..(53.9,-134.5)
				..controls(53.6,-133.8)and(53,-132.7)..(52.6,-132)
				..controls(51.6,-130.5)and(51.4,-130.2)..(51.1,-129)
				..controls(50.6,-127.5)and(50.6,-127.4)..(49.7,-126.3)
				..controls(49.2,-125.7)and(48.6,-124.8)..(48.3,-124.2)
				..controls(48,-123.6)and(47.5,-122.7)..(47.3,-122.2)
				..controls(47,-121.8)and(46.5,-120.7)..(46.1,-119.8)
				..controls(45.8,-118.9)and(45.2,-117.8)..(44.9,-117.2)
				..controls(44.2,-116.1)and(43.3,-114.4)..(42.3,-112.2)
				..controls(41.8,-111.3)and(41.3,-110.4)..(41.1,-110.1)
				..controls(40.5,-109.4)and(40.3,-109)..(40,-107.8)
				..controls(39.8,-107)and(39.6,-106.5)..(39.2,-106)
				..controls(38.9,-105.6)and(38.4,-104.7)..(38,-104)
				..controls(37.7,-103.4)and(37.2,-102.5)..(37,-102.1)
				..controls(36.7,-101.7)and(36.3,-100.8)..(36,-100.1)
				..controls(35.6,-98.9)and(35.3,-98.4)..(34.6,-97.6)
				..controls(34.2,-97.1)and(34.2,-96.3)..(34.6,-96)
				..controls(35,-95.7)and(35.5,-95.9)..(36.2,-96.7)
				..controls(36.5,-97)and(37.2,-97.5)..(37.6,-97.7)..controls(38,-98)and(38.6,-98.4)..(39.1,-98.9)
				..controls(39.6,-99.4)and(40.3,-100)..(40.6,-100.3)
				..controls(41.3,-100.8)and(42.4,-101.8)..(43.7,-102.9)
				..controls(44,-103.2)and(44.6,-103.7)..(45,-103.9)
				..controls(45.4,-104.2)and(45.9,-104.6)..(46.1,-104.8)
				..controls(47.1,-105.9)and(48.8,-107.1)..(50.5,-107.9)
				..controls(51.8,-108.6)and(52.3,-109)..(53.4,-110)
				..controls(54.8,-111.2)and(55.6,-111.6)..(57,-112.1)
				..controls(57.5,-112.3)and(58.5,-112.8)..(59.4,-113.3)
				..controls(60.2,-113.7)and(61.3,-114.3)..(61.9,-114.5)
				..controls(62.4,-114.7)and(63.3,-115.1)..(63.8,-115.3)
				..controls(64.4,-115.6)and(65.4,-115.9)..(66.2,-116.1)
				..controls(67,-116.3)and(68.4,-116.6)..(69.2,-116.9)
				..controls(72.3,-117.9)and(72.8,-118)..(78,-118.1)
				..controls(84,-118.3)and(86.1,-118.1)..(88.1,-117.1)
				..controls(89.4,-116.5)and(93,-115)..(93.7,-114.9)
				..controls(94.8,-114.6)and(95.4,-114.2)..(95.9,-113.4)
				..controls(96.2,-113)and(96.7,-112.5)..(97.1,-112.2)
				..controls(98,-111.6)and(101.3,-108.4)..(102.1,-107.2)
				..controls(102.4,-106.8)and(102.9,-106.2)..(103.2,-105.8)
				..controls(103.5,-105.5)and(104,-104.8)..(104.2,-104.2)
				..controls(104.5,-103.6)and(105,-102.8)..(105.3,-102.4)
				..controls(105.6,-102)and(106.1,-101.2)..(106.5,-100.6)
				..controls(107.1,-99.5)and(108.6,-97.5)..(109.5,-96.5)
				..controls(109.9,-96.1)and(110.1,-96)..(110.4,-96)
				..controls(111.1,-96)and(111.6,-96.5)..(111.8,-97.6)
				..controls(112,-98.6)and(112.3,-99.3)..(112.8,-99.9)
				..controls(113.7,-100.9)and(115.2,-103.1)..(115.7,-104.2)
				..controls(116.1,-105)and(116.5,-105.5)..(116.9,-106)
				..controls(117.3,-106.4)and(117.8,-106.9)..(118,-107.3)
				..controls(118.2,-107.6)and(118.6,-108)..(119,-108.3)
				..controls(119.6,-108.8)and(121.1,-110.3)..(121.8,-111.2)
				..controls(122.1,-111.7)and(122.5,-112)..(123.4,-112.4)
				..controls(124.8,-113.2)and(125.8,-113.8)..(127,-114.6)
				..controls(127.5,-114.9)and(128.1,-115.3)..(128.2,-115.4)
				..controls(128.4,-115.5)and(129.3,-115.8)..(130.2,-116)
				..controls(131.1,-116.2)and(132.5,-116.6)..(133.3,-116.9)
				..controls(135.7,-117.9)and(136.1,-117.9)..(142.5,-117.8)
				..controls(145.8,-117.8)and(148.4,-117.7)..(148.8,-117.6)
				..controls(149.2,-117.6)and(150.2,-117.2)..(151,-116.8)
				..controls(152,-116.4)and(152.8,-116)..(153.6,-115.9)
				..controls(154.5,-115.7)and(155.2,-115.4)..(156.3,-114.9)
				..controls(157.1,-114.5)and(158.2,-114)..(158.7,-113.8)
				..controls(159.1,-113.7)and(159.8,-113.3)..(160.2,-113)
				..controls(160.6,-112.7)and(161.3,-112.4)..(161.7,-112.2)
				..controls(162.7,-111.9)and(163.2,-111.7)..(164.1,-111.1)
				..controls(164.5,-110.8)and(165.3,-110.4)..(165.9,-110)
				..controls(166.5,-109.7)and(167.3,-109.2)..(167.6,-108.9)
				..controls(167.9,-108.6)and(168.6,-108.2)..(169,-108)
				..controls(169.4,-107.8)and(170.2,-107.3)..(170.7,-106.9)
				..controls(171.2,-106.5)and(172.1,-105.8)..(172.7,-105.4)
				..controls(173.2,-105)and(174.1,-104.4)..(174.5,-104)
				..controls(174.9,-103.7)and(175.7,-103.1)..(176.2,-102.6)
				..controls(176.7,-102.2)and(177.4,-101.5)..(177.9,-101.2)
				..controls(178.3,-100.8)and(178.8,-100.4)..(179,-100.2)
				..controls(179.2,-100)and(179.6,-99.6)..(180,-99.3)
				..controls(180.4,-99.1)and(180.9,-98.6)..(181.2,-98.4)
				..controls(181.5,-98.1)and(182.1,-97.7)..(182.5,-97.5)
				..controls(183.2,-97.2)and(183.5,-96.9)..(184.3,-95.9)
				..controls(184.6,-95.6)and(184.6,-95.6)..(184.8,-95.7)
				..controls(185,-95.8)and(185.1,-96)..(185.1,-96.2)
				..controls(185.1,-96.5)and(185,-96.7)..(184.6,-97.2)
				..controls(184.1,-97.8)and(184,-97.9)..(183.9,-98.8)
				..controls(183.7,-100.3)and(183.3,-101.3)..(182.9,-101.8)
				..controls(182.6,-102.1)and(182.3,-102.6)..(182.2,-102.9)
				..controls(182,-103.3)and(181.6,-104)..(181.2,-104.6)
				..controls(180.9,-105.2)and(180.5,-106.1)..(180.4,-106.6)
				..controls(179.9,-108)and(179.7,-108.5)..(179,-109.4)
				..controls(178.6,-109.8)and(178.2,-110.5)..(177.9,-111)
				..controls(177.7,-111.5)and(177.3,-112.3)..(177,-112.7)
				..controls(176.6,-113.2)and(176.4,-113.7)..(176.3,-114.2)
				..controls(176.2,-114.7)and(175.9,-115.2)..(175.4,-116)
				..controls(175,-116.6)and(174.4,-117.7)..(174,-118.5)
				..controls(173.7,-119.2)and(173.3,-120.1)..(173,-120.5)
				..controls(172.8,-120.9)and(172.4,-121.8)..(172.1,-122.5)
				..controls(171.8,-123.2)and(171.3,-124.1)..(171,-124.6)
				..controls(170.8,-125.1)and(170.3,-126)..(169.9,-126.6)
				..controls(169.6,-127.3)and(169.2,-128.1)..(169,-128.3)
				..controls(168.8,-128.6)and(168.5,-129.3)..(168.3,-129.8)
				..controls(168.1,-130.3)and(167.8,-131.1)..(167.5,-131.6)
				..controls(166.9,-132.5)and(164.2,-137.9)..(164.2,-138.1)
				..controls(164.2,-138.1)and(163.9,-138.8)..(163.4,-139.7)
				..controls(162.9,-140.5)and(162.4,-141.5)..(162.3,-141.9)
				..controls(161.9,-143)and(161.6,-143.4)..(160.9,-144.1)
				..controls(160.4,-144.6)and(160.3,-144.8)..(160.1,-145.4)
				..controls(160,-145.8)and(159.6,-146.8)..(159.2,-147.8)
				..controls(158.7,-148.8)and(158.1,-150)..(157.9,-150.6)
				..controls(157.6,-151.1)and(157.2,-151.9)..(156.9,-152.4)
				..controls(156.5,-152.8)and(156.1,-153.7)..(155.8,-154.4)
				..controls(155.6,-155)and(155.2,-155.9)..(154.9,-156.3)
				..controls(154.7,-156.7)and(154.4,-157.4)..(154.2,-157.9)
				..controls(154.1,-158.4)and(153.8,-159)..(153.5,-159.3)
				..controls(152.8,-160.1)and(153,-160.1)..(142.4,-160.1)
				..controls(137.2,-160.1)and(131.8,-160.1)..(130.4,-160)--cycle
				(141.5,-103.7)..controls(140.8,-103)and(139.8,-102.1)..(139.2,-101.6)
				..controls(138.5,-101.2)and(137.9,-100.6)..(137.8,-100.4)
				..controls(137,-99.5)and(134,-96.5)..(133.2,-96)
				..controls(132.7,-95.7)and(131.8,-94.9)..(131.1,-94.3)
				..controls(129.4,-92.7)and(127.9,-91.6)..(126.6,-91.1)
				..controls(125.2,-90.5)and(124.1,-90)..(122.9,-89.1)
				..controls(122.1,-88.6)and(121.3,-88.2)..(119.9,-87.7)
				..controls(118.8,-87.3)and(117.5,-86.8)..(116.9,-86.5)
				..controls(114.3,-85.5)and(113.7,-85.4)..(109.7,-85.5)
				..controls(105.3,-85.7)and(103.9,-85.9)..(101.6,-87)
				..controls(100.7,-87.4)and(99.4,-87.8)..(97.9,-88)
				..controls(97,-88.2)and(96.6,-88.3)..(95.8,-88.9)
				..controls(95.4,-89.2)and(94.7,-89.6)..(94.2,-89.8)
				..controls(92.9,-90.3)and(92.3,-90.7)..(90.9,-92.1)
				..controls(90.2,-92.7)and(89.4,-93.3)..(89,-93.5)..controls(87,-94.5)and(85.2,-96)..(83.6,-97.9)
				..controls(83.3,-98.2)and(82.5,-98.9)..(81.8,-99.4)
				..controls(80.7,-100.1)and(80.5,-100.3)..(80,-101.1)
				..controls(78.7,-103.1)and(77.2,-103.9)..(76.4,-103)
				..controls(75.7,-102.2)and(75.7,-103.2)..(75.7,-83.6)
				..controls(75.6,-73.6)and(75.6,-64.8)..(75.5,-64)
				..controls(75.5,-63.2)and(75.5,-61.4)..(75.6,-59.8)--(75.7,-57.1)--(76,-56.8)
				..controls(76.8,-56.2)and(78,-56.3)..(79.8,-57.2)
				..controls(80.4,-57.5)and(81.3,-57.8)..(81.8,-57.9)
				..controls(82.3,-58.1)and(83,-58.4)..(83.6,-58.8)..controls(84.3,-59.2)and(85,-59.5)..(86,-59.8)
				..controls(86.8,-60.1)and(88,-60.6)..(88.7,-60.9)
				..controls(89.3,-61.2)and(90.6,-61.7)..(91.4,-62)
				..controls(92.2,-62.3)and(93.7,-62.9)..(94.6,-63.3)
				..controls(95.6,-63.7)and(96.8,-64.1)..(97.4,-64.3)
				..controls(98,-64.5)and(99,-64.9)..(99.7,-65.2)
				..controls(100.4,-65.6)and(101.5,-66)..(102.3,-66.2)
				..controls(104.2,-66.8)and(105,-67.1)..(106.2,-67.8)
				..controls(107.9,-68.7)and(110.2,-69.4)..(111.3,-69.2)
				..controls(112.5,-69)and(115,-68.2)..(115.6,-67.8)
				..controls(117.3,-66.7)and(117.6,-66.6)..(119.5,-66.2)
				..controls(121.9,-65.8)and(122.1,-65.7)..(123,-65.1)
				..controls(124.1,-64.4)and(124.4,-64.3)..(126,-63.9)
				..controls(127,-63.7)and(127.7,-63.4)..(128.4,-63.1)
				..controls(129,-62.8)and(129.7,-62.5)..(130.6,-62.4)
				..controls(132.1,-62.1)and(134,-61.5)..(135.2,-61)
				..controls(135.7,-60.8)and(136.7,-60.5)..(137.3,-60.3)
				..controls(138.4,-60)and(139.3,-59.7)..(141.1,-59)
				..controls(141.6,-58.7)and(142,-58.6)..(142.5,-58.6)
				..controls(143.2,-58.6)and(143.3,-58.7)..(143.6,-59)
				..controls(144,-59.5)and(144.1,-59.9)..(144,-60.8)
				..controls(144,-61.2)and(144,-71.1)..(144,-82.8)--(143.9,-104.1)--(143.7,-104.5)
				..controls(143.5,-104.7)and(143.3,-104.8)..(143,-104.9)
				..controls(142.6,-104.9)and(142.5,-104.8)..(141.5,-103.7)--cycle
				(27.8,-91.4)..controls(27.1,-90.9)and(26.3,-90)..(26.1,-89.5)
				..controls(26,-89.2)and(25.8,-88.6)..(25.7,-87.9)
				..controls(25.4,-86.2)and(25.6,-84.8)..(26.5,-83.5)
				..controls(27.4,-82.3)and(28.3,-82)..(29.4,-82.3)
				..controls(31.1,-82.8)and(32.1,-84.9)..(32,-87.4)
				..controls(31.9,-89.3)and(31.2,-90.6)..(29.9,-91.4)
				..controls(29.1,-91.9)and(28.6,-91.9)..(27.8,-91.4)--cycle
				(190.1,-91.3)..controls(189.7,-91.2)and(189.3,-90.9)..(189.1,-90.7)
				..controls(188.3,-89.8)and(187.6,-87.7)..(187.6,-86.2)
				..controls(187.6,-85.5)and(187.7,-85.2)..(188.2,-84.4)
				..controls(189.1,-82.6)and(190.1,-81.8)..(191.3,-82.1)
				..controls(192,-82.3)and(193.5,-83.7)..(193.8,-84.4)
				..controls(194.1,-85.1)and(194.2,-87.1)..(194,-88.3)
				..controls(193.8,-89.8)and(193.3,-90.5)..(192.3,-91.3)
				..controls(191.7,-91.7)and(190.9,-91.7)..(190.1,-91.3)--cycle
				(173.7,-83.4)..controls(172.2,-82.5)and(172,-82.5)..(170.6,-82.4)
				..controls(169.7,-82.4)and(168.8,-82.4)..(168.3,-82.5)
				..controls(167.2,-82.7)and(166.9,-82.7)..(166.6,-82.2)
				..controls(166.3,-81.8)and(166.3,-81)..(166.8,-80)
				..controls(167.2,-78.9)and(167.4,-77.4)..(167.7,-74.1)
				..controls(167.8,-72.8)and(167.9,-71.8)..(167.9,-71.7)
				..controls(168,-71.7)and(168,-70.2)..(168,-68.4)..controls(168,-65.4)and(168,-64.9)..(167.7,-63)
				..controls(167.3,-60.5)and(167.1,-59.9)..(166.4,-58.4)
				..controls(166.1,-57.7)and(165.8,-56.9)..(165.7,-56.6)
				..controls(165.7,-56.2)and(165.4,-55.6)..(165.2,-55.2)
				..controls(164.6,-53.8)and(161.9,-50.7)..(161,-50.3)
				..controls(159.8,-49.7)and(157.9,-50)..(155.2,-51.2)
				..controls(154.7,-51.4)and(153.4,-51.8)..(152.2,-52.1)
				..controls(150.8,-52.4)and(149.4,-52.9)..(148.6,-53.2)
				..controls(147.8,-53.6)and(146.7,-54)..(146,-54.2)
				..controls(145.4,-54.4)and(144.5,-54.7)..(144,-54.9)
				..controls(143.6,-55.2)and(142.7,-55.5)..(141.8,-55.7)
				..controls(141,-55.9)and(139.6,-56.3)..(138.7,-56.6)
				..controls(137.8,-57)and(136.3,-57.4)..(135.3,-57.6)
				..controls(133.8,-57.9)and(133.2,-58.1)..(132.1,-58.7)
				..controls(131.1,-59.2)and(130.3,-59.4)..(129.4,-59.6)
				..controls(128.7,-59.8)and(127.8,-60)..(127.4,-60.1)
				..controls(127,-60.3)and(126,-60.5)..(125.1,-60.7)
				..controls(123.2,-61.1)and(121.7,-61.5)..(120.2,-62.1)
				..controls(119.7,-62.3)and(118.1,-62.8)..(116.7,-63.2)
				..controls(114.6,-63.9)and(114,-64.1)..(113,-64.7)
				..controls(111.9,-65.4)and(111.8,-65.4)..(111,-65.5)
				..controls(109.9,-65.5)and(109.3,-65.4)..(107.7,-64.7)
				..controls(107,-64.4)and(106,-64)..(105.3,-63.9)
				..controls(103.5,-63.5)and(102.2,-63)..(101.2,-62.5)
				..controls(100.1,-61.8)and(99.3,-61.5)..(97.8,-61.1)
				..controls(97.2,-61)and(96.2,-60.7)..(95.7,-60.4)
				..controls(95.1,-60.1)and(94.2,-59.8)..(93.5,-59.6)
				..controls(92.9,-59.4)and(91.7,-59)..(91,-58.7)..controls(90.2,-58.4)and(88.9,-58)..(88,-57.7)
				..controls(87.2,-57.5)and(86.1,-57.1)..(85.5,-56.8)
				..controls(85,-56.6)and(84.1,-56.2)..(83.5,-56)..controls(82.9,-55.7)and(81.8,-55.2)..(81,-54.9)
				..controls(80.1,-54.5)and(79.1,-54.2)..(78.7,-54.1)
				..controls(77.9,-53.8)and(77.3,-53.6)..(76.2,-52.9)
				..controls(75.7,-52.6)and(75.2,-52.4)..(74.1,-52.2)
				..controls(72.5,-51.9)and(72.4,-51.8)..(71.3,-51.1)
				..controls(69.9,-50.1)and(69.4,-49.9)..(67.9,-49.8)
				..controls(67,-49.7)and(66.6,-49.6)..(66.3,-49.4)
				..controls(66.1,-49.3)and(65.7,-49)..(65.4,-48.8)
				..controls(65.1,-48.7)and(64.2,-48.4)..(63.6,-48.2)
				..controls(62.9,-48)and(61.7,-47.6)..(60.9,-47.2)
				..controls(60.1,-46.8)and(58.9,-46.4)..(58.2,-46.2)
				..controls(57.6,-46)and(56.8,-45.7)..(56.5,-45.5)
				..controls(54.9,-44.6)and(54.6,-44.5)..(53.2,-44.2)
				..controls(51.7,-43.8)and(51,-43.6)..(49.9,-42.9)
				..controls(49.5,-42.6)and(48.9,-42.4)..(47.5,-42)
				..controls(45.7,-41.6)and(45.6,-41.5)..(44.3,-40.6)
				..controls(43.5,-40.1)and(42.8,-39.6)..(42.7,-39.4)
				..controls(42.4,-38.7)and(43,-38.1)..(44.5,-37.6)..controls(44.9,-37.5)and(46,-37)..(46.9,-36.6)
				..controls(47.9,-36.1)and(48.8,-35.8)..(49.6,-35.6)
				..controls(50.5,-35.4)and(51.3,-35.2)..(52.4,-34.6)
				..controls(53.3,-34.1)and(54.5,-33.7)..(54.9,-33.5)
				..controls(55.4,-33.4)and(56.4,-33)..(57.2,-32.6)
				..controls(58.6,-32)and(60.1,-31.5)..(61.6,-31.2)
				..controls(62.3,-31.1)and(62.7,-31)..(63.2,-30.7)
				..controls(64.1,-30)and(64.3,-29.9)..(65.3,-29.7)
				..controls(65.7,-29.6)and(66.5,-29.2)..(67.1,-28.9)
				..controls(67.6,-28.6)and(68.8,-28.2)..(69.6,-27.9)
				..controls(70.5,-27.7)and(71.6,-27.3)..(72.1,-27)
				..controls(72.6,-26.7)and(73.8,-26.2)..(74.8,-25.9)
				..controls(75.7,-25.5)and(76.7,-25.1)..(77.1,-25)
				..controls(77.4,-24.8)and(78.5,-24.4)..(79.6,-24.1)
				..controls(80.6,-23.8)and(82,-23.2)..(82.8,-22.9)
				..controls(83.5,-22.5)and(84.8,-22)..(85.8,-21.8)
				..controls(86.7,-21.5)and(87.8,-21.1)..(88.1,-20.9)
				..controls(88.4,-20.7)and(89.3,-20.3)..(90.1,-20.1)
				..controls(90.8,-19.8)and(91.9,-19.4)..(92.4,-19.1)
				..controls(93.5,-18.5)and(94.4,-18.2)..(95.4,-18.1)
				..controls(95.7,-18)and(96.4,-17.8)..(96.9,-17.5)
				..controls(97.5,-17.3)and(98.4,-16.9)..(99,-16.6)
				..controls(99.7,-16.3)and(100.9,-15.9)..(101.9,-15.6)
				..controls(102.9,-15.3)and(104.4,-14.7)..(105.3,-14.3)
				..controls(108.1,-13.2)and(108.5,-13.1)..(110.1,-13.1)
				..controls(111.4,-13.1)and(111.5,-13.1)..(112.4,-13.5)
				..controls(113.2,-14)and(115,-14.8)..(116.6,-15.4)
				..controls(116.9,-15.5)and(117.9,-15.8)..(118.8,-16.1)
				..controls(120.9,-16.6)and(121.5,-16.9)..(122.6,-17.9)
				..controls(123,-18.3)and(123.7,-18.8)..(123.9,-18.9)
				..controls(124.6,-19.4)and(126,-19.8)..(127.9,-20.2)
				..controls(129.6,-20.5)and(129.8,-20.6)..(131.1,-21.4)
				..controls(131.5,-21.7)and(132.4,-22.1)..(133.1,-22.3)
				..controls(136.1,-23.2)and(138.3,-24.1)..(140.2,-25)
				..controls(141,-25.5)and(141.5,-25.6)..(142.5,-25.8)
				..controls(145.1,-26.3)and(146,-26.6)..(149.4,-28.1)
				..controls(151.3,-28.9)and(153.5,-29.8)..(154.3,-30.1)
				..controls(155.1,-30.4)and(156.2,-30.8)..(156.6,-31)
				..controls(157.1,-31.2)and(158.2,-31.6)..(159.1,-31.8)
				..controls(160.6,-32.2)and(160.9,-32.3)..(161.7,-32.9)
				..controls(162.4,-33.3)and(162.9,-33.6)..(163.7,-33.8)
				..controls(164.3,-34)and(164.8,-34.1)..(164.9,-34.1)
				..controls(165,-34.1)and(165.6,-34.4)..(166.4,-34.8)
				..controls(167.3,-35.2)and(168.1,-35.4)..(169,-35.6)
				..controls(170,-35.8)and(170.5,-36)..(171.7,-36.6)
				..controls(172.4,-37)and(173.3,-37.4)..(173.6,-37.4)
				..controls(175.6,-37.7)and(176.2,-38)..(176.7,-39.1)
				..controls(176.9,-39.6)and(176.9,-39.7)..(176.7,-40.4)
				..controls(176.4,-41.6)and(175.8,-42.7)..(175.2,-43.4)
				..controls(174.5,-44.1)and(172.6,-45.5)..(171.7,-45.8)
				..controls(170.6,-46.2)and(169.1,-46)..(168.2,-45.4)
				..controls(167.9,-45.2)and(167.3,-44.8)..(166.9,-44.4)
				..controls(166.4,-44.1)and(165.8,-43.6)..(165.4,-43.2)
				..controls(165.1,-42.9)and(164.4,-42.3)..(163.9,-41.9)
				..controls(163.3,-41.5)and(162.7,-40.9)..(162.4,-40.5)
				..controls(161.4,-39.2)and(159.9,-38)..(158.8,-37.8)
				..controls(158.4,-37.8)and(158.3,-37.7)..(158.2,-37.2)
				..controls(158,-36.8)and(157.4,-36.1)..(157.1,-36.1)
				..controls(157,-36.1)and(156.7,-36.1)..(156.5,-36)
				..controls(156.1,-35.9)and(155.9,-35.7)..(155.7,-35.2)
				..controls(155.3,-34.6)and(155.1,-34.4)..(154,-34.1)
				..controls(153.6,-34)and(153.2,-33.7)..(152.9,-33.5)
				..controls(152.4,-32.9)and(151.6,-32.4)..(151.3,-32.4)
				..controls(150.9,-32.4)and(150.4,-32.6)..(150.4,-32.9)
				..controls(150.4,-33.2)and(150.8,-33.4)..(151.4,-33.6)
				..controls(152.1,-33.7)and(152.8,-34.1)..(153.5,-34.9)
				..controls(154.1,-35.5)and(154.3,-35.6)..(155,-35.8)
				..controls(155.7,-36)and(155.7,-36)..(156,-36.6)
				..controls(156.3,-37.3)and(156.7,-37.7)..(157.5,-37.8)
				..controls(158.1,-37.9)and(158.2,-38)..(158.4,-38.8)
				..controls(158.5,-39.4)and(158.6,-39.6)..(159.6,-40.4)
				..controls(160.2,-41)and(160.9,-41.6)..(161.3,-41.8)
				..controls(161.7,-42)and(161.9,-42.3)..(162.2,-42.8)
				..controls(162.4,-43.2)and(162.8,-43.8)..(163.1,-44.1)
				..controls(163.9,-44.8)and(165.1,-46.7)..(165.7,-48)
				..controls(166.2,-49.3)and(166.4,-49.7)..(167.2,-50.5)
				..controls(167.5,-50.8)and(167.9,-51.4)..(168.1,-51.8)
				..controls(168.3,-52.2)and(168.8,-53)..(169.1,-53.6)
				..controls(169.6,-54.3)and(169.8,-54.9)..(169.9,-55.6)
				..controls(170.4,-57.8)and(170.5,-58.1)..(171.2,-59.5)
				..controls(171.6,-60.3)and(172,-61.5)..(172.2,-62.1)
				..controls(172.3,-62.8)and(172.7,-63.7)..(172.9,-64.1)
				..controls(173.2,-64.6)and(173.4,-65.3)..(173.4,-65.6)
				..controls(173.8,-68.8)and(174.2,-70.6)..(174.9,-72.6)
				..controls(175.5,-74.5)and(175.9,-77.1)..(176.1,-80.2)
				..controls(176.3,-83.4)and(175.5,-84.4)..(173.7,-83.4)--cycle;
				\definecolor{ezs}{RGB}{247,179,26}\fill[ezs](-5.7,-156.6)--(-3.4,-157.5)--(61.9,12.7)--(59.7,13.6)--cycle;
				\definecolor{ezs}{RGB}{247,179,22}\fill[ezs](-3.8,-157.3)--(-1.5,-158.2)--(63.9,12)--(61.6,12.9)--cycle;
				\definecolor{ezs}{RGB}{248,179,19}\fill[ezs](-1.8,-158.1)--(0.4,-158.9)--(65.8,11.2)--(63.5,12.1)--cycle;
				\definecolor{ezs}{RGB}{248,178,16}\fill[ezs](0.1,-158.8)--(2.4,-159.7)--(67.7,10.5)--(65.4,11.4)--cycle;
				\definecolor{ezs}{RGB}{249,178,15}\fill[ezs](2,-159.5)--(4.3,-160.4)--(69.6,9.8)--(67.3,10.7)--cycle;
				\definecolor{ezs}{RGB}{251,177,17}\fill[ezs](3.9,-160.3)--(6.2,-161.1)--(71.5,9)--(69.2,9.9)--cycle;
				\definecolor{ezs}{RGB}{252,179,18}\fill[ezs](5.8,-161)--(8.1,-161.9)--(73.4,8.3)--(71.1,9.2)--cycle;
				\definecolor{ezs}{RGB}{252,182,17}\fill[ezs](7.7,-161.7)--(10,-162.6)--(75.3,7.6)--(73,8.5)--cycle;
				\definecolor{ezs}{RGB}{252,185,17}\fill[ezs](9.6,-162.5)--(11.9,-163.4)--(77.2,6.8)--(75,7.7)--cycle;
				\definecolor{ezs}{RGB}{251,189,18}\fill[ezs](11.5,-163.2)--(13.8,-164.1)--(79.2,6.1)--(76.9,7)--cycle;
				\definecolor{ezs}{RGB}{251,191,22}\fill[ezs](13.4,-163.9)--(15.7,-164.8)--(81.1,5.4)--(78.8,6.2)--cycle;
				\definecolor{ezs}{RGB}{252,193,25}\fill[ezs](15.4,-164.7)--(17.6,-165.6)--(83,4.6)--(80.7,5.5)--cycle;
				\definecolor{ezs}{RGB}{253,195,27}\fill[ezs](17.3,-165.4)--(19.6,-166.3)--(84.9,3.9)--(82.6,4.8)--cycle;
				\definecolor{ezs}{RGB}{253,196,31}\fill[ezs](19.2,-166.1)--(21.5,-167)--(86.8,3.2)--(84.5,4)--cycle;
				\definecolor{ezs}{RGB}{254,199,35}\fill[ezs](21.1,-166.9)--(23.4,-167.8)--(88.7,2.4)--(86.4,3.3)--cycle;
				\definecolor{ezs}{RGB}{253,201,38}\fill[ezs](23,-167.6)--(25.3,-168.5)--(90.6,1.7)--(88.3,2.6)--cycle;
				\definecolor{ezs}{RGB}{253,203,40}\fill[ezs](24.9,-168.3)--(27.2,-169.2)--(92.5,1)--(90.3,1.8)--cycle;
				\definecolor{ezs}{RGB}{253,205,42}\fill[ezs](26.8,-169.1)--(29.1,-170)--(94.4,0.2)--(92.2,1.1)--cycle;
				\definecolor{ezs}{RGB}{253,207,45}\fill[ezs](28.7,-169.8)--(31,-170.7)--(96.4,-0.5)--(94.1,0.4)--cycle;
				\definecolor{ezs}{RGB}{253,210,49}\fill[ezs](30.7,-170.5)--(32.9,-171.4)--(98.3,-1.2)--(96,-0.4)--cycle;
				\definecolor{ezs}{RGB}{253,212,54}\fill[ezs](32.6,-171.3)--(34.9,-172.2)--(100.2,-2)--(97.9,-1.1)--cycle;
				\definecolor{ezs}{RGB}{253,214,59}\fill[ezs](34.5,-172)--(36.8,-172.9)--(102.1,-2.7)--(99.8,-1.8)--cycle;
				\definecolor{ezs}{RGB}{253,216,63}\fill[ezs](36.4,-172.7)--(38.7,-173.6)--(104,-3.4)--(101.7,-2.6)--cycle;
				\definecolor{ezs}{RGB}{253,218,65}\fill[ezs](38.3,-173.5)--(40.6,-174.4)--(105.9,-4.2)--(103.6,-3.3)--cycle;
				\definecolor{ezs}{RGB}{254,220,66}\fill[ezs](40.2,-174.2)--(42.5,-175.1)--(107.8,-4.9)--(105.5,-4)--cycle;
				\definecolor{ezs}{RGB}{254,222,68}\fill[ezs](42.1,-175)--(44.4,-175.8)--(109.7,-5.6)--(107.5,-4.8)--cycle;
				\definecolor{ezs}{RGB}{254,224,71}\fill[ezs](44,-175.7)--(46.3,-176.6)--(111.7,-6.4)--(109.4,-5.5)--cycle;
				\definecolor{ezs}{RGB}{254,226,75}\fill[ezs](46,-176.4)--(48.2,-177.3)--(113.6,-7.1)--(111.3,-6.2)--cycle;
				\definecolor{ezs}{RGB}{254,228,79}\fill[ezs](47.9,-177.2)--(50.1,-178)--(115.5,-7.8)--(113.2,-7)--cycle;
				\definecolor{ezs}{RGB}{254,230,82}\fill[ezs](49.8,-177.9)--(52.1,-178.8)--(117.4,-8.6)--(115.1,-7.7)--cycle;
				\definecolor{ezs}{RGB}{254,232,85}\fill[ezs](51.7,-178.6)--(54,-179.5)--(119.3,-9.3)--(117,-8.4)--cycle;
				\definecolor{ezs}{RGB}{254,234,88}\fill[ezs](53.6,-179.4)--(55.9,-180.2)--(121.2,-10)--(118.9,-9.2)--cycle;
				\definecolor{ezs}{RGB}{254,237,91}\fill[ezs](55.5,-180.1)--(57.8,-181)--(123.1,-10.8)--(120.8,-9.9)--cycle;
				\definecolor{ezs}{RGB}{254,239,94}\fill[ezs](57.4,-180.8)--(59.7,-181.7)--(125,-11.5)--(122.8,-10.6)--cycle;
				\definecolor{ezs}{RGB}{254,241,97}\fill[ezs](59.3,-181.6)--(61.6,-182.4)--(127,-12.2)--(124.7,-11.4)--cycle;
				\definecolor{ezs}{RGB}{254,243,101}\fill[ezs](61.2,-182.3)--(63.5,-183.2)--(128.9,-13)--(126.6,-12.1)--cycle;
				\definecolor{ezs}{RGB}{254,245,103}\fill[ezs](63.2,-183)--(65.4,-183.9)--(130.8,-13.7)--(128.5,-12.8)--cycle;
				\definecolor{ezs}{RGB}{254,245,105}\fill[ezs](65.1,-183.8)--(67.4,-184.6)--(132.7,-14.4)--(130.4,-13.6)--cycle;
				\definecolor{ezs}{RGB}{254,245,105}\fill[ezs](67,-184.5)--(69.3,-185.4)--(134.6,-15.2)--(132.3,-14.3)--cycle;
				\definecolor{ezs}{RGB}{254,244,103}\fill[ezs](68.9,-185.2)--(71.2,-186.1)--(136.5,-15.9)--(134.2,-15)--cycle;
				\definecolor{ezs}{RGB}{253,243,101}\fill[ezs](70.8,-186)--(73.1,-186.8)--(138.4,-16.6)--(136.1,-15.8)--cycle;
				\definecolor{ezs}{RGB}{253,241,98}\fill[ezs](72.7,-186.7)--(75,-187.6)--(140.3,-17.4)--(138,-16.5)--cycle;
				\definecolor{ezs}{RGB}{253,239,95}\fill[ezs](74.6,-187.4)--(76.9,-188.3)--(142.2,-18.1)--(140,-17.2)--cycle;
				\definecolor{ezs}{RGB}{253,236,92}\fill[ezs](76.5,-188.2)--(78.8,-189)--(144.2,-18.8)--(141.9,-18)--cycle;
				\definecolor{ezs}{RGB}{253,234,88}\fill[ezs](78.5,-188.9)--(80.7,-189.8)--(146.1,-19.6)--(143.8,-18.7)--cycle;
				\definecolor{ezs}{RGB}{253,232,85}\fill[ezs](80.4,-189.6)--(82.7,-190.5)--(148,-20.3)--(145.7,-19.4)--cycle;
				\definecolor{ezs}{RGB}{253,230,82}\fill[ezs](82.3,-190.4)--(84.6,-191.2)--(149.9,-21)--(147.6,-20.2)--cycle;
				\definecolor{ezs}{RGB}{253,228,78}\fill[ezs](84.2,-191.1)--(86.5,-192)--(151.8,-21.8)--(149.5,-20.9)--cycle;
				\definecolor{ezs}{RGB}{253,226,74}\fill[ezs](86.1,-191.8)--(88.4,-192.7)--(153.7,-22.5)--(151.4,-21.6)--cycle;
				\definecolor{ezs}{RGB}{253,224,70}\fill[ezs](88,-192.6)--(90.3,-193.4)--(155.6,-23.3)--(153.3,-22.4)--cycle;
				\definecolor{ezs}{RGB}{253,222,67}\fill[ezs](89.9,-193.3)--(92.2,-194.2)--(157.5,-24)--(155.3,-23.1)--cycle;
				\definecolor{ezs}{RGB}{253,220,65}\fill[ezs](91.8,-194)--(94.1,-194.9)--(159.5,-24.7)--(157.2,-23.8)--cycle;
				\definecolor{ezs}{RGB}{253,217,63}\fill[ezs](93.7,-194.8)--(96,-195.6)--(161.4,-25.5)--(159.1,-24.6)--cycle;
				\definecolor{ezs}{RGB}{253,215,60}\fill[ezs](95.7,-195.5)--(97.9,-196.4)--(163.3,-26.2)--(161,-25.3)--cycle;
				\definecolor{ezs}{RGB}{253,213,56}\fill[ezs](97.6,-196.2)--(99.9,-197.1)--(165.2,-26.9)--(162.9,-26)--cycle;
				\definecolor{ezs}{RGB}{252,211,53}\fill[ezs](99.5,-197)--(101.8,-197.8)--(167.1,-27.7)--(164.8,-26.8)--cycle;
				\definecolor{ezs}{RGB}{252,209,49}\fill[ezs](101.4,-197.7)--(103.7,-198.6)--(169,-28.4)--(166.7,-27.5)--cycle;
				\definecolor{ezs}{RGB}{252,207,45}\fill[ezs](103.3,-198.4)--(105.6,-199.3)--(170.9,-29.1)--(168.6,-28.2)--cycle;
				\definecolor{ezs}{RGB}{253,205,42}\fill[ezs](105.2,-199.2)--(107.5,-200)--(172.8,-29.9)--(170.6,-29)--cycle;
				\definecolor{ezs}{RGB}{253,202,40}\fill[ezs](107.1,-199.9)--(109.4,-200.8)--(174.7,-30.6)--(172.5,-29.7)--cycle;
				\definecolor{ezs}{RGB}{253,200,38}\fill[ezs](109,-200.6)--(111.3,-201.5)--(176.7,-31.3)--(174.4,-30.4)--cycle;
				\definecolor{ezs}{RGB}{253,198,35}\fill[ezs](111,-201.4)--(113.2,-202.2)--(178.6,-32.1)--(176.3,-31.2)--cycle;
				\definecolor{ezs}{RGB}{253,196,31}\fill[ezs](112.9,-202.1)--(115.2,-203)--(180.5,-32.8)--(178.2,-31.9)--cycle;
				\definecolor{ezs}{RGB}{253,194,28}\fill[ezs](114.8,-202.8)--(117.1,-203.7)--(182.4,-33.5)--(180.1,-32.6)--cycle;
				\definecolor{ezs}{RGB}{253,192,26}\fill[ezs](116.7,-203.6)--(119,-204.5)--(184.3,-34.3)--(182,-33.4)--cycle;
				\definecolor{ezs}{RGB}{253,189,24}\fill[ezs](118.6,-204.3)--(120.9,-205.2)--(186.2,-35)--(183.9,-34.1)--cycle;
				\definecolor{ezs}{RGB}{253,186,20}\fill[ezs](120.5,-205)--(122.8,-205.9)--(188.1,-35.7)--(185.8,-34.8)--cycle;
				\definecolor{ezs}{RGB}{253,184,17}\fill[ezs](122.4,-205.8)--(124.7,-206.7)--(190,-36.5)--(187.8,-35.6)--cycle;
				\definecolor{ezs}{RGB}{252,183,15}\fill[ezs](124.3,-206.5)--(126.6,-207.4)--(192,-37.2)--(189.7,-36.3)--cycle;
				\definecolor{ezs}{RGB}{252,181,14}\fill[ezs](126.3,-207.2)--(128.5,-208.1)--(193.9,-37.9)--(191.6,-37.1)--cycle;
				\definecolor{ezs}{RGB}{252,179,13}\fill[ezs](128.2,-208)--(130.4,-208.9)--(195.8,-38.7)--(193.5,-37.8)--cycle;
				\definecolor{ezs}{RGB}{252,178,12}\fill[ezs](130.1,-208.7)--(132.4,-209.6)--(197.7,-39.4)--(195.4,-38.5)--cycle;
				\definecolor{ezs}{RGB}{251,178,10}\fill[ezs](132,-209.4)--(134.3,-210.3)--(199.6,-40.1)--(197.3,-39.3)--cycle;
				\definecolor{ezs}{RGB}{251,178,10}\fill[ezs](133.9,-210.2)--(136.2,-211.1)--(201.5,-40.9)--(199.2,-40)--cycle;
				\definecolor{ezs}{RGB}{250,177,11}\fill[ezs](135.8,-210.9)--(138.1,-211.8)--(203.4,-41.6)--(201.1,-40.7)--cycle;
				\definecolor{ezs}{RGB}{250,178,11}\fill[ezs](137.7,-211.6)--(140,-212.5)--(205.3,-42.3)--(203.1,-41.5)--cycle;
				\definecolor{ezs}{RGB}{248,180,10}\fill[ezs](139.6,-212.4)--(141.9,-213.3)--(207.3,-43.1)--(205,-42.2)--cycle;
				\definecolor{ezs}{RGB}{247,181,9}\fill[ezs](141.5,-213.1)--(143.8,-214)--(209.2,-43.8)--(206.9,-42.9)--cycle;
				\definecolor{ezs}{RGB}{246,182,9}\fill[ezs](143.5,-213.8)--(145.7,-214.7)--(211.1,-44.5)--(208.8,-43.7)--cycle;
				\definecolor{ezs}{RGB}{246,182,10}\fill[ezs](145.4,-214.6)--(147.7,-215.5)--(213,-45.3)--(210.7,-44.4)--cycle;
			\end{scope}
			% chữ EDUZY (Times đậm, làm dày nét bằng 5 lớp lệch nhẹ; căn đúng khung chữ logo gốc)
			\foreach \dx/\dy in {-0.7/0,0.7/0,0/-0.3,0/0.3,0/0}
			\node[anchor=center,inner sep=0pt,outer sep=0pt,text=ezText,align=center] at (109.6+\dx,-193.6+\dy)
			{\resizebox{1.01cm}{0.22cm}{\fontencoding{T1}\fontfamily{ptm}\fontseries{b}\selectfont EDUZY}};
		\end{tikzpicture}%
	}%
}
%====================== THIẾT LẬP THÂN ĐỀ ======================
%--- Môi trường tự luận {bt}: ex_test 3.3.4 chưa khai báo sẵn nên tạo ở đây
%    (ex_test đã tự móc \fixshowans và \listenumerate cho bt)
\makeatletter
\@ifundefined{bt}{\newtheorem{bt}{\color{gold1}\namebt}}{}
\makeatother
\renewcommand{\sectionmark}[1]{\markright{\thesection\ #1}}
%--- Khung tiêu đề (logo + tối đa 4 dòng; dòng 3, 4 để trống {} thì bỏ qua)
%  Đề thi:    \khungtieude{THI THỬ THPT QUỐC GIA NĂM 2026}{Môn Toán -- Đề số 21}
%                         {(Đề thi gồm: ...)}{Thời gian làm bài: 90 phút (không kể thời gian phát đề).}
%  Chuyên đề: \khungtieude{CHUYÊN ĐỀ 3 -- DÃY SỐ VÀ CÔNG THỨC TRUY HỒI}
%                         {Bồi dưỡng học sinh giỏi Toán THPT}{...}{}
\newcommand{\khungtieude}[4]{%
	\begin{tcolorbox}[colframe=gold1,colback=gold5!45,boxrule=0.5mm]
		\begin{minipage}[c]{3cm}
			\centerline{\eduzylogo[2.4cm]}
		\end{minipage}\hspace{0.5cm}
		\begin{minipage}[c]{12.5cm}
			\vskip 0.15 cm
			\centerline{\ind{#1}}
			\vskip 0.15 cm
			\centerline{\textbf{#2}}
			\if\relax\detokenize{#3}\relax\else
				\vskip 0.15 cm
				\centerline{\textit{#3}}
			\fi
			\if\relax\detokenize{#4}\relax\else
				\vskip 0.15 cm
				\centerline{\textit{#4}}
			\fi
			\vskip 0.15 cm
		\end{minipage}
	\end{tcolorbox}
}
%====================== BỔ SUNG CHO BÀI HỌC, CHUYÊN ĐỀ ======================
\usepackage[shortlabels]{enumitem}
\usepackage{tabularx}
%--- Danh sách gọn: \begin{enumerate}[gon] ... \end{enumerate}
\SetEnumitemKey{gon}{leftmargin=*,itemsep=1pt,parsep=0pt,topsep=3pt}
%--- Khung lý thuyết: \begin{khung}{} ... \end{khung}; tiêu đề chỉ dùng khi thật ngắn ("Mẹo", "Lỗi thường gặp")
\newtcolorbox{khung}[1]{breakable,enhanced,colframe=gold1,colback=gold5!45,
	boxrule=0.4mm,fonttitle=\bfseries\sffamily,coltitle=white,colbacktitle=gold1,
	attach boxed title to top left={xshift=6mm,yshift=-3mm},
	boxed title style={arc=1mm,boxrule=0.3mm},title={#1}}
%--- Tiêu đề mục nhỏ trong một dạng: \muc{Dấu hiệu}, \muc{Phương pháp}, \muc{Bài tự luyện}
\makeatletter
\newcommand{\muc}[1]{\par\medskip\noindent\indam{#1}\par\nobreak\smallskip\@afterindentfalse\@afterheading}
\makeatother
%--- Ô công thức cần nhớ (căn giữa, viền vàng, vừa khít nội dung)
%     \begin{center}\congthuc{$Q(n)=(n-1)!$}\end{center}
\newtcbox{\congthuc}{on line,colframe=gold1,colback=gold5!45,boxrule=0.5pt,arc=2pt,
	left=8pt,right=8pt,top=4pt,bottom=4pt}
%--- Công tắc lời giải bài tự luyện (ví dụ mẫu luôn dùng \loigiai)
%     \loigiaiTL[{đáp số}]{lời giải đầy đủ}
%     \TLfull -> in lời giải đầy đủ;  \TLdapso -> chỉ in một dòng "Đáp số";  \TLan -> ẩn hoàn toàn
%     (\CheDo là dethi thì mọi lời giải và đáp số đều ẩn)
\makeatletter
\newif\ifTLhien   \TLhientrue
\newif\ifTLdapso  \TLdapsofalse
\newcommand{\TLfull}{\TLhientrue\TLdapsofalse}
\newcommand{\TLdapso}{\TLhientrue\TLdapsotrue}
\newcommand{\TLan}{\TLhienfalse}
\newcommand{\loigiaiTL}[2][]{%
	\ifTLhien
		\ifTLdapso
			\def\@tlarg{#1}%
			\ifx\@tlarg\@empty\else
				\ifdefstring{\kindOPTN}{dethi}{}{\par\noindent{\bfseries\color{gold1}Đáp số:} #1\par}%
			\fi
		\else
			\loigiai{#2}%
		\fi
	\fi}
\makeatother
\TLdapso% mặc định: chỉ in đáp số bài tự luyện
%====================== BẢNG ĐÚNG/SAI (\choiceTFt) ======================
%  Hàng tiêu đề nền nâu đồng chữ trắng (Phát biểu | Đ | S), các hàng tô xen kẽ kem / trắng,
%  đường kẻ vàng kim, nhãn a) b) c) d) màu nâu đồng. Bản có đáp án: cột được chọn có vòng tròn
%  nâu đồng với dấu tích trắng (câu đúng tích ở cột Đúng, câu sai tích ở cột Sai).
%  Tiêu đề ghi đủ "Đúng | Sai": \OPTN{viettat=0} sau \input (cột Đúng/Sai rộng ra, phát biểu hẹp lại).
%  Dùng bảng gốc của ex_test: \def\BangDSgoc{1} trước \input{preamble-eduzy-v2}.
\providecommand{\BangDSgoc}{0}
\ifnum\BangDSgoc=0
\usepackage{colortbl,pifont}
\colorlet{ezDSke}{gold2!60!gold1}% đường kẻ bảng (nét 0,4pt như ex_test để không đổi chiều cao bảng)
\colorlet{ezDSle}{gold5!60}% hàng a), c)
\colorlet{ezDSchan}{white}% hàng b), d)
%--- dấu chọn ở cột Đúng/Sai
%    (bề rộng 0, căn giữa ô: không làm đổi bề rộng cột của ex_test, không báo tràn lề)
\newcommand{\EZdauchon}{\makebox[0pt]{\tikz[baseline=-1ex]\node[circle,fill=gold1,inner sep=0pt,
	minimum size=1.2em,text=white,font=\small]{\ding{51}};}}
\def\tickT{\EZdauchon}\def\tickF{\EZdauchon}
\AtBeginDocument{\sbox{\ccTFTab}{\EZdauchon}}
%--- nhãn a) b) c) d) màu nâu đồng (ở bản lời giải, nhãn câu đúng vẫn đóng khung như ex_test)
\renewcommand{\FalseTF}{\stepcounter{dapan}\textbf{\color{gold1}\DapAnTF\sepTF}\ignorespaces}
%--- trước khi dựng bảng: kẻ màu vàng kim (bề rộng cột giữ nguyên như ex_test); sau bảng trả lại màu đen
\newcount\EZhangTF
\newcommand{\EZbatdauTF}{%
	\global\EZhangTF=0\relax
	\arrayrulecolor{ezDSke}}
\newcommand{\EZketthucTF}{\arrayrulecolor{black}}
%--- hàng tiêu đề nền nâu đồng, chữ trắng
\makeatletter
\newcommand{\EZo@TF}[1]{\centering\bfseries\color{white}#1}
\newcommand{\EZtieudeTF}{%
	\ifcase\cotDS
		\ifnum\cotSTT=1\relax
			\def\rowHeader{\hline\rowcolor{gold1}\EZo@TF{TT}&\EZo@TF{\phatbieu}\tabularnewline}%
		\else
			\def\rowHeader{\hline\rowcolor{gold1}\EZo@TF{\phatbieu}\tabularnewline}%
		\fi
	\or
		\ifnum\cotSTT=1\relax
			\def\rowHeader{\hline\rowcolor{gold1}\EZo@TF{TT}&\EZo@TF{\phatbieu}&\EZo@TF{\Dung}\tabularnewline}%
		\else
			\def\rowHeader{\hline\rowcolor{gold1}\EZo@TF{\phatbieu}&\EZo@TF{\Dung}\tabularnewline}%
		\fi
	\or
		\ifnum\cotSTT=1\relax
			\def\rowHeader{\hline\rowcolor{gold1}\EZo@TF{TT}&\EZo@TF{\phatbieu}&\EZo@TF{\Dung}&\EZo@TF{\Sai}\tabularnewline}%
		\else
			\def\rowHeader{\hline\rowcolor{gold1}\EZo@TF{\phatbieu}&\EZo@TF{\Dung}&\EZo@TF{\Sai}\tabularnewline}%
		\fi
	\fi}
%--- tô xen kẽ từng hàng (gọi ngay sau \hline đầu mỗi hàng)
\newcommand{\EZmauhangTF}{%
	\noalign{\global\advance\EZhangTF\@ne}%
	\ifodd\EZhangTF\rowcolor{ezDSle}\else\rowcolor{ezDSchan}\fi}
%--- gắn vào \choiceTFt / \writeChoiceTFt của ex_test
\xpretocmd{\choiceTFt}{\EZbatdauTF}{}{\typeout{[preamble-eduzy-v2] Không gắn được bảng Đúng/Sai (đầu bảng)}}
\xpatchcmd{\choiceTFt}{\ifnum\dongPB=0\def\rowHeader{}\fi}{\EZtieudeTF\ifnum\dongPB=0\def\rowHeader{}\fi}
	{}{\typeout{[preamble-eduzy-v2] Không gắn được bảng Đúng/Sai (hàng tiêu đề)}}
\xpatchcmd{\writeChoiceTFt}{\hline\xdef\dotEX{\sdotEX}}{\hline\EZmauhangTF\xdef\dotEX{\sdotEX}}
	{}{\typeout{[preamble-eduzy-v2] Không gắn được bảng Đúng/Sai (tô hàng)}}
\xpatchcmd{\writeChoiceTFt}{{\end{longtable}}}{{\end{longtable}\EZketthucTF}}
	{}{\typeout{[preamble-eduzy-v2] Không gắn được bảng Đúng/Sai (cuối longtable)}}
\xpatchcmd{\writeChoiceTFt}{{\end{xtabular}\end{center}\medskip\noindent}}{{\end{xtabular}\end{center}\EZketthucTF\medskip\noindent}}
	{}{\typeout{[preamble-eduzy-v2] Không gắn được bảng Đúng/Sai (cuối xtabular)}}
\fi
%=====================================================================
\renewcommand{\parallel}{\mathrel{/\!\!/}}

%      \TLdapso                      % tài liệu: \TLfull | \TLdapso | \TLan
%      \begin{document} ... \end{document}
%  Logo Eduzy vẽ bằng TikZ ngay trong file (\eduzylogo), không cần file ảnh.
%  Biên dịch: pdfLaTeX (máy có gói vietnam) hoặc XeLaTeX.
%============================================================
\documentclass[12pt,a4paper,twoside]{article}%<-- twoside: bat dau trang/chan trang chan-le
\usepackage{iftex}
\usepackage{amsmath,amssymb,tikz,mathrsfs,tkz-tab}
\usepackage{graphicx}
\usepackage{pgfplots}% đồ thị hàm số
\pgfplotsset{compat=1.17}
%\usepackage{mathptmx}
\usepackage{mathpazo}
\ifXeTeX% XeLaTeX: font chữ có đủ dấu tiếng Việt, giữ nguyên font toán mathpazo
	\usepackage[no-math]{fontspec}
	\setmainfont{TeX Gyre Pagella}
	\newfontfamily\EZqag{TeX Gyre Adventor}[NFSSFamily=qag,Ligatures=TeX]
\fi
\usepackage[many]{tcolorbox}
\usetikzlibrary{math,through,calc,intersections,angles,quotes,shapes,shapes.geometric,arrows,patterns,matrix,chains,arrows.meta,decorations.shapes,decorations.fractals,decorations.markings,shadows,decorations.pathreplacing}
\usepackage[top=1.6cm, bottom=1.8cm, left=1.8cm, right=1.8cm]{geometry}% v2: top 1.6cm (bản cũ 1.5cm) để chữ không sát dải đầu trang
\providecommand{\CheDo}{loigiai}% loigiai | dethi | solcolor | color | book
\expandafter\usepackage\expandafter[\CheDo]{ex_test}
% \tenchuong: khai báo ở file chính (trước \input) hoặc \renewcommand sau \input
%============================================================
%  [GỘP TỪ eduzy-de.tex] EDUZY - MÀU NÂU ĐỒNG / VÀNG KIM
%  + ĐẦU TRANG, CHÂN TRANG KIỂU DẢI NGANG (v2), ĐỐI XỨNG TRANG CHẴN / TRANG LẺ.
%  (không cần file eduzy-de.tex riêng nữa)
%  Khối này phải đặt SAU \usepackage[...]{ex_test} và SAU \def\tenchuong{...}
%============================================================
\usepackage{lastpage}
\usepackage{needspace}
%--- DÀN TRANG ĐỀ THI: tách câu (THÊM 05/11/2026; chỉ thêm lệnh mới, KHÔNG đổi hành vi mặc định) ---
%  Hai công tắc, mặc định tắt (= preamble chạy y như cũ). Bật bằng cách đặt TRƯỚC %============================================================
%  PREAMBLE EDUZY v2 - dùng chung cho đề thi, bài học, chuyên đề (gói ex_test)
%  Khác bản preamble-eduzy.tex ở đầu trang, chân trang: dải nền nâu đồng tràn ngang,
%  logo + Eduzy.com.vn, ô chữ vàng kem cạnh xiên, chân trang "Trang (số)";
%  logo Eduzy chìm giữa mọi trang (tắt: \logochimtat, bật lại: \logochimbat).
%  Cách dùng ở file chính (xem mau-chuan.tex, mau-chuyen-de.tex):
%      \def\CheDo{loigiai}          % loigiai | dethi | solcolor | color | book
%      \def\tenchuong{<dòng chữ trên đầu trang>}
%      \def\thuonghieu{<chữ ở chân trang lẻ>}   % không khai báo thì dùng mặc định
%      \def\tacgia{<chữ ở chân trang chẵn>}      % không khai báo thì dùng mặc định
%      \def\khauhieu{<chữ cạnh logo ở đầu trang>}  % tùy chọn, mặc định Eduzy.com.vn
%      \input{preamble-eduzy-v2}
%      \TLdapso                      % tài liệu: \TLfull | \TLdapso | \TLan
%      \begin{document} ... \end{document}
%  Logo Eduzy vẽ bằng TikZ ngay trong file (\eduzylogo), không cần file ảnh.
%  Biên dịch: pdfLaTeX (máy có gói vietnam) hoặc XeLaTeX.
%============================================================
\documentclass[12pt,a4paper,twoside]{article}%<-- twoside: bat dau trang/chan trang chan-le
\usepackage{iftex}
\usepackage{amsmath,amssymb,tikz,mathrsfs,tkz-tab}
\usepackage{graphicx}
\usepackage{pgfplots}% đồ thị hàm số
\pgfplotsset{compat=1.17}
%\usepackage{mathptmx}
\usepackage{mathpazo}
\ifXeTeX% XeLaTeX: font chữ có đủ dấu tiếng Việt, giữ nguyên font toán mathpazo
	\usepackage[no-math]{fontspec}
	\setmainfont{TeX Gyre Pagella}
	\newfontfamily\EZqag{TeX Gyre Adventor}[NFSSFamily=qag,Ligatures=TeX]
\fi
\usepackage[many]{tcolorbox}
\usetikzlibrary{math,through,calc,intersections,angles,quotes,shapes,shapes.geometric,arrows,patterns,matrix,chains,arrows.meta,decorations.shapes,decorations.fractals,decorations.markings,shadows,decorations.pathreplacing}
\usepackage[top=1.6cm, bottom=1.8cm, left=1.8cm, right=1.8cm]{geometry}% v2: top 1.6cm (bản cũ 1.5cm) để chữ không sát dải đầu trang
\providecommand{\CheDo}{loigiai}% loigiai | dethi | solcolor | color | book
\expandafter\usepackage\expandafter[\CheDo]{ex_test}
% \tenchuong: khai báo ở file chính (trước \input) hoặc \renewcommand sau \input
%============================================================
%  [GỘP TỪ eduzy-de.tex] EDUZY - MÀU NÂU ĐỒNG / VÀNG KIM
%  + ĐẦU TRANG, CHÂN TRANG KIỂU DẢI NGANG (v2), ĐỐI XỨNG TRANG CHẴN / TRANG LẺ.
%  (không cần file eduzy-de.tex riêng nữa)
%  Khối này phải đặt SAU \usepackage[...]{ex_test} và SAU \def\tenchuong{...}
%============================================================
\usepackage{lastpage}
\usepackage{needspace}
%--- DÀN TRANG ĐỀ THI: tách câu (THÊM 05/11/2026; chỉ thêm lệnh mới, KHÔNG đổi hành vi mặc định) ---
%  Hai công tắc, mặc định tắt (= preamble chạy y như cũ). Bật bằng cách đặt TRƯỚC \input{preamble-eduzy-v2}
%  trong file đề (tools/dan_trang/dan_trang.py tự ghi khi đặt \NeedspaceCau):
%    \def\TachCauDS{1}   câu Đúng/Sai được tách: câu dẫn (kèm hình) ở cuối trang, bảng Phát biểu sang trang sau
%                        (câu dẫn và hình không tách nhau, bảng không tách giữa hai trang)
%    \def\TachCauLG{1}   bản lời giải (\CheDo khác dethi): phần đề của câu cao hơn 1/3 trang được phép tách
%  Lệnh dùng trong file đề (giá trị do tools/dan_trang/dan_trang.py đo ở bản dethi):
%    \NeedspaceCau[<cao câu dẫn + hình>]{<cao cả khối câu>}   ngay trên \begin{ex} (câu đầu phần: trên tiêu đề phần)
%        công tắc tắt: đúng bằng \Needspace{<cao cả khối câu>} như cũ;
%        \TachCauDS=1 và có [..] (câu Đúng/Sai được chọn tách): chỉ giữ câu dẫn + hình;
%        \TachCauLG=1 ở bản lời giải: đòi tối đa 1/3 \textheight.
%    \NeedspaceBang{<cao bảng Phát biểu>}%   ngay trước \choiceTFt: \TachCauDS=1 thì giữ trọn bảng trên một trang;
%        công tắc tắt: không làm gì.
\providecommand{\TachCauDS}{0}
\providecommand{\TachCauLG}{0}
\newlength{\EZnsCau}
\newcommand{\NeedspaceCau}[2][]{%
	\setlength{\EZnsCau}{#2}%
	\ifnum\TachCauDS=1 \if\relax\detokenize{#1}\relax\else\setlength{\EZnsCau}{#1}\fi\fi
	\ifnum\TachCauLG=1 \ifdefstring{\kindOPTN}{dethi}{}{\ifdim\EZnsCau>0.3333\textheight\setlength{\EZnsCau}{0.3333\textheight}\fi}\fi
	\Needspace{\EZnsCau}}
\newcommand{\NeedspaceBang}[1]{\ifnum\TachCauDS=1 \Needspace{#1}\fi}

%--- Thang màu: đậm -> nhạt
\definecolor{gold1}{RGB}{126,96,0}    % nâu đồng đậm  #7E6000
\definecolor{gold2}{RGB}{192,144,3}   % vàng kim      #C09003
\definecolor{gold3}{RGB}{255,215,99}  % vàng nhạt     #FFD763
\definecolor{gold4}{RGB}{255,231,153} % vàng rất nhạt #FFE799
\definecolor{gold5}{RGB}{255,243,205} % vàng kem      #FFF3CD
%--- Ghi đè tông xanh/đỏ cũ
\definecolor{myblue}{RGB}{126,96,0}
\colorlet{blue}{gold1}
\colorlet{cyan}{gold2}
\colorlet{teal}{gold2}

%--- Tên chương mặc định (file đề nên tự khai báo lại)
\providecommand{\tenchuong}{}
%--- Chân trang
\providecommand{\thuonghieu}{TOÁN 12 -- Eduzy.com.vn}
\providecommand{\tacgia}{Tô Gia Bảo -- 0353703783}

%--- Chữ cạnh logo ở đầu trang
\providecommand{\khauhieu}{Eduzy.com.vn}

%====================== THÔNG SỐ ĐẦU TRANG, CHÂN TRANG (v2) ======================
%  Dải nền đậm tràn ngang trang; ô chữ vàng kem có cạnh xiên, kèm 2 vạch xiên chuyển màu.
%  Đầu trang lẻ : logo + \khauhieu (trái)   | ô chữ \tenchuong (phải)
%  Chân trang lẻ: ô chữ \thuonghieu (trái)  | Trang (số) (phải)
%  Trang chẵn soi gương trang lẻ; chân trang chẵn ghi \tacgia.
%--- Màu: chỉ dùng thang vàng kim Eduzy
\colorlet{ezNen}{gold1}% dải nền đậm
\colorlet{ezVach1}{gold3}% vạch xiên sát dải nền (đậm hơn)
\colorlet{ezVach2}{gold4}% vạch xiên sát ô chữ (nhạt hơn)
\colorlet{ezO}{gold5}% ô chữ vàng kem
%--- Kích thước
\newlength{\EZdai}\setlength{\EZdai}{1.1cm}% chiều cao dải nền
\newlength{\EZo}\setlength{\EZo}{0.68cm}% chiều cao ô chữ
\newlength{\EZxien}\setlength{\EZxien}{0.55cm}% độ lệch ngang của cạnh xiên
\newlength{\EZvach}\setlength{\EZvach}{0.75cm}% bề rộng mỗi vạch xiên
\newlength{\EZkhe}\setlength{\EZkhe}{0.12cm}% khe trắng giữa vạch và ô chữ
\newlength{\EZdem}\setlength{\EZdem}{0.4cm}% khoảng đệm từ chữ tới cạnh xiên
\newlength{\EZle}\setlength{\EZle}{1.8cm}% lề trái/phải của chữ (bằng lề thân bài)
\newlength{\EZrongmin}\setlength{\EZrongmin}{4.5cm}% ô chữ rộng tối thiểu
\newlength{\EZrong}% bề rộng chữ trong ô (tính tự động)
%--- Font
\newcommand{\EZfontO}{\fontsize{10.5pt}{12pt}\selectfont\bfseries}% chữ trong ô
\newcommand{\EZfontKH}{\fontfamily{qag}\fontsize{11pt}{12pt}\selectfont\bfseries}% \khauhieu cạnh logo
\newcommand{\EZfontTrang}{\fontsize{11.5pt}{12pt}\selectfont\bfseries}% chữ "Trang" và số trang
%--- Đo bề rộng chữ trong ô
\newcommand{\EZdorong}[1]{%
	\settowidth{\EZrong}{\EZfontO #1}%
	\ifdim\EZrong<\EZrongmin\setlength{\EZrong}{\EZrongmin}\fi}
%--- Logo + chữ trắng \khauhieu (dựng sẵn một lần khi vào \begin{document})
\newsavebox{\EZlogobox}
\AtBeginDocument{\sbox{\EZlogobox}{%
		\raisebox{-0.5\height}{\eduzylogo[1cm]}\hspace{0.15cm}%
		\raisebox{-0.5\height}{\color{white}\EZfontKH\khauhieu}}}
%--- Hai vạch xiên + khe trắng. #1: điểm gốc ở mép dưới, #2: hướng (1 sang phải, -1 sang trái),
%    #3: hướng nghiêng (1 kiểu "/", -1 kiểu "\")
\newcommand{\EZvachxien}[3]{%
	\fill[white] (#1) --++({(#3)*\EZxien},\EZo) --++({(#2)*\EZkhe},0) --++({-(#3)*\EZxien},-\EZo) --cycle;
	\fill[ezVach2] ([xshift=(#2)*\EZkhe]#1) --++({(#3)*\EZxien},\EZo) --++({(#2)*\EZvach},0) --++({-(#3)*\EZxien},-\EZo) --cycle;
	\fill[ezVach1] ([xshift=(#2)*(\EZkhe+\EZvach)]#1) --++({(#3)*\EZxien},\EZo) --++({(#2)*\EZvach},0) --++({-(#3)*\EZxien},-\EZo) --cycle;
	\draw[white,line width=0.6pt] ([xshift=(#2)*(\EZkhe+\EZvach)]#1) --++({(#3)*\EZxien},\EZo);
}
%--- Chữ trong ô (căn giữa theo chiều cao ô, không lệch vì dấu tiếng Việt)
\tikzset{ezchu/.style={inner sep=0pt,text=ezNen,text height=0.72em,text depth=0.22em}}

%====================== LOGO CHÌM GIỮA TRANG ======================
%  Vẽ cùng đầu trang nên nằm dưới chữ của thân bài, có trên mọi trang.
%  Tắt: \logochimtat (từ trang đó trở đi), bật lại: \logochimbat.
%  Đổi cỡ, độ đậm: khai báo trước \input, ví dụ \def\EZchimcao{9cm} \def\EZchimdo{0.1}
\providecommand{\EZchimcao}{10cm}% chiều cao logo chìm
\providecommand{\EZchimdo}{0.12}% độ đậm (0 = trong suốt, 1 = rõ hẳn)
\newif\ifEZchim \EZchimtrue
\newcommand{\logochimtat}{\EZchimfalse}
\newcommand{\logochimbat}{\EZchimtrue}
%  Dựng sẵn logo mờ một lần khi vào \begin{document}; transparency group làm mờ cả logo
%  một lượt nên không lộ vệt ở chỗ các dải màu chồng nhau.
\newsavebox{\EZchimtho}
\newsavebox{\EZchimbox}
\AtBeginDocument{%
	\sbox{\EZchimtho}{\eduzylogo[\EZchimcao]}%
	\sbox{\EZchimbox}{\begin{tikzpicture}
			\begin{scope}[transparency group,opacity=\EZchimdo]
				\node[inner sep=0pt] {\usebox{\EZchimtho}};
			\end{scope}
	\end{tikzpicture}}}
\newcommand{\EZvechim}{%
	\ifEZchim
		\node[inner sep=0pt] at (current page.center) {\usebox{\EZchimbox}};
	\fi}

%====================== ĐẦU TRANG - TRANG LẺ ======================
%  (logo bên trái, ô chữ \tenchuong bên phải)
\newcommand{\myfancyheadOdd}{%
	\EZdorong{\tenchuong}%
	\begin{tikzpicture}[remember picture,overlay]
		\coordinate (T) at (current page.north west);
		\coordinate (P) at (current page.north east);
		%--- logo chìm giữa trang
		\EZvechim
		%--- dải nền
		\fill[ezNen] (T) rectangle ([yshift=-\EZdai]P);
		%--- ô chữ: X là góc dưới bên trái, cạnh trái xiên kiểu "/"
		\coordinate (X) at ([xshift=-\EZle-\EZrong-\EZdem-\EZxien,yshift=-\EZdai]P);
		\fill[ezO] (X) --++(\EZxien,\EZo) -- ([yshift=-\EZdai+\EZo]P) -- ([yshift=-\EZdai]P) --cycle;
		\EZvachxien{X}{-1}{1}
		\node[ezchu,anchor=east] at ([xshift=-\EZle,yshift=-\EZdai+0.5*\EZo]P) {\EZfontO\tenchuong};
		%--- logo + tên
		\node[inner sep=0pt,anchor=west] at ([xshift=\EZle,yshift=-0.5*\EZdai]T) {\usebox{\EZlogobox}};
	\end{tikzpicture}%
}

%====================== ĐẦU TRANG - TRANG CHẴN =====================
%  (soi gương: ô chữ \tenchuong bên trái, logo bên phải)
\newcommand{\myfancyheadEven}{%
	\EZdorong{\tenchuong}%
	\begin{tikzpicture}[remember picture,overlay]
		\coordinate (T) at (current page.north west);
		\coordinate (P) at (current page.north east);
		%--- logo chìm giữa trang
		\EZvechim
		%--- dải nền
		\fill[ezNen] (T) rectangle ([yshift=-\EZdai]P);
		%--- ô chữ: X là góc dưới bên phải, cạnh phải xiên kiểu "\"
		\coordinate (X) at ([xshift=\EZle+\EZrong+\EZdem+\EZxien,yshift=-\EZdai]T);
		\fill[ezO] (X) --++(-\EZxien,\EZo) -- ([yshift=-\EZdai+\EZo]T) -- ([yshift=-\EZdai]T) --cycle;
		\EZvachxien{X}{1}{-1}
		\node[ezchu,anchor=west] at ([xshift=\EZle,yshift=-\EZdai+0.5*\EZo]T) {\EZfontO\tenchuong};
		%--- logo + tên
		\node[inner sep=0pt,anchor=east] at ([xshift=-\EZle,yshift=-0.5*\EZdai]P) {\usebox{\EZlogobox}};
	\end{tikzpicture}%
}

%====================== CHÂN TRANG - TRANG LẺ ======================
%  (ô chữ \thuonghieu bên trái, "Trang (số)" bên phải)
\newcommand{\myfancyfootOdd}{%
	\EZdorong{\thuonghieu}%
	\begin{tikzpicture}[remember picture,overlay]
		\coordinate (T) at (current page.south west);
		\coordinate (P) at (current page.south east);
		%--- dải nền
		\fill[ezNen] (T) rectangle ([yshift=\EZdai]P);
		%--- ô chữ sát mép trên dải: X là góc dưới bên phải, cạnh phải xiên kiểu "/"
		\coordinate (X) at ([xshift=\EZle+\EZrong+\EZdem,yshift=\EZdai-\EZo]T);
		\fill[ezO] ([yshift=\EZdai-\EZo]T) -- (X) --++(\EZxien,\EZo) -- ([yshift=\EZdai]T) --cycle;
		\EZvachxien{X}{1}{1}
		\node[ezchu,anchor=west] at ([xshift=\EZle,yshift=\EZdai-0.5*\EZo]T) {\EZfontO\thuonghieu};
		%--- số trang
		\node[circle,draw=white,line width=0.9pt,text=white,inner sep=0pt,minimum size=0.62cm,anchor=east]
		(so) at ([xshift=-\EZle,yshift=0.5*\EZdai]P) {\EZfontTrang\thepage};
		\node[inner sep=0pt,text=white,anchor=base east] (tr) at ([xshift=-0.15cm]so.base west) {\EZfontTrang Trang};
	\end{tikzpicture}%
}

%====================== CHÂN TRANG - TRANG CHẴN ====================
%  (soi gương: "Trang (số)" bên trái, ô chữ \tacgia bên phải)
\newcommand{\myfancyfootEven}{%
	\EZdorong{\tacgia}%
	\begin{tikzpicture}[remember picture,overlay]
		\coordinate (T) at (current page.south west);
		\coordinate (P) at (current page.south east);
		%--- dải nền
		\fill[ezNen] (T) rectangle ([yshift=\EZdai]P);
		%--- ô chữ sát mép trên dải: X là góc dưới bên trái, cạnh trái xiên kiểu "\"
		\coordinate (X) at ([xshift=-\EZle-\EZrong-\EZdem,yshift=\EZdai-\EZo]P);
		\fill[ezO] ([yshift=\EZdai-\EZo]P) -- (X) --++(-\EZxien,\EZo) -- ([yshift=\EZdai]P) --cycle;
		\EZvachxien{X}{-1}{-1}
		\node[ezchu,anchor=east] at ([xshift=-\EZle,yshift=\EZdai-0.5*\EZo]P) {\EZfontO\tacgia};
		%--- số trang
		%    (vòng tròn ẩn so0 nằm giữa dải để lấy đường chân chữ, giống trang lẻ)
		\node[circle,inner sep=0pt,minimum size=0.62cm] (so0) at ([yshift=0.5*\EZdai]T) {\EZfontTrang\phantom{\thepage}};
		\node[inner sep=0pt,text=white,anchor=base west] (tr) at ($(so0.base -| T)+(\EZle,0)$) {\EZfontTrang Trang};
		\node[circle,draw=white,line width=0.9pt,text=white,inner sep=0pt,minimum size=0.62cm,anchor=base west]
		(so) at ([xshift=0.15cm]tr.base east) {\EZfontTrang\thepage};
	\end{tikzpicture}%
}

%--- Giữ tương thích với các file cũ đang gọi \myfancyhead, \myfancyfoot
\newcommand{\myfancyhead}{\myfancyheadOdd}
\newcommand{\myfancyfoot}{\myfancyfootOdd}

%======================= GẮN VÀO TRANG =============================
\pagestyle{fancy}
\fancyhf{}
\fancyhead[LO]{\myfancyheadOdd}
\fancyhead[RE]{\myfancyheadEven}
\fancyfoot[LO]{\myfancyfootOdd}
\fancyfoot[RE]{\myfancyfootEven}
\renewcommand{\headrulewidth}{0pt}
\renewcommand{\footrulewidth}{0pt}
\fancypagestyle{plain}{%
	\fancyhf{}
	\fancyhead[LO]{\myfancyheadOdd}
	\fancyhead[RE]{\myfancyheadEven}
	\fancyfoot[LO]{\myfancyfootOdd}
	\fancyfoot[RE]{\myfancyfootEven}
	\renewcommand{\headrulewidth}{0pt}
	\renewcommand{\footrulewidth}{0pt}
}
%=================== HẾT PHẦN GỘP TỪ eduzy-de.tex ===================
\newcommand{\indam}[1]{
	{\fontsize{12pt}{9pt}\sffamily \selectfont\bfseries\color{gold1} #1}
}
\newcommand{\ind}[1]{
	{\fontsize{14pt}{9pt}\fontfamily{qag} \selectfont\bfseries\color{gold1} #1}
}
\renewcommand{\FalseEX}{\stepcounter{dapan}\textbf{\selectfont\color{gold1}\Alph{dapan}.}\ignorespaces}
\renewtheorem{ex}{\color{gold1}Câu}
\renewcommand{\baselinestretch}{1.1}
\newcommand{\hoac}[1]{
	\left[\begin{aligned}#1\end{aligned}\right.}
\newcommand{\heva}[1]{
	\left\{\begin{aligned}#1\end{aligned}\right.}
\usepackage{esvect}
\def\vec{\vv}
\def\vt{\vv}
\def\overrightarrow{\vv}
%====================== LOGO EDUZY VẼ BẰNG TIKZ ======================
%  Đường viền dò từ logo gốc; dải màu vàng kim đo trực tiếp trên ảnh gốc
%  (nghiêng 21 độ, tô bằng 80 dải hẹp nên hiển thị giống nhau trên mọi trình xem PDF).
%  Dùng: \eduzylogo          -> cao 2.4cm (mặc định)
%        \eduzylogo[1.5cm]   -> cao 1.5cm
\definecolor{ezText}{RGB}{245,200,46}
\newcommand{\eduzylogo}[1][2.4cm]{%
	\resizebox{!}{#1}{%
		\begin{tikzpicture}[x=0.01cm,y=0.01cm]
			\path[use as bounding box] (0,0) rectangle (224,-224);
			% mũ tốt nghiệp + vương miện
			\begin{scope}
				\clip
				(119.4,-176.1)..controls(118.9,-176.1)and(116.8,-176)..(114.6,-176)
				..controls(101.4,-176)and(85.8,-175.9)..(83.7,-175.7)
				..controls(81.8,-175.6)and(81,-175.6)..(80.2,-175.7)
				..controls(78.8,-175.9)and(77.7,-175.9)..(76.2,-175.7)
				..controls(74.9,-175.5)and(70.8,-175.5)..(64.3,-175.7)
				..controls(59.1,-175.9)and(58.3,-175.7)..(58.6,-173.7)
				..controls(58.6,-172.8)and(59.6,-169.6)..(60.3,-168)
				..controls(60.8,-166.9)and(61.4,-165.8)..(62.3,-165)
				..controls(63.1,-164.2)and(64.3,-163.4)..(65,-163.2)
				..controls(65.5,-163.1)and(72.3,-163)..(76.6,-163.1)
				..controls(77.9,-163.1)and(80.1,-163.1)..(81.4,-163)
				..controls(82.8,-162.9)and(86.3,-162.9)..(89.2,-162.9)
				..controls(93,-162.9)and(94.9,-162.8)..(96,-162.7)
				..controls(97.2,-162.5)and(101,-162.5)..(114,-162.4)
				..controls(123,-162.3)and(134.9,-162.2)..(140.2,-162.1)
				..controls(150.1,-162)and(152.2,-162.1)..(153.6,-162.4)
				..controls(154.4,-162.6)and(157.1,-163.8)..(158.1,-164.4)
				..controls(158.8,-164.8)and(159.7,-165.7)..(159.9,-166.3)
				..controls(160,-166.5)and(160.2,-167.1)..(160.2,-167.6)
				..controls(160.3,-168)and(160.6,-169.4)..(160.9,-170.6)
				..controls(161.2,-171.7)and(161.5,-173.1)..(161.6,-173.6)
				..controls(161.8,-174.4)and(161.8,-174.5)..(161.5,-174.9)
				..controls(161.2,-175.6)and(160.8,-175.9)..(159.9,-176)
				..controls(159.2,-176.1)and(154.1,-176.2)..(151.5,-176.2)
				..controls(148.1,-176.1)and(122.6,-176.1)..(121.6,-176.2)
				..controls(120.8,-176.2)and(119.8,-176.2)..(119.4,-176.1)--cycle
				(130.4,-160)..controls(129,-160)and(119.8,-160)..(110.1,-159.9)
				..controls(100.3,-159.8)and(88.5,-159.7)..(83.9,-159.6)
				..controls(77.8,-159.5)and(75.3,-159.5)..(74.5,-159.6)
				..controls(73.9,-159.7)and(72.5,-159.8)..(71.3,-159.8)
				..controls(69.6,-159.8)and(69,-159.8)..(68.5,-159.6)
				..controls(67.6,-159.3)and(66.8,-158.6)..(66.2,-157.6)
				..controls(66,-157.3)and(65.6,-156.5)..(65.3,-156.1)
				..controls(64.6,-155.1)and(64.5,-154.9)..(63.7,-152.9)
				..controls(63.2,-151.6)and(62.9,-150.9)..(61.8,-149.6)
				..controls(61.5,-149.2)and(61.1,-148.5)..(60.9,-147.9)
				..controls(60.7,-147.4)and(60.3,-146.6)..(60.1,-146.1)
				..controls(59.9,-145.5)and(59.5,-144.9)..(59,-144.3)
				..controls(58.5,-143.7)and(58.2,-143.2)..(57.9,-142.5)
				..controls(57.7,-141.9)and(57.3,-141.1)..(57,-140.7)
				..controls(56.6,-140.3)and(56.1,-139.3)..(55.8,-138.5)
				..controls(55.5,-137.8)and(55.1,-136.8)..(54.9,-136.5)
				..controls(54.7,-136.1)and(54.3,-135.2)..(53.9,-134.5)
				..controls(53.6,-133.8)and(53,-132.7)..(52.6,-132)
				..controls(51.6,-130.5)and(51.4,-130.2)..(51.1,-129)
				..controls(50.6,-127.5)and(50.6,-127.4)..(49.7,-126.3)
				..controls(49.2,-125.7)and(48.6,-124.8)..(48.3,-124.2)
				..controls(48,-123.6)and(47.5,-122.7)..(47.3,-122.2)
				..controls(47,-121.8)and(46.5,-120.7)..(46.1,-119.8)
				..controls(45.8,-118.9)and(45.2,-117.8)..(44.9,-117.2)
				..controls(44.2,-116.1)and(43.3,-114.4)..(42.3,-112.2)
				..controls(41.8,-111.3)and(41.3,-110.4)..(41.1,-110.1)
				..controls(40.5,-109.4)and(40.3,-109)..(40,-107.8)
				..controls(39.8,-107)and(39.6,-106.5)..(39.2,-106)
				..controls(38.9,-105.6)and(38.4,-104.7)..(38,-104)
				..controls(37.7,-103.4)and(37.2,-102.5)..(37,-102.1)
				..controls(36.7,-101.7)and(36.3,-100.8)..(36,-100.1)
				..controls(35.6,-98.9)and(35.3,-98.4)..(34.6,-97.6)
				..controls(34.2,-97.1)and(34.2,-96.3)..(34.6,-96)
				..controls(35,-95.7)and(35.5,-95.9)..(36.2,-96.7)
				..controls(36.5,-97)and(37.2,-97.5)..(37.6,-97.7)..controls(38,-98)and(38.6,-98.4)..(39.1,-98.9)
				..controls(39.6,-99.4)and(40.3,-100)..(40.6,-100.3)
				..controls(41.3,-100.8)and(42.4,-101.8)..(43.7,-102.9)
				..controls(44,-103.2)and(44.6,-103.7)..(45,-103.9)
				..controls(45.4,-104.2)and(45.9,-104.6)..(46.1,-104.8)
				..controls(47.1,-105.9)and(48.8,-107.1)..(50.5,-107.9)
				..controls(51.8,-108.6)and(52.3,-109)..(53.4,-110)
				..controls(54.8,-111.2)and(55.6,-111.6)..(57,-112.1)
				..controls(57.5,-112.3)and(58.5,-112.8)..(59.4,-113.3)
				..controls(60.2,-113.7)and(61.3,-114.3)..(61.9,-114.5)
				..controls(62.4,-114.7)and(63.3,-115.1)..(63.8,-115.3)
				..controls(64.4,-115.6)and(65.4,-115.9)..(66.2,-116.1)
				..controls(67,-116.3)and(68.4,-116.6)..(69.2,-116.9)
				..controls(72.3,-117.9)and(72.8,-118)..(78,-118.1)
				..controls(84,-118.3)and(86.1,-118.1)..(88.1,-117.1)
				..controls(89.4,-116.5)and(93,-115)..(93.7,-114.9)
				..controls(94.8,-114.6)and(95.4,-114.2)..(95.9,-113.4)
				..controls(96.2,-113)and(96.7,-112.5)..(97.1,-112.2)
				..controls(98,-111.6)and(101.3,-108.4)..(102.1,-107.2)
				..controls(102.4,-106.8)and(102.9,-106.2)..(103.2,-105.8)
				..controls(103.5,-105.5)and(104,-104.8)..(104.2,-104.2)
				..controls(104.5,-103.6)and(105,-102.8)..(105.3,-102.4)
				..controls(105.6,-102)and(106.1,-101.2)..(106.5,-100.6)
				..controls(107.1,-99.5)and(108.6,-97.5)..(109.5,-96.5)
				..controls(109.9,-96.1)and(110.1,-96)..(110.4,-96)
				..controls(111.1,-96)and(111.6,-96.5)..(111.8,-97.6)
				..controls(112,-98.6)and(112.3,-99.3)..(112.8,-99.9)
				..controls(113.7,-100.9)and(115.2,-103.1)..(115.7,-104.2)
				..controls(116.1,-105)and(116.5,-105.5)..(116.9,-106)
				..controls(117.3,-106.4)and(117.8,-106.9)..(118,-107.3)
				..controls(118.2,-107.6)and(118.6,-108)..(119,-108.3)
				..controls(119.6,-108.8)and(121.1,-110.3)..(121.8,-111.2)
				..controls(122.1,-111.7)and(122.5,-112)..(123.4,-112.4)
				..controls(124.8,-113.2)and(125.8,-113.8)..(127,-114.6)
				..controls(127.5,-114.9)and(128.1,-115.3)..(128.2,-115.4)
				..controls(128.4,-115.5)and(129.3,-115.8)..(130.2,-116)
				..controls(131.1,-116.2)and(132.5,-116.6)..(133.3,-116.9)
				..controls(135.7,-117.9)and(136.1,-117.9)..(142.5,-117.8)
				..controls(145.8,-117.8)and(148.4,-117.7)..(148.8,-117.6)
				..controls(149.2,-117.6)and(150.2,-117.2)..(151,-116.8)
				..controls(152,-116.4)and(152.8,-116)..(153.6,-115.9)
				..controls(154.5,-115.7)and(155.2,-115.4)..(156.3,-114.9)
				..controls(157.1,-114.5)and(158.2,-114)..(158.7,-113.8)
				..controls(159.1,-113.7)and(159.8,-113.3)..(160.2,-113)
				..controls(160.6,-112.7)and(161.3,-112.4)..(161.7,-112.2)
				..controls(162.7,-111.9)and(163.2,-111.7)..(164.1,-111.1)
				..controls(164.5,-110.8)and(165.3,-110.4)..(165.9,-110)
				..controls(166.5,-109.7)and(167.3,-109.2)..(167.6,-108.9)
				..controls(167.9,-108.6)and(168.6,-108.2)..(169,-108)
				..controls(169.4,-107.8)and(170.2,-107.3)..(170.7,-106.9)
				..controls(171.2,-106.5)and(172.1,-105.8)..(172.7,-105.4)
				..controls(173.2,-105)and(174.1,-104.4)..(174.5,-104)
				..controls(174.9,-103.7)and(175.7,-103.1)..(176.2,-102.6)
				..controls(176.7,-102.2)and(177.4,-101.5)..(177.9,-101.2)
				..controls(178.3,-100.8)and(178.8,-100.4)..(179,-100.2)
				..controls(179.2,-100)and(179.6,-99.6)..(180,-99.3)
				..controls(180.4,-99.1)and(180.9,-98.6)..(181.2,-98.4)
				..controls(181.5,-98.1)and(182.1,-97.7)..(182.5,-97.5)
				..controls(183.2,-97.2)and(183.5,-96.9)..(184.3,-95.9)
				..controls(184.6,-95.6)and(184.6,-95.6)..(184.8,-95.7)
				..controls(185,-95.8)and(185.1,-96)..(185.1,-96.2)
				..controls(185.1,-96.5)and(185,-96.7)..(184.6,-97.2)
				..controls(184.1,-97.8)and(184,-97.9)..(183.9,-98.8)
				..controls(183.7,-100.3)and(183.3,-101.3)..(182.9,-101.8)
				..controls(182.6,-102.1)and(182.3,-102.6)..(182.2,-102.9)
				..controls(182,-103.3)and(181.6,-104)..(181.2,-104.6)
				..controls(180.9,-105.2)and(180.5,-106.1)..(180.4,-106.6)
				..controls(179.9,-108)and(179.7,-108.5)..(179,-109.4)
				..controls(178.6,-109.8)and(178.2,-110.5)..(177.9,-111)
				..controls(177.7,-111.5)and(177.3,-112.3)..(177,-112.7)
				..controls(176.6,-113.2)and(176.4,-113.7)..(176.3,-114.2)
				..controls(176.2,-114.7)and(175.9,-115.2)..(175.4,-116)
				..controls(175,-116.6)and(174.4,-117.7)..(174,-118.5)
				..controls(173.7,-119.2)and(173.3,-120.1)..(173,-120.5)
				..controls(172.8,-120.9)and(172.4,-121.8)..(172.1,-122.5)
				..controls(171.8,-123.2)and(171.3,-124.1)..(171,-124.6)
				..controls(170.8,-125.1)and(170.3,-126)..(169.9,-126.6)
				..controls(169.6,-127.3)and(169.2,-128.1)..(169,-128.3)
				..controls(168.8,-128.6)and(168.5,-129.3)..(168.3,-129.8)
				..controls(168.1,-130.3)and(167.8,-131.1)..(167.5,-131.6)
				..controls(166.9,-132.5)and(164.2,-137.9)..(164.2,-138.1)
				..controls(164.2,-138.1)and(163.9,-138.8)..(163.4,-139.7)
				..controls(162.9,-140.5)and(162.4,-141.5)..(162.3,-141.9)
				..controls(161.9,-143)and(161.6,-143.4)..(160.9,-144.1)
				..controls(160.4,-144.6)and(160.3,-144.8)..(160.1,-145.4)
				..controls(160,-145.8)and(159.6,-146.8)..(159.2,-147.8)
				..controls(158.7,-148.8)and(158.1,-150)..(157.9,-150.6)
				..controls(157.6,-151.1)and(157.2,-151.9)..(156.9,-152.4)
				..controls(156.5,-152.8)and(156.1,-153.7)..(155.8,-154.4)
				..controls(155.6,-155)and(155.2,-155.9)..(154.9,-156.3)
				..controls(154.7,-156.7)and(154.4,-157.4)..(154.2,-157.9)
				..controls(154.1,-158.4)and(153.8,-159)..(153.5,-159.3)
				..controls(152.8,-160.1)and(153,-160.1)..(142.4,-160.1)
				..controls(137.2,-160.1)and(131.8,-160.1)..(130.4,-160)--cycle
				(141.5,-103.7)..controls(140.8,-103)and(139.8,-102.1)..(139.2,-101.6)
				..controls(138.5,-101.2)and(137.9,-100.6)..(137.8,-100.4)
				..controls(137,-99.5)and(134,-96.5)..(133.2,-96)
				..controls(132.7,-95.7)and(131.8,-94.9)..(131.1,-94.3)
				..controls(129.4,-92.7)and(127.9,-91.6)..(126.6,-91.1)
				..controls(125.2,-90.5)and(124.1,-90)..(122.9,-89.1)
				..controls(122.1,-88.6)and(121.3,-88.2)..(119.9,-87.7)
				..controls(118.8,-87.3)and(117.5,-86.8)..(116.9,-86.5)
				..controls(114.3,-85.5)and(113.7,-85.4)..(109.7,-85.5)
				..controls(105.3,-85.7)and(103.9,-85.9)..(101.6,-87)
				..controls(100.7,-87.4)and(99.4,-87.8)..(97.9,-88)
				..controls(97,-88.2)and(96.6,-88.3)..(95.8,-88.9)
				..controls(95.4,-89.2)and(94.7,-89.6)..(94.2,-89.8)
				..controls(92.9,-90.3)and(92.3,-90.7)..(90.9,-92.1)
				..controls(90.2,-92.7)and(89.4,-93.3)..(89,-93.5)..controls(87,-94.5)and(85.2,-96)..(83.6,-97.9)
				..controls(83.3,-98.2)and(82.5,-98.9)..(81.8,-99.4)
				..controls(80.7,-100.1)and(80.5,-100.3)..(80,-101.1)
				..controls(78.7,-103.1)and(77.2,-103.9)..(76.4,-103)
				..controls(75.7,-102.2)and(75.7,-103.2)..(75.7,-83.6)
				..controls(75.6,-73.6)and(75.6,-64.8)..(75.5,-64)
				..controls(75.5,-63.2)and(75.5,-61.4)..(75.6,-59.8)--(75.7,-57.1)--(76,-56.8)
				..controls(76.8,-56.2)and(78,-56.3)..(79.8,-57.2)
				..controls(80.4,-57.5)and(81.3,-57.8)..(81.8,-57.9)
				..controls(82.3,-58.1)and(83,-58.4)..(83.6,-58.8)..controls(84.3,-59.2)and(85,-59.5)..(86,-59.8)
				..controls(86.8,-60.1)and(88,-60.6)..(88.7,-60.9)
				..controls(89.3,-61.2)and(90.6,-61.7)..(91.4,-62)
				..controls(92.2,-62.3)and(93.7,-62.9)..(94.6,-63.3)
				..controls(95.6,-63.7)and(96.8,-64.1)..(97.4,-64.3)
				..controls(98,-64.5)and(99,-64.9)..(99.7,-65.2)
				..controls(100.4,-65.6)and(101.5,-66)..(102.3,-66.2)
				..controls(104.2,-66.8)and(105,-67.1)..(106.2,-67.8)
				..controls(107.9,-68.7)and(110.2,-69.4)..(111.3,-69.2)
				..controls(112.5,-69)and(115,-68.2)..(115.6,-67.8)
				..controls(117.3,-66.7)and(117.6,-66.6)..(119.5,-66.2)
				..controls(121.9,-65.8)and(122.1,-65.7)..(123,-65.1)
				..controls(124.1,-64.4)and(124.4,-64.3)..(126,-63.9)
				..controls(127,-63.7)and(127.7,-63.4)..(128.4,-63.1)
				..controls(129,-62.8)and(129.7,-62.5)..(130.6,-62.4)
				..controls(132.1,-62.1)and(134,-61.5)..(135.2,-61)
				..controls(135.7,-60.8)and(136.7,-60.5)..(137.3,-60.3)
				..controls(138.4,-60)and(139.3,-59.7)..(141.1,-59)
				..controls(141.6,-58.7)and(142,-58.6)..(142.5,-58.6)
				..controls(143.2,-58.6)and(143.3,-58.7)..(143.6,-59)
				..controls(144,-59.5)and(144.1,-59.9)..(144,-60.8)
				..controls(144,-61.2)and(144,-71.1)..(144,-82.8)--(143.9,-104.1)--(143.7,-104.5)
				..controls(143.5,-104.7)and(143.3,-104.8)..(143,-104.9)
				..controls(142.6,-104.9)and(142.5,-104.8)..(141.5,-103.7)--cycle
				(27.8,-91.4)..controls(27.1,-90.9)and(26.3,-90)..(26.1,-89.5)
				..controls(26,-89.2)and(25.8,-88.6)..(25.7,-87.9)
				..controls(25.4,-86.2)and(25.6,-84.8)..(26.5,-83.5)
				..controls(27.4,-82.3)and(28.3,-82)..(29.4,-82.3)
				..controls(31.1,-82.8)and(32.1,-84.9)..(32,-87.4)
				..controls(31.9,-89.3)and(31.2,-90.6)..(29.9,-91.4)
				..controls(29.1,-91.9)and(28.6,-91.9)..(27.8,-91.4)--cycle
				(190.1,-91.3)..controls(189.7,-91.2)and(189.3,-90.9)..(189.1,-90.7)
				..controls(188.3,-89.8)and(187.6,-87.7)..(187.6,-86.2)
				..controls(187.6,-85.5)and(187.7,-85.2)..(188.2,-84.4)
				..controls(189.1,-82.6)and(190.1,-81.8)..(191.3,-82.1)
				..controls(192,-82.3)and(193.5,-83.7)..(193.8,-84.4)
				..controls(194.1,-85.1)and(194.2,-87.1)..(194,-88.3)
				..controls(193.8,-89.8)and(193.3,-90.5)..(192.3,-91.3)
				..controls(191.7,-91.7)and(190.9,-91.7)..(190.1,-91.3)--cycle
				(173.7,-83.4)..controls(172.2,-82.5)and(172,-82.5)..(170.6,-82.4)
				..controls(169.7,-82.4)and(168.8,-82.4)..(168.3,-82.5)
				..controls(167.2,-82.7)and(166.9,-82.7)..(166.6,-82.2)
				..controls(166.3,-81.8)and(166.3,-81)..(166.8,-80)
				..controls(167.2,-78.9)and(167.4,-77.4)..(167.7,-74.1)
				..controls(167.8,-72.8)and(167.9,-71.8)..(167.9,-71.7)
				..controls(168,-71.7)and(168,-70.2)..(168,-68.4)..controls(168,-65.4)and(168,-64.9)..(167.7,-63)
				..controls(167.3,-60.5)and(167.1,-59.9)..(166.4,-58.4)
				..controls(166.1,-57.7)and(165.8,-56.9)..(165.7,-56.6)
				..controls(165.7,-56.2)and(165.4,-55.6)..(165.2,-55.2)
				..controls(164.6,-53.8)and(161.9,-50.7)..(161,-50.3)
				..controls(159.8,-49.7)and(157.9,-50)..(155.2,-51.2)
				..controls(154.7,-51.4)and(153.4,-51.8)..(152.2,-52.1)
				..controls(150.8,-52.4)and(149.4,-52.9)..(148.6,-53.2)
				..controls(147.8,-53.6)and(146.7,-54)..(146,-54.2)
				..controls(145.4,-54.4)and(144.5,-54.7)..(144,-54.9)
				..controls(143.6,-55.2)and(142.7,-55.5)..(141.8,-55.7)
				..controls(141,-55.9)and(139.6,-56.3)..(138.7,-56.6)
				..controls(137.8,-57)and(136.3,-57.4)..(135.3,-57.6)
				..controls(133.8,-57.9)and(133.2,-58.1)..(132.1,-58.7)
				..controls(131.1,-59.2)and(130.3,-59.4)..(129.4,-59.6)
				..controls(128.7,-59.8)and(127.8,-60)..(127.4,-60.1)
				..controls(127,-60.3)and(126,-60.5)..(125.1,-60.7)
				..controls(123.2,-61.1)and(121.7,-61.5)..(120.2,-62.1)
				..controls(119.7,-62.3)and(118.1,-62.8)..(116.7,-63.2)
				..controls(114.6,-63.9)and(114,-64.1)..(113,-64.7)
				..controls(111.9,-65.4)and(111.8,-65.4)..(111,-65.5)
				..controls(109.9,-65.5)and(109.3,-65.4)..(107.7,-64.7)
				..controls(107,-64.4)and(106,-64)..(105.3,-63.9)
				..controls(103.5,-63.5)and(102.2,-63)..(101.2,-62.5)
				..controls(100.1,-61.8)and(99.3,-61.5)..(97.8,-61.1)
				..controls(97.2,-61)and(96.2,-60.7)..(95.7,-60.4)
				..controls(95.1,-60.1)and(94.2,-59.8)..(93.5,-59.6)
				..controls(92.9,-59.4)and(91.7,-59)..(91,-58.7)..controls(90.2,-58.4)and(88.9,-58)..(88,-57.7)
				..controls(87.2,-57.5)and(86.1,-57.1)..(85.5,-56.8)
				..controls(85,-56.6)and(84.1,-56.2)..(83.5,-56)..controls(82.9,-55.7)and(81.8,-55.2)..(81,-54.9)
				..controls(80.1,-54.5)and(79.1,-54.2)..(78.7,-54.1)
				..controls(77.9,-53.8)and(77.3,-53.6)..(76.2,-52.9)
				..controls(75.7,-52.6)and(75.2,-52.4)..(74.1,-52.2)
				..controls(72.5,-51.9)and(72.4,-51.8)..(71.3,-51.1)
				..controls(69.9,-50.1)and(69.4,-49.9)..(67.9,-49.8)
				..controls(67,-49.7)and(66.6,-49.6)..(66.3,-49.4)
				..controls(66.1,-49.3)and(65.7,-49)..(65.4,-48.8)
				..controls(65.1,-48.7)and(64.2,-48.4)..(63.6,-48.2)
				..controls(62.9,-48)and(61.7,-47.6)..(60.9,-47.2)
				..controls(60.1,-46.8)and(58.9,-46.4)..(58.2,-46.2)
				..controls(57.6,-46)and(56.8,-45.7)..(56.5,-45.5)
				..controls(54.9,-44.6)and(54.6,-44.5)..(53.2,-44.2)
				..controls(51.7,-43.8)and(51,-43.6)..(49.9,-42.9)
				..controls(49.5,-42.6)and(48.9,-42.4)..(47.5,-42)
				..controls(45.7,-41.6)and(45.6,-41.5)..(44.3,-40.6)
				..controls(43.5,-40.1)and(42.8,-39.6)..(42.7,-39.4)
				..controls(42.4,-38.7)and(43,-38.1)..(44.5,-37.6)..controls(44.9,-37.5)and(46,-37)..(46.9,-36.6)
				..controls(47.9,-36.1)and(48.8,-35.8)..(49.6,-35.6)
				..controls(50.5,-35.4)and(51.3,-35.2)..(52.4,-34.6)
				..controls(53.3,-34.1)and(54.5,-33.7)..(54.9,-33.5)
				..controls(55.4,-33.4)and(56.4,-33)..(57.2,-32.6)
				..controls(58.6,-32)and(60.1,-31.5)..(61.6,-31.2)
				..controls(62.3,-31.1)and(62.7,-31)..(63.2,-30.7)
				..controls(64.1,-30)and(64.3,-29.9)..(65.3,-29.7)
				..controls(65.7,-29.6)and(66.5,-29.2)..(67.1,-28.9)
				..controls(67.6,-28.6)and(68.8,-28.2)..(69.6,-27.9)
				..controls(70.5,-27.7)and(71.6,-27.3)..(72.1,-27)
				..controls(72.6,-26.7)and(73.8,-26.2)..(74.8,-25.9)
				..controls(75.7,-25.5)and(76.7,-25.1)..(77.1,-25)
				..controls(77.4,-24.8)and(78.5,-24.4)..(79.6,-24.1)
				..controls(80.6,-23.8)and(82,-23.2)..(82.8,-22.9)
				..controls(83.5,-22.5)and(84.8,-22)..(85.8,-21.8)
				..controls(86.7,-21.5)and(87.8,-21.1)..(88.1,-20.9)
				..controls(88.4,-20.7)and(89.3,-20.3)..(90.1,-20.1)
				..controls(90.8,-19.8)and(91.9,-19.4)..(92.4,-19.1)
				..controls(93.5,-18.5)and(94.4,-18.2)..(95.4,-18.1)
				..controls(95.7,-18)and(96.4,-17.8)..(96.9,-17.5)
				..controls(97.5,-17.3)and(98.4,-16.9)..(99,-16.6)
				..controls(99.7,-16.3)and(100.9,-15.9)..(101.9,-15.6)
				..controls(102.9,-15.3)and(104.4,-14.7)..(105.3,-14.3)
				..controls(108.1,-13.2)and(108.5,-13.1)..(110.1,-13.1)
				..controls(111.4,-13.1)and(111.5,-13.1)..(112.4,-13.5)
				..controls(113.2,-14)and(115,-14.8)..(116.6,-15.4)
				..controls(116.9,-15.5)and(117.9,-15.8)..(118.8,-16.1)
				..controls(120.9,-16.6)and(121.5,-16.9)..(122.6,-17.9)
				..controls(123,-18.3)and(123.7,-18.8)..(123.9,-18.9)
				..controls(124.6,-19.4)and(126,-19.8)..(127.9,-20.2)
				..controls(129.6,-20.5)and(129.8,-20.6)..(131.1,-21.4)
				..controls(131.5,-21.7)and(132.4,-22.1)..(133.1,-22.3)
				..controls(136.1,-23.2)and(138.3,-24.1)..(140.2,-25)
				..controls(141,-25.5)and(141.5,-25.6)..(142.5,-25.8)
				..controls(145.1,-26.3)and(146,-26.6)..(149.4,-28.1)
				..controls(151.3,-28.9)and(153.5,-29.8)..(154.3,-30.1)
				..controls(155.1,-30.4)and(156.2,-30.8)..(156.6,-31)
				..controls(157.1,-31.2)and(158.2,-31.6)..(159.1,-31.8)
				..controls(160.6,-32.2)and(160.9,-32.3)..(161.7,-32.9)
				..controls(162.4,-33.3)and(162.9,-33.6)..(163.7,-33.8)
				..controls(164.3,-34)and(164.8,-34.1)..(164.9,-34.1)
				..controls(165,-34.1)and(165.6,-34.4)..(166.4,-34.8)
				..controls(167.3,-35.2)and(168.1,-35.4)..(169,-35.6)
				..controls(170,-35.8)and(170.5,-36)..(171.7,-36.6)
				..controls(172.4,-37)and(173.3,-37.4)..(173.6,-37.4)
				..controls(175.6,-37.7)and(176.2,-38)..(176.7,-39.1)
				..controls(176.9,-39.6)and(176.9,-39.7)..(176.7,-40.4)
				..controls(176.4,-41.6)and(175.8,-42.7)..(175.2,-43.4)
				..controls(174.5,-44.1)and(172.6,-45.5)..(171.7,-45.8)
				..controls(170.6,-46.2)and(169.1,-46)..(168.2,-45.4)
				..controls(167.9,-45.2)and(167.3,-44.8)..(166.9,-44.4)
				..controls(166.4,-44.1)and(165.8,-43.6)..(165.4,-43.2)
				..controls(165.1,-42.9)and(164.4,-42.3)..(163.9,-41.9)
				..controls(163.3,-41.5)and(162.7,-40.9)..(162.4,-40.5)
				..controls(161.4,-39.2)and(159.9,-38)..(158.8,-37.8)
				..controls(158.4,-37.8)and(158.3,-37.7)..(158.2,-37.2)
				..controls(158,-36.8)and(157.4,-36.1)..(157.1,-36.1)
				..controls(157,-36.1)and(156.7,-36.1)..(156.5,-36)
				..controls(156.1,-35.9)and(155.9,-35.7)..(155.7,-35.2)
				..controls(155.3,-34.6)and(155.1,-34.4)..(154,-34.1)
				..controls(153.6,-34)and(153.2,-33.7)..(152.9,-33.5)
				..controls(152.4,-32.9)and(151.6,-32.4)..(151.3,-32.4)
				..controls(150.9,-32.4)and(150.4,-32.6)..(150.4,-32.9)
				..controls(150.4,-33.2)and(150.8,-33.4)..(151.4,-33.6)
				..controls(152.1,-33.7)and(152.8,-34.1)..(153.5,-34.9)
				..controls(154.1,-35.5)and(154.3,-35.6)..(155,-35.8)
				..controls(155.7,-36)and(155.7,-36)..(156,-36.6)
				..controls(156.3,-37.3)and(156.7,-37.7)..(157.5,-37.8)
				..controls(158.1,-37.9)and(158.2,-38)..(158.4,-38.8)
				..controls(158.5,-39.4)and(158.6,-39.6)..(159.6,-40.4)
				..controls(160.2,-41)and(160.9,-41.6)..(161.3,-41.8)
				..controls(161.7,-42)and(161.9,-42.3)..(162.2,-42.8)
				..controls(162.4,-43.2)and(162.8,-43.8)..(163.1,-44.1)
				..controls(163.9,-44.8)and(165.1,-46.7)..(165.7,-48)
				..controls(166.2,-49.3)and(166.4,-49.7)..(167.2,-50.5)
				..controls(167.5,-50.8)and(167.9,-51.4)..(168.1,-51.8)
				..controls(168.3,-52.2)and(168.8,-53)..(169.1,-53.6)
				..controls(169.6,-54.3)and(169.8,-54.9)..(169.9,-55.6)
				..controls(170.4,-57.8)and(170.5,-58.1)..(171.2,-59.5)
				..controls(171.6,-60.3)and(172,-61.5)..(172.2,-62.1)
				..controls(172.3,-62.8)and(172.7,-63.7)..(172.9,-64.1)
				..controls(173.2,-64.6)and(173.4,-65.3)..(173.4,-65.6)
				..controls(173.8,-68.8)and(174.2,-70.6)..(174.9,-72.6)
				..controls(175.5,-74.5)and(175.9,-77.1)..(176.1,-80.2)
				..controls(176.3,-83.4)and(175.5,-84.4)..(173.7,-83.4)--cycle;
				\definecolor{ezs}{RGB}{247,179,26}\fill[ezs](-5.7,-156.6)--(-3.4,-157.5)--(61.9,12.7)--(59.7,13.6)--cycle;
				\definecolor{ezs}{RGB}{247,179,22}\fill[ezs](-3.8,-157.3)--(-1.5,-158.2)--(63.9,12)--(61.6,12.9)--cycle;
				\definecolor{ezs}{RGB}{248,179,19}\fill[ezs](-1.8,-158.1)--(0.4,-158.9)--(65.8,11.2)--(63.5,12.1)--cycle;
				\definecolor{ezs}{RGB}{248,178,16}\fill[ezs](0.1,-158.8)--(2.4,-159.7)--(67.7,10.5)--(65.4,11.4)--cycle;
				\definecolor{ezs}{RGB}{249,178,15}\fill[ezs](2,-159.5)--(4.3,-160.4)--(69.6,9.8)--(67.3,10.7)--cycle;
				\definecolor{ezs}{RGB}{251,177,17}\fill[ezs](3.9,-160.3)--(6.2,-161.1)--(71.5,9)--(69.2,9.9)--cycle;
				\definecolor{ezs}{RGB}{252,179,18}\fill[ezs](5.8,-161)--(8.1,-161.9)--(73.4,8.3)--(71.1,9.2)--cycle;
				\definecolor{ezs}{RGB}{252,182,17}\fill[ezs](7.7,-161.7)--(10,-162.6)--(75.3,7.6)--(73,8.5)--cycle;
				\definecolor{ezs}{RGB}{252,185,17}\fill[ezs](9.6,-162.5)--(11.9,-163.4)--(77.2,6.8)--(75,7.7)--cycle;
				\definecolor{ezs}{RGB}{251,189,18}\fill[ezs](11.5,-163.2)--(13.8,-164.1)--(79.2,6.1)--(76.9,7)--cycle;
				\definecolor{ezs}{RGB}{251,191,22}\fill[ezs](13.4,-163.9)--(15.7,-164.8)--(81.1,5.4)--(78.8,6.2)--cycle;
				\definecolor{ezs}{RGB}{252,193,25}\fill[ezs](15.4,-164.7)--(17.6,-165.6)--(83,4.6)--(80.7,5.5)--cycle;
				\definecolor{ezs}{RGB}{253,195,27}\fill[ezs](17.3,-165.4)--(19.6,-166.3)--(84.9,3.9)--(82.6,4.8)--cycle;
				\definecolor{ezs}{RGB}{253,196,31}\fill[ezs](19.2,-166.1)--(21.5,-167)--(86.8,3.2)--(84.5,4)--cycle;
				\definecolor{ezs}{RGB}{254,199,35}\fill[ezs](21.1,-166.9)--(23.4,-167.8)--(88.7,2.4)--(86.4,3.3)--cycle;
				\definecolor{ezs}{RGB}{253,201,38}\fill[ezs](23,-167.6)--(25.3,-168.5)--(90.6,1.7)--(88.3,2.6)--cycle;
				\definecolor{ezs}{RGB}{253,203,40}\fill[ezs](24.9,-168.3)--(27.2,-169.2)--(92.5,1)--(90.3,1.8)--cycle;
				\definecolor{ezs}{RGB}{253,205,42}\fill[ezs](26.8,-169.1)--(29.1,-170)--(94.4,0.2)--(92.2,1.1)--cycle;
				\definecolor{ezs}{RGB}{253,207,45}\fill[ezs](28.7,-169.8)--(31,-170.7)--(96.4,-0.5)--(94.1,0.4)--cycle;
				\definecolor{ezs}{RGB}{253,210,49}\fill[ezs](30.7,-170.5)--(32.9,-171.4)--(98.3,-1.2)--(96,-0.4)--cycle;
				\definecolor{ezs}{RGB}{253,212,54}\fill[ezs](32.6,-171.3)--(34.9,-172.2)--(100.2,-2)--(97.9,-1.1)--cycle;
				\definecolor{ezs}{RGB}{253,214,59}\fill[ezs](34.5,-172)--(36.8,-172.9)--(102.1,-2.7)--(99.8,-1.8)--cycle;
				\definecolor{ezs}{RGB}{253,216,63}\fill[ezs](36.4,-172.7)--(38.7,-173.6)--(104,-3.4)--(101.7,-2.6)--cycle;
				\definecolor{ezs}{RGB}{253,218,65}\fill[ezs](38.3,-173.5)--(40.6,-174.4)--(105.9,-4.2)--(103.6,-3.3)--cycle;
				\definecolor{ezs}{RGB}{254,220,66}\fill[ezs](40.2,-174.2)--(42.5,-175.1)--(107.8,-4.9)--(105.5,-4)--cycle;
				\definecolor{ezs}{RGB}{254,222,68}\fill[ezs](42.1,-175)--(44.4,-175.8)--(109.7,-5.6)--(107.5,-4.8)--cycle;
				\definecolor{ezs}{RGB}{254,224,71}\fill[ezs](44,-175.7)--(46.3,-176.6)--(111.7,-6.4)--(109.4,-5.5)--cycle;
				\definecolor{ezs}{RGB}{254,226,75}\fill[ezs](46,-176.4)--(48.2,-177.3)--(113.6,-7.1)--(111.3,-6.2)--cycle;
				\definecolor{ezs}{RGB}{254,228,79}\fill[ezs](47.9,-177.2)--(50.1,-178)--(115.5,-7.8)--(113.2,-7)--cycle;
				\definecolor{ezs}{RGB}{254,230,82}\fill[ezs](49.8,-177.9)--(52.1,-178.8)--(117.4,-8.6)--(115.1,-7.7)--cycle;
				\definecolor{ezs}{RGB}{254,232,85}\fill[ezs](51.7,-178.6)--(54,-179.5)--(119.3,-9.3)--(117,-8.4)--cycle;
				\definecolor{ezs}{RGB}{254,234,88}\fill[ezs](53.6,-179.4)--(55.9,-180.2)--(121.2,-10)--(118.9,-9.2)--cycle;
				\definecolor{ezs}{RGB}{254,237,91}\fill[ezs](55.5,-180.1)--(57.8,-181)--(123.1,-10.8)--(120.8,-9.9)--cycle;
				\definecolor{ezs}{RGB}{254,239,94}\fill[ezs](57.4,-180.8)--(59.7,-181.7)--(125,-11.5)--(122.8,-10.6)--cycle;
				\definecolor{ezs}{RGB}{254,241,97}\fill[ezs](59.3,-181.6)--(61.6,-182.4)--(127,-12.2)--(124.7,-11.4)--cycle;
				\definecolor{ezs}{RGB}{254,243,101}\fill[ezs](61.2,-182.3)--(63.5,-183.2)--(128.9,-13)--(126.6,-12.1)--cycle;
				\definecolor{ezs}{RGB}{254,245,103}\fill[ezs](63.2,-183)--(65.4,-183.9)--(130.8,-13.7)--(128.5,-12.8)--cycle;
				\definecolor{ezs}{RGB}{254,245,105}\fill[ezs](65.1,-183.8)--(67.4,-184.6)--(132.7,-14.4)--(130.4,-13.6)--cycle;
				\definecolor{ezs}{RGB}{254,245,105}\fill[ezs](67,-184.5)--(69.3,-185.4)--(134.6,-15.2)--(132.3,-14.3)--cycle;
				\definecolor{ezs}{RGB}{254,244,103}\fill[ezs](68.9,-185.2)--(71.2,-186.1)--(136.5,-15.9)--(134.2,-15)--cycle;
				\definecolor{ezs}{RGB}{253,243,101}\fill[ezs](70.8,-186)--(73.1,-186.8)--(138.4,-16.6)--(136.1,-15.8)--cycle;
				\definecolor{ezs}{RGB}{253,241,98}\fill[ezs](72.7,-186.7)--(75,-187.6)--(140.3,-17.4)--(138,-16.5)--cycle;
				\definecolor{ezs}{RGB}{253,239,95}\fill[ezs](74.6,-187.4)--(76.9,-188.3)--(142.2,-18.1)--(140,-17.2)--cycle;
				\definecolor{ezs}{RGB}{253,236,92}\fill[ezs](76.5,-188.2)--(78.8,-189)--(144.2,-18.8)--(141.9,-18)--cycle;
				\definecolor{ezs}{RGB}{253,234,88}\fill[ezs](78.5,-188.9)--(80.7,-189.8)--(146.1,-19.6)--(143.8,-18.7)--cycle;
				\definecolor{ezs}{RGB}{253,232,85}\fill[ezs](80.4,-189.6)--(82.7,-190.5)--(148,-20.3)--(145.7,-19.4)--cycle;
				\definecolor{ezs}{RGB}{253,230,82}\fill[ezs](82.3,-190.4)--(84.6,-191.2)--(149.9,-21)--(147.6,-20.2)--cycle;
				\definecolor{ezs}{RGB}{253,228,78}\fill[ezs](84.2,-191.1)--(86.5,-192)--(151.8,-21.8)--(149.5,-20.9)--cycle;
				\definecolor{ezs}{RGB}{253,226,74}\fill[ezs](86.1,-191.8)--(88.4,-192.7)--(153.7,-22.5)--(151.4,-21.6)--cycle;
				\definecolor{ezs}{RGB}{253,224,70}\fill[ezs](88,-192.6)--(90.3,-193.4)--(155.6,-23.3)--(153.3,-22.4)--cycle;
				\definecolor{ezs}{RGB}{253,222,67}\fill[ezs](89.9,-193.3)--(92.2,-194.2)--(157.5,-24)--(155.3,-23.1)--cycle;
				\definecolor{ezs}{RGB}{253,220,65}\fill[ezs](91.8,-194)--(94.1,-194.9)--(159.5,-24.7)--(157.2,-23.8)--cycle;
				\definecolor{ezs}{RGB}{253,217,63}\fill[ezs](93.7,-194.8)--(96,-195.6)--(161.4,-25.5)--(159.1,-24.6)--cycle;
				\definecolor{ezs}{RGB}{253,215,60}\fill[ezs](95.7,-195.5)--(97.9,-196.4)--(163.3,-26.2)--(161,-25.3)--cycle;
				\definecolor{ezs}{RGB}{253,213,56}\fill[ezs](97.6,-196.2)--(99.9,-197.1)--(165.2,-26.9)--(162.9,-26)--cycle;
				\definecolor{ezs}{RGB}{252,211,53}\fill[ezs](99.5,-197)--(101.8,-197.8)--(167.1,-27.7)--(164.8,-26.8)--cycle;
				\definecolor{ezs}{RGB}{252,209,49}\fill[ezs](101.4,-197.7)--(103.7,-198.6)--(169,-28.4)--(166.7,-27.5)--cycle;
				\definecolor{ezs}{RGB}{252,207,45}\fill[ezs](103.3,-198.4)--(105.6,-199.3)--(170.9,-29.1)--(168.6,-28.2)--cycle;
				\definecolor{ezs}{RGB}{253,205,42}\fill[ezs](105.2,-199.2)--(107.5,-200)--(172.8,-29.9)--(170.6,-29)--cycle;
				\definecolor{ezs}{RGB}{253,202,40}\fill[ezs](107.1,-199.9)--(109.4,-200.8)--(174.7,-30.6)--(172.5,-29.7)--cycle;
				\definecolor{ezs}{RGB}{253,200,38}\fill[ezs](109,-200.6)--(111.3,-201.5)--(176.7,-31.3)--(174.4,-30.4)--cycle;
				\definecolor{ezs}{RGB}{253,198,35}\fill[ezs](111,-201.4)--(113.2,-202.2)--(178.6,-32.1)--(176.3,-31.2)--cycle;
				\definecolor{ezs}{RGB}{253,196,31}\fill[ezs](112.9,-202.1)--(115.2,-203)--(180.5,-32.8)--(178.2,-31.9)--cycle;
				\definecolor{ezs}{RGB}{253,194,28}\fill[ezs](114.8,-202.8)--(117.1,-203.7)--(182.4,-33.5)--(180.1,-32.6)--cycle;
				\definecolor{ezs}{RGB}{253,192,26}\fill[ezs](116.7,-203.6)--(119,-204.5)--(184.3,-34.3)--(182,-33.4)--cycle;
				\definecolor{ezs}{RGB}{253,189,24}\fill[ezs](118.6,-204.3)--(120.9,-205.2)--(186.2,-35)--(183.9,-34.1)--cycle;
				\definecolor{ezs}{RGB}{253,186,20}\fill[ezs](120.5,-205)--(122.8,-205.9)--(188.1,-35.7)--(185.8,-34.8)--cycle;
				\definecolor{ezs}{RGB}{253,184,17}\fill[ezs](122.4,-205.8)--(124.7,-206.7)--(190,-36.5)--(187.8,-35.6)--cycle;
				\definecolor{ezs}{RGB}{252,183,15}\fill[ezs](124.3,-206.5)--(126.6,-207.4)--(192,-37.2)--(189.7,-36.3)--cycle;
				\definecolor{ezs}{RGB}{252,181,14}\fill[ezs](126.3,-207.2)--(128.5,-208.1)--(193.9,-37.9)--(191.6,-37.1)--cycle;
				\definecolor{ezs}{RGB}{252,179,13}\fill[ezs](128.2,-208)--(130.4,-208.9)--(195.8,-38.7)--(193.5,-37.8)--cycle;
				\definecolor{ezs}{RGB}{252,178,12}\fill[ezs](130.1,-208.7)--(132.4,-209.6)--(197.7,-39.4)--(195.4,-38.5)--cycle;
				\definecolor{ezs}{RGB}{251,178,10}\fill[ezs](132,-209.4)--(134.3,-210.3)--(199.6,-40.1)--(197.3,-39.3)--cycle;
				\definecolor{ezs}{RGB}{251,178,10}\fill[ezs](133.9,-210.2)--(136.2,-211.1)--(201.5,-40.9)--(199.2,-40)--cycle;
				\definecolor{ezs}{RGB}{250,177,11}\fill[ezs](135.8,-210.9)--(138.1,-211.8)--(203.4,-41.6)--(201.1,-40.7)--cycle;
				\definecolor{ezs}{RGB}{250,178,11}\fill[ezs](137.7,-211.6)--(140,-212.5)--(205.3,-42.3)--(203.1,-41.5)--cycle;
				\definecolor{ezs}{RGB}{248,180,10}\fill[ezs](139.6,-212.4)--(141.9,-213.3)--(207.3,-43.1)--(205,-42.2)--cycle;
				\definecolor{ezs}{RGB}{247,181,9}\fill[ezs](141.5,-213.1)--(143.8,-214)--(209.2,-43.8)--(206.9,-42.9)--cycle;
				\definecolor{ezs}{RGB}{246,182,9}\fill[ezs](143.5,-213.8)--(145.7,-214.7)--(211.1,-44.5)--(208.8,-43.7)--cycle;
				\definecolor{ezs}{RGB}{246,182,10}\fill[ezs](145.4,-214.6)--(147.7,-215.5)--(213,-45.3)--(210.7,-44.4)--cycle;
			\end{scope}
			% chữ EDUZY (Times đậm, làm dày nét bằng 5 lớp lệch nhẹ; căn đúng khung chữ logo gốc)
			\foreach \dx/\dy in {-0.7/0,0.7/0,0/-0.3,0/0.3,0/0}
			\node[anchor=center,inner sep=0pt,outer sep=0pt,text=ezText,align=center] at (109.6+\dx,-193.6+\dy)
			{\resizebox{1.01cm}{0.22cm}{\fontencoding{T1}\fontfamily{ptm}\fontseries{b}\selectfont EDUZY}};
		\end{tikzpicture}%
	}%
}
%====================== THIẾT LẬP THÂN ĐỀ ======================
%--- Môi trường tự luận {bt}: ex_test 3.3.4 chưa khai báo sẵn nên tạo ở đây
%    (ex_test đã tự móc \fixshowans và \listenumerate cho bt)
\makeatletter
\@ifundefined{bt}{\newtheorem{bt}{\color{gold1}\namebt}}{}
\makeatother
\renewcommand{\sectionmark}[1]{\markright{\thesection\ #1}}
%--- Khung tiêu đề (logo + tối đa 4 dòng; dòng 3, 4 để trống {} thì bỏ qua)
%  Đề thi:    \khungtieude{THI THỬ THPT QUỐC GIA NĂM 2026}{Môn Toán -- Đề số 21}
%                         {(Đề thi gồm: ...)}{Thời gian làm bài: 90 phút (không kể thời gian phát đề).}
%  Chuyên đề: \khungtieude{CHUYÊN ĐỀ 3 -- DÃY SỐ VÀ CÔNG THỨC TRUY HỒI}
%                         {Bồi dưỡng học sinh giỏi Toán THPT}{...}{}
\newcommand{\khungtieude}[4]{%
	\begin{tcolorbox}[colframe=gold1,colback=gold5!45,boxrule=0.5mm]
		\begin{minipage}[c]{3cm}
			\centerline{\eduzylogo[2.4cm]}
		\end{minipage}\hspace{0.5cm}
		\begin{minipage}[c]{12.5cm}
			\vskip 0.15 cm
			\centerline{\ind{#1}}
			\vskip 0.15 cm
			\centerline{\textbf{#2}}
			\if\relax\detokenize{#3}\relax\else
				\vskip 0.15 cm
				\centerline{\textit{#3}}
			\fi
			\if\relax\detokenize{#4}\relax\else
				\vskip 0.15 cm
				\centerline{\textit{#4}}
			\fi
			\vskip 0.15 cm
		\end{minipage}
	\end{tcolorbox}
}
%====================== BỔ SUNG CHO BÀI HỌC, CHUYÊN ĐỀ ======================
\usepackage[shortlabels]{enumitem}
\usepackage{tabularx}
%--- Danh sách gọn: \begin{enumerate}[gon] ... \end{enumerate}
\SetEnumitemKey{gon}{leftmargin=*,itemsep=1pt,parsep=0pt,topsep=3pt}
%--- Khung lý thuyết: \begin{khung}{} ... \end{khung}; tiêu đề chỉ dùng khi thật ngắn ("Mẹo", "Lỗi thường gặp")
\newtcolorbox{khung}[1]{breakable,enhanced,colframe=gold1,colback=gold5!45,
	boxrule=0.4mm,fonttitle=\bfseries\sffamily,coltitle=white,colbacktitle=gold1,
	attach boxed title to top left={xshift=6mm,yshift=-3mm},
	boxed title style={arc=1mm,boxrule=0.3mm},title={#1}}
%--- Tiêu đề mục nhỏ trong một dạng: \muc{Dấu hiệu}, \muc{Phương pháp}, \muc{Bài tự luyện}
\makeatletter
\newcommand{\muc}[1]{\par\medskip\noindent\indam{#1}\par\nobreak\smallskip\@afterindentfalse\@afterheading}
\makeatother
%--- Ô công thức cần nhớ (căn giữa, viền vàng, vừa khít nội dung)
%     \begin{center}\congthuc{$Q(n)=(n-1)!$}\end{center}
\newtcbox{\congthuc}{on line,colframe=gold1,colback=gold5!45,boxrule=0.5pt,arc=2pt,
	left=8pt,right=8pt,top=4pt,bottom=4pt}
%--- Công tắc lời giải bài tự luyện (ví dụ mẫu luôn dùng \loigiai)
%     \loigiaiTL[{đáp số}]{lời giải đầy đủ}
%     \TLfull -> in lời giải đầy đủ;  \TLdapso -> chỉ in một dòng "Đáp số";  \TLan -> ẩn hoàn toàn
%     (\CheDo là dethi thì mọi lời giải và đáp số đều ẩn)
\makeatletter
\newif\ifTLhien   \TLhientrue
\newif\ifTLdapso  \TLdapsofalse
\newcommand{\TLfull}{\TLhientrue\TLdapsofalse}
\newcommand{\TLdapso}{\TLhientrue\TLdapsotrue}
\newcommand{\TLan}{\TLhienfalse}
\newcommand{\loigiaiTL}[2][]{%
	\ifTLhien
		\ifTLdapso
			\def\@tlarg{#1}%
			\ifx\@tlarg\@empty\else
				\ifdefstring{\kindOPTN}{dethi}{}{\par\noindent{\bfseries\color{gold1}Đáp số:} #1\par}%
			\fi
		\else
			\loigiai{#2}%
		\fi
	\fi}
\makeatother
\TLdapso% mặc định: chỉ in đáp số bài tự luyện
%====================== BẢNG ĐÚNG/SAI (\choiceTFt) ======================
%  Hàng tiêu đề nền nâu đồng chữ trắng (Phát biểu | Đ | S), các hàng tô xen kẽ kem / trắng,
%  đường kẻ vàng kim, nhãn a) b) c) d) màu nâu đồng. Bản có đáp án: cột được chọn có vòng tròn
%  nâu đồng với dấu tích trắng (câu đúng tích ở cột Đúng, câu sai tích ở cột Sai).
%  Tiêu đề ghi đủ "Đúng | Sai": \OPTN{viettat=0} sau \input (cột Đúng/Sai rộng ra, phát biểu hẹp lại).
%  Dùng bảng gốc của ex_test: \def\BangDSgoc{1} trước \input{preamble-eduzy-v2}.
\providecommand{\BangDSgoc}{0}
\ifnum\BangDSgoc=0
\usepackage{colortbl,pifont}
\colorlet{ezDSke}{gold2!60!gold1}% đường kẻ bảng (nét 0,4pt như ex_test để không đổi chiều cao bảng)
\colorlet{ezDSle}{gold5!60}% hàng a), c)
\colorlet{ezDSchan}{white}% hàng b), d)
%--- dấu chọn ở cột Đúng/Sai
%    (bề rộng 0, căn giữa ô: không làm đổi bề rộng cột của ex_test, không báo tràn lề)
\newcommand{\EZdauchon}{\makebox[0pt]{\tikz[baseline=-1ex]\node[circle,fill=gold1,inner sep=0pt,
	minimum size=1.2em,text=white,font=\small]{\ding{51}};}}
\def\tickT{\EZdauchon}\def\tickF{\EZdauchon}
\AtBeginDocument{\sbox{\ccTFTab}{\EZdauchon}}
%--- nhãn a) b) c) d) màu nâu đồng (ở bản lời giải, nhãn câu đúng vẫn đóng khung như ex_test)
\renewcommand{\FalseTF}{\stepcounter{dapan}\textbf{\color{gold1}\DapAnTF\sepTF}\ignorespaces}
%--- trước khi dựng bảng: kẻ màu vàng kim (bề rộng cột giữ nguyên như ex_test); sau bảng trả lại màu đen
\newcount\EZhangTF
\newcommand{\EZbatdauTF}{%
	\global\EZhangTF=0\relax
	\arrayrulecolor{ezDSke}}
\newcommand{\EZketthucTF}{\arrayrulecolor{black}}
%--- hàng tiêu đề nền nâu đồng, chữ trắng
\makeatletter
\newcommand{\EZo@TF}[1]{\centering\bfseries\color{white}#1}
\newcommand{\EZtieudeTF}{%
	\ifcase\cotDS
		\ifnum\cotSTT=1\relax
			\def\rowHeader{\hline\rowcolor{gold1}\EZo@TF{TT}&\EZo@TF{\phatbieu}\tabularnewline}%
		\else
			\def\rowHeader{\hline\rowcolor{gold1}\EZo@TF{\phatbieu}\tabularnewline}%
		\fi
	\or
		\ifnum\cotSTT=1\relax
			\def\rowHeader{\hline\rowcolor{gold1}\EZo@TF{TT}&\EZo@TF{\phatbieu}&\EZo@TF{\Dung}\tabularnewline}%
		\else
			\def\rowHeader{\hline\rowcolor{gold1}\EZo@TF{\phatbieu}&\EZo@TF{\Dung}\tabularnewline}%
		\fi
	\or
		\ifnum\cotSTT=1\relax
			\def\rowHeader{\hline\rowcolor{gold1}\EZo@TF{TT}&\EZo@TF{\phatbieu}&\EZo@TF{\Dung}&\EZo@TF{\Sai}\tabularnewline}%
		\else
			\def\rowHeader{\hline\rowcolor{gold1}\EZo@TF{\phatbieu}&\EZo@TF{\Dung}&\EZo@TF{\Sai}\tabularnewline}%
		\fi
	\fi}
%--- tô xen kẽ từng hàng (gọi ngay sau \hline đầu mỗi hàng)
\newcommand{\EZmauhangTF}{%
	\noalign{\global\advance\EZhangTF\@ne}%
	\ifodd\EZhangTF\rowcolor{ezDSle}\else\rowcolor{ezDSchan}\fi}
%--- gắn vào \choiceTFt / \writeChoiceTFt của ex_test
\xpretocmd{\choiceTFt}{\EZbatdauTF}{}{\typeout{[preamble-eduzy-v2] Không gắn được bảng Đúng/Sai (đầu bảng)}}
\xpatchcmd{\choiceTFt}{\ifnum\dongPB=0\def\rowHeader{}\fi}{\EZtieudeTF\ifnum\dongPB=0\def\rowHeader{}\fi}
	{}{\typeout{[preamble-eduzy-v2] Không gắn được bảng Đúng/Sai (hàng tiêu đề)}}
\xpatchcmd{\writeChoiceTFt}{\hline\xdef\dotEX{\sdotEX}}{\hline\EZmauhangTF\xdef\dotEX{\sdotEX}}
	{}{\typeout{[preamble-eduzy-v2] Không gắn được bảng Đúng/Sai (tô hàng)}}
\xpatchcmd{\writeChoiceTFt}{{\end{longtable}}}{{\end{longtable}\EZketthucTF}}
	{}{\typeout{[preamble-eduzy-v2] Không gắn được bảng Đúng/Sai (cuối longtable)}}
\xpatchcmd{\writeChoiceTFt}{{\end{xtabular}\end{center}\medskip\noindent}}{{\end{xtabular}\end{center}\EZketthucTF\medskip\noindent}}
	{}{\typeout{[preamble-eduzy-v2] Không gắn được bảng Đúng/Sai (cuối xtabular)}}
\fi
%=====================================================================
\renewcommand{\parallel}{\mathrel{/\!\!/}}

%  trong file đề (tools/dan_trang/dan_trang.py tự ghi khi đặt \NeedspaceCau):
%    \def\TachCauDS{1}   câu Đúng/Sai được tách: câu dẫn (kèm hình) ở cuối trang, bảng Phát biểu sang trang sau
%                        (câu dẫn và hình không tách nhau, bảng không tách giữa hai trang)
%    \def\TachCauLG{1}   bản lời giải (\CheDo khác dethi): phần đề của câu cao hơn 1/3 trang được phép tách
%  Lệnh dùng trong file đề (giá trị do tools/dan_trang/dan_trang.py đo ở bản dethi):
%    \NeedspaceCau[<cao câu dẫn + hình>]{<cao cả khối câu>}   ngay trên \begin{ex} (câu đầu phần: trên tiêu đề phần)
%        công tắc tắt: đúng bằng \Needspace{<cao cả khối câu>} như cũ;
%        \TachCauDS=1 và có [..] (câu Đúng/Sai được chọn tách): chỉ giữ câu dẫn + hình;
%        \TachCauLG=1 ở bản lời giải: đòi tối đa 1/3 \textheight.
%    \NeedspaceBang{<cao bảng Phát biểu>}%   ngay trước \choiceTFt: \TachCauDS=1 thì giữ trọn bảng trên một trang;
%        công tắc tắt: không làm gì.
\providecommand{\TachCauDS}{0}
\providecommand{\TachCauLG}{0}
\newlength{\EZnsCau}
\newcommand{\NeedspaceCau}[2][]{%
	\setlength{\EZnsCau}{#2}%
	\ifnum\TachCauDS=1 \if\relax\detokenize{#1}\relax\else\setlength{\EZnsCau}{#1}\fi\fi
	\ifnum\TachCauLG=1 \ifdefstring{\kindOPTN}{dethi}{}{\ifdim\EZnsCau>0.3333\textheight\setlength{\EZnsCau}{0.3333\textheight}\fi}\fi
	\Needspace{\EZnsCau}}
\newcommand{\NeedspaceBang}[1]{\ifnum\TachCauDS=1 \Needspace{#1}\fi}

%--- Thang màu: đậm -> nhạt
\definecolor{gold1}{RGB}{126,96,0}    % nâu đồng đậm  #7E6000
\definecolor{gold2}{RGB}{192,144,3}   % vàng kim      #C09003
\definecolor{gold3}{RGB}{255,215,99}  % vàng nhạt     #FFD763
\definecolor{gold4}{RGB}{255,231,153} % vàng rất nhạt #FFE799
\definecolor{gold5}{RGB}{255,243,205} % vàng kem      #FFF3CD
%--- Ghi đè tông xanh/đỏ cũ
\definecolor{myblue}{RGB}{126,96,0}
\colorlet{blue}{gold1}
\colorlet{cyan}{gold2}
\colorlet{teal}{gold2}

%--- Tên chương mặc định (file đề nên tự khai báo lại)
\providecommand{\tenchuong}{}
%--- Chân trang
\providecommand{\thuonghieu}{TOÁN 12 -- Eduzy.com.vn}
\providecommand{\tacgia}{Tô Gia Bảo -- 0353703783}

%--- Chữ cạnh logo ở đầu trang
\providecommand{\khauhieu}{Eduzy.com.vn}

%====================== THÔNG SỐ ĐẦU TRANG, CHÂN TRANG (v2) ======================
%  Dải nền đậm tràn ngang trang; ô chữ vàng kem có cạnh xiên, kèm 2 vạch xiên chuyển màu.
%  Đầu trang lẻ : logo + \khauhieu (trái)   | ô chữ \tenchuong (phải)
%  Chân trang lẻ: ô chữ \thuonghieu (trái)  | Trang (số) (phải)
%  Trang chẵn soi gương trang lẻ; chân trang chẵn ghi \tacgia.
%--- Màu: chỉ dùng thang vàng kim Eduzy
\colorlet{ezNen}{gold1}% dải nền đậm
\colorlet{ezVach1}{gold3}% vạch xiên sát dải nền (đậm hơn)
\colorlet{ezVach2}{gold4}% vạch xiên sát ô chữ (nhạt hơn)
\colorlet{ezO}{gold5}% ô chữ vàng kem
%--- Kích thước
\newlength{\EZdai}\setlength{\EZdai}{1.1cm}% chiều cao dải nền
\newlength{\EZo}\setlength{\EZo}{0.68cm}% chiều cao ô chữ
\newlength{\EZxien}\setlength{\EZxien}{0.55cm}% độ lệch ngang của cạnh xiên
\newlength{\EZvach}\setlength{\EZvach}{0.75cm}% bề rộng mỗi vạch xiên
\newlength{\EZkhe}\setlength{\EZkhe}{0.12cm}% khe trắng giữa vạch và ô chữ
\newlength{\EZdem}\setlength{\EZdem}{0.4cm}% khoảng đệm từ chữ tới cạnh xiên
\newlength{\EZle}\setlength{\EZle}{1.8cm}% lề trái/phải của chữ (bằng lề thân bài)
\newlength{\EZrongmin}\setlength{\EZrongmin}{4.5cm}% ô chữ rộng tối thiểu
\newlength{\EZrong}% bề rộng chữ trong ô (tính tự động)
%--- Font
\newcommand{\EZfontO}{\fontsize{10.5pt}{12pt}\selectfont\bfseries}% chữ trong ô
\newcommand{\EZfontKH}{\fontfamily{qag}\fontsize{11pt}{12pt}\selectfont\bfseries}% \khauhieu cạnh logo
\newcommand{\EZfontTrang}{\fontsize{11.5pt}{12pt}\selectfont\bfseries}% chữ "Trang" và số trang
%--- Đo bề rộng chữ trong ô
\newcommand{\EZdorong}[1]{%
	\settowidth{\EZrong}{\EZfontO #1}%
	\ifdim\EZrong<\EZrongmin\setlength{\EZrong}{\EZrongmin}\fi}
%--- Logo + chữ trắng \khauhieu (dựng sẵn một lần khi vào \begin{document})
\newsavebox{\EZlogobox}
\AtBeginDocument{\sbox{\EZlogobox}{%
		\raisebox{-0.5\height}{\eduzylogo[1cm]}\hspace{0.15cm}%
		\raisebox{-0.5\height}{\color{white}\EZfontKH\khauhieu}}}
%--- Hai vạch xiên + khe trắng. #1: điểm gốc ở mép dưới, #2: hướng (1 sang phải, -1 sang trái),
%    #3: hướng nghiêng (1 kiểu "/", -1 kiểu "\")
\newcommand{\EZvachxien}[3]{%
	\fill[white] (#1) --++({(#3)*\EZxien},\EZo) --++({(#2)*\EZkhe},0) --++({-(#3)*\EZxien},-\EZo) --cycle;
	\fill[ezVach2] ([xshift=(#2)*\EZkhe]#1) --++({(#3)*\EZxien},\EZo) --++({(#2)*\EZvach},0) --++({-(#3)*\EZxien},-\EZo) --cycle;
	\fill[ezVach1] ([xshift=(#2)*(\EZkhe+\EZvach)]#1) --++({(#3)*\EZxien},\EZo) --++({(#2)*\EZvach},0) --++({-(#3)*\EZxien},-\EZo) --cycle;
	\draw[white,line width=0.6pt] ([xshift=(#2)*(\EZkhe+\EZvach)]#1) --++({(#3)*\EZxien},\EZo);
}
%--- Chữ trong ô (căn giữa theo chiều cao ô, không lệch vì dấu tiếng Việt)
\tikzset{ezchu/.style={inner sep=0pt,text=ezNen,text height=0.72em,text depth=0.22em}}

%====================== LOGO CHÌM GIỮA TRANG ======================
%  Vẽ cùng đầu trang nên nằm dưới chữ của thân bài, có trên mọi trang.
%  Tắt: \logochimtat (từ trang đó trở đi), bật lại: \logochimbat.
%  Đổi cỡ, độ đậm: khai báo trước \input, ví dụ \def\EZchimcao{9cm} \def\EZchimdo{0.1}
\providecommand{\EZchimcao}{10cm}% chiều cao logo chìm
\providecommand{\EZchimdo}{0.12}% độ đậm (0 = trong suốt, 1 = rõ hẳn)
\newif\ifEZchim \EZchimtrue
\newcommand{\logochimtat}{\EZchimfalse}
\newcommand{\logochimbat}{\EZchimtrue}
%  Dựng sẵn logo mờ một lần khi vào \begin{document}; transparency group làm mờ cả logo
%  một lượt nên không lộ vệt ở chỗ các dải màu chồng nhau.
\newsavebox{\EZchimtho}
\newsavebox{\EZchimbox}
\AtBeginDocument{%
	\sbox{\EZchimtho}{\eduzylogo[\EZchimcao]}%
	\sbox{\EZchimbox}{\begin{tikzpicture}
			\begin{scope}[transparency group,opacity=\EZchimdo]
				\node[inner sep=0pt] {\usebox{\EZchimtho}};
			\end{scope}
	\end{tikzpicture}}}
\newcommand{\EZvechim}{%
	\ifEZchim
		\node[inner sep=0pt] at (current page.center) {\usebox{\EZchimbox}};
	\fi}

%====================== ĐẦU TRANG - TRANG LẺ ======================
%  (logo bên trái, ô chữ \tenchuong bên phải)
\newcommand{\myfancyheadOdd}{%
	\EZdorong{\tenchuong}%
	\begin{tikzpicture}[remember picture,overlay]
		\coordinate (T) at (current page.north west);
		\coordinate (P) at (current page.north east);
		%--- logo chìm giữa trang
		\EZvechim
		%--- dải nền
		\fill[ezNen] (T) rectangle ([yshift=-\EZdai]P);
		%--- ô chữ: X là góc dưới bên trái, cạnh trái xiên kiểu "/"
		\coordinate (X) at ([xshift=-\EZle-\EZrong-\EZdem-\EZxien,yshift=-\EZdai]P);
		\fill[ezO] (X) --++(\EZxien,\EZo) -- ([yshift=-\EZdai+\EZo]P) -- ([yshift=-\EZdai]P) --cycle;
		\EZvachxien{X}{-1}{1}
		\node[ezchu,anchor=east] at ([xshift=-\EZle,yshift=-\EZdai+0.5*\EZo]P) {\EZfontO\tenchuong};
		%--- logo + tên
		\node[inner sep=0pt,anchor=west] at ([xshift=\EZle,yshift=-0.5*\EZdai]T) {\usebox{\EZlogobox}};
	\end{tikzpicture}%
}

%====================== ĐẦU TRANG - TRANG CHẴN =====================
%  (soi gương: ô chữ \tenchuong bên trái, logo bên phải)
\newcommand{\myfancyheadEven}{%
	\EZdorong{\tenchuong}%
	\begin{tikzpicture}[remember picture,overlay]
		\coordinate (T) at (current page.north west);
		\coordinate (P) at (current page.north east);
		%--- logo chìm giữa trang
		\EZvechim
		%--- dải nền
		\fill[ezNen] (T) rectangle ([yshift=-\EZdai]P);
		%--- ô chữ: X là góc dưới bên phải, cạnh phải xiên kiểu "\"
		\coordinate (X) at ([xshift=\EZle+\EZrong+\EZdem+\EZxien,yshift=-\EZdai]T);
		\fill[ezO] (X) --++(-\EZxien,\EZo) -- ([yshift=-\EZdai+\EZo]T) -- ([yshift=-\EZdai]T) --cycle;
		\EZvachxien{X}{1}{-1}
		\node[ezchu,anchor=west] at ([xshift=\EZle,yshift=-\EZdai+0.5*\EZo]T) {\EZfontO\tenchuong};
		%--- logo + tên
		\node[inner sep=0pt,anchor=east] at ([xshift=-\EZle,yshift=-0.5*\EZdai]P) {\usebox{\EZlogobox}};
	\end{tikzpicture}%
}

%====================== CHÂN TRANG - TRANG LẺ ======================
%  (ô chữ \thuonghieu bên trái, "Trang (số)" bên phải)
\newcommand{\myfancyfootOdd}{%
	\EZdorong{\thuonghieu}%
	\begin{tikzpicture}[remember picture,overlay]
		\coordinate (T) at (current page.south west);
		\coordinate (P) at (current page.south east);
		%--- dải nền
		\fill[ezNen] (T) rectangle ([yshift=\EZdai]P);
		%--- ô chữ sát mép trên dải: X là góc dưới bên phải, cạnh phải xiên kiểu "/"
		\coordinate (X) at ([xshift=\EZle+\EZrong+\EZdem,yshift=\EZdai-\EZo]T);
		\fill[ezO] ([yshift=\EZdai-\EZo]T) -- (X) --++(\EZxien,\EZo) -- ([yshift=\EZdai]T) --cycle;
		\EZvachxien{X}{1}{1}
		\node[ezchu,anchor=west] at ([xshift=\EZle,yshift=\EZdai-0.5*\EZo]T) {\EZfontO\thuonghieu};
		%--- số trang
		\node[circle,draw=white,line width=0.9pt,text=white,inner sep=0pt,minimum size=0.62cm,anchor=east]
		(so) at ([xshift=-\EZle,yshift=0.5*\EZdai]P) {\EZfontTrang\thepage};
		\node[inner sep=0pt,text=white,anchor=base east] (tr) at ([xshift=-0.15cm]so.base west) {\EZfontTrang Trang};
	\end{tikzpicture}%
}

%====================== CHÂN TRANG - TRANG CHẴN ====================
%  (soi gương: "Trang (số)" bên trái, ô chữ \tacgia bên phải)
\newcommand{\myfancyfootEven}{%
	\EZdorong{\tacgia}%
	\begin{tikzpicture}[remember picture,overlay]
		\coordinate (T) at (current page.south west);
		\coordinate (P) at (current page.south east);
		%--- dải nền
		\fill[ezNen] (T) rectangle ([yshift=\EZdai]P);
		%--- ô chữ sát mép trên dải: X là góc dưới bên trái, cạnh trái xiên kiểu "\"
		\coordinate (X) at ([xshift=-\EZle-\EZrong-\EZdem,yshift=\EZdai-\EZo]P);
		\fill[ezO] ([yshift=\EZdai-\EZo]P) -- (X) --++(-\EZxien,\EZo) -- ([yshift=\EZdai]P) --cycle;
		\EZvachxien{X}{-1}{-1}
		\node[ezchu,anchor=east] at ([xshift=-\EZle,yshift=\EZdai-0.5*\EZo]P) {\EZfontO\tacgia};
		%--- số trang
		%    (vòng tròn ẩn so0 nằm giữa dải để lấy đường chân chữ, giống trang lẻ)
		\node[circle,inner sep=0pt,minimum size=0.62cm] (so0) at ([yshift=0.5*\EZdai]T) {\EZfontTrang\phantom{\thepage}};
		\node[inner sep=0pt,text=white,anchor=base west] (tr) at ($(so0.base -| T)+(\EZle,0)$) {\EZfontTrang Trang};
		\node[circle,draw=white,line width=0.9pt,text=white,inner sep=0pt,minimum size=0.62cm,anchor=base west]
		(so) at ([xshift=0.15cm]tr.base east) {\EZfontTrang\thepage};
	\end{tikzpicture}%
}

%--- Giữ tương thích với các file cũ đang gọi \myfancyhead, \myfancyfoot
\newcommand{\myfancyhead}{\myfancyheadOdd}
\newcommand{\myfancyfoot}{\myfancyfootOdd}

%======================= GẮN VÀO TRANG =============================
\pagestyle{fancy}
\fancyhf{}
\fancyhead[LO]{\myfancyheadOdd}
\fancyhead[RE]{\myfancyheadEven}
\fancyfoot[LO]{\myfancyfootOdd}
\fancyfoot[RE]{\myfancyfootEven}
\renewcommand{\headrulewidth}{0pt}
\renewcommand{\footrulewidth}{0pt}
\fancypagestyle{plain}{%
	\fancyhf{}
	\fancyhead[LO]{\myfancyheadOdd}
	\fancyhead[RE]{\myfancyheadEven}
	\fancyfoot[LO]{\myfancyfootOdd}
	\fancyfoot[RE]{\myfancyfootEven}
	\renewcommand{\headrulewidth}{0pt}
	\renewcommand{\footrulewidth}{0pt}
}
%=================== HẾT PHẦN GỘP TỪ eduzy-de.tex ===================
\newcommand{\indam}[1]{
	{\fontsize{12pt}{9pt}\sffamily \selectfont\bfseries\color{gold1} #1}
}
\newcommand{\ind}[1]{
	{\fontsize{14pt}{9pt}\fontfamily{qag} \selectfont\bfseries\color{gold1} #1}
}
\renewcommand{\FalseEX}{\stepcounter{dapan}\textbf{\selectfont\color{gold1}\Alph{dapan}.}\ignorespaces}
\renewtheorem{ex}{\color{gold1}Câu}
\renewcommand{\baselinestretch}{1.1}
\newcommand{\hoac}[1]{
	\left[\begin{aligned}#1\end{aligned}\right.}
\newcommand{\heva}[1]{
	\left\{\begin{aligned}#1\end{aligned}\right.}
\usepackage{esvect}
\def\vec{\vv}
\def\vt{\vv}
\def\overrightarrow{\vv}
%====================== LOGO EDUZY VẼ BẰNG TIKZ ======================
%  Đường viền dò từ logo gốc; dải màu vàng kim đo trực tiếp trên ảnh gốc
%  (nghiêng 21 độ, tô bằng 80 dải hẹp nên hiển thị giống nhau trên mọi trình xem PDF).
%  Dùng: \eduzylogo          -> cao 2.4cm (mặc định)
%        \eduzylogo[1.5cm]   -> cao 1.5cm
\definecolor{ezText}{RGB}{245,200,46}
\newcommand{\eduzylogo}[1][2.4cm]{%
	\resizebox{!}{#1}{%
		\begin{tikzpicture}[x=0.01cm,y=0.01cm]
			\path[use as bounding box] (0,0) rectangle (224,-224);
			% mũ tốt nghiệp + vương miện
			\begin{scope}
				\clip
				(119.4,-176.1)..controls(118.9,-176.1)and(116.8,-176)..(114.6,-176)
				..controls(101.4,-176)and(85.8,-175.9)..(83.7,-175.7)
				..controls(81.8,-175.6)and(81,-175.6)..(80.2,-175.7)
				..controls(78.8,-175.9)and(77.7,-175.9)..(76.2,-175.7)
				..controls(74.9,-175.5)and(70.8,-175.5)..(64.3,-175.7)
				..controls(59.1,-175.9)and(58.3,-175.7)..(58.6,-173.7)
				..controls(58.6,-172.8)and(59.6,-169.6)..(60.3,-168)
				..controls(60.8,-166.9)and(61.4,-165.8)..(62.3,-165)
				..controls(63.1,-164.2)and(64.3,-163.4)..(65,-163.2)
				..controls(65.5,-163.1)and(72.3,-163)..(76.6,-163.1)
				..controls(77.9,-163.1)and(80.1,-163.1)..(81.4,-163)
				..controls(82.8,-162.9)and(86.3,-162.9)..(89.2,-162.9)
				..controls(93,-162.9)and(94.9,-162.8)..(96,-162.7)
				..controls(97.2,-162.5)and(101,-162.5)..(114,-162.4)
				..controls(123,-162.3)and(134.9,-162.2)..(140.2,-162.1)
				..controls(150.1,-162)and(152.2,-162.1)..(153.6,-162.4)
				..controls(154.4,-162.6)and(157.1,-163.8)..(158.1,-164.4)
				..controls(158.8,-164.8)and(159.7,-165.7)..(159.9,-166.3)
				..controls(160,-166.5)and(160.2,-167.1)..(160.2,-167.6)
				..controls(160.3,-168)and(160.6,-169.4)..(160.9,-170.6)
				..controls(161.2,-171.7)and(161.5,-173.1)..(161.6,-173.6)
				..controls(161.8,-174.4)and(161.8,-174.5)..(161.5,-174.9)
				..controls(161.2,-175.6)and(160.8,-175.9)..(159.9,-176)
				..controls(159.2,-176.1)and(154.1,-176.2)..(151.5,-176.2)
				..controls(148.1,-176.1)and(122.6,-176.1)..(121.6,-176.2)
				..controls(120.8,-176.2)and(119.8,-176.2)..(119.4,-176.1)--cycle
				(130.4,-160)..controls(129,-160)and(119.8,-160)..(110.1,-159.9)
				..controls(100.3,-159.8)and(88.5,-159.7)..(83.9,-159.6)
				..controls(77.8,-159.5)and(75.3,-159.5)..(74.5,-159.6)
				..controls(73.9,-159.7)and(72.5,-159.8)..(71.3,-159.8)
				..controls(69.6,-159.8)and(69,-159.8)..(68.5,-159.6)
				..controls(67.6,-159.3)and(66.8,-158.6)..(66.2,-157.6)
				..controls(66,-157.3)and(65.6,-156.5)..(65.3,-156.1)
				..controls(64.6,-155.1)and(64.5,-154.9)..(63.7,-152.9)
				..controls(63.2,-151.6)and(62.9,-150.9)..(61.8,-149.6)
				..controls(61.5,-149.2)and(61.1,-148.5)..(60.9,-147.9)
				..controls(60.7,-147.4)and(60.3,-146.6)..(60.1,-146.1)
				..controls(59.9,-145.5)and(59.5,-144.9)..(59,-144.3)
				..controls(58.5,-143.7)and(58.2,-143.2)..(57.9,-142.5)
				..controls(57.7,-141.9)and(57.3,-141.1)..(57,-140.7)
				..controls(56.6,-140.3)and(56.1,-139.3)..(55.8,-138.5)
				..controls(55.5,-137.8)and(55.1,-136.8)..(54.9,-136.5)
				..controls(54.7,-136.1)and(54.3,-135.2)..(53.9,-134.5)
				..controls(53.6,-133.8)and(53,-132.7)..(52.6,-132)
				..controls(51.6,-130.5)and(51.4,-130.2)..(51.1,-129)
				..controls(50.6,-127.5)and(50.6,-127.4)..(49.7,-126.3)
				..controls(49.2,-125.7)and(48.6,-124.8)..(48.3,-124.2)
				..controls(48,-123.6)and(47.5,-122.7)..(47.3,-122.2)
				..controls(47,-121.8)and(46.5,-120.7)..(46.1,-119.8)
				..controls(45.8,-118.9)and(45.2,-117.8)..(44.9,-117.2)
				..controls(44.2,-116.1)and(43.3,-114.4)..(42.3,-112.2)
				..controls(41.8,-111.3)and(41.3,-110.4)..(41.1,-110.1)
				..controls(40.5,-109.4)and(40.3,-109)..(40,-107.8)
				..controls(39.8,-107)and(39.6,-106.5)..(39.2,-106)
				..controls(38.9,-105.6)and(38.4,-104.7)..(38,-104)
				..controls(37.7,-103.4)and(37.2,-102.5)..(37,-102.1)
				..controls(36.7,-101.7)and(36.3,-100.8)..(36,-100.1)
				..controls(35.6,-98.9)and(35.3,-98.4)..(34.6,-97.6)
				..controls(34.2,-97.1)and(34.2,-96.3)..(34.6,-96)
				..controls(35,-95.7)and(35.5,-95.9)..(36.2,-96.7)
				..controls(36.5,-97)and(37.2,-97.5)..(37.6,-97.7)..controls(38,-98)and(38.6,-98.4)..(39.1,-98.9)
				..controls(39.6,-99.4)and(40.3,-100)..(40.6,-100.3)
				..controls(41.3,-100.8)and(42.4,-101.8)..(43.7,-102.9)
				..controls(44,-103.2)and(44.6,-103.7)..(45,-103.9)
				..controls(45.4,-104.2)and(45.9,-104.6)..(46.1,-104.8)
				..controls(47.1,-105.9)and(48.8,-107.1)..(50.5,-107.9)
				..controls(51.8,-108.6)and(52.3,-109)..(53.4,-110)
				..controls(54.8,-111.2)and(55.6,-111.6)..(57,-112.1)
				..controls(57.5,-112.3)and(58.5,-112.8)..(59.4,-113.3)
				..controls(60.2,-113.7)and(61.3,-114.3)..(61.9,-114.5)
				..controls(62.4,-114.7)and(63.3,-115.1)..(63.8,-115.3)
				..controls(64.4,-115.6)and(65.4,-115.9)..(66.2,-116.1)
				..controls(67,-116.3)and(68.4,-116.6)..(69.2,-116.9)
				..controls(72.3,-117.9)and(72.8,-118)..(78,-118.1)
				..controls(84,-118.3)and(86.1,-118.1)..(88.1,-117.1)
				..controls(89.4,-116.5)and(93,-115)..(93.7,-114.9)
				..controls(94.8,-114.6)and(95.4,-114.2)..(95.9,-113.4)
				..controls(96.2,-113)and(96.7,-112.5)..(97.1,-112.2)
				..controls(98,-111.6)and(101.3,-108.4)..(102.1,-107.2)
				..controls(102.4,-106.8)and(102.9,-106.2)..(103.2,-105.8)
				..controls(103.5,-105.5)and(104,-104.8)..(104.2,-104.2)
				..controls(104.5,-103.6)and(105,-102.8)..(105.3,-102.4)
				..controls(105.6,-102)and(106.1,-101.2)..(106.5,-100.6)
				..controls(107.1,-99.5)and(108.6,-97.5)..(109.5,-96.5)
				..controls(109.9,-96.1)and(110.1,-96)..(110.4,-96)
				..controls(111.1,-96)and(111.6,-96.5)..(111.8,-97.6)
				..controls(112,-98.6)and(112.3,-99.3)..(112.8,-99.9)
				..controls(113.7,-100.9)and(115.2,-103.1)..(115.7,-104.2)
				..controls(116.1,-105)and(116.5,-105.5)..(116.9,-106)
				..controls(117.3,-106.4)and(117.8,-106.9)..(118,-107.3)
				..controls(118.2,-107.6)and(118.6,-108)..(119,-108.3)
				..controls(119.6,-108.8)and(121.1,-110.3)..(121.8,-111.2)
				..controls(122.1,-111.7)and(122.5,-112)..(123.4,-112.4)
				..controls(124.8,-113.2)and(125.8,-113.8)..(127,-114.6)
				..controls(127.5,-114.9)and(128.1,-115.3)..(128.2,-115.4)
				..controls(128.4,-115.5)and(129.3,-115.8)..(130.2,-116)
				..controls(131.1,-116.2)and(132.5,-116.6)..(133.3,-116.9)
				..controls(135.7,-117.9)and(136.1,-117.9)..(142.5,-117.8)
				..controls(145.8,-117.8)and(148.4,-117.7)..(148.8,-117.6)
				..controls(149.2,-117.6)and(150.2,-117.2)..(151,-116.8)
				..controls(152,-116.4)and(152.8,-116)..(153.6,-115.9)
				..controls(154.5,-115.7)and(155.2,-115.4)..(156.3,-114.9)
				..controls(157.1,-114.5)and(158.2,-114)..(158.7,-113.8)
				..controls(159.1,-113.7)and(159.8,-113.3)..(160.2,-113)
				..controls(160.6,-112.7)and(161.3,-112.4)..(161.7,-112.2)
				..controls(162.7,-111.9)and(163.2,-111.7)..(164.1,-111.1)
				..controls(164.5,-110.8)and(165.3,-110.4)..(165.9,-110)
				..controls(166.5,-109.7)and(167.3,-109.2)..(167.6,-108.9)
				..controls(167.9,-108.6)and(168.6,-108.2)..(169,-108)
				..controls(169.4,-107.8)and(170.2,-107.3)..(170.7,-106.9)
				..controls(171.2,-106.5)and(172.1,-105.8)..(172.7,-105.4)
				..controls(173.2,-105)and(174.1,-104.4)..(174.5,-104)
				..controls(174.9,-103.7)and(175.7,-103.1)..(176.2,-102.6)
				..controls(176.7,-102.2)and(177.4,-101.5)..(177.9,-101.2)
				..controls(178.3,-100.8)and(178.8,-100.4)..(179,-100.2)
				..controls(179.2,-100)and(179.6,-99.6)..(180,-99.3)
				..controls(180.4,-99.1)and(180.9,-98.6)..(181.2,-98.4)
				..controls(181.5,-98.1)and(182.1,-97.7)..(182.5,-97.5)
				..controls(183.2,-97.2)and(183.5,-96.9)..(184.3,-95.9)
				..controls(184.6,-95.6)and(184.6,-95.6)..(184.8,-95.7)
				..controls(185,-95.8)and(185.1,-96)..(185.1,-96.2)
				..controls(185.1,-96.5)and(185,-96.7)..(184.6,-97.2)
				..controls(184.1,-97.8)and(184,-97.9)..(183.9,-98.8)
				..controls(183.7,-100.3)and(183.3,-101.3)..(182.9,-101.8)
				..controls(182.6,-102.1)and(182.3,-102.6)..(182.2,-102.9)
				..controls(182,-103.3)and(181.6,-104)..(181.2,-104.6)
				..controls(180.9,-105.2)and(180.5,-106.1)..(180.4,-106.6)
				..controls(179.9,-108)and(179.7,-108.5)..(179,-109.4)
				..controls(178.6,-109.8)and(178.2,-110.5)..(177.9,-111)
				..controls(177.7,-111.5)and(177.3,-112.3)..(177,-112.7)
				..controls(176.6,-113.2)and(176.4,-113.7)..(176.3,-114.2)
				..controls(176.2,-114.7)and(175.9,-115.2)..(175.4,-116)
				..controls(175,-116.6)and(174.4,-117.7)..(174,-118.5)
				..controls(173.7,-119.2)and(173.3,-120.1)..(173,-120.5)
				..controls(172.8,-120.9)and(172.4,-121.8)..(172.1,-122.5)
				..controls(171.8,-123.2)and(171.3,-124.1)..(171,-124.6)
				..controls(170.8,-125.1)and(170.3,-126)..(169.9,-126.6)
				..controls(169.6,-127.3)and(169.2,-128.1)..(169,-128.3)
				..controls(168.8,-128.6)and(168.5,-129.3)..(168.3,-129.8)
				..controls(168.1,-130.3)and(167.8,-131.1)..(167.5,-131.6)
				..controls(166.9,-132.5)and(164.2,-137.9)..(164.2,-138.1)
				..controls(164.2,-138.1)and(163.9,-138.8)..(163.4,-139.7)
				..controls(162.9,-140.5)and(162.4,-141.5)..(162.3,-141.9)
				..controls(161.9,-143)and(161.6,-143.4)..(160.9,-144.1)
				..controls(160.4,-144.6)and(160.3,-144.8)..(160.1,-145.4)
				..controls(160,-145.8)and(159.6,-146.8)..(159.2,-147.8)
				..controls(158.7,-148.8)and(158.1,-150)..(157.9,-150.6)
				..controls(157.6,-151.1)and(157.2,-151.9)..(156.9,-152.4)
				..controls(156.5,-152.8)and(156.1,-153.7)..(155.8,-154.4)
				..controls(155.6,-155)and(155.2,-155.9)..(154.9,-156.3)
				..controls(154.7,-156.7)and(154.4,-157.4)..(154.2,-157.9)
				..controls(154.1,-158.4)and(153.8,-159)..(153.5,-159.3)
				..controls(152.8,-160.1)and(153,-160.1)..(142.4,-160.1)
				..controls(137.2,-160.1)and(131.8,-160.1)..(130.4,-160)--cycle
				(141.5,-103.7)..controls(140.8,-103)and(139.8,-102.1)..(139.2,-101.6)
				..controls(138.5,-101.2)and(137.9,-100.6)..(137.8,-100.4)
				..controls(137,-99.5)and(134,-96.5)..(133.2,-96)
				..controls(132.7,-95.7)and(131.8,-94.9)..(131.1,-94.3)
				..controls(129.4,-92.7)and(127.9,-91.6)..(126.6,-91.1)
				..controls(125.2,-90.5)and(124.1,-90)..(122.9,-89.1)
				..controls(122.1,-88.6)and(121.3,-88.2)..(119.9,-87.7)
				..controls(118.8,-87.3)and(117.5,-86.8)..(116.9,-86.5)
				..controls(114.3,-85.5)and(113.7,-85.4)..(109.7,-85.5)
				..controls(105.3,-85.7)and(103.9,-85.9)..(101.6,-87)
				..controls(100.7,-87.4)and(99.4,-87.8)..(97.9,-88)
				..controls(97,-88.2)and(96.6,-88.3)..(95.8,-88.9)
				..controls(95.4,-89.2)and(94.7,-89.6)..(94.2,-89.8)
				..controls(92.9,-90.3)and(92.3,-90.7)..(90.9,-92.1)
				..controls(90.2,-92.7)and(89.4,-93.3)..(89,-93.5)..controls(87,-94.5)and(85.2,-96)..(83.6,-97.9)
				..controls(83.3,-98.2)and(82.5,-98.9)..(81.8,-99.4)
				..controls(80.7,-100.1)and(80.5,-100.3)..(80,-101.1)
				..controls(78.7,-103.1)and(77.2,-103.9)..(76.4,-103)
				..controls(75.7,-102.2)and(75.7,-103.2)..(75.7,-83.6)
				..controls(75.6,-73.6)and(75.6,-64.8)..(75.5,-64)
				..controls(75.5,-63.2)and(75.5,-61.4)..(75.6,-59.8)--(75.7,-57.1)--(76,-56.8)
				..controls(76.8,-56.2)and(78,-56.3)..(79.8,-57.2)
				..controls(80.4,-57.5)and(81.3,-57.8)..(81.8,-57.9)
				..controls(82.3,-58.1)and(83,-58.4)..(83.6,-58.8)..controls(84.3,-59.2)and(85,-59.5)..(86,-59.8)
				..controls(86.8,-60.1)and(88,-60.6)..(88.7,-60.9)
				..controls(89.3,-61.2)and(90.6,-61.7)..(91.4,-62)
				..controls(92.2,-62.3)and(93.7,-62.9)..(94.6,-63.3)
				..controls(95.6,-63.7)and(96.8,-64.1)..(97.4,-64.3)
				..controls(98,-64.5)and(99,-64.9)..(99.7,-65.2)
				..controls(100.4,-65.6)and(101.5,-66)..(102.3,-66.2)
				..controls(104.2,-66.8)and(105,-67.1)..(106.2,-67.8)
				..controls(107.9,-68.7)and(110.2,-69.4)..(111.3,-69.2)
				..controls(112.5,-69)and(115,-68.2)..(115.6,-67.8)
				..controls(117.3,-66.7)and(117.6,-66.6)..(119.5,-66.2)
				..controls(121.9,-65.8)and(122.1,-65.7)..(123,-65.1)
				..controls(124.1,-64.4)and(124.4,-64.3)..(126,-63.9)
				..controls(127,-63.7)and(127.7,-63.4)..(128.4,-63.1)
				..controls(129,-62.8)and(129.7,-62.5)..(130.6,-62.4)
				..controls(132.1,-62.1)and(134,-61.5)..(135.2,-61)
				..controls(135.7,-60.8)and(136.7,-60.5)..(137.3,-60.3)
				..controls(138.4,-60)and(139.3,-59.7)..(141.1,-59)
				..controls(141.6,-58.7)and(142,-58.6)..(142.5,-58.6)
				..controls(143.2,-58.6)and(143.3,-58.7)..(143.6,-59)
				..controls(144,-59.5)and(144.1,-59.9)..(144,-60.8)
				..controls(144,-61.2)and(144,-71.1)..(144,-82.8)--(143.9,-104.1)--(143.7,-104.5)
				..controls(143.5,-104.7)and(143.3,-104.8)..(143,-104.9)
				..controls(142.6,-104.9)and(142.5,-104.8)..(141.5,-103.7)--cycle
				(27.8,-91.4)..controls(27.1,-90.9)and(26.3,-90)..(26.1,-89.5)
				..controls(26,-89.2)and(25.8,-88.6)..(25.7,-87.9)
				..controls(25.4,-86.2)and(25.6,-84.8)..(26.5,-83.5)
				..controls(27.4,-82.3)and(28.3,-82)..(29.4,-82.3)
				..controls(31.1,-82.8)and(32.1,-84.9)..(32,-87.4)
				..controls(31.9,-89.3)and(31.2,-90.6)..(29.9,-91.4)
				..controls(29.1,-91.9)and(28.6,-91.9)..(27.8,-91.4)--cycle
				(190.1,-91.3)..controls(189.7,-91.2)and(189.3,-90.9)..(189.1,-90.7)
				..controls(188.3,-89.8)and(187.6,-87.7)..(187.6,-86.2)
				..controls(187.6,-85.5)and(187.7,-85.2)..(188.2,-84.4)
				..controls(189.1,-82.6)and(190.1,-81.8)..(191.3,-82.1)
				..controls(192,-82.3)and(193.5,-83.7)..(193.8,-84.4)
				..controls(194.1,-85.1)and(194.2,-87.1)..(194,-88.3)
				..controls(193.8,-89.8)and(193.3,-90.5)..(192.3,-91.3)
				..controls(191.7,-91.7)and(190.9,-91.7)..(190.1,-91.3)--cycle
				(173.7,-83.4)..controls(172.2,-82.5)and(172,-82.5)..(170.6,-82.4)
				..controls(169.7,-82.4)and(168.8,-82.4)..(168.3,-82.5)
				..controls(167.2,-82.7)and(166.9,-82.7)..(166.6,-82.2)
				..controls(166.3,-81.8)and(166.3,-81)..(166.8,-80)
				..controls(167.2,-78.9)and(167.4,-77.4)..(167.7,-74.1)
				..controls(167.8,-72.8)and(167.9,-71.8)..(167.9,-71.7)
				..controls(168,-71.7)and(168,-70.2)..(168,-68.4)..controls(168,-65.4)and(168,-64.9)..(167.7,-63)
				..controls(167.3,-60.5)and(167.1,-59.9)..(166.4,-58.4)
				..controls(166.1,-57.7)and(165.8,-56.9)..(165.7,-56.6)
				..controls(165.7,-56.2)and(165.4,-55.6)..(165.2,-55.2)
				..controls(164.6,-53.8)and(161.9,-50.7)..(161,-50.3)
				..controls(159.8,-49.7)and(157.9,-50)..(155.2,-51.2)
				..controls(154.7,-51.4)and(153.4,-51.8)..(152.2,-52.1)
				..controls(150.8,-52.4)and(149.4,-52.9)..(148.6,-53.2)
				..controls(147.8,-53.6)and(146.7,-54)..(146,-54.2)
				..controls(145.4,-54.4)and(144.5,-54.7)..(144,-54.9)
				..controls(143.6,-55.2)and(142.7,-55.5)..(141.8,-55.7)
				..controls(141,-55.9)and(139.6,-56.3)..(138.7,-56.6)
				..controls(137.8,-57)and(136.3,-57.4)..(135.3,-57.6)
				..controls(133.8,-57.9)and(133.2,-58.1)..(132.1,-58.7)
				..controls(131.1,-59.2)and(130.3,-59.4)..(129.4,-59.6)
				..controls(128.7,-59.8)and(127.8,-60)..(127.4,-60.1)
				..controls(127,-60.3)and(126,-60.5)..(125.1,-60.7)
				..controls(123.2,-61.1)and(121.7,-61.5)..(120.2,-62.1)
				..controls(119.7,-62.3)and(118.1,-62.8)..(116.7,-63.2)
				..controls(114.6,-63.9)and(114,-64.1)..(113,-64.7)
				..controls(111.9,-65.4)and(111.8,-65.4)..(111,-65.5)
				..controls(109.9,-65.5)and(109.3,-65.4)..(107.7,-64.7)
				..controls(107,-64.4)and(106,-64)..(105.3,-63.9)
				..controls(103.5,-63.5)and(102.2,-63)..(101.2,-62.5)
				..controls(100.1,-61.8)and(99.3,-61.5)..(97.8,-61.1)
				..controls(97.2,-61)and(96.2,-60.7)..(95.7,-60.4)
				..controls(95.1,-60.1)and(94.2,-59.8)..(93.5,-59.6)
				..controls(92.9,-59.4)and(91.7,-59)..(91,-58.7)..controls(90.2,-58.4)and(88.9,-58)..(88,-57.7)
				..controls(87.2,-57.5)and(86.1,-57.1)..(85.5,-56.8)
				..controls(85,-56.6)and(84.1,-56.2)..(83.5,-56)..controls(82.9,-55.7)and(81.8,-55.2)..(81,-54.9)
				..controls(80.1,-54.5)and(79.1,-54.2)..(78.7,-54.1)
				..controls(77.9,-53.8)and(77.3,-53.6)..(76.2,-52.9)
				..controls(75.7,-52.6)and(75.2,-52.4)..(74.1,-52.2)
				..controls(72.5,-51.9)and(72.4,-51.8)..(71.3,-51.1)
				..controls(69.9,-50.1)and(69.4,-49.9)..(67.9,-49.8)
				..controls(67,-49.7)and(66.6,-49.6)..(66.3,-49.4)
				..controls(66.1,-49.3)and(65.7,-49)..(65.4,-48.8)
				..controls(65.1,-48.7)and(64.2,-48.4)..(63.6,-48.2)
				..controls(62.9,-48)and(61.7,-47.6)..(60.9,-47.2)
				..controls(60.1,-46.8)and(58.9,-46.4)..(58.2,-46.2)
				..controls(57.6,-46)and(56.8,-45.7)..(56.5,-45.5)
				..controls(54.9,-44.6)and(54.6,-44.5)..(53.2,-44.2)
				..controls(51.7,-43.8)and(51,-43.6)..(49.9,-42.9)
				..controls(49.5,-42.6)and(48.9,-42.4)..(47.5,-42)
				..controls(45.7,-41.6)and(45.6,-41.5)..(44.3,-40.6)
				..controls(43.5,-40.1)and(42.8,-39.6)..(42.7,-39.4)
				..controls(42.4,-38.7)and(43,-38.1)..(44.5,-37.6)..controls(44.9,-37.5)and(46,-37)..(46.9,-36.6)
				..controls(47.9,-36.1)and(48.8,-35.8)..(49.6,-35.6)
				..controls(50.5,-35.4)and(51.3,-35.2)..(52.4,-34.6)
				..controls(53.3,-34.1)and(54.5,-33.7)..(54.9,-33.5)
				..controls(55.4,-33.4)and(56.4,-33)..(57.2,-32.6)
				..controls(58.6,-32)and(60.1,-31.5)..(61.6,-31.2)
				..controls(62.3,-31.1)and(62.7,-31)..(63.2,-30.7)
				..controls(64.1,-30)and(64.3,-29.9)..(65.3,-29.7)
				..controls(65.7,-29.6)and(66.5,-29.2)..(67.1,-28.9)
				..controls(67.6,-28.6)and(68.8,-28.2)..(69.6,-27.9)
				..controls(70.5,-27.7)and(71.6,-27.3)..(72.1,-27)
				..controls(72.6,-26.7)and(73.8,-26.2)..(74.8,-25.9)
				..controls(75.7,-25.5)and(76.7,-25.1)..(77.1,-25)
				..controls(77.4,-24.8)and(78.5,-24.4)..(79.6,-24.1)
				..controls(80.6,-23.8)and(82,-23.2)..(82.8,-22.9)
				..controls(83.5,-22.5)and(84.8,-22)..(85.8,-21.8)
				..controls(86.7,-21.5)and(87.8,-21.1)..(88.1,-20.9)
				..controls(88.4,-20.7)and(89.3,-20.3)..(90.1,-20.1)
				..controls(90.8,-19.8)and(91.9,-19.4)..(92.4,-19.1)
				..controls(93.5,-18.5)and(94.4,-18.2)..(95.4,-18.1)
				..controls(95.7,-18)and(96.4,-17.8)..(96.9,-17.5)
				..controls(97.5,-17.3)and(98.4,-16.9)..(99,-16.6)
				..controls(99.7,-16.3)and(100.9,-15.9)..(101.9,-15.6)
				..controls(102.9,-15.3)and(104.4,-14.7)..(105.3,-14.3)
				..controls(108.1,-13.2)and(108.5,-13.1)..(110.1,-13.1)
				..controls(111.4,-13.1)and(111.5,-13.1)..(112.4,-13.5)
				..controls(113.2,-14)and(115,-14.8)..(116.6,-15.4)
				..controls(116.9,-15.5)and(117.9,-15.8)..(118.8,-16.1)
				..controls(120.9,-16.6)and(121.5,-16.9)..(122.6,-17.9)
				..controls(123,-18.3)and(123.7,-18.8)..(123.9,-18.9)
				..controls(124.6,-19.4)and(126,-19.8)..(127.9,-20.2)
				..controls(129.6,-20.5)and(129.8,-20.6)..(131.1,-21.4)
				..controls(131.5,-21.7)and(132.4,-22.1)..(133.1,-22.3)
				..controls(136.1,-23.2)and(138.3,-24.1)..(140.2,-25)
				..controls(141,-25.5)and(141.5,-25.6)..(142.5,-25.8)
				..controls(145.1,-26.3)and(146,-26.6)..(149.4,-28.1)
				..controls(151.3,-28.9)and(153.5,-29.8)..(154.3,-30.1)
				..controls(155.1,-30.4)and(156.2,-30.8)..(156.6,-31)
				..controls(157.1,-31.2)and(158.2,-31.6)..(159.1,-31.8)
				..controls(160.6,-32.2)and(160.9,-32.3)..(161.7,-32.9)
				..controls(162.4,-33.3)and(162.9,-33.6)..(163.7,-33.8)
				..controls(164.3,-34)and(164.8,-34.1)..(164.9,-34.1)
				..controls(165,-34.1)and(165.6,-34.4)..(166.4,-34.8)
				..controls(167.3,-35.2)and(168.1,-35.4)..(169,-35.6)
				..controls(170,-35.8)and(170.5,-36)..(171.7,-36.6)
				..controls(172.4,-37)and(173.3,-37.4)..(173.6,-37.4)
				..controls(175.6,-37.7)and(176.2,-38)..(176.7,-39.1)
				..controls(176.9,-39.6)and(176.9,-39.7)..(176.7,-40.4)
				..controls(176.4,-41.6)and(175.8,-42.7)..(175.2,-43.4)
				..controls(174.5,-44.1)and(172.6,-45.5)..(171.7,-45.8)
				..controls(170.6,-46.2)and(169.1,-46)..(168.2,-45.4)
				..controls(167.9,-45.2)and(167.3,-44.8)..(166.9,-44.4)
				..controls(166.4,-44.1)and(165.8,-43.6)..(165.4,-43.2)
				..controls(165.1,-42.9)and(164.4,-42.3)..(163.9,-41.9)
				..controls(163.3,-41.5)and(162.7,-40.9)..(162.4,-40.5)
				..controls(161.4,-39.2)and(159.9,-38)..(158.8,-37.8)
				..controls(158.4,-37.8)and(158.3,-37.7)..(158.2,-37.2)
				..controls(158,-36.8)and(157.4,-36.1)..(157.1,-36.1)
				..controls(157,-36.1)and(156.7,-36.1)..(156.5,-36)
				..controls(156.1,-35.9)and(155.9,-35.7)..(155.7,-35.2)
				..controls(155.3,-34.6)and(155.1,-34.4)..(154,-34.1)
				..controls(153.6,-34)and(153.2,-33.7)..(152.9,-33.5)
				..controls(152.4,-32.9)and(151.6,-32.4)..(151.3,-32.4)
				..controls(150.9,-32.4)and(150.4,-32.6)..(150.4,-32.9)
				..controls(150.4,-33.2)and(150.8,-33.4)..(151.4,-33.6)
				..controls(152.1,-33.7)and(152.8,-34.1)..(153.5,-34.9)
				..controls(154.1,-35.5)and(154.3,-35.6)..(155,-35.8)
				..controls(155.7,-36)and(155.7,-36)..(156,-36.6)
				..controls(156.3,-37.3)and(156.7,-37.7)..(157.5,-37.8)
				..controls(158.1,-37.9)and(158.2,-38)..(158.4,-38.8)
				..controls(158.5,-39.4)and(158.6,-39.6)..(159.6,-40.4)
				..controls(160.2,-41)and(160.9,-41.6)..(161.3,-41.8)
				..controls(161.7,-42)and(161.9,-42.3)..(162.2,-42.8)
				..controls(162.4,-43.2)and(162.8,-43.8)..(163.1,-44.1)
				..controls(163.9,-44.8)and(165.1,-46.7)..(165.7,-48)
				..controls(166.2,-49.3)and(166.4,-49.7)..(167.2,-50.5)
				..controls(167.5,-50.8)and(167.9,-51.4)..(168.1,-51.8)
				..controls(168.3,-52.2)and(168.8,-53)..(169.1,-53.6)
				..controls(169.6,-54.3)and(169.8,-54.9)..(169.9,-55.6)
				..controls(170.4,-57.8)and(170.5,-58.1)..(171.2,-59.5)
				..controls(171.6,-60.3)and(172,-61.5)..(172.2,-62.1)
				..controls(172.3,-62.8)and(172.7,-63.7)..(172.9,-64.1)
				..controls(173.2,-64.6)and(173.4,-65.3)..(173.4,-65.6)
				..controls(173.8,-68.8)and(174.2,-70.6)..(174.9,-72.6)
				..controls(175.5,-74.5)and(175.9,-77.1)..(176.1,-80.2)
				..controls(176.3,-83.4)and(175.5,-84.4)..(173.7,-83.4)--cycle;
				\definecolor{ezs}{RGB}{247,179,26}\fill[ezs](-5.7,-156.6)--(-3.4,-157.5)--(61.9,12.7)--(59.7,13.6)--cycle;
				\definecolor{ezs}{RGB}{247,179,22}\fill[ezs](-3.8,-157.3)--(-1.5,-158.2)--(63.9,12)--(61.6,12.9)--cycle;
				\definecolor{ezs}{RGB}{248,179,19}\fill[ezs](-1.8,-158.1)--(0.4,-158.9)--(65.8,11.2)--(63.5,12.1)--cycle;
				\definecolor{ezs}{RGB}{248,178,16}\fill[ezs](0.1,-158.8)--(2.4,-159.7)--(67.7,10.5)--(65.4,11.4)--cycle;
				\definecolor{ezs}{RGB}{249,178,15}\fill[ezs](2,-159.5)--(4.3,-160.4)--(69.6,9.8)--(67.3,10.7)--cycle;
				\definecolor{ezs}{RGB}{251,177,17}\fill[ezs](3.9,-160.3)--(6.2,-161.1)--(71.5,9)--(69.2,9.9)--cycle;
				\definecolor{ezs}{RGB}{252,179,18}\fill[ezs](5.8,-161)--(8.1,-161.9)--(73.4,8.3)--(71.1,9.2)--cycle;
				\definecolor{ezs}{RGB}{252,182,17}\fill[ezs](7.7,-161.7)--(10,-162.6)--(75.3,7.6)--(73,8.5)--cycle;
				\definecolor{ezs}{RGB}{252,185,17}\fill[ezs](9.6,-162.5)--(11.9,-163.4)--(77.2,6.8)--(75,7.7)--cycle;
				\definecolor{ezs}{RGB}{251,189,18}\fill[ezs](11.5,-163.2)--(13.8,-164.1)--(79.2,6.1)--(76.9,7)--cycle;
				\definecolor{ezs}{RGB}{251,191,22}\fill[ezs](13.4,-163.9)--(15.7,-164.8)--(81.1,5.4)--(78.8,6.2)--cycle;
				\definecolor{ezs}{RGB}{252,193,25}\fill[ezs](15.4,-164.7)--(17.6,-165.6)--(83,4.6)--(80.7,5.5)--cycle;
				\definecolor{ezs}{RGB}{253,195,27}\fill[ezs](17.3,-165.4)--(19.6,-166.3)--(84.9,3.9)--(82.6,4.8)--cycle;
				\definecolor{ezs}{RGB}{253,196,31}\fill[ezs](19.2,-166.1)--(21.5,-167)--(86.8,3.2)--(84.5,4)--cycle;
				\definecolor{ezs}{RGB}{254,199,35}\fill[ezs](21.1,-166.9)--(23.4,-167.8)--(88.7,2.4)--(86.4,3.3)--cycle;
				\definecolor{ezs}{RGB}{253,201,38}\fill[ezs](23,-167.6)--(25.3,-168.5)--(90.6,1.7)--(88.3,2.6)--cycle;
				\definecolor{ezs}{RGB}{253,203,40}\fill[ezs](24.9,-168.3)--(27.2,-169.2)--(92.5,1)--(90.3,1.8)--cycle;
				\definecolor{ezs}{RGB}{253,205,42}\fill[ezs](26.8,-169.1)--(29.1,-170)--(94.4,0.2)--(92.2,1.1)--cycle;
				\definecolor{ezs}{RGB}{253,207,45}\fill[ezs](28.7,-169.8)--(31,-170.7)--(96.4,-0.5)--(94.1,0.4)--cycle;
				\definecolor{ezs}{RGB}{253,210,49}\fill[ezs](30.7,-170.5)--(32.9,-171.4)--(98.3,-1.2)--(96,-0.4)--cycle;
				\definecolor{ezs}{RGB}{253,212,54}\fill[ezs](32.6,-171.3)--(34.9,-172.2)--(100.2,-2)--(97.9,-1.1)--cycle;
				\definecolor{ezs}{RGB}{253,214,59}\fill[ezs](34.5,-172)--(36.8,-172.9)--(102.1,-2.7)--(99.8,-1.8)--cycle;
				\definecolor{ezs}{RGB}{253,216,63}\fill[ezs](36.4,-172.7)--(38.7,-173.6)--(104,-3.4)--(101.7,-2.6)--cycle;
				\definecolor{ezs}{RGB}{253,218,65}\fill[ezs](38.3,-173.5)--(40.6,-174.4)--(105.9,-4.2)--(103.6,-3.3)--cycle;
				\definecolor{ezs}{RGB}{254,220,66}\fill[ezs](40.2,-174.2)--(42.5,-175.1)--(107.8,-4.9)--(105.5,-4)--cycle;
				\definecolor{ezs}{RGB}{254,222,68}\fill[ezs](42.1,-175)--(44.4,-175.8)--(109.7,-5.6)--(107.5,-4.8)--cycle;
				\definecolor{ezs}{RGB}{254,224,71}\fill[ezs](44,-175.7)--(46.3,-176.6)--(111.7,-6.4)--(109.4,-5.5)--cycle;
				\definecolor{ezs}{RGB}{254,226,75}\fill[ezs](46,-176.4)--(48.2,-177.3)--(113.6,-7.1)--(111.3,-6.2)--cycle;
				\definecolor{ezs}{RGB}{254,228,79}\fill[ezs](47.9,-177.2)--(50.1,-178)--(115.5,-7.8)--(113.2,-7)--cycle;
				\definecolor{ezs}{RGB}{254,230,82}\fill[ezs](49.8,-177.9)--(52.1,-178.8)--(117.4,-8.6)--(115.1,-7.7)--cycle;
				\definecolor{ezs}{RGB}{254,232,85}\fill[ezs](51.7,-178.6)--(54,-179.5)--(119.3,-9.3)--(117,-8.4)--cycle;
				\definecolor{ezs}{RGB}{254,234,88}\fill[ezs](53.6,-179.4)--(55.9,-180.2)--(121.2,-10)--(118.9,-9.2)--cycle;
				\definecolor{ezs}{RGB}{254,237,91}\fill[ezs](55.5,-180.1)--(57.8,-181)--(123.1,-10.8)--(120.8,-9.9)--cycle;
				\definecolor{ezs}{RGB}{254,239,94}\fill[ezs](57.4,-180.8)--(59.7,-181.7)--(125,-11.5)--(122.8,-10.6)--cycle;
				\definecolor{ezs}{RGB}{254,241,97}\fill[ezs](59.3,-181.6)--(61.6,-182.4)--(127,-12.2)--(124.7,-11.4)--cycle;
				\definecolor{ezs}{RGB}{254,243,101}\fill[ezs](61.2,-182.3)--(63.5,-183.2)--(128.9,-13)--(126.6,-12.1)--cycle;
				\definecolor{ezs}{RGB}{254,245,103}\fill[ezs](63.2,-183)--(65.4,-183.9)--(130.8,-13.7)--(128.5,-12.8)--cycle;
				\definecolor{ezs}{RGB}{254,245,105}\fill[ezs](65.1,-183.8)--(67.4,-184.6)--(132.7,-14.4)--(130.4,-13.6)--cycle;
				\definecolor{ezs}{RGB}{254,245,105}\fill[ezs](67,-184.5)--(69.3,-185.4)--(134.6,-15.2)--(132.3,-14.3)--cycle;
				\definecolor{ezs}{RGB}{254,244,103}\fill[ezs](68.9,-185.2)--(71.2,-186.1)--(136.5,-15.9)--(134.2,-15)--cycle;
				\definecolor{ezs}{RGB}{253,243,101}\fill[ezs](70.8,-186)--(73.1,-186.8)--(138.4,-16.6)--(136.1,-15.8)--cycle;
				\definecolor{ezs}{RGB}{253,241,98}\fill[ezs](72.7,-186.7)--(75,-187.6)--(140.3,-17.4)--(138,-16.5)--cycle;
				\definecolor{ezs}{RGB}{253,239,95}\fill[ezs](74.6,-187.4)--(76.9,-188.3)--(142.2,-18.1)--(140,-17.2)--cycle;
				\definecolor{ezs}{RGB}{253,236,92}\fill[ezs](76.5,-188.2)--(78.8,-189)--(144.2,-18.8)--(141.9,-18)--cycle;
				\definecolor{ezs}{RGB}{253,234,88}\fill[ezs](78.5,-188.9)--(80.7,-189.8)--(146.1,-19.6)--(143.8,-18.7)--cycle;
				\definecolor{ezs}{RGB}{253,232,85}\fill[ezs](80.4,-189.6)--(82.7,-190.5)--(148,-20.3)--(145.7,-19.4)--cycle;
				\definecolor{ezs}{RGB}{253,230,82}\fill[ezs](82.3,-190.4)--(84.6,-191.2)--(149.9,-21)--(147.6,-20.2)--cycle;
				\definecolor{ezs}{RGB}{253,228,78}\fill[ezs](84.2,-191.1)--(86.5,-192)--(151.8,-21.8)--(149.5,-20.9)--cycle;
				\definecolor{ezs}{RGB}{253,226,74}\fill[ezs](86.1,-191.8)--(88.4,-192.7)--(153.7,-22.5)--(151.4,-21.6)--cycle;
				\definecolor{ezs}{RGB}{253,224,70}\fill[ezs](88,-192.6)--(90.3,-193.4)--(155.6,-23.3)--(153.3,-22.4)--cycle;
				\definecolor{ezs}{RGB}{253,222,67}\fill[ezs](89.9,-193.3)--(92.2,-194.2)--(157.5,-24)--(155.3,-23.1)--cycle;
				\definecolor{ezs}{RGB}{253,220,65}\fill[ezs](91.8,-194)--(94.1,-194.9)--(159.5,-24.7)--(157.2,-23.8)--cycle;
				\definecolor{ezs}{RGB}{253,217,63}\fill[ezs](93.7,-194.8)--(96,-195.6)--(161.4,-25.5)--(159.1,-24.6)--cycle;
				\definecolor{ezs}{RGB}{253,215,60}\fill[ezs](95.7,-195.5)--(97.9,-196.4)--(163.3,-26.2)--(161,-25.3)--cycle;
				\definecolor{ezs}{RGB}{253,213,56}\fill[ezs](97.6,-196.2)--(99.9,-197.1)--(165.2,-26.9)--(162.9,-26)--cycle;
				\definecolor{ezs}{RGB}{252,211,53}\fill[ezs](99.5,-197)--(101.8,-197.8)--(167.1,-27.7)--(164.8,-26.8)--cycle;
				\definecolor{ezs}{RGB}{252,209,49}\fill[ezs](101.4,-197.7)--(103.7,-198.6)--(169,-28.4)--(166.7,-27.5)--cycle;
				\definecolor{ezs}{RGB}{252,207,45}\fill[ezs](103.3,-198.4)--(105.6,-199.3)--(170.9,-29.1)--(168.6,-28.2)--cycle;
				\definecolor{ezs}{RGB}{253,205,42}\fill[ezs](105.2,-199.2)--(107.5,-200)--(172.8,-29.9)--(170.6,-29)--cycle;
				\definecolor{ezs}{RGB}{253,202,40}\fill[ezs](107.1,-199.9)--(109.4,-200.8)--(174.7,-30.6)--(172.5,-29.7)--cycle;
				\definecolor{ezs}{RGB}{253,200,38}\fill[ezs](109,-200.6)--(111.3,-201.5)--(176.7,-31.3)--(174.4,-30.4)--cycle;
				\definecolor{ezs}{RGB}{253,198,35}\fill[ezs](111,-201.4)--(113.2,-202.2)--(178.6,-32.1)--(176.3,-31.2)--cycle;
				\definecolor{ezs}{RGB}{253,196,31}\fill[ezs](112.9,-202.1)--(115.2,-203)--(180.5,-32.8)--(178.2,-31.9)--cycle;
				\definecolor{ezs}{RGB}{253,194,28}\fill[ezs](114.8,-202.8)--(117.1,-203.7)--(182.4,-33.5)--(180.1,-32.6)--cycle;
				\definecolor{ezs}{RGB}{253,192,26}\fill[ezs](116.7,-203.6)--(119,-204.5)--(184.3,-34.3)--(182,-33.4)--cycle;
				\definecolor{ezs}{RGB}{253,189,24}\fill[ezs](118.6,-204.3)--(120.9,-205.2)--(186.2,-35)--(183.9,-34.1)--cycle;
				\definecolor{ezs}{RGB}{253,186,20}\fill[ezs](120.5,-205)--(122.8,-205.9)--(188.1,-35.7)--(185.8,-34.8)--cycle;
				\definecolor{ezs}{RGB}{253,184,17}\fill[ezs](122.4,-205.8)--(124.7,-206.7)--(190,-36.5)--(187.8,-35.6)--cycle;
				\definecolor{ezs}{RGB}{252,183,15}\fill[ezs](124.3,-206.5)--(126.6,-207.4)--(192,-37.2)--(189.7,-36.3)--cycle;
				\definecolor{ezs}{RGB}{252,181,14}\fill[ezs](126.3,-207.2)--(128.5,-208.1)--(193.9,-37.9)--(191.6,-37.1)--cycle;
				\definecolor{ezs}{RGB}{252,179,13}\fill[ezs](128.2,-208)--(130.4,-208.9)--(195.8,-38.7)--(193.5,-37.8)--cycle;
				\definecolor{ezs}{RGB}{252,178,12}\fill[ezs](130.1,-208.7)--(132.4,-209.6)--(197.7,-39.4)--(195.4,-38.5)--cycle;
				\definecolor{ezs}{RGB}{251,178,10}\fill[ezs](132,-209.4)--(134.3,-210.3)--(199.6,-40.1)--(197.3,-39.3)--cycle;
				\definecolor{ezs}{RGB}{251,178,10}\fill[ezs](133.9,-210.2)--(136.2,-211.1)--(201.5,-40.9)--(199.2,-40)--cycle;
				\definecolor{ezs}{RGB}{250,177,11}\fill[ezs](135.8,-210.9)--(138.1,-211.8)--(203.4,-41.6)--(201.1,-40.7)--cycle;
				\definecolor{ezs}{RGB}{250,178,11}\fill[ezs](137.7,-211.6)--(140,-212.5)--(205.3,-42.3)--(203.1,-41.5)--cycle;
				\definecolor{ezs}{RGB}{248,180,10}\fill[ezs](139.6,-212.4)--(141.9,-213.3)--(207.3,-43.1)--(205,-42.2)--cycle;
				\definecolor{ezs}{RGB}{247,181,9}\fill[ezs](141.5,-213.1)--(143.8,-214)--(209.2,-43.8)--(206.9,-42.9)--cycle;
				\definecolor{ezs}{RGB}{246,182,9}\fill[ezs](143.5,-213.8)--(145.7,-214.7)--(211.1,-44.5)--(208.8,-43.7)--cycle;
				\definecolor{ezs}{RGB}{246,182,10}\fill[ezs](145.4,-214.6)--(147.7,-215.5)--(213,-45.3)--(210.7,-44.4)--cycle;
			\end{scope}
			% chữ EDUZY (Times đậm, làm dày nét bằng 5 lớp lệch nhẹ; căn đúng khung chữ logo gốc)
			\foreach \dx/\dy in {-0.7/0,0.7/0,0/-0.3,0/0.3,0/0}
			\node[anchor=center,inner sep=0pt,outer sep=0pt,text=ezText,align=center] at (109.6+\dx,-193.6+\dy)
			{\resizebox{1.01cm}{0.22cm}{\fontencoding{T1}\fontfamily{ptm}\fontseries{b}\selectfont EDUZY}};
		\end{tikzpicture}%
	}%
}
%====================== THIẾT LẬP THÂN ĐỀ ======================
%--- Môi trường tự luận {bt}: ex_test 3.3.4 chưa khai báo sẵn nên tạo ở đây
%    (ex_test đã tự móc \fixshowans và \listenumerate cho bt)
\makeatletter
\@ifundefined{bt}{\newtheorem{bt}{\color{gold1}\namebt}}{}
\makeatother
\renewcommand{\sectionmark}[1]{\markright{\thesection\ #1}}
%--- Khung tiêu đề (logo + tối đa 4 dòng; dòng 3, 4 để trống {} thì bỏ qua)
%  Đề thi:    \khungtieude{THI THỬ THPT QUỐC GIA NĂM 2026}{Môn Toán -- Đề số 21}
%                         {(Đề thi gồm: ...)}{Thời gian làm bài: 90 phút (không kể thời gian phát đề).}
%  Chuyên đề: \khungtieude{CHUYÊN ĐỀ 3 -- DÃY SỐ VÀ CÔNG THỨC TRUY HỒI}
%                         {Bồi dưỡng học sinh giỏi Toán THPT}{...}{}
\newcommand{\khungtieude}[4]{%
	\begin{tcolorbox}[colframe=gold1,colback=gold5!45,boxrule=0.5mm]
		\begin{minipage}[c]{3cm}
			\centerline{\eduzylogo[2.4cm]}
		\end{minipage}\hspace{0.5cm}
		\begin{minipage}[c]{12.5cm}
			\vskip 0.15 cm
			\centerline{\ind{#1}}
			\vskip 0.15 cm
			\centerline{\textbf{#2}}
			\if\relax\detokenize{#3}\relax\else
				\vskip 0.15 cm
				\centerline{\textit{#3}}
			\fi
			\if\relax\detokenize{#4}\relax\else
				\vskip 0.15 cm
				\centerline{\textit{#4}}
			\fi
			\vskip 0.15 cm
		\end{minipage}
	\end{tcolorbox}
}
%====================== BỔ SUNG CHO BÀI HỌC, CHUYÊN ĐỀ ======================
\usepackage[shortlabels]{enumitem}
\usepackage{tabularx}
%--- Danh sách gọn: \begin{enumerate}[gon] ... \end{enumerate}
\SetEnumitemKey{gon}{leftmargin=*,itemsep=1pt,parsep=0pt,topsep=3pt}
%--- Khung lý thuyết: \begin{khung}{} ... \end{khung}; tiêu đề chỉ dùng khi thật ngắn ("Mẹo", "Lỗi thường gặp")
\newtcolorbox{khung}[1]{breakable,enhanced,colframe=gold1,colback=gold5!45,
	boxrule=0.4mm,fonttitle=\bfseries\sffamily,coltitle=white,colbacktitle=gold1,
	attach boxed title to top left={xshift=6mm,yshift=-3mm},
	boxed title style={arc=1mm,boxrule=0.3mm},title={#1}}
%--- Tiêu đề mục nhỏ trong một dạng: \muc{Dấu hiệu}, \muc{Phương pháp}, \muc{Bài tự luyện}
\makeatletter
\newcommand{\muc}[1]{\par\medskip\noindent\indam{#1}\par\nobreak\smallskip\@afterindentfalse\@afterheading}
\makeatother
%--- Ô công thức cần nhớ (căn giữa, viền vàng, vừa khít nội dung)
%     \begin{center}\congthuc{$Q(n)=(n-1)!$}\end{center}
\newtcbox{\congthuc}{on line,colframe=gold1,colback=gold5!45,boxrule=0.5pt,arc=2pt,
	left=8pt,right=8pt,top=4pt,bottom=4pt}
%--- Công tắc lời giải bài tự luyện (ví dụ mẫu luôn dùng \loigiai)
%     \loigiaiTL[{đáp số}]{lời giải đầy đủ}
%     \TLfull -> in lời giải đầy đủ;  \TLdapso -> chỉ in một dòng "Đáp số";  \TLan -> ẩn hoàn toàn
%     (\CheDo là dethi thì mọi lời giải và đáp số đều ẩn)
\makeatletter
\newif\ifTLhien   \TLhientrue
\newif\ifTLdapso  \TLdapsofalse
\newcommand{\TLfull}{\TLhientrue\TLdapsofalse}
\newcommand{\TLdapso}{\TLhientrue\TLdapsotrue}
\newcommand{\TLan}{\TLhienfalse}
\newcommand{\loigiaiTL}[2][]{%
	\ifTLhien
		\ifTLdapso
			\def\@tlarg{#1}%
			\ifx\@tlarg\@empty\else
				\ifdefstring{\kindOPTN}{dethi}{}{\par\noindent{\bfseries\color{gold1}Đáp số:} #1\par}%
			\fi
		\else
			\loigiai{#2}%
		\fi
	\fi}
\makeatother
\TLdapso% mặc định: chỉ in đáp số bài tự luyện
%====================== BẢNG ĐÚNG/SAI (\choiceTFt) ======================
%  Hàng tiêu đề nền nâu đồng chữ trắng (Phát biểu | Đ | S), các hàng tô xen kẽ kem / trắng,
%  đường kẻ vàng kim, nhãn a) b) c) d) màu nâu đồng. Bản có đáp án: cột được chọn có vòng tròn
%  nâu đồng với dấu tích trắng (câu đúng tích ở cột Đúng, câu sai tích ở cột Sai).
%  Tiêu đề ghi đủ "Đúng | Sai": \OPTN{viettat=0} sau \input (cột Đúng/Sai rộng ra, phát biểu hẹp lại).
%  Dùng bảng gốc của ex_test: \def\BangDSgoc{1} trước %============================================================
%  PREAMBLE EDUZY v2 - dùng chung cho đề thi, bài học, chuyên đề (gói ex_test)
%  Khác bản preamble-eduzy.tex ở đầu trang, chân trang: dải nền nâu đồng tràn ngang,
%  logo + Eduzy.com.vn, ô chữ vàng kem cạnh xiên, chân trang "Trang (số)";
%  logo Eduzy chìm giữa mọi trang (tắt: \logochimtat, bật lại: \logochimbat).
%  Cách dùng ở file chính (xem mau-chuan.tex, mau-chuyen-de.tex):
%      \def\CheDo{loigiai}          % loigiai | dethi | solcolor | color | book
%      \def\tenchuong{<dòng chữ trên đầu trang>}
%      \def\thuonghieu{<chữ ở chân trang lẻ>}   % không khai báo thì dùng mặc định
%      \def\tacgia{<chữ ở chân trang chẵn>}      % không khai báo thì dùng mặc định
%      \def\khauhieu{<chữ cạnh logo ở đầu trang>}  % tùy chọn, mặc định Eduzy.com.vn
%      \input{preamble-eduzy-v2}
%      \TLdapso                      % tài liệu: \TLfull | \TLdapso | \TLan
%      \begin{document} ... \end{document}
%  Logo Eduzy vẽ bằng TikZ ngay trong file (\eduzylogo), không cần file ảnh.
%  Biên dịch: pdfLaTeX (máy có gói vietnam) hoặc XeLaTeX.
%============================================================
\documentclass[12pt,a4paper,twoside]{article}%<-- twoside: bat dau trang/chan trang chan-le
\usepackage{iftex}
\usepackage{amsmath,amssymb,tikz,mathrsfs,tkz-tab}
\usepackage{graphicx}
\usepackage{pgfplots}% đồ thị hàm số
\pgfplotsset{compat=1.17}
%\usepackage{mathptmx}
\usepackage{mathpazo}
\ifXeTeX% XeLaTeX: font chữ có đủ dấu tiếng Việt, giữ nguyên font toán mathpazo
	\usepackage[no-math]{fontspec}
	\setmainfont{TeX Gyre Pagella}
	\newfontfamily\EZqag{TeX Gyre Adventor}[NFSSFamily=qag,Ligatures=TeX]
\fi
\usepackage[many]{tcolorbox}
\usetikzlibrary{math,through,calc,intersections,angles,quotes,shapes,shapes.geometric,arrows,patterns,matrix,chains,arrows.meta,decorations.shapes,decorations.fractals,decorations.markings,shadows,decorations.pathreplacing}
\usepackage[top=1.6cm, bottom=1.8cm, left=1.8cm, right=1.8cm]{geometry}% v2: top 1.6cm (bản cũ 1.5cm) để chữ không sát dải đầu trang
\providecommand{\CheDo}{loigiai}% loigiai | dethi | solcolor | color | book
\expandafter\usepackage\expandafter[\CheDo]{ex_test}
% \tenchuong: khai báo ở file chính (trước \input) hoặc \renewcommand sau \input
%============================================================
%  [GỘP TỪ eduzy-de.tex] EDUZY - MÀU NÂU ĐỒNG / VÀNG KIM
%  + ĐẦU TRANG, CHÂN TRANG KIỂU DẢI NGANG (v2), ĐỐI XỨNG TRANG CHẴN / TRANG LẺ.
%  (không cần file eduzy-de.tex riêng nữa)
%  Khối này phải đặt SAU \usepackage[...]{ex_test} và SAU \def\tenchuong{...}
%============================================================
\usepackage{lastpage}
\usepackage{needspace}
%--- DÀN TRANG ĐỀ THI: tách câu (THÊM 05/11/2026; chỉ thêm lệnh mới, KHÔNG đổi hành vi mặc định) ---
%  Hai công tắc, mặc định tắt (= preamble chạy y như cũ). Bật bằng cách đặt TRƯỚC \input{preamble-eduzy-v2}
%  trong file đề (tools/dan_trang/dan_trang.py tự ghi khi đặt \NeedspaceCau):
%    \def\TachCauDS{1}   câu Đúng/Sai được tách: câu dẫn (kèm hình) ở cuối trang, bảng Phát biểu sang trang sau
%                        (câu dẫn và hình không tách nhau, bảng không tách giữa hai trang)
%    \def\TachCauLG{1}   bản lời giải (\CheDo khác dethi): phần đề của câu cao hơn 1/3 trang được phép tách
%  Lệnh dùng trong file đề (giá trị do tools/dan_trang/dan_trang.py đo ở bản dethi):
%    \NeedspaceCau[<cao câu dẫn + hình>]{<cao cả khối câu>}   ngay trên \begin{ex} (câu đầu phần: trên tiêu đề phần)
%        công tắc tắt: đúng bằng \Needspace{<cao cả khối câu>} như cũ;
%        \TachCauDS=1 và có [..] (câu Đúng/Sai được chọn tách): chỉ giữ câu dẫn + hình;
%        \TachCauLG=1 ở bản lời giải: đòi tối đa 1/3 \textheight.
%    \NeedspaceBang{<cao bảng Phát biểu>}%   ngay trước \choiceTFt: \TachCauDS=1 thì giữ trọn bảng trên một trang;
%        công tắc tắt: không làm gì.
\providecommand{\TachCauDS}{0}
\providecommand{\TachCauLG}{0}
\newlength{\EZnsCau}
\newcommand{\NeedspaceCau}[2][]{%
	\setlength{\EZnsCau}{#2}%
	\ifnum\TachCauDS=1 \if\relax\detokenize{#1}\relax\else\setlength{\EZnsCau}{#1}\fi\fi
	\ifnum\TachCauLG=1 \ifdefstring{\kindOPTN}{dethi}{}{\ifdim\EZnsCau>0.3333\textheight\setlength{\EZnsCau}{0.3333\textheight}\fi}\fi
	\Needspace{\EZnsCau}}
\newcommand{\NeedspaceBang}[1]{\ifnum\TachCauDS=1 \Needspace{#1}\fi}

%--- Thang màu: đậm -> nhạt
\definecolor{gold1}{RGB}{126,96,0}    % nâu đồng đậm  #7E6000
\definecolor{gold2}{RGB}{192,144,3}   % vàng kim      #C09003
\definecolor{gold3}{RGB}{255,215,99}  % vàng nhạt     #FFD763
\definecolor{gold4}{RGB}{255,231,153} % vàng rất nhạt #FFE799
\definecolor{gold5}{RGB}{255,243,205} % vàng kem      #FFF3CD
%--- Ghi đè tông xanh/đỏ cũ
\definecolor{myblue}{RGB}{126,96,0}
\colorlet{blue}{gold1}
\colorlet{cyan}{gold2}
\colorlet{teal}{gold2}

%--- Tên chương mặc định (file đề nên tự khai báo lại)
\providecommand{\tenchuong}{}
%--- Chân trang
\providecommand{\thuonghieu}{TOÁN 12 -- Eduzy.com.vn}
\providecommand{\tacgia}{Tô Gia Bảo -- 0353703783}

%--- Chữ cạnh logo ở đầu trang
\providecommand{\khauhieu}{Eduzy.com.vn}

%====================== THÔNG SỐ ĐẦU TRANG, CHÂN TRANG (v2) ======================
%  Dải nền đậm tràn ngang trang; ô chữ vàng kem có cạnh xiên, kèm 2 vạch xiên chuyển màu.
%  Đầu trang lẻ : logo + \khauhieu (trái)   | ô chữ \tenchuong (phải)
%  Chân trang lẻ: ô chữ \thuonghieu (trái)  | Trang (số) (phải)
%  Trang chẵn soi gương trang lẻ; chân trang chẵn ghi \tacgia.
%--- Màu: chỉ dùng thang vàng kim Eduzy
\colorlet{ezNen}{gold1}% dải nền đậm
\colorlet{ezVach1}{gold3}% vạch xiên sát dải nền (đậm hơn)
\colorlet{ezVach2}{gold4}% vạch xiên sát ô chữ (nhạt hơn)
\colorlet{ezO}{gold5}% ô chữ vàng kem
%--- Kích thước
\newlength{\EZdai}\setlength{\EZdai}{1.1cm}% chiều cao dải nền
\newlength{\EZo}\setlength{\EZo}{0.68cm}% chiều cao ô chữ
\newlength{\EZxien}\setlength{\EZxien}{0.55cm}% độ lệch ngang của cạnh xiên
\newlength{\EZvach}\setlength{\EZvach}{0.75cm}% bề rộng mỗi vạch xiên
\newlength{\EZkhe}\setlength{\EZkhe}{0.12cm}% khe trắng giữa vạch và ô chữ
\newlength{\EZdem}\setlength{\EZdem}{0.4cm}% khoảng đệm từ chữ tới cạnh xiên
\newlength{\EZle}\setlength{\EZle}{1.8cm}% lề trái/phải của chữ (bằng lề thân bài)
\newlength{\EZrongmin}\setlength{\EZrongmin}{4.5cm}% ô chữ rộng tối thiểu
\newlength{\EZrong}% bề rộng chữ trong ô (tính tự động)
%--- Font
\newcommand{\EZfontO}{\fontsize{10.5pt}{12pt}\selectfont\bfseries}% chữ trong ô
\newcommand{\EZfontKH}{\fontfamily{qag}\fontsize{11pt}{12pt}\selectfont\bfseries}% \khauhieu cạnh logo
\newcommand{\EZfontTrang}{\fontsize{11.5pt}{12pt}\selectfont\bfseries}% chữ "Trang" và số trang
%--- Đo bề rộng chữ trong ô
\newcommand{\EZdorong}[1]{%
	\settowidth{\EZrong}{\EZfontO #1}%
	\ifdim\EZrong<\EZrongmin\setlength{\EZrong}{\EZrongmin}\fi}
%--- Logo + chữ trắng \khauhieu (dựng sẵn một lần khi vào \begin{document})
\newsavebox{\EZlogobox}
\AtBeginDocument{\sbox{\EZlogobox}{%
		\raisebox{-0.5\height}{\eduzylogo[1cm]}\hspace{0.15cm}%
		\raisebox{-0.5\height}{\color{white}\EZfontKH\khauhieu}}}
%--- Hai vạch xiên + khe trắng. #1: điểm gốc ở mép dưới, #2: hướng (1 sang phải, -1 sang trái),
%    #3: hướng nghiêng (1 kiểu "/", -1 kiểu "\")
\newcommand{\EZvachxien}[3]{%
	\fill[white] (#1) --++({(#3)*\EZxien},\EZo) --++({(#2)*\EZkhe},0) --++({-(#3)*\EZxien},-\EZo) --cycle;
	\fill[ezVach2] ([xshift=(#2)*\EZkhe]#1) --++({(#3)*\EZxien},\EZo) --++({(#2)*\EZvach},0) --++({-(#3)*\EZxien},-\EZo) --cycle;
	\fill[ezVach1] ([xshift=(#2)*(\EZkhe+\EZvach)]#1) --++({(#3)*\EZxien},\EZo) --++({(#2)*\EZvach},0) --++({-(#3)*\EZxien},-\EZo) --cycle;
	\draw[white,line width=0.6pt] ([xshift=(#2)*(\EZkhe+\EZvach)]#1) --++({(#3)*\EZxien},\EZo);
}
%--- Chữ trong ô (căn giữa theo chiều cao ô, không lệch vì dấu tiếng Việt)
\tikzset{ezchu/.style={inner sep=0pt,text=ezNen,text height=0.72em,text depth=0.22em}}

%====================== LOGO CHÌM GIỮA TRANG ======================
%  Vẽ cùng đầu trang nên nằm dưới chữ của thân bài, có trên mọi trang.
%  Tắt: \logochimtat (từ trang đó trở đi), bật lại: \logochimbat.
%  Đổi cỡ, độ đậm: khai báo trước \input, ví dụ \def\EZchimcao{9cm} \def\EZchimdo{0.1}
\providecommand{\EZchimcao}{10cm}% chiều cao logo chìm
\providecommand{\EZchimdo}{0.12}% độ đậm (0 = trong suốt, 1 = rõ hẳn)
\newif\ifEZchim \EZchimtrue
\newcommand{\logochimtat}{\EZchimfalse}
\newcommand{\logochimbat}{\EZchimtrue}
%  Dựng sẵn logo mờ một lần khi vào \begin{document}; transparency group làm mờ cả logo
%  một lượt nên không lộ vệt ở chỗ các dải màu chồng nhau.
\newsavebox{\EZchimtho}
\newsavebox{\EZchimbox}
\AtBeginDocument{%
	\sbox{\EZchimtho}{\eduzylogo[\EZchimcao]}%
	\sbox{\EZchimbox}{\begin{tikzpicture}
			\begin{scope}[transparency group,opacity=\EZchimdo]
				\node[inner sep=0pt] {\usebox{\EZchimtho}};
			\end{scope}
	\end{tikzpicture}}}
\newcommand{\EZvechim}{%
	\ifEZchim
		\node[inner sep=0pt] at (current page.center) {\usebox{\EZchimbox}};
	\fi}

%====================== ĐẦU TRANG - TRANG LẺ ======================
%  (logo bên trái, ô chữ \tenchuong bên phải)
\newcommand{\myfancyheadOdd}{%
	\EZdorong{\tenchuong}%
	\begin{tikzpicture}[remember picture,overlay]
		\coordinate (T) at (current page.north west);
		\coordinate (P) at (current page.north east);
		%--- logo chìm giữa trang
		\EZvechim
		%--- dải nền
		\fill[ezNen] (T) rectangle ([yshift=-\EZdai]P);
		%--- ô chữ: X là góc dưới bên trái, cạnh trái xiên kiểu "/"
		\coordinate (X) at ([xshift=-\EZle-\EZrong-\EZdem-\EZxien,yshift=-\EZdai]P);
		\fill[ezO] (X) --++(\EZxien,\EZo) -- ([yshift=-\EZdai+\EZo]P) -- ([yshift=-\EZdai]P) --cycle;
		\EZvachxien{X}{-1}{1}
		\node[ezchu,anchor=east] at ([xshift=-\EZle,yshift=-\EZdai+0.5*\EZo]P) {\EZfontO\tenchuong};
		%--- logo + tên
		\node[inner sep=0pt,anchor=west] at ([xshift=\EZle,yshift=-0.5*\EZdai]T) {\usebox{\EZlogobox}};
	\end{tikzpicture}%
}

%====================== ĐẦU TRANG - TRANG CHẴN =====================
%  (soi gương: ô chữ \tenchuong bên trái, logo bên phải)
\newcommand{\myfancyheadEven}{%
	\EZdorong{\tenchuong}%
	\begin{tikzpicture}[remember picture,overlay]
		\coordinate (T) at (current page.north west);
		\coordinate (P) at (current page.north east);
		%--- logo chìm giữa trang
		\EZvechim
		%--- dải nền
		\fill[ezNen] (T) rectangle ([yshift=-\EZdai]P);
		%--- ô chữ: X là góc dưới bên phải, cạnh phải xiên kiểu "\"
		\coordinate (X) at ([xshift=\EZle+\EZrong+\EZdem+\EZxien,yshift=-\EZdai]T);
		\fill[ezO] (X) --++(-\EZxien,\EZo) -- ([yshift=-\EZdai+\EZo]T) -- ([yshift=-\EZdai]T) --cycle;
		\EZvachxien{X}{1}{-1}
		\node[ezchu,anchor=west] at ([xshift=\EZle,yshift=-\EZdai+0.5*\EZo]T) {\EZfontO\tenchuong};
		%--- logo + tên
		\node[inner sep=0pt,anchor=east] at ([xshift=-\EZle,yshift=-0.5*\EZdai]P) {\usebox{\EZlogobox}};
	\end{tikzpicture}%
}

%====================== CHÂN TRANG - TRANG LẺ ======================
%  (ô chữ \thuonghieu bên trái, "Trang (số)" bên phải)
\newcommand{\myfancyfootOdd}{%
	\EZdorong{\thuonghieu}%
	\begin{tikzpicture}[remember picture,overlay]
		\coordinate (T) at (current page.south west);
		\coordinate (P) at (current page.south east);
		%--- dải nền
		\fill[ezNen] (T) rectangle ([yshift=\EZdai]P);
		%--- ô chữ sát mép trên dải: X là góc dưới bên phải, cạnh phải xiên kiểu "/"
		\coordinate (X) at ([xshift=\EZle+\EZrong+\EZdem,yshift=\EZdai-\EZo]T);
		\fill[ezO] ([yshift=\EZdai-\EZo]T) -- (X) --++(\EZxien,\EZo) -- ([yshift=\EZdai]T) --cycle;
		\EZvachxien{X}{1}{1}
		\node[ezchu,anchor=west] at ([xshift=\EZle,yshift=\EZdai-0.5*\EZo]T) {\EZfontO\thuonghieu};
		%--- số trang
		\node[circle,draw=white,line width=0.9pt,text=white,inner sep=0pt,minimum size=0.62cm,anchor=east]
		(so) at ([xshift=-\EZle,yshift=0.5*\EZdai]P) {\EZfontTrang\thepage};
		\node[inner sep=0pt,text=white,anchor=base east] (tr) at ([xshift=-0.15cm]so.base west) {\EZfontTrang Trang};
	\end{tikzpicture}%
}

%====================== CHÂN TRANG - TRANG CHẴN ====================
%  (soi gương: "Trang (số)" bên trái, ô chữ \tacgia bên phải)
\newcommand{\myfancyfootEven}{%
	\EZdorong{\tacgia}%
	\begin{tikzpicture}[remember picture,overlay]
		\coordinate (T) at (current page.south west);
		\coordinate (P) at (current page.south east);
		%--- dải nền
		\fill[ezNen] (T) rectangle ([yshift=\EZdai]P);
		%--- ô chữ sát mép trên dải: X là góc dưới bên trái, cạnh trái xiên kiểu "\"
		\coordinate (X) at ([xshift=-\EZle-\EZrong-\EZdem,yshift=\EZdai-\EZo]P);
		\fill[ezO] ([yshift=\EZdai-\EZo]P) -- (X) --++(-\EZxien,\EZo) -- ([yshift=\EZdai]P) --cycle;
		\EZvachxien{X}{-1}{-1}
		\node[ezchu,anchor=east] at ([xshift=-\EZle,yshift=\EZdai-0.5*\EZo]P) {\EZfontO\tacgia};
		%--- số trang
		%    (vòng tròn ẩn so0 nằm giữa dải để lấy đường chân chữ, giống trang lẻ)
		\node[circle,inner sep=0pt,minimum size=0.62cm] (so0) at ([yshift=0.5*\EZdai]T) {\EZfontTrang\phantom{\thepage}};
		\node[inner sep=0pt,text=white,anchor=base west] (tr) at ($(so0.base -| T)+(\EZle,0)$) {\EZfontTrang Trang};
		\node[circle,draw=white,line width=0.9pt,text=white,inner sep=0pt,minimum size=0.62cm,anchor=base west]
		(so) at ([xshift=0.15cm]tr.base east) {\EZfontTrang\thepage};
	\end{tikzpicture}%
}

%--- Giữ tương thích với các file cũ đang gọi \myfancyhead, \myfancyfoot
\newcommand{\myfancyhead}{\myfancyheadOdd}
\newcommand{\myfancyfoot}{\myfancyfootOdd}

%======================= GẮN VÀO TRANG =============================
\pagestyle{fancy}
\fancyhf{}
\fancyhead[LO]{\myfancyheadOdd}
\fancyhead[RE]{\myfancyheadEven}
\fancyfoot[LO]{\myfancyfootOdd}
\fancyfoot[RE]{\myfancyfootEven}
\renewcommand{\headrulewidth}{0pt}
\renewcommand{\footrulewidth}{0pt}
\fancypagestyle{plain}{%
	\fancyhf{}
	\fancyhead[LO]{\myfancyheadOdd}
	\fancyhead[RE]{\myfancyheadEven}
	\fancyfoot[LO]{\myfancyfootOdd}
	\fancyfoot[RE]{\myfancyfootEven}
	\renewcommand{\headrulewidth}{0pt}
	\renewcommand{\footrulewidth}{0pt}
}
%=================== HẾT PHẦN GỘP TỪ eduzy-de.tex ===================
\newcommand{\indam}[1]{
	{\fontsize{12pt}{9pt}\sffamily \selectfont\bfseries\color{gold1} #1}
}
\newcommand{\ind}[1]{
	{\fontsize{14pt}{9pt}\fontfamily{qag} \selectfont\bfseries\color{gold1} #1}
}
\renewcommand{\FalseEX}{\stepcounter{dapan}\textbf{\selectfont\color{gold1}\Alph{dapan}.}\ignorespaces}
\renewtheorem{ex}{\color{gold1}Câu}
\renewcommand{\baselinestretch}{1.1}
\newcommand{\hoac}[1]{
	\left[\begin{aligned}#1\end{aligned}\right.}
\newcommand{\heva}[1]{
	\left\{\begin{aligned}#1\end{aligned}\right.}
\usepackage{esvect}
\def\vec{\vv}
\def\vt{\vv}
\def\overrightarrow{\vv}
%====================== LOGO EDUZY VẼ BẰNG TIKZ ======================
%  Đường viền dò từ logo gốc; dải màu vàng kim đo trực tiếp trên ảnh gốc
%  (nghiêng 21 độ, tô bằng 80 dải hẹp nên hiển thị giống nhau trên mọi trình xem PDF).
%  Dùng: \eduzylogo          -> cao 2.4cm (mặc định)
%        \eduzylogo[1.5cm]   -> cao 1.5cm
\definecolor{ezText}{RGB}{245,200,46}
\newcommand{\eduzylogo}[1][2.4cm]{%
	\resizebox{!}{#1}{%
		\begin{tikzpicture}[x=0.01cm,y=0.01cm]
			\path[use as bounding box] (0,0) rectangle (224,-224);
			% mũ tốt nghiệp + vương miện
			\begin{scope}
				\clip
				(119.4,-176.1)..controls(118.9,-176.1)and(116.8,-176)..(114.6,-176)
				..controls(101.4,-176)and(85.8,-175.9)..(83.7,-175.7)
				..controls(81.8,-175.6)and(81,-175.6)..(80.2,-175.7)
				..controls(78.8,-175.9)and(77.7,-175.9)..(76.2,-175.7)
				..controls(74.9,-175.5)and(70.8,-175.5)..(64.3,-175.7)
				..controls(59.1,-175.9)and(58.3,-175.7)..(58.6,-173.7)
				..controls(58.6,-172.8)and(59.6,-169.6)..(60.3,-168)
				..controls(60.8,-166.9)and(61.4,-165.8)..(62.3,-165)
				..controls(63.1,-164.2)and(64.3,-163.4)..(65,-163.2)
				..controls(65.5,-163.1)and(72.3,-163)..(76.6,-163.1)
				..controls(77.9,-163.1)and(80.1,-163.1)..(81.4,-163)
				..controls(82.8,-162.9)and(86.3,-162.9)..(89.2,-162.9)
				..controls(93,-162.9)and(94.9,-162.8)..(96,-162.7)
				..controls(97.2,-162.5)and(101,-162.5)..(114,-162.4)
				..controls(123,-162.3)and(134.9,-162.2)..(140.2,-162.1)
				..controls(150.1,-162)and(152.2,-162.1)..(153.6,-162.4)
				..controls(154.4,-162.6)and(157.1,-163.8)..(158.1,-164.4)
				..controls(158.8,-164.8)and(159.7,-165.7)..(159.9,-166.3)
				..controls(160,-166.5)and(160.2,-167.1)..(160.2,-167.6)
				..controls(160.3,-168)and(160.6,-169.4)..(160.9,-170.6)
				..controls(161.2,-171.7)and(161.5,-173.1)..(161.6,-173.6)
				..controls(161.8,-174.4)and(161.8,-174.5)..(161.5,-174.9)
				..controls(161.2,-175.6)and(160.8,-175.9)..(159.9,-176)
				..controls(159.2,-176.1)and(154.1,-176.2)..(151.5,-176.2)
				..controls(148.1,-176.1)and(122.6,-176.1)..(121.6,-176.2)
				..controls(120.8,-176.2)and(119.8,-176.2)..(119.4,-176.1)--cycle
				(130.4,-160)..controls(129,-160)and(119.8,-160)..(110.1,-159.9)
				..controls(100.3,-159.8)and(88.5,-159.7)..(83.9,-159.6)
				..controls(77.8,-159.5)and(75.3,-159.5)..(74.5,-159.6)
				..controls(73.9,-159.7)and(72.5,-159.8)..(71.3,-159.8)
				..controls(69.6,-159.8)and(69,-159.8)..(68.5,-159.6)
				..controls(67.6,-159.3)and(66.8,-158.6)..(66.2,-157.6)
				..controls(66,-157.3)and(65.6,-156.5)..(65.3,-156.1)
				..controls(64.6,-155.1)and(64.5,-154.9)..(63.7,-152.9)
				..controls(63.2,-151.6)and(62.9,-150.9)..(61.8,-149.6)
				..controls(61.5,-149.2)and(61.1,-148.5)..(60.9,-147.9)
				..controls(60.7,-147.4)and(60.3,-146.6)..(60.1,-146.1)
				..controls(59.9,-145.5)and(59.5,-144.9)..(59,-144.3)
				..controls(58.5,-143.7)and(58.2,-143.2)..(57.9,-142.5)
				..controls(57.7,-141.9)and(57.3,-141.1)..(57,-140.7)
				..controls(56.6,-140.3)and(56.1,-139.3)..(55.8,-138.5)
				..controls(55.5,-137.8)and(55.1,-136.8)..(54.9,-136.5)
				..controls(54.7,-136.1)and(54.3,-135.2)..(53.9,-134.5)
				..controls(53.6,-133.8)and(53,-132.7)..(52.6,-132)
				..controls(51.6,-130.5)and(51.4,-130.2)..(51.1,-129)
				..controls(50.6,-127.5)and(50.6,-127.4)..(49.7,-126.3)
				..controls(49.2,-125.7)and(48.6,-124.8)..(48.3,-124.2)
				..controls(48,-123.6)and(47.5,-122.7)..(47.3,-122.2)
				..controls(47,-121.8)and(46.5,-120.7)..(46.1,-119.8)
				..controls(45.8,-118.9)and(45.2,-117.8)..(44.9,-117.2)
				..controls(44.2,-116.1)and(43.3,-114.4)..(42.3,-112.2)
				..controls(41.8,-111.3)and(41.3,-110.4)..(41.1,-110.1)
				..controls(40.5,-109.4)and(40.3,-109)..(40,-107.8)
				..controls(39.8,-107)and(39.6,-106.5)..(39.2,-106)
				..controls(38.9,-105.6)and(38.4,-104.7)..(38,-104)
				..controls(37.7,-103.4)and(37.2,-102.5)..(37,-102.1)
				..controls(36.7,-101.7)and(36.3,-100.8)..(36,-100.1)
				..controls(35.6,-98.9)and(35.3,-98.4)..(34.6,-97.6)
				..controls(34.2,-97.1)and(34.2,-96.3)..(34.6,-96)
				..controls(35,-95.7)and(35.5,-95.9)..(36.2,-96.7)
				..controls(36.5,-97)and(37.2,-97.5)..(37.6,-97.7)..controls(38,-98)and(38.6,-98.4)..(39.1,-98.9)
				..controls(39.6,-99.4)and(40.3,-100)..(40.6,-100.3)
				..controls(41.3,-100.8)and(42.4,-101.8)..(43.7,-102.9)
				..controls(44,-103.2)and(44.6,-103.7)..(45,-103.9)
				..controls(45.4,-104.2)and(45.9,-104.6)..(46.1,-104.8)
				..controls(47.1,-105.9)and(48.8,-107.1)..(50.5,-107.9)
				..controls(51.8,-108.6)and(52.3,-109)..(53.4,-110)
				..controls(54.8,-111.2)and(55.6,-111.6)..(57,-112.1)
				..controls(57.5,-112.3)and(58.5,-112.8)..(59.4,-113.3)
				..controls(60.2,-113.7)and(61.3,-114.3)..(61.9,-114.5)
				..controls(62.4,-114.7)and(63.3,-115.1)..(63.8,-115.3)
				..controls(64.4,-115.6)and(65.4,-115.9)..(66.2,-116.1)
				..controls(67,-116.3)and(68.4,-116.6)..(69.2,-116.9)
				..controls(72.3,-117.9)and(72.8,-118)..(78,-118.1)
				..controls(84,-118.3)and(86.1,-118.1)..(88.1,-117.1)
				..controls(89.4,-116.5)and(93,-115)..(93.7,-114.9)
				..controls(94.8,-114.6)and(95.4,-114.2)..(95.9,-113.4)
				..controls(96.2,-113)and(96.7,-112.5)..(97.1,-112.2)
				..controls(98,-111.6)and(101.3,-108.4)..(102.1,-107.2)
				..controls(102.4,-106.8)and(102.9,-106.2)..(103.2,-105.8)
				..controls(103.5,-105.5)and(104,-104.8)..(104.2,-104.2)
				..controls(104.5,-103.6)and(105,-102.8)..(105.3,-102.4)
				..controls(105.6,-102)and(106.1,-101.2)..(106.5,-100.6)
				..controls(107.1,-99.5)and(108.6,-97.5)..(109.5,-96.5)
				..controls(109.9,-96.1)and(110.1,-96)..(110.4,-96)
				..controls(111.1,-96)and(111.6,-96.5)..(111.8,-97.6)
				..controls(112,-98.6)and(112.3,-99.3)..(112.8,-99.9)
				..controls(113.7,-100.9)and(115.2,-103.1)..(115.7,-104.2)
				..controls(116.1,-105)and(116.5,-105.5)..(116.9,-106)
				..controls(117.3,-106.4)and(117.8,-106.9)..(118,-107.3)
				..controls(118.2,-107.6)and(118.6,-108)..(119,-108.3)
				..controls(119.6,-108.8)and(121.1,-110.3)..(121.8,-111.2)
				..controls(122.1,-111.7)and(122.5,-112)..(123.4,-112.4)
				..controls(124.8,-113.2)and(125.8,-113.8)..(127,-114.6)
				..controls(127.5,-114.9)and(128.1,-115.3)..(128.2,-115.4)
				..controls(128.4,-115.5)and(129.3,-115.8)..(130.2,-116)
				..controls(131.1,-116.2)and(132.5,-116.6)..(133.3,-116.9)
				..controls(135.7,-117.9)and(136.1,-117.9)..(142.5,-117.8)
				..controls(145.8,-117.8)and(148.4,-117.7)..(148.8,-117.6)
				..controls(149.2,-117.6)and(150.2,-117.2)..(151,-116.8)
				..controls(152,-116.4)and(152.8,-116)..(153.6,-115.9)
				..controls(154.5,-115.7)and(155.2,-115.4)..(156.3,-114.9)
				..controls(157.1,-114.5)and(158.2,-114)..(158.7,-113.8)
				..controls(159.1,-113.7)and(159.8,-113.3)..(160.2,-113)
				..controls(160.6,-112.7)and(161.3,-112.4)..(161.7,-112.2)
				..controls(162.7,-111.9)and(163.2,-111.7)..(164.1,-111.1)
				..controls(164.5,-110.8)and(165.3,-110.4)..(165.9,-110)
				..controls(166.5,-109.7)and(167.3,-109.2)..(167.6,-108.9)
				..controls(167.9,-108.6)and(168.6,-108.2)..(169,-108)
				..controls(169.4,-107.8)and(170.2,-107.3)..(170.7,-106.9)
				..controls(171.2,-106.5)and(172.1,-105.8)..(172.7,-105.4)
				..controls(173.2,-105)and(174.1,-104.4)..(174.5,-104)
				..controls(174.9,-103.7)and(175.7,-103.1)..(176.2,-102.6)
				..controls(176.7,-102.2)and(177.4,-101.5)..(177.9,-101.2)
				..controls(178.3,-100.8)and(178.8,-100.4)..(179,-100.2)
				..controls(179.2,-100)and(179.6,-99.6)..(180,-99.3)
				..controls(180.4,-99.1)and(180.9,-98.6)..(181.2,-98.4)
				..controls(181.5,-98.1)and(182.1,-97.7)..(182.5,-97.5)
				..controls(183.2,-97.2)and(183.5,-96.9)..(184.3,-95.9)
				..controls(184.6,-95.6)and(184.6,-95.6)..(184.8,-95.7)
				..controls(185,-95.8)and(185.1,-96)..(185.1,-96.2)
				..controls(185.1,-96.5)and(185,-96.7)..(184.6,-97.2)
				..controls(184.1,-97.8)and(184,-97.9)..(183.9,-98.8)
				..controls(183.7,-100.3)and(183.3,-101.3)..(182.9,-101.8)
				..controls(182.6,-102.1)and(182.3,-102.6)..(182.2,-102.9)
				..controls(182,-103.3)and(181.6,-104)..(181.2,-104.6)
				..controls(180.9,-105.2)and(180.5,-106.1)..(180.4,-106.6)
				..controls(179.9,-108)and(179.7,-108.5)..(179,-109.4)
				..controls(178.6,-109.8)and(178.2,-110.5)..(177.9,-111)
				..controls(177.7,-111.5)and(177.3,-112.3)..(177,-112.7)
				..controls(176.6,-113.2)and(176.4,-113.7)..(176.3,-114.2)
				..controls(176.2,-114.7)and(175.9,-115.2)..(175.4,-116)
				..controls(175,-116.6)and(174.4,-117.7)..(174,-118.5)
				..controls(173.7,-119.2)and(173.3,-120.1)..(173,-120.5)
				..controls(172.8,-120.9)and(172.4,-121.8)..(172.1,-122.5)
				..controls(171.8,-123.2)and(171.3,-124.1)..(171,-124.6)
				..controls(170.8,-125.1)and(170.3,-126)..(169.9,-126.6)
				..controls(169.6,-127.3)and(169.2,-128.1)..(169,-128.3)
				..controls(168.8,-128.6)and(168.5,-129.3)..(168.3,-129.8)
				..controls(168.1,-130.3)and(167.8,-131.1)..(167.5,-131.6)
				..controls(166.9,-132.5)and(164.2,-137.9)..(164.2,-138.1)
				..controls(164.2,-138.1)and(163.9,-138.8)..(163.4,-139.7)
				..controls(162.9,-140.5)and(162.4,-141.5)..(162.3,-141.9)
				..controls(161.9,-143)and(161.6,-143.4)..(160.9,-144.1)
				..controls(160.4,-144.6)and(160.3,-144.8)..(160.1,-145.4)
				..controls(160,-145.8)and(159.6,-146.8)..(159.2,-147.8)
				..controls(158.7,-148.8)and(158.1,-150)..(157.9,-150.6)
				..controls(157.6,-151.1)and(157.2,-151.9)..(156.9,-152.4)
				..controls(156.5,-152.8)and(156.1,-153.7)..(155.8,-154.4)
				..controls(155.6,-155)and(155.2,-155.9)..(154.9,-156.3)
				..controls(154.7,-156.7)and(154.4,-157.4)..(154.2,-157.9)
				..controls(154.1,-158.4)and(153.8,-159)..(153.5,-159.3)
				..controls(152.8,-160.1)and(153,-160.1)..(142.4,-160.1)
				..controls(137.2,-160.1)and(131.8,-160.1)..(130.4,-160)--cycle
				(141.5,-103.7)..controls(140.8,-103)and(139.8,-102.1)..(139.2,-101.6)
				..controls(138.5,-101.2)and(137.9,-100.6)..(137.8,-100.4)
				..controls(137,-99.5)and(134,-96.5)..(133.2,-96)
				..controls(132.7,-95.7)and(131.8,-94.9)..(131.1,-94.3)
				..controls(129.4,-92.7)and(127.9,-91.6)..(126.6,-91.1)
				..controls(125.2,-90.5)and(124.1,-90)..(122.9,-89.1)
				..controls(122.1,-88.6)and(121.3,-88.2)..(119.9,-87.7)
				..controls(118.8,-87.3)and(117.5,-86.8)..(116.9,-86.5)
				..controls(114.3,-85.5)and(113.7,-85.4)..(109.7,-85.5)
				..controls(105.3,-85.7)and(103.9,-85.9)..(101.6,-87)
				..controls(100.7,-87.4)and(99.4,-87.8)..(97.9,-88)
				..controls(97,-88.2)and(96.6,-88.3)..(95.8,-88.9)
				..controls(95.4,-89.2)and(94.7,-89.6)..(94.2,-89.8)
				..controls(92.9,-90.3)and(92.3,-90.7)..(90.9,-92.1)
				..controls(90.2,-92.7)and(89.4,-93.3)..(89,-93.5)..controls(87,-94.5)and(85.2,-96)..(83.6,-97.9)
				..controls(83.3,-98.2)and(82.5,-98.9)..(81.8,-99.4)
				..controls(80.7,-100.1)and(80.5,-100.3)..(80,-101.1)
				..controls(78.7,-103.1)and(77.2,-103.9)..(76.4,-103)
				..controls(75.7,-102.2)and(75.7,-103.2)..(75.7,-83.6)
				..controls(75.6,-73.6)and(75.6,-64.8)..(75.5,-64)
				..controls(75.5,-63.2)and(75.5,-61.4)..(75.6,-59.8)--(75.7,-57.1)--(76,-56.8)
				..controls(76.8,-56.2)and(78,-56.3)..(79.8,-57.2)
				..controls(80.4,-57.5)and(81.3,-57.8)..(81.8,-57.9)
				..controls(82.3,-58.1)and(83,-58.4)..(83.6,-58.8)..controls(84.3,-59.2)and(85,-59.5)..(86,-59.8)
				..controls(86.8,-60.1)and(88,-60.6)..(88.7,-60.9)
				..controls(89.3,-61.2)and(90.6,-61.7)..(91.4,-62)
				..controls(92.2,-62.3)and(93.7,-62.9)..(94.6,-63.3)
				..controls(95.6,-63.7)and(96.8,-64.1)..(97.4,-64.3)
				..controls(98,-64.5)and(99,-64.9)..(99.7,-65.2)
				..controls(100.4,-65.6)and(101.5,-66)..(102.3,-66.2)
				..controls(104.2,-66.8)and(105,-67.1)..(106.2,-67.8)
				..controls(107.9,-68.7)and(110.2,-69.4)..(111.3,-69.2)
				..controls(112.5,-69)and(115,-68.2)..(115.6,-67.8)
				..controls(117.3,-66.7)and(117.6,-66.6)..(119.5,-66.2)
				..controls(121.9,-65.8)and(122.1,-65.7)..(123,-65.1)
				..controls(124.1,-64.4)and(124.4,-64.3)..(126,-63.9)
				..controls(127,-63.7)and(127.7,-63.4)..(128.4,-63.1)
				..controls(129,-62.8)and(129.7,-62.5)..(130.6,-62.4)
				..controls(132.1,-62.1)and(134,-61.5)..(135.2,-61)
				..controls(135.7,-60.8)and(136.7,-60.5)..(137.3,-60.3)
				..controls(138.4,-60)and(139.3,-59.7)..(141.1,-59)
				..controls(141.6,-58.7)and(142,-58.6)..(142.5,-58.6)
				..controls(143.2,-58.6)and(143.3,-58.7)..(143.6,-59)
				..controls(144,-59.5)and(144.1,-59.9)..(144,-60.8)
				..controls(144,-61.2)and(144,-71.1)..(144,-82.8)--(143.9,-104.1)--(143.7,-104.5)
				..controls(143.5,-104.7)and(143.3,-104.8)..(143,-104.9)
				..controls(142.6,-104.9)and(142.5,-104.8)..(141.5,-103.7)--cycle
				(27.8,-91.4)..controls(27.1,-90.9)and(26.3,-90)..(26.1,-89.5)
				..controls(26,-89.2)and(25.8,-88.6)..(25.7,-87.9)
				..controls(25.4,-86.2)and(25.6,-84.8)..(26.5,-83.5)
				..controls(27.4,-82.3)and(28.3,-82)..(29.4,-82.3)
				..controls(31.1,-82.8)and(32.1,-84.9)..(32,-87.4)
				..controls(31.9,-89.3)and(31.2,-90.6)..(29.9,-91.4)
				..controls(29.1,-91.9)and(28.6,-91.9)..(27.8,-91.4)--cycle
				(190.1,-91.3)..controls(189.7,-91.2)and(189.3,-90.9)..(189.1,-90.7)
				..controls(188.3,-89.8)and(187.6,-87.7)..(187.6,-86.2)
				..controls(187.6,-85.5)and(187.7,-85.2)..(188.2,-84.4)
				..controls(189.1,-82.6)and(190.1,-81.8)..(191.3,-82.1)
				..controls(192,-82.3)and(193.5,-83.7)..(193.8,-84.4)
				..controls(194.1,-85.1)and(194.2,-87.1)..(194,-88.3)
				..controls(193.8,-89.8)and(193.3,-90.5)..(192.3,-91.3)
				..controls(191.7,-91.7)and(190.9,-91.7)..(190.1,-91.3)--cycle
				(173.7,-83.4)..controls(172.2,-82.5)and(172,-82.5)..(170.6,-82.4)
				..controls(169.7,-82.4)and(168.8,-82.4)..(168.3,-82.5)
				..controls(167.2,-82.7)and(166.9,-82.7)..(166.6,-82.2)
				..controls(166.3,-81.8)and(166.3,-81)..(166.8,-80)
				..controls(167.2,-78.9)and(167.4,-77.4)..(167.7,-74.1)
				..controls(167.8,-72.8)and(167.9,-71.8)..(167.9,-71.7)
				..controls(168,-71.7)and(168,-70.2)..(168,-68.4)..controls(168,-65.4)and(168,-64.9)..(167.7,-63)
				..controls(167.3,-60.5)and(167.1,-59.9)..(166.4,-58.4)
				..controls(166.1,-57.7)and(165.8,-56.9)..(165.7,-56.6)
				..controls(165.7,-56.2)and(165.4,-55.6)..(165.2,-55.2)
				..controls(164.6,-53.8)and(161.9,-50.7)..(161,-50.3)
				..controls(159.8,-49.7)and(157.9,-50)..(155.2,-51.2)
				..controls(154.7,-51.4)and(153.4,-51.8)..(152.2,-52.1)
				..controls(150.8,-52.4)and(149.4,-52.9)..(148.6,-53.2)
				..controls(147.8,-53.6)and(146.7,-54)..(146,-54.2)
				..controls(145.4,-54.4)and(144.5,-54.7)..(144,-54.9)
				..controls(143.6,-55.2)and(142.7,-55.5)..(141.8,-55.7)
				..controls(141,-55.9)and(139.6,-56.3)..(138.7,-56.6)
				..controls(137.8,-57)and(136.3,-57.4)..(135.3,-57.6)
				..controls(133.8,-57.9)and(133.2,-58.1)..(132.1,-58.7)
				..controls(131.1,-59.2)and(130.3,-59.4)..(129.4,-59.6)
				..controls(128.7,-59.8)and(127.8,-60)..(127.4,-60.1)
				..controls(127,-60.3)and(126,-60.5)..(125.1,-60.7)
				..controls(123.2,-61.1)and(121.7,-61.5)..(120.2,-62.1)
				..controls(119.7,-62.3)and(118.1,-62.8)..(116.7,-63.2)
				..controls(114.6,-63.9)and(114,-64.1)..(113,-64.7)
				..controls(111.9,-65.4)and(111.8,-65.4)..(111,-65.5)
				..controls(109.9,-65.5)and(109.3,-65.4)..(107.7,-64.7)
				..controls(107,-64.4)and(106,-64)..(105.3,-63.9)
				..controls(103.5,-63.5)and(102.2,-63)..(101.2,-62.5)
				..controls(100.1,-61.8)and(99.3,-61.5)..(97.8,-61.1)
				..controls(97.2,-61)and(96.2,-60.7)..(95.7,-60.4)
				..controls(95.1,-60.1)and(94.2,-59.8)..(93.5,-59.6)
				..controls(92.9,-59.4)and(91.7,-59)..(91,-58.7)..controls(90.2,-58.4)and(88.9,-58)..(88,-57.7)
				..controls(87.2,-57.5)and(86.1,-57.1)..(85.5,-56.8)
				..controls(85,-56.6)and(84.1,-56.2)..(83.5,-56)..controls(82.9,-55.7)and(81.8,-55.2)..(81,-54.9)
				..controls(80.1,-54.5)and(79.1,-54.2)..(78.7,-54.1)
				..controls(77.9,-53.8)and(77.3,-53.6)..(76.2,-52.9)
				..controls(75.7,-52.6)and(75.2,-52.4)..(74.1,-52.2)
				..controls(72.5,-51.9)and(72.4,-51.8)..(71.3,-51.1)
				..controls(69.9,-50.1)and(69.4,-49.9)..(67.9,-49.8)
				..controls(67,-49.7)and(66.6,-49.6)..(66.3,-49.4)
				..controls(66.1,-49.3)and(65.7,-49)..(65.4,-48.8)
				..controls(65.1,-48.7)and(64.2,-48.4)..(63.6,-48.2)
				..controls(62.9,-48)and(61.7,-47.6)..(60.9,-47.2)
				..controls(60.1,-46.8)and(58.9,-46.4)..(58.2,-46.2)
				..controls(57.6,-46)and(56.8,-45.7)..(56.5,-45.5)
				..controls(54.9,-44.6)and(54.6,-44.5)..(53.2,-44.2)
				..controls(51.7,-43.8)and(51,-43.6)..(49.9,-42.9)
				..controls(49.5,-42.6)and(48.9,-42.4)..(47.5,-42)
				..controls(45.7,-41.6)and(45.6,-41.5)..(44.3,-40.6)
				..controls(43.5,-40.1)and(42.8,-39.6)..(42.7,-39.4)
				..controls(42.4,-38.7)and(43,-38.1)..(44.5,-37.6)..controls(44.9,-37.5)and(46,-37)..(46.9,-36.6)
				..controls(47.9,-36.1)and(48.8,-35.8)..(49.6,-35.6)
				..controls(50.5,-35.4)and(51.3,-35.2)..(52.4,-34.6)
				..controls(53.3,-34.1)and(54.5,-33.7)..(54.9,-33.5)
				..controls(55.4,-33.4)and(56.4,-33)..(57.2,-32.6)
				..controls(58.6,-32)and(60.1,-31.5)..(61.6,-31.2)
				..controls(62.3,-31.1)and(62.7,-31)..(63.2,-30.7)
				..controls(64.1,-30)and(64.3,-29.9)..(65.3,-29.7)
				..controls(65.7,-29.6)and(66.5,-29.2)..(67.1,-28.9)
				..controls(67.6,-28.6)and(68.8,-28.2)..(69.6,-27.9)
				..controls(70.5,-27.7)and(71.6,-27.3)..(72.1,-27)
				..controls(72.6,-26.7)and(73.8,-26.2)..(74.8,-25.9)
				..controls(75.7,-25.5)and(76.7,-25.1)..(77.1,-25)
				..controls(77.4,-24.8)and(78.5,-24.4)..(79.6,-24.1)
				..controls(80.6,-23.8)and(82,-23.2)..(82.8,-22.9)
				..controls(83.5,-22.5)and(84.8,-22)..(85.8,-21.8)
				..controls(86.7,-21.5)and(87.8,-21.1)..(88.1,-20.9)
				..controls(88.4,-20.7)and(89.3,-20.3)..(90.1,-20.1)
				..controls(90.8,-19.8)and(91.9,-19.4)..(92.4,-19.1)
				..controls(93.5,-18.5)and(94.4,-18.2)..(95.4,-18.1)
				..controls(95.7,-18)and(96.4,-17.8)..(96.9,-17.5)
				..controls(97.5,-17.3)and(98.4,-16.9)..(99,-16.6)
				..controls(99.7,-16.3)and(100.9,-15.9)..(101.9,-15.6)
				..controls(102.9,-15.3)and(104.4,-14.7)..(105.3,-14.3)
				..controls(108.1,-13.2)and(108.5,-13.1)..(110.1,-13.1)
				..controls(111.4,-13.1)and(111.5,-13.1)..(112.4,-13.5)
				..controls(113.2,-14)and(115,-14.8)..(116.6,-15.4)
				..controls(116.9,-15.5)and(117.9,-15.8)..(118.8,-16.1)
				..controls(120.9,-16.6)and(121.5,-16.9)..(122.6,-17.9)
				..controls(123,-18.3)and(123.7,-18.8)..(123.9,-18.9)
				..controls(124.6,-19.4)and(126,-19.8)..(127.9,-20.2)
				..controls(129.6,-20.5)and(129.8,-20.6)..(131.1,-21.4)
				..controls(131.5,-21.7)and(132.4,-22.1)..(133.1,-22.3)
				..controls(136.1,-23.2)and(138.3,-24.1)..(140.2,-25)
				..controls(141,-25.5)and(141.5,-25.6)..(142.5,-25.8)
				..controls(145.1,-26.3)and(146,-26.6)..(149.4,-28.1)
				..controls(151.3,-28.9)and(153.5,-29.8)..(154.3,-30.1)
				..controls(155.1,-30.4)and(156.2,-30.8)..(156.6,-31)
				..controls(157.1,-31.2)and(158.2,-31.6)..(159.1,-31.8)
				..controls(160.6,-32.2)and(160.9,-32.3)..(161.7,-32.9)
				..controls(162.4,-33.3)and(162.9,-33.6)..(163.7,-33.8)
				..controls(164.3,-34)and(164.8,-34.1)..(164.9,-34.1)
				..controls(165,-34.1)and(165.6,-34.4)..(166.4,-34.8)
				..controls(167.3,-35.2)and(168.1,-35.4)..(169,-35.6)
				..controls(170,-35.8)and(170.5,-36)..(171.7,-36.6)
				..controls(172.4,-37)and(173.3,-37.4)..(173.6,-37.4)
				..controls(175.6,-37.7)and(176.2,-38)..(176.7,-39.1)
				..controls(176.9,-39.6)and(176.9,-39.7)..(176.7,-40.4)
				..controls(176.4,-41.6)and(175.8,-42.7)..(175.2,-43.4)
				..controls(174.5,-44.1)and(172.6,-45.5)..(171.7,-45.8)
				..controls(170.6,-46.2)and(169.1,-46)..(168.2,-45.4)
				..controls(167.9,-45.2)and(167.3,-44.8)..(166.9,-44.4)
				..controls(166.4,-44.1)and(165.8,-43.6)..(165.4,-43.2)
				..controls(165.1,-42.9)and(164.4,-42.3)..(163.9,-41.9)
				..controls(163.3,-41.5)and(162.7,-40.9)..(162.4,-40.5)
				..controls(161.4,-39.2)and(159.9,-38)..(158.8,-37.8)
				..controls(158.4,-37.8)and(158.3,-37.7)..(158.2,-37.2)
				..controls(158,-36.8)and(157.4,-36.1)..(157.1,-36.1)
				..controls(157,-36.1)and(156.7,-36.1)..(156.5,-36)
				..controls(156.1,-35.9)and(155.9,-35.7)..(155.7,-35.2)
				..controls(155.3,-34.6)and(155.1,-34.4)..(154,-34.1)
				..controls(153.6,-34)and(153.2,-33.7)..(152.9,-33.5)
				..controls(152.4,-32.9)and(151.6,-32.4)..(151.3,-32.4)
				..controls(150.9,-32.4)and(150.4,-32.6)..(150.4,-32.9)
				..controls(150.4,-33.2)and(150.8,-33.4)..(151.4,-33.6)
				..controls(152.1,-33.7)and(152.8,-34.1)..(153.5,-34.9)
				..controls(154.1,-35.5)and(154.3,-35.6)..(155,-35.8)
				..controls(155.7,-36)and(155.7,-36)..(156,-36.6)
				..controls(156.3,-37.3)and(156.7,-37.7)..(157.5,-37.8)
				..controls(158.1,-37.9)and(158.2,-38)..(158.4,-38.8)
				..controls(158.5,-39.4)and(158.6,-39.6)..(159.6,-40.4)
				..controls(160.2,-41)and(160.9,-41.6)..(161.3,-41.8)
				..controls(161.7,-42)and(161.9,-42.3)..(162.2,-42.8)
				..controls(162.4,-43.2)and(162.8,-43.8)..(163.1,-44.1)
				..controls(163.9,-44.8)and(165.1,-46.7)..(165.7,-48)
				..controls(166.2,-49.3)and(166.4,-49.7)..(167.2,-50.5)
				..controls(167.5,-50.8)and(167.9,-51.4)..(168.1,-51.8)
				..controls(168.3,-52.2)and(168.8,-53)..(169.1,-53.6)
				..controls(169.6,-54.3)and(169.8,-54.9)..(169.9,-55.6)
				..controls(170.4,-57.8)and(170.5,-58.1)..(171.2,-59.5)
				..controls(171.6,-60.3)and(172,-61.5)..(172.2,-62.1)
				..controls(172.3,-62.8)and(172.7,-63.7)..(172.9,-64.1)
				..controls(173.2,-64.6)and(173.4,-65.3)..(173.4,-65.6)
				..controls(173.8,-68.8)and(174.2,-70.6)..(174.9,-72.6)
				..controls(175.5,-74.5)and(175.9,-77.1)..(176.1,-80.2)
				..controls(176.3,-83.4)and(175.5,-84.4)..(173.7,-83.4)--cycle;
				\definecolor{ezs}{RGB}{247,179,26}\fill[ezs](-5.7,-156.6)--(-3.4,-157.5)--(61.9,12.7)--(59.7,13.6)--cycle;
				\definecolor{ezs}{RGB}{247,179,22}\fill[ezs](-3.8,-157.3)--(-1.5,-158.2)--(63.9,12)--(61.6,12.9)--cycle;
				\definecolor{ezs}{RGB}{248,179,19}\fill[ezs](-1.8,-158.1)--(0.4,-158.9)--(65.8,11.2)--(63.5,12.1)--cycle;
				\definecolor{ezs}{RGB}{248,178,16}\fill[ezs](0.1,-158.8)--(2.4,-159.7)--(67.7,10.5)--(65.4,11.4)--cycle;
				\definecolor{ezs}{RGB}{249,178,15}\fill[ezs](2,-159.5)--(4.3,-160.4)--(69.6,9.8)--(67.3,10.7)--cycle;
				\definecolor{ezs}{RGB}{251,177,17}\fill[ezs](3.9,-160.3)--(6.2,-161.1)--(71.5,9)--(69.2,9.9)--cycle;
				\definecolor{ezs}{RGB}{252,179,18}\fill[ezs](5.8,-161)--(8.1,-161.9)--(73.4,8.3)--(71.1,9.2)--cycle;
				\definecolor{ezs}{RGB}{252,182,17}\fill[ezs](7.7,-161.7)--(10,-162.6)--(75.3,7.6)--(73,8.5)--cycle;
				\definecolor{ezs}{RGB}{252,185,17}\fill[ezs](9.6,-162.5)--(11.9,-163.4)--(77.2,6.8)--(75,7.7)--cycle;
				\definecolor{ezs}{RGB}{251,189,18}\fill[ezs](11.5,-163.2)--(13.8,-164.1)--(79.2,6.1)--(76.9,7)--cycle;
				\definecolor{ezs}{RGB}{251,191,22}\fill[ezs](13.4,-163.9)--(15.7,-164.8)--(81.1,5.4)--(78.8,6.2)--cycle;
				\definecolor{ezs}{RGB}{252,193,25}\fill[ezs](15.4,-164.7)--(17.6,-165.6)--(83,4.6)--(80.7,5.5)--cycle;
				\definecolor{ezs}{RGB}{253,195,27}\fill[ezs](17.3,-165.4)--(19.6,-166.3)--(84.9,3.9)--(82.6,4.8)--cycle;
				\definecolor{ezs}{RGB}{253,196,31}\fill[ezs](19.2,-166.1)--(21.5,-167)--(86.8,3.2)--(84.5,4)--cycle;
				\definecolor{ezs}{RGB}{254,199,35}\fill[ezs](21.1,-166.9)--(23.4,-167.8)--(88.7,2.4)--(86.4,3.3)--cycle;
				\definecolor{ezs}{RGB}{253,201,38}\fill[ezs](23,-167.6)--(25.3,-168.5)--(90.6,1.7)--(88.3,2.6)--cycle;
				\definecolor{ezs}{RGB}{253,203,40}\fill[ezs](24.9,-168.3)--(27.2,-169.2)--(92.5,1)--(90.3,1.8)--cycle;
				\definecolor{ezs}{RGB}{253,205,42}\fill[ezs](26.8,-169.1)--(29.1,-170)--(94.4,0.2)--(92.2,1.1)--cycle;
				\definecolor{ezs}{RGB}{253,207,45}\fill[ezs](28.7,-169.8)--(31,-170.7)--(96.4,-0.5)--(94.1,0.4)--cycle;
				\definecolor{ezs}{RGB}{253,210,49}\fill[ezs](30.7,-170.5)--(32.9,-171.4)--(98.3,-1.2)--(96,-0.4)--cycle;
				\definecolor{ezs}{RGB}{253,212,54}\fill[ezs](32.6,-171.3)--(34.9,-172.2)--(100.2,-2)--(97.9,-1.1)--cycle;
				\definecolor{ezs}{RGB}{253,214,59}\fill[ezs](34.5,-172)--(36.8,-172.9)--(102.1,-2.7)--(99.8,-1.8)--cycle;
				\definecolor{ezs}{RGB}{253,216,63}\fill[ezs](36.4,-172.7)--(38.7,-173.6)--(104,-3.4)--(101.7,-2.6)--cycle;
				\definecolor{ezs}{RGB}{253,218,65}\fill[ezs](38.3,-173.5)--(40.6,-174.4)--(105.9,-4.2)--(103.6,-3.3)--cycle;
				\definecolor{ezs}{RGB}{254,220,66}\fill[ezs](40.2,-174.2)--(42.5,-175.1)--(107.8,-4.9)--(105.5,-4)--cycle;
				\definecolor{ezs}{RGB}{254,222,68}\fill[ezs](42.1,-175)--(44.4,-175.8)--(109.7,-5.6)--(107.5,-4.8)--cycle;
				\definecolor{ezs}{RGB}{254,224,71}\fill[ezs](44,-175.7)--(46.3,-176.6)--(111.7,-6.4)--(109.4,-5.5)--cycle;
				\definecolor{ezs}{RGB}{254,226,75}\fill[ezs](46,-176.4)--(48.2,-177.3)--(113.6,-7.1)--(111.3,-6.2)--cycle;
				\definecolor{ezs}{RGB}{254,228,79}\fill[ezs](47.9,-177.2)--(50.1,-178)--(115.5,-7.8)--(113.2,-7)--cycle;
				\definecolor{ezs}{RGB}{254,230,82}\fill[ezs](49.8,-177.9)--(52.1,-178.8)--(117.4,-8.6)--(115.1,-7.7)--cycle;
				\definecolor{ezs}{RGB}{254,232,85}\fill[ezs](51.7,-178.6)--(54,-179.5)--(119.3,-9.3)--(117,-8.4)--cycle;
				\definecolor{ezs}{RGB}{254,234,88}\fill[ezs](53.6,-179.4)--(55.9,-180.2)--(121.2,-10)--(118.9,-9.2)--cycle;
				\definecolor{ezs}{RGB}{254,237,91}\fill[ezs](55.5,-180.1)--(57.8,-181)--(123.1,-10.8)--(120.8,-9.9)--cycle;
				\definecolor{ezs}{RGB}{254,239,94}\fill[ezs](57.4,-180.8)--(59.7,-181.7)--(125,-11.5)--(122.8,-10.6)--cycle;
				\definecolor{ezs}{RGB}{254,241,97}\fill[ezs](59.3,-181.6)--(61.6,-182.4)--(127,-12.2)--(124.7,-11.4)--cycle;
				\definecolor{ezs}{RGB}{254,243,101}\fill[ezs](61.2,-182.3)--(63.5,-183.2)--(128.9,-13)--(126.6,-12.1)--cycle;
				\definecolor{ezs}{RGB}{254,245,103}\fill[ezs](63.2,-183)--(65.4,-183.9)--(130.8,-13.7)--(128.5,-12.8)--cycle;
				\definecolor{ezs}{RGB}{254,245,105}\fill[ezs](65.1,-183.8)--(67.4,-184.6)--(132.7,-14.4)--(130.4,-13.6)--cycle;
				\definecolor{ezs}{RGB}{254,245,105}\fill[ezs](67,-184.5)--(69.3,-185.4)--(134.6,-15.2)--(132.3,-14.3)--cycle;
				\definecolor{ezs}{RGB}{254,244,103}\fill[ezs](68.9,-185.2)--(71.2,-186.1)--(136.5,-15.9)--(134.2,-15)--cycle;
				\definecolor{ezs}{RGB}{253,243,101}\fill[ezs](70.8,-186)--(73.1,-186.8)--(138.4,-16.6)--(136.1,-15.8)--cycle;
				\definecolor{ezs}{RGB}{253,241,98}\fill[ezs](72.7,-186.7)--(75,-187.6)--(140.3,-17.4)--(138,-16.5)--cycle;
				\definecolor{ezs}{RGB}{253,239,95}\fill[ezs](74.6,-187.4)--(76.9,-188.3)--(142.2,-18.1)--(140,-17.2)--cycle;
				\definecolor{ezs}{RGB}{253,236,92}\fill[ezs](76.5,-188.2)--(78.8,-189)--(144.2,-18.8)--(141.9,-18)--cycle;
				\definecolor{ezs}{RGB}{253,234,88}\fill[ezs](78.5,-188.9)--(80.7,-189.8)--(146.1,-19.6)--(143.8,-18.7)--cycle;
				\definecolor{ezs}{RGB}{253,232,85}\fill[ezs](80.4,-189.6)--(82.7,-190.5)--(148,-20.3)--(145.7,-19.4)--cycle;
				\definecolor{ezs}{RGB}{253,230,82}\fill[ezs](82.3,-190.4)--(84.6,-191.2)--(149.9,-21)--(147.6,-20.2)--cycle;
				\definecolor{ezs}{RGB}{253,228,78}\fill[ezs](84.2,-191.1)--(86.5,-192)--(151.8,-21.8)--(149.5,-20.9)--cycle;
				\definecolor{ezs}{RGB}{253,226,74}\fill[ezs](86.1,-191.8)--(88.4,-192.7)--(153.7,-22.5)--(151.4,-21.6)--cycle;
				\definecolor{ezs}{RGB}{253,224,70}\fill[ezs](88,-192.6)--(90.3,-193.4)--(155.6,-23.3)--(153.3,-22.4)--cycle;
				\definecolor{ezs}{RGB}{253,222,67}\fill[ezs](89.9,-193.3)--(92.2,-194.2)--(157.5,-24)--(155.3,-23.1)--cycle;
				\definecolor{ezs}{RGB}{253,220,65}\fill[ezs](91.8,-194)--(94.1,-194.9)--(159.5,-24.7)--(157.2,-23.8)--cycle;
				\definecolor{ezs}{RGB}{253,217,63}\fill[ezs](93.7,-194.8)--(96,-195.6)--(161.4,-25.5)--(159.1,-24.6)--cycle;
				\definecolor{ezs}{RGB}{253,215,60}\fill[ezs](95.7,-195.5)--(97.9,-196.4)--(163.3,-26.2)--(161,-25.3)--cycle;
				\definecolor{ezs}{RGB}{253,213,56}\fill[ezs](97.6,-196.2)--(99.9,-197.1)--(165.2,-26.9)--(162.9,-26)--cycle;
				\definecolor{ezs}{RGB}{252,211,53}\fill[ezs](99.5,-197)--(101.8,-197.8)--(167.1,-27.7)--(164.8,-26.8)--cycle;
				\definecolor{ezs}{RGB}{252,209,49}\fill[ezs](101.4,-197.7)--(103.7,-198.6)--(169,-28.4)--(166.7,-27.5)--cycle;
				\definecolor{ezs}{RGB}{252,207,45}\fill[ezs](103.3,-198.4)--(105.6,-199.3)--(170.9,-29.1)--(168.6,-28.2)--cycle;
				\definecolor{ezs}{RGB}{253,205,42}\fill[ezs](105.2,-199.2)--(107.5,-200)--(172.8,-29.9)--(170.6,-29)--cycle;
				\definecolor{ezs}{RGB}{253,202,40}\fill[ezs](107.1,-199.9)--(109.4,-200.8)--(174.7,-30.6)--(172.5,-29.7)--cycle;
				\definecolor{ezs}{RGB}{253,200,38}\fill[ezs](109,-200.6)--(111.3,-201.5)--(176.7,-31.3)--(174.4,-30.4)--cycle;
				\definecolor{ezs}{RGB}{253,198,35}\fill[ezs](111,-201.4)--(113.2,-202.2)--(178.6,-32.1)--(176.3,-31.2)--cycle;
				\definecolor{ezs}{RGB}{253,196,31}\fill[ezs](112.9,-202.1)--(115.2,-203)--(180.5,-32.8)--(178.2,-31.9)--cycle;
				\definecolor{ezs}{RGB}{253,194,28}\fill[ezs](114.8,-202.8)--(117.1,-203.7)--(182.4,-33.5)--(180.1,-32.6)--cycle;
				\definecolor{ezs}{RGB}{253,192,26}\fill[ezs](116.7,-203.6)--(119,-204.5)--(184.3,-34.3)--(182,-33.4)--cycle;
				\definecolor{ezs}{RGB}{253,189,24}\fill[ezs](118.6,-204.3)--(120.9,-205.2)--(186.2,-35)--(183.9,-34.1)--cycle;
				\definecolor{ezs}{RGB}{253,186,20}\fill[ezs](120.5,-205)--(122.8,-205.9)--(188.1,-35.7)--(185.8,-34.8)--cycle;
				\definecolor{ezs}{RGB}{253,184,17}\fill[ezs](122.4,-205.8)--(124.7,-206.7)--(190,-36.5)--(187.8,-35.6)--cycle;
				\definecolor{ezs}{RGB}{252,183,15}\fill[ezs](124.3,-206.5)--(126.6,-207.4)--(192,-37.2)--(189.7,-36.3)--cycle;
				\definecolor{ezs}{RGB}{252,181,14}\fill[ezs](126.3,-207.2)--(128.5,-208.1)--(193.9,-37.9)--(191.6,-37.1)--cycle;
				\definecolor{ezs}{RGB}{252,179,13}\fill[ezs](128.2,-208)--(130.4,-208.9)--(195.8,-38.7)--(193.5,-37.8)--cycle;
				\definecolor{ezs}{RGB}{252,178,12}\fill[ezs](130.1,-208.7)--(132.4,-209.6)--(197.7,-39.4)--(195.4,-38.5)--cycle;
				\definecolor{ezs}{RGB}{251,178,10}\fill[ezs](132,-209.4)--(134.3,-210.3)--(199.6,-40.1)--(197.3,-39.3)--cycle;
				\definecolor{ezs}{RGB}{251,178,10}\fill[ezs](133.9,-210.2)--(136.2,-211.1)--(201.5,-40.9)--(199.2,-40)--cycle;
				\definecolor{ezs}{RGB}{250,177,11}\fill[ezs](135.8,-210.9)--(138.1,-211.8)--(203.4,-41.6)--(201.1,-40.7)--cycle;
				\definecolor{ezs}{RGB}{250,178,11}\fill[ezs](137.7,-211.6)--(140,-212.5)--(205.3,-42.3)--(203.1,-41.5)--cycle;
				\definecolor{ezs}{RGB}{248,180,10}\fill[ezs](139.6,-212.4)--(141.9,-213.3)--(207.3,-43.1)--(205,-42.2)--cycle;
				\definecolor{ezs}{RGB}{247,181,9}\fill[ezs](141.5,-213.1)--(143.8,-214)--(209.2,-43.8)--(206.9,-42.9)--cycle;
				\definecolor{ezs}{RGB}{246,182,9}\fill[ezs](143.5,-213.8)--(145.7,-214.7)--(211.1,-44.5)--(208.8,-43.7)--cycle;
				\definecolor{ezs}{RGB}{246,182,10}\fill[ezs](145.4,-214.6)--(147.7,-215.5)--(213,-45.3)--(210.7,-44.4)--cycle;
			\end{scope}
			% chữ EDUZY (Times đậm, làm dày nét bằng 5 lớp lệch nhẹ; căn đúng khung chữ logo gốc)
			\foreach \dx/\dy in {-0.7/0,0.7/0,0/-0.3,0/0.3,0/0}
			\node[anchor=center,inner sep=0pt,outer sep=0pt,text=ezText,align=center] at (109.6+\dx,-193.6+\dy)
			{\resizebox{1.01cm}{0.22cm}{\fontencoding{T1}\fontfamily{ptm}\fontseries{b}\selectfont EDUZY}};
		\end{tikzpicture}%
	}%
}
%====================== THIẾT LẬP THÂN ĐỀ ======================
%--- Môi trường tự luận {bt}: ex_test 3.3.4 chưa khai báo sẵn nên tạo ở đây
%    (ex_test đã tự móc \fixshowans và \listenumerate cho bt)
\makeatletter
\@ifundefined{bt}{\newtheorem{bt}{\color{gold1}\namebt}}{}
\makeatother
\renewcommand{\sectionmark}[1]{\markright{\thesection\ #1}}
%--- Khung tiêu đề (logo + tối đa 4 dòng; dòng 3, 4 để trống {} thì bỏ qua)
%  Đề thi:    \khungtieude{THI THỬ THPT QUỐC GIA NĂM 2026}{Môn Toán -- Đề số 21}
%                         {(Đề thi gồm: ...)}{Thời gian làm bài: 90 phút (không kể thời gian phát đề).}
%  Chuyên đề: \khungtieude{CHUYÊN ĐỀ 3 -- DÃY SỐ VÀ CÔNG THỨC TRUY HỒI}
%                         {Bồi dưỡng học sinh giỏi Toán THPT}{...}{}
\newcommand{\khungtieude}[4]{%
	\begin{tcolorbox}[colframe=gold1,colback=gold5!45,boxrule=0.5mm]
		\begin{minipage}[c]{3cm}
			\centerline{\eduzylogo[2.4cm]}
		\end{minipage}\hspace{0.5cm}
		\begin{minipage}[c]{12.5cm}
			\vskip 0.15 cm
			\centerline{\ind{#1}}
			\vskip 0.15 cm
			\centerline{\textbf{#2}}
			\if\relax\detokenize{#3}\relax\else
				\vskip 0.15 cm
				\centerline{\textit{#3}}
			\fi
			\if\relax\detokenize{#4}\relax\else
				\vskip 0.15 cm
				\centerline{\textit{#4}}
			\fi
			\vskip 0.15 cm
		\end{minipage}
	\end{tcolorbox}
}
%====================== BỔ SUNG CHO BÀI HỌC, CHUYÊN ĐỀ ======================
\usepackage[shortlabels]{enumitem}
\usepackage{tabularx}
%--- Danh sách gọn: \begin{enumerate}[gon] ... \end{enumerate}
\SetEnumitemKey{gon}{leftmargin=*,itemsep=1pt,parsep=0pt,topsep=3pt}
%--- Khung lý thuyết: \begin{khung}{} ... \end{khung}; tiêu đề chỉ dùng khi thật ngắn ("Mẹo", "Lỗi thường gặp")
\newtcolorbox{khung}[1]{breakable,enhanced,colframe=gold1,colback=gold5!45,
	boxrule=0.4mm,fonttitle=\bfseries\sffamily,coltitle=white,colbacktitle=gold1,
	attach boxed title to top left={xshift=6mm,yshift=-3mm},
	boxed title style={arc=1mm,boxrule=0.3mm},title={#1}}
%--- Tiêu đề mục nhỏ trong một dạng: \muc{Dấu hiệu}, \muc{Phương pháp}, \muc{Bài tự luyện}
\makeatletter
\newcommand{\muc}[1]{\par\medskip\noindent\indam{#1}\par\nobreak\smallskip\@afterindentfalse\@afterheading}
\makeatother
%--- Ô công thức cần nhớ (căn giữa, viền vàng, vừa khít nội dung)
%     \begin{center}\congthuc{$Q(n)=(n-1)!$}\end{center}
\newtcbox{\congthuc}{on line,colframe=gold1,colback=gold5!45,boxrule=0.5pt,arc=2pt,
	left=8pt,right=8pt,top=4pt,bottom=4pt}
%--- Công tắc lời giải bài tự luyện (ví dụ mẫu luôn dùng \loigiai)
%     \loigiaiTL[{đáp số}]{lời giải đầy đủ}
%     \TLfull -> in lời giải đầy đủ;  \TLdapso -> chỉ in một dòng "Đáp số";  \TLan -> ẩn hoàn toàn
%     (\CheDo là dethi thì mọi lời giải và đáp số đều ẩn)
\makeatletter
\newif\ifTLhien   \TLhientrue
\newif\ifTLdapso  \TLdapsofalse
\newcommand{\TLfull}{\TLhientrue\TLdapsofalse}
\newcommand{\TLdapso}{\TLhientrue\TLdapsotrue}
\newcommand{\TLan}{\TLhienfalse}
\newcommand{\loigiaiTL}[2][]{%
	\ifTLhien
		\ifTLdapso
			\def\@tlarg{#1}%
			\ifx\@tlarg\@empty\else
				\ifdefstring{\kindOPTN}{dethi}{}{\par\noindent{\bfseries\color{gold1}Đáp số:} #1\par}%
			\fi
		\else
			\loigiai{#2}%
		\fi
	\fi}
\makeatother
\TLdapso% mặc định: chỉ in đáp số bài tự luyện
%====================== BẢNG ĐÚNG/SAI (\choiceTFt) ======================
%  Hàng tiêu đề nền nâu đồng chữ trắng (Phát biểu | Đ | S), các hàng tô xen kẽ kem / trắng,
%  đường kẻ vàng kim, nhãn a) b) c) d) màu nâu đồng. Bản có đáp án: cột được chọn có vòng tròn
%  nâu đồng với dấu tích trắng (câu đúng tích ở cột Đúng, câu sai tích ở cột Sai).
%  Tiêu đề ghi đủ "Đúng | Sai": \OPTN{viettat=0} sau \input (cột Đúng/Sai rộng ra, phát biểu hẹp lại).
%  Dùng bảng gốc của ex_test: \def\BangDSgoc{1} trước \input{preamble-eduzy-v2}.
\providecommand{\BangDSgoc}{0}
\ifnum\BangDSgoc=0
\usepackage{colortbl,pifont}
\colorlet{ezDSke}{gold2!60!gold1}% đường kẻ bảng (nét 0,4pt như ex_test để không đổi chiều cao bảng)
\colorlet{ezDSle}{gold5!60}% hàng a), c)
\colorlet{ezDSchan}{white}% hàng b), d)
%--- dấu chọn ở cột Đúng/Sai
%    (bề rộng 0, căn giữa ô: không làm đổi bề rộng cột của ex_test, không báo tràn lề)
\newcommand{\EZdauchon}{\makebox[0pt]{\tikz[baseline=-1ex]\node[circle,fill=gold1,inner sep=0pt,
	minimum size=1.2em,text=white,font=\small]{\ding{51}};}}
\def\tickT{\EZdauchon}\def\tickF{\EZdauchon}
\AtBeginDocument{\sbox{\ccTFTab}{\EZdauchon}}
%--- nhãn a) b) c) d) màu nâu đồng (ở bản lời giải, nhãn câu đúng vẫn đóng khung như ex_test)
\renewcommand{\FalseTF}{\stepcounter{dapan}\textbf{\color{gold1}\DapAnTF\sepTF}\ignorespaces}
%--- trước khi dựng bảng: kẻ màu vàng kim (bề rộng cột giữ nguyên như ex_test); sau bảng trả lại màu đen
\newcount\EZhangTF
\newcommand{\EZbatdauTF}{%
	\global\EZhangTF=0\relax
	\arrayrulecolor{ezDSke}}
\newcommand{\EZketthucTF}{\arrayrulecolor{black}}
%--- hàng tiêu đề nền nâu đồng, chữ trắng
\makeatletter
\newcommand{\EZo@TF}[1]{\centering\bfseries\color{white}#1}
\newcommand{\EZtieudeTF}{%
	\ifcase\cotDS
		\ifnum\cotSTT=1\relax
			\def\rowHeader{\hline\rowcolor{gold1}\EZo@TF{TT}&\EZo@TF{\phatbieu}\tabularnewline}%
		\else
			\def\rowHeader{\hline\rowcolor{gold1}\EZo@TF{\phatbieu}\tabularnewline}%
		\fi
	\or
		\ifnum\cotSTT=1\relax
			\def\rowHeader{\hline\rowcolor{gold1}\EZo@TF{TT}&\EZo@TF{\phatbieu}&\EZo@TF{\Dung}\tabularnewline}%
		\else
			\def\rowHeader{\hline\rowcolor{gold1}\EZo@TF{\phatbieu}&\EZo@TF{\Dung}\tabularnewline}%
		\fi
	\or
		\ifnum\cotSTT=1\relax
			\def\rowHeader{\hline\rowcolor{gold1}\EZo@TF{TT}&\EZo@TF{\phatbieu}&\EZo@TF{\Dung}&\EZo@TF{\Sai}\tabularnewline}%
		\else
			\def\rowHeader{\hline\rowcolor{gold1}\EZo@TF{\phatbieu}&\EZo@TF{\Dung}&\EZo@TF{\Sai}\tabularnewline}%
		\fi
	\fi}
%--- tô xen kẽ từng hàng (gọi ngay sau \hline đầu mỗi hàng)
\newcommand{\EZmauhangTF}{%
	\noalign{\global\advance\EZhangTF\@ne}%
	\ifodd\EZhangTF\rowcolor{ezDSle}\else\rowcolor{ezDSchan}\fi}
%--- gắn vào \choiceTFt / \writeChoiceTFt của ex_test
\xpretocmd{\choiceTFt}{\EZbatdauTF}{}{\typeout{[preamble-eduzy-v2] Không gắn được bảng Đúng/Sai (đầu bảng)}}
\xpatchcmd{\choiceTFt}{\ifnum\dongPB=0\def\rowHeader{}\fi}{\EZtieudeTF\ifnum\dongPB=0\def\rowHeader{}\fi}
	{}{\typeout{[preamble-eduzy-v2] Không gắn được bảng Đúng/Sai (hàng tiêu đề)}}
\xpatchcmd{\writeChoiceTFt}{\hline\xdef\dotEX{\sdotEX}}{\hline\EZmauhangTF\xdef\dotEX{\sdotEX}}
	{}{\typeout{[preamble-eduzy-v2] Không gắn được bảng Đúng/Sai (tô hàng)}}
\xpatchcmd{\writeChoiceTFt}{{\end{longtable}}}{{\end{longtable}\EZketthucTF}}
	{}{\typeout{[preamble-eduzy-v2] Không gắn được bảng Đúng/Sai (cuối longtable)}}
\xpatchcmd{\writeChoiceTFt}{{\end{xtabular}\end{center}\medskip\noindent}}{{\end{xtabular}\end{center}\EZketthucTF\medskip\noindent}}
	{}{\typeout{[preamble-eduzy-v2] Không gắn được bảng Đúng/Sai (cuối xtabular)}}
\fi
%=====================================================================
\renewcommand{\parallel}{\mathrel{/\!\!/}}
.
\providecommand{\BangDSgoc}{0}
\ifnum\BangDSgoc=0
\usepackage{colortbl,pifont}
\colorlet{ezDSke}{gold2!60!gold1}% đường kẻ bảng (nét 0,4pt như ex_test để không đổi chiều cao bảng)
\colorlet{ezDSle}{gold5!60}% hàng a), c)
\colorlet{ezDSchan}{white}% hàng b), d)
%--- dấu chọn ở cột Đúng/Sai
%    (bề rộng 0, căn giữa ô: không làm đổi bề rộng cột của ex_test, không báo tràn lề)
\newcommand{\EZdauchon}{\makebox[0pt]{\tikz[baseline=-1ex]\node[circle,fill=gold1,inner sep=0pt,
	minimum size=1.2em,text=white,font=\small]{\ding{51}};}}
\def\tickT{\EZdauchon}\def\tickF{\EZdauchon}
\AtBeginDocument{\sbox{\ccTFTab}{\EZdauchon}}
%--- nhãn a) b) c) d) màu nâu đồng (ở bản lời giải, nhãn câu đúng vẫn đóng khung như ex_test)
\renewcommand{\FalseTF}{\stepcounter{dapan}\textbf{\color{gold1}\DapAnTF\sepTF}\ignorespaces}
%--- trước khi dựng bảng: kẻ màu vàng kim (bề rộng cột giữ nguyên như ex_test); sau bảng trả lại màu đen
\newcount\EZhangTF
\newcommand{\EZbatdauTF}{%
	\global\EZhangTF=0\relax
	\arrayrulecolor{ezDSke}}
\newcommand{\EZketthucTF}{\arrayrulecolor{black}}
%--- hàng tiêu đề nền nâu đồng, chữ trắng
\makeatletter
\newcommand{\EZo@TF}[1]{\centering\bfseries\color{white}#1}
\newcommand{\EZtieudeTF}{%
	\ifcase\cotDS
		\ifnum\cotSTT=1\relax
			\def\rowHeader{\hline\rowcolor{gold1}\EZo@TF{TT}&\EZo@TF{\phatbieu}\tabularnewline}%
		\else
			\def\rowHeader{\hline\rowcolor{gold1}\EZo@TF{\phatbieu}\tabularnewline}%
		\fi
	\or
		\ifnum\cotSTT=1\relax
			\def\rowHeader{\hline\rowcolor{gold1}\EZo@TF{TT}&\EZo@TF{\phatbieu}&\EZo@TF{\Dung}\tabularnewline}%
		\else
			\def\rowHeader{\hline\rowcolor{gold1}\EZo@TF{\phatbieu}&\EZo@TF{\Dung}\tabularnewline}%
		\fi
	\or
		\ifnum\cotSTT=1\relax
			\def\rowHeader{\hline\rowcolor{gold1}\EZo@TF{TT}&\EZo@TF{\phatbieu}&\EZo@TF{\Dung}&\EZo@TF{\Sai}\tabularnewline}%
		\else
			\def\rowHeader{\hline\rowcolor{gold1}\EZo@TF{\phatbieu}&\EZo@TF{\Dung}&\EZo@TF{\Sai}\tabularnewline}%
		\fi
	\fi}
%--- tô xen kẽ từng hàng (gọi ngay sau \hline đầu mỗi hàng)
\newcommand{\EZmauhangTF}{%
	\noalign{\global\advance\EZhangTF\@ne}%
	\ifodd\EZhangTF\rowcolor{ezDSle}\else\rowcolor{ezDSchan}\fi}
%--- gắn vào \choiceTFt / \writeChoiceTFt của ex_test
\xpretocmd{\choiceTFt}{\EZbatdauTF}{}{\typeout{[preamble-eduzy-v2] Không gắn được bảng Đúng/Sai (đầu bảng)}}
\xpatchcmd{\choiceTFt}{\ifnum\dongPB=0\def\rowHeader{}\fi}{\EZtieudeTF\ifnum\dongPB=0\def\rowHeader{}\fi}
	{}{\typeout{[preamble-eduzy-v2] Không gắn được bảng Đúng/Sai (hàng tiêu đề)}}
\xpatchcmd{\writeChoiceTFt}{\hline\xdef\dotEX{\sdotEX}}{\hline\EZmauhangTF\xdef\dotEX{\sdotEX}}
	{}{\typeout{[preamble-eduzy-v2] Không gắn được bảng Đúng/Sai (tô hàng)}}
\xpatchcmd{\writeChoiceTFt}{{\end{longtable}}}{{\end{longtable}\EZketthucTF}}
	{}{\typeout{[preamble-eduzy-v2] Không gắn được bảng Đúng/Sai (cuối longtable)}}
\xpatchcmd{\writeChoiceTFt}{{\end{xtabular}\end{center}\medskip\noindent}}{{\end{xtabular}\end{center}\EZketthucTF\medskip\noindent}}
	{}{\typeout{[preamble-eduzy-v2] Không gắn được bảng Đúng/Sai (cuối xtabular)}}
\fi
%=====================================================================
\renewcommand{\parallel}{\mathrel{/\!\!/}}
.
\providecommand{\BangDSgoc}{0}
\ifnum\BangDSgoc=0
\usepackage{colortbl,pifont}
\colorlet{ezDSke}{gold2!60!gold1}% đường kẻ bảng (nét 0,4pt như ex_test để không đổi chiều cao bảng)
\colorlet{ezDSle}{gold5!60}% hàng a), c)
\colorlet{ezDSchan}{white}% hàng b), d)
%--- dấu chọn ở cột Đúng/Sai
%    (bề rộng 0, căn giữa ô: không làm đổi bề rộng cột của ex_test, không báo tràn lề)
\newcommand{\EZdauchon}{\makebox[0pt]{\tikz[baseline=-1ex]\node[circle,fill=gold1,inner sep=0pt,
	minimum size=1.2em,text=white,font=\small]{\ding{51}};}}
\def\tickT{\EZdauchon}\def\tickF{\EZdauchon}
\AtBeginDocument{\sbox{\ccTFTab}{\EZdauchon}}
%--- nhãn a) b) c) d) màu nâu đồng (ở bản lời giải, nhãn câu đúng vẫn đóng khung như ex_test)
\renewcommand{\FalseTF}{\stepcounter{dapan}\textbf{\color{gold1}\DapAnTF\sepTF}\ignorespaces}
%--- trước khi dựng bảng: kẻ màu vàng kim (bề rộng cột giữ nguyên như ex_test); sau bảng trả lại màu đen
\newcount\EZhangTF
\newcommand{\EZbatdauTF}{%
	\global\EZhangTF=0\relax
	\arrayrulecolor{ezDSke}}
\newcommand{\EZketthucTF}{\arrayrulecolor{black}}
%--- hàng tiêu đề nền nâu đồng, chữ trắng
\makeatletter
\newcommand{\EZo@TF}[1]{\centering\bfseries\color{white}#1}
\newcommand{\EZtieudeTF}{%
	\ifcase\cotDS
		\ifnum\cotSTT=1\relax
			\def\rowHeader{\hline\rowcolor{gold1}\EZo@TF{TT}&\EZo@TF{\phatbieu}\tabularnewline}%
		\else
			\def\rowHeader{\hline\rowcolor{gold1}\EZo@TF{\phatbieu}\tabularnewline}%
		\fi
	\or
		\ifnum\cotSTT=1\relax
			\def\rowHeader{\hline\rowcolor{gold1}\EZo@TF{TT}&\EZo@TF{\phatbieu}&\EZo@TF{\Dung}\tabularnewline}%
		\else
			\def\rowHeader{\hline\rowcolor{gold1}\EZo@TF{\phatbieu}&\EZo@TF{\Dung}\tabularnewline}%
		\fi
	\or
		\ifnum\cotSTT=1\relax
			\def\rowHeader{\hline\rowcolor{gold1}\EZo@TF{TT}&\EZo@TF{\phatbieu}&\EZo@TF{\Dung}&\EZo@TF{\Sai}\tabularnewline}%
		\else
			\def\rowHeader{\hline\rowcolor{gold1}\EZo@TF{\phatbieu}&\EZo@TF{\Dung}&\EZo@TF{\Sai}\tabularnewline}%
		\fi
	\fi}
%--- tô xen kẽ từng hàng (gọi ngay sau \hline đầu mỗi hàng)
\newcommand{\EZmauhangTF}{%
	\noalign{\global\advance\EZhangTF\@ne}%
	\ifodd\EZhangTF\rowcolor{ezDSle}\else\rowcolor{ezDSchan}\fi}
%--- gắn vào \choiceTFt / \writeChoiceTFt của ex_test
\xpretocmd{\choiceTFt}{\EZbatdauTF}{}{\typeout{[preamble-eduzy-v2] Không gắn được bảng Đúng/Sai (đầu bảng)}}
\xpatchcmd{\choiceTFt}{\ifnum\dongPB=0\def\rowHeader{}\fi}{\EZtieudeTF\ifnum\dongPB=0\def\rowHeader{}\fi}
	{}{\typeout{[preamble-eduzy-v2] Không gắn được bảng Đúng/Sai (hàng tiêu đề)}}
\xpatchcmd{\writeChoiceTFt}{\hline\xdef\dotEX{\sdotEX}}{\hline\EZmauhangTF\xdef\dotEX{\sdotEX}}
	{}{\typeout{[preamble-eduzy-v2] Không gắn được bảng Đúng/Sai (tô hàng)}}
\xpatchcmd{\writeChoiceTFt}{{\end{longtable}}}{{\end{longtable}\EZketthucTF}}
	{}{\typeout{[preamble-eduzy-v2] Không gắn được bảng Đúng/Sai (cuối longtable)}}
\xpatchcmd{\writeChoiceTFt}{{\end{xtabular}\end{center}\medskip\noindent}}{{\end{xtabular}\end{center}\EZketthucTF\medskip\noindent}}
	{}{\typeout{[preamble-eduzy-v2] Không gắn được bảng Đúng/Sai (cuối xtabular)}}
\fi
%=====================================================================
\renewcommand{\parallel}{\mathrel{/\!\!/}}

%=====================================================================
% preamble-vatli-hinh.tex -- BỘ MACRO HÌNH VẬT LÍ (dùng chung cho mọi đề Vật lí)
% Gọi ngay sau preamble chính:
%     %============================================================
%  PREAMBLE EDUZY v2 - dùng chung cho đề thi, bài học, chuyên đề (gói ex_test)
%  Khác bản preamble-eduzy.tex ở đầu trang, chân trang: dải nền nâu đồng tràn ngang,
%  logo + Eduzy.com.vn, ô chữ vàng kem cạnh xiên, chân trang "Trang (số)";
%  logo Eduzy chìm giữa mọi trang (tắt: \logochimtat, bật lại: \logochimbat).
%  Cách dùng ở file chính (xem mau-chuan.tex, mau-chuyen-de.tex):
%      \def\CheDo{loigiai}          % loigiai | dethi | solcolor | color | book
%      \def\tenchuong{<dòng chữ trên đầu trang>}
%      \def\thuonghieu{<chữ ở chân trang lẻ>}   % không khai báo thì dùng mặc định
%      \def\tacgia{<chữ ở chân trang chẵn>}      % không khai báo thì dùng mặc định
%      \def\khauhieu{<chữ cạnh logo ở đầu trang>}  % tùy chọn, mặc định Eduzy.com.vn
%      %============================================================
%  PREAMBLE EDUZY v2 - dùng chung cho đề thi, bài học, chuyên đề (gói ex_test)
%  Khác bản preamble-eduzy.tex ở đầu trang, chân trang: dải nền nâu đồng tràn ngang,
%  logo + Eduzy.com.vn, ô chữ vàng kem cạnh xiên, chân trang "Trang (số)";
%  logo Eduzy chìm giữa mọi trang (tắt: \logochimtat, bật lại: \logochimbat).
%  Cách dùng ở file chính (xem mau-chuan.tex, mau-chuyen-de.tex):
%      \def\CheDo{loigiai}          % loigiai | dethi | solcolor | color | book
%      \def\tenchuong{<dòng chữ trên đầu trang>}
%      \def\thuonghieu{<chữ ở chân trang lẻ>}   % không khai báo thì dùng mặc định
%      \def\tacgia{<chữ ở chân trang chẵn>}      % không khai báo thì dùng mặc định
%      \def\khauhieu{<chữ cạnh logo ở đầu trang>}  % tùy chọn, mặc định Eduzy.com.vn
%      \input{preamble-eduzy-v2}
%      \TLdapso                      % tài liệu: \TLfull | \TLdapso | \TLan
%      \begin{document} ... \end{document}
%  Logo Eduzy vẽ bằng TikZ ngay trong file (\eduzylogo), không cần file ảnh.
%  Biên dịch: pdfLaTeX (máy có gói vietnam) hoặc XeLaTeX.
%============================================================
\documentclass[12pt,a4paper,twoside]{article}%<-- twoside: bat dau trang/chan trang chan-le
\usepackage{iftex}
\usepackage{amsmath,amssymb,tikz,mathrsfs,tkz-tab}
\usepackage{graphicx}
\usepackage{pgfplots}% đồ thị hàm số
\pgfplotsset{compat=1.17}
%\usepackage{mathptmx}
\usepackage{mathpazo}
\ifXeTeX% XeLaTeX: font chữ có đủ dấu tiếng Việt, giữ nguyên font toán mathpazo
	\usepackage[no-math]{fontspec}
	\setmainfont{TeX Gyre Pagella}
	\newfontfamily\EZqag{TeX Gyre Adventor}[NFSSFamily=qag,Ligatures=TeX]
\fi
\usepackage[many]{tcolorbox}
\usetikzlibrary{math,through,calc,intersections,angles,quotes,shapes,shapes.geometric,arrows,patterns,matrix,chains,arrows.meta,decorations.shapes,decorations.fractals,decorations.markings,shadows,decorations.pathreplacing}
\usepackage[top=1.6cm, bottom=1.8cm, left=1.8cm, right=1.8cm]{geometry}% v2: top 1.6cm (bản cũ 1.5cm) để chữ không sát dải đầu trang
\providecommand{\CheDo}{loigiai}% loigiai | dethi | solcolor | color | book
\expandafter\usepackage\expandafter[\CheDo]{ex_test}
% \tenchuong: khai báo ở file chính (trước \input) hoặc \renewcommand sau \input
%============================================================
%  [GỘP TỪ eduzy-de.tex] EDUZY - MÀU NÂU ĐỒNG / VÀNG KIM
%  + ĐẦU TRANG, CHÂN TRANG KIỂU DẢI NGANG (v2), ĐỐI XỨNG TRANG CHẴN / TRANG LẺ.
%  (không cần file eduzy-de.tex riêng nữa)
%  Khối này phải đặt SAU \usepackage[...]{ex_test} và SAU \def\tenchuong{...}
%============================================================
\usepackage{lastpage}
\usepackage{needspace}
%--- DÀN TRANG ĐỀ THI: tách câu (THÊM 05/11/2026; chỉ thêm lệnh mới, KHÔNG đổi hành vi mặc định) ---
%  Hai công tắc, mặc định tắt (= preamble chạy y như cũ). Bật bằng cách đặt TRƯỚC \input{preamble-eduzy-v2}
%  trong file đề (tools/dan_trang/dan_trang.py tự ghi khi đặt \NeedspaceCau):
%    \def\TachCauDS{1}   câu Đúng/Sai được tách: câu dẫn (kèm hình) ở cuối trang, bảng Phát biểu sang trang sau
%                        (câu dẫn và hình không tách nhau, bảng không tách giữa hai trang)
%    \def\TachCauLG{1}   bản lời giải (\CheDo khác dethi): phần đề của câu cao hơn 1/3 trang được phép tách
%  Lệnh dùng trong file đề (giá trị do tools/dan_trang/dan_trang.py đo ở bản dethi):
%    \NeedspaceCau[<cao câu dẫn + hình>]{<cao cả khối câu>}   ngay trên \begin{ex} (câu đầu phần: trên tiêu đề phần)
%        công tắc tắt: đúng bằng \Needspace{<cao cả khối câu>} như cũ;
%        \TachCauDS=1 và có [..] (câu Đúng/Sai được chọn tách): chỉ giữ câu dẫn + hình;
%        \TachCauLG=1 ở bản lời giải: đòi tối đa 1/3 \textheight.
%    \NeedspaceBang{<cao bảng Phát biểu>}%   ngay trước \choiceTFt: \TachCauDS=1 thì giữ trọn bảng trên một trang;
%        công tắc tắt: không làm gì.
\providecommand{\TachCauDS}{0}
\providecommand{\TachCauLG}{0}
\newlength{\EZnsCau}
\newcommand{\NeedspaceCau}[2][]{%
	\setlength{\EZnsCau}{#2}%
	\ifnum\TachCauDS=1 \if\relax\detokenize{#1}\relax\else\setlength{\EZnsCau}{#1}\fi\fi
	\ifnum\TachCauLG=1 \ifdefstring{\kindOPTN}{dethi}{}{\ifdim\EZnsCau>0.3333\textheight\setlength{\EZnsCau}{0.3333\textheight}\fi}\fi
	\Needspace{\EZnsCau}}
\newcommand{\NeedspaceBang}[1]{\ifnum\TachCauDS=1 \Needspace{#1}\fi}

%--- Thang màu: đậm -> nhạt
\definecolor{gold1}{RGB}{126,96,0}    % nâu đồng đậm  #7E6000
\definecolor{gold2}{RGB}{192,144,3}   % vàng kim      #C09003
\definecolor{gold3}{RGB}{255,215,99}  % vàng nhạt     #FFD763
\definecolor{gold4}{RGB}{255,231,153} % vàng rất nhạt #FFE799
\definecolor{gold5}{RGB}{255,243,205} % vàng kem      #FFF3CD
%--- Ghi đè tông xanh/đỏ cũ
\definecolor{myblue}{RGB}{126,96,0}
\colorlet{blue}{gold1}
\colorlet{cyan}{gold2}
\colorlet{teal}{gold2}

%--- Tên chương mặc định (file đề nên tự khai báo lại)
\providecommand{\tenchuong}{}
%--- Chân trang
\providecommand{\thuonghieu}{TOÁN 12 -- Eduzy.com.vn}
\providecommand{\tacgia}{Tô Gia Bảo -- 0353703783}

%--- Chữ cạnh logo ở đầu trang
\providecommand{\khauhieu}{Eduzy.com.vn}

%====================== THÔNG SỐ ĐẦU TRANG, CHÂN TRANG (v2) ======================
%  Dải nền đậm tràn ngang trang; ô chữ vàng kem có cạnh xiên, kèm 2 vạch xiên chuyển màu.
%  Đầu trang lẻ : logo + \khauhieu (trái)   | ô chữ \tenchuong (phải)
%  Chân trang lẻ: ô chữ \thuonghieu (trái)  | Trang (số) (phải)
%  Trang chẵn soi gương trang lẻ; chân trang chẵn ghi \tacgia.
%--- Màu: chỉ dùng thang vàng kim Eduzy
\colorlet{ezNen}{gold1}% dải nền đậm
\colorlet{ezVach1}{gold3}% vạch xiên sát dải nền (đậm hơn)
\colorlet{ezVach2}{gold4}% vạch xiên sát ô chữ (nhạt hơn)
\colorlet{ezO}{gold5}% ô chữ vàng kem
%--- Kích thước
\newlength{\EZdai}\setlength{\EZdai}{1.1cm}% chiều cao dải nền
\newlength{\EZo}\setlength{\EZo}{0.68cm}% chiều cao ô chữ
\newlength{\EZxien}\setlength{\EZxien}{0.55cm}% độ lệch ngang của cạnh xiên
\newlength{\EZvach}\setlength{\EZvach}{0.75cm}% bề rộng mỗi vạch xiên
\newlength{\EZkhe}\setlength{\EZkhe}{0.12cm}% khe trắng giữa vạch và ô chữ
\newlength{\EZdem}\setlength{\EZdem}{0.4cm}% khoảng đệm từ chữ tới cạnh xiên
\newlength{\EZle}\setlength{\EZle}{1.8cm}% lề trái/phải của chữ (bằng lề thân bài)
\newlength{\EZrongmin}\setlength{\EZrongmin}{4.5cm}% ô chữ rộng tối thiểu
\newlength{\EZrong}% bề rộng chữ trong ô (tính tự động)
%--- Font
\newcommand{\EZfontO}{\fontsize{10.5pt}{12pt}\selectfont\bfseries}% chữ trong ô
\newcommand{\EZfontKH}{\fontfamily{qag}\fontsize{11pt}{12pt}\selectfont\bfseries}% \khauhieu cạnh logo
\newcommand{\EZfontTrang}{\fontsize{11.5pt}{12pt}\selectfont\bfseries}% chữ "Trang" và số trang
%--- Đo bề rộng chữ trong ô
\newcommand{\EZdorong}[1]{%
	\settowidth{\EZrong}{\EZfontO #1}%
	\ifdim\EZrong<\EZrongmin\setlength{\EZrong}{\EZrongmin}\fi}
%--- Logo + chữ trắng \khauhieu (dựng sẵn một lần khi vào \begin{document})
\newsavebox{\EZlogobox}
\AtBeginDocument{\sbox{\EZlogobox}{%
		\raisebox{-0.5\height}{\eduzylogo[1cm]}\hspace{0.15cm}%
		\raisebox{-0.5\height}{\color{white}\EZfontKH\khauhieu}}}
%--- Hai vạch xiên + khe trắng. #1: điểm gốc ở mép dưới, #2: hướng (1 sang phải, -1 sang trái),
%    #3: hướng nghiêng (1 kiểu "/", -1 kiểu "\")
\newcommand{\EZvachxien}[3]{%
	\fill[white] (#1) --++({(#3)*\EZxien},\EZo) --++({(#2)*\EZkhe},0) --++({-(#3)*\EZxien},-\EZo) --cycle;
	\fill[ezVach2] ([xshift=(#2)*\EZkhe]#1) --++({(#3)*\EZxien},\EZo) --++({(#2)*\EZvach},0) --++({-(#3)*\EZxien},-\EZo) --cycle;
	\fill[ezVach1] ([xshift=(#2)*(\EZkhe+\EZvach)]#1) --++({(#3)*\EZxien},\EZo) --++({(#2)*\EZvach},0) --++({-(#3)*\EZxien},-\EZo) --cycle;
	\draw[white,line width=0.6pt] ([xshift=(#2)*(\EZkhe+\EZvach)]#1) --++({(#3)*\EZxien},\EZo);
}
%--- Chữ trong ô (căn giữa theo chiều cao ô, không lệch vì dấu tiếng Việt)
\tikzset{ezchu/.style={inner sep=0pt,text=ezNen,text height=0.72em,text depth=0.22em}}

%====================== LOGO CHÌM GIỮA TRANG ======================
%  Vẽ cùng đầu trang nên nằm dưới chữ của thân bài, có trên mọi trang.
%  Tắt: \logochimtat (từ trang đó trở đi), bật lại: \logochimbat.
%  Đổi cỡ, độ đậm: khai báo trước \input, ví dụ \def\EZchimcao{9cm} \def\EZchimdo{0.1}
\providecommand{\EZchimcao}{10cm}% chiều cao logo chìm
\providecommand{\EZchimdo}{0.12}% độ đậm (0 = trong suốt, 1 = rõ hẳn)
\newif\ifEZchim \EZchimtrue
\newcommand{\logochimtat}{\EZchimfalse}
\newcommand{\logochimbat}{\EZchimtrue}
%  Dựng sẵn logo mờ một lần khi vào \begin{document}; transparency group làm mờ cả logo
%  một lượt nên không lộ vệt ở chỗ các dải màu chồng nhau.
\newsavebox{\EZchimtho}
\newsavebox{\EZchimbox}
\AtBeginDocument{%
	\sbox{\EZchimtho}{\eduzylogo[\EZchimcao]}%
	\sbox{\EZchimbox}{\begin{tikzpicture}
			\begin{scope}[transparency group,opacity=\EZchimdo]
				\node[inner sep=0pt] {\usebox{\EZchimtho}};
			\end{scope}
	\end{tikzpicture}}}
\newcommand{\EZvechim}{%
	\ifEZchim
		\node[inner sep=0pt] at (current page.center) {\usebox{\EZchimbox}};
	\fi}

%====================== ĐẦU TRANG - TRANG LẺ ======================
%  (logo bên trái, ô chữ \tenchuong bên phải)
\newcommand{\myfancyheadOdd}{%
	\EZdorong{\tenchuong}%
	\begin{tikzpicture}[remember picture,overlay]
		\coordinate (T) at (current page.north west);
		\coordinate (P) at (current page.north east);
		%--- logo chìm giữa trang
		\EZvechim
		%--- dải nền
		\fill[ezNen] (T) rectangle ([yshift=-\EZdai]P);
		%--- ô chữ: X là góc dưới bên trái, cạnh trái xiên kiểu "/"
		\coordinate (X) at ([xshift=-\EZle-\EZrong-\EZdem-\EZxien,yshift=-\EZdai]P);
		\fill[ezO] (X) --++(\EZxien,\EZo) -- ([yshift=-\EZdai+\EZo]P) -- ([yshift=-\EZdai]P) --cycle;
		\EZvachxien{X}{-1}{1}
		\node[ezchu,anchor=east] at ([xshift=-\EZle,yshift=-\EZdai+0.5*\EZo]P) {\EZfontO\tenchuong};
		%--- logo + tên
		\node[inner sep=0pt,anchor=west] at ([xshift=\EZle,yshift=-0.5*\EZdai]T) {\usebox{\EZlogobox}};
	\end{tikzpicture}%
}

%====================== ĐẦU TRANG - TRANG CHẴN =====================
%  (soi gương: ô chữ \tenchuong bên trái, logo bên phải)
\newcommand{\myfancyheadEven}{%
	\EZdorong{\tenchuong}%
	\begin{tikzpicture}[remember picture,overlay]
		\coordinate (T) at (current page.north west);
		\coordinate (P) at (current page.north east);
		%--- logo chìm giữa trang
		\EZvechim
		%--- dải nền
		\fill[ezNen] (T) rectangle ([yshift=-\EZdai]P);
		%--- ô chữ: X là góc dưới bên phải, cạnh phải xiên kiểu "\"
		\coordinate (X) at ([xshift=\EZle+\EZrong+\EZdem+\EZxien,yshift=-\EZdai]T);
		\fill[ezO] (X) --++(-\EZxien,\EZo) -- ([yshift=-\EZdai+\EZo]T) -- ([yshift=-\EZdai]T) --cycle;
		\EZvachxien{X}{1}{-1}
		\node[ezchu,anchor=west] at ([xshift=\EZle,yshift=-\EZdai+0.5*\EZo]T) {\EZfontO\tenchuong};
		%--- logo + tên
		\node[inner sep=0pt,anchor=east] at ([xshift=-\EZle,yshift=-0.5*\EZdai]P) {\usebox{\EZlogobox}};
	\end{tikzpicture}%
}

%====================== CHÂN TRANG - TRANG LẺ ======================
%  (ô chữ \thuonghieu bên trái, "Trang (số)" bên phải)
\newcommand{\myfancyfootOdd}{%
	\EZdorong{\thuonghieu}%
	\begin{tikzpicture}[remember picture,overlay]
		\coordinate (T) at (current page.south west);
		\coordinate (P) at (current page.south east);
		%--- dải nền
		\fill[ezNen] (T) rectangle ([yshift=\EZdai]P);
		%--- ô chữ sát mép trên dải: X là góc dưới bên phải, cạnh phải xiên kiểu "/"
		\coordinate (X) at ([xshift=\EZle+\EZrong+\EZdem,yshift=\EZdai-\EZo]T);
		\fill[ezO] ([yshift=\EZdai-\EZo]T) -- (X) --++(\EZxien,\EZo) -- ([yshift=\EZdai]T) --cycle;
		\EZvachxien{X}{1}{1}
		\node[ezchu,anchor=west] at ([xshift=\EZle,yshift=\EZdai-0.5*\EZo]T) {\EZfontO\thuonghieu};
		%--- số trang
		\node[circle,draw=white,line width=0.9pt,text=white,inner sep=0pt,minimum size=0.62cm,anchor=east]
		(so) at ([xshift=-\EZle,yshift=0.5*\EZdai]P) {\EZfontTrang\thepage};
		\node[inner sep=0pt,text=white,anchor=base east] (tr) at ([xshift=-0.15cm]so.base west) {\EZfontTrang Trang};
	\end{tikzpicture}%
}

%====================== CHÂN TRANG - TRANG CHẴN ====================
%  (soi gương: "Trang (số)" bên trái, ô chữ \tacgia bên phải)
\newcommand{\myfancyfootEven}{%
	\EZdorong{\tacgia}%
	\begin{tikzpicture}[remember picture,overlay]
		\coordinate (T) at (current page.south west);
		\coordinate (P) at (current page.south east);
		%--- dải nền
		\fill[ezNen] (T) rectangle ([yshift=\EZdai]P);
		%--- ô chữ sát mép trên dải: X là góc dưới bên trái, cạnh trái xiên kiểu "\"
		\coordinate (X) at ([xshift=-\EZle-\EZrong-\EZdem,yshift=\EZdai-\EZo]P);
		\fill[ezO] ([yshift=\EZdai-\EZo]P) -- (X) --++(-\EZxien,\EZo) -- ([yshift=\EZdai]P) --cycle;
		\EZvachxien{X}{-1}{-1}
		\node[ezchu,anchor=east] at ([xshift=-\EZle,yshift=\EZdai-0.5*\EZo]P) {\EZfontO\tacgia};
		%--- số trang
		%    (vòng tròn ẩn so0 nằm giữa dải để lấy đường chân chữ, giống trang lẻ)
		\node[circle,inner sep=0pt,minimum size=0.62cm] (so0) at ([yshift=0.5*\EZdai]T) {\EZfontTrang\phantom{\thepage}};
		\node[inner sep=0pt,text=white,anchor=base west] (tr) at ($(so0.base -| T)+(\EZle,0)$) {\EZfontTrang Trang};
		\node[circle,draw=white,line width=0.9pt,text=white,inner sep=0pt,minimum size=0.62cm,anchor=base west]
		(so) at ([xshift=0.15cm]tr.base east) {\EZfontTrang\thepage};
	\end{tikzpicture}%
}

%--- Giữ tương thích với các file cũ đang gọi \myfancyhead, \myfancyfoot
\newcommand{\myfancyhead}{\myfancyheadOdd}
\newcommand{\myfancyfoot}{\myfancyfootOdd}

%======================= GẮN VÀO TRANG =============================
\pagestyle{fancy}
\fancyhf{}
\fancyhead[LO]{\myfancyheadOdd}
\fancyhead[RE]{\myfancyheadEven}
\fancyfoot[LO]{\myfancyfootOdd}
\fancyfoot[RE]{\myfancyfootEven}
\renewcommand{\headrulewidth}{0pt}
\renewcommand{\footrulewidth}{0pt}
\fancypagestyle{plain}{%
	\fancyhf{}
	\fancyhead[LO]{\myfancyheadOdd}
	\fancyhead[RE]{\myfancyheadEven}
	\fancyfoot[LO]{\myfancyfootOdd}
	\fancyfoot[RE]{\myfancyfootEven}
	\renewcommand{\headrulewidth}{0pt}
	\renewcommand{\footrulewidth}{0pt}
}
%=================== HẾT PHẦN GỘP TỪ eduzy-de.tex ===================
\newcommand{\indam}[1]{
	{\fontsize{12pt}{9pt}\sffamily \selectfont\bfseries\color{gold1} #1}
}
\newcommand{\ind}[1]{
	{\fontsize{14pt}{9pt}\fontfamily{qag} \selectfont\bfseries\color{gold1} #1}
}
\renewcommand{\FalseEX}{\stepcounter{dapan}\textbf{\selectfont\color{gold1}\Alph{dapan}.}\ignorespaces}
\renewtheorem{ex}{\color{gold1}Câu}
\renewcommand{\baselinestretch}{1.1}
\newcommand{\hoac}[1]{
	\left[\begin{aligned}#1\end{aligned}\right.}
\newcommand{\heva}[1]{
	\left\{\begin{aligned}#1\end{aligned}\right.}
\usepackage{esvect}
\def\vec{\vv}
\def\vt{\vv}
\def\overrightarrow{\vv}
%====================== LOGO EDUZY VẼ BẰNG TIKZ ======================
%  Đường viền dò từ logo gốc; dải màu vàng kim đo trực tiếp trên ảnh gốc
%  (nghiêng 21 độ, tô bằng 80 dải hẹp nên hiển thị giống nhau trên mọi trình xem PDF).
%  Dùng: \eduzylogo          -> cao 2.4cm (mặc định)
%        \eduzylogo[1.5cm]   -> cao 1.5cm
\definecolor{ezText}{RGB}{245,200,46}
\newcommand{\eduzylogo}[1][2.4cm]{%
	\resizebox{!}{#1}{%
		\begin{tikzpicture}[x=0.01cm,y=0.01cm]
			\path[use as bounding box] (0,0) rectangle (224,-224);
			% mũ tốt nghiệp + vương miện
			\begin{scope}
				\clip
				(119.4,-176.1)..controls(118.9,-176.1)and(116.8,-176)..(114.6,-176)
				..controls(101.4,-176)and(85.8,-175.9)..(83.7,-175.7)
				..controls(81.8,-175.6)and(81,-175.6)..(80.2,-175.7)
				..controls(78.8,-175.9)and(77.7,-175.9)..(76.2,-175.7)
				..controls(74.9,-175.5)and(70.8,-175.5)..(64.3,-175.7)
				..controls(59.1,-175.9)and(58.3,-175.7)..(58.6,-173.7)
				..controls(58.6,-172.8)and(59.6,-169.6)..(60.3,-168)
				..controls(60.8,-166.9)and(61.4,-165.8)..(62.3,-165)
				..controls(63.1,-164.2)and(64.3,-163.4)..(65,-163.2)
				..controls(65.5,-163.1)and(72.3,-163)..(76.6,-163.1)
				..controls(77.9,-163.1)and(80.1,-163.1)..(81.4,-163)
				..controls(82.8,-162.9)and(86.3,-162.9)..(89.2,-162.9)
				..controls(93,-162.9)and(94.9,-162.8)..(96,-162.7)
				..controls(97.2,-162.5)and(101,-162.5)..(114,-162.4)
				..controls(123,-162.3)and(134.9,-162.2)..(140.2,-162.1)
				..controls(150.1,-162)and(152.2,-162.1)..(153.6,-162.4)
				..controls(154.4,-162.6)and(157.1,-163.8)..(158.1,-164.4)
				..controls(158.8,-164.8)and(159.7,-165.7)..(159.9,-166.3)
				..controls(160,-166.5)and(160.2,-167.1)..(160.2,-167.6)
				..controls(160.3,-168)and(160.6,-169.4)..(160.9,-170.6)
				..controls(161.2,-171.7)and(161.5,-173.1)..(161.6,-173.6)
				..controls(161.8,-174.4)and(161.8,-174.5)..(161.5,-174.9)
				..controls(161.2,-175.6)and(160.8,-175.9)..(159.9,-176)
				..controls(159.2,-176.1)and(154.1,-176.2)..(151.5,-176.2)
				..controls(148.1,-176.1)and(122.6,-176.1)..(121.6,-176.2)
				..controls(120.8,-176.2)and(119.8,-176.2)..(119.4,-176.1)--cycle
				(130.4,-160)..controls(129,-160)and(119.8,-160)..(110.1,-159.9)
				..controls(100.3,-159.8)and(88.5,-159.7)..(83.9,-159.6)
				..controls(77.8,-159.5)and(75.3,-159.5)..(74.5,-159.6)
				..controls(73.9,-159.7)and(72.5,-159.8)..(71.3,-159.8)
				..controls(69.6,-159.8)and(69,-159.8)..(68.5,-159.6)
				..controls(67.6,-159.3)and(66.8,-158.6)..(66.2,-157.6)
				..controls(66,-157.3)and(65.6,-156.5)..(65.3,-156.1)
				..controls(64.6,-155.1)and(64.5,-154.9)..(63.7,-152.9)
				..controls(63.2,-151.6)and(62.9,-150.9)..(61.8,-149.6)
				..controls(61.5,-149.2)and(61.1,-148.5)..(60.9,-147.9)
				..controls(60.7,-147.4)and(60.3,-146.6)..(60.1,-146.1)
				..controls(59.9,-145.5)and(59.5,-144.9)..(59,-144.3)
				..controls(58.5,-143.7)and(58.2,-143.2)..(57.9,-142.5)
				..controls(57.7,-141.9)and(57.3,-141.1)..(57,-140.7)
				..controls(56.6,-140.3)and(56.1,-139.3)..(55.8,-138.5)
				..controls(55.5,-137.8)and(55.1,-136.8)..(54.9,-136.5)
				..controls(54.7,-136.1)and(54.3,-135.2)..(53.9,-134.5)
				..controls(53.6,-133.8)and(53,-132.7)..(52.6,-132)
				..controls(51.6,-130.5)and(51.4,-130.2)..(51.1,-129)
				..controls(50.6,-127.5)and(50.6,-127.4)..(49.7,-126.3)
				..controls(49.2,-125.7)and(48.6,-124.8)..(48.3,-124.2)
				..controls(48,-123.6)and(47.5,-122.7)..(47.3,-122.2)
				..controls(47,-121.8)and(46.5,-120.7)..(46.1,-119.8)
				..controls(45.8,-118.9)and(45.2,-117.8)..(44.9,-117.2)
				..controls(44.2,-116.1)and(43.3,-114.4)..(42.3,-112.2)
				..controls(41.8,-111.3)and(41.3,-110.4)..(41.1,-110.1)
				..controls(40.5,-109.4)and(40.3,-109)..(40,-107.8)
				..controls(39.8,-107)and(39.6,-106.5)..(39.2,-106)
				..controls(38.9,-105.6)and(38.4,-104.7)..(38,-104)
				..controls(37.7,-103.4)and(37.2,-102.5)..(37,-102.1)
				..controls(36.7,-101.7)and(36.3,-100.8)..(36,-100.1)
				..controls(35.6,-98.9)and(35.3,-98.4)..(34.6,-97.6)
				..controls(34.2,-97.1)and(34.2,-96.3)..(34.6,-96)
				..controls(35,-95.7)and(35.5,-95.9)..(36.2,-96.7)
				..controls(36.5,-97)and(37.2,-97.5)..(37.6,-97.7)..controls(38,-98)and(38.6,-98.4)..(39.1,-98.9)
				..controls(39.6,-99.4)and(40.3,-100)..(40.6,-100.3)
				..controls(41.3,-100.8)and(42.4,-101.8)..(43.7,-102.9)
				..controls(44,-103.2)and(44.6,-103.7)..(45,-103.9)
				..controls(45.4,-104.2)and(45.9,-104.6)..(46.1,-104.8)
				..controls(47.1,-105.9)and(48.8,-107.1)..(50.5,-107.9)
				..controls(51.8,-108.6)and(52.3,-109)..(53.4,-110)
				..controls(54.8,-111.2)and(55.6,-111.6)..(57,-112.1)
				..controls(57.5,-112.3)and(58.5,-112.8)..(59.4,-113.3)
				..controls(60.2,-113.7)and(61.3,-114.3)..(61.9,-114.5)
				..controls(62.4,-114.7)and(63.3,-115.1)..(63.8,-115.3)
				..controls(64.4,-115.6)and(65.4,-115.9)..(66.2,-116.1)
				..controls(67,-116.3)and(68.4,-116.6)..(69.2,-116.9)
				..controls(72.3,-117.9)and(72.8,-118)..(78,-118.1)
				..controls(84,-118.3)and(86.1,-118.1)..(88.1,-117.1)
				..controls(89.4,-116.5)and(93,-115)..(93.7,-114.9)
				..controls(94.8,-114.6)and(95.4,-114.2)..(95.9,-113.4)
				..controls(96.2,-113)and(96.7,-112.5)..(97.1,-112.2)
				..controls(98,-111.6)and(101.3,-108.4)..(102.1,-107.2)
				..controls(102.4,-106.8)and(102.9,-106.2)..(103.2,-105.8)
				..controls(103.5,-105.5)and(104,-104.8)..(104.2,-104.2)
				..controls(104.5,-103.6)and(105,-102.8)..(105.3,-102.4)
				..controls(105.6,-102)and(106.1,-101.2)..(106.5,-100.6)
				..controls(107.1,-99.5)and(108.6,-97.5)..(109.5,-96.5)
				..controls(109.9,-96.1)and(110.1,-96)..(110.4,-96)
				..controls(111.1,-96)and(111.6,-96.5)..(111.8,-97.6)
				..controls(112,-98.6)and(112.3,-99.3)..(112.8,-99.9)
				..controls(113.7,-100.9)and(115.2,-103.1)..(115.7,-104.2)
				..controls(116.1,-105)and(116.5,-105.5)..(116.9,-106)
				..controls(117.3,-106.4)and(117.8,-106.9)..(118,-107.3)
				..controls(118.2,-107.6)and(118.6,-108)..(119,-108.3)
				..controls(119.6,-108.8)and(121.1,-110.3)..(121.8,-111.2)
				..controls(122.1,-111.7)and(122.5,-112)..(123.4,-112.4)
				..controls(124.8,-113.2)and(125.8,-113.8)..(127,-114.6)
				..controls(127.5,-114.9)and(128.1,-115.3)..(128.2,-115.4)
				..controls(128.4,-115.5)and(129.3,-115.8)..(130.2,-116)
				..controls(131.1,-116.2)and(132.5,-116.6)..(133.3,-116.9)
				..controls(135.7,-117.9)and(136.1,-117.9)..(142.5,-117.8)
				..controls(145.8,-117.8)and(148.4,-117.7)..(148.8,-117.6)
				..controls(149.2,-117.6)and(150.2,-117.2)..(151,-116.8)
				..controls(152,-116.4)and(152.8,-116)..(153.6,-115.9)
				..controls(154.5,-115.7)and(155.2,-115.4)..(156.3,-114.9)
				..controls(157.1,-114.5)and(158.2,-114)..(158.7,-113.8)
				..controls(159.1,-113.7)and(159.8,-113.3)..(160.2,-113)
				..controls(160.6,-112.7)and(161.3,-112.4)..(161.7,-112.2)
				..controls(162.7,-111.9)and(163.2,-111.7)..(164.1,-111.1)
				..controls(164.5,-110.8)and(165.3,-110.4)..(165.9,-110)
				..controls(166.5,-109.7)and(167.3,-109.2)..(167.6,-108.9)
				..controls(167.9,-108.6)and(168.6,-108.2)..(169,-108)
				..controls(169.4,-107.8)and(170.2,-107.3)..(170.7,-106.9)
				..controls(171.2,-106.5)and(172.1,-105.8)..(172.7,-105.4)
				..controls(173.2,-105)and(174.1,-104.4)..(174.5,-104)
				..controls(174.9,-103.7)and(175.7,-103.1)..(176.2,-102.6)
				..controls(176.7,-102.2)and(177.4,-101.5)..(177.9,-101.2)
				..controls(178.3,-100.8)and(178.8,-100.4)..(179,-100.2)
				..controls(179.2,-100)and(179.6,-99.6)..(180,-99.3)
				..controls(180.4,-99.1)and(180.9,-98.6)..(181.2,-98.4)
				..controls(181.5,-98.1)and(182.1,-97.7)..(182.5,-97.5)
				..controls(183.2,-97.2)and(183.5,-96.9)..(184.3,-95.9)
				..controls(184.6,-95.6)and(184.6,-95.6)..(184.8,-95.7)
				..controls(185,-95.8)and(185.1,-96)..(185.1,-96.2)
				..controls(185.1,-96.5)and(185,-96.7)..(184.6,-97.2)
				..controls(184.1,-97.8)and(184,-97.9)..(183.9,-98.8)
				..controls(183.7,-100.3)and(183.3,-101.3)..(182.9,-101.8)
				..controls(182.6,-102.1)and(182.3,-102.6)..(182.2,-102.9)
				..controls(182,-103.3)and(181.6,-104)..(181.2,-104.6)
				..controls(180.9,-105.2)and(180.5,-106.1)..(180.4,-106.6)
				..controls(179.9,-108)and(179.7,-108.5)..(179,-109.4)
				..controls(178.6,-109.8)and(178.2,-110.5)..(177.9,-111)
				..controls(177.7,-111.5)and(177.3,-112.3)..(177,-112.7)
				..controls(176.6,-113.2)and(176.4,-113.7)..(176.3,-114.2)
				..controls(176.2,-114.7)and(175.9,-115.2)..(175.4,-116)
				..controls(175,-116.6)and(174.4,-117.7)..(174,-118.5)
				..controls(173.7,-119.2)and(173.3,-120.1)..(173,-120.5)
				..controls(172.8,-120.9)and(172.4,-121.8)..(172.1,-122.5)
				..controls(171.8,-123.2)and(171.3,-124.1)..(171,-124.6)
				..controls(170.8,-125.1)and(170.3,-126)..(169.9,-126.6)
				..controls(169.6,-127.3)and(169.2,-128.1)..(169,-128.3)
				..controls(168.8,-128.6)and(168.5,-129.3)..(168.3,-129.8)
				..controls(168.1,-130.3)and(167.8,-131.1)..(167.5,-131.6)
				..controls(166.9,-132.5)and(164.2,-137.9)..(164.2,-138.1)
				..controls(164.2,-138.1)and(163.9,-138.8)..(163.4,-139.7)
				..controls(162.9,-140.5)and(162.4,-141.5)..(162.3,-141.9)
				..controls(161.9,-143)and(161.6,-143.4)..(160.9,-144.1)
				..controls(160.4,-144.6)and(160.3,-144.8)..(160.1,-145.4)
				..controls(160,-145.8)and(159.6,-146.8)..(159.2,-147.8)
				..controls(158.7,-148.8)and(158.1,-150)..(157.9,-150.6)
				..controls(157.6,-151.1)and(157.2,-151.9)..(156.9,-152.4)
				..controls(156.5,-152.8)and(156.1,-153.7)..(155.8,-154.4)
				..controls(155.6,-155)and(155.2,-155.9)..(154.9,-156.3)
				..controls(154.7,-156.7)and(154.4,-157.4)..(154.2,-157.9)
				..controls(154.1,-158.4)and(153.8,-159)..(153.5,-159.3)
				..controls(152.8,-160.1)and(153,-160.1)..(142.4,-160.1)
				..controls(137.2,-160.1)and(131.8,-160.1)..(130.4,-160)--cycle
				(141.5,-103.7)..controls(140.8,-103)and(139.8,-102.1)..(139.2,-101.6)
				..controls(138.5,-101.2)and(137.9,-100.6)..(137.8,-100.4)
				..controls(137,-99.5)and(134,-96.5)..(133.2,-96)
				..controls(132.7,-95.7)and(131.8,-94.9)..(131.1,-94.3)
				..controls(129.4,-92.7)and(127.9,-91.6)..(126.6,-91.1)
				..controls(125.2,-90.5)and(124.1,-90)..(122.9,-89.1)
				..controls(122.1,-88.6)and(121.3,-88.2)..(119.9,-87.7)
				..controls(118.8,-87.3)and(117.5,-86.8)..(116.9,-86.5)
				..controls(114.3,-85.5)and(113.7,-85.4)..(109.7,-85.5)
				..controls(105.3,-85.7)and(103.9,-85.9)..(101.6,-87)
				..controls(100.7,-87.4)and(99.4,-87.8)..(97.9,-88)
				..controls(97,-88.2)and(96.6,-88.3)..(95.8,-88.9)
				..controls(95.4,-89.2)and(94.7,-89.6)..(94.2,-89.8)
				..controls(92.9,-90.3)and(92.3,-90.7)..(90.9,-92.1)
				..controls(90.2,-92.7)and(89.4,-93.3)..(89,-93.5)..controls(87,-94.5)and(85.2,-96)..(83.6,-97.9)
				..controls(83.3,-98.2)and(82.5,-98.9)..(81.8,-99.4)
				..controls(80.7,-100.1)and(80.5,-100.3)..(80,-101.1)
				..controls(78.7,-103.1)and(77.2,-103.9)..(76.4,-103)
				..controls(75.7,-102.2)and(75.7,-103.2)..(75.7,-83.6)
				..controls(75.6,-73.6)and(75.6,-64.8)..(75.5,-64)
				..controls(75.5,-63.2)and(75.5,-61.4)..(75.6,-59.8)--(75.7,-57.1)--(76,-56.8)
				..controls(76.8,-56.2)and(78,-56.3)..(79.8,-57.2)
				..controls(80.4,-57.5)and(81.3,-57.8)..(81.8,-57.9)
				..controls(82.3,-58.1)and(83,-58.4)..(83.6,-58.8)..controls(84.3,-59.2)and(85,-59.5)..(86,-59.8)
				..controls(86.8,-60.1)and(88,-60.6)..(88.7,-60.9)
				..controls(89.3,-61.2)and(90.6,-61.7)..(91.4,-62)
				..controls(92.2,-62.3)and(93.7,-62.9)..(94.6,-63.3)
				..controls(95.6,-63.7)and(96.8,-64.1)..(97.4,-64.3)
				..controls(98,-64.5)and(99,-64.9)..(99.7,-65.2)
				..controls(100.4,-65.6)and(101.5,-66)..(102.3,-66.2)
				..controls(104.2,-66.8)and(105,-67.1)..(106.2,-67.8)
				..controls(107.9,-68.7)and(110.2,-69.4)..(111.3,-69.2)
				..controls(112.5,-69)and(115,-68.2)..(115.6,-67.8)
				..controls(117.3,-66.7)and(117.6,-66.6)..(119.5,-66.2)
				..controls(121.9,-65.8)and(122.1,-65.7)..(123,-65.1)
				..controls(124.1,-64.4)and(124.4,-64.3)..(126,-63.9)
				..controls(127,-63.7)and(127.7,-63.4)..(128.4,-63.1)
				..controls(129,-62.8)and(129.7,-62.5)..(130.6,-62.4)
				..controls(132.1,-62.1)and(134,-61.5)..(135.2,-61)
				..controls(135.7,-60.8)and(136.7,-60.5)..(137.3,-60.3)
				..controls(138.4,-60)and(139.3,-59.7)..(141.1,-59)
				..controls(141.6,-58.7)and(142,-58.6)..(142.5,-58.6)
				..controls(143.2,-58.6)and(143.3,-58.7)..(143.6,-59)
				..controls(144,-59.5)and(144.1,-59.9)..(144,-60.8)
				..controls(144,-61.2)and(144,-71.1)..(144,-82.8)--(143.9,-104.1)--(143.7,-104.5)
				..controls(143.5,-104.7)and(143.3,-104.8)..(143,-104.9)
				..controls(142.6,-104.9)and(142.5,-104.8)..(141.5,-103.7)--cycle
				(27.8,-91.4)..controls(27.1,-90.9)and(26.3,-90)..(26.1,-89.5)
				..controls(26,-89.2)and(25.8,-88.6)..(25.7,-87.9)
				..controls(25.4,-86.2)and(25.6,-84.8)..(26.5,-83.5)
				..controls(27.4,-82.3)and(28.3,-82)..(29.4,-82.3)
				..controls(31.1,-82.8)and(32.1,-84.9)..(32,-87.4)
				..controls(31.9,-89.3)and(31.2,-90.6)..(29.9,-91.4)
				..controls(29.1,-91.9)and(28.6,-91.9)..(27.8,-91.4)--cycle
				(190.1,-91.3)..controls(189.7,-91.2)and(189.3,-90.9)..(189.1,-90.7)
				..controls(188.3,-89.8)and(187.6,-87.7)..(187.6,-86.2)
				..controls(187.6,-85.5)and(187.7,-85.2)..(188.2,-84.4)
				..controls(189.1,-82.6)and(190.1,-81.8)..(191.3,-82.1)
				..controls(192,-82.3)and(193.5,-83.7)..(193.8,-84.4)
				..controls(194.1,-85.1)and(194.2,-87.1)..(194,-88.3)
				..controls(193.8,-89.8)and(193.3,-90.5)..(192.3,-91.3)
				..controls(191.7,-91.7)and(190.9,-91.7)..(190.1,-91.3)--cycle
				(173.7,-83.4)..controls(172.2,-82.5)and(172,-82.5)..(170.6,-82.4)
				..controls(169.7,-82.4)and(168.8,-82.4)..(168.3,-82.5)
				..controls(167.2,-82.7)and(166.9,-82.7)..(166.6,-82.2)
				..controls(166.3,-81.8)and(166.3,-81)..(166.8,-80)
				..controls(167.2,-78.9)and(167.4,-77.4)..(167.7,-74.1)
				..controls(167.8,-72.8)and(167.9,-71.8)..(167.9,-71.7)
				..controls(168,-71.7)and(168,-70.2)..(168,-68.4)..controls(168,-65.4)and(168,-64.9)..(167.7,-63)
				..controls(167.3,-60.5)and(167.1,-59.9)..(166.4,-58.4)
				..controls(166.1,-57.7)and(165.8,-56.9)..(165.7,-56.6)
				..controls(165.7,-56.2)and(165.4,-55.6)..(165.2,-55.2)
				..controls(164.6,-53.8)and(161.9,-50.7)..(161,-50.3)
				..controls(159.8,-49.7)and(157.9,-50)..(155.2,-51.2)
				..controls(154.7,-51.4)and(153.4,-51.8)..(152.2,-52.1)
				..controls(150.8,-52.4)and(149.4,-52.9)..(148.6,-53.2)
				..controls(147.8,-53.6)and(146.7,-54)..(146,-54.2)
				..controls(145.4,-54.4)and(144.5,-54.7)..(144,-54.9)
				..controls(143.6,-55.2)and(142.7,-55.5)..(141.8,-55.7)
				..controls(141,-55.9)and(139.6,-56.3)..(138.7,-56.6)
				..controls(137.8,-57)and(136.3,-57.4)..(135.3,-57.6)
				..controls(133.8,-57.9)and(133.2,-58.1)..(132.1,-58.7)
				..controls(131.1,-59.2)and(130.3,-59.4)..(129.4,-59.6)
				..controls(128.7,-59.8)and(127.8,-60)..(127.4,-60.1)
				..controls(127,-60.3)and(126,-60.5)..(125.1,-60.7)
				..controls(123.2,-61.1)and(121.7,-61.5)..(120.2,-62.1)
				..controls(119.7,-62.3)and(118.1,-62.8)..(116.7,-63.2)
				..controls(114.6,-63.9)and(114,-64.1)..(113,-64.7)
				..controls(111.9,-65.4)and(111.8,-65.4)..(111,-65.5)
				..controls(109.9,-65.5)and(109.3,-65.4)..(107.7,-64.7)
				..controls(107,-64.4)and(106,-64)..(105.3,-63.9)
				..controls(103.5,-63.5)and(102.2,-63)..(101.2,-62.5)
				..controls(100.1,-61.8)and(99.3,-61.5)..(97.8,-61.1)
				..controls(97.2,-61)and(96.2,-60.7)..(95.7,-60.4)
				..controls(95.1,-60.1)and(94.2,-59.8)..(93.5,-59.6)
				..controls(92.9,-59.4)and(91.7,-59)..(91,-58.7)..controls(90.2,-58.4)and(88.9,-58)..(88,-57.7)
				..controls(87.2,-57.5)and(86.1,-57.1)..(85.5,-56.8)
				..controls(85,-56.6)and(84.1,-56.2)..(83.5,-56)..controls(82.9,-55.7)and(81.8,-55.2)..(81,-54.9)
				..controls(80.1,-54.5)and(79.1,-54.2)..(78.7,-54.1)
				..controls(77.9,-53.8)and(77.3,-53.6)..(76.2,-52.9)
				..controls(75.7,-52.6)and(75.2,-52.4)..(74.1,-52.2)
				..controls(72.5,-51.9)and(72.4,-51.8)..(71.3,-51.1)
				..controls(69.9,-50.1)and(69.4,-49.9)..(67.9,-49.8)
				..controls(67,-49.7)and(66.6,-49.6)..(66.3,-49.4)
				..controls(66.1,-49.3)and(65.7,-49)..(65.4,-48.8)
				..controls(65.1,-48.7)and(64.2,-48.4)..(63.6,-48.2)
				..controls(62.9,-48)and(61.7,-47.6)..(60.9,-47.2)
				..controls(60.1,-46.8)and(58.9,-46.4)..(58.2,-46.2)
				..controls(57.6,-46)and(56.8,-45.7)..(56.5,-45.5)
				..controls(54.9,-44.6)and(54.6,-44.5)..(53.2,-44.2)
				..controls(51.7,-43.8)and(51,-43.6)..(49.9,-42.9)
				..controls(49.5,-42.6)and(48.9,-42.4)..(47.5,-42)
				..controls(45.7,-41.6)and(45.6,-41.5)..(44.3,-40.6)
				..controls(43.5,-40.1)and(42.8,-39.6)..(42.7,-39.4)
				..controls(42.4,-38.7)and(43,-38.1)..(44.5,-37.6)..controls(44.9,-37.5)and(46,-37)..(46.9,-36.6)
				..controls(47.9,-36.1)and(48.8,-35.8)..(49.6,-35.6)
				..controls(50.5,-35.4)and(51.3,-35.2)..(52.4,-34.6)
				..controls(53.3,-34.1)and(54.5,-33.7)..(54.9,-33.5)
				..controls(55.4,-33.4)and(56.4,-33)..(57.2,-32.6)
				..controls(58.6,-32)and(60.1,-31.5)..(61.6,-31.2)
				..controls(62.3,-31.1)and(62.7,-31)..(63.2,-30.7)
				..controls(64.1,-30)and(64.3,-29.9)..(65.3,-29.7)
				..controls(65.7,-29.6)and(66.5,-29.2)..(67.1,-28.9)
				..controls(67.6,-28.6)and(68.8,-28.2)..(69.6,-27.9)
				..controls(70.5,-27.7)and(71.6,-27.3)..(72.1,-27)
				..controls(72.6,-26.7)and(73.8,-26.2)..(74.8,-25.9)
				..controls(75.7,-25.5)and(76.7,-25.1)..(77.1,-25)
				..controls(77.4,-24.8)and(78.5,-24.4)..(79.6,-24.1)
				..controls(80.6,-23.8)and(82,-23.2)..(82.8,-22.9)
				..controls(83.5,-22.5)and(84.8,-22)..(85.8,-21.8)
				..controls(86.7,-21.5)and(87.8,-21.1)..(88.1,-20.9)
				..controls(88.4,-20.7)and(89.3,-20.3)..(90.1,-20.1)
				..controls(90.8,-19.8)and(91.9,-19.4)..(92.4,-19.1)
				..controls(93.5,-18.5)and(94.4,-18.2)..(95.4,-18.1)
				..controls(95.7,-18)and(96.4,-17.8)..(96.9,-17.5)
				..controls(97.5,-17.3)and(98.4,-16.9)..(99,-16.6)
				..controls(99.7,-16.3)and(100.9,-15.9)..(101.9,-15.6)
				..controls(102.9,-15.3)and(104.4,-14.7)..(105.3,-14.3)
				..controls(108.1,-13.2)and(108.5,-13.1)..(110.1,-13.1)
				..controls(111.4,-13.1)and(111.5,-13.1)..(112.4,-13.5)
				..controls(113.2,-14)and(115,-14.8)..(116.6,-15.4)
				..controls(116.9,-15.5)and(117.9,-15.8)..(118.8,-16.1)
				..controls(120.9,-16.6)and(121.5,-16.9)..(122.6,-17.9)
				..controls(123,-18.3)and(123.7,-18.8)..(123.9,-18.9)
				..controls(124.6,-19.4)and(126,-19.8)..(127.9,-20.2)
				..controls(129.6,-20.5)and(129.8,-20.6)..(131.1,-21.4)
				..controls(131.5,-21.7)and(132.4,-22.1)..(133.1,-22.3)
				..controls(136.1,-23.2)and(138.3,-24.1)..(140.2,-25)
				..controls(141,-25.5)and(141.5,-25.6)..(142.5,-25.8)
				..controls(145.1,-26.3)and(146,-26.6)..(149.4,-28.1)
				..controls(151.3,-28.9)and(153.5,-29.8)..(154.3,-30.1)
				..controls(155.1,-30.4)and(156.2,-30.8)..(156.6,-31)
				..controls(157.1,-31.2)and(158.2,-31.6)..(159.1,-31.8)
				..controls(160.6,-32.2)and(160.9,-32.3)..(161.7,-32.9)
				..controls(162.4,-33.3)and(162.9,-33.6)..(163.7,-33.8)
				..controls(164.3,-34)and(164.8,-34.1)..(164.9,-34.1)
				..controls(165,-34.1)and(165.6,-34.4)..(166.4,-34.8)
				..controls(167.3,-35.2)and(168.1,-35.4)..(169,-35.6)
				..controls(170,-35.8)and(170.5,-36)..(171.7,-36.6)
				..controls(172.4,-37)and(173.3,-37.4)..(173.6,-37.4)
				..controls(175.6,-37.7)and(176.2,-38)..(176.7,-39.1)
				..controls(176.9,-39.6)and(176.9,-39.7)..(176.7,-40.4)
				..controls(176.4,-41.6)and(175.8,-42.7)..(175.2,-43.4)
				..controls(174.5,-44.1)and(172.6,-45.5)..(171.7,-45.8)
				..controls(170.6,-46.2)and(169.1,-46)..(168.2,-45.4)
				..controls(167.9,-45.2)and(167.3,-44.8)..(166.9,-44.4)
				..controls(166.4,-44.1)and(165.8,-43.6)..(165.4,-43.2)
				..controls(165.1,-42.9)and(164.4,-42.3)..(163.9,-41.9)
				..controls(163.3,-41.5)and(162.7,-40.9)..(162.4,-40.5)
				..controls(161.4,-39.2)and(159.9,-38)..(158.8,-37.8)
				..controls(158.4,-37.8)and(158.3,-37.7)..(158.2,-37.2)
				..controls(158,-36.8)and(157.4,-36.1)..(157.1,-36.1)
				..controls(157,-36.1)and(156.7,-36.1)..(156.5,-36)
				..controls(156.1,-35.9)and(155.9,-35.7)..(155.7,-35.2)
				..controls(155.3,-34.6)and(155.1,-34.4)..(154,-34.1)
				..controls(153.6,-34)and(153.2,-33.7)..(152.9,-33.5)
				..controls(152.4,-32.9)and(151.6,-32.4)..(151.3,-32.4)
				..controls(150.9,-32.4)and(150.4,-32.6)..(150.4,-32.9)
				..controls(150.4,-33.2)and(150.8,-33.4)..(151.4,-33.6)
				..controls(152.1,-33.7)and(152.8,-34.1)..(153.5,-34.9)
				..controls(154.1,-35.5)and(154.3,-35.6)..(155,-35.8)
				..controls(155.7,-36)and(155.7,-36)..(156,-36.6)
				..controls(156.3,-37.3)and(156.7,-37.7)..(157.5,-37.8)
				..controls(158.1,-37.9)and(158.2,-38)..(158.4,-38.8)
				..controls(158.5,-39.4)and(158.6,-39.6)..(159.6,-40.4)
				..controls(160.2,-41)and(160.9,-41.6)..(161.3,-41.8)
				..controls(161.7,-42)and(161.9,-42.3)..(162.2,-42.8)
				..controls(162.4,-43.2)and(162.8,-43.8)..(163.1,-44.1)
				..controls(163.9,-44.8)and(165.1,-46.7)..(165.7,-48)
				..controls(166.2,-49.3)and(166.4,-49.7)..(167.2,-50.5)
				..controls(167.5,-50.8)and(167.9,-51.4)..(168.1,-51.8)
				..controls(168.3,-52.2)and(168.8,-53)..(169.1,-53.6)
				..controls(169.6,-54.3)and(169.8,-54.9)..(169.9,-55.6)
				..controls(170.4,-57.8)and(170.5,-58.1)..(171.2,-59.5)
				..controls(171.6,-60.3)and(172,-61.5)..(172.2,-62.1)
				..controls(172.3,-62.8)and(172.7,-63.7)..(172.9,-64.1)
				..controls(173.2,-64.6)and(173.4,-65.3)..(173.4,-65.6)
				..controls(173.8,-68.8)and(174.2,-70.6)..(174.9,-72.6)
				..controls(175.5,-74.5)and(175.9,-77.1)..(176.1,-80.2)
				..controls(176.3,-83.4)and(175.5,-84.4)..(173.7,-83.4)--cycle;
				\definecolor{ezs}{RGB}{247,179,26}\fill[ezs](-5.7,-156.6)--(-3.4,-157.5)--(61.9,12.7)--(59.7,13.6)--cycle;
				\definecolor{ezs}{RGB}{247,179,22}\fill[ezs](-3.8,-157.3)--(-1.5,-158.2)--(63.9,12)--(61.6,12.9)--cycle;
				\definecolor{ezs}{RGB}{248,179,19}\fill[ezs](-1.8,-158.1)--(0.4,-158.9)--(65.8,11.2)--(63.5,12.1)--cycle;
				\definecolor{ezs}{RGB}{248,178,16}\fill[ezs](0.1,-158.8)--(2.4,-159.7)--(67.7,10.5)--(65.4,11.4)--cycle;
				\definecolor{ezs}{RGB}{249,178,15}\fill[ezs](2,-159.5)--(4.3,-160.4)--(69.6,9.8)--(67.3,10.7)--cycle;
				\definecolor{ezs}{RGB}{251,177,17}\fill[ezs](3.9,-160.3)--(6.2,-161.1)--(71.5,9)--(69.2,9.9)--cycle;
				\definecolor{ezs}{RGB}{252,179,18}\fill[ezs](5.8,-161)--(8.1,-161.9)--(73.4,8.3)--(71.1,9.2)--cycle;
				\definecolor{ezs}{RGB}{252,182,17}\fill[ezs](7.7,-161.7)--(10,-162.6)--(75.3,7.6)--(73,8.5)--cycle;
				\definecolor{ezs}{RGB}{252,185,17}\fill[ezs](9.6,-162.5)--(11.9,-163.4)--(77.2,6.8)--(75,7.7)--cycle;
				\definecolor{ezs}{RGB}{251,189,18}\fill[ezs](11.5,-163.2)--(13.8,-164.1)--(79.2,6.1)--(76.9,7)--cycle;
				\definecolor{ezs}{RGB}{251,191,22}\fill[ezs](13.4,-163.9)--(15.7,-164.8)--(81.1,5.4)--(78.8,6.2)--cycle;
				\definecolor{ezs}{RGB}{252,193,25}\fill[ezs](15.4,-164.7)--(17.6,-165.6)--(83,4.6)--(80.7,5.5)--cycle;
				\definecolor{ezs}{RGB}{253,195,27}\fill[ezs](17.3,-165.4)--(19.6,-166.3)--(84.9,3.9)--(82.6,4.8)--cycle;
				\definecolor{ezs}{RGB}{253,196,31}\fill[ezs](19.2,-166.1)--(21.5,-167)--(86.8,3.2)--(84.5,4)--cycle;
				\definecolor{ezs}{RGB}{254,199,35}\fill[ezs](21.1,-166.9)--(23.4,-167.8)--(88.7,2.4)--(86.4,3.3)--cycle;
				\definecolor{ezs}{RGB}{253,201,38}\fill[ezs](23,-167.6)--(25.3,-168.5)--(90.6,1.7)--(88.3,2.6)--cycle;
				\definecolor{ezs}{RGB}{253,203,40}\fill[ezs](24.9,-168.3)--(27.2,-169.2)--(92.5,1)--(90.3,1.8)--cycle;
				\definecolor{ezs}{RGB}{253,205,42}\fill[ezs](26.8,-169.1)--(29.1,-170)--(94.4,0.2)--(92.2,1.1)--cycle;
				\definecolor{ezs}{RGB}{253,207,45}\fill[ezs](28.7,-169.8)--(31,-170.7)--(96.4,-0.5)--(94.1,0.4)--cycle;
				\definecolor{ezs}{RGB}{253,210,49}\fill[ezs](30.7,-170.5)--(32.9,-171.4)--(98.3,-1.2)--(96,-0.4)--cycle;
				\definecolor{ezs}{RGB}{253,212,54}\fill[ezs](32.6,-171.3)--(34.9,-172.2)--(100.2,-2)--(97.9,-1.1)--cycle;
				\definecolor{ezs}{RGB}{253,214,59}\fill[ezs](34.5,-172)--(36.8,-172.9)--(102.1,-2.7)--(99.8,-1.8)--cycle;
				\definecolor{ezs}{RGB}{253,216,63}\fill[ezs](36.4,-172.7)--(38.7,-173.6)--(104,-3.4)--(101.7,-2.6)--cycle;
				\definecolor{ezs}{RGB}{253,218,65}\fill[ezs](38.3,-173.5)--(40.6,-174.4)--(105.9,-4.2)--(103.6,-3.3)--cycle;
				\definecolor{ezs}{RGB}{254,220,66}\fill[ezs](40.2,-174.2)--(42.5,-175.1)--(107.8,-4.9)--(105.5,-4)--cycle;
				\definecolor{ezs}{RGB}{254,222,68}\fill[ezs](42.1,-175)--(44.4,-175.8)--(109.7,-5.6)--(107.5,-4.8)--cycle;
				\definecolor{ezs}{RGB}{254,224,71}\fill[ezs](44,-175.7)--(46.3,-176.6)--(111.7,-6.4)--(109.4,-5.5)--cycle;
				\definecolor{ezs}{RGB}{254,226,75}\fill[ezs](46,-176.4)--(48.2,-177.3)--(113.6,-7.1)--(111.3,-6.2)--cycle;
				\definecolor{ezs}{RGB}{254,228,79}\fill[ezs](47.9,-177.2)--(50.1,-178)--(115.5,-7.8)--(113.2,-7)--cycle;
				\definecolor{ezs}{RGB}{254,230,82}\fill[ezs](49.8,-177.9)--(52.1,-178.8)--(117.4,-8.6)--(115.1,-7.7)--cycle;
				\definecolor{ezs}{RGB}{254,232,85}\fill[ezs](51.7,-178.6)--(54,-179.5)--(119.3,-9.3)--(117,-8.4)--cycle;
				\definecolor{ezs}{RGB}{254,234,88}\fill[ezs](53.6,-179.4)--(55.9,-180.2)--(121.2,-10)--(118.9,-9.2)--cycle;
				\definecolor{ezs}{RGB}{254,237,91}\fill[ezs](55.5,-180.1)--(57.8,-181)--(123.1,-10.8)--(120.8,-9.9)--cycle;
				\definecolor{ezs}{RGB}{254,239,94}\fill[ezs](57.4,-180.8)--(59.7,-181.7)--(125,-11.5)--(122.8,-10.6)--cycle;
				\definecolor{ezs}{RGB}{254,241,97}\fill[ezs](59.3,-181.6)--(61.6,-182.4)--(127,-12.2)--(124.7,-11.4)--cycle;
				\definecolor{ezs}{RGB}{254,243,101}\fill[ezs](61.2,-182.3)--(63.5,-183.2)--(128.9,-13)--(126.6,-12.1)--cycle;
				\definecolor{ezs}{RGB}{254,245,103}\fill[ezs](63.2,-183)--(65.4,-183.9)--(130.8,-13.7)--(128.5,-12.8)--cycle;
				\definecolor{ezs}{RGB}{254,245,105}\fill[ezs](65.1,-183.8)--(67.4,-184.6)--(132.7,-14.4)--(130.4,-13.6)--cycle;
				\definecolor{ezs}{RGB}{254,245,105}\fill[ezs](67,-184.5)--(69.3,-185.4)--(134.6,-15.2)--(132.3,-14.3)--cycle;
				\definecolor{ezs}{RGB}{254,244,103}\fill[ezs](68.9,-185.2)--(71.2,-186.1)--(136.5,-15.9)--(134.2,-15)--cycle;
				\definecolor{ezs}{RGB}{253,243,101}\fill[ezs](70.8,-186)--(73.1,-186.8)--(138.4,-16.6)--(136.1,-15.8)--cycle;
				\definecolor{ezs}{RGB}{253,241,98}\fill[ezs](72.7,-186.7)--(75,-187.6)--(140.3,-17.4)--(138,-16.5)--cycle;
				\definecolor{ezs}{RGB}{253,239,95}\fill[ezs](74.6,-187.4)--(76.9,-188.3)--(142.2,-18.1)--(140,-17.2)--cycle;
				\definecolor{ezs}{RGB}{253,236,92}\fill[ezs](76.5,-188.2)--(78.8,-189)--(144.2,-18.8)--(141.9,-18)--cycle;
				\definecolor{ezs}{RGB}{253,234,88}\fill[ezs](78.5,-188.9)--(80.7,-189.8)--(146.1,-19.6)--(143.8,-18.7)--cycle;
				\definecolor{ezs}{RGB}{253,232,85}\fill[ezs](80.4,-189.6)--(82.7,-190.5)--(148,-20.3)--(145.7,-19.4)--cycle;
				\definecolor{ezs}{RGB}{253,230,82}\fill[ezs](82.3,-190.4)--(84.6,-191.2)--(149.9,-21)--(147.6,-20.2)--cycle;
				\definecolor{ezs}{RGB}{253,228,78}\fill[ezs](84.2,-191.1)--(86.5,-192)--(151.8,-21.8)--(149.5,-20.9)--cycle;
				\definecolor{ezs}{RGB}{253,226,74}\fill[ezs](86.1,-191.8)--(88.4,-192.7)--(153.7,-22.5)--(151.4,-21.6)--cycle;
				\definecolor{ezs}{RGB}{253,224,70}\fill[ezs](88,-192.6)--(90.3,-193.4)--(155.6,-23.3)--(153.3,-22.4)--cycle;
				\definecolor{ezs}{RGB}{253,222,67}\fill[ezs](89.9,-193.3)--(92.2,-194.2)--(157.5,-24)--(155.3,-23.1)--cycle;
				\definecolor{ezs}{RGB}{253,220,65}\fill[ezs](91.8,-194)--(94.1,-194.9)--(159.5,-24.7)--(157.2,-23.8)--cycle;
				\definecolor{ezs}{RGB}{253,217,63}\fill[ezs](93.7,-194.8)--(96,-195.6)--(161.4,-25.5)--(159.1,-24.6)--cycle;
				\definecolor{ezs}{RGB}{253,215,60}\fill[ezs](95.7,-195.5)--(97.9,-196.4)--(163.3,-26.2)--(161,-25.3)--cycle;
				\definecolor{ezs}{RGB}{253,213,56}\fill[ezs](97.6,-196.2)--(99.9,-197.1)--(165.2,-26.9)--(162.9,-26)--cycle;
				\definecolor{ezs}{RGB}{252,211,53}\fill[ezs](99.5,-197)--(101.8,-197.8)--(167.1,-27.7)--(164.8,-26.8)--cycle;
				\definecolor{ezs}{RGB}{252,209,49}\fill[ezs](101.4,-197.7)--(103.7,-198.6)--(169,-28.4)--(166.7,-27.5)--cycle;
				\definecolor{ezs}{RGB}{252,207,45}\fill[ezs](103.3,-198.4)--(105.6,-199.3)--(170.9,-29.1)--(168.6,-28.2)--cycle;
				\definecolor{ezs}{RGB}{253,205,42}\fill[ezs](105.2,-199.2)--(107.5,-200)--(172.8,-29.9)--(170.6,-29)--cycle;
				\definecolor{ezs}{RGB}{253,202,40}\fill[ezs](107.1,-199.9)--(109.4,-200.8)--(174.7,-30.6)--(172.5,-29.7)--cycle;
				\definecolor{ezs}{RGB}{253,200,38}\fill[ezs](109,-200.6)--(111.3,-201.5)--(176.7,-31.3)--(174.4,-30.4)--cycle;
				\definecolor{ezs}{RGB}{253,198,35}\fill[ezs](111,-201.4)--(113.2,-202.2)--(178.6,-32.1)--(176.3,-31.2)--cycle;
				\definecolor{ezs}{RGB}{253,196,31}\fill[ezs](112.9,-202.1)--(115.2,-203)--(180.5,-32.8)--(178.2,-31.9)--cycle;
				\definecolor{ezs}{RGB}{253,194,28}\fill[ezs](114.8,-202.8)--(117.1,-203.7)--(182.4,-33.5)--(180.1,-32.6)--cycle;
				\definecolor{ezs}{RGB}{253,192,26}\fill[ezs](116.7,-203.6)--(119,-204.5)--(184.3,-34.3)--(182,-33.4)--cycle;
				\definecolor{ezs}{RGB}{253,189,24}\fill[ezs](118.6,-204.3)--(120.9,-205.2)--(186.2,-35)--(183.9,-34.1)--cycle;
				\definecolor{ezs}{RGB}{253,186,20}\fill[ezs](120.5,-205)--(122.8,-205.9)--(188.1,-35.7)--(185.8,-34.8)--cycle;
				\definecolor{ezs}{RGB}{253,184,17}\fill[ezs](122.4,-205.8)--(124.7,-206.7)--(190,-36.5)--(187.8,-35.6)--cycle;
				\definecolor{ezs}{RGB}{252,183,15}\fill[ezs](124.3,-206.5)--(126.6,-207.4)--(192,-37.2)--(189.7,-36.3)--cycle;
				\definecolor{ezs}{RGB}{252,181,14}\fill[ezs](126.3,-207.2)--(128.5,-208.1)--(193.9,-37.9)--(191.6,-37.1)--cycle;
				\definecolor{ezs}{RGB}{252,179,13}\fill[ezs](128.2,-208)--(130.4,-208.9)--(195.8,-38.7)--(193.5,-37.8)--cycle;
				\definecolor{ezs}{RGB}{252,178,12}\fill[ezs](130.1,-208.7)--(132.4,-209.6)--(197.7,-39.4)--(195.4,-38.5)--cycle;
				\definecolor{ezs}{RGB}{251,178,10}\fill[ezs](132,-209.4)--(134.3,-210.3)--(199.6,-40.1)--(197.3,-39.3)--cycle;
				\definecolor{ezs}{RGB}{251,178,10}\fill[ezs](133.9,-210.2)--(136.2,-211.1)--(201.5,-40.9)--(199.2,-40)--cycle;
				\definecolor{ezs}{RGB}{250,177,11}\fill[ezs](135.8,-210.9)--(138.1,-211.8)--(203.4,-41.6)--(201.1,-40.7)--cycle;
				\definecolor{ezs}{RGB}{250,178,11}\fill[ezs](137.7,-211.6)--(140,-212.5)--(205.3,-42.3)--(203.1,-41.5)--cycle;
				\definecolor{ezs}{RGB}{248,180,10}\fill[ezs](139.6,-212.4)--(141.9,-213.3)--(207.3,-43.1)--(205,-42.2)--cycle;
				\definecolor{ezs}{RGB}{247,181,9}\fill[ezs](141.5,-213.1)--(143.8,-214)--(209.2,-43.8)--(206.9,-42.9)--cycle;
				\definecolor{ezs}{RGB}{246,182,9}\fill[ezs](143.5,-213.8)--(145.7,-214.7)--(211.1,-44.5)--(208.8,-43.7)--cycle;
				\definecolor{ezs}{RGB}{246,182,10}\fill[ezs](145.4,-214.6)--(147.7,-215.5)--(213,-45.3)--(210.7,-44.4)--cycle;
			\end{scope}
			% chữ EDUZY (Times đậm, làm dày nét bằng 5 lớp lệch nhẹ; căn đúng khung chữ logo gốc)
			\foreach \dx/\dy in {-0.7/0,0.7/0,0/-0.3,0/0.3,0/0}
			\node[anchor=center,inner sep=0pt,outer sep=0pt,text=ezText,align=center] at (109.6+\dx,-193.6+\dy)
			{\resizebox{1.01cm}{0.22cm}{\fontencoding{T1}\fontfamily{ptm}\fontseries{b}\selectfont EDUZY}};
		\end{tikzpicture}%
	}%
}
%====================== THIẾT LẬP THÂN ĐỀ ======================
%--- Môi trường tự luận {bt}: ex_test 3.3.4 chưa khai báo sẵn nên tạo ở đây
%    (ex_test đã tự móc \fixshowans và \listenumerate cho bt)
\makeatletter
\@ifundefined{bt}{\newtheorem{bt}{\color{gold1}\namebt}}{}
\makeatother
\renewcommand{\sectionmark}[1]{\markright{\thesection\ #1}}
%--- Khung tiêu đề (logo + tối đa 4 dòng; dòng 3, 4 để trống {} thì bỏ qua)
%  Đề thi:    \khungtieude{THI THỬ THPT QUỐC GIA NĂM 2026}{Môn Toán -- Đề số 21}
%                         {(Đề thi gồm: ...)}{Thời gian làm bài: 90 phút (không kể thời gian phát đề).}
%  Chuyên đề: \khungtieude{CHUYÊN ĐỀ 3 -- DÃY SỐ VÀ CÔNG THỨC TRUY HỒI}
%                         {Bồi dưỡng học sinh giỏi Toán THPT}{...}{}
\newcommand{\khungtieude}[4]{%
	\begin{tcolorbox}[colframe=gold1,colback=gold5!45,boxrule=0.5mm]
		\begin{minipage}[c]{3cm}
			\centerline{\eduzylogo[2.4cm]}
		\end{minipage}\hspace{0.5cm}
		\begin{minipage}[c]{12.5cm}
			\vskip 0.15 cm
			\centerline{\ind{#1}}
			\vskip 0.15 cm
			\centerline{\textbf{#2}}
			\if\relax\detokenize{#3}\relax\else
				\vskip 0.15 cm
				\centerline{\textit{#3}}
			\fi
			\if\relax\detokenize{#4}\relax\else
				\vskip 0.15 cm
				\centerline{\textit{#4}}
			\fi
			\vskip 0.15 cm
		\end{minipage}
	\end{tcolorbox}
}
%====================== BỔ SUNG CHO BÀI HỌC, CHUYÊN ĐỀ ======================
\usepackage[shortlabels]{enumitem}
\usepackage{tabularx}
%--- Danh sách gọn: \begin{enumerate}[gon] ... \end{enumerate}
\SetEnumitemKey{gon}{leftmargin=*,itemsep=1pt,parsep=0pt,topsep=3pt}
%--- Khung lý thuyết: \begin{khung}{} ... \end{khung}; tiêu đề chỉ dùng khi thật ngắn ("Mẹo", "Lỗi thường gặp")
\newtcolorbox{khung}[1]{breakable,enhanced,colframe=gold1,colback=gold5!45,
	boxrule=0.4mm,fonttitle=\bfseries\sffamily,coltitle=white,colbacktitle=gold1,
	attach boxed title to top left={xshift=6mm,yshift=-3mm},
	boxed title style={arc=1mm,boxrule=0.3mm},title={#1}}
%--- Tiêu đề mục nhỏ trong một dạng: \muc{Dấu hiệu}, \muc{Phương pháp}, \muc{Bài tự luyện}
\makeatletter
\newcommand{\muc}[1]{\par\medskip\noindent\indam{#1}\par\nobreak\smallskip\@afterindentfalse\@afterheading}
\makeatother
%--- Ô công thức cần nhớ (căn giữa, viền vàng, vừa khít nội dung)
%     \begin{center}\congthuc{$Q(n)=(n-1)!$}\end{center}
\newtcbox{\congthuc}{on line,colframe=gold1,colback=gold5!45,boxrule=0.5pt,arc=2pt,
	left=8pt,right=8pt,top=4pt,bottom=4pt}
%--- Công tắc lời giải bài tự luyện (ví dụ mẫu luôn dùng \loigiai)
%     \loigiaiTL[{đáp số}]{lời giải đầy đủ}
%     \TLfull -> in lời giải đầy đủ;  \TLdapso -> chỉ in một dòng "Đáp số";  \TLan -> ẩn hoàn toàn
%     (\CheDo là dethi thì mọi lời giải và đáp số đều ẩn)
\makeatletter
\newif\ifTLhien   \TLhientrue
\newif\ifTLdapso  \TLdapsofalse
\newcommand{\TLfull}{\TLhientrue\TLdapsofalse}
\newcommand{\TLdapso}{\TLhientrue\TLdapsotrue}
\newcommand{\TLan}{\TLhienfalse}
\newcommand{\loigiaiTL}[2][]{%
	\ifTLhien
		\ifTLdapso
			\def\@tlarg{#1}%
			\ifx\@tlarg\@empty\else
				\ifdefstring{\kindOPTN}{dethi}{}{\par\noindent{\bfseries\color{gold1}Đáp số:} #1\par}%
			\fi
		\else
			\loigiai{#2}%
		\fi
	\fi}
\makeatother
\TLdapso% mặc định: chỉ in đáp số bài tự luyện
%====================== BẢNG ĐÚNG/SAI (\choiceTFt) ======================
%  Hàng tiêu đề nền nâu đồng chữ trắng (Phát biểu | Đ | S), các hàng tô xen kẽ kem / trắng,
%  đường kẻ vàng kim, nhãn a) b) c) d) màu nâu đồng. Bản có đáp án: cột được chọn có vòng tròn
%  nâu đồng với dấu tích trắng (câu đúng tích ở cột Đúng, câu sai tích ở cột Sai).
%  Tiêu đề ghi đủ "Đúng | Sai": \OPTN{viettat=0} sau \input (cột Đúng/Sai rộng ra, phát biểu hẹp lại).
%  Dùng bảng gốc của ex_test: \def\BangDSgoc{1} trước \input{preamble-eduzy-v2}.
\providecommand{\BangDSgoc}{0}
\ifnum\BangDSgoc=0
\usepackage{colortbl,pifont}
\colorlet{ezDSke}{gold2!60!gold1}% đường kẻ bảng (nét 0,4pt như ex_test để không đổi chiều cao bảng)
\colorlet{ezDSle}{gold5!60}% hàng a), c)
\colorlet{ezDSchan}{white}% hàng b), d)
%--- dấu chọn ở cột Đúng/Sai
%    (bề rộng 0, căn giữa ô: không làm đổi bề rộng cột của ex_test, không báo tràn lề)
\newcommand{\EZdauchon}{\makebox[0pt]{\tikz[baseline=-1ex]\node[circle,fill=gold1,inner sep=0pt,
	minimum size=1.2em,text=white,font=\small]{\ding{51}};}}
\def\tickT{\EZdauchon}\def\tickF{\EZdauchon}
\AtBeginDocument{\sbox{\ccTFTab}{\EZdauchon}}
%--- nhãn a) b) c) d) màu nâu đồng (ở bản lời giải, nhãn câu đúng vẫn đóng khung như ex_test)
\renewcommand{\FalseTF}{\stepcounter{dapan}\textbf{\color{gold1}\DapAnTF\sepTF}\ignorespaces}
%--- trước khi dựng bảng: kẻ màu vàng kim (bề rộng cột giữ nguyên như ex_test); sau bảng trả lại màu đen
\newcount\EZhangTF
\newcommand{\EZbatdauTF}{%
	\global\EZhangTF=0\relax
	\arrayrulecolor{ezDSke}}
\newcommand{\EZketthucTF}{\arrayrulecolor{black}}
%--- hàng tiêu đề nền nâu đồng, chữ trắng
\makeatletter
\newcommand{\EZo@TF}[1]{\centering\bfseries\color{white}#1}
\newcommand{\EZtieudeTF}{%
	\ifcase\cotDS
		\ifnum\cotSTT=1\relax
			\def\rowHeader{\hline\rowcolor{gold1}\EZo@TF{TT}&\EZo@TF{\phatbieu}\tabularnewline}%
		\else
			\def\rowHeader{\hline\rowcolor{gold1}\EZo@TF{\phatbieu}\tabularnewline}%
		\fi
	\or
		\ifnum\cotSTT=1\relax
			\def\rowHeader{\hline\rowcolor{gold1}\EZo@TF{TT}&\EZo@TF{\phatbieu}&\EZo@TF{\Dung}\tabularnewline}%
		\else
			\def\rowHeader{\hline\rowcolor{gold1}\EZo@TF{\phatbieu}&\EZo@TF{\Dung}\tabularnewline}%
		\fi
	\or
		\ifnum\cotSTT=1\relax
			\def\rowHeader{\hline\rowcolor{gold1}\EZo@TF{TT}&\EZo@TF{\phatbieu}&\EZo@TF{\Dung}&\EZo@TF{\Sai}\tabularnewline}%
		\else
			\def\rowHeader{\hline\rowcolor{gold1}\EZo@TF{\phatbieu}&\EZo@TF{\Dung}&\EZo@TF{\Sai}\tabularnewline}%
		\fi
	\fi}
%--- tô xen kẽ từng hàng (gọi ngay sau \hline đầu mỗi hàng)
\newcommand{\EZmauhangTF}{%
	\noalign{\global\advance\EZhangTF\@ne}%
	\ifodd\EZhangTF\rowcolor{ezDSle}\else\rowcolor{ezDSchan}\fi}
%--- gắn vào \choiceTFt / \writeChoiceTFt của ex_test
\xpretocmd{\choiceTFt}{\EZbatdauTF}{}{\typeout{[preamble-eduzy-v2] Không gắn được bảng Đúng/Sai (đầu bảng)}}
\xpatchcmd{\choiceTFt}{\ifnum\dongPB=0\def\rowHeader{}\fi}{\EZtieudeTF\ifnum\dongPB=0\def\rowHeader{}\fi}
	{}{\typeout{[preamble-eduzy-v2] Không gắn được bảng Đúng/Sai (hàng tiêu đề)}}
\xpatchcmd{\writeChoiceTFt}{\hline\xdef\dotEX{\sdotEX}}{\hline\EZmauhangTF\xdef\dotEX{\sdotEX}}
	{}{\typeout{[preamble-eduzy-v2] Không gắn được bảng Đúng/Sai (tô hàng)}}
\xpatchcmd{\writeChoiceTFt}{{\end{longtable}}}{{\end{longtable}\EZketthucTF}}
	{}{\typeout{[preamble-eduzy-v2] Không gắn được bảng Đúng/Sai (cuối longtable)}}
\xpatchcmd{\writeChoiceTFt}{{\end{xtabular}\end{center}\medskip\noindent}}{{\end{xtabular}\end{center}\EZketthucTF\medskip\noindent}}
	{}{\typeout{[preamble-eduzy-v2] Không gắn được bảng Đúng/Sai (cuối xtabular)}}
\fi
%=====================================================================
\renewcommand{\parallel}{\mathrel{/\!\!/}}

%      \TLdapso                      % tài liệu: \TLfull | \TLdapso | \TLan
%      \begin{document} ... \end{document}
%  Logo Eduzy vẽ bằng TikZ ngay trong file (\eduzylogo), không cần file ảnh.
%  Biên dịch: pdfLaTeX (máy có gói vietnam) hoặc XeLaTeX.
%============================================================
\documentclass[12pt,a4paper,twoside]{article}%<-- twoside: bat dau trang/chan trang chan-le
\usepackage{iftex}
\usepackage{amsmath,amssymb,tikz,mathrsfs,tkz-tab}
\usepackage{graphicx}
\usepackage{pgfplots}% đồ thị hàm số
\pgfplotsset{compat=1.17}
%\usepackage{mathptmx}
\usepackage{mathpazo}
\ifXeTeX% XeLaTeX: font chữ có đủ dấu tiếng Việt, giữ nguyên font toán mathpazo
	\usepackage[no-math]{fontspec}
	\setmainfont{TeX Gyre Pagella}
	\newfontfamily\EZqag{TeX Gyre Adventor}[NFSSFamily=qag,Ligatures=TeX]
\fi
\usepackage[many]{tcolorbox}
\usetikzlibrary{math,through,calc,intersections,angles,quotes,shapes,shapes.geometric,arrows,patterns,matrix,chains,arrows.meta,decorations.shapes,decorations.fractals,decorations.markings,shadows,decorations.pathreplacing}
\usepackage[top=1.6cm, bottom=1.8cm, left=1.8cm, right=1.8cm]{geometry}% v2: top 1.6cm (bản cũ 1.5cm) để chữ không sát dải đầu trang
\providecommand{\CheDo}{loigiai}% loigiai | dethi | solcolor | color | book
\expandafter\usepackage\expandafter[\CheDo]{ex_test}
% \tenchuong: khai báo ở file chính (trước \input) hoặc \renewcommand sau \input
%============================================================
%  [GỘP TỪ eduzy-de.tex] EDUZY - MÀU NÂU ĐỒNG / VÀNG KIM
%  + ĐẦU TRANG, CHÂN TRANG KIỂU DẢI NGANG (v2), ĐỐI XỨNG TRANG CHẴN / TRANG LẺ.
%  (không cần file eduzy-de.tex riêng nữa)
%  Khối này phải đặt SAU \usepackage[...]{ex_test} và SAU \def\tenchuong{...}
%============================================================
\usepackage{lastpage}
\usepackage{needspace}
%--- DÀN TRANG ĐỀ THI: tách câu (THÊM 05/11/2026; chỉ thêm lệnh mới, KHÔNG đổi hành vi mặc định) ---
%  Hai công tắc, mặc định tắt (= preamble chạy y như cũ). Bật bằng cách đặt TRƯỚC %============================================================
%  PREAMBLE EDUZY v2 - dùng chung cho đề thi, bài học, chuyên đề (gói ex_test)
%  Khác bản preamble-eduzy.tex ở đầu trang, chân trang: dải nền nâu đồng tràn ngang,
%  logo + Eduzy.com.vn, ô chữ vàng kem cạnh xiên, chân trang "Trang (số)";
%  logo Eduzy chìm giữa mọi trang (tắt: \logochimtat, bật lại: \logochimbat).
%  Cách dùng ở file chính (xem mau-chuan.tex, mau-chuyen-de.tex):
%      \def\CheDo{loigiai}          % loigiai | dethi | solcolor | color | book
%      \def\tenchuong{<dòng chữ trên đầu trang>}
%      \def\thuonghieu{<chữ ở chân trang lẻ>}   % không khai báo thì dùng mặc định
%      \def\tacgia{<chữ ở chân trang chẵn>}      % không khai báo thì dùng mặc định
%      \def\khauhieu{<chữ cạnh logo ở đầu trang>}  % tùy chọn, mặc định Eduzy.com.vn
%      \input{preamble-eduzy-v2}
%      \TLdapso                      % tài liệu: \TLfull | \TLdapso | \TLan
%      \begin{document} ... \end{document}
%  Logo Eduzy vẽ bằng TikZ ngay trong file (\eduzylogo), không cần file ảnh.
%  Biên dịch: pdfLaTeX (máy có gói vietnam) hoặc XeLaTeX.
%============================================================
\documentclass[12pt,a4paper,twoside]{article}%<-- twoside: bat dau trang/chan trang chan-le
\usepackage{iftex}
\usepackage{amsmath,amssymb,tikz,mathrsfs,tkz-tab}
\usepackage{graphicx}
\usepackage{pgfplots}% đồ thị hàm số
\pgfplotsset{compat=1.17}
%\usepackage{mathptmx}
\usepackage{mathpazo}
\ifXeTeX% XeLaTeX: font chữ có đủ dấu tiếng Việt, giữ nguyên font toán mathpazo
	\usepackage[no-math]{fontspec}
	\setmainfont{TeX Gyre Pagella}
	\newfontfamily\EZqag{TeX Gyre Adventor}[NFSSFamily=qag,Ligatures=TeX]
\fi
\usepackage[many]{tcolorbox}
\usetikzlibrary{math,through,calc,intersections,angles,quotes,shapes,shapes.geometric,arrows,patterns,matrix,chains,arrows.meta,decorations.shapes,decorations.fractals,decorations.markings,shadows,decorations.pathreplacing}
\usepackage[top=1.6cm, bottom=1.8cm, left=1.8cm, right=1.8cm]{geometry}% v2: top 1.6cm (bản cũ 1.5cm) để chữ không sát dải đầu trang
\providecommand{\CheDo}{loigiai}% loigiai | dethi | solcolor | color | book
\expandafter\usepackage\expandafter[\CheDo]{ex_test}
% \tenchuong: khai báo ở file chính (trước \input) hoặc \renewcommand sau \input
%============================================================
%  [GỘP TỪ eduzy-de.tex] EDUZY - MÀU NÂU ĐỒNG / VÀNG KIM
%  + ĐẦU TRANG, CHÂN TRANG KIỂU DẢI NGANG (v2), ĐỐI XỨNG TRANG CHẴN / TRANG LẺ.
%  (không cần file eduzy-de.tex riêng nữa)
%  Khối này phải đặt SAU \usepackage[...]{ex_test} và SAU \def\tenchuong{...}
%============================================================
\usepackage{lastpage}
\usepackage{needspace}
%--- DÀN TRANG ĐỀ THI: tách câu (THÊM 05/11/2026; chỉ thêm lệnh mới, KHÔNG đổi hành vi mặc định) ---
%  Hai công tắc, mặc định tắt (= preamble chạy y như cũ). Bật bằng cách đặt TRƯỚC \input{preamble-eduzy-v2}
%  trong file đề (tools/dan_trang/dan_trang.py tự ghi khi đặt \NeedspaceCau):
%    \def\TachCauDS{1}   câu Đúng/Sai được tách: câu dẫn (kèm hình) ở cuối trang, bảng Phát biểu sang trang sau
%                        (câu dẫn và hình không tách nhau, bảng không tách giữa hai trang)
%    \def\TachCauLG{1}   bản lời giải (\CheDo khác dethi): phần đề của câu cao hơn 1/3 trang được phép tách
%  Lệnh dùng trong file đề (giá trị do tools/dan_trang/dan_trang.py đo ở bản dethi):
%    \NeedspaceCau[<cao câu dẫn + hình>]{<cao cả khối câu>}   ngay trên \begin{ex} (câu đầu phần: trên tiêu đề phần)
%        công tắc tắt: đúng bằng \Needspace{<cao cả khối câu>} như cũ;
%        \TachCauDS=1 và có [..] (câu Đúng/Sai được chọn tách): chỉ giữ câu dẫn + hình;
%        \TachCauLG=1 ở bản lời giải: đòi tối đa 1/3 \textheight.
%    \NeedspaceBang{<cao bảng Phát biểu>}%   ngay trước \choiceTFt: \TachCauDS=1 thì giữ trọn bảng trên một trang;
%        công tắc tắt: không làm gì.
\providecommand{\TachCauDS}{0}
\providecommand{\TachCauLG}{0}
\newlength{\EZnsCau}
\newcommand{\NeedspaceCau}[2][]{%
	\setlength{\EZnsCau}{#2}%
	\ifnum\TachCauDS=1 \if\relax\detokenize{#1}\relax\else\setlength{\EZnsCau}{#1}\fi\fi
	\ifnum\TachCauLG=1 \ifdefstring{\kindOPTN}{dethi}{}{\ifdim\EZnsCau>0.3333\textheight\setlength{\EZnsCau}{0.3333\textheight}\fi}\fi
	\Needspace{\EZnsCau}}
\newcommand{\NeedspaceBang}[1]{\ifnum\TachCauDS=1 \Needspace{#1}\fi}

%--- Thang màu: đậm -> nhạt
\definecolor{gold1}{RGB}{126,96,0}    % nâu đồng đậm  #7E6000
\definecolor{gold2}{RGB}{192,144,3}   % vàng kim      #C09003
\definecolor{gold3}{RGB}{255,215,99}  % vàng nhạt     #FFD763
\definecolor{gold4}{RGB}{255,231,153} % vàng rất nhạt #FFE799
\definecolor{gold5}{RGB}{255,243,205} % vàng kem      #FFF3CD
%--- Ghi đè tông xanh/đỏ cũ
\definecolor{myblue}{RGB}{126,96,0}
\colorlet{blue}{gold1}
\colorlet{cyan}{gold2}
\colorlet{teal}{gold2}

%--- Tên chương mặc định (file đề nên tự khai báo lại)
\providecommand{\tenchuong}{}
%--- Chân trang
\providecommand{\thuonghieu}{TOÁN 12 -- Eduzy.com.vn}
\providecommand{\tacgia}{Tô Gia Bảo -- 0353703783}

%--- Chữ cạnh logo ở đầu trang
\providecommand{\khauhieu}{Eduzy.com.vn}

%====================== THÔNG SỐ ĐẦU TRANG, CHÂN TRANG (v2) ======================
%  Dải nền đậm tràn ngang trang; ô chữ vàng kem có cạnh xiên, kèm 2 vạch xiên chuyển màu.
%  Đầu trang lẻ : logo + \khauhieu (trái)   | ô chữ \tenchuong (phải)
%  Chân trang lẻ: ô chữ \thuonghieu (trái)  | Trang (số) (phải)
%  Trang chẵn soi gương trang lẻ; chân trang chẵn ghi \tacgia.
%--- Màu: chỉ dùng thang vàng kim Eduzy
\colorlet{ezNen}{gold1}% dải nền đậm
\colorlet{ezVach1}{gold3}% vạch xiên sát dải nền (đậm hơn)
\colorlet{ezVach2}{gold4}% vạch xiên sát ô chữ (nhạt hơn)
\colorlet{ezO}{gold5}% ô chữ vàng kem
%--- Kích thước
\newlength{\EZdai}\setlength{\EZdai}{1.1cm}% chiều cao dải nền
\newlength{\EZo}\setlength{\EZo}{0.68cm}% chiều cao ô chữ
\newlength{\EZxien}\setlength{\EZxien}{0.55cm}% độ lệch ngang của cạnh xiên
\newlength{\EZvach}\setlength{\EZvach}{0.75cm}% bề rộng mỗi vạch xiên
\newlength{\EZkhe}\setlength{\EZkhe}{0.12cm}% khe trắng giữa vạch và ô chữ
\newlength{\EZdem}\setlength{\EZdem}{0.4cm}% khoảng đệm từ chữ tới cạnh xiên
\newlength{\EZle}\setlength{\EZle}{1.8cm}% lề trái/phải của chữ (bằng lề thân bài)
\newlength{\EZrongmin}\setlength{\EZrongmin}{4.5cm}% ô chữ rộng tối thiểu
\newlength{\EZrong}% bề rộng chữ trong ô (tính tự động)
%--- Font
\newcommand{\EZfontO}{\fontsize{10.5pt}{12pt}\selectfont\bfseries}% chữ trong ô
\newcommand{\EZfontKH}{\fontfamily{qag}\fontsize{11pt}{12pt}\selectfont\bfseries}% \khauhieu cạnh logo
\newcommand{\EZfontTrang}{\fontsize{11.5pt}{12pt}\selectfont\bfseries}% chữ "Trang" và số trang
%--- Đo bề rộng chữ trong ô
\newcommand{\EZdorong}[1]{%
	\settowidth{\EZrong}{\EZfontO #1}%
	\ifdim\EZrong<\EZrongmin\setlength{\EZrong}{\EZrongmin}\fi}
%--- Logo + chữ trắng \khauhieu (dựng sẵn một lần khi vào \begin{document})
\newsavebox{\EZlogobox}
\AtBeginDocument{\sbox{\EZlogobox}{%
		\raisebox{-0.5\height}{\eduzylogo[1cm]}\hspace{0.15cm}%
		\raisebox{-0.5\height}{\color{white}\EZfontKH\khauhieu}}}
%--- Hai vạch xiên + khe trắng. #1: điểm gốc ở mép dưới, #2: hướng (1 sang phải, -1 sang trái),
%    #3: hướng nghiêng (1 kiểu "/", -1 kiểu "\")
\newcommand{\EZvachxien}[3]{%
	\fill[white] (#1) --++({(#3)*\EZxien},\EZo) --++({(#2)*\EZkhe},0) --++({-(#3)*\EZxien},-\EZo) --cycle;
	\fill[ezVach2] ([xshift=(#2)*\EZkhe]#1) --++({(#3)*\EZxien},\EZo) --++({(#2)*\EZvach},0) --++({-(#3)*\EZxien},-\EZo) --cycle;
	\fill[ezVach1] ([xshift=(#2)*(\EZkhe+\EZvach)]#1) --++({(#3)*\EZxien},\EZo) --++({(#2)*\EZvach},0) --++({-(#3)*\EZxien},-\EZo) --cycle;
	\draw[white,line width=0.6pt] ([xshift=(#2)*(\EZkhe+\EZvach)]#1) --++({(#3)*\EZxien},\EZo);
}
%--- Chữ trong ô (căn giữa theo chiều cao ô, không lệch vì dấu tiếng Việt)
\tikzset{ezchu/.style={inner sep=0pt,text=ezNen,text height=0.72em,text depth=0.22em}}

%====================== LOGO CHÌM GIỮA TRANG ======================
%  Vẽ cùng đầu trang nên nằm dưới chữ của thân bài, có trên mọi trang.
%  Tắt: \logochimtat (từ trang đó trở đi), bật lại: \logochimbat.
%  Đổi cỡ, độ đậm: khai báo trước \input, ví dụ \def\EZchimcao{9cm} \def\EZchimdo{0.1}
\providecommand{\EZchimcao}{10cm}% chiều cao logo chìm
\providecommand{\EZchimdo}{0.12}% độ đậm (0 = trong suốt, 1 = rõ hẳn)
\newif\ifEZchim \EZchimtrue
\newcommand{\logochimtat}{\EZchimfalse}
\newcommand{\logochimbat}{\EZchimtrue}
%  Dựng sẵn logo mờ một lần khi vào \begin{document}; transparency group làm mờ cả logo
%  một lượt nên không lộ vệt ở chỗ các dải màu chồng nhau.
\newsavebox{\EZchimtho}
\newsavebox{\EZchimbox}
\AtBeginDocument{%
	\sbox{\EZchimtho}{\eduzylogo[\EZchimcao]}%
	\sbox{\EZchimbox}{\begin{tikzpicture}
			\begin{scope}[transparency group,opacity=\EZchimdo]
				\node[inner sep=0pt] {\usebox{\EZchimtho}};
			\end{scope}
	\end{tikzpicture}}}
\newcommand{\EZvechim}{%
	\ifEZchim
		\node[inner sep=0pt] at (current page.center) {\usebox{\EZchimbox}};
	\fi}

%====================== ĐẦU TRANG - TRANG LẺ ======================
%  (logo bên trái, ô chữ \tenchuong bên phải)
\newcommand{\myfancyheadOdd}{%
	\EZdorong{\tenchuong}%
	\begin{tikzpicture}[remember picture,overlay]
		\coordinate (T) at (current page.north west);
		\coordinate (P) at (current page.north east);
		%--- logo chìm giữa trang
		\EZvechim
		%--- dải nền
		\fill[ezNen] (T) rectangle ([yshift=-\EZdai]P);
		%--- ô chữ: X là góc dưới bên trái, cạnh trái xiên kiểu "/"
		\coordinate (X) at ([xshift=-\EZle-\EZrong-\EZdem-\EZxien,yshift=-\EZdai]P);
		\fill[ezO] (X) --++(\EZxien,\EZo) -- ([yshift=-\EZdai+\EZo]P) -- ([yshift=-\EZdai]P) --cycle;
		\EZvachxien{X}{-1}{1}
		\node[ezchu,anchor=east] at ([xshift=-\EZle,yshift=-\EZdai+0.5*\EZo]P) {\EZfontO\tenchuong};
		%--- logo + tên
		\node[inner sep=0pt,anchor=west] at ([xshift=\EZle,yshift=-0.5*\EZdai]T) {\usebox{\EZlogobox}};
	\end{tikzpicture}%
}

%====================== ĐẦU TRANG - TRANG CHẴN =====================
%  (soi gương: ô chữ \tenchuong bên trái, logo bên phải)
\newcommand{\myfancyheadEven}{%
	\EZdorong{\tenchuong}%
	\begin{tikzpicture}[remember picture,overlay]
		\coordinate (T) at (current page.north west);
		\coordinate (P) at (current page.north east);
		%--- logo chìm giữa trang
		\EZvechim
		%--- dải nền
		\fill[ezNen] (T) rectangle ([yshift=-\EZdai]P);
		%--- ô chữ: X là góc dưới bên phải, cạnh phải xiên kiểu "\"
		\coordinate (X) at ([xshift=\EZle+\EZrong+\EZdem+\EZxien,yshift=-\EZdai]T);
		\fill[ezO] (X) --++(-\EZxien,\EZo) -- ([yshift=-\EZdai+\EZo]T) -- ([yshift=-\EZdai]T) --cycle;
		\EZvachxien{X}{1}{-1}
		\node[ezchu,anchor=west] at ([xshift=\EZle,yshift=-\EZdai+0.5*\EZo]T) {\EZfontO\tenchuong};
		%--- logo + tên
		\node[inner sep=0pt,anchor=east] at ([xshift=-\EZle,yshift=-0.5*\EZdai]P) {\usebox{\EZlogobox}};
	\end{tikzpicture}%
}

%====================== CHÂN TRANG - TRANG LẺ ======================
%  (ô chữ \thuonghieu bên trái, "Trang (số)" bên phải)
\newcommand{\myfancyfootOdd}{%
	\EZdorong{\thuonghieu}%
	\begin{tikzpicture}[remember picture,overlay]
		\coordinate (T) at (current page.south west);
		\coordinate (P) at (current page.south east);
		%--- dải nền
		\fill[ezNen] (T) rectangle ([yshift=\EZdai]P);
		%--- ô chữ sát mép trên dải: X là góc dưới bên phải, cạnh phải xiên kiểu "/"
		\coordinate (X) at ([xshift=\EZle+\EZrong+\EZdem,yshift=\EZdai-\EZo]T);
		\fill[ezO] ([yshift=\EZdai-\EZo]T) -- (X) --++(\EZxien,\EZo) -- ([yshift=\EZdai]T) --cycle;
		\EZvachxien{X}{1}{1}
		\node[ezchu,anchor=west] at ([xshift=\EZle,yshift=\EZdai-0.5*\EZo]T) {\EZfontO\thuonghieu};
		%--- số trang
		\node[circle,draw=white,line width=0.9pt,text=white,inner sep=0pt,minimum size=0.62cm,anchor=east]
		(so) at ([xshift=-\EZle,yshift=0.5*\EZdai]P) {\EZfontTrang\thepage};
		\node[inner sep=0pt,text=white,anchor=base east] (tr) at ([xshift=-0.15cm]so.base west) {\EZfontTrang Trang};
	\end{tikzpicture}%
}

%====================== CHÂN TRANG - TRANG CHẴN ====================
%  (soi gương: "Trang (số)" bên trái, ô chữ \tacgia bên phải)
\newcommand{\myfancyfootEven}{%
	\EZdorong{\tacgia}%
	\begin{tikzpicture}[remember picture,overlay]
		\coordinate (T) at (current page.south west);
		\coordinate (P) at (current page.south east);
		%--- dải nền
		\fill[ezNen] (T) rectangle ([yshift=\EZdai]P);
		%--- ô chữ sát mép trên dải: X là góc dưới bên trái, cạnh trái xiên kiểu "\"
		\coordinate (X) at ([xshift=-\EZle-\EZrong-\EZdem,yshift=\EZdai-\EZo]P);
		\fill[ezO] ([yshift=\EZdai-\EZo]P) -- (X) --++(-\EZxien,\EZo) -- ([yshift=\EZdai]P) --cycle;
		\EZvachxien{X}{-1}{-1}
		\node[ezchu,anchor=east] at ([xshift=-\EZle,yshift=\EZdai-0.5*\EZo]P) {\EZfontO\tacgia};
		%--- số trang
		%    (vòng tròn ẩn so0 nằm giữa dải để lấy đường chân chữ, giống trang lẻ)
		\node[circle,inner sep=0pt,minimum size=0.62cm] (so0) at ([yshift=0.5*\EZdai]T) {\EZfontTrang\phantom{\thepage}};
		\node[inner sep=0pt,text=white,anchor=base west] (tr) at ($(so0.base -| T)+(\EZle,0)$) {\EZfontTrang Trang};
		\node[circle,draw=white,line width=0.9pt,text=white,inner sep=0pt,minimum size=0.62cm,anchor=base west]
		(so) at ([xshift=0.15cm]tr.base east) {\EZfontTrang\thepage};
	\end{tikzpicture}%
}

%--- Giữ tương thích với các file cũ đang gọi \myfancyhead, \myfancyfoot
\newcommand{\myfancyhead}{\myfancyheadOdd}
\newcommand{\myfancyfoot}{\myfancyfootOdd}

%======================= GẮN VÀO TRANG =============================
\pagestyle{fancy}
\fancyhf{}
\fancyhead[LO]{\myfancyheadOdd}
\fancyhead[RE]{\myfancyheadEven}
\fancyfoot[LO]{\myfancyfootOdd}
\fancyfoot[RE]{\myfancyfootEven}
\renewcommand{\headrulewidth}{0pt}
\renewcommand{\footrulewidth}{0pt}
\fancypagestyle{plain}{%
	\fancyhf{}
	\fancyhead[LO]{\myfancyheadOdd}
	\fancyhead[RE]{\myfancyheadEven}
	\fancyfoot[LO]{\myfancyfootOdd}
	\fancyfoot[RE]{\myfancyfootEven}
	\renewcommand{\headrulewidth}{0pt}
	\renewcommand{\footrulewidth}{0pt}
}
%=================== HẾT PHẦN GỘP TỪ eduzy-de.tex ===================
\newcommand{\indam}[1]{
	{\fontsize{12pt}{9pt}\sffamily \selectfont\bfseries\color{gold1} #1}
}
\newcommand{\ind}[1]{
	{\fontsize{14pt}{9pt}\fontfamily{qag} \selectfont\bfseries\color{gold1} #1}
}
\renewcommand{\FalseEX}{\stepcounter{dapan}\textbf{\selectfont\color{gold1}\Alph{dapan}.}\ignorespaces}
\renewtheorem{ex}{\color{gold1}Câu}
\renewcommand{\baselinestretch}{1.1}
\newcommand{\hoac}[1]{
	\left[\begin{aligned}#1\end{aligned}\right.}
\newcommand{\heva}[1]{
	\left\{\begin{aligned}#1\end{aligned}\right.}
\usepackage{esvect}
\def\vec{\vv}
\def\vt{\vv}
\def\overrightarrow{\vv}
%====================== LOGO EDUZY VẼ BẰNG TIKZ ======================
%  Đường viền dò từ logo gốc; dải màu vàng kim đo trực tiếp trên ảnh gốc
%  (nghiêng 21 độ, tô bằng 80 dải hẹp nên hiển thị giống nhau trên mọi trình xem PDF).
%  Dùng: \eduzylogo          -> cao 2.4cm (mặc định)
%        \eduzylogo[1.5cm]   -> cao 1.5cm
\definecolor{ezText}{RGB}{245,200,46}
\newcommand{\eduzylogo}[1][2.4cm]{%
	\resizebox{!}{#1}{%
		\begin{tikzpicture}[x=0.01cm,y=0.01cm]
			\path[use as bounding box] (0,0) rectangle (224,-224);
			% mũ tốt nghiệp + vương miện
			\begin{scope}
				\clip
				(119.4,-176.1)..controls(118.9,-176.1)and(116.8,-176)..(114.6,-176)
				..controls(101.4,-176)and(85.8,-175.9)..(83.7,-175.7)
				..controls(81.8,-175.6)and(81,-175.6)..(80.2,-175.7)
				..controls(78.8,-175.9)and(77.7,-175.9)..(76.2,-175.7)
				..controls(74.9,-175.5)and(70.8,-175.5)..(64.3,-175.7)
				..controls(59.1,-175.9)and(58.3,-175.7)..(58.6,-173.7)
				..controls(58.6,-172.8)and(59.6,-169.6)..(60.3,-168)
				..controls(60.8,-166.9)and(61.4,-165.8)..(62.3,-165)
				..controls(63.1,-164.2)and(64.3,-163.4)..(65,-163.2)
				..controls(65.5,-163.1)and(72.3,-163)..(76.6,-163.1)
				..controls(77.9,-163.1)and(80.1,-163.1)..(81.4,-163)
				..controls(82.8,-162.9)and(86.3,-162.9)..(89.2,-162.9)
				..controls(93,-162.9)and(94.9,-162.8)..(96,-162.7)
				..controls(97.2,-162.5)and(101,-162.5)..(114,-162.4)
				..controls(123,-162.3)and(134.9,-162.2)..(140.2,-162.1)
				..controls(150.1,-162)and(152.2,-162.1)..(153.6,-162.4)
				..controls(154.4,-162.6)and(157.1,-163.8)..(158.1,-164.4)
				..controls(158.8,-164.8)and(159.7,-165.7)..(159.9,-166.3)
				..controls(160,-166.5)and(160.2,-167.1)..(160.2,-167.6)
				..controls(160.3,-168)and(160.6,-169.4)..(160.9,-170.6)
				..controls(161.2,-171.7)and(161.5,-173.1)..(161.6,-173.6)
				..controls(161.8,-174.4)and(161.8,-174.5)..(161.5,-174.9)
				..controls(161.2,-175.6)and(160.8,-175.9)..(159.9,-176)
				..controls(159.2,-176.1)and(154.1,-176.2)..(151.5,-176.2)
				..controls(148.1,-176.1)and(122.6,-176.1)..(121.6,-176.2)
				..controls(120.8,-176.2)and(119.8,-176.2)..(119.4,-176.1)--cycle
				(130.4,-160)..controls(129,-160)and(119.8,-160)..(110.1,-159.9)
				..controls(100.3,-159.8)and(88.5,-159.7)..(83.9,-159.6)
				..controls(77.8,-159.5)and(75.3,-159.5)..(74.5,-159.6)
				..controls(73.9,-159.7)and(72.5,-159.8)..(71.3,-159.8)
				..controls(69.6,-159.8)and(69,-159.8)..(68.5,-159.6)
				..controls(67.6,-159.3)and(66.8,-158.6)..(66.2,-157.6)
				..controls(66,-157.3)and(65.6,-156.5)..(65.3,-156.1)
				..controls(64.6,-155.1)and(64.5,-154.9)..(63.7,-152.9)
				..controls(63.2,-151.6)and(62.9,-150.9)..(61.8,-149.6)
				..controls(61.5,-149.2)and(61.1,-148.5)..(60.9,-147.9)
				..controls(60.7,-147.4)and(60.3,-146.6)..(60.1,-146.1)
				..controls(59.9,-145.5)and(59.5,-144.9)..(59,-144.3)
				..controls(58.5,-143.7)and(58.2,-143.2)..(57.9,-142.5)
				..controls(57.7,-141.9)and(57.3,-141.1)..(57,-140.7)
				..controls(56.6,-140.3)and(56.1,-139.3)..(55.8,-138.5)
				..controls(55.5,-137.8)and(55.1,-136.8)..(54.9,-136.5)
				..controls(54.7,-136.1)and(54.3,-135.2)..(53.9,-134.5)
				..controls(53.6,-133.8)and(53,-132.7)..(52.6,-132)
				..controls(51.6,-130.5)and(51.4,-130.2)..(51.1,-129)
				..controls(50.6,-127.5)and(50.6,-127.4)..(49.7,-126.3)
				..controls(49.2,-125.7)and(48.6,-124.8)..(48.3,-124.2)
				..controls(48,-123.6)and(47.5,-122.7)..(47.3,-122.2)
				..controls(47,-121.8)and(46.5,-120.7)..(46.1,-119.8)
				..controls(45.8,-118.9)and(45.2,-117.8)..(44.9,-117.2)
				..controls(44.2,-116.1)and(43.3,-114.4)..(42.3,-112.2)
				..controls(41.8,-111.3)and(41.3,-110.4)..(41.1,-110.1)
				..controls(40.5,-109.4)and(40.3,-109)..(40,-107.8)
				..controls(39.8,-107)and(39.6,-106.5)..(39.2,-106)
				..controls(38.9,-105.6)and(38.4,-104.7)..(38,-104)
				..controls(37.7,-103.4)and(37.2,-102.5)..(37,-102.1)
				..controls(36.7,-101.7)and(36.3,-100.8)..(36,-100.1)
				..controls(35.6,-98.9)and(35.3,-98.4)..(34.6,-97.6)
				..controls(34.2,-97.1)and(34.2,-96.3)..(34.6,-96)
				..controls(35,-95.7)and(35.5,-95.9)..(36.2,-96.7)
				..controls(36.5,-97)and(37.2,-97.5)..(37.6,-97.7)..controls(38,-98)and(38.6,-98.4)..(39.1,-98.9)
				..controls(39.6,-99.4)and(40.3,-100)..(40.6,-100.3)
				..controls(41.3,-100.8)and(42.4,-101.8)..(43.7,-102.9)
				..controls(44,-103.2)and(44.6,-103.7)..(45,-103.9)
				..controls(45.4,-104.2)and(45.9,-104.6)..(46.1,-104.8)
				..controls(47.1,-105.9)and(48.8,-107.1)..(50.5,-107.9)
				..controls(51.8,-108.6)and(52.3,-109)..(53.4,-110)
				..controls(54.8,-111.2)and(55.6,-111.6)..(57,-112.1)
				..controls(57.5,-112.3)and(58.5,-112.8)..(59.4,-113.3)
				..controls(60.2,-113.7)and(61.3,-114.3)..(61.9,-114.5)
				..controls(62.4,-114.7)and(63.3,-115.1)..(63.8,-115.3)
				..controls(64.4,-115.6)and(65.4,-115.9)..(66.2,-116.1)
				..controls(67,-116.3)and(68.4,-116.6)..(69.2,-116.9)
				..controls(72.3,-117.9)and(72.8,-118)..(78,-118.1)
				..controls(84,-118.3)and(86.1,-118.1)..(88.1,-117.1)
				..controls(89.4,-116.5)and(93,-115)..(93.7,-114.9)
				..controls(94.8,-114.6)and(95.4,-114.2)..(95.9,-113.4)
				..controls(96.2,-113)and(96.7,-112.5)..(97.1,-112.2)
				..controls(98,-111.6)and(101.3,-108.4)..(102.1,-107.2)
				..controls(102.4,-106.8)and(102.9,-106.2)..(103.2,-105.8)
				..controls(103.5,-105.5)and(104,-104.8)..(104.2,-104.2)
				..controls(104.5,-103.6)and(105,-102.8)..(105.3,-102.4)
				..controls(105.6,-102)and(106.1,-101.2)..(106.5,-100.6)
				..controls(107.1,-99.5)and(108.6,-97.5)..(109.5,-96.5)
				..controls(109.9,-96.1)and(110.1,-96)..(110.4,-96)
				..controls(111.1,-96)and(111.6,-96.5)..(111.8,-97.6)
				..controls(112,-98.6)and(112.3,-99.3)..(112.8,-99.9)
				..controls(113.7,-100.9)and(115.2,-103.1)..(115.7,-104.2)
				..controls(116.1,-105)and(116.5,-105.5)..(116.9,-106)
				..controls(117.3,-106.4)and(117.8,-106.9)..(118,-107.3)
				..controls(118.2,-107.6)and(118.6,-108)..(119,-108.3)
				..controls(119.6,-108.8)and(121.1,-110.3)..(121.8,-111.2)
				..controls(122.1,-111.7)and(122.5,-112)..(123.4,-112.4)
				..controls(124.8,-113.2)and(125.8,-113.8)..(127,-114.6)
				..controls(127.5,-114.9)and(128.1,-115.3)..(128.2,-115.4)
				..controls(128.4,-115.5)and(129.3,-115.8)..(130.2,-116)
				..controls(131.1,-116.2)and(132.5,-116.6)..(133.3,-116.9)
				..controls(135.7,-117.9)and(136.1,-117.9)..(142.5,-117.8)
				..controls(145.8,-117.8)and(148.4,-117.7)..(148.8,-117.6)
				..controls(149.2,-117.6)and(150.2,-117.2)..(151,-116.8)
				..controls(152,-116.4)and(152.8,-116)..(153.6,-115.9)
				..controls(154.5,-115.7)and(155.2,-115.4)..(156.3,-114.9)
				..controls(157.1,-114.5)and(158.2,-114)..(158.7,-113.8)
				..controls(159.1,-113.7)and(159.8,-113.3)..(160.2,-113)
				..controls(160.6,-112.7)and(161.3,-112.4)..(161.7,-112.2)
				..controls(162.7,-111.9)and(163.2,-111.7)..(164.1,-111.1)
				..controls(164.5,-110.8)and(165.3,-110.4)..(165.9,-110)
				..controls(166.5,-109.7)and(167.3,-109.2)..(167.6,-108.9)
				..controls(167.9,-108.6)and(168.6,-108.2)..(169,-108)
				..controls(169.4,-107.8)and(170.2,-107.3)..(170.7,-106.9)
				..controls(171.2,-106.5)and(172.1,-105.8)..(172.7,-105.4)
				..controls(173.2,-105)and(174.1,-104.4)..(174.5,-104)
				..controls(174.9,-103.7)and(175.7,-103.1)..(176.2,-102.6)
				..controls(176.7,-102.2)and(177.4,-101.5)..(177.9,-101.2)
				..controls(178.3,-100.8)and(178.8,-100.4)..(179,-100.2)
				..controls(179.2,-100)and(179.6,-99.6)..(180,-99.3)
				..controls(180.4,-99.1)and(180.9,-98.6)..(181.2,-98.4)
				..controls(181.5,-98.1)and(182.1,-97.7)..(182.5,-97.5)
				..controls(183.2,-97.2)and(183.5,-96.9)..(184.3,-95.9)
				..controls(184.6,-95.6)and(184.6,-95.6)..(184.8,-95.7)
				..controls(185,-95.8)and(185.1,-96)..(185.1,-96.2)
				..controls(185.1,-96.5)and(185,-96.7)..(184.6,-97.2)
				..controls(184.1,-97.8)and(184,-97.9)..(183.9,-98.8)
				..controls(183.7,-100.3)and(183.3,-101.3)..(182.9,-101.8)
				..controls(182.6,-102.1)and(182.3,-102.6)..(182.2,-102.9)
				..controls(182,-103.3)and(181.6,-104)..(181.2,-104.6)
				..controls(180.9,-105.2)and(180.5,-106.1)..(180.4,-106.6)
				..controls(179.9,-108)and(179.7,-108.5)..(179,-109.4)
				..controls(178.6,-109.8)and(178.2,-110.5)..(177.9,-111)
				..controls(177.7,-111.5)and(177.3,-112.3)..(177,-112.7)
				..controls(176.6,-113.2)and(176.4,-113.7)..(176.3,-114.2)
				..controls(176.2,-114.7)and(175.9,-115.2)..(175.4,-116)
				..controls(175,-116.6)and(174.4,-117.7)..(174,-118.5)
				..controls(173.7,-119.2)and(173.3,-120.1)..(173,-120.5)
				..controls(172.8,-120.9)and(172.4,-121.8)..(172.1,-122.5)
				..controls(171.8,-123.2)and(171.3,-124.1)..(171,-124.6)
				..controls(170.8,-125.1)and(170.3,-126)..(169.9,-126.6)
				..controls(169.6,-127.3)and(169.2,-128.1)..(169,-128.3)
				..controls(168.8,-128.6)and(168.5,-129.3)..(168.3,-129.8)
				..controls(168.1,-130.3)and(167.8,-131.1)..(167.5,-131.6)
				..controls(166.9,-132.5)and(164.2,-137.9)..(164.2,-138.1)
				..controls(164.2,-138.1)and(163.9,-138.8)..(163.4,-139.7)
				..controls(162.9,-140.5)and(162.4,-141.5)..(162.3,-141.9)
				..controls(161.9,-143)and(161.6,-143.4)..(160.9,-144.1)
				..controls(160.4,-144.6)and(160.3,-144.8)..(160.1,-145.4)
				..controls(160,-145.8)and(159.6,-146.8)..(159.2,-147.8)
				..controls(158.7,-148.8)and(158.1,-150)..(157.9,-150.6)
				..controls(157.6,-151.1)and(157.2,-151.9)..(156.9,-152.4)
				..controls(156.5,-152.8)and(156.1,-153.7)..(155.8,-154.4)
				..controls(155.6,-155)and(155.2,-155.9)..(154.9,-156.3)
				..controls(154.7,-156.7)and(154.4,-157.4)..(154.2,-157.9)
				..controls(154.1,-158.4)and(153.8,-159)..(153.5,-159.3)
				..controls(152.8,-160.1)and(153,-160.1)..(142.4,-160.1)
				..controls(137.2,-160.1)and(131.8,-160.1)..(130.4,-160)--cycle
				(141.5,-103.7)..controls(140.8,-103)and(139.8,-102.1)..(139.2,-101.6)
				..controls(138.5,-101.2)and(137.9,-100.6)..(137.8,-100.4)
				..controls(137,-99.5)and(134,-96.5)..(133.2,-96)
				..controls(132.7,-95.7)and(131.8,-94.9)..(131.1,-94.3)
				..controls(129.4,-92.7)and(127.9,-91.6)..(126.6,-91.1)
				..controls(125.2,-90.5)and(124.1,-90)..(122.9,-89.1)
				..controls(122.1,-88.6)and(121.3,-88.2)..(119.9,-87.7)
				..controls(118.8,-87.3)and(117.5,-86.8)..(116.9,-86.5)
				..controls(114.3,-85.5)and(113.7,-85.4)..(109.7,-85.5)
				..controls(105.3,-85.7)and(103.9,-85.9)..(101.6,-87)
				..controls(100.7,-87.4)and(99.4,-87.8)..(97.9,-88)
				..controls(97,-88.2)and(96.6,-88.3)..(95.8,-88.9)
				..controls(95.4,-89.2)and(94.7,-89.6)..(94.2,-89.8)
				..controls(92.9,-90.3)and(92.3,-90.7)..(90.9,-92.1)
				..controls(90.2,-92.7)and(89.4,-93.3)..(89,-93.5)..controls(87,-94.5)and(85.2,-96)..(83.6,-97.9)
				..controls(83.3,-98.2)and(82.5,-98.9)..(81.8,-99.4)
				..controls(80.7,-100.1)and(80.5,-100.3)..(80,-101.1)
				..controls(78.7,-103.1)and(77.2,-103.9)..(76.4,-103)
				..controls(75.7,-102.2)and(75.7,-103.2)..(75.7,-83.6)
				..controls(75.6,-73.6)and(75.6,-64.8)..(75.5,-64)
				..controls(75.5,-63.2)and(75.5,-61.4)..(75.6,-59.8)--(75.7,-57.1)--(76,-56.8)
				..controls(76.8,-56.2)and(78,-56.3)..(79.8,-57.2)
				..controls(80.4,-57.5)and(81.3,-57.8)..(81.8,-57.9)
				..controls(82.3,-58.1)and(83,-58.4)..(83.6,-58.8)..controls(84.3,-59.2)and(85,-59.5)..(86,-59.8)
				..controls(86.8,-60.1)and(88,-60.6)..(88.7,-60.9)
				..controls(89.3,-61.2)and(90.6,-61.7)..(91.4,-62)
				..controls(92.2,-62.3)and(93.7,-62.9)..(94.6,-63.3)
				..controls(95.6,-63.7)and(96.8,-64.1)..(97.4,-64.3)
				..controls(98,-64.5)and(99,-64.9)..(99.7,-65.2)
				..controls(100.4,-65.6)and(101.5,-66)..(102.3,-66.2)
				..controls(104.2,-66.8)and(105,-67.1)..(106.2,-67.8)
				..controls(107.9,-68.7)and(110.2,-69.4)..(111.3,-69.2)
				..controls(112.5,-69)and(115,-68.2)..(115.6,-67.8)
				..controls(117.3,-66.7)and(117.6,-66.6)..(119.5,-66.2)
				..controls(121.9,-65.8)and(122.1,-65.7)..(123,-65.1)
				..controls(124.1,-64.4)and(124.4,-64.3)..(126,-63.9)
				..controls(127,-63.7)and(127.7,-63.4)..(128.4,-63.1)
				..controls(129,-62.8)and(129.7,-62.5)..(130.6,-62.4)
				..controls(132.1,-62.1)and(134,-61.5)..(135.2,-61)
				..controls(135.7,-60.8)and(136.7,-60.5)..(137.3,-60.3)
				..controls(138.4,-60)and(139.3,-59.7)..(141.1,-59)
				..controls(141.6,-58.7)and(142,-58.6)..(142.5,-58.6)
				..controls(143.2,-58.6)and(143.3,-58.7)..(143.6,-59)
				..controls(144,-59.5)and(144.1,-59.9)..(144,-60.8)
				..controls(144,-61.2)and(144,-71.1)..(144,-82.8)--(143.9,-104.1)--(143.7,-104.5)
				..controls(143.5,-104.7)and(143.3,-104.8)..(143,-104.9)
				..controls(142.6,-104.9)and(142.5,-104.8)..(141.5,-103.7)--cycle
				(27.8,-91.4)..controls(27.1,-90.9)and(26.3,-90)..(26.1,-89.5)
				..controls(26,-89.2)and(25.8,-88.6)..(25.7,-87.9)
				..controls(25.4,-86.2)and(25.6,-84.8)..(26.5,-83.5)
				..controls(27.4,-82.3)and(28.3,-82)..(29.4,-82.3)
				..controls(31.1,-82.8)and(32.1,-84.9)..(32,-87.4)
				..controls(31.9,-89.3)and(31.2,-90.6)..(29.9,-91.4)
				..controls(29.1,-91.9)and(28.6,-91.9)..(27.8,-91.4)--cycle
				(190.1,-91.3)..controls(189.7,-91.2)and(189.3,-90.9)..(189.1,-90.7)
				..controls(188.3,-89.8)and(187.6,-87.7)..(187.6,-86.2)
				..controls(187.6,-85.5)and(187.7,-85.2)..(188.2,-84.4)
				..controls(189.1,-82.6)and(190.1,-81.8)..(191.3,-82.1)
				..controls(192,-82.3)and(193.5,-83.7)..(193.8,-84.4)
				..controls(194.1,-85.1)and(194.2,-87.1)..(194,-88.3)
				..controls(193.8,-89.8)and(193.3,-90.5)..(192.3,-91.3)
				..controls(191.7,-91.7)and(190.9,-91.7)..(190.1,-91.3)--cycle
				(173.7,-83.4)..controls(172.2,-82.5)and(172,-82.5)..(170.6,-82.4)
				..controls(169.7,-82.4)and(168.8,-82.4)..(168.3,-82.5)
				..controls(167.2,-82.7)and(166.9,-82.7)..(166.6,-82.2)
				..controls(166.3,-81.8)and(166.3,-81)..(166.8,-80)
				..controls(167.2,-78.9)and(167.4,-77.4)..(167.7,-74.1)
				..controls(167.8,-72.8)and(167.9,-71.8)..(167.9,-71.7)
				..controls(168,-71.7)and(168,-70.2)..(168,-68.4)..controls(168,-65.4)and(168,-64.9)..(167.7,-63)
				..controls(167.3,-60.5)and(167.1,-59.9)..(166.4,-58.4)
				..controls(166.1,-57.7)and(165.8,-56.9)..(165.7,-56.6)
				..controls(165.7,-56.2)and(165.4,-55.6)..(165.2,-55.2)
				..controls(164.6,-53.8)and(161.9,-50.7)..(161,-50.3)
				..controls(159.8,-49.7)and(157.9,-50)..(155.2,-51.2)
				..controls(154.7,-51.4)and(153.4,-51.8)..(152.2,-52.1)
				..controls(150.8,-52.4)and(149.4,-52.9)..(148.6,-53.2)
				..controls(147.8,-53.6)and(146.7,-54)..(146,-54.2)
				..controls(145.4,-54.4)and(144.5,-54.7)..(144,-54.9)
				..controls(143.6,-55.2)and(142.7,-55.5)..(141.8,-55.7)
				..controls(141,-55.9)and(139.6,-56.3)..(138.7,-56.6)
				..controls(137.8,-57)and(136.3,-57.4)..(135.3,-57.6)
				..controls(133.8,-57.9)and(133.2,-58.1)..(132.1,-58.7)
				..controls(131.1,-59.2)and(130.3,-59.4)..(129.4,-59.6)
				..controls(128.7,-59.8)and(127.8,-60)..(127.4,-60.1)
				..controls(127,-60.3)and(126,-60.5)..(125.1,-60.7)
				..controls(123.2,-61.1)and(121.7,-61.5)..(120.2,-62.1)
				..controls(119.7,-62.3)and(118.1,-62.8)..(116.7,-63.2)
				..controls(114.6,-63.9)and(114,-64.1)..(113,-64.7)
				..controls(111.9,-65.4)and(111.8,-65.4)..(111,-65.5)
				..controls(109.9,-65.5)and(109.3,-65.4)..(107.7,-64.7)
				..controls(107,-64.4)and(106,-64)..(105.3,-63.9)
				..controls(103.5,-63.5)and(102.2,-63)..(101.2,-62.5)
				..controls(100.1,-61.8)and(99.3,-61.5)..(97.8,-61.1)
				..controls(97.2,-61)and(96.2,-60.7)..(95.7,-60.4)
				..controls(95.1,-60.1)and(94.2,-59.8)..(93.5,-59.6)
				..controls(92.9,-59.4)and(91.7,-59)..(91,-58.7)..controls(90.2,-58.4)and(88.9,-58)..(88,-57.7)
				..controls(87.2,-57.5)and(86.1,-57.1)..(85.5,-56.8)
				..controls(85,-56.6)and(84.1,-56.2)..(83.5,-56)..controls(82.9,-55.7)and(81.8,-55.2)..(81,-54.9)
				..controls(80.1,-54.5)and(79.1,-54.2)..(78.7,-54.1)
				..controls(77.9,-53.8)and(77.3,-53.6)..(76.2,-52.9)
				..controls(75.7,-52.6)and(75.2,-52.4)..(74.1,-52.2)
				..controls(72.5,-51.9)and(72.4,-51.8)..(71.3,-51.1)
				..controls(69.9,-50.1)and(69.4,-49.9)..(67.9,-49.8)
				..controls(67,-49.7)and(66.6,-49.6)..(66.3,-49.4)
				..controls(66.1,-49.3)and(65.7,-49)..(65.4,-48.8)
				..controls(65.1,-48.7)and(64.2,-48.4)..(63.6,-48.2)
				..controls(62.9,-48)and(61.7,-47.6)..(60.9,-47.2)
				..controls(60.1,-46.8)and(58.9,-46.4)..(58.2,-46.2)
				..controls(57.6,-46)and(56.8,-45.7)..(56.5,-45.5)
				..controls(54.9,-44.6)and(54.6,-44.5)..(53.2,-44.2)
				..controls(51.7,-43.8)and(51,-43.6)..(49.9,-42.9)
				..controls(49.5,-42.6)and(48.9,-42.4)..(47.5,-42)
				..controls(45.7,-41.6)and(45.6,-41.5)..(44.3,-40.6)
				..controls(43.5,-40.1)and(42.8,-39.6)..(42.7,-39.4)
				..controls(42.4,-38.7)and(43,-38.1)..(44.5,-37.6)..controls(44.9,-37.5)and(46,-37)..(46.9,-36.6)
				..controls(47.9,-36.1)and(48.8,-35.8)..(49.6,-35.6)
				..controls(50.5,-35.4)and(51.3,-35.2)..(52.4,-34.6)
				..controls(53.3,-34.1)and(54.5,-33.7)..(54.9,-33.5)
				..controls(55.4,-33.4)and(56.4,-33)..(57.2,-32.6)
				..controls(58.6,-32)and(60.1,-31.5)..(61.6,-31.2)
				..controls(62.3,-31.1)and(62.7,-31)..(63.2,-30.7)
				..controls(64.1,-30)and(64.3,-29.9)..(65.3,-29.7)
				..controls(65.7,-29.6)and(66.5,-29.2)..(67.1,-28.9)
				..controls(67.6,-28.6)and(68.8,-28.2)..(69.6,-27.9)
				..controls(70.5,-27.7)and(71.6,-27.3)..(72.1,-27)
				..controls(72.6,-26.7)and(73.8,-26.2)..(74.8,-25.9)
				..controls(75.7,-25.5)and(76.7,-25.1)..(77.1,-25)
				..controls(77.4,-24.8)and(78.5,-24.4)..(79.6,-24.1)
				..controls(80.6,-23.8)and(82,-23.2)..(82.8,-22.9)
				..controls(83.5,-22.5)and(84.8,-22)..(85.8,-21.8)
				..controls(86.7,-21.5)and(87.8,-21.1)..(88.1,-20.9)
				..controls(88.4,-20.7)and(89.3,-20.3)..(90.1,-20.1)
				..controls(90.8,-19.8)and(91.9,-19.4)..(92.4,-19.1)
				..controls(93.5,-18.5)and(94.4,-18.2)..(95.4,-18.1)
				..controls(95.7,-18)and(96.4,-17.8)..(96.9,-17.5)
				..controls(97.5,-17.3)and(98.4,-16.9)..(99,-16.6)
				..controls(99.7,-16.3)and(100.9,-15.9)..(101.9,-15.6)
				..controls(102.9,-15.3)and(104.4,-14.7)..(105.3,-14.3)
				..controls(108.1,-13.2)and(108.5,-13.1)..(110.1,-13.1)
				..controls(111.4,-13.1)and(111.5,-13.1)..(112.4,-13.5)
				..controls(113.2,-14)and(115,-14.8)..(116.6,-15.4)
				..controls(116.9,-15.5)and(117.9,-15.8)..(118.8,-16.1)
				..controls(120.9,-16.6)and(121.5,-16.9)..(122.6,-17.9)
				..controls(123,-18.3)and(123.7,-18.8)..(123.9,-18.9)
				..controls(124.6,-19.4)and(126,-19.8)..(127.9,-20.2)
				..controls(129.6,-20.5)and(129.8,-20.6)..(131.1,-21.4)
				..controls(131.5,-21.7)and(132.4,-22.1)..(133.1,-22.3)
				..controls(136.1,-23.2)and(138.3,-24.1)..(140.2,-25)
				..controls(141,-25.5)and(141.5,-25.6)..(142.5,-25.8)
				..controls(145.1,-26.3)and(146,-26.6)..(149.4,-28.1)
				..controls(151.3,-28.9)and(153.5,-29.8)..(154.3,-30.1)
				..controls(155.1,-30.4)and(156.2,-30.8)..(156.6,-31)
				..controls(157.1,-31.2)and(158.2,-31.6)..(159.1,-31.8)
				..controls(160.6,-32.2)and(160.9,-32.3)..(161.7,-32.9)
				..controls(162.4,-33.3)and(162.9,-33.6)..(163.7,-33.8)
				..controls(164.3,-34)and(164.8,-34.1)..(164.9,-34.1)
				..controls(165,-34.1)and(165.6,-34.4)..(166.4,-34.8)
				..controls(167.3,-35.2)and(168.1,-35.4)..(169,-35.6)
				..controls(170,-35.8)and(170.5,-36)..(171.7,-36.6)
				..controls(172.4,-37)and(173.3,-37.4)..(173.6,-37.4)
				..controls(175.6,-37.7)and(176.2,-38)..(176.7,-39.1)
				..controls(176.9,-39.6)and(176.9,-39.7)..(176.7,-40.4)
				..controls(176.4,-41.6)and(175.8,-42.7)..(175.2,-43.4)
				..controls(174.5,-44.1)and(172.6,-45.5)..(171.7,-45.8)
				..controls(170.6,-46.2)and(169.1,-46)..(168.2,-45.4)
				..controls(167.9,-45.2)and(167.3,-44.8)..(166.9,-44.4)
				..controls(166.4,-44.1)and(165.8,-43.6)..(165.4,-43.2)
				..controls(165.1,-42.9)and(164.4,-42.3)..(163.9,-41.9)
				..controls(163.3,-41.5)and(162.7,-40.9)..(162.4,-40.5)
				..controls(161.4,-39.2)and(159.9,-38)..(158.8,-37.8)
				..controls(158.4,-37.8)and(158.3,-37.7)..(158.2,-37.2)
				..controls(158,-36.8)and(157.4,-36.1)..(157.1,-36.1)
				..controls(157,-36.1)and(156.7,-36.1)..(156.5,-36)
				..controls(156.1,-35.9)and(155.9,-35.7)..(155.7,-35.2)
				..controls(155.3,-34.6)and(155.1,-34.4)..(154,-34.1)
				..controls(153.6,-34)and(153.2,-33.7)..(152.9,-33.5)
				..controls(152.4,-32.9)and(151.6,-32.4)..(151.3,-32.4)
				..controls(150.9,-32.4)and(150.4,-32.6)..(150.4,-32.9)
				..controls(150.4,-33.2)and(150.8,-33.4)..(151.4,-33.6)
				..controls(152.1,-33.7)and(152.8,-34.1)..(153.5,-34.9)
				..controls(154.1,-35.5)and(154.3,-35.6)..(155,-35.8)
				..controls(155.7,-36)and(155.7,-36)..(156,-36.6)
				..controls(156.3,-37.3)and(156.7,-37.7)..(157.5,-37.8)
				..controls(158.1,-37.9)and(158.2,-38)..(158.4,-38.8)
				..controls(158.5,-39.4)and(158.6,-39.6)..(159.6,-40.4)
				..controls(160.2,-41)and(160.9,-41.6)..(161.3,-41.8)
				..controls(161.7,-42)and(161.9,-42.3)..(162.2,-42.8)
				..controls(162.4,-43.2)and(162.8,-43.8)..(163.1,-44.1)
				..controls(163.9,-44.8)and(165.1,-46.7)..(165.7,-48)
				..controls(166.2,-49.3)and(166.4,-49.7)..(167.2,-50.5)
				..controls(167.5,-50.8)and(167.9,-51.4)..(168.1,-51.8)
				..controls(168.3,-52.2)and(168.8,-53)..(169.1,-53.6)
				..controls(169.6,-54.3)and(169.8,-54.9)..(169.9,-55.6)
				..controls(170.4,-57.8)and(170.5,-58.1)..(171.2,-59.5)
				..controls(171.6,-60.3)and(172,-61.5)..(172.2,-62.1)
				..controls(172.3,-62.8)and(172.7,-63.7)..(172.9,-64.1)
				..controls(173.2,-64.6)and(173.4,-65.3)..(173.4,-65.6)
				..controls(173.8,-68.8)and(174.2,-70.6)..(174.9,-72.6)
				..controls(175.5,-74.5)and(175.9,-77.1)..(176.1,-80.2)
				..controls(176.3,-83.4)and(175.5,-84.4)..(173.7,-83.4)--cycle;
				\definecolor{ezs}{RGB}{247,179,26}\fill[ezs](-5.7,-156.6)--(-3.4,-157.5)--(61.9,12.7)--(59.7,13.6)--cycle;
				\definecolor{ezs}{RGB}{247,179,22}\fill[ezs](-3.8,-157.3)--(-1.5,-158.2)--(63.9,12)--(61.6,12.9)--cycle;
				\definecolor{ezs}{RGB}{248,179,19}\fill[ezs](-1.8,-158.1)--(0.4,-158.9)--(65.8,11.2)--(63.5,12.1)--cycle;
				\definecolor{ezs}{RGB}{248,178,16}\fill[ezs](0.1,-158.8)--(2.4,-159.7)--(67.7,10.5)--(65.4,11.4)--cycle;
				\definecolor{ezs}{RGB}{249,178,15}\fill[ezs](2,-159.5)--(4.3,-160.4)--(69.6,9.8)--(67.3,10.7)--cycle;
				\definecolor{ezs}{RGB}{251,177,17}\fill[ezs](3.9,-160.3)--(6.2,-161.1)--(71.5,9)--(69.2,9.9)--cycle;
				\definecolor{ezs}{RGB}{252,179,18}\fill[ezs](5.8,-161)--(8.1,-161.9)--(73.4,8.3)--(71.1,9.2)--cycle;
				\definecolor{ezs}{RGB}{252,182,17}\fill[ezs](7.7,-161.7)--(10,-162.6)--(75.3,7.6)--(73,8.5)--cycle;
				\definecolor{ezs}{RGB}{252,185,17}\fill[ezs](9.6,-162.5)--(11.9,-163.4)--(77.2,6.8)--(75,7.7)--cycle;
				\definecolor{ezs}{RGB}{251,189,18}\fill[ezs](11.5,-163.2)--(13.8,-164.1)--(79.2,6.1)--(76.9,7)--cycle;
				\definecolor{ezs}{RGB}{251,191,22}\fill[ezs](13.4,-163.9)--(15.7,-164.8)--(81.1,5.4)--(78.8,6.2)--cycle;
				\definecolor{ezs}{RGB}{252,193,25}\fill[ezs](15.4,-164.7)--(17.6,-165.6)--(83,4.6)--(80.7,5.5)--cycle;
				\definecolor{ezs}{RGB}{253,195,27}\fill[ezs](17.3,-165.4)--(19.6,-166.3)--(84.9,3.9)--(82.6,4.8)--cycle;
				\definecolor{ezs}{RGB}{253,196,31}\fill[ezs](19.2,-166.1)--(21.5,-167)--(86.8,3.2)--(84.5,4)--cycle;
				\definecolor{ezs}{RGB}{254,199,35}\fill[ezs](21.1,-166.9)--(23.4,-167.8)--(88.7,2.4)--(86.4,3.3)--cycle;
				\definecolor{ezs}{RGB}{253,201,38}\fill[ezs](23,-167.6)--(25.3,-168.5)--(90.6,1.7)--(88.3,2.6)--cycle;
				\definecolor{ezs}{RGB}{253,203,40}\fill[ezs](24.9,-168.3)--(27.2,-169.2)--(92.5,1)--(90.3,1.8)--cycle;
				\definecolor{ezs}{RGB}{253,205,42}\fill[ezs](26.8,-169.1)--(29.1,-170)--(94.4,0.2)--(92.2,1.1)--cycle;
				\definecolor{ezs}{RGB}{253,207,45}\fill[ezs](28.7,-169.8)--(31,-170.7)--(96.4,-0.5)--(94.1,0.4)--cycle;
				\definecolor{ezs}{RGB}{253,210,49}\fill[ezs](30.7,-170.5)--(32.9,-171.4)--(98.3,-1.2)--(96,-0.4)--cycle;
				\definecolor{ezs}{RGB}{253,212,54}\fill[ezs](32.6,-171.3)--(34.9,-172.2)--(100.2,-2)--(97.9,-1.1)--cycle;
				\definecolor{ezs}{RGB}{253,214,59}\fill[ezs](34.5,-172)--(36.8,-172.9)--(102.1,-2.7)--(99.8,-1.8)--cycle;
				\definecolor{ezs}{RGB}{253,216,63}\fill[ezs](36.4,-172.7)--(38.7,-173.6)--(104,-3.4)--(101.7,-2.6)--cycle;
				\definecolor{ezs}{RGB}{253,218,65}\fill[ezs](38.3,-173.5)--(40.6,-174.4)--(105.9,-4.2)--(103.6,-3.3)--cycle;
				\definecolor{ezs}{RGB}{254,220,66}\fill[ezs](40.2,-174.2)--(42.5,-175.1)--(107.8,-4.9)--(105.5,-4)--cycle;
				\definecolor{ezs}{RGB}{254,222,68}\fill[ezs](42.1,-175)--(44.4,-175.8)--(109.7,-5.6)--(107.5,-4.8)--cycle;
				\definecolor{ezs}{RGB}{254,224,71}\fill[ezs](44,-175.7)--(46.3,-176.6)--(111.7,-6.4)--(109.4,-5.5)--cycle;
				\definecolor{ezs}{RGB}{254,226,75}\fill[ezs](46,-176.4)--(48.2,-177.3)--(113.6,-7.1)--(111.3,-6.2)--cycle;
				\definecolor{ezs}{RGB}{254,228,79}\fill[ezs](47.9,-177.2)--(50.1,-178)--(115.5,-7.8)--(113.2,-7)--cycle;
				\definecolor{ezs}{RGB}{254,230,82}\fill[ezs](49.8,-177.9)--(52.1,-178.8)--(117.4,-8.6)--(115.1,-7.7)--cycle;
				\definecolor{ezs}{RGB}{254,232,85}\fill[ezs](51.7,-178.6)--(54,-179.5)--(119.3,-9.3)--(117,-8.4)--cycle;
				\definecolor{ezs}{RGB}{254,234,88}\fill[ezs](53.6,-179.4)--(55.9,-180.2)--(121.2,-10)--(118.9,-9.2)--cycle;
				\definecolor{ezs}{RGB}{254,237,91}\fill[ezs](55.5,-180.1)--(57.8,-181)--(123.1,-10.8)--(120.8,-9.9)--cycle;
				\definecolor{ezs}{RGB}{254,239,94}\fill[ezs](57.4,-180.8)--(59.7,-181.7)--(125,-11.5)--(122.8,-10.6)--cycle;
				\definecolor{ezs}{RGB}{254,241,97}\fill[ezs](59.3,-181.6)--(61.6,-182.4)--(127,-12.2)--(124.7,-11.4)--cycle;
				\definecolor{ezs}{RGB}{254,243,101}\fill[ezs](61.2,-182.3)--(63.5,-183.2)--(128.9,-13)--(126.6,-12.1)--cycle;
				\definecolor{ezs}{RGB}{254,245,103}\fill[ezs](63.2,-183)--(65.4,-183.9)--(130.8,-13.7)--(128.5,-12.8)--cycle;
				\definecolor{ezs}{RGB}{254,245,105}\fill[ezs](65.1,-183.8)--(67.4,-184.6)--(132.7,-14.4)--(130.4,-13.6)--cycle;
				\definecolor{ezs}{RGB}{254,245,105}\fill[ezs](67,-184.5)--(69.3,-185.4)--(134.6,-15.2)--(132.3,-14.3)--cycle;
				\definecolor{ezs}{RGB}{254,244,103}\fill[ezs](68.9,-185.2)--(71.2,-186.1)--(136.5,-15.9)--(134.2,-15)--cycle;
				\definecolor{ezs}{RGB}{253,243,101}\fill[ezs](70.8,-186)--(73.1,-186.8)--(138.4,-16.6)--(136.1,-15.8)--cycle;
				\definecolor{ezs}{RGB}{253,241,98}\fill[ezs](72.7,-186.7)--(75,-187.6)--(140.3,-17.4)--(138,-16.5)--cycle;
				\definecolor{ezs}{RGB}{253,239,95}\fill[ezs](74.6,-187.4)--(76.9,-188.3)--(142.2,-18.1)--(140,-17.2)--cycle;
				\definecolor{ezs}{RGB}{253,236,92}\fill[ezs](76.5,-188.2)--(78.8,-189)--(144.2,-18.8)--(141.9,-18)--cycle;
				\definecolor{ezs}{RGB}{253,234,88}\fill[ezs](78.5,-188.9)--(80.7,-189.8)--(146.1,-19.6)--(143.8,-18.7)--cycle;
				\definecolor{ezs}{RGB}{253,232,85}\fill[ezs](80.4,-189.6)--(82.7,-190.5)--(148,-20.3)--(145.7,-19.4)--cycle;
				\definecolor{ezs}{RGB}{253,230,82}\fill[ezs](82.3,-190.4)--(84.6,-191.2)--(149.9,-21)--(147.6,-20.2)--cycle;
				\definecolor{ezs}{RGB}{253,228,78}\fill[ezs](84.2,-191.1)--(86.5,-192)--(151.8,-21.8)--(149.5,-20.9)--cycle;
				\definecolor{ezs}{RGB}{253,226,74}\fill[ezs](86.1,-191.8)--(88.4,-192.7)--(153.7,-22.5)--(151.4,-21.6)--cycle;
				\definecolor{ezs}{RGB}{253,224,70}\fill[ezs](88,-192.6)--(90.3,-193.4)--(155.6,-23.3)--(153.3,-22.4)--cycle;
				\definecolor{ezs}{RGB}{253,222,67}\fill[ezs](89.9,-193.3)--(92.2,-194.2)--(157.5,-24)--(155.3,-23.1)--cycle;
				\definecolor{ezs}{RGB}{253,220,65}\fill[ezs](91.8,-194)--(94.1,-194.9)--(159.5,-24.7)--(157.2,-23.8)--cycle;
				\definecolor{ezs}{RGB}{253,217,63}\fill[ezs](93.7,-194.8)--(96,-195.6)--(161.4,-25.5)--(159.1,-24.6)--cycle;
				\definecolor{ezs}{RGB}{253,215,60}\fill[ezs](95.7,-195.5)--(97.9,-196.4)--(163.3,-26.2)--(161,-25.3)--cycle;
				\definecolor{ezs}{RGB}{253,213,56}\fill[ezs](97.6,-196.2)--(99.9,-197.1)--(165.2,-26.9)--(162.9,-26)--cycle;
				\definecolor{ezs}{RGB}{252,211,53}\fill[ezs](99.5,-197)--(101.8,-197.8)--(167.1,-27.7)--(164.8,-26.8)--cycle;
				\definecolor{ezs}{RGB}{252,209,49}\fill[ezs](101.4,-197.7)--(103.7,-198.6)--(169,-28.4)--(166.7,-27.5)--cycle;
				\definecolor{ezs}{RGB}{252,207,45}\fill[ezs](103.3,-198.4)--(105.6,-199.3)--(170.9,-29.1)--(168.6,-28.2)--cycle;
				\definecolor{ezs}{RGB}{253,205,42}\fill[ezs](105.2,-199.2)--(107.5,-200)--(172.8,-29.9)--(170.6,-29)--cycle;
				\definecolor{ezs}{RGB}{253,202,40}\fill[ezs](107.1,-199.9)--(109.4,-200.8)--(174.7,-30.6)--(172.5,-29.7)--cycle;
				\definecolor{ezs}{RGB}{253,200,38}\fill[ezs](109,-200.6)--(111.3,-201.5)--(176.7,-31.3)--(174.4,-30.4)--cycle;
				\definecolor{ezs}{RGB}{253,198,35}\fill[ezs](111,-201.4)--(113.2,-202.2)--(178.6,-32.1)--(176.3,-31.2)--cycle;
				\definecolor{ezs}{RGB}{253,196,31}\fill[ezs](112.9,-202.1)--(115.2,-203)--(180.5,-32.8)--(178.2,-31.9)--cycle;
				\definecolor{ezs}{RGB}{253,194,28}\fill[ezs](114.8,-202.8)--(117.1,-203.7)--(182.4,-33.5)--(180.1,-32.6)--cycle;
				\definecolor{ezs}{RGB}{253,192,26}\fill[ezs](116.7,-203.6)--(119,-204.5)--(184.3,-34.3)--(182,-33.4)--cycle;
				\definecolor{ezs}{RGB}{253,189,24}\fill[ezs](118.6,-204.3)--(120.9,-205.2)--(186.2,-35)--(183.9,-34.1)--cycle;
				\definecolor{ezs}{RGB}{253,186,20}\fill[ezs](120.5,-205)--(122.8,-205.9)--(188.1,-35.7)--(185.8,-34.8)--cycle;
				\definecolor{ezs}{RGB}{253,184,17}\fill[ezs](122.4,-205.8)--(124.7,-206.7)--(190,-36.5)--(187.8,-35.6)--cycle;
				\definecolor{ezs}{RGB}{252,183,15}\fill[ezs](124.3,-206.5)--(126.6,-207.4)--(192,-37.2)--(189.7,-36.3)--cycle;
				\definecolor{ezs}{RGB}{252,181,14}\fill[ezs](126.3,-207.2)--(128.5,-208.1)--(193.9,-37.9)--(191.6,-37.1)--cycle;
				\definecolor{ezs}{RGB}{252,179,13}\fill[ezs](128.2,-208)--(130.4,-208.9)--(195.8,-38.7)--(193.5,-37.8)--cycle;
				\definecolor{ezs}{RGB}{252,178,12}\fill[ezs](130.1,-208.7)--(132.4,-209.6)--(197.7,-39.4)--(195.4,-38.5)--cycle;
				\definecolor{ezs}{RGB}{251,178,10}\fill[ezs](132,-209.4)--(134.3,-210.3)--(199.6,-40.1)--(197.3,-39.3)--cycle;
				\definecolor{ezs}{RGB}{251,178,10}\fill[ezs](133.9,-210.2)--(136.2,-211.1)--(201.5,-40.9)--(199.2,-40)--cycle;
				\definecolor{ezs}{RGB}{250,177,11}\fill[ezs](135.8,-210.9)--(138.1,-211.8)--(203.4,-41.6)--(201.1,-40.7)--cycle;
				\definecolor{ezs}{RGB}{250,178,11}\fill[ezs](137.7,-211.6)--(140,-212.5)--(205.3,-42.3)--(203.1,-41.5)--cycle;
				\definecolor{ezs}{RGB}{248,180,10}\fill[ezs](139.6,-212.4)--(141.9,-213.3)--(207.3,-43.1)--(205,-42.2)--cycle;
				\definecolor{ezs}{RGB}{247,181,9}\fill[ezs](141.5,-213.1)--(143.8,-214)--(209.2,-43.8)--(206.9,-42.9)--cycle;
				\definecolor{ezs}{RGB}{246,182,9}\fill[ezs](143.5,-213.8)--(145.7,-214.7)--(211.1,-44.5)--(208.8,-43.7)--cycle;
				\definecolor{ezs}{RGB}{246,182,10}\fill[ezs](145.4,-214.6)--(147.7,-215.5)--(213,-45.3)--(210.7,-44.4)--cycle;
			\end{scope}
			% chữ EDUZY (Times đậm, làm dày nét bằng 5 lớp lệch nhẹ; căn đúng khung chữ logo gốc)
			\foreach \dx/\dy in {-0.7/0,0.7/0,0/-0.3,0/0.3,0/0}
			\node[anchor=center,inner sep=0pt,outer sep=0pt,text=ezText,align=center] at (109.6+\dx,-193.6+\dy)
			{\resizebox{1.01cm}{0.22cm}{\fontencoding{T1}\fontfamily{ptm}\fontseries{b}\selectfont EDUZY}};
		\end{tikzpicture}%
	}%
}
%====================== THIẾT LẬP THÂN ĐỀ ======================
%--- Môi trường tự luận {bt}: ex_test 3.3.4 chưa khai báo sẵn nên tạo ở đây
%    (ex_test đã tự móc \fixshowans và \listenumerate cho bt)
\makeatletter
\@ifundefined{bt}{\newtheorem{bt}{\color{gold1}\namebt}}{}
\makeatother
\renewcommand{\sectionmark}[1]{\markright{\thesection\ #1}}
%--- Khung tiêu đề (logo + tối đa 4 dòng; dòng 3, 4 để trống {} thì bỏ qua)
%  Đề thi:    \khungtieude{THI THỬ THPT QUỐC GIA NĂM 2026}{Môn Toán -- Đề số 21}
%                         {(Đề thi gồm: ...)}{Thời gian làm bài: 90 phút (không kể thời gian phát đề).}
%  Chuyên đề: \khungtieude{CHUYÊN ĐỀ 3 -- DÃY SỐ VÀ CÔNG THỨC TRUY HỒI}
%                         {Bồi dưỡng học sinh giỏi Toán THPT}{...}{}
\newcommand{\khungtieude}[4]{%
	\begin{tcolorbox}[colframe=gold1,colback=gold5!45,boxrule=0.5mm]
		\begin{minipage}[c]{3cm}
			\centerline{\eduzylogo[2.4cm]}
		\end{minipage}\hspace{0.5cm}
		\begin{minipage}[c]{12.5cm}
			\vskip 0.15 cm
			\centerline{\ind{#1}}
			\vskip 0.15 cm
			\centerline{\textbf{#2}}
			\if\relax\detokenize{#3}\relax\else
				\vskip 0.15 cm
				\centerline{\textit{#3}}
			\fi
			\if\relax\detokenize{#4}\relax\else
				\vskip 0.15 cm
				\centerline{\textit{#4}}
			\fi
			\vskip 0.15 cm
		\end{minipage}
	\end{tcolorbox}
}
%====================== BỔ SUNG CHO BÀI HỌC, CHUYÊN ĐỀ ======================
\usepackage[shortlabels]{enumitem}
\usepackage{tabularx}
%--- Danh sách gọn: \begin{enumerate}[gon] ... \end{enumerate}
\SetEnumitemKey{gon}{leftmargin=*,itemsep=1pt,parsep=0pt,topsep=3pt}
%--- Khung lý thuyết: \begin{khung}{} ... \end{khung}; tiêu đề chỉ dùng khi thật ngắn ("Mẹo", "Lỗi thường gặp")
\newtcolorbox{khung}[1]{breakable,enhanced,colframe=gold1,colback=gold5!45,
	boxrule=0.4mm,fonttitle=\bfseries\sffamily,coltitle=white,colbacktitle=gold1,
	attach boxed title to top left={xshift=6mm,yshift=-3mm},
	boxed title style={arc=1mm,boxrule=0.3mm},title={#1}}
%--- Tiêu đề mục nhỏ trong một dạng: \muc{Dấu hiệu}, \muc{Phương pháp}, \muc{Bài tự luyện}
\makeatletter
\newcommand{\muc}[1]{\par\medskip\noindent\indam{#1}\par\nobreak\smallskip\@afterindentfalse\@afterheading}
\makeatother
%--- Ô công thức cần nhớ (căn giữa, viền vàng, vừa khít nội dung)
%     \begin{center}\congthuc{$Q(n)=(n-1)!$}\end{center}
\newtcbox{\congthuc}{on line,colframe=gold1,colback=gold5!45,boxrule=0.5pt,arc=2pt,
	left=8pt,right=8pt,top=4pt,bottom=4pt}
%--- Công tắc lời giải bài tự luyện (ví dụ mẫu luôn dùng \loigiai)
%     \loigiaiTL[{đáp số}]{lời giải đầy đủ}
%     \TLfull -> in lời giải đầy đủ;  \TLdapso -> chỉ in một dòng "Đáp số";  \TLan -> ẩn hoàn toàn
%     (\CheDo là dethi thì mọi lời giải và đáp số đều ẩn)
\makeatletter
\newif\ifTLhien   \TLhientrue
\newif\ifTLdapso  \TLdapsofalse
\newcommand{\TLfull}{\TLhientrue\TLdapsofalse}
\newcommand{\TLdapso}{\TLhientrue\TLdapsotrue}
\newcommand{\TLan}{\TLhienfalse}
\newcommand{\loigiaiTL}[2][]{%
	\ifTLhien
		\ifTLdapso
			\def\@tlarg{#1}%
			\ifx\@tlarg\@empty\else
				\ifdefstring{\kindOPTN}{dethi}{}{\par\noindent{\bfseries\color{gold1}Đáp số:} #1\par}%
			\fi
		\else
			\loigiai{#2}%
		\fi
	\fi}
\makeatother
\TLdapso% mặc định: chỉ in đáp số bài tự luyện
%====================== BẢNG ĐÚNG/SAI (\choiceTFt) ======================
%  Hàng tiêu đề nền nâu đồng chữ trắng (Phát biểu | Đ | S), các hàng tô xen kẽ kem / trắng,
%  đường kẻ vàng kim, nhãn a) b) c) d) màu nâu đồng. Bản có đáp án: cột được chọn có vòng tròn
%  nâu đồng với dấu tích trắng (câu đúng tích ở cột Đúng, câu sai tích ở cột Sai).
%  Tiêu đề ghi đủ "Đúng | Sai": \OPTN{viettat=0} sau \input (cột Đúng/Sai rộng ra, phát biểu hẹp lại).
%  Dùng bảng gốc của ex_test: \def\BangDSgoc{1} trước \input{preamble-eduzy-v2}.
\providecommand{\BangDSgoc}{0}
\ifnum\BangDSgoc=0
\usepackage{colortbl,pifont}
\colorlet{ezDSke}{gold2!60!gold1}% đường kẻ bảng (nét 0,4pt như ex_test để không đổi chiều cao bảng)
\colorlet{ezDSle}{gold5!60}% hàng a), c)
\colorlet{ezDSchan}{white}% hàng b), d)
%--- dấu chọn ở cột Đúng/Sai
%    (bề rộng 0, căn giữa ô: không làm đổi bề rộng cột của ex_test, không báo tràn lề)
\newcommand{\EZdauchon}{\makebox[0pt]{\tikz[baseline=-1ex]\node[circle,fill=gold1,inner sep=0pt,
	minimum size=1.2em,text=white,font=\small]{\ding{51}};}}
\def\tickT{\EZdauchon}\def\tickF{\EZdauchon}
\AtBeginDocument{\sbox{\ccTFTab}{\EZdauchon}}
%--- nhãn a) b) c) d) màu nâu đồng (ở bản lời giải, nhãn câu đúng vẫn đóng khung như ex_test)
\renewcommand{\FalseTF}{\stepcounter{dapan}\textbf{\color{gold1}\DapAnTF\sepTF}\ignorespaces}
%--- trước khi dựng bảng: kẻ màu vàng kim (bề rộng cột giữ nguyên như ex_test); sau bảng trả lại màu đen
\newcount\EZhangTF
\newcommand{\EZbatdauTF}{%
	\global\EZhangTF=0\relax
	\arrayrulecolor{ezDSke}}
\newcommand{\EZketthucTF}{\arrayrulecolor{black}}
%--- hàng tiêu đề nền nâu đồng, chữ trắng
\makeatletter
\newcommand{\EZo@TF}[1]{\centering\bfseries\color{white}#1}
\newcommand{\EZtieudeTF}{%
	\ifcase\cotDS
		\ifnum\cotSTT=1\relax
			\def\rowHeader{\hline\rowcolor{gold1}\EZo@TF{TT}&\EZo@TF{\phatbieu}\tabularnewline}%
		\else
			\def\rowHeader{\hline\rowcolor{gold1}\EZo@TF{\phatbieu}\tabularnewline}%
		\fi
	\or
		\ifnum\cotSTT=1\relax
			\def\rowHeader{\hline\rowcolor{gold1}\EZo@TF{TT}&\EZo@TF{\phatbieu}&\EZo@TF{\Dung}\tabularnewline}%
		\else
			\def\rowHeader{\hline\rowcolor{gold1}\EZo@TF{\phatbieu}&\EZo@TF{\Dung}\tabularnewline}%
		\fi
	\or
		\ifnum\cotSTT=1\relax
			\def\rowHeader{\hline\rowcolor{gold1}\EZo@TF{TT}&\EZo@TF{\phatbieu}&\EZo@TF{\Dung}&\EZo@TF{\Sai}\tabularnewline}%
		\else
			\def\rowHeader{\hline\rowcolor{gold1}\EZo@TF{\phatbieu}&\EZo@TF{\Dung}&\EZo@TF{\Sai}\tabularnewline}%
		\fi
	\fi}
%--- tô xen kẽ từng hàng (gọi ngay sau \hline đầu mỗi hàng)
\newcommand{\EZmauhangTF}{%
	\noalign{\global\advance\EZhangTF\@ne}%
	\ifodd\EZhangTF\rowcolor{ezDSle}\else\rowcolor{ezDSchan}\fi}
%--- gắn vào \choiceTFt / \writeChoiceTFt của ex_test
\xpretocmd{\choiceTFt}{\EZbatdauTF}{}{\typeout{[preamble-eduzy-v2] Không gắn được bảng Đúng/Sai (đầu bảng)}}
\xpatchcmd{\choiceTFt}{\ifnum\dongPB=0\def\rowHeader{}\fi}{\EZtieudeTF\ifnum\dongPB=0\def\rowHeader{}\fi}
	{}{\typeout{[preamble-eduzy-v2] Không gắn được bảng Đúng/Sai (hàng tiêu đề)}}
\xpatchcmd{\writeChoiceTFt}{\hline\xdef\dotEX{\sdotEX}}{\hline\EZmauhangTF\xdef\dotEX{\sdotEX}}
	{}{\typeout{[preamble-eduzy-v2] Không gắn được bảng Đúng/Sai (tô hàng)}}
\xpatchcmd{\writeChoiceTFt}{{\end{longtable}}}{{\end{longtable}\EZketthucTF}}
	{}{\typeout{[preamble-eduzy-v2] Không gắn được bảng Đúng/Sai (cuối longtable)}}
\xpatchcmd{\writeChoiceTFt}{{\end{xtabular}\end{center}\medskip\noindent}}{{\end{xtabular}\end{center}\EZketthucTF\medskip\noindent}}
	{}{\typeout{[preamble-eduzy-v2] Không gắn được bảng Đúng/Sai (cuối xtabular)}}
\fi
%=====================================================================
\renewcommand{\parallel}{\mathrel{/\!\!/}}

%  trong file đề (tools/dan_trang/dan_trang.py tự ghi khi đặt \NeedspaceCau):
%    \def\TachCauDS{1}   câu Đúng/Sai được tách: câu dẫn (kèm hình) ở cuối trang, bảng Phát biểu sang trang sau
%                        (câu dẫn và hình không tách nhau, bảng không tách giữa hai trang)
%    \def\TachCauLG{1}   bản lời giải (\CheDo khác dethi): phần đề của câu cao hơn 1/3 trang được phép tách
%  Lệnh dùng trong file đề (giá trị do tools/dan_trang/dan_trang.py đo ở bản dethi):
%    \NeedspaceCau[<cao câu dẫn + hình>]{<cao cả khối câu>}   ngay trên \begin{ex} (câu đầu phần: trên tiêu đề phần)
%        công tắc tắt: đúng bằng \Needspace{<cao cả khối câu>} như cũ;
%        \TachCauDS=1 và có [..] (câu Đúng/Sai được chọn tách): chỉ giữ câu dẫn + hình;
%        \TachCauLG=1 ở bản lời giải: đòi tối đa 1/3 \textheight.
%    \NeedspaceBang{<cao bảng Phát biểu>}%   ngay trước \choiceTFt: \TachCauDS=1 thì giữ trọn bảng trên một trang;
%        công tắc tắt: không làm gì.
\providecommand{\TachCauDS}{0}
\providecommand{\TachCauLG}{0}
\newlength{\EZnsCau}
\newcommand{\NeedspaceCau}[2][]{%
	\setlength{\EZnsCau}{#2}%
	\ifnum\TachCauDS=1 \if\relax\detokenize{#1}\relax\else\setlength{\EZnsCau}{#1}\fi\fi
	\ifnum\TachCauLG=1 \ifdefstring{\kindOPTN}{dethi}{}{\ifdim\EZnsCau>0.3333\textheight\setlength{\EZnsCau}{0.3333\textheight}\fi}\fi
	\Needspace{\EZnsCau}}
\newcommand{\NeedspaceBang}[1]{\ifnum\TachCauDS=1 \Needspace{#1}\fi}

%--- Thang màu: đậm -> nhạt
\definecolor{gold1}{RGB}{126,96,0}    % nâu đồng đậm  #7E6000
\definecolor{gold2}{RGB}{192,144,3}   % vàng kim      #C09003
\definecolor{gold3}{RGB}{255,215,99}  % vàng nhạt     #FFD763
\definecolor{gold4}{RGB}{255,231,153} % vàng rất nhạt #FFE799
\definecolor{gold5}{RGB}{255,243,205} % vàng kem      #FFF3CD
%--- Ghi đè tông xanh/đỏ cũ
\definecolor{myblue}{RGB}{126,96,0}
\colorlet{blue}{gold1}
\colorlet{cyan}{gold2}
\colorlet{teal}{gold2}

%--- Tên chương mặc định (file đề nên tự khai báo lại)
\providecommand{\tenchuong}{}
%--- Chân trang
\providecommand{\thuonghieu}{TOÁN 12 -- Eduzy.com.vn}
\providecommand{\tacgia}{Tô Gia Bảo -- 0353703783}

%--- Chữ cạnh logo ở đầu trang
\providecommand{\khauhieu}{Eduzy.com.vn}

%====================== THÔNG SỐ ĐẦU TRANG, CHÂN TRANG (v2) ======================
%  Dải nền đậm tràn ngang trang; ô chữ vàng kem có cạnh xiên, kèm 2 vạch xiên chuyển màu.
%  Đầu trang lẻ : logo + \khauhieu (trái)   | ô chữ \tenchuong (phải)
%  Chân trang lẻ: ô chữ \thuonghieu (trái)  | Trang (số) (phải)
%  Trang chẵn soi gương trang lẻ; chân trang chẵn ghi \tacgia.
%--- Màu: chỉ dùng thang vàng kim Eduzy
\colorlet{ezNen}{gold1}% dải nền đậm
\colorlet{ezVach1}{gold3}% vạch xiên sát dải nền (đậm hơn)
\colorlet{ezVach2}{gold4}% vạch xiên sát ô chữ (nhạt hơn)
\colorlet{ezO}{gold5}% ô chữ vàng kem
%--- Kích thước
\newlength{\EZdai}\setlength{\EZdai}{1.1cm}% chiều cao dải nền
\newlength{\EZo}\setlength{\EZo}{0.68cm}% chiều cao ô chữ
\newlength{\EZxien}\setlength{\EZxien}{0.55cm}% độ lệch ngang của cạnh xiên
\newlength{\EZvach}\setlength{\EZvach}{0.75cm}% bề rộng mỗi vạch xiên
\newlength{\EZkhe}\setlength{\EZkhe}{0.12cm}% khe trắng giữa vạch và ô chữ
\newlength{\EZdem}\setlength{\EZdem}{0.4cm}% khoảng đệm từ chữ tới cạnh xiên
\newlength{\EZle}\setlength{\EZle}{1.8cm}% lề trái/phải của chữ (bằng lề thân bài)
\newlength{\EZrongmin}\setlength{\EZrongmin}{4.5cm}% ô chữ rộng tối thiểu
\newlength{\EZrong}% bề rộng chữ trong ô (tính tự động)
%--- Font
\newcommand{\EZfontO}{\fontsize{10.5pt}{12pt}\selectfont\bfseries}% chữ trong ô
\newcommand{\EZfontKH}{\fontfamily{qag}\fontsize{11pt}{12pt}\selectfont\bfseries}% \khauhieu cạnh logo
\newcommand{\EZfontTrang}{\fontsize{11.5pt}{12pt}\selectfont\bfseries}% chữ "Trang" và số trang
%--- Đo bề rộng chữ trong ô
\newcommand{\EZdorong}[1]{%
	\settowidth{\EZrong}{\EZfontO #1}%
	\ifdim\EZrong<\EZrongmin\setlength{\EZrong}{\EZrongmin}\fi}
%--- Logo + chữ trắng \khauhieu (dựng sẵn một lần khi vào \begin{document})
\newsavebox{\EZlogobox}
\AtBeginDocument{\sbox{\EZlogobox}{%
		\raisebox{-0.5\height}{\eduzylogo[1cm]}\hspace{0.15cm}%
		\raisebox{-0.5\height}{\color{white}\EZfontKH\khauhieu}}}
%--- Hai vạch xiên + khe trắng. #1: điểm gốc ở mép dưới, #2: hướng (1 sang phải, -1 sang trái),
%    #3: hướng nghiêng (1 kiểu "/", -1 kiểu "\")
\newcommand{\EZvachxien}[3]{%
	\fill[white] (#1) --++({(#3)*\EZxien},\EZo) --++({(#2)*\EZkhe},0) --++({-(#3)*\EZxien},-\EZo) --cycle;
	\fill[ezVach2] ([xshift=(#2)*\EZkhe]#1) --++({(#3)*\EZxien},\EZo) --++({(#2)*\EZvach},0) --++({-(#3)*\EZxien},-\EZo) --cycle;
	\fill[ezVach1] ([xshift=(#2)*(\EZkhe+\EZvach)]#1) --++({(#3)*\EZxien},\EZo) --++({(#2)*\EZvach},0) --++({-(#3)*\EZxien},-\EZo) --cycle;
	\draw[white,line width=0.6pt] ([xshift=(#2)*(\EZkhe+\EZvach)]#1) --++({(#3)*\EZxien},\EZo);
}
%--- Chữ trong ô (căn giữa theo chiều cao ô, không lệch vì dấu tiếng Việt)
\tikzset{ezchu/.style={inner sep=0pt,text=ezNen,text height=0.72em,text depth=0.22em}}

%====================== LOGO CHÌM GIỮA TRANG ======================
%  Vẽ cùng đầu trang nên nằm dưới chữ của thân bài, có trên mọi trang.
%  Tắt: \logochimtat (từ trang đó trở đi), bật lại: \logochimbat.
%  Đổi cỡ, độ đậm: khai báo trước \input, ví dụ \def\EZchimcao{9cm} \def\EZchimdo{0.1}
\providecommand{\EZchimcao}{10cm}% chiều cao logo chìm
\providecommand{\EZchimdo}{0.12}% độ đậm (0 = trong suốt, 1 = rõ hẳn)
\newif\ifEZchim \EZchimtrue
\newcommand{\logochimtat}{\EZchimfalse}
\newcommand{\logochimbat}{\EZchimtrue}
%  Dựng sẵn logo mờ một lần khi vào \begin{document}; transparency group làm mờ cả logo
%  một lượt nên không lộ vệt ở chỗ các dải màu chồng nhau.
\newsavebox{\EZchimtho}
\newsavebox{\EZchimbox}
\AtBeginDocument{%
	\sbox{\EZchimtho}{\eduzylogo[\EZchimcao]}%
	\sbox{\EZchimbox}{\begin{tikzpicture}
			\begin{scope}[transparency group,opacity=\EZchimdo]
				\node[inner sep=0pt] {\usebox{\EZchimtho}};
			\end{scope}
	\end{tikzpicture}}}
\newcommand{\EZvechim}{%
	\ifEZchim
		\node[inner sep=0pt] at (current page.center) {\usebox{\EZchimbox}};
	\fi}

%====================== ĐẦU TRANG - TRANG LẺ ======================
%  (logo bên trái, ô chữ \tenchuong bên phải)
\newcommand{\myfancyheadOdd}{%
	\EZdorong{\tenchuong}%
	\begin{tikzpicture}[remember picture,overlay]
		\coordinate (T) at (current page.north west);
		\coordinate (P) at (current page.north east);
		%--- logo chìm giữa trang
		\EZvechim
		%--- dải nền
		\fill[ezNen] (T) rectangle ([yshift=-\EZdai]P);
		%--- ô chữ: X là góc dưới bên trái, cạnh trái xiên kiểu "/"
		\coordinate (X) at ([xshift=-\EZle-\EZrong-\EZdem-\EZxien,yshift=-\EZdai]P);
		\fill[ezO] (X) --++(\EZxien,\EZo) -- ([yshift=-\EZdai+\EZo]P) -- ([yshift=-\EZdai]P) --cycle;
		\EZvachxien{X}{-1}{1}
		\node[ezchu,anchor=east] at ([xshift=-\EZle,yshift=-\EZdai+0.5*\EZo]P) {\EZfontO\tenchuong};
		%--- logo + tên
		\node[inner sep=0pt,anchor=west] at ([xshift=\EZle,yshift=-0.5*\EZdai]T) {\usebox{\EZlogobox}};
	\end{tikzpicture}%
}

%====================== ĐẦU TRANG - TRANG CHẴN =====================
%  (soi gương: ô chữ \tenchuong bên trái, logo bên phải)
\newcommand{\myfancyheadEven}{%
	\EZdorong{\tenchuong}%
	\begin{tikzpicture}[remember picture,overlay]
		\coordinate (T) at (current page.north west);
		\coordinate (P) at (current page.north east);
		%--- logo chìm giữa trang
		\EZvechim
		%--- dải nền
		\fill[ezNen] (T) rectangle ([yshift=-\EZdai]P);
		%--- ô chữ: X là góc dưới bên phải, cạnh phải xiên kiểu "\"
		\coordinate (X) at ([xshift=\EZle+\EZrong+\EZdem+\EZxien,yshift=-\EZdai]T);
		\fill[ezO] (X) --++(-\EZxien,\EZo) -- ([yshift=-\EZdai+\EZo]T) -- ([yshift=-\EZdai]T) --cycle;
		\EZvachxien{X}{1}{-1}
		\node[ezchu,anchor=west] at ([xshift=\EZle,yshift=-\EZdai+0.5*\EZo]T) {\EZfontO\tenchuong};
		%--- logo + tên
		\node[inner sep=0pt,anchor=east] at ([xshift=-\EZle,yshift=-0.5*\EZdai]P) {\usebox{\EZlogobox}};
	\end{tikzpicture}%
}

%====================== CHÂN TRANG - TRANG LẺ ======================
%  (ô chữ \thuonghieu bên trái, "Trang (số)" bên phải)
\newcommand{\myfancyfootOdd}{%
	\EZdorong{\thuonghieu}%
	\begin{tikzpicture}[remember picture,overlay]
		\coordinate (T) at (current page.south west);
		\coordinate (P) at (current page.south east);
		%--- dải nền
		\fill[ezNen] (T) rectangle ([yshift=\EZdai]P);
		%--- ô chữ sát mép trên dải: X là góc dưới bên phải, cạnh phải xiên kiểu "/"
		\coordinate (X) at ([xshift=\EZle+\EZrong+\EZdem,yshift=\EZdai-\EZo]T);
		\fill[ezO] ([yshift=\EZdai-\EZo]T) -- (X) --++(\EZxien,\EZo) -- ([yshift=\EZdai]T) --cycle;
		\EZvachxien{X}{1}{1}
		\node[ezchu,anchor=west] at ([xshift=\EZle,yshift=\EZdai-0.5*\EZo]T) {\EZfontO\thuonghieu};
		%--- số trang
		\node[circle,draw=white,line width=0.9pt,text=white,inner sep=0pt,minimum size=0.62cm,anchor=east]
		(so) at ([xshift=-\EZle,yshift=0.5*\EZdai]P) {\EZfontTrang\thepage};
		\node[inner sep=0pt,text=white,anchor=base east] (tr) at ([xshift=-0.15cm]so.base west) {\EZfontTrang Trang};
	\end{tikzpicture}%
}

%====================== CHÂN TRANG - TRANG CHẴN ====================
%  (soi gương: "Trang (số)" bên trái, ô chữ \tacgia bên phải)
\newcommand{\myfancyfootEven}{%
	\EZdorong{\tacgia}%
	\begin{tikzpicture}[remember picture,overlay]
		\coordinate (T) at (current page.south west);
		\coordinate (P) at (current page.south east);
		%--- dải nền
		\fill[ezNen] (T) rectangle ([yshift=\EZdai]P);
		%--- ô chữ sát mép trên dải: X là góc dưới bên trái, cạnh trái xiên kiểu "\"
		\coordinate (X) at ([xshift=-\EZle-\EZrong-\EZdem,yshift=\EZdai-\EZo]P);
		\fill[ezO] ([yshift=\EZdai-\EZo]P) -- (X) --++(-\EZxien,\EZo) -- ([yshift=\EZdai]P) --cycle;
		\EZvachxien{X}{-1}{-1}
		\node[ezchu,anchor=east] at ([xshift=-\EZle,yshift=\EZdai-0.5*\EZo]P) {\EZfontO\tacgia};
		%--- số trang
		%    (vòng tròn ẩn so0 nằm giữa dải để lấy đường chân chữ, giống trang lẻ)
		\node[circle,inner sep=0pt,minimum size=0.62cm] (so0) at ([yshift=0.5*\EZdai]T) {\EZfontTrang\phantom{\thepage}};
		\node[inner sep=0pt,text=white,anchor=base west] (tr) at ($(so0.base -| T)+(\EZle,0)$) {\EZfontTrang Trang};
		\node[circle,draw=white,line width=0.9pt,text=white,inner sep=0pt,minimum size=0.62cm,anchor=base west]
		(so) at ([xshift=0.15cm]tr.base east) {\EZfontTrang\thepage};
	\end{tikzpicture}%
}

%--- Giữ tương thích với các file cũ đang gọi \myfancyhead, \myfancyfoot
\newcommand{\myfancyhead}{\myfancyheadOdd}
\newcommand{\myfancyfoot}{\myfancyfootOdd}

%======================= GẮN VÀO TRANG =============================
\pagestyle{fancy}
\fancyhf{}
\fancyhead[LO]{\myfancyheadOdd}
\fancyhead[RE]{\myfancyheadEven}
\fancyfoot[LO]{\myfancyfootOdd}
\fancyfoot[RE]{\myfancyfootEven}
\renewcommand{\headrulewidth}{0pt}
\renewcommand{\footrulewidth}{0pt}
\fancypagestyle{plain}{%
	\fancyhf{}
	\fancyhead[LO]{\myfancyheadOdd}
	\fancyhead[RE]{\myfancyheadEven}
	\fancyfoot[LO]{\myfancyfootOdd}
	\fancyfoot[RE]{\myfancyfootEven}
	\renewcommand{\headrulewidth}{0pt}
	\renewcommand{\footrulewidth}{0pt}
}
%=================== HẾT PHẦN GỘP TỪ eduzy-de.tex ===================
\newcommand{\indam}[1]{
	{\fontsize{12pt}{9pt}\sffamily \selectfont\bfseries\color{gold1} #1}
}
\newcommand{\ind}[1]{
	{\fontsize{14pt}{9pt}\fontfamily{qag} \selectfont\bfseries\color{gold1} #1}
}
\renewcommand{\FalseEX}{\stepcounter{dapan}\textbf{\selectfont\color{gold1}\Alph{dapan}.}\ignorespaces}
\renewtheorem{ex}{\color{gold1}Câu}
\renewcommand{\baselinestretch}{1.1}
\newcommand{\hoac}[1]{
	\left[\begin{aligned}#1\end{aligned}\right.}
\newcommand{\heva}[1]{
	\left\{\begin{aligned}#1\end{aligned}\right.}
\usepackage{esvect}
\def\vec{\vv}
\def\vt{\vv}
\def\overrightarrow{\vv}
%====================== LOGO EDUZY VẼ BẰNG TIKZ ======================
%  Đường viền dò từ logo gốc; dải màu vàng kim đo trực tiếp trên ảnh gốc
%  (nghiêng 21 độ, tô bằng 80 dải hẹp nên hiển thị giống nhau trên mọi trình xem PDF).
%  Dùng: \eduzylogo          -> cao 2.4cm (mặc định)
%        \eduzylogo[1.5cm]   -> cao 1.5cm
\definecolor{ezText}{RGB}{245,200,46}
\newcommand{\eduzylogo}[1][2.4cm]{%
	\resizebox{!}{#1}{%
		\begin{tikzpicture}[x=0.01cm,y=0.01cm]
			\path[use as bounding box] (0,0) rectangle (224,-224);
			% mũ tốt nghiệp + vương miện
			\begin{scope}
				\clip
				(119.4,-176.1)..controls(118.9,-176.1)and(116.8,-176)..(114.6,-176)
				..controls(101.4,-176)and(85.8,-175.9)..(83.7,-175.7)
				..controls(81.8,-175.6)and(81,-175.6)..(80.2,-175.7)
				..controls(78.8,-175.9)and(77.7,-175.9)..(76.2,-175.7)
				..controls(74.9,-175.5)and(70.8,-175.5)..(64.3,-175.7)
				..controls(59.1,-175.9)and(58.3,-175.7)..(58.6,-173.7)
				..controls(58.6,-172.8)and(59.6,-169.6)..(60.3,-168)
				..controls(60.8,-166.9)and(61.4,-165.8)..(62.3,-165)
				..controls(63.1,-164.2)and(64.3,-163.4)..(65,-163.2)
				..controls(65.5,-163.1)and(72.3,-163)..(76.6,-163.1)
				..controls(77.9,-163.1)and(80.1,-163.1)..(81.4,-163)
				..controls(82.8,-162.9)and(86.3,-162.9)..(89.2,-162.9)
				..controls(93,-162.9)and(94.9,-162.8)..(96,-162.7)
				..controls(97.2,-162.5)and(101,-162.5)..(114,-162.4)
				..controls(123,-162.3)and(134.9,-162.2)..(140.2,-162.1)
				..controls(150.1,-162)and(152.2,-162.1)..(153.6,-162.4)
				..controls(154.4,-162.6)and(157.1,-163.8)..(158.1,-164.4)
				..controls(158.8,-164.8)and(159.7,-165.7)..(159.9,-166.3)
				..controls(160,-166.5)and(160.2,-167.1)..(160.2,-167.6)
				..controls(160.3,-168)and(160.6,-169.4)..(160.9,-170.6)
				..controls(161.2,-171.7)and(161.5,-173.1)..(161.6,-173.6)
				..controls(161.8,-174.4)and(161.8,-174.5)..(161.5,-174.9)
				..controls(161.2,-175.6)and(160.8,-175.9)..(159.9,-176)
				..controls(159.2,-176.1)and(154.1,-176.2)..(151.5,-176.2)
				..controls(148.1,-176.1)and(122.6,-176.1)..(121.6,-176.2)
				..controls(120.8,-176.2)and(119.8,-176.2)..(119.4,-176.1)--cycle
				(130.4,-160)..controls(129,-160)and(119.8,-160)..(110.1,-159.9)
				..controls(100.3,-159.8)and(88.5,-159.7)..(83.9,-159.6)
				..controls(77.8,-159.5)and(75.3,-159.5)..(74.5,-159.6)
				..controls(73.9,-159.7)and(72.5,-159.8)..(71.3,-159.8)
				..controls(69.6,-159.8)and(69,-159.8)..(68.5,-159.6)
				..controls(67.6,-159.3)and(66.8,-158.6)..(66.2,-157.6)
				..controls(66,-157.3)and(65.6,-156.5)..(65.3,-156.1)
				..controls(64.6,-155.1)and(64.5,-154.9)..(63.7,-152.9)
				..controls(63.2,-151.6)and(62.9,-150.9)..(61.8,-149.6)
				..controls(61.5,-149.2)and(61.1,-148.5)..(60.9,-147.9)
				..controls(60.7,-147.4)and(60.3,-146.6)..(60.1,-146.1)
				..controls(59.9,-145.5)and(59.5,-144.9)..(59,-144.3)
				..controls(58.5,-143.7)and(58.2,-143.2)..(57.9,-142.5)
				..controls(57.7,-141.9)and(57.3,-141.1)..(57,-140.7)
				..controls(56.6,-140.3)and(56.1,-139.3)..(55.8,-138.5)
				..controls(55.5,-137.8)and(55.1,-136.8)..(54.9,-136.5)
				..controls(54.7,-136.1)and(54.3,-135.2)..(53.9,-134.5)
				..controls(53.6,-133.8)and(53,-132.7)..(52.6,-132)
				..controls(51.6,-130.5)and(51.4,-130.2)..(51.1,-129)
				..controls(50.6,-127.5)and(50.6,-127.4)..(49.7,-126.3)
				..controls(49.2,-125.7)and(48.6,-124.8)..(48.3,-124.2)
				..controls(48,-123.6)and(47.5,-122.7)..(47.3,-122.2)
				..controls(47,-121.8)and(46.5,-120.7)..(46.1,-119.8)
				..controls(45.8,-118.9)and(45.2,-117.8)..(44.9,-117.2)
				..controls(44.2,-116.1)and(43.3,-114.4)..(42.3,-112.2)
				..controls(41.8,-111.3)and(41.3,-110.4)..(41.1,-110.1)
				..controls(40.5,-109.4)and(40.3,-109)..(40,-107.8)
				..controls(39.8,-107)and(39.6,-106.5)..(39.2,-106)
				..controls(38.9,-105.6)and(38.4,-104.7)..(38,-104)
				..controls(37.7,-103.4)and(37.2,-102.5)..(37,-102.1)
				..controls(36.7,-101.7)and(36.3,-100.8)..(36,-100.1)
				..controls(35.6,-98.9)and(35.3,-98.4)..(34.6,-97.6)
				..controls(34.2,-97.1)and(34.2,-96.3)..(34.6,-96)
				..controls(35,-95.7)and(35.5,-95.9)..(36.2,-96.7)
				..controls(36.5,-97)and(37.2,-97.5)..(37.6,-97.7)..controls(38,-98)and(38.6,-98.4)..(39.1,-98.9)
				..controls(39.6,-99.4)and(40.3,-100)..(40.6,-100.3)
				..controls(41.3,-100.8)and(42.4,-101.8)..(43.7,-102.9)
				..controls(44,-103.2)and(44.6,-103.7)..(45,-103.9)
				..controls(45.4,-104.2)and(45.9,-104.6)..(46.1,-104.8)
				..controls(47.1,-105.9)and(48.8,-107.1)..(50.5,-107.9)
				..controls(51.8,-108.6)and(52.3,-109)..(53.4,-110)
				..controls(54.8,-111.2)and(55.6,-111.6)..(57,-112.1)
				..controls(57.5,-112.3)and(58.5,-112.8)..(59.4,-113.3)
				..controls(60.2,-113.7)and(61.3,-114.3)..(61.9,-114.5)
				..controls(62.4,-114.7)and(63.3,-115.1)..(63.8,-115.3)
				..controls(64.4,-115.6)and(65.4,-115.9)..(66.2,-116.1)
				..controls(67,-116.3)and(68.4,-116.6)..(69.2,-116.9)
				..controls(72.3,-117.9)and(72.8,-118)..(78,-118.1)
				..controls(84,-118.3)and(86.1,-118.1)..(88.1,-117.1)
				..controls(89.4,-116.5)and(93,-115)..(93.7,-114.9)
				..controls(94.8,-114.6)and(95.4,-114.2)..(95.9,-113.4)
				..controls(96.2,-113)and(96.7,-112.5)..(97.1,-112.2)
				..controls(98,-111.6)and(101.3,-108.4)..(102.1,-107.2)
				..controls(102.4,-106.8)and(102.9,-106.2)..(103.2,-105.8)
				..controls(103.5,-105.5)and(104,-104.8)..(104.2,-104.2)
				..controls(104.5,-103.6)and(105,-102.8)..(105.3,-102.4)
				..controls(105.6,-102)and(106.1,-101.2)..(106.5,-100.6)
				..controls(107.1,-99.5)and(108.6,-97.5)..(109.5,-96.5)
				..controls(109.9,-96.1)and(110.1,-96)..(110.4,-96)
				..controls(111.1,-96)and(111.6,-96.5)..(111.8,-97.6)
				..controls(112,-98.6)and(112.3,-99.3)..(112.8,-99.9)
				..controls(113.7,-100.9)and(115.2,-103.1)..(115.7,-104.2)
				..controls(116.1,-105)and(116.5,-105.5)..(116.9,-106)
				..controls(117.3,-106.4)and(117.8,-106.9)..(118,-107.3)
				..controls(118.2,-107.6)and(118.6,-108)..(119,-108.3)
				..controls(119.6,-108.8)and(121.1,-110.3)..(121.8,-111.2)
				..controls(122.1,-111.7)and(122.5,-112)..(123.4,-112.4)
				..controls(124.8,-113.2)and(125.8,-113.8)..(127,-114.6)
				..controls(127.5,-114.9)and(128.1,-115.3)..(128.2,-115.4)
				..controls(128.4,-115.5)and(129.3,-115.8)..(130.2,-116)
				..controls(131.1,-116.2)and(132.5,-116.6)..(133.3,-116.9)
				..controls(135.7,-117.9)and(136.1,-117.9)..(142.5,-117.8)
				..controls(145.8,-117.8)and(148.4,-117.7)..(148.8,-117.6)
				..controls(149.2,-117.6)and(150.2,-117.2)..(151,-116.8)
				..controls(152,-116.4)and(152.8,-116)..(153.6,-115.9)
				..controls(154.5,-115.7)and(155.2,-115.4)..(156.3,-114.9)
				..controls(157.1,-114.5)and(158.2,-114)..(158.7,-113.8)
				..controls(159.1,-113.7)and(159.8,-113.3)..(160.2,-113)
				..controls(160.6,-112.7)and(161.3,-112.4)..(161.7,-112.2)
				..controls(162.7,-111.9)and(163.2,-111.7)..(164.1,-111.1)
				..controls(164.5,-110.8)and(165.3,-110.4)..(165.9,-110)
				..controls(166.5,-109.7)and(167.3,-109.2)..(167.6,-108.9)
				..controls(167.9,-108.6)and(168.6,-108.2)..(169,-108)
				..controls(169.4,-107.8)and(170.2,-107.3)..(170.7,-106.9)
				..controls(171.2,-106.5)and(172.1,-105.8)..(172.7,-105.4)
				..controls(173.2,-105)and(174.1,-104.4)..(174.5,-104)
				..controls(174.9,-103.7)and(175.7,-103.1)..(176.2,-102.6)
				..controls(176.7,-102.2)and(177.4,-101.5)..(177.9,-101.2)
				..controls(178.3,-100.8)and(178.8,-100.4)..(179,-100.2)
				..controls(179.2,-100)and(179.6,-99.6)..(180,-99.3)
				..controls(180.4,-99.1)and(180.9,-98.6)..(181.2,-98.4)
				..controls(181.5,-98.1)and(182.1,-97.7)..(182.5,-97.5)
				..controls(183.2,-97.2)and(183.5,-96.9)..(184.3,-95.9)
				..controls(184.6,-95.6)and(184.6,-95.6)..(184.8,-95.7)
				..controls(185,-95.8)and(185.1,-96)..(185.1,-96.2)
				..controls(185.1,-96.5)and(185,-96.7)..(184.6,-97.2)
				..controls(184.1,-97.8)and(184,-97.9)..(183.9,-98.8)
				..controls(183.7,-100.3)and(183.3,-101.3)..(182.9,-101.8)
				..controls(182.6,-102.1)and(182.3,-102.6)..(182.2,-102.9)
				..controls(182,-103.3)and(181.6,-104)..(181.2,-104.6)
				..controls(180.9,-105.2)and(180.5,-106.1)..(180.4,-106.6)
				..controls(179.9,-108)and(179.7,-108.5)..(179,-109.4)
				..controls(178.6,-109.8)and(178.2,-110.5)..(177.9,-111)
				..controls(177.7,-111.5)and(177.3,-112.3)..(177,-112.7)
				..controls(176.6,-113.2)and(176.4,-113.7)..(176.3,-114.2)
				..controls(176.2,-114.7)and(175.9,-115.2)..(175.4,-116)
				..controls(175,-116.6)and(174.4,-117.7)..(174,-118.5)
				..controls(173.7,-119.2)and(173.3,-120.1)..(173,-120.5)
				..controls(172.8,-120.9)and(172.4,-121.8)..(172.1,-122.5)
				..controls(171.8,-123.2)and(171.3,-124.1)..(171,-124.6)
				..controls(170.8,-125.1)and(170.3,-126)..(169.9,-126.6)
				..controls(169.6,-127.3)and(169.2,-128.1)..(169,-128.3)
				..controls(168.8,-128.6)and(168.5,-129.3)..(168.3,-129.8)
				..controls(168.1,-130.3)and(167.8,-131.1)..(167.5,-131.6)
				..controls(166.9,-132.5)and(164.2,-137.9)..(164.2,-138.1)
				..controls(164.2,-138.1)and(163.9,-138.8)..(163.4,-139.7)
				..controls(162.9,-140.5)and(162.4,-141.5)..(162.3,-141.9)
				..controls(161.9,-143)and(161.6,-143.4)..(160.9,-144.1)
				..controls(160.4,-144.6)and(160.3,-144.8)..(160.1,-145.4)
				..controls(160,-145.8)and(159.6,-146.8)..(159.2,-147.8)
				..controls(158.7,-148.8)and(158.1,-150)..(157.9,-150.6)
				..controls(157.6,-151.1)and(157.2,-151.9)..(156.9,-152.4)
				..controls(156.5,-152.8)and(156.1,-153.7)..(155.8,-154.4)
				..controls(155.6,-155)and(155.2,-155.9)..(154.9,-156.3)
				..controls(154.7,-156.7)and(154.4,-157.4)..(154.2,-157.9)
				..controls(154.1,-158.4)and(153.8,-159)..(153.5,-159.3)
				..controls(152.8,-160.1)and(153,-160.1)..(142.4,-160.1)
				..controls(137.2,-160.1)and(131.8,-160.1)..(130.4,-160)--cycle
				(141.5,-103.7)..controls(140.8,-103)and(139.8,-102.1)..(139.2,-101.6)
				..controls(138.5,-101.2)and(137.9,-100.6)..(137.8,-100.4)
				..controls(137,-99.5)and(134,-96.5)..(133.2,-96)
				..controls(132.7,-95.7)and(131.8,-94.9)..(131.1,-94.3)
				..controls(129.4,-92.7)and(127.9,-91.6)..(126.6,-91.1)
				..controls(125.2,-90.5)and(124.1,-90)..(122.9,-89.1)
				..controls(122.1,-88.6)and(121.3,-88.2)..(119.9,-87.7)
				..controls(118.8,-87.3)and(117.5,-86.8)..(116.9,-86.5)
				..controls(114.3,-85.5)and(113.7,-85.4)..(109.7,-85.5)
				..controls(105.3,-85.7)and(103.9,-85.9)..(101.6,-87)
				..controls(100.7,-87.4)and(99.4,-87.8)..(97.9,-88)
				..controls(97,-88.2)and(96.6,-88.3)..(95.8,-88.9)
				..controls(95.4,-89.2)and(94.7,-89.6)..(94.2,-89.8)
				..controls(92.9,-90.3)and(92.3,-90.7)..(90.9,-92.1)
				..controls(90.2,-92.7)and(89.4,-93.3)..(89,-93.5)..controls(87,-94.5)and(85.2,-96)..(83.6,-97.9)
				..controls(83.3,-98.2)and(82.5,-98.9)..(81.8,-99.4)
				..controls(80.7,-100.1)and(80.5,-100.3)..(80,-101.1)
				..controls(78.7,-103.1)and(77.2,-103.9)..(76.4,-103)
				..controls(75.7,-102.2)and(75.7,-103.2)..(75.7,-83.6)
				..controls(75.6,-73.6)and(75.6,-64.8)..(75.5,-64)
				..controls(75.5,-63.2)and(75.5,-61.4)..(75.6,-59.8)--(75.7,-57.1)--(76,-56.8)
				..controls(76.8,-56.2)and(78,-56.3)..(79.8,-57.2)
				..controls(80.4,-57.5)and(81.3,-57.8)..(81.8,-57.9)
				..controls(82.3,-58.1)and(83,-58.4)..(83.6,-58.8)..controls(84.3,-59.2)and(85,-59.5)..(86,-59.8)
				..controls(86.8,-60.1)and(88,-60.6)..(88.7,-60.9)
				..controls(89.3,-61.2)and(90.6,-61.7)..(91.4,-62)
				..controls(92.2,-62.3)and(93.7,-62.9)..(94.6,-63.3)
				..controls(95.6,-63.7)and(96.8,-64.1)..(97.4,-64.3)
				..controls(98,-64.5)and(99,-64.9)..(99.7,-65.2)
				..controls(100.4,-65.6)and(101.5,-66)..(102.3,-66.2)
				..controls(104.2,-66.8)and(105,-67.1)..(106.2,-67.8)
				..controls(107.9,-68.7)and(110.2,-69.4)..(111.3,-69.2)
				..controls(112.5,-69)and(115,-68.2)..(115.6,-67.8)
				..controls(117.3,-66.7)and(117.6,-66.6)..(119.5,-66.2)
				..controls(121.9,-65.8)and(122.1,-65.7)..(123,-65.1)
				..controls(124.1,-64.4)and(124.4,-64.3)..(126,-63.9)
				..controls(127,-63.7)and(127.7,-63.4)..(128.4,-63.1)
				..controls(129,-62.8)and(129.7,-62.5)..(130.6,-62.4)
				..controls(132.1,-62.1)and(134,-61.5)..(135.2,-61)
				..controls(135.7,-60.8)and(136.7,-60.5)..(137.3,-60.3)
				..controls(138.4,-60)and(139.3,-59.7)..(141.1,-59)
				..controls(141.6,-58.7)and(142,-58.6)..(142.5,-58.6)
				..controls(143.2,-58.6)and(143.3,-58.7)..(143.6,-59)
				..controls(144,-59.5)and(144.1,-59.9)..(144,-60.8)
				..controls(144,-61.2)and(144,-71.1)..(144,-82.8)--(143.9,-104.1)--(143.7,-104.5)
				..controls(143.5,-104.7)and(143.3,-104.8)..(143,-104.9)
				..controls(142.6,-104.9)and(142.5,-104.8)..(141.5,-103.7)--cycle
				(27.8,-91.4)..controls(27.1,-90.9)and(26.3,-90)..(26.1,-89.5)
				..controls(26,-89.2)and(25.8,-88.6)..(25.7,-87.9)
				..controls(25.4,-86.2)and(25.6,-84.8)..(26.5,-83.5)
				..controls(27.4,-82.3)and(28.3,-82)..(29.4,-82.3)
				..controls(31.1,-82.8)and(32.1,-84.9)..(32,-87.4)
				..controls(31.9,-89.3)and(31.2,-90.6)..(29.9,-91.4)
				..controls(29.1,-91.9)and(28.6,-91.9)..(27.8,-91.4)--cycle
				(190.1,-91.3)..controls(189.7,-91.2)and(189.3,-90.9)..(189.1,-90.7)
				..controls(188.3,-89.8)and(187.6,-87.7)..(187.6,-86.2)
				..controls(187.6,-85.5)and(187.7,-85.2)..(188.2,-84.4)
				..controls(189.1,-82.6)and(190.1,-81.8)..(191.3,-82.1)
				..controls(192,-82.3)and(193.5,-83.7)..(193.8,-84.4)
				..controls(194.1,-85.1)and(194.2,-87.1)..(194,-88.3)
				..controls(193.8,-89.8)and(193.3,-90.5)..(192.3,-91.3)
				..controls(191.7,-91.7)and(190.9,-91.7)..(190.1,-91.3)--cycle
				(173.7,-83.4)..controls(172.2,-82.5)and(172,-82.5)..(170.6,-82.4)
				..controls(169.7,-82.4)and(168.8,-82.4)..(168.3,-82.5)
				..controls(167.2,-82.7)and(166.9,-82.7)..(166.6,-82.2)
				..controls(166.3,-81.8)and(166.3,-81)..(166.8,-80)
				..controls(167.2,-78.9)and(167.4,-77.4)..(167.7,-74.1)
				..controls(167.8,-72.8)and(167.9,-71.8)..(167.9,-71.7)
				..controls(168,-71.7)and(168,-70.2)..(168,-68.4)..controls(168,-65.4)and(168,-64.9)..(167.7,-63)
				..controls(167.3,-60.5)and(167.1,-59.9)..(166.4,-58.4)
				..controls(166.1,-57.7)and(165.8,-56.9)..(165.7,-56.6)
				..controls(165.7,-56.2)and(165.4,-55.6)..(165.2,-55.2)
				..controls(164.6,-53.8)and(161.9,-50.7)..(161,-50.3)
				..controls(159.8,-49.7)and(157.9,-50)..(155.2,-51.2)
				..controls(154.7,-51.4)and(153.4,-51.8)..(152.2,-52.1)
				..controls(150.8,-52.4)and(149.4,-52.9)..(148.6,-53.2)
				..controls(147.8,-53.6)and(146.7,-54)..(146,-54.2)
				..controls(145.4,-54.4)and(144.5,-54.7)..(144,-54.9)
				..controls(143.6,-55.2)and(142.7,-55.5)..(141.8,-55.7)
				..controls(141,-55.9)and(139.6,-56.3)..(138.7,-56.6)
				..controls(137.8,-57)and(136.3,-57.4)..(135.3,-57.6)
				..controls(133.8,-57.9)and(133.2,-58.1)..(132.1,-58.7)
				..controls(131.1,-59.2)and(130.3,-59.4)..(129.4,-59.6)
				..controls(128.7,-59.8)and(127.8,-60)..(127.4,-60.1)
				..controls(127,-60.3)and(126,-60.5)..(125.1,-60.7)
				..controls(123.2,-61.1)and(121.7,-61.5)..(120.2,-62.1)
				..controls(119.7,-62.3)and(118.1,-62.8)..(116.7,-63.2)
				..controls(114.6,-63.9)and(114,-64.1)..(113,-64.7)
				..controls(111.9,-65.4)and(111.8,-65.4)..(111,-65.5)
				..controls(109.9,-65.5)and(109.3,-65.4)..(107.7,-64.7)
				..controls(107,-64.4)and(106,-64)..(105.3,-63.9)
				..controls(103.5,-63.5)and(102.2,-63)..(101.2,-62.5)
				..controls(100.1,-61.8)and(99.3,-61.5)..(97.8,-61.1)
				..controls(97.2,-61)and(96.2,-60.7)..(95.7,-60.4)
				..controls(95.1,-60.1)and(94.2,-59.8)..(93.5,-59.6)
				..controls(92.9,-59.4)and(91.7,-59)..(91,-58.7)..controls(90.2,-58.4)and(88.9,-58)..(88,-57.7)
				..controls(87.2,-57.5)and(86.1,-57.1)..(85.5,-56.8)
				..controls(85,-56.6)and(84.1,-56.2)..(83.5,-56)..controls(82.9,-55.7)and(81.8,-55.2)..(81,-54.9)
				..controls(80.1,-54.5)and(79.1,-54.2)..(78.7,-54.1)
				..controls(77.9,-53.8)and(77.3,-53.6)..(76.2,-52.9)
				..controls(75.7,-52.6)and(75.2,-52.4)..(74.1,-52.2)
				..controls(72.5,-51.9)and(72.4,-51.8)..(71.3,-51.1)
				..controls(69.9,-50.1)and(69.4,-49.9)..(67.9,-49.8)
				..controls(67,-49.7)and(66.6,-49.6)..(66.3,-49.4)
				..controls(66.1,-49.3)and(65.7,-49)..(65.4,-48.8)
				..controls(65.1,-48.7)and(64.2,-48.4)..(63.6,-48.2)
				..controls(62.9,-48)and(61.7,-47.6)..(60.9,-47.2)
				..controls(60.1,-46.8)and(58.9,-46.4)..(58.2,-46.2)
				..controls(57.6,-46)and(56.8,-45.7)..(56.5,-45.5)
				..controls(54.9,-44.6)and(54.6,-44.5)..(53.2,-44.2)
				..controls(51.7,-43.8)and(51,-43.6)..(49.9,-42.9)
				..controls(49.5,-42.6)and(48.9,-42.4)..(47.5,-42)
				..controls(45.7,-41.6)and(45.6,-41.5)..(44.3,-40.6)
				..controls(43.5,-40.1)and(42.8,-39.6)..(42.7,-39.4)
				..controls(42.4,-38.7)and(43,-38.1)..(44.5,-37.6)..controls(44.9,-37.5)and(46,-37)..(46.9,-36.6)
				..controls(47.9,-36.1)and(48.8,-35.8)..(49.6,-35.6)
				..controls(50.5,-35.4)and(51.3,-35.2)..(52.4,-34.6)
				..controls(53.3,-34.1)and(54.5,-33.7)..(54.9,-33.5)
				..controls(55.4,-33.4)and(56.4,-33)..(57.2,-32.6)
				..controls(58.6,-32)and(60.1,-31.5)..(61.6,-31.2)
				..controls(62.3,-31.1)and(62.7,-31)..(63.2,-30.7)
				..controls(64.1,-30)and(64.3,-29.9)..(65.3,-29.7)
				..controls(65.7,-29.6)and(66.5,-29.2)..(67.1,-28.9)
				..controls(67.6,-28.6)and(68.8,-28.2)..(69.6,-27.9)
				..controls(70.5,-27.7)and(71.6,-27.3)..(72.1,-27)
				..controls(72.6,-26.7)and(73.8,-26.2)..(74.8,-25.9)
				..controls(75.7,-25.5)and(76.7,-25.1)..(77.1,-25)
				..controls(77.4,-24.8)and(78.5,-24.4)..(79.6,-24.1)
				..controls(80.6,-23.8)and(82,-23.2)..(82.8,-22.9)
				..controls(83.5,-22.5)and(84.8,-22)..(85.8,-21.8)
				..controls(86.7,-21.5)and(87.8,-21.1)..(88.1,-20.9)
				..controls(88.4,-20.7)and(89.3,-20.3)..(90.1,-20.1)
				..controls(90.8,-19.8)and(91.9,-19.4)..(92.4,-19.1)
				..controls(93.5,-18.5)and(94.4,-18.2)..(95.4,-18.1)
				..controls(95.7,-18)and(96.4,-17.8)..(96.9,-17.5)
				..controls(97.5,-17.3)and(98.4,-16.9)..(99,-16.6)
				..controls(99.7,-16.3)and(100.9,-15.9)..(101.9,-15.6)
				..controls(102.9,-15.3)and(104.4,-14.7)..(105.3,-14.3)
				..controls(108.1,-13.2)and(108.5,-13.1)..(110.1,-13.1)
				..controls(111.4,-13.1)and(111.5,-13.1)..(112.4,-13.5)
				..controls(113.2,-14)and(115,-14.8)..(116.6,-15.4)
				..controls(116.9,-15.5)and(117.9,-15.8)..(118.8,-16.1)
				..controls(120.9,-16.6)and(121.5,-16.9)..(122.6,-17.9)
				..controls(123,-18.3)and(123.7,-18.8)..(123.9,-18.9)
				..controls(124.6,-19.4)and(126,-19.8)..(127.9,-20.2)
				..controls(129.6,-20.5)and(129.8,-20.6)..(131.1,-21.4)
				..controls(131.5,-21.7)and(132.4,-22.1)..(133.1,-22.3)
				..controls(136.1,-23.2)and(138.3,-24.1)..(140.2,-25)
				..controls(141,-25.5)and(141.5,-25.6)..(142.5,-25.8)
				..controls(145.1,-26.3)and(146,-26.6)..(149.4,-28.1)
				..controls(151.3,-28.9)and(153.5,-29.8)..(154.3,-30.1)
				..controls(155.1,-30.4)and(156.2,-30.8)..(156.6,-31)
				..controls(157.1,-31.2)and(158.2,-31.6)..(159.1,-31.8)
				..controls(160.6,-32.2)and(160.9,-32.3)..(161.7,-32.9)
				..controls(162.4,-33.3)and(162.9,-33.6)..(163.7,-33.8)
				..controls(164.3,-34)and(164.8,-34.1)..(164.9,-34.1)
				..controls(165,-34.1)and(165.6,-34.4)..(166.4,-34.8)
				..controls(167.3,-35.2)and(168.1,-35.4)..(169,-35.6)
				..controls(170,-35.8)and(170.5,-36)..(171.7,-36.6)
				..controls(172.4,-37)and(173.3,-37.4)..(173.6,-37.4)
				..controls(175.6,-37.7)and(176.2,-38)..(176.7,-39.1)
				..controls(176.9,-39.6)and(176.9,-39.7)..(176.7,-40.4)
				..controls(176.4,-41.6)and(175.8,-42.7)..(175.2,-43.4)
				..controls(174.5,-44.1)and(172.6,-45.5)..(171.7,-45.8)
				..controls(170.6,-46.2)and(169.1,-46)..(168.2,-45.4)
				..controls(167.9,-45.2)and(167.3,-44.8)..(166.9,-44.4)
				..controls(166.4,-44.1)and(165.8,-43.6)..(165.4,-43.2)
				..controls(165.1,-42.9)and(164.4,-42.3)..(163.9,-41.9)
				..controls(163.3,-41.5)and(162.7,-40.9)..(162.4,-40.5)
				..controls(161.4,-39.2)and(159.9,-38)..(158.8,-37.8)
				..controls(158.4,-37.8)and(158.3,-37.7)..(158.2,-37.2)
				..controls(158,-36.8)and(157.4,-36.1)..(157.1,-36.1)
				..controls(157,-36.1)and(156.7,-36.1)..(156.5,-36)
				..controls(156.1,-35.9)and(155.9,-35.7)..(155.7,-35.2)
				..controls(155.3,-34.6)and(155.1,-34.4)..(154,-34.1)
				..controls(153.6,-34)and(153.2,-33.7)..(152.9,-33.5)
				..controls(152.4,-32.9)and(151.6,-32.4)..(151.3,-32.4)
				..controls(150.9,-32.4)and(150.4,-32.6)..(150.4,-32.9)
				..controls(150.4,-33.2)and(150.8,-33.4)..(151.4,-33.6)
				..controls(152.1,-33.7)and(152.8,-34.1)..(153.5,-34.9)
				..controls(154.1,-35.5)and(154.3,-35.6)..(155,-35.8)
				..controls(155.7,-36)and(155.7,-36)..(156,-36.6)
				..controls(156.3,-37.3)and(156.7,-37.7)..(157.5,-37.8)
				..controls(158.1,-37.9)and(158.2,-38)..(158.4,-38.8)
				..controls(158.5,-39.4)and(158.6,-39.6)..(159.6,-40.4)
				..controls(160.2,-41)and(160.9,-41.6)..(161.3,-41.8)
				..controls(161.7,-42)and(161.9,-42.3)..(162.2,-42.8)
				..controls(162.4,-43.2)and(162.8,-43.8)..(163.1,-44.1)
				..controls(163.9,-44.8)and(165.1,-46.7)..(165.7,-48)
				..controls(166.2,-49.3)and(166.4,-49.7)..(167.2,-50.5)
				..controls(167.5,-50.8)and(167.9,-51.4)..(168.1,-51.8)
				..controls(168.3,-52.2)and(168.8,-53)..(169.1,-53.6)
				..controls(169.6,-54.3)and(169.8,-54.9)..(169.9,-55.6)
				..controls(170.4,-57.8)and(170.5,-58.1)..(171.2,-59.5)
				..controls(171.6,-60.3)and(172,-61.5)..(172.2,-62.1)
				..controls(172.3,-62.8)and(172.7,-63.7)..(172.9,-64.1)
				..controls(173.2,-64.6)and(173.4,-65.3)..(173.4,-65.6)
				..controls(173.8,-68.8)and(174.2,-70.6)..(174.9,-72.6)
				..controls(175.5,-74.5)and(175.9,-77.1)..(176.1,-80.2)
				..controls(176.3,-83.4)and(175.5,-84.4)..(173.7,-83.4)--cycle;
				\definecolor{ezs}{RGB}{247,179,26}\fill[ezs](-5.7,-156.6)--(-3.4,-157.5)--(61.9,12.7)--(59.7,13.6)--cycle;
				\definecolor{ezs}{RGB}{247,179,22}\fill[ezs](-3.8,-157.3)--(-1.5,-158.2)--(63.9,12)--(61.6,12.9)--cycle;
				\definecolor{ezs}{RGB}{248,179,19}\fill[ezs](-1.8,-158.1)--(0.4,-158.9)--(65.8,11.2)--(63.5,12.1)--cycle;
				\definecolor{ezs}{RGB}{248,178,16}\fill[ezs](0.1,-158.8)--(2.4,-159.7)--(67.7,10.5)--(65.4,11.4)--cycle;
				\definecolor{ezs}{RGB}{249,178,15}\fill[ezs](2,-159.5)--(4.3,-160.4)--(69.6,9.8)--(67.3,10.7)--cycle;
				\definecolor{ezs}{RGB}{251,177,17}\fill[ezs](3.9,-160.3)--(6.2,-161.1)--(71.5,9)--(69.2,9.9)--cycle;
				\definecolor{ezs}{RGB}{252,179,18}\fill[ezs](5.8,-161)--(8.1,-161.9)--(73.4,8.3)--(71.1,9.2)--cycle;
				\definecolor{ezs}{RGB}{252,182,17}\fill[ezs](7.7,-161.7)--(10,-162.6)--(75.3,7.6)--(73,8.5)--cycle;
				\definecolor{ezs}{RGB}{252,185,17}\fill[ezs](9.6,-162.5)--(11.9,-163.4)--(77.2,6.8)--(75,7.7)--cycle;
				\definecolor{ezs}{RGB}{251,189,18}\fill[ezs](11.5,-163.2)--(13.8,-164.1)--(79.2,6.1)--(76.9,7)--cycle;
				\definecolor{ezs}{RGB}{251,191,22}\fill[ezs](13.4,-163.9)--(15.7,-164.8)--(81.1,5.4)--(78.8,6.2)--cycle;
				\definecolor{ezs}{RGB}{252,193,25}\fill[ezs](15.4,-164.7)--(17.6,-165.6)--(83,4.6)--(80.7,5.5)--cycle;
				\definecolor{ezs}{RGB}{253,195,27}\fill[ezs](17.3,-165.4)--(19.6,-166.3)--(84.9,3.9)--(82.6,4.8)--cycle;
				\definecolor{ezs}{RGB}{253,196,31}\fill[ezs](19.2,-166.1)--(21.5,-167)--(86.8,3.2)--(84.5,4)--cycle;
				\definecolor{ezs}{RGB}{254,199,35}\fill[ezs](21.1,-166.9)--(23.4,-167.8)--(88.7,2.4)--(86.4,3.3)--cycle;
				\definecolor{ezs}{RGB}{253,201,38}\fill[ezs](23,-167.6)--(25.3,-168.5)--(90.6,1.7)--(88.3,2.6)--cycle;
				\definecolor{ezs}{RGB}{253,203,40}\fill[ezs](24.9,-168.3)--(27.2,-169.2)--(92.5,1)--(90.3,1.8)--cycle;
				\definecolor{ezs}{RGB}{253,205,42}\fill[ezs](26.8,-169.1)--(29.1,-170)--(94.4,0.2)--(92.2,1.1)--cycle;
				\definecolor{ezs}{RGB}{253,207,45}\fill[ezs](28.7,-169.8)--(31,-170.7)--(96.4,-0.5)--(94.1,0.4)--cycle;
				\definecolor{ezs}{RGB}{253,210,49}\fill[ezs](30.7,-170.5)--(32.9,-171.4)--(98.3,-1.2)--(96,-0.4)--cycle;
				\definecolor{ezs}{RGB}{253,212,54}\fill[ezs](32.6,-171.3)--(34.9,-172.2)--(100.2,-2)--(97.9,-1.1)--cycle;
				\definecolor{ezs}{RGB}{253,214,59}\fill[ezs](34.5,-172)--(36.8,-172.9)--(102.1,-2.7)--(99.8,-1.8)--cycle;
				\definecolor{ezs}{RGB}{253,216,63}\fill[ezs](36.4,-172.7)--(38.7,-173.6)--(104,-3.4)--(101.7,-2.6)--cycle;
				\definecolor{ezs}{RGB}{253,218,65}\fill[ezs](38.3,-173.5)--(40.6,-174.4)--(105.9,-4.2)--(103.6,-3.3)--cycle;
				\definecolor{ezs}{RGB}{254,220,66}\fill[ezs](40.2,-174.2)--(42.5,-175.1)--(107.8,-4.9)--(105.5,-4)--cycle;
				\definecolor{ezs}{RGB}{254,222,68}\fill[ezs](42.1,-175)--(44.4,-175.8)--(109.7,-5.6)--(107.5,-4.8)--cycle;
				\definecolor{ezs}{RGB}{254,224,71}\fill[ezs](44,-175.7)--(46.3,-176.6)--(111.7,-6.4)--(109.4,-5.5)--cycle;
				\definecolor{ezs}{RGB}{254,226,75}\fill[ezs](46,-176.4)--(48.2,-177.3)--(113.6,-7.1)--(111.3,-6.2)--cycle;
				\definecolor{ezs}{RGB}{254,228,79}\fill[ezs](47.9,-177.2)--(50.1,-178)--(115.5,-7.8)--(113.2,-7)--cycle;
				\definecolor{ezs}{RGB}{254,230,82}\fill[ezs](49.8,-177.9)--(52.1,-178.8)--(117.4,-8.6)--(115.1,-7.7)--cycle;
				\definecolor{ezs}{RGB}{254,232,85}\fill[ezs](51.7,-178.6)--(54,-179.5)--(119.3,-9.3)--(117,-8.4)--cycle;
				\definecolor{ezs}{RGB}{254,234,88}\fill[ezs](53.6,-179.4)--(55.9,-180.2)--(121.2,-10)--(118.9,-9.2)--cycle;
				\definecolor{ezs}{RGB}{254,237,91}\fill[ezs](55.5,-180.1)--(57.8,-181)--(123.1,-10.8)--(120.8,-9.9)--cycle;
				\definecolor{ezs}{RGB}{254,239,94}\fill[ezs](57.4,-180.8)--(59.7,-181.7)--(125,-11.5)--(122.8,-10.6)--cycle;
				\definecolor{ezs}{RGB}{254,241,97}\fill[ezs](59.3,-181.6)--(61.6,-182.4)--(127,-12.2)--(124.7,-11.4)--cycle;
				\definecolor{ezs}{RGB}{254,243,101}\fill[ezs](61.2,-182.3)--(63.5,-183.2)--(128.9,-13)--(126.6,-12.1)--cycle;
				\definecolor{ezs}{RGB}{254,245,103}\fill[ezs](63.2,-183)--(65.4,-183.9)--(130.8,-13.7)--(128.5,-12.8)--cycle;
				\definecolor{ezs}{RGB}{254,245,105}\fill[ezs](65.1,-183.8)--(67.4,-184.6)--(132.7,-14.4)--(130.4,-13.6)--cycle;
				\definecolor{ezs}{RGB}{254,245,105}\fill[ezs](67,-184.5)--(69.3,-185.4)--(134.6,-15.2)--(132.3,-14.3)--cycle;
				\definecolor{ezs}{RGB}{254,244,103}\fill[ezs](68.9,-185.2)--(71.2,-186.1)--(136.5,-15.9)--(134.2,-15)--cycle;
				\definecolor{ezs}{RGB}{253,243,101}\fill[ezs](70.8,-186)--(73.1,-186.8)--(138.4,-16.6)--(136.1,-15.8)--cycle;
				\definecolor{ezs}{RGB}{253,241,98}\fill[ezs](72.7,-186.7)--(75,-187.6)--(140.3,-17.4)--(138,-16.5)--cycle;
				\definecolor{ezs}{RGB}{253,239,95}\fill[ezs](74.6,-187.4)--(76.9,-188.3)--(142.2,-18.1)--(140,-17.2)--cycle;
				\definecolor{ezs}{RGB}{253,236,92}\fill[ezs](76.5,-188.2)--(78.8,-189)--(144.2,-18.8)--(141.9,-18)--cycle;
				\definecolor{ezs}{RGB}{253,234,88}\fill[ezs](78.5,-188.9)--(80.7,-189.8)--(146.1,-19.6)--(143.8,-18.7)--cycle;
				\definecolor{ezs}{RGB}{253,232,85}\fill[ezs](80.4,-189.6)--(82.7,-190.5)--(148,-20.3)--(145.7,-19.4)--cycle;
				\definecolor{ezs}{RGB}{253,230,82}\fill[ezs](82.3,-190.4)--(84.6,-191.2)--(149.9,-21)--(147.6,-20.2)--cycle;
				\definecolor{ezs}{RGB}{253,228,78}\fill[ezs](84.2,-191.1)--(86.5,-192)--(151.8,-21.8)--(149.5,-20.9)--cycle;
				\definecolor{ezs}{RGB}{253,226,74}\fill[ezs](86.1,-191.8)--(88.4,-192.7)--(153.7,-22.5)--(151.4,-21.6)--cycle;
				\definecolor{ezs}{RGB}{253,224,70}\fill[ezs](88,-192.6)--(90.3,-193.4)--(155.6,-23.3)--(153.3,-22.4)--cycle;
				\definecolor{ezs}{RGB}{253,222,67}\fill[ezs](89.9,-193.3)--(92.2,-194.2)--(157.5,-24)--(155.3,-23.1)--cycle;
				\definecolor{ezs}{RGB}{253,220,65}\fill[ezs](91.8,-194)--(94.1,-194.9)--(159.5,-24.7)--(157.2,-23.8)--cycle;
				\definecolor{ezs}{RGB}{253,217,63}\fill[ezs](93.7,-194.8)--(96,-195.6)--(161.4,-25.5)--(159.1,-24.6)--cycle;
				\definecolor{ezs}{RGB}{253,215,60}\fill[ezs](95.7,-195.5)--(97.9,-196.4)--(163.3,-26.2)--(161,-25.3)--cycle;
				\definecolor{ezs}{RGB}{253,213,56}\fill[ezs](97.6,-196.2)--(99.9,-197.1)--(165.2,-26.9)--(162.9,-26)--cycle;
				\definecolor{ezs}{RGB}{252,211,53}\fill[ezs](99.5,-197)--(101.8,-197.8)--(167.1,-27.7)--(164.8,-26.8)--cycle;
				\definecolor{ezs}{RGB}{252,209,49}\fill[ezs](101.4,-197.7)--(103.7,-198.6)--(169,-28.4)--(166.7,-27.5)--cycle;
				\definecolor{ezs}{RGB}{252,207,45}\fill[ezs](103.3,-198.4)--(105.6,-199.3)--(170.9,-29.1)--(168.6,-28.2)--cycle;
				\definecolor{ezs}{RGB}{253,205,42}\fill[ezs](105.2,-199.2)--(107.5,-200)--(172.8,-29.9)--(170.6,-29)--cycle;
				\definecolor{ezs}{RGB}{253,202,40}\fill[ezs](107.1,-199.9)--(109.4,-200.8)--(174.7,-30.6)--(172.5,-29.7)--cycle;
				\definecolor{ezs}{RGB}{253,200,38}\fill[ezs](109,-200.6)--(111.3,-201.5)--(176.7,-31.3)--(174.4,-30.4)--cycle;
				\definecolor{ezs}{RGB}{253,198,35}\fill[ezs](111,-201.4)--(113.2,-202.2)--(178.6,-32.1)--(176.3,-31.2)--cycle;
				\definecolor{ezs}{RGB}{253,196,31}\fill[ezs](112.9,-202.1)--(115.2,-203)--(180.5,-32.8)--(178.2,-31.9)--cycle;
				\definecolor{ezs}{RGB}{253,194,28}\fill[ezs](114.8,-202.8)--(117.1,-203.7)--(182.4,-33.5)--(180.1,-32.6)--cycle;
				\definecolor{ezs}{RGB}{253,192,26}\fill[ezs](116.7,-203.6)--(119,-204.5)--(184.3,-34.3)--(182,-33.4)--cycle;
				\definecolor{ezs}{RGB}{253,189,24}\fill[ezs](118.6,-204.3)--(120.9,-205.2)--(186.2,-35)--(183.9,-34.1)--cycle;
				\definecolor{ezs}{RGB}{253,186,20}\fill[ezs](120.5,-205)--(122.8,-205.9)--(188.1,-35.7)--(185.8,-34.8)--cycle;
				\definecolor{ezs}{RGB}{253,184,17}\fill[ezs](122.4,-205.8)--(124.7,-206.7)--(190,-36.5)--(187.8,-35.6)--cycle;
				\definecolor{ezs}{RGB}{252,183,15}\fill[ezs](124.3,-206.5)--(126.6,-207.4)--(192,-37.2)--(189.7,-36.3)--cycle;
				\definecolor{ezs}{RGB}{252,181,14}\fill[ezs](126.3,-207.2)--(128.5,-208.1)--(193.9,-37.9)--(191.6,-37.1)--cycle;
				\definecolor{ezs}{RGB}{252,179,13}\fill[ezs](128.2,-208)--(130.4,-208.9)--(195.8,-38.7)--(193.5,-37.8)--cycle;
				\definecolor{ezs}{RGB}{252,178,12}\fill[ezs](130.1,-208.7)--(132.4,-209.6)--(197.7,-39.4)--(195.4,-38.5)--cycle;
				\definecolor{ezs}{RGB}{251,178,10}\fill[ezs](132,-209.4)--(134.3,-210.3)--(199.6,-40.1)--(197.3,-39.3)--cycle;
				\definecolor{ezs}{RGB}{251,178,10}\fill[ezs](133.9,-210.2)--(136.2,-211.1)--(201.5,-40.9)--(199.2,-40)--cycle;
				\definecolor{ezs}{RGB}{250,177,11}\fill[ezs](135.8,-210.9)--(138.1,-211.8)--(203.4,-41.6)--(201.1,-40.7)--cycle;
				\definecolor{ezs}{RGB}{250,178,11}\fill[ezs](137.7,-211.6)--(140,-212.5)--(205.3,-42.3)--(203.1,-41.5)--cycle;
				\definecolor{ezs}{RGB}{248,180,10}\fill[ezs](139.6,-212.4)--(141.9,-213.3)--(207.3,-43.1)--(205,-42.2)--cycle;
				\definecolor{ezs}{RGB}{247,181,9}\fill[ezs](141.5,-213.1)--(143.8,-214)--(209.2,-43.8)--(206.9,-42.9)--cycle;
				\definecolor{ezs}{RGB}{246,182,9}\fill[ezs](143.5,-213.8)--(145.7,-214.7)--(211.1,-44.5)--(208.8,-43.7)--cycle;
				\definecolor{ezs}{RGB}{246,182,10}\fill[ezs](145.4,-214.6)--(147.7,-215.5)--(213,-45.3)--(210.7,-44.4)--cycle;
			\end{scope}
			% chữ EDUZY (Times đậm, làm dày nét bằng 5 lớp lệch nhẹ; căn đúng khung chữ logo gốc)
			\foreach \dx/\dy in {-0.7/0,0.7/0,0/-0.3,0/0.3,0/0}
			\node[anchor=center,inner sep=0pt,outer sep=0pt,text=ezText,align=center] at (109.6+\dx,-193.6+\dy)
			{\resizebox{1.01cm}{0.22cm}{\fontencoding{T1}\fontfamily{ptm}\fontseries{b}\selectfont EDUZY}};
		\end{tikzpicture}%
	}%
}
%====================== THIẾT LẬP THÂN ĐỀ ======================
%--- Môi trường tự luận {bt}: ex_test 3.3.4 chưa khai báo sẵn nên tạo ở đây
%    (ex_test đã tự móc \fixshowans và \listenumerate cho bt)
\makeatletter
\@ifundefined{bt}{\newtheorem{bt}{\color{gold1}\namebt}}{}
\makeatother
\renewcommand{\sectionmark}[1]{\markright{\thesection\ #1}}
%--- Khung tiêu đề (logo + tối đa 4 dòng; dòng 3, 4 để trống {} thì bỏ qua)
%  Đề thi:    \khungtieude{THI THỬ THPT QUỐC GIA NĂM 2026}{Môn Toán -- Đề số 21}
%                         {(Đề thi gồm: ...)}{Thời gian làm bài: 90 phút (không kể thời gian phát đề).}
%  Chuyên đề: \khungtieude{CHUYÊN ĐỀ 3 -- DÃY SỐ VÀ CÔNG THỨC TRUY HỒI}
%                         {Bồi dưỡng học sinh giỏi Toán THPT}{...}{}
\newcommand{\khungtieude}[4]{%
	\begin{tcolorbox}[colframe=gold1,colback=gold5!45,boxrule=0.5mm]
		\begin{minipage}[c]{3cm}
			\centerline{\eduzylogo[2.4cm]}
		\end{minipage}\hspace{0.5cm}
		\begin{minipage}[c]{12.5cm}
			\vskip 0.15 cm
			\centerline{\ind{#1}}
			\vskip 0.15 cm
			\centerline{\textbf{#2}}
			\if\relax\detokenize{#3}\relax\else
				\vskip 0.15 cm
				\centerline{\textit{#3}}
			\fi
			\if\relax\detokenize{#4}\relax\else
				\vskip 0.15 cm
				\centerline{\textit{#4}}
			\fi
			\vskip 0.15 cm
		\end{minipage}
	\end{tcolorbox}
}
%====================== BỔ SUNG CHO BÀI HỌC, CHUYÊN ĐỀ ======================
\usepackage[shortlabels]{enumitem}
\usepackage{tabularx}
%--- Danh sách gọn: \begin{enumerate}[gon] ... \end{enumerate}
\SetEnumitemKey{gon}{leftmargin=*,itemsep=1pt,parsep=0pt,topsep=3pt}
%--- Khung lý thuyết: \begin{khung}{} ... \end{khung}; tiêu đề chỉ dùng khi thật ngắn ("Mẹo", "Lỗi thường gặp")
\newtcolorbox{khung}[1]{breakable,enhanced,colframe=gold1,colback=gold5!45,
	boxrule=0.4mm,fonttitle=\bfseries\sffamily,coltitle=white,colbacktitle=gold1,
	attach boxed title to top left={xshift=6mm,yshift=-3mm},
	boxed title style={arc=1mm,boxrule=0.3mm},title={#1}}
%--- Tiêu đề mục nhỏ trong một dạng: \muc{Dấu hiệu}, \muc{Phương pháp}, \muc{Bài tự luyện}
\makeatletter
\newcommand{\muc}[1]{\par\medskip\noindent\indam{#1}\par\nobreak\smallskip\@afterindentfalse\@afterheading}
\makeatother
%--- Ô công thức cần nhớ (căn giữa, viền vàng, vừa khít nội dung)
%     \begin{center}\congthuc{$Q(n)=(n-1)!$}\end{center}
\newtcbox{\congthuc}{on line,colframe=gold1,colback=gold5!45,boxrule=0.5pt,arc=2pt,
	left=8pt,right=8pt,top=4pt,bottom=4pt}
%--- Công tắc lời giải bài tự luyện (ví dụ mẫu luôn dùng \loigiai)
%     \loigiaiTL[{đáp số}]{lời giải đầy đủ}
%     \TLfull -> in lời giải đầy đủ;  \TLdapso -> chỉ in một dòng "Đáp số";  \TLan -> ẩn hoàn toàn
%     (\CheDo là dethi thì mọi lời giải và đáp số đều ẩn)
\makeatletter
\newif\ifTLhien   \TLhientrue
\newif\ifTLdapso  \TLdapsofalse
\newcommand{\TLfull}{\TLhientrue\TLdapsofalse}
\newcommand{\TLdapso}{\TLhientrue\TLdapsotrue}
\newcommand{\TLan}{\TLhienfalse}
\newcommand{\loigiaiTL}[2][]{%
	\ifTLhien
		\ifTLdapso
			\def\@tlarg{#1}%
			\ifx\@tlarg\@empty\else
				\ifdefstring{\kindOPTN}{dethi}{}{\par\noindent{\bfseries\color{gold1}Đáp số:} #1\par}%
			\fi
		\else
			\loigiai{#2}%
		\fi
	\fi}
\makeatother
\TLdapso% mặc định: chỉ in đáp số bài tự luyện
%====================== BẢNG ĐÚNG/SAI (\choiceTFt) ======================
%  Hàng tiêu đề nền nâu đồng chữ trắng (Phát biểu | Đ | S), các hàng tô xen kẽ kem / trắng,
%  đường kẻ vàng kim, nhãn a) b) c) d) màu nâu đồng. Bản có đáp án: cột được chọn có vòng tròn
%  nâu đồng với dấu tích trắng (câu đúng tích ở cột Đúng, câu sai tích ở cột Sai).
%  Tiêu đề ghi đủ "Đúng | Sai": \OPTN{viettat=0} sau \input (cột Đúng/Sai rộng ra, phát biểu hẹp lại).
%  Dùng bảng gốc của ex_test: \def\BangDSgoc{1} trước %============================================================
%  PREAMBLE EDUZY v2 - dùng chung cho đề thi, bài học, chuyên đề (gói ex_test)
%  Khác bản preamble-eduzy.tex ở đầu trang, chân trang: dải nền nâu đồng tràn ngang,
%  logo + Eduzy.com.vn, ô chữ vàng kem cạnh xiên, chân trang "Trang (số)";
%  logo Eduzy chìm giữa mọi trang (tắt: \logochimtat, bật lại: \logochimbat).
%  Cách dùng ở file chính (xem mau-chuan.tex, mau-chuyen-de.tex):
%      \def\CheDo{loigiai}          % loigiai | dethi | solcolor | color | book
%      \def\tenchuong{<dòng chữ trên đầu trang>}
%      \def\thuonghieu{<chữ ở chân trang lẻ>}   % không khai báo thì dùng mặc định
%      \def\tacgia{<chữ ở chân trang chẵn>}      % không khai báo thì dùng mặc định
%      \def\khauhieu{<chữ cạnh logo ở đầu trang>}  % tùy chọn, mặc định Eduzy.com.vn
%      \input{preamble-eduzy-v2}
%      \TLdapso                      % tài liệu: \TLfull | \TLdapso | \TLan
%      \begin{document} ... \end{document}
%  Logo Eduzy vẽ bằng TikZ ngay trong file (\eduzylogo), không cần file ảnh.
%  Biên dịch: pdfLaTeX (máy có gói vietnam) hoặc XeLaTeX.
%============================================================
\documentclass[12pt,a4paper,twoside]{article}%<-- twoside: bat dau trang/chan trang chan-le
\usepackage{iftex}
\usepackage{amsmath,amssymb,tikz,mathrsfs,tkz-tab}
\usepackage{graphicx}
\usepackage{pgfplots}% đồ thị hàm số
\pgfplotsset{compat=1.17}
%\usepackage{mathptmx}
\usepackage{mathpazo}
\ifXeTeX% XeLaTeX: font chữ có đủ dấu tiếng Việt, giữ nguyên font toán mathpazo
	\usepackage[no-math]{fontspec}
	\setmainfont{TeX Gyre Pagella}
	\newfontfamily\EZqag{TeX Gyre Adventor}[NFSSFamily=qag,Ligatures=TeX]
\fi
\usepackage[many]{tcolorbox}
\usetikzlibrary{math,through,calc,intersections,angles,quotes,shapes,shapes.geometric,arrows,patterns,matrix,chains,arrows.meta,decorations.shapes,decorations.fractals,decorations.markings,shadows,decorations.pathreplacing}
\usepackage[top=1.6cm, bottom=1.8cm, left=1.8cm, right=1.8cm]{geometry}% v2: top 1.6cm (bản cũ 1.5cm) để chữ không sát dải đầu trang
\providecommand{\CheDo}{loigiai}% loigiai | dethi | solcolor | color | book
\expandafter\usepackage\expandafter[\CheDo]{ex_test}
% \tenchuong: khai báo ở file chính (trước \input) hoặc \renewcommand sau \input
%============================================================
%  [GỘP TỪ eduzy-de.tex] EDUZY - MÀU NÂU ĐỒNG / VÀNG KIM
%  + ĐẦU TRANG, CHÂN TRANG KIỂU DẢI NGANG (v2), ĐỐI XỨNG TRANG CHẴN / TRANG LẺ.
%  (không cần file eduzy-de.tex riêng nữa)
%  Khối này phải đặt SAU \usepackage[...]{ex_test} và SAU \def\tenchuong{...}
%============================================================
\usepackage{lastpage}
\usepackage{needspace}
%--- DÀN TRANG ĐỀ THI: tách câu (THÊM 05/11/2026; chỉ thêm lệnh mới, KHÔNG đổi hành vi mặc định) ---
%  Hai công tắc, mặc định tắt (= preamble chạy y như cũ). Bật bằng cách đặt TRƯỚC \input{preamble-eduzy-v2}
%  trong file đề (tools/dan_trang/dan_trang.py tự ghi khi đặt \NeedspaceCau):
%    \def\TachCauDS{1}   câu Đúng/Sai được tách: câu dẫn (kèm hình) ở cuối trang, bảng Phát biểu sang trang sau
%                        (câu dẫn và hình không tách nhau, bảng không tách giữa hai trang)
%    \def\TachCauLG{1}   bản lời giải (\CheDo khác dethi): phần đề của câu cao hơn 1/3 trang được phép tách
%  Lệnh dùng trong file đề (giá trị do tools/dan_trang/dan_trang.py đo ở bản dethi):
%    \NeedspaceCau[<cao câu dẫn + hình>]{<cao cả khối câu>}   ngay trên \begin{ex} (câu đầu phần: trên tiêu đề phần)
%        công tắc tắt: đúng bằng \Needspace{<cao cả khối câu>} như cũ;
%        \TachCauDS=1 và có [..] (câu Đúng/Sai được chọn tách): chỉ giữ câu dẫn + hình;
%        \TachCauLG=1 ở bản lời giải: đòi tối đa 1/3 \textheight.
%    \NeedspaceBang{<cao bảng Phát biểu>}%   ngay trước \choiceTFt: \TachCauDS=1 thì giữ trọn bảng trên một trang;
%        công tắc tắt: không làm gì.
\providecommand{\TachCauDS}{0}
\providecommand{\TachCauLG}{0}
\newlength{\EZnsCau}
\newcommand{\NeedspaceCau}[2][]{%
	\setlength{\EZnsCau}{#2}%
	\ifnum\TachCauDS=1 \if\relax\detokenize{#1}\relax\else\setlength{\EZnsCau}{#1}\fi\fi
	\ifnum\TachCauLG=1 \ifdefstring{\kindOPTN}{dethi}{}{\ifdim\EZnsCau>0.3333\textheight\setlength{\EZnsCau}{0.3333\textheight}\fi}\fi
	\Needspace{\EZnsCau}}
\newcommand{\NeedspaceBang}[1]{\ifnum\TachCauDS=1 \Needspace{#1}\fi}

%--- Thang màu: đậm -> nhạt
\definecolor{gold1}{RGB}{126,96,0}    % nâu đồng đậm  #7E6000
\definecolor{gold2}{RGB}{192,144,3}   % vàng kim      #C09003
\definecolor{gold3}{RGB}{255,215,99}  % vàng nhạt     #FFD763
\definecolor{gold4}{RGB}{255,231,153} % vàng rất nhạt #FFE799
\definecolor{gold5}{RGB}{255,243,205} % vàng kem      #FFF3CD
%--- Ghi đè tông xanh/đỏ cũ
\definecolor{myblue}{RGB}{126,96,0}
\colorlet{blue}{gold1}
\colorlet{cyan}{gold2}
\colorlet{teal}{gold2}

%--- Tên chương mặc định (file đề nên tự khai báo lại)
\providecommand{\tenchuong}{}
%--- Chân trang
\providecommand{\thuonghieu}{TOÁN 12 -- Eduzy.com.vn}
\providecommand{\tacgia}{Tô Gia Bảo -- 0353703783}

%--- Chữ cạnh logo ở đầu trang
\providecommand{\khauhieu}{Eduzy.com.vn}

%====================== THÔNG SỐ ĐẦU TRANG, CHÂN TRANG (v2) ======================
%  Dải nền đậm tràn ngang trang; ô chữ vàng kem có cạnh xiên, kèm 2 vạch xiên chuyển màu.
%  Đầu trang lẻ : logo + \khauhieu (trái)   | ô chữ \tenchuong (phải)
%  Chân trang lẻ: ô chữ \thuonghieu (trái)  | Trang (số) (phải)
%  Trang chẵn soi gương trang lẻ; chân trang chẵn ghi \tacgia.
%--- Màu: chỉ dùng thang vàng kim Eduzy
\colorlet{ezNen}{gold1}% dải nền đậm
\colorlet{ezVach1}{gold3}% vạch xiên sát dải nền (đậm hơn)
\colorlet{ezVach2}{gold4}% vạch xiên sát ô chữ (nhạt hơn)
\colorlet{ezO}{gold5}% ô chữ vàng kem
%--- Kích thước
\newlength{\EZdai}\setlength{\EZdai}{1.1cm}% chiều cao dải nền
\newlength{\EZo}\setlength{\EZo}{0.68cm}% chiều cao ô chữ
\newlength{\EZxien}\setlength{\EZxien}{0.55cm}% độ lệch ngang của cạnh xiên
\newlength{\EZvach}\setlength{\EZvach}{0.75cm}% bề rộng mỗi vạch xiên
\newlength{\EZkhe}\setlength{\EZkhe}{0.12cm}% khe trắng giữa vạch và ô chữ
\newlength{\EZdem}\setlength{\EZdem}{0.4cm}% khoảng đệm từ chữ tới cạnh xiên
\newlength{\EZle}\setlength{\EZle}{1.8cm}% lề trái/phải của chữ (bằng lề thân bài)
\newlength{\EZrongmin}\setlength{\EZrongmin}{4.5cm}% ô chữ rộng tối thiểu
\newlength{\EZrong}% bề rộng chữ trong ô (tính tự động)
%--- Font
\newcommand{\EZfontO}{\fontsize{10.5pt}{12pt}\selectfont\bfseries}% chữ trong ô
\newcommand{\EZfontKH}{\fontfamily{qag}\fontsize{11pt}{12pt}\selectfont\bfseries}% \khauhieu cạnh logo
\newcommand{\EZfontTrang}{\fontsize{11.5pt}{12pt}\selectfont\bfseries}% chữ "Trang" và số trang
%--- Đo bề rộng chữ trong ô
\newcommand{\EZdorong}[1]{%
	\settowidth{\EZrong}{\EZfontO #1}%
	\ifdim\EZrong<\EZrongmin\setlength{\EZrong}{\EZrongmin}\fi}
%--- Logo + chữ trắng \khauhieu (dựng sẵn một lần khi vào \begin{document})
\newsavebox{\EZlogobox}
\AtBeginDocument{\sbox{\EZlogobox}{%
		\raisebox{-0.5\height}{\eduzylogo[1cm]}\hspace{0.15cm}%
		\raisebox{-0.5\height}{\color{white}\EZfontKH\khauhieu}}}
%--- Hai vạch xiên + khe trắng. #1: điểm gốc ở mép dưới, #2: hướng (1 sang phải, -1 sang trái),
%    #3: hướng nghiêng (1 kiểu "/", -1 kiểu "\")
\newcommand{\EZvachxien}[3]{%
	\fill[white] (#1) --++({(#3)*\EZxien},\EZo) --++({(#2)*\EZkhe},0) --++({-(#3)*\EZxien},-\EZo) --cycle;
	\fill[ezVach2] ([xshift=(#2)*\EZkhe]#1) --++({(#3)*\EZxien},\EZo) --++({(#2)*\EZvach},0) --++({-(#3)*\EZxien},-\EZo) --cycle;
	\fill[ezVach1] ([xshift=(#2)*(\EZkhe+\EZvach)]#1) --++({(#3)*\EZxien},\EZo) --++({(#2)*\EZvach},0) --++({-(#3)*\EZxien},-\EZo) --cycle;
	\draw[white,line width=0.6pt] ([xshift=(#2)*(\EZkhe+\EZvach)]#1) --++({(#3)*\EZxien},\EZo);
}
%--- Chữ trong ô (căn giữa theo chiều cao ô, không lệch vì dấu tiếng Việt)
\tikzset{ezchu/.style={inner sep=0pt,text=ezNen,text height=0.72em,text depth=0.22em}}

%====================== LOGO CHÌM GIỮA TRANG ======================
%  Vẽ cùng đầu trang nên nằm dưới chữ của thân bài, có trên mọi trang.
%  Tắt: \logochimtat (từ trang đó trở đi), bật lại: \logochimbat.
%  Đổi cỡ, độ đậm: khai báo trước \input, ví dụ \def\EZchimcao{9cm} \def\EZchimdo{0.1}
\providecommand{\EZchimcao}{10cm}% chiều cao logo chìm
\providecommand{\EZchimdo}{0.12}% độ đậm (0 = trong suốt, 1 = rõ hẳn)
\newif\ifEZchim \EZchimtrue
\newcommand{\logochimtat}{\EZchimfalse}
\newcommand{\logochimbat}{\EZchimtrue}
%  Dựng sẵn logo mờ một lần khi vào \begin{document}; transparency group làm mờ cả logo
%  một lượt nên không lộ vệt ở chỗ các dải màu chồng nhau.
\newsavebox{\EZchimtho}
\newsavebox{\EZchimbox}
\AtBeginDocument{%
	\sbox{\EZchimtho}{\eduzylogo[\EZchimcao]}%
	\sbox{\EZchimbox}{\begin{tikzpicture}
			\begin{scope}[transparency group,opacity=\EZchimdo]
				\node[inner sep=0pt] {\usebox{\EZchimtho}};
			\end{scope}
	\end{tikzpicture}}}
\newcommand{\EZvechim}{%
	\ifEZchim
		\node[inner sep=0pt] at (current page.center) {\usebox{\EZchimbox}};
	\fi}

%====================== ĐẦU TRANG - TRANG LẺ ======================
%  (logo bên trái, ô chữ \tenchuong bên phải)
\newcommand{\myfancyheadOdd}{%
	\EZdorong{\tenchuong}%
	\begin{tikzpicture}[remember picture,overlay]
		\coordinate (T) at (current page.north west);
		\coordinate (P) at (current page.north east);
		%--- logo chìm giữa trang
		\EZvechim
		%--- dải nền
		\fill[ezNen] (T) rectangle ([yshift=-\EZdai]P);
		%--- ô chữ: X là góc dưới bên trái, cạnh trái xiên kiểu "/"
		\coordinate (X) at ([xshift=-\EZle-\EZrong-\EZdem-\EZxien,yshift=-\EZdai]P);
		\fill[ezO] (X) --++(\EZxien,\EZo) -- ([yshift=-\EZdai+\EZo]P) -- ([yshift=-\EZdai]P) --cycle;
		\EZvachxien{X}{-1}{1}
		\node[ezchu,anchor=east] at ([xshift=-\EZle,yshift=-\EZdai+0.5*\EZo]P) {\EZfontO\tenchuong};
		%--- logo + tên
		\node[inner sep=0pt,anchor=west] at ([xshift=\EZle,yshift=-0.5*\EZdai]T) {\usebox{\EZlogobox}};
	\end{tikzpicture}%
}

%====================== ĐẦU TRANG - TRANG CHẴN =====================
%  (soi gương: ô chữ \tenchuong bên trái, logo bên phải)
\newcommand{\myfancyheadEven}{%
	\EZdorong{\tenchuong}%
	\begin{tikzpicture}[remember picture,overlay]
		\coordinate (T) at (current page.north west);
		\coordinate (P) at (current page.north east);
		%--- logo chìm giữa trang
		\EZvechim
		%--- dải nền
		\fill[ezNen] (T) rectangle ([yshift=-\EZdai]P);
		%--- ô chữ: X là góc dưới bên phải, cạnh phải xiên kiểu "\"
		\coordinate (X) at ([xshift=\EZle+\EZrong+\EZdem+\EZxien,yshift=-\EZdai]T);
		\fill[ezO] (X) --++(-\EZxien,\EZo) -- ([yshift=-\EZdai+\EZo]T) -- ([yshift=-\EZdai]T) --cycle;
		\EZvachxien{X}{1}{-1}
		\node[ezchu,anchor=west] at ([xshift=\EZle,yshift=-\EZdai+0.5*\EZo]T) {\EZfontO\tenchuong};
		%--- logo + tên
		\node[inner sep=0pt,anchor=east] at ([xshift=-\EZle,yshift=-0.5*\EZdai]P) {\usebox{\EZlogobox}};
	\end{tikzpicture}%
}

%====================== CHÂN TRANG - TRANG LẺ ======================
%  (ô chữ \thuonghieu bên trái, "Trang (số)" bên phải)
\newcommand{\myfancyfootOdd}{%
	\EZdorong{\thuonghieu}%
	\begin{tikzpicture}[remember picture,overlay]
		\coordinate (T) at (current page.south west);
		\coordinate (P) at (current page.south east);
		%--- dải nền
		\fill[ezNen] (T) rectangle ([yshift=\EZdai]P);
		%--- ô chữ sát mép trên dải: X là góc dưới bên phải, cạnh phải xiên kiểu "/"
		\coordinate (X) at ([xshift=\EZle+\EZrong+\EZdem,yshift=\EZdai-\EZo]T);
		\fill[ezO] ([yshift=\EZdai-\EZo]T) -- (X) --++(\EZxien,\EZo) -- ([yshift=\EZdai]T) --cycle;
		\EZvachxien{X}{1}{1}
		\node[ezchu,anchor=west] at ([xshift=\EZle,yshift=\EZdai-0.5*\EZo]T) {\EZfontO\thuonghieu};
		%--- số trang
		\node[circle,draw=white,line width=0.9pt,text=white,inner sep=0pt,minimum size=0.62cm,anchor=east]
		(so) at ([xshift=-\EZle,yshift=0.5*\EZdai]P) {\EZfontTrang\thepage};
		\node[inner sep=0pt,text=white,anchor=base east] (tr) at ([xshift=-0.15cm]so.base west) {\EZfontTrang Trang};
	\end{tikzpicture}%
}

%====================== CHÂN TRANG - TRANG CHẴN ====================
%  (soi gương: "Trang (số)" bên trái, ô chữ \tacgia bên phải)
\newcommand{\myfancyfootEven}{%
	\EZdorong{\tacgia}%
	\begin{tikzpicture}[remember picture,overlay]
		\coordinate (T) at (current page.south west);
		\coordinate (P) at (current page.south east);
		%--- dải nền
		\fill[ezNen] (T) rectangle ([yshift=\EZdai]P);
		%--- ô chữ sát mép trên dải: X là góc dưới bên trái, cạnh trái xiên kiểu "\"
		\coordinate (X) at ([xshift=-\EZle-\EZrong-\EZdem,yshift=\EZdai-\EZo]P);
		\fill[ezO] ([yshift=\EZdai-\EZo]P) -- (X) --++(-\EZxien,\EZo) -- ([yshift=\EZdai]P) --cycle;
		\EZvachxien{X}{-1}{-1}
		\node[ezchu,anchor=east] at ([xshift=-\EZle,yshift=\EZdai-0.5*\EZo]P) {\EZfontO\tacgia};
		%--- số trang
		%    (vòng tròn ẩn so0 nằm giữa dải để lấy đường chân chữ, giống trang lẻ)
		\node[circle,inner sep=0pt,minimum size=0.62cm] (so0) at ([yshift=0.5*\EZdai]T) {\EZfontTrang\phantom{\thepage}};
		\node[inner sep=0pt,text=white,anchor=base west] (tr) at ($(so0.base -| T)+(\EZle,0)$) {\EZfontTrang Trang};
		\node[circle,draw=white,line width=0.9pt,text=white,inner sep=0pt,minimum size=0.62cm,anchor=base west]
		(so) at ([xshift=0.15cm]tr.base east) {\EZfontTrang\thepage};
	\end{tikzpicture}%
}

%--- Giữ tương thích với các file cũ đang gọi \myfancyhead, \myfancyfoot
\newcommand{\myfancyhead}{\myfancyheadOdd}
\newcommand{\myfancyfoot}{\myfancyfootOdd}

%======================= GẮN VÀO TRANG =============================
\pagestyle{fancy}
\fancyhf{}
\fancyhead[LO]{\myfancyheadOdd}
\fancyhead[RE]{\myfancyheadEven}
\fancyfoot[LO]{\myfancyfootOdd}
\fancyfoot[RE]{\myfancyfootEven}
\renewcommand{\headrulewidth}{0pt}
\renewcommand{\footrulewidth}{0pt}
\fancypagestyle{plain}{%
	\fancyhf{}
	\fancyhead[LO]{\myfancyheadOdd}
	\fancyhead[RE]{\myfancyheadEven}
	\fancyfoot[LO]{\myfancyfootOdd}
	\fancyfoot[RE]{\myfancyfootEven}
	\renewcommand{\headrulewidth}{0pt}
	\renewcommand{\footrulewidth}{0pt}
}
%=================== HẾT PHẦN GỘP TỪ eduzy-de.tex ===================
\newcommand{\indam}[1]{
	{\fontsize{12pt}{9pt}\sffamily \selectfont\bfseries\color{gold1} #1}
}
\newcommand{\ind}[1]{
	{\fontsize{14pt}{9pt}\fontfamily{qag} \selectfont\bfseries\color{gold1} #1}
}
\renewcommand{\FalseEX}{\stepcounter{dapan}\textbf{\selectfont\color{gold1}\Alph{dapan}.}\ignorespaces}
\renewtheorem{ex}{\color{gold1}Câu}
\renewcommand{\baselinestretch}{1.1}
\newcommand{\hoac}[1]{
	\left[\begin{aligned}#1\end{aligned}\right.}
\newcommand{\heva}[1]{
	\left\{\begin{aligned}#1\end{aligned}\right.}
\usepackage{esvect}
\def\vec{\vv}
\def\vt{\vv}
\def\overrightarrow{\vv}
%====================== LOGO EDUZY VẼ BẰNG TIKZ ======================
%  Đường viền dò từ logo gốc; dải màu vàng kim đo trực tiếp trên ảnh gốc
%  (nghiêng 21 độ, tô bằng 80 dải hẹp nên hiển thị giống nhau trên mọi trình xem PDF).
%  Dùng: \eduzylogo          -> cao 2.4cm (mặc định)
%        \eduzylogo[1.5cm]   -> cao 1.5cm
\definecolor{ezText}{RGB}{245,200,46}
\newcommand{\eduzylogo}[1][2.4cm]{%
	\resizebox{!}{#1}{%
		\begin{tikzpicture}[x=0.01cm,y=0.01cm]
			\path[use as bounding box] (0,0) rectangle (224,-224);
			% mũ tốt nghiệp + vương miện
			\begin{scope}
				\clip
				(119.4,-176.1)..controls(118.9,-176.1)and(116.8,-176)..(114.6,-176)
				..controls(101.4,-176)and(85.8,-175.9)..(83.7,-175.7)
				..controls(81.8,-175.6)and(81,-175.6)..(80.2,-175.7)
				..controls(78.8,-175.9)and(77.7,-175.9)..(76.2,-175.7)
				..controls(74.9,-175.5)and(70.8,-175.5)..(64.3,-175.7)
				..controls(59.1,-175.9)and(58.3,-175.7)..(58.6,-173.7)
				..controls(58.6,-172.8)and(59.6,-169.6)..(60.3,-168)
				..controls(60.8,-166.9)and(61.4,-165.8)..(62.3,-165)
				..controls(63.1,-164.2)and(64.3,-163.4)..(65,-163.2)
				..controls(65.5,-163.1)and(72.3,-163)..(76.6,-163.1)
				..controls(77.9,-163.1)and(80.1,-163.1)..(81.4,-163)
				..controls(82.8,-162.9)and(86.3,-162.9)..(89.2,-162.9)
				..controls(93,-162.9)and(94.9,-162.8)..(96,-162.7)
				..controls(97.2,-162.5)and(101,-162.5)..(114,-162.4)
				..controls(123,-162.3)and(134.9,-162.2)..(140.2,-162.1)
				..controls(150.1,-162)and(152.2,-162.1)..(153.6,-162.4)
				..controls(154.4,-162.6)and(157.1,-163.8)..(158.1,-164.4)
				..controls(158.8,-164.8)and(159.7,-165.7)..(159.9,-166.3)
				..controls(160,-166.5)and(160.2,-167.1)..(160.2,-167.6)
				..controls(160.3,-168)and(160.6,-169.4)..(160.9,-170.6)
				..controls(161.2,-171.7)and(161.5,-173.1)..(161.6,-173.6)
				..controls(161.8,-174.4)and(161.8,-174.5)..(161.5,-174.9)
				..controls(161.2,-175.6)and(160.8,-175.9)..(159.9,-176)
				..controls(159.2,-176.1)and(154.1,-176.2)..(151.5,-176.2)
				..controls(148.1,-176.1)and(122.6,-176.1)..(121.6,-176.2)
				..controls(120.8,-176.2)and(119.8,-176.2)..(119.4,-176.1)--cycle
				(130.4,-160)..controls(129,-160)and(119.8,-160)..(110.1,-159.9)
				..controls(100.3,-159.8)and(88.5,-159.7)..(83.9,-159.6)
				..controls(77.8,-159.5)and(75.3,-159.5)..(74.5,-159.6)
				..controls(73.9,-159.7)and(72.5,-159.8)..(71.3,-159.8)
				..controls(69.6,-159.8)and(69,-159.8)..(68.5,-159.6)
				..controls(67.6,-159.3)and(66.8,-158.6)..(66.2,-157.6)
				..controls(66,-157.3)and(65.6,-156.5)..(65.3,-156.1)
				..controls(64.6,-155.1)and(64.5,-154.9)..(63.7,-152.9)
				..controls(63.2,-151.6)and(62.9,-150.9)..(61.8,-149.6)
				..controls(61.5,-149.2)and(61.1,-148.5)..(60.9,-147.9)
				..controls(60.7,-147.4)and(60.3,-146.6)..(60.1,-146.1)
				..controls(59.9,-145.5)and(59.5,-144.9)..(59,-144.3)
				..controls(58.5,-143.7)and(58.2,-143.2)..(57.9,-142.5)
				..controls(57.7,-141.9)and(57.3,-141.1)..(57,-140.7)
				..controls(56.6,-140.3)and(56.1,-139.3)..(55.8,-138.5)
				..controls(55.5,-137.8)and(55.1,-136.8)..(54.9,-136.5)
				..controls(54.7,-136.1)and(54.3,-135.2)..(53.9,-134.5)
				..controls(53.6,-133.8)and(53,-132.7)..(52.6,-132)
				..controls(51.6,-130.5)and(51.4,-130.2)..(51.1,-129)
				..controls(50.6,-127.5)and(50.6,-127.4)..(49.7,-126.3)
				..controls(49.2,-125.7)and(48.6,-124.8)..(48.3,-124.2)
				..controls(48,-123.6)and(47.5,-122.7)..(47.3,-122.2)
				..controls(47,-121.8)and(46.5,-120.7)..(46.1,-119.8)
				..controls(45.8,-118.9)and(45.2,-117.8)..(44.9,-117.2)
				..controls(44.2,-116.1)and(43.3,-114.4)..(42.3,-112.2)
				..controls(41.8,-111.3)and(41.3,-110.4)..(41.1,-110.1)
				..controls(40.5,-109.4)and(40.3,-109)..(40,-107.8)
				..controls(39.8,-107)and(39.6,-106.5)..(39.2,-106)
				..controls(38.9,-105.6)and(38.4,-104.7)..(38,-104)
				..controls(37.7,-103.4)and(37.2,-102.5)..(37,-102.1)
				..controls(36.7,-101.7)and(36.3,-100.8)..(36,-100.1)
				..controls(35.6,-98.9)and(35.3,-98.4)..(34.6,-97.6)
				..controls(34.2,-97.1)and(34.2,-96.3)..(34.6,-96)
				..controls(35,-95.7)and(35.5,-95.9)..(36.2,-96.7)
				..controls(36.5,-97)and(37.2,-97.5)..(37.6,-97.7)..controls(38,-98)and(38.6,-98.4)..(39.1,-98.9)
				..controls(39.6,-99.4)and(40.3,-100)..(40.6,-100.3)
				..controls(41.3,-100.8)and(42.4,-101.8)..(43.7,-102.9)
				..controls(44,-103.2)and(44.6,-103.7)..(45,-103.9)
				..controls(45.4,-104.2)and(45.9,-104.6)..(46.1,-104.8)
				..controls(47.1,-105.9)and(48.8,-107.1)..(50.5,-107.9)
				..controls(51.8,-108.6)and(52.3,-109)..(53.4,-110)
				..controls(54.8,-111.2)and(55.6,-111.6)..(57,-112.1)
				..controls(57.5,-112.3)and(58.5,-112.8)..(59.4,-113.3)
				..controls(60.2,-113.7)and(61.3,-114.3)..(61.9,-114.5)
				..controls(62.4,-114.7)and(63.3,-115.1)..(63.8,-115.3)
				..controls(64.4,-115.6)and(65.4,-115.9)..(66.2,-116.1)
				..controls(67,-116.3)and(68.4,-116.6)..(69.2,-116.9)
				..controls(72.3,-117.9)and(72.8,-118)..(78,-118.1)
				..controls(84,-118.3)and(86.1,-118.1)..(88.1,-117.1)
				..controls(89.4,-116.5)and(93,-115)..(93.7,-114.9)
				..controls(94.8,-114.6)and(95.4,-114.2)..(95.9,-113.4)
				..controls(96.2,-113)and(96.7,-112.5)..(97.1,-112.2)
				..controls(98,-111.6)and(101.3,-108.4)..(102.1,-107.2)
				..controls(102.4,-106.8)and(102.9,-106.2)..(103.2,-105.8)
				..controls(103.5,-105.5)and(104,-104.8)..(104.2,-104.2)
				..controls(104.5,-103.6)and(105,-102.8)..(105.3,-102.4)
				..controls(105.6,-102)and(106.1,-101.2)..(106.5,-100.6)
				..controls(107.1,-99.5)and(108.6,-97.5)..(109.5,-96.5)
				..controls(109.9,-96.1)and(110.1,-96)..(110.4,-96)
				..controls(111.1,-96)and(111.6,-96.5)..(111.8,-97.6)
				..controls(112,-98.6)and(112.3,-99.3)..(112.8,-99.9)
				..controls(113.7,-100.9)and(115.2,-103.1)..(115.7,-104.2)
				..controls(116.1,-105)and(116.5,-105.5)..(116.9,-106)
				..controls(117.3,-106.4)and(117.8,-106.9)..(118,-107.3)
				..controls(118.2,-107.6)and(118.6,-108)..(119,-108.3)
				..controls(119.6,-108.8)and(121.1,-110.3)..(121.8,-111.2)
				..controls(122.1,-111.7)and(122.5,-112)..(123.4,-112.4)
				..controls(124.8,-113.2)and(125.8,-113.8)..(127,-114.6)
				..controls(127.5,-114.9)and(128.1,-115.3)..(128.2,-115.4)
				..controls(128.4,-115.5)and(129.3,-115.8)..(130.2,-116)
				..controls(131.1,-116.2)and(132.5,-116.6)..(133.3,-116.9)
				..controls(135.7,-117.9)and(136.1,-117.9)..(142.5,-117.8)
				..controls(145.8,-117.8)and(148.4,-117.7)..(148.8,-117.6)
				..controls(149.2,-117.6)and(150.2,-117.2)..(151,-116.8)
				..controls(152,-116.4)and(152.8,-116)..(153.6,-115.9)
				..controls(154.5,-115.7)and(155.2,-115.4)..(156.3,-114.9)
				..controls(157.1,-114.5)and(158.2,-114)..(158.7,-113.8)
				..controls(159.1,-113.7)and(159.8,-113.3)..(160.2,-113)
				..controls(160.6,-112.7)and(161.3,-112.4)..(161.7,-112.2)
				..controls(162.7,-111.9)and(163.2,-111.7)..(164.1,-111.1)
				..controls(164.5,-110.8)and(165.3,-110.4)..(165.9,-110)
				..controls(166.5,-109.7)and(167.3,-109.2)..(167.6,-108.9)
				..controls(167.9,-108.6)and(168.6,-108.2)..(169,-108)
				..controls(169.4,-107.8)and(170.2,-107.3)..(170.7,-106.9)
				..controls(171.2,-106.5)and(172.1,-105.8)..(172.7,-105.4)
				..controls(173.2,-105)and(174.1,-104.4)..(174.5,-104)
				..controls(174.9,-103.7)and(175.7,-103.1)..(176.2,-102.6)
				..controls(176.7,-102.2)and(177.4,-101.5)..(177.9,-101.2)
				..controls(178.3,-100.8)and(178.8,-100.4)..(179,-100.2)
				..controls(179.2,-100)and(179.6,-99.6)..(180,-99.3)
				..controls(180.4,-99.1)and(180.9,-98.6)..(181.2,-98.4)
				..controls(181.5,-98.1)and(182.1,-97.7)..(182.5,-97.5)
				..controls(183.2,-97.2)and(183.5,-96.9)..(184.3,-95.9)
				..controls(184.6,-95.6)and(184.6,-95.6)..(184.8,-95.7)
				..controls(185,-95.8)and(185.1,-96)..(185.1,-96.2)
				..controls(185.1,-96.5)and(185,-96.7)..(184.6,-97.2)
				..controls(184.1,-97.8)and(184,-97.9)..(183.9,-98.8)
				..controls(183.7,-100.3)and(183.3,-101.3)..(182.9,-101.8)
				..controls(182.6,-102.1)and(182.3,-102.6)..(182.2,-102.9)
				..controls(182,-103.3)and(181.6,-104)..(181.2,-104.6)
				..controls(180.9,-105.2)and(180.5,-106.1)..(180.4,-106.6)
				..controls(179.9,-108)and(179.7,-108.5)..(179,-109.4)
				..controls(178.6,-109.8)and(178.2,-110.5)..(177.9,-111)
				..controls(177.7,-111.5)and(177.3,-112.3)..(177,-112.7)
				..controls(176.6,-113.2)and(176.4,-113.7)..(176.3,-114.2)
				..controls(176.2,-114.7)and(175.9,-115.2)..(175.4,-116)
				..controls(175,-116.6)and(174.4,-117.7)..(174,-118.5)
				..controls(173.7,-119.2)and(173.3,-120.1)..(173,-120.5)
				..controls(172.8,-120.9)and(172.4,-121.8)..(172.1,-122.5)
				..controls(171.8,-123.2)and(171.3,-124.1)..(171,-124.6)
				..controls(170.8,-125.1)and(170.3,-126)..(169.9,-126.6)
				..controls(169.6,-127.3)and(169.2,-128.1)..(169,-128.3)
				..controls(168.8,-128.6)and(168.5,-129.3)..(168.3,-129.8)
				..controls(168.1,-130.3)and(167.8,-131.1)..(167.5,-131.6)
				..controls(166.9,-132.5)and(164.2,-137.9)..(164.2,-138.1)
				..controls(164.2,-138.1)and(163.9,-138.8)..(163.4,-139.7)
				..controls(162.9,-140.5)and(162.4,-141.5)..(162.3,-141.9)
				..controls(161.9,-143)and(161.6,-143.4)..(160.9,-144.1)
				..controls(160.4,-144.6)and(160.3,-144.8)..(160.1,-145.4)
				..controls(160,-145.8)and(159.6,-146.8)..(159.2,-147.8)
				..controls(158.7,-148.8)and(158.1,-150)..(157.9,-150.6)
				..controls(157.6,-151.1)and(157.2,-151.9)..(156.9,-152.4)
				..controls(156.5,-152.8)and(156.1,-153.7)..(155.8,-154.4)
				..controls(155.6,-155)and(155.2,-155.9)..(154.9,-156.3)
				..controls(154.7,-156.7)and(154.4,-157.4)..(154.2,-157.9)
				..controls(154.1,-158.4)and(153.8,-159)..(153.5,-159.3)
				..controls(152.8,-160.1)and(153,-160.1)..(142.4,-160.1)
				..controls(137.2,-160.1)and(131.8,-160.1)..(130.4,-160)--cycle
				(141.5,-103.7)..controls(140.8,-103)and(139.8,-102.1)..(139.2,-101.6)
				..controls(138.5,-101.2)and(137.9,-100.6)..(137.8,-100.4)
				..controls(137,-99.5)and(134,-96.5)..(133.2,-96)
				..controls(132.7,-95.7)and(131.8,-94.9)..(131.1,-94.3)
				..controls(129.4,-92.7)and(127.9,-91.6)..(126.6,-91.1)
				..controls(125.2,-90.5)and(124.1,-90)..(122.9,-89.1)
				..controls(122.1,-88.6)and(121.3,-88.2)..(119.9,-87.7)
				..controls(118.8,-87.3)and(117.5,-86.8)..(116.9,-86.5)
				..controls(114.3,-85.5)and(113.7,-85.4)..(109.7,-85.5)
				..controls(105.3,-85.7)and(103.9,-85.9)..(101.6,-87)
				..controls(100.7,-87.4)and(99.4,-87.8)..(97.9,-88)
				..controls(97,-88.2)and(96.6,-88.3)..(95.8,-88.9)
				..controls(95.4,-89.2)and(94.7,-89.6)..(94.2,-89.8)
				..controls(92.9,-90.3)and(92.3,-90.7)..(90.9,-92.1)
				..controls(90.2,-92.7)and(89.4,-93.3)..(89,-93.5)..controls(87,-94.5)and(85.2,-96)..(83.6,-97.9)
				..controls(83.3,-98.2)and(82.5,-98.9)..(81.8,-99.4)
				..controls(80.7,-100.1)and(80.5,-100.3)..(80,-101.1)
				..controls(78.7,-103.1)and(77.2,-103.9)..(76.4,-103)
				..controls(75.7,-102.2)and(75.7,-103.2)..(75.7,-83.6)
				..controls(75.6,-73.6)and(75.6,-64.8)..(75.5,-64)
				..controls(75.5,-63.2)and(75.5,-61.4)..(75.6,-59.8)--(75.7,-57.1)--(76,-56.8)
				..controls(76.8,-56.2)and(78,-56.3)..(79.8,-57.2)
				..controls(80.4,-57.5)and(81.3,-57.8)..(81.8,-57.9)
				..controls(82.3,-58.1)and(83,-58.4)..(83.6,-58.8)..controls(84.3,-59.2)and(85,-59.5)..(86,-59.8)
				..controls(86.8,-60.1)and(88,-60.6)..(88.7,-60.9)
				..controls(89.3,-61.2)and(90.6,-61.7)..(91.4,-62)
				..controls(92.2,-62.3)and(93.7,-62.9)..(94.6,-63.3)
				..controls(95.6,-63.7)and(96.8,-64.1)..(97.4,-64.3)
				..controls(98,-64.5)and(99,-64.9)..(99.7,-65.2)
				..controls(100.4,-65.6)and(101.5,-66)..(102.3,-66.2)
				..controls(104.2,-66.8)and(105,-67.1)..(106.2,-67.8)
				..controls(107.9,-68.7)and(110.2,-69.4)..(111.3,-69.2)
				..controls(112.5,-69)and(115,-68.2)..(115.6,-67.8)
				..controls(117.3,-66.7)and(117.6,-66.6)..(119.5,-66.2)
				..controls(121.9,-65.8)and(122.1,-65.7)..(123,-65.1)
				..controls(124.1,-64.4)and(124.4,-64.3)..(126,-63.9)
				..controls(127,-63.7)and(127.7,-63.4)..(128.4,-63.1)
				..controls(129,-62.8)and(129.7,-62.5)..(130.6,-62.4)
				..controls(132.1,-62.1)and(134,-61.5)..(135.2,-61)
				..controls(135.7,-60.8)and(136.7,-60.5)..(137.3,-60.3)
				..controls(138.4,-60)and(139.3,-59.7)..(141.1,-59)
				..controls(141.6,-58.7)and(142,-58.6)..(142.5,-58.6)
				..controls(143.2,-58.6)and(143.3,-58.7)..(143.6,-59)
				..controls(144,-59.5)and(144.1,-59.9)..(144,-60.8)
				..controls(144,-61.2)and(144,-71.1)..(144,-82.8)--(143.9,-104.1)--(143.7,-104.5)
				..controls(143.5,-104.7)and(143.3,-104.8)..(143,-104.9)
				..controls(142.6,-104.9)and(142.5,-104.8)..(141.5,-103.7)--cycle
				(27.8,-91.4)..controls(27.1,-90.9)and(26.3,-90)..(26.1,-89.5)
				..controls(26,-89.2)and(25.8,-88.6)..(25.7,-87.9)
				..controls(25.4,-86.2)and(25.6,-84.8)..(26.5,-83.5)
				..controls(27.4,-82.3)and(28.3,-82)..(29.4,-82.3)
				..controls(31.1,-82.8)and(32.1,-84.9)..(32,-87.4)
				..controls(31.9,-89.3)and(31.2,-90.6)..(29.9,-91.4)
				..controls(29.1,-91.9)and(28.6,-91.9)..(27.8,-91.4)--cycle
				(190.1,-91.3)..controls(189.7,-91.2)and(189.3,-90.9)..(189.1,-90.7)
				..controls(188.3,-89.8)and(187.6,-87.7)..(187.6,-86.2)
				..controls(187.6,-85.5)and(187.7,-85.2)..(188.2,-84.4)
				..controls(189.1,-82.6)and(190.1,-81.8)..(191.3,-82.1)
				..controls(192,-82.3)and(193.5,-83.7)..(193.8,-84.4)
				..controls(194.1,-85.1)and(194.2,-87.1)..(194,-88.3)
				..controls(193.8,-89.8)and(193.3,-90.5)..(192.3,-91.3)
				..controls(191.7,-91.7)and(190.9,-91.7)..(190.1,-91.3)--cycle
				(173.7,-83.4)..controls(172.2,-82.5)and(172,-82.5)..(170.6,-82.4)
				..controls(169.7,-82.4)and(168.8,-82.4)..(168.3,-82.5)
				..controls(167.2,-82.7)and(166.9,-82.7)..(166.6,-82.2)
				..controls(166.3,-81.8)and(166.3,-81)..(166.8,-80)
				..controls(167.2,-78.9)and(167.4,-77.4)..(167.7,-74.1)
				..controls(167.8,-72.8)and(167.9,-71.8)..(167.9,-71.7)
				..controls(168,-71.7)and(168,-70.2)..(168,-68.4)..controls(168,-65.4)and(168,-64.9)..(167.7,-63)
				..controls(167.3,-60.5)and(167.1,-59.9)..(166.4,-58.4)
				..controls(166.1,-57.7)and(165.8,-56.9)..(165.7,-56.6)
				..controls(165.7,-56.2)and(165.4,-55.6)..(165.2,-55.2)
				..controls(164.6,-53.8)and(161.9,-50.7)..(161,-50.3)
				..controls(159.8,-49.7)and(157.9,-50)..(155.2,-51.2)
				..controls(154.7,-51.4)and(153.4,-51.8)..(152.2,-52.1)
				..controls(150.8,-52.4)and(149.4,-52.9)..(148.6,-53.2)
				..controls(147.8,-53.6)and(146.7,-54)..(146,-54.2)
				..controls(145.4,-54.4)and(144.5,-54.7)..(144,-54.9)
				..controls(143.6,-55.2)and(142.7,-55.5)..(141.8,-55.7)
				..controls(141,-55.9)and(139.6,-56.3)..(138.7,-56.6)
				..controls(137.8,-57)and(136.3,-57.4)..(135.3,-57.6)
				..controls(133.8,-57.9)and(133.2,-58.1)..(132.1,-58.7)
				..controls(131.1,-59.2)and(130.3,-59.4)..(129.4,-59.6)
				..controls(128.7,-59.8)and(127.8,-60)..(127.4,-60.1)
				..controls(127,-60.3)and(126,-60.5)..(125.1,-60.7)
				..controls(123.2,-61.1)and(121.7,-61.5)..(120.2,-62.1)
				..controls(119.7,-62.3)and(118.1,-62.8)..(116.7,-63.2)
				..controls(114.6,-63.9)and(114,-64.1)..(113,-64.7)
				..controls(111.9,-65.4)and(111.8,-65.4)..(111,-65.5)
				..controls(109.9,-65.5)and(109.3,-65.4)..(107.7,-64.7)
				..controls(107,-64.4)and(106,-64)..(105.3,-63.9)
				..controls(103.5,-63.5)and(102.2,-63)..(101.2,-62.5)
				..controls(100.1,-61.8)and(99.3,-61.5)..(97.8,-61.1)
				..controls(97.2,-61)and(96.2,-60.7)..(95.7,-60.4)
				..controls(95.1,-60.1)and(94.2,-59.8)..(93.5,-59.6)
				..controls(92.9,-59.4)and(91.7,-59)..(91,-58.7)..controls(90.2,-58.4)and(88.9,-58)..(88,-57.7)
				..controls(87.2,-57.5)and(86.1,-57.1)..(85.5,-56.8)
				..controls(85,-56.6)and(84.1,-56.2)..(83.5,-56)..controls(82.9,-55.7)and(81.8,-55.2)..(81,-54.9)
				..controls(80.1,-54.5)and(79.1,-54.2)..(78.7,-54.1)
				..controls(77.9,-53.8)and(77.3,-53.6)..(76.2,-52.9)
				..controls(75.7,-52.6)and(75.2,-52.4)..(74.1,-52.2)
				..controls(72.5,-51.9)and(72.4,-51.8)..(71.3,-51.1)
				..controls(69.9,-50.1)and(69.4,-49.9)..(67.9,-49.8)
				..controls(67,-49.7)and(66.6,-49.6)..(66.3,-49.4)
				..controls(66.1,-49.3)and(65.7,-49)..(65.4,-48.8)
				..controls(65.1,-48.7)and(64.2,-48.4)..(63.6,-48.2)
				..controls(62.9,-48)and(61.7,-47.6)..(60.9,-47.2)
				..controls(60.1,-46.8)and(58.9,-46.4)..(58.2,-46.2)
				..controls(57.6,-46)and(56.8,-45.7)..(56.5,-45.5)
				..controls(54.9,-44.6)and(54.6,-44.5)..(53.2,-44.2)
				..controls(51.7,-43.8)and(51,-43.6)..(49.9,-42.9)
				..controls(49.5,-42.6)and(48.9,-42.4)..(47.5,-42)
				..controls(45.7,-41.6)and(45.6,-41.5)..(44.3,-40.6)
				..controls(43.5,-40.1)and(42.8,-39.6)..(42.7,-39.4)
				..controls(42.4,-38.7)and(43,-38.1)..(44.5,-37.6)..controls(44.9,-37.5)and(46,-37)..(46.9,-36.6)
				..controls(47.9,-36.1)and(48.8,-35.8)..(49.6,-35.6)
				..controls(50.5,-35.4)and(51.3,-35.2)..(52.4,-34.6)
				..controls(53.3,-34.1)and(54.5,-33.7)..(54.9,-33.5)
				..controls(55.4,-33.4)and(56.4,-33)..(57.2,-32.6)
				..controls(58.6,-32)and(60.1,-31.5)..(61.6,-31.2)
				..controls(62.3,-31.1)and(62.7,-31)..(63.2,-30.7)
				..controls(64.1,-30)and(64.3,-29.9)..(65.3,-29.7)
				..controls(65.7,-29.6)and(66.5,-29.2)..(67.1,-28.9)
				..controls(67.6,-28.6)and(68.8,-28.2)..(69.6,-27.9)
				..controls(70.5,-27.7)and(71.6,-27.3)..(72.1,-27)
				..controls(72.6,-26.7)and(73.8,-26.2)..(74.8,-25.9)
				..controls(75.7,-25.5)and(76.7,-25.1)..(77.1,-25)
				..controls(77.4,-24.8)and(78.5,-24.4)..(79.6,-24.1)
				..controls(80.6,-23.8)and(82,-23.2)..(82.8,-22.9)
				..controls(83.5,-22.5)and(84.8,-22)..(85.8,-21.8)
				..controls(86.7,-21.5)and(87.8,-21.1)..(88.1,-20.9)
				..controls(88.4,-20.7)and(89.3,-20.3)..(90.1,-20.1)
				..controls(90.8,-19.8)and(91.9,-19.4)..(92.4,-19.1)
				..controls(93.5,-18.5)and(94.4,-18.2)..(95.4,-18.1)
				..controls(95.7,-18)and(96.4,-17.8)..(96.9,-17.5)
				..controls(97.5,-17.3)and(98.4,-16.9)..(99,-16.6)
				..controls(99.7,-16.3)and(100.9,-15.9)..(101.9,-15.6)
				..controls(102.9,-15.3)and(104.4,-14.7)..(105.3,-14.3)
				..controls(108.1,-13.2)and(108.5,-13.1)..(110.1,-13.1)
				..controls(111.4,-13.1)and(111.5,-13.1)..(112.4,-13.5)
				..controls(113.2,-14)and(115,-14.8)..(116.6,-15.4)
				..controls(116.9,-15.5)and(117.9,-15.8)..(118.8,-16.1)
				..controls(120.9,-16.6)and(121.5,-16.9)..(122.6,-17.9)
				..controls(123,-18.3)and(123.7,-18.8)..(123.9,-18.9)
				..controls(124.6,-19.4)and(126,-19.8)..(127.9,-20.2)
				..controls(129.6,-20.5)and(129.8,-20.6)..(131.1,-21.4)
				..controls(131.5,-21.7)and(132.4,-22.1)..(133.1,-22.3)
				..controls(136.1,-23.2)and(138.3,-24.1)..(140.2,-25)
				..controls(141,-25.5)and(141.5,-25.6)..(142.5,-25.8)
				..controls(145.1,-26.3)and(146,-26.6)..(149.4,-28.1)
				..controls(151.3,-28.9)and(153.5,-29.8)..(154.3,-30.1)
				..controls(155.1,-30.4)and(156.2,-30.8)..(156.6,-31)
				..controls(157.1,-31.2)and(158.2,-31.6)..(159.1,-31.8)
				..controls(160.6,-32.2)and(160.9,-32.3)..(161.7,-32.9)
				..controls(162.4,-33.3)and(162.9,-33.6)..(163.7,-33.8)
				..controls(164.3,-34)and(164.8,-34.1)..(164.9,-34.1)
				..controls(165,-34.1)and(165.6,-34.4)..(166.4,-34.8)
				..controls(167.3,-35.2)and(168.1,-35.4)..(169,-35.6)
				..controls(170,-35.8)and(170.5,-36)..(171.7,-36.6)
				..controls(172.4,-37)and(173.3,-37.4)..(173.6,-37.4)
				..controls(175.6,-37.7)and(176.2,-38)..(176.7,-39.1)
				..controls(176.9,-39.6)and(176.9,-39.7)..(176.7,-40.4)
				..controls(176.4,-41.6)and(175.8,-42.7)..(175.2,-43.4)
				..controls(174.5,-44.1)and(172.6,-45.5)..(171.7,-45.8)
				..controls(170.6,-46.2)and(169.1,-46)..(168.2,-45.4)
				..controls(167.9,-45.2)and(167.3,-44.8)..(166.9,-44.4)
				..controls(166.4,-44.1)and(165.8,-43.6)..(165.4,-43.2)
				..controls(165.1,-42.9)and(164.4,-42.3)..(163.9,-41.9)
				..controls(163.3,-41.5)and(162.7,-40.9)..(162.4,-40.5)
				..controls(161.4,-39.2)and(159.9,-38)..(158.8,-37.8)
				..controls(158.4,-37.8)and(158.3,-37.7)..(158.2,-37.2)
				..controls(158,-36.8)and(157.4,-36.1)..(157.1,-36.1)
				..controls(157,-36.1)and(156.7,-36.1)..(156.5,-36)
				..controls(156.1,-35.9)and(155.9,-35.7)..(155.7,-35.2)
				..controls(155.3,-34.6)and(155.1,-34.4)..(154,-34.1)
				..controls(153.6,-34)and(153.2,-33.7)..(152.9,-33.5)
				..controls(152.4,-32.9)and(151.6,-32.4)..(151.3,-32.4)
				..controls(150.9,-32.4)and(150.4,-32.6)..(150.4,-32.9)
				..controls(150.4,-33.2)and(150.8,-33.4)..(151.4,-33.6)
				..controls(152.1,-33.7)and(152.8,-34.1)..(153.5,-34.9)
				..controls(154.1,-35.5)and(154.3,-35.6)..(155,-35.8)
				..controls(155.7,-36)and(155.7,-36)..(156,-36.6)
				..controls(156.3,-37.3)and(156.7,-37.7)..(157.5,-37.8)
				..controls(158.1,-37.9)and(158.2,-38)..(158.4,-38.8)
				..controls(158.5,-39.4)and(158.6,-39.6)..(159.6,-40.4)
				..controls(160.2,-41)and(160.9,-41.6)..(161.3,-41.8)
				..controls(161.7,-42)and(161.9,-42.3)..(162.2,-42.8)
				..controls(162.4,-43.2)and(162.8,-43.8)..(163.1,-44.1)
				..controls(163.9,-44.8)and(165.1,-46.7)..(165.7,-48)
				..controls(166.2,-49.3)and(166.4,-49.7)..(167.2,-50.5)
				..controls(167.5,-50.8)and(167.9,-51.4)..(168.1,-51.8)
				..controls(168.3,-52.2)and(168.8,-53)..(169.1,-53.6)
				..controls(169.6,-54.3)and(169.8,-54.9)..(169.9,-55.6)
				..controls(170.4,-57.8)and(170.5,-58.1)..(171.2,-59.5)
				..controls(171.6,-60.3)and(172,-61.5)..(172.2,-62.1)
				..controls(172.3,-62.8)and(172.7,-63.7)..(172.9,-64.1)
				..controls(173.2,-64.6)and(173.4,-65.3)..(173.4,-65.6)
				..controls(173.8,-68.8)and(174.2,-70.6)..(174.9,-72.6)
				..controls(175.5,-74.5)and(175.9,-77.1)..(176.1,-80.2)
				..controls(176.3,-83.4)and(175.5,-84.4)..(173.7,-83.4)--cycle;
				\definecolor{ezs}{RGB}{247,179,26}\fill[ezs](-5.7,-156.6)--(-3.4,-157.5)--(61.9,12.7)--(59.7,13.6)--cycle;
				\definecolor{ezs}{RGB}{247,179,22}\fill[ezs](-3.8,-157.3)--(-1.5,-158.2)--(63.9,12)--(61.6,12.9)--cycle;
				\definecolor{ezs}{RGB}{248,179,19}\fill[ezs](-1.8,-158.1)--(0.4,-158.9)--(65.8,11.2)--(63.5,12.1)--cycle;
				\definecolor{ezs}{RGB}{248,178,16}\fill[ezs](0.1,-158.8)--(2.4,-159.7)--(67.7,10.5)--(65.4,11.4)--cycle;
				\definecolor{ezs}{RGB}{249,178,15}\fill[ezs](2,-159.5)--(4.3,-160.4)--(69.6,9.8)--(67.3,10.7)--cycle;
				\definecolor{ezs}{RGB}{251,177,17}\fill[ezs](3.9,-160.3)--(6.2,-161.1)--(71.5,9)--(69.2,9.9)--cycle;
				\definecolor{ezs}{RGB}{252,179,18}\fill[ezs](5.8,-161)--(8.1,-161.9)--(73.4,8.3)--(71.1,9.2)--cycle;
				\definecolor{ezs}{RGB}{252,182,17}\fill[ezs](7.7,-161.7)--(10,-162.6)--(75.3,7.6)--(73,8.5)--cycle;
				\definecolor{ezs}{RGB}{252,185,17}\fill[ezs](9.6,-162.5)--(11.9,-163.4)--(77.2,6.8)--(75,7.7)--cycle;
				\definecolor{ezs}{RGB}{251,189,18}\fill[ezs](11.5,-163.2)--(13.8,-164.1)--(79.2,6.1)--(76.9,7)--cycle;
				\definecolor{ezs}{RGB}{251,191,22}\fill[ezs](13.4,-163.9)--(15.7,-164.8)--(81.1,5.4)--(78.8,6.2)--cycle;
				\definecolor{ezs}{RGB}{252,193,25}\fill[ezs](15.4,-164.7)--(17.6,-165.6)--(83,4.6)--(80.7,5.5)--cycle;
				\definecolor{ezs}{RGB}{253,195,27}\fill[ezs](17.3,-165.4)--(19.6,-166.3)--(84.9,3.9)--(82.6,4.8)--cycle;
				\definecolor{ezs}{RGB}{253,196,31}\fill[ezs](19.2,-166.1)--(21.5,-167)--(86.8,3.2)--(84.5,4)--cycle;
				\definecolor{ezs}{RGB}{254,199,35}\fill[ezs](21.1,-166.9)--(23.4,-167.8)--(88.7,2.4)--(86.4,3.3)--cycle;
				\definecolor{ezs}{RGB}{253,201,38}\fill[ezs](23,-167.6)--(25.3,-168.5)--(90.6,1.7)--(88.3,2.6)--cycle;
				\definecolor{ezs}{RGB}{253,203,40}\fill[ezs](24.9,-168.3)--(27.2,-169.2)--(92.5,1)--(90.3,1.8)--cycle;
				\definecolor{ezs}{RGB}{253,205,42}\fill[ezs](26.8,-169.1)--(29.1,-170)--(94.4,0.2)--(92.2,1.1)--cycle;
				\definecolor{ezs}{RGB}{253,207,45}\fill[ezs](28.7,-169.8)--(31,-170.7)--(96.4,-0.5)--(94.1,0.4)--cycle;
				\definecolor{ezs}{RGB}{253,210,49}\fill[ezs](30.7,-170.5)--(32.9,-171.4)--(98.3,-1.2)--(96,-0.4)--cycle;
				\definecolor{ezs}{RGB}{253,212,54}\fill[ezs](32.6,-171.3)--(34.9,-172.2)--(100.2,-2)--(97.9,-1.1)--cycle;
				\definecolor{ezs}{RGB}{253,214,59}\fill[ezs](34.5,-172)--(36.8,-172.9)--(102.1,-2.7)--(99.8,-1.8)--cycle;
				\definecolor{ezs}{RGB}{253,216,63}\fill[ezs](36.4,-172.7)--(38.7,-173.6)--(104,-3.4)--(101.7,-2.6)--cycle;
				\definecolor{ezs}{RGB}{253,218,65}\fill[ezs](38.3,-173.5)--(40.6,-174.4)--(105.9,-4.2)--(103.6,-3.3)--cycle;
				\definecolor{ezs}{RGB}{254,220,66}\fill[ezs](40.2,-174.2)--(42.5,-175.1)--(107.8,-4.9)--(105.5,-4)--cycle;
				\definecolor{ezs}{RGB}{254,222,68}\fill[ezs](42.1,-175)--(44.4,-175.8)--(109.7,-5.6)--(107.5,-4.8)--cycle;
				\definecolor{ezs}{RGB}{254,224,71}\fill[ezs](44,-175.7)--(46.3,-176.6)--(111.7,-6.4)--(109.4,-5.5)--cycle;
				\definecolor{ezs}{RGB}{254,226,75}\fill[ezs](46,-176.4)--(48.2,-177.3)--(113.6,-7.1)--(111.3,-6.2)--cycle;
				\definecolor{ezs}{RGB}{254,228,79}\fill[ezs](47.9,-177.2)--(50.1,-178)--(115.5,-7.8)--(113.2,-7)--cycle;
				\definecolor{ezs}{RGB}{254,230,82}\fill[ezs](49.8,-177.9)--(52.1,-178.8)--(117.4,-8.6)--(115.1,-7.7)--cycle;
				\definecolor{ezs}{RGB}{254,232,85}\fill[ezs](51.7,-178.6)--(54,-179.5)--(119.3,-9.3)--(117,-8.4)--cycle;
				\definecolor{ezs}{RGB}{254,234,88}\fill[ezs](53.6,-179.4)--(55.9,-180.2)--(121.2,-10)--(118.9,-9.2)--cycle;
				\definecolor{ezs}{RGB}{254,237,91}\fill[ezs](55.5,-180.1)--(57.8,-181)--(123.1,-10.8)--(120.8,-9.9)--cycle;
				\definecolor{ezs}{RGB}{254,239,94}\fill[ezs](57.4,-180.8)--(59.7,-181.7)--(125,-11.5)--(122.8,-10.6)--cycle;
				\definecolor{ezs}{RGB}{254,241,97}\fill[ezs](59.3,-181.6)--(61.6,-182.4)--(127,-12.2)--(124.7,-11.4)--cycle;
				\definecolor{ezs}{RGB}{254,243,101}\fill[ezs](61.2,-182.3)--(63.5,-183.2)--(128.9,-13)--(126.6,-12.1)--cycle;
				\definecolor{ezs}{RGB}{254,245,103}\fill[ezs](63.2,-183)--(65.4,-183.9)--(130.8,-13.7)--(128.5,-12.8)--cycle;
				\definecolor{ezs}{RGB}{254,245,105}\fill[ezs](65.1,-183.8)--(67.4,-184.6)--(132.7,-14.4)--(130.4,-13.6)--cycle;
				\definecolor{ezs}{RGB}{254,245,105}\fill[ezs](67,-184.5)--(69.3,-185.4)--(134.6,-15.2)--(132.3,-14.3)--cycle;
				\definecolor{ezs}{RGB}{254,244,103}\fill[ezs](68.9,-185.2)--(71.2,-186.1)--(136.5,-15.9)--(134.2,-15)--cycle;
				\definecolor{ezs}{RGB}{253,243,101}\fill[ezs](70.8,-186)--(73.1,-186.8)--(138.4,-16.6)--(136.1,-15.8)--cycle;
				\definecolor{ezs}{RGB}{253,241,98}\fill[ezs](72.7,-186.7)--(75,-187.6)--(140.3,-17.4)--(138,-16.5)--cycle;
				\definecolor{ezs}{RGB}{253,239,95}\fill[ezs](74.6,-187.4)--(76.9,-188.3)--(142.2,-18.1)--(140,-17.2)--cycle;
				\definecolor{ezs}{RGB}{253,236,92}\fill[ezs](76.5,-188.2)--(78.8,-189)--(144.2,-18.8)--(141.9,-18)--cycle;
				\definecolor{ezs}{RGB}{253,234,88}\fill[ezs](78.5,-188.9)--(80.7,-189.8)--(146.1,-19.6)--(143.8,-18.7)--cycle;
				\definecolor{ezs}{RGB}{253,232,85}\fill[ezs](80.4,-189.6)--(82.7,-190.5)--(148,-20.3)--(145.7,-19.4)--cycle;
				\definecolor{ezs}{RGB}{253,230,82}\fill[ezs](82.3,-190.4)--(84.6,-191.2)--(149.9,-21)--(147.6,-20.2)--cycle;
				\definecolor{ezs}{RGB}{253,228,78}\fill[ezs](84.2,-191.1)--(86.5,-192)--(151.8,-21.8)--(149.5,-20.9)--cycle;
				\definecolor{ezs}{RGB}{253,226,74}\fill[ezs](86.1,-191.8)--(88.4,-192.7)--(153.7,-22.5)--(151.4,-21.6)--cycle;
				\definecolor{ezs}{RGB}{253,224,70}\fill[ezs](88,-192.6)--(90.3,-193.4)--(155.6,-23.3)--(153.3,-22.4)--cycle;
				\definecolor{ezs}{RGB}{253,222,67}\fill[ezs](89.9,-193.3)--(92.2,-194.2)--(157.5,-24)--(155.3,-23.1)--cycle;
				\definecolor{ezs}{RGB}{253,220,65}\fill[ezs](91.8,-194)--(94.1,-194.9)--(159.5,-24.7)--(157.2,-23.8)--cycle;
				\definecolor{ezs}{RGB}{253,217,63}\fill[ezs](93.7,-194.8)--(96,-195.6)--(161.4,-25.5)--(159.1,-24.6)--cycle;
				\definecolor{ezs}{RGB}{253,215,60}\fill[ezs](95.7,-195.5)--(97.9,-196.4)--(163.3,-26.2)--(161,-25.3)--cycle;
				\definecolor{ezs}{RGB}{253,213,56}\fill[ezs](97.6,-196.2)--(99.9,-197.1)--(165.2,-26.9)--(162.9,-26)--cycle;
				\definecolor{ezs}{RGB}{252,211,53}\fill[ezs](99.5,-197)--(101.8,-197.8)--(167.1,-27.7)--(164.8,-26.8)--cycle;
				\definecolor{ezs}{RGB}{252,209,49}\fill[ezs](101.4,-197.7)--(103.7,-198.6)--(169,-28.4)--(166.7,-27.5)--cycle;
				\definecolor{ezs}{RGB}{252,207,45}\fill[ezs](103.3,-198.4)--(105.6,-199.3)--(170.9,-29.1)--(168.6,-28.2)--cycle;
				\definecolor{ezs}{RGB}{253,205,42}\fill[ezs](105.2,-199.2)--(107.5,-200)--(172.8,-29.9)--(170.6,-29)--cycle;
				\definecolor{ezs}{RGB}{253,202,40}\fill[ezs](107.1,-199.9)--(109.4,-200.8)--(174.7,-30.6)--(172.5,-29.7)--cycle;
				\definecolor{ezs}{RGB}{253,200,38}\fill[ezs](109,-200.6)--(111.3,-201.5)--(176.7,-31.3)--(174.4,-30.4)--cycle;
				\definecolor{ezs}{RGB}{253,198,35}\fill[ezs](111,-201.4)--(113.2,-202.2)--(178.6,-32.1)--(176.3,-31.2)--cycle;
				\definecolor{ezs}{RGB}{253,196,31}\fill[ezs](112.9,-202.1)--(115.2,-203)--(180.5,-32.8)--(178.2,-31.9)--cycle;
				\definecolor{ezs}{RGB}{253,194,28}\fill[ezs](114.8,-202.8)--(117.1,-203.7)--(182.4,-33.5)--(180.1,-32.6)--cycle;
				\definecolor{ezs}{RGB}{253,192,26}\fill[ezs](116.7,-203.6)--(119,-204.5)--(184.3,-34.3)--(182,-33.4)--cycle;
				\definecolor{ezs}{RGB}{253,189,24}\fill[ezs](118.6,-204.3)--(120.9,-205.2)--(186.2,-35)--(183.9,-34.1)--cycle;
				\definecolor{ezs}{RGB}{253,186,20}\fill[ezs](120.5,-205)--(122.8,-205.9)--(188.1,-35.7)--(185.8,-34.8)--cycle;
				\definecolor{ezs}{RGB}{253,184,17}\fill[ezs](122.4,-205.8)--(124.7,-206.7)--(190,-36.5)--(187.8,-35.6)--cycle;
				\definecolor{ezs}{RGB}{252,183,15}\fill[ezs](124.3,-206.5)--(126.6,-207.4)--(192,-37.2)--(189.7,-36.3)--cycle;
				\definecolor{ezs}{RGB}{252,181,14}\fill[ezs](126.3,-207.2)--(128.5,-208.1)--(193.9,-37.9)--(191.6,-37.1)--cycle;
				\definecolor{ezs}{RGB}{252,179,13}\fill[ezs](128.2,-208)--(130.4,-208.9)--(195.8,-38.7)--(193.5,-37.8)--cycle;
				\definecolor{ezs}{RGB}{252,178,12}\fill[ezs](130.1,-208.7)--(132.4,-209.6)--(197.7,-39.4)--(195.4,-38.5)--cycle;
				\definecolor{ezs}{RGB}{251,178,10}\fill[ezs](132,-209.4)--(134.3,-210.3)--(199.6,-40.1)--(197.3,-39.3)--cycle;
				\definecolor{ezs}{RGB}{251,178,10}\fill[ezs](133.9,-210.2)--(136.2,-211.1)--(201.5,-40.9)--(199.2,-40)--cycle;
				\definecolor{ezs}{RGB}{250,177,11}\fill[ezs](135.8,-210.9)--(138.1,-211.8)--(203.4,-41.6)--(201.1,-40.7)--cycle;
				\definecolor{ezs}{RGB}{250,178,11}\fill[ezs](137.7,-211.6)--(140,-212.5)--(205.3,-42.3)--(203.1,-41.5)--cycle;
				\definecolor{ezs}{RGB}{248,180,10}\fill[ezs](139.6,-212.4)--(141.9,-213.3)--(207.3,-43.1)--(205,-42.2)--cycle;
				\definecolor{ezs}{RGB}{247,181,9}\fill[ezs](141.5,-213.1)--(143.8,-214)--(209.2,-43.8)--(206.9,-42.9)--cycle;
				\definecolor{ezs}{RGB}{246,182,9}\fill[ezs](143.5,-213.8)--(145.7,-214.7)--(211.1,-44.5)--(208.8,-43.7)--cycle;
				\definecolor{ezs}{RGB}{246,182,10}\fill[ezs](145.4,-214.6)--(147.7,-215.5)--(213,-45.3)--(210.7,-44.4)--cycle;
			\end{scope}
			% chữ EDUZY (Times đậm, làm dày nét bằng 5 lớp lệch nhẹ; căn đúng khung chữ logo gốc)
			\foreach \dx/\dy in {-0.7/0,0.7/0,0/-0.3,0/0.3,0/0}
			\node[anchor=center,inner sep=0pt,outer sep=0pt,text=ezText,align=center] at (109.6+\dx,-193.6+\dy)
			{\resizebox{1.01cm}{0.22cm}{\fontencoding{T1}\fontfamily{ptm}\fontseries{b}\selectfont EDUZY}};
		\end{tikzpicture}%
	}%
}
%====================== THIẾT LẬP THÂN ĐỀ ======================
%--- Môi trường tự luận {bt}: ex_test 3.3.4 chưa khai báo sẵn nên tạo ở đây
%    (ex_test đã tự móc \fixshowans và \listenumerate cho bt)
\makeatletter
\@ifundefined{bt}{\newtheorem{bt}{\color{gold1}\namebt}}{}
\makeatother
\renewcommand{\sectionmark}[1]{\markright{\thesection\ #1}}
%--- Khung tiêu đề (logo + tối đa 4 dòng; dòng 3, 4 để trống {} thì bỏ qua)
%  Đề thi:    \khungtieude{THI THỬ THPT QUỐC GIA NĂM 2026}{Môn Toán -- Đề số 21}
%                         {(Đề thi gồm: ...)}{Thời gian làm bài: 90 phút (không kể thời gian phát đề).}
%  Chuyên đề: \khungtieude{CHUYÊN ĐỀ 3 -- DÃY SỐ VÀ CÔNG THỨC TRUY HỒI}
%                         {Bồi dưỡng học sinh giỏi Toán THPT}{...}{}
\newcommand{\khungtieude}[4]{%
	\begin{tcolorbox}[colframe=gold1,colback=gold5!45,boxrule=0.5mm]
		\begin{minipage}[c]{3cm}
			\centerline{\eduzylogo[2.4cm]}
		\end{minipage}\hspace{0.5cm}
		\begin{minipage}[c]{12.5cm}
			\vskip 0.15 cm
			\centerline{\ind{#1}}
			\vskip 0.15 cm
			\centerline{\textbf{#2}}
			\if\relax\detokenize{#3}\relax\else
				\vskip 0.15 cm
				\centerline{\textit{#3}}
			\fi
			\if\relax\detokenize{#4}\relax\else
				\vskip 0.15 cm
				\centerline{\textit{#4}}
			\fi
			\vskip 0.15 cm
		\end{minipage}
	\end{tcolorbox}
}
%====================== BỔ SUNG CHO BÀI HỌC, CHUYÊN ĐỀ ======================
\usepackage[shortlabels]{enumitem}
\usepackage{tabularx}
%--- Danh sách gọn: \begin{enumerate}[gon] ... \end{enumerate}
\SetEnumitemKey{gon}{leftmargin=*,itemsep=1pt,parsep=0pt,topsep=3pt}
%--- Khung lý thuyết: \begin{khung}{} ... \end{khung}; tiêu đề chỉ dùng khi thật ngắn ("Mẹo", "Lỗi thường gặp")
\newtcolorbox{khung}[1]{breakable,enhanced,colframe=gold1,colback=gold5!45,
	boxrule=0.4mm,fonttitle=\bfseries\sffamily,coltitle=white,colbacktitle=gold1,
	attach boxed title to top left={xshift=6mm,yshift=-3mm},
	boxed title style={arc=1mm,boxrule=0.3mm},title={#1}}
%--- Tiêu đề mục nhỏ trong một dạng: \muc{Dấu hiệu}, \muc{Phương pháp}, \muc{Bài tự luyện}
\makeatletter
\newcommand{\muc}[1]{\par\medskip\noindent\indam{#1}\par\nobreak\smallskip\@afterindentfalse\@afterheading}
\makeatother
%--- Ô công thức cần nhớ (căn giữa, viền vàng, vừa khít nội dung)
%     \begin{center}\congthuc{$Q(n)=(n-1)!$}\end{center}
\newtcbox{\congthuc}{on line,colframe=gold1,colback=gold5!45,boxrule=0.5pt,arc=2pt,
	left=8pt,right=8pt,top=4pt,bottom=4pt}
%--- Công tắc lời giải bài tự luyện (ví dụ mẫu luôn dùng \loigiai)
%     \loigiaiTL[{đáp số}]{lời giải đầy đủ}
%     \TLfull -> in lời giải đầy đủ;  \TLdapso -> chỉ in một dòng "Đáp số";  \TLan -> ẩn hoàn toàn
%     (\CheDo là dethi thì mọi lời giải và đáp số đều ẩn)
\makeatletter
\newif\ifTLhien   \TLhientrue
\newif\ifTLdapso  \TLdapsofalse
\newcommand{\TLfull}{\TLhientrue\TLdapsofalse}
\newcommand{\TLdapso}{\TLhientrue\TLdapsotrue}
\newcommand{\TLan}{\TLhienfalse}
\newcommand{\loigiaiTL}[2][]{%
	\ifTLhien
		\ifTLdapso
			\def\@tlarg{#1}%
			\ifx\@tlarg\@empty\else
				\ifdefstring{\kindOPTN}{dethi}{}{\par\noindent{\bfseries\color{gold1}Đáp số:} #1\par}%
			\fi
		\else
			\loigiai{#2}%
		\fi
	\fi}
\makeatother
\TLdapso% mặc định: chỉ in đáp số bài tự luyện
%====================== BẢNG ĐÚNG/SAI (\choiceTFt) ======================
%  Hàng tiêu đề nền nâu đồng chữ trắng (Phát biểu | Đ | S), các hàng tô xen kẽ kem / trắng,
%  đường kẻ vàng kim, nhãn a) b) c) d) màu nâu đồng. Bản có đáp án: cột được chọn có vòng tròn
%  nâu đồng với dấu tích trắng (câu đúng tích ở cột Đúng, câu sai tích ở cột Sai).
%  Tiêu đề ghi đủ "Đúng | Sai": \OPTN{viettat=0} sau \input (cột Đúng/Sai rộng ra, phát biểu hẹp lại).
%  Dùng bảng gốc của ex_test: \def\BangDSgoc{1} trước \input{preamble-eduzy-v2}.
\providecommand{\BangDSgoc}{0}
\ifnum\BangDSgoc=0
\usepackage{colortbl,pifont}
\colorlet{ezDSke}{gold2!60!gold1}% đường kẻ bảng (nét 0,4pt như ex_test để không đổi chiều cao bảng)
\colorlet{ezDSle}{gold5!60}% hàng a), c)
\colorlet{ezDSchan}{white}% hàng b), d)
%--- dấu chọn ở cột Đúng/Sai
%    (bề rộng 0, căn giữa ô: không làm đổi bề rộng cột của ex_test, không báo tràn lề)
\newcommand{\EZdauchon}{\makebox[0pt]{\tikz[baseline=-1ex]\node[circle,fill=gold1,inner sep=0pt,
	minimum size=1.2em,text=white,font=\small]{\ding{51}};}}
\def\tickT{\EZdauchon}\def\tickF{\EZdauchon}
\AtBeginDocument{\sbox{\ccTFTab}{\EZdauchon}}
%--- nhãn a) b) c) d) màu nâu đồng (ở bản lời giải, nhãn câu đúng vẫn đóng khung như ex_test)
\renewcommand{\FalseTF}{\stepcounter{dapan}\textbf{\color{gold1}\DapAnTF\sepTF}\ignorespaces}
%--- trước khi dựng bảng: kẻ màu vàng kim (bề rộng cột giữ nguyên như ex_test); sau bảng trả lại màu đen
\newcount\EZhangTF
\newcommand{\EZbatdauTF}{%
	\global\EZhangTF=0\relax
	\arrayrulecolor{ezDSke}}
\newcommand{\EZketthucTF}{\arrayrulecolor{black}}
%--- hàng tiêu đề nền nâu đồng, chữ trắng
\makeatletter
\newcommand{\EZo@TF}[1]{\centering\bfseries\color{white}#1}
\newcommand{\EZtieudeTF}{%
	\ifcase\cotDS
		\ifnum\cotSTT=1\relax
			\def\rowHeader{\hline\rowcolor{gold1}\EZo@TF{TT}&\EZo@TF{\phatbieu}\tabularnewline}%
		\else
			\def\rowHeader{\hline\rowcolor{gold1}\EZo@TF{\phatbieu}\tabularnewline}%
		\fi
	\or
		\ifnum\cotSTT=1\relax
			\def\rowHeader{\hline\rowcolor{gold1}\EZo@TF{TT}&\EZo@TF{\phatbieu}&\EZo@TF{\Dung}\tabularnewline}%
		\else
			\def\rowHeader{\hline\rowcolor{gold1}\EZo@TF{\phatbieu}&\EZo@TF{\Dung}\tabularnewline}%
		\fi
	\or
		\ifnum\cotSTT=1\relax
			\def\rowHeader{\hline\rowcolor{gold1}\EZo@TF{TT}&\EZo@TF{\phatbieu}&\EZo@TF{\Dung}&\EZo@TF{\Sai}\tabularnewline}%
		\else
			\def\rowHeader{\hline\rowcolor{gold1}\EZo@TF{\phatbieu}&\EZo@TF{\Dung}&\EZo@TF{\Sai}\tabularnewline}%
		\fi
	\fi}
%--- tô xen kẽ từng hàng (gọi ngay sau \hline đầu mỗi hàng)
\newcommand{\EZmauhangTF}{%
	\noalign{\global\advance\EZhangTF\@ne}%
	\ifodd\EZhangTF\rowcolor{ezDSle}\else\rowcolor{ezDSchan}\fi}
%--- gắn vào \choiceTFt / \writeChoiceTFt của ex_test
\xpretocmd{\choiceTFt}{\EZbatdauTF}{}{\typeout{[preamble-eduzy-v2] Không gắn được bảng Đúng/Sai (đầu bảng)}}
\xpatchcmd{\choiceTFt}{\ifnum\dongPB=0\def\rowHeader{}\fi}{\EZtieudeTF\ifnum\dongPB=0\def\rowHeader{}\fi}
	{}{\typeout{[preamble-eduzy-v2] Không gắn được bảng Đúng/Sai (hàng tiêu đề)}}
\xpatchcmd{\writeChoiceTFt}{\hline\xdef\dotEX{\sdotEX}}{\hline\EZmauhangTF\xdef\dotEX{\sdotEX}}
	{}{\typeout{[preamble-eduzy-v2] Không gắn được bảng Đúng/Sai (tô hàng)}}
\xpatchcmd{\writeChoiceTFt}{{\end{longtable}}}{{\end{longtable}\EZketthucTF}}
	{}{\typeout{[preamble-eduzy-v2] Không gắn được bảng Đúng/Sai (cuối longtable)}}
\xpatchcmd{\writeChoiceTFt}{{\end{xtabular}\end{center}\medskip\noindent}}{{\end{xtabular}\end{center}\EZketthucTF\medskip\noindent}}
	{}{\typeout{[preamble-eduzy-v2] Không gắn được bảng Đúng/Sai (cuối xtabular)}}
\fi
%=====================================================================
\renewcommand{\parallel}{\mathrel{/\!\!/}}
.
\providecommand{\BangDSgoc}{0}
\ifnum\BangDSgoc=0
\usepackage{colortbl,pifont}
\colorlet{ezDSke}{gold2!60!gold1}% đường kẻ bảng (nét 0,4pt như ex_test để không đổi chiều cao bảng)
\colorlet{ezDSle}{gold5!60}% hàng a), c)
\colorlet{ezDSchan}{white}% hàng b), d)
%--- dấu chọn ở cột Đúng/Sai
%    (bề rộng 0, căn giữa ô: không làm đổi bề rộng cột của ex_test, không báo tràn lề)
\newcommand{\EZdauchon}{\makebox[0pt]{\tikz[baseline=-1ex]\node[circle,fill=gold1,inner sep=0pt,
	minimum size=1.2em,text=white,font=\small]{\ding{51}};}}
\def\tickT{\EZdauchon}\def\tickF{\EZdauchon}
\AtBeginDocument{\sbox{\ccTFTab}{\EZdauchon}}
%--- nhãn a) b) c) d) màu nâu đồng (ở bản lời giải, nhãn câu đúng vẫn đóng khung như ex_test)
\renewcommand{\FalseTF}{\stepcounter{dapan}\textbf{\color{gold1}\DapAnTF\sepTF}\ignorespaces}
%--- trước khi dựng bảng: kẻ màu vàng kim (bề rộng cột giữ nguyên như ex_test); sau bảng trả lại màu đen
\newcount\EZhangTF
\newcommand{\EZbatdauTF}{%
	\global\EZhangTF=0\relax
	\arrayrulecolor{ezDSke}}
\newcommand{\EZketthucTF}{\arrayrulecolor{black}}
%--- hàng tiêu đề nền nâu đồng, chữ trắng
\makeatletter
\newcommand{\EZo@TF}[1]{\centering\bfseries\color{white}#1}
\newcommand{\EZtieudeTF}{%
	\ifcase\cotDS
		\ifnum\cotSTT=1\relax
			\def\rowHeader{\hline\rowcolor{gold1}\EZo@TF{TT}&\EZo@TF{\phatbieu}\tabularnewline}%
		\else
			\def\rowHeader{\hline\rowcolor{gold1}\EZo@TF{\phatbieu}\tabularnewline}%
		\fi
	\or
		\ifnum\cotSTT=1\relax
			\def\rowHeader{\hline\rowcolor{gold1}\EZo@TF{TT}&\EZo@TF{\phatbieu}&\EZo@TF{\Dung}\tabularnewline}%
		\else
			\def\rowHeader{\hline\rowcolor{gold1}\EZo@TF{\phatbieu}&\EZo@TF{\Dung}\tabularnewline}%
		\fi
	\or
		\ifnum\cotSTT=1\relax
			\def\rowHeader{\hline\rowcolor{gold1}\EZo@TF{TT}&\EZo@TF{\phatbieu}&\EZo@TF{\Dung}&\EZo@TF{\Sai}\tabularnewline}%
		\else
			\def\rowHeader{\hline\rowcolor{gold1}\EZo@TF{\phatbieu}&\EZo@TF{\Dung}&\EZo@TF{\Sai}\tabularnewline}%
		\fi
	\fi}
%--- tô xen kẽ từng hàng (gọi ngay sau \hline đầu mỗi hàng)
\newcommand{\EZmauhangTF}{%
	\noalign{\global\advance\EZhangTF\@ne}%
	\ifodd\EZhangTF\rowcolor{ezDSle}\else\rowcolor{ezDSchan}\fi}
%--- gắn vào \choiceTFt / \writeChoiceTFt của ex_test
\xpretocmd{\choiceTFt}{\EZbatdauTF}{}{\typeout{[preamble-eduzy-v2] Không gắn được bảng Đúng/Sai (đầu bảng)}}
\xpatchcmd{\choiceTFt}{\ifnum\dongPB=0\def\rowHeader{}\fi}{\EZtieudeTF\ifnum\dongPB=0\def\rowHeader{}\fi}
	{}{\typeout{[preamble-eduzy-v2] Không gắn được bảng Đúng/Sai (hàng tiêu đề)}}
\xpatchcmd{\writeChoiceTFt}{\hline\xdef\dotEX{\sdotEX}}{\hline\EZmauhangTF\xdef\dotEX{\sdotEX}}
	{}{\typeout{[preamble-eduzy-v2] Không gắn được bảng Đúng/Sai (tô hàng)}}
\xpatchcmd{\writeChoiceTFt}{{\end{longtable}}}{{\end{longtable}\EZketthucTF}}
	{}{\typeout{[preamble-eduzy-v2] Không gắn được bảng Đúng/Sai (cuối longtable)}}
\xpatchcmd{\writeChoiceTFt}{{\end{xtabular}\end{center}\medskip\noindent}}{{\end{xtabular}\end{center}\EZketthucTF\medskip\noindent}}
	{}{\typeout{[preamble-eduzy-v2] Không gắn được bảng Đúng/Sai (cuối xtabular)}}
\fi
%=====================================================================
\renewcommand{\parallel}{\mathrel{/\!\!/}}

%     %=====================================================================
% preamble-vatli-hinh.tex -- BỘ MACRO HÌNH VẬT LÍ (dùng chung cho mọi đề Vật lí)
% Gọi ngay sau preamble chính:
%     %============================================================
%  PREAMBLE EDUZY v2 - dùng chung cho đề thi, bài học, chuyên đề (gói ex_test)
%  Khác bản preamble-eduzy.tex ở đầu trang, chân trang: dải nền nâu đồng tràn ngang,
%  logo + Eduzy.com.vn, ô chữ vàng kem cạnh xiên, chân trang "Trang (số)";
%  logo Eduzy chìm giữa mọi trang (tắt: \logochimtat, bật lại: \logochimbat).
%  Cách dùng ở file chính (xem mau-chuan.tex, mau-chuyen-de.tex):
%      \def\CheDo{loigiai}          % loigiai | dethi | solcolor | color | book
%      \def\tenchuong{<dòng chữ trên đầu trang>}
%      \def\thuonghieu{<chữ ở chân trang lẻ>}   % không khai báo thì dùng mặc định
%      \def\tacgia{<chữ ở chân trang chẵn>}      % không khai báo thì dùng mặc định
%      \def\khauhieu{<chữ cạnh logo ở đầu trang>}  % tùy chọn, mặc định Eduzy.com.vn
%      \input{preamble-eduzy-v2}
%      \TLdapso                      % tài liệu: \TLfull | \TLdapso | \TLan
%      \begin{document} ... \end{document}
%  Logo Eduzy vẽ bằng TikZ ngay trong file (\eduzylogo), không cần file ảnh.
%  Biên dịch: pdfLaTeX (máy có gói vietnam) hoặc XeLaTeX.
%============================================================
\documentclass[12pt,a4paper,twoside]{article}%<-- twoside: bat dau trang/chan trang chan-le
\usepackage{iftex}
\usepackage{amsmath,amssymb,tikz,mathrsfs,tkz-tab}
\usepackage{graphicx}
\usepackage{pgfplots}% đồ thị hàm số
\pgfplotsset{compat=1.17}
%\usepackage{mathptmx}
\usepackage{mathpazo}
\ifXeTeX% XeLaTeX: font chữ có đủ dấu tiếng Việt, giữ nguyên font toán mathpazo
	\usepackage[no-math]{fontspec}
	\setmainfont{TeX Gyre Pagella}
	\newfontfamily\EZqag{TeX Gyre Adventor}[NFSSFamily=qag,Ligatures=TeX]
\fi
\usepackage[many]{tcolorbox}
\usetikzlibrary{math,through,calc,intersections,angles,quotes,shapes,shapes.geometric,arrows,patterns,matrix,chains,arrows.meta,decorations.shapes,decorations.fractals,decorations.markings,shadows,decorations.pathreplacing}
\usepackage[top=1.6cm, bottom=1.8cm, left=1.8cm, right=1.8cm]{geometry}% v2: top 1.6cm (bản cũ 1.5cm) để chữ không sát dải đầu trang
\providecommand{\CheDo}{loigiai}% loigiai | dethi | solcolor | color | book
\expandafter\usepackage\expandafter[\CheDo]{ex_test}
% \tenchuong: khai báo ở file chính (trước \input) hoặc \renewcommand sau \input
%============================================================
%  [GỘP TỪ eduzy-de.tex] EDUZY - MÀU NÂU ĐỒNG / VÀNG KIM
%  + ĐẦU TRANG, CHÂN TRANG KIỂU DẢI NGANG (v2), ĐỐI XỨNG TRANG CHẴN / TRANG LẺ.
%  (không cần file eduzy-de.tex riêng nữa)
%  Khối này phải đặt SAU \usepackage[...]{ex_test} và SAU \def\tenchuong{...}
%============================================================
\usepackage{lastpage}
\usepackage{needspace}
%--- DÀN TRANG ĐỀ THI: tách câu (THÊM 05/11/2026; chỉ thêm lệnh mới, KHÔNG đổi hành vi mặc định) ---
%  Hai công tắc, mặc định tắt (= preamble chạy y như cũ). Bật bằng cách đặt TRƯỚC \input{preamble-eduzy-v2}
%  trong file đề (tools/dan_trang/dan_trang.py tự ghi khi đặt \NeedspaceCau):
%    \def\TachCauDS{1}   câu Đúng/Sai được tách: câu dẫn (kèm hình) ở cuối trang, bảng Phát biểu sang trang sau
%                        (câu dẫn và hình không tách nhau, bảng không tách giữa hai trang)
%    \def\TachCauLG{1}   bản lời giải (\CheDo khác dethi): phần đề của câu cao hơn 1/3 trang được phép tách
%  Lệnh dùng trong file đề (giá trị do tools/dan_trang/dan_trang.py đo ở bản dethi):
%    \NeedspaceCau[<cao câu dẫn + hình>]{<cao cả khối câu>}   ngay trên \begin{ex} (câu đầu phần: trên tiêu đề phần)
%        công tắc tắt: đúng bằng \Needspace{<cao cả khối câu>} như cũ;
%        \TachCauDS=1 và có [..] (câu Đúng/Sai được chọn tách): chỉ giữ câu dẫn + hình;
%        \TachCauLG=1 ở bản lời giải: đòi tối đa 1/3 \textheight.
%    \NeedspaceBang{<cao bảng Phát biểu>}%   ngay trước \choiceTFt: \TachCauDS=1 thì giữ trọn bảng trên một trang;
%        công tắc tắt: không làm gì.
\providecommand{\TachCauDS}{0}
\providecommand{\TachCauLG}{0}
\newlength{\EZnsCau}
\newcommand{\NeedspaceCau}[2][]{%
	\setlength{\EZnsCau}{#2}%
	\ifnum\TachCauDS=1 \if\relax\detokenize{#1}\relax\else\setlength{\EZnsCau}{#1}\fi\fi
	\ifnum\TachCauLG=1 \ifdefstring{\kindOPTN}{dethi}{}{\ifdim\EZnsCau>0.3333\textheight\setlength{\EZnsCau}{0.3333\textheight}\fi}\fi
	\Needspace{\EZnsCau}}
\newcommand{\NeedspaceBang}[1]{\ifnum\TachCauDS=1 \Needspace{#1}\fi}

%--- Thang màu: đậm -> nhạt
\definecolor{gold1}{RGB}{126,96,0}    % nâu đồng đậm  #7E6000
\definecolor{gold2}{RGB}{192,144,3}   % vàng kim      #C09003
\definecolor{gold3}{RGB}{255,215,99}  % vàng nhạt     #FFD763
\definecolor{gold4}{RGB}{255,231,153} % vàng rất nhạt #FFE799
\definecolor{gold5}{RGB}{255,243,205} % vàng kem      #FFF3CD
%--- Ghi đè tông xanh/đỏ cũ
\definecolor{myblue}{RGB}{126,96,0}
\colorlet{blue}{gold1}
\colorlet{cyan}{gold2}
\colorlet{teal}{gold2}

%--- Tên chương mặc định (file đề nên tự khai báo lại)
\providecommand{\tenchuong}{}
%--- Chân trang
\providecommand{\thuonghieu}{TOÁN 12 -- Eduzy.com.vn}
\providecommand{\tacgia}{Tô Gia Bảo -- 0353703783}

%--- Chữ cạnh logo ở đầu trang
\providecommand{\khauhieu}{Eduzy.com.vn}

%====================== THÔNG SỐ ĐẦU TRANG, CHÂN TRANG (v2) ======================
%  Dải nền đậm tràn ngang trang; ô chữ vàng kem có cạnh xiên, kèm 2 vạch xiên chuyển màu.
%  Đầu trang lẻ : logo + \khauhieu (trái)   | ô chữ \tenchuong (phải)
%  Chân trang lẻ: ô chữ \thuonghieu (trái)  | Trang (số) (phải)
%  Trang chẵn soi gương trang lẻ; chân trang chẵn ghi \tacgia.
%--- Màu: chỉ dùng thang vàng kim Eduzy
\colorlet{ezNen}{gold1}% dải nền đậm
\colorlet{ezVach1}{gold3}% vạch xiên sát dải nền (đậm hơn)
\colorlet{ezVach2}{gold4}% vạch xiên sát ô chữ (nhạt hơn)
\colorlet{ezO}{gold5}% ô chữ vàng kem
%--- Kích thước
\newlength{\EZdai}\setlength{\EZdai}{1.1cm}% chiều cao dải nền
\newlength{\EZo}\setlength{\EZo}{0.68cm}% chiều cao ô chữ
\newlength{\EZxien}\setlength{\EZxien}{0.55cm}% độ lệch ngang của cạnh xiên
\newlength{\EZvach}\setlength{\EZvach}{0.75cm}% bề rộng mỗi vạch xiên
\newlength{\EZkhe}\setlength{\EZkhe}{0.12cm}% khe trắng giữa vạch và ô chữ
\newlength{\EZdem}\setlength{\EZdem}{0.4cm}% khoảng đệm từ chữ tới cạnh xiên
\newlength{\EZle}\setlength{\EZle}{1.8cm}% lề trái/phải của chữ (bằng lề thân bài)
\newlength{\EZrongmin}\setlength{\EZrongmin}{4.5cm}% ô chữ rộng tối thiểu
\newlength{\EZrong}% bề rộng chữ trong ô (tính tự động)
%--- Font
\newcommand{\EZfontO}{\fontsize{10.5pt}{12pt}\selectfont\bfseries}% chữ trong ô
\newcommand{\EZfontKH}{\fontfamily{qag}\fontsize{11pt}{12pt}\selectfont\bfseries}% \khauhieu cạnh logo
\newcommand{\EZfontTrang}{\fontsize{11.5pt}{12pt}\selectfont\bfseries}% chữ "Trang" và số trang
%--- Đo bề rộng chữ trong ô
\newcommand{\EZdorong}[1]{%
	\settowidth{\EZrong}{\EZfontO #1}%
	\ifdim\EZrong<\EZrongmin\setlength{\EZrong}{\EZrongmin}\fi}
%--- Logo + chữ trắng \khauhieu (dựng sẵn một lần khi vào \begin{document})
\newsavebox{\EZlogobox}
\AtBeginDocument{\sbox{\EZlogobox}{%
		\raisebox{-0.5\height}{\eduzylogo[1cm]}\hspace{0.15cm}%
		\raisebox{-0.5\height}{\color{white}\EZfontKH\khauhieu}}}
%--- Hai vạch xiên + khe trắng. #1: điểm gốc ở mép dưới, #2: hướng (1 sang phải, -1 sang trái),
%    #3: hướng nghiêng (1 kiểu "/", -1 kiểu "\")
\newcommand{\EZvachxien}[3]{%
	\fill[white] (#1) --++({(#3)*\EZxien},\EZo) --++({(#2)*\EZkhe},0) --++({-(#3)*\EZxien},-\EZo) --cycle;
	\fill[ezVach2] ([xshift=(#2)*\EZkhe]#1) --++({(#3)*\EZxien},\EZo) --++({(#2)*\EZvach},0) --++({-(#3)*\EZxien},-\EZo) --cycle;
	\fill[ezVach1] ([xshift=(#2)*(\EZkhe+\EZvach)]#1) --++({(#3)*\EZxien},\EZo) --++({(#2)*\EZvach},0) --++({-(#3)*\EZxien},-\EZo) --cycle;
	\draw[white,line width=0.6pt] ([xshift=(#2)*(\EZkhe+\EZvach)]#1) --++({(#3)*\EZxien},\EZo);
}
%--- Chữ trong ô (căn giữa theo chiều cao ô, không lệch vì dấu tiếng Việt)
\tikzset{ezchu/.style={inner sep=0pt,text=ezNen,text height=0.72em,text depth=0.22em}}

%====================== LOGO CHÌM GIỮA TRANG ======================
%  Vẽ cùng đầu trang nên nằm dưới chữ của thân bài, có trên mọi trang.
%  Tắt: \logochimtat (từ trang đó trở đi), bật lại: \logochimbat.
%  Đổi cỡ, độ đậm: khai báo trước \input, ví dụ \def\EZchimcao{9cm} \def\EZchimdo{0.1}
\providecommand{\EZchimcao}{10cm}% chiều cao logo chìm
\providecommand{\EZchimdo}{0.12}% độ đậm (0 = trong suốt, 1 = rõ hẳn)
\newif\ifEZchim \EZchimtrue
\newcommand{\logochimtat}{\EZchimfalse}
\newcommand{\logochimbat}{\EZchimtrue}
%  Dựng sẵn logo mờ một lần khi vào \begin{document}; transparency group làm mờ cả logo
%  một lượt nên không lộ vệt ở chỗ các dải màu chồng nhau.
\newsavebox{\EZchimtho}
\newsavebox{\EZchimbox}
\AtBeginDocument{%
	\sbox{\EZchimtho}{\eduzylogo[\EZchimcao]}%
	\sbox{\EZchimbox}{\begin{tikzpicture}
			\begin{scope}[transparency group,opacity=\EZchimdo]
				\node[inner sep=0pt] {\usebox{\EZchimtho}};
			\end{scope}
	\end{tikzpicture}}}
\newcommand{\EZvechim}{%
	\ifEZchim
		\node[inner sep=0pt] at (current page.center) {\usebox{\EZchimbox}};
	\fi}

%====================== ĐẦU TRANG - TRANG LẺ ======================
%  (logo bên trái, ô chữ \tenchuong bên phải)
\newcommand{\myfancyheadOdd}{%
	\EZdorong{\tenchuong}%
	\begin{tikzpicture}[remember picture,overlay]
		\coordinate (T) at (current page.north west);
		\coordinate (P) at (current page.north east);
		%--- logo chìm giữa trang
		\EZvechim
		%--- dải nền
		\fill[ezNen] (T) rectangle ([yshift=-\EZdai]P);
		%--- ô chữ: X là góc dưới bên trái, cạnh trái xiên kiểu "/"
		\coordinate (X) at ([xshift=-\EZle-\EZrong-\EZdem-\EZxien,yshift=-\EZdai]P);
		\fill[ezO] (X) --++(\EZxien,\EZo) -- ([yshift=-\EZdai+\EZo]P) -- ([yshift=-\EZdai]P) --cycle;
		\EZvachxien{X}{-1}{1}
		\node[ezchu,anchor=east] at ([xshift=-\EZle,yshift=-\EZdai+0.5*\EZo]P) {\EZfontO\tenchuong};
		%--- logo + tên
		\node[inner sep=0pt,anchor=west] at ([xshift=\EZle,yshift=-0.5*\EZdai]T) {\usebox{\EZlogobox}};
	\end{tikzpicture}%
}

%====================== ĐẦU TRANG - TRANG CHẴN =====================
%  (soi gương: ô chữ \tenchuong bên trái, logo bên phải)
\newcommand{\myfancyheadEven}{%
	\EZdorong{\tenchuong}%
	\begin{tikzpicture}[remember picture,overlay]
		\coordinate (T) at (current page.north west);
		\coordinate (P) at (current page.north east);
		%--- logo chìm giữa trang
		\EZvechim
		%--- dải nền
		\fill[ezNen] (T) rectangle ([yshift=-\EZdai]P);
		%--- ô chữ: X là góc dưới bên phải, cạnh phải xiên kiểu "\"
		\coordinate (X) at ([xshift=\EZle+\EZrong+\EZdem+\EZxien,yshift=-\EZdai]T);
		\fill[ezO] (X) --++(-\EZxien,\EZo) -- ([yshift=-\EZdai+\EZo]T) -- ([yshift=-\EZdai]T) --cycle;
		\EZvachxien{X}{1}{-1}
		\node[ezchu,anchor=west] at ([xshift=\EZle,yshift=-\EZdai+0.5*\EZo]T) {\EZfontO\tenchuong};
		%--- logo + tên
		\node[inner sep=0pt,anchor=east] at ([xshift=-\EZle,yshift=-0.5*\EZdai]P) {\usebox{\EZlogobox}};
	\end{tikzpicture}%
}

%====================== CHÂN TRANG - TRANG LẺ ======================
%  (ô chữ \thuonghieu bên trái, "Trang (số)" bên phải)
\newcommand{\myfancyfootOdd}{%
	\EZdorong{\thuonghieu}%
	\begin{tikzpicture}[remember picture,overlay]
		\coordinate (T) at (current page.south west);
		\coordinate (P) at (current page.south east);
		%--- dải nền
		\fill[ezNen] (T) rectangle ([yshift=\EZdai]P);
		%--- ô chữ sát mép trên dải: X là góc dưới bên phải, cạnh phải xiên kiểu "/"
		\coordinate (X) at ([xshift=\EZle+\EZrong+\EZdem,yshift=\EZdai-\EZo]T);
		\fill[ezO] ([yshift=\EZdai-\EZo]T) -- (X) --++(\EZxien,\EZo) -- ([yshift=\EZdai]T) --cycle;
		\EZvachxien{X}{1}{1}
		\node[ezchu,anchor=west] at ([xshift=\EZle,yshift=\EZdai-0.5*\EZo]T) {\EZfontO\thuonghieu};
		%--- số trang
		\node[circle,draw=white,line width=0.9pt,text=white,inner sep=0pt,minimum size=0.62cm,anchor=east]
		(so) at ([xshift=-\EZle,yshift=0.5*\EZdai]P) {\EZfontTrang\thepage};
		\node[inner sep=0pt,text=white,anchor=base east] (tr) at ([xshift=-0.15cm]so.base west) {\EZfontTrang Trang};
	\end{tikzpicture}%
}

%====================== CHÂN TRANG - TRANG CHẴN ====================
%  (soi gương: "Trang (số)" bên trái, ô chữ \tacgia bên phải)
\newcommand{\myfancyfootEven}{%
	\EZdorong{\tacgia}%
	\begin{tikzpicture}[remember picture,overlay]
		\coordinate (T) at (current page.south west);
		\coordinate (P) at (current page.south east);
		%--- dải nền
		\fill[ezNen] (T) rectangle ([yshift=\EZdai]P);
		%--- ô chữ sát mép trên dải: X là góc dưới bên trái, cạnh trái xiên kiểu "\"
		\coordinate (X) at ([xshift=-\EZle-\EZrong-\EZdem,yshift=\EZdai-\EZo]P);
		\fill[ezO] ([yshift=\EZdai-\EZo]P) -- (X) --++(-\EZxien,\EZo) -- ([yshift=\EZdai]P) --cycle;
		\EZvachxien{X}{-1}{-1}
		\node[ezchu,anchor=east] at ([xshift=-\EZle,yshift=\EZdai-0.5*\EZo]P) {\EZfontO\tacgia};
		%--- số trang
		%    (vòng tròn ẩn so0 nằm giữa dải để lấy đường chân chữ, giống trang lẻ)
		\node[circle,inner sep=0pt,minimum size=0.62cm] (so0) at ([yshift=0.5*\EZdai]T) {\EZfontTrang\phantom{\thepage}};
		\node[inner sep=0pt,text=white,anchor=base west] (tr) at ($(so0.base -| T)+(\EZle,0)$) {\EZfontTrang Trang};
		\node[circle,draw=white,line width=0.9pt,text=white,inner sep=0pt,minimum size=0.62cm,anchor=base west]
		(so) at ([xshift=0.15cm]tr.base east) {\EZfontTrang\thepage};
	\end{tikzpicture}%
}

%--- Giữ tương thích với các file cũ đang gọi \myfancyhead, \myfancyfoot
\newcommand{\myfancyhead}{\myfancyheadOdd}
\newcommand{\myfancyfoot}{\myfancyfootOdd}

%======================= GẮN VÀO TRANG =============================
\pagestyle{fancy}
\fancyhf{}
\fancyhead[LO]{\myfancyheadOdd}
\fancyhead[RE]{\myfancyheadEven}
\fancyfoot[LO]{\myfancyfootOdd}
\fancyfoot[RE]{\myfancyfootEven}
\renewcommand{\headrulewidth}{0pt}
\renewcommand{\footrulewidth}{0pt}
\fancypagestyle{plain}{%
	\fancyhf{}
	\fancyhead[LO]{\myfancyheadOdd}
	\fancyhead[RE]{\myfancyheadEven}
	\fancyfoot[LO]{\myfancyfootOdd}
	\fancyfoot[RE]{\myfancyfootEven}
	\renewcommand{\headrulewidth}{0pt}
	\renewcommand{\footrulewidth}{0pt}
}
%=================== HẾT PHẦN GỘP TỪ eduzy-de.tex ===================
\newcommand{\indam}[1]{
	{\fontsize{12pt}{9pt}\sffamily \selectfont\bfseries\color{gold1} #1}
}
\newcommand{\ind}[1]{
	{\fontsize{14pt}{9pt}\fontfamily{qag} \selectfont\bfseries\color{gold1} #1}
}
\renewcommand{\FalseEX}{\stepcounter{dapan}\textbf{\selectfont\color{gold1}\Alph{dapan}.}\ignorespaces}
\renewtheorem{ex}{\color{gold1}Câu}
\renewcommand{\baselinestretch}{1.1}
\newcommand{\hoac}[1]{
	\left[\begin{aligned}#1\end{aligned}\right.}
\newcommand{\heva}[1]{
	\left\{\begin{aligned}#1\end{aligned}\right.}
\usepackage{esvect}
\def\vec{\vv}
\def\vt{\vv}
\def\overrightarrow{\vv}
%====================== LOGO EDUZY VẼ BẰNG TIKZ ======================
%  Đường viền dò từ logo gốc; dải màu vàng kim đo trực tiếp trên ảnh gốc
%  (nghiêng 21 độ, tô bằng 80 dải hẹp nên hiển thị giống nhau trên mọi trình xem PDF).
%  Dùng: \eduzylogo          -> cao 2.4cm (mặc định)
%        \eduzylogo[1.5cm]   -> cao 1.5cm
\definecolor{ezText}{RGB}{245,200,46}
\newcommand{\eduzylogo}[1][2.4cm]{%
	\resizebox{!}{#1}{%
		\begin{tikzpicture}[x=0.01cm,y=0.01cm]
			\path[use as bounding box] (0,0) rectangle (224,-224);
			% mũ tốt nghiệp + vương miện
			\begin{scope}
				\clip
				(119.4,-176.1)..controls(118.9,-176.1)and(116.8,-176)..(114.6,-176)
				..controls(101.4,-176)and(85.8,-175.9)..(83.7,-175.7)
				..controls(81.8,-175.6)and(81,-175.6)..(80.2,-175.7)
				..controls(78.8,-175.9)and(77.7,-175.9)..(76.2,-175.7)
				..controls(74.9,-175.5)and(70.8,-175.5)..(64.3,-175.7)
				..controls(59.1,-175.9)and(58.3,-175.7)..(58.6,-173.7)
				..controls(58.6,-172.8)and(59.6,-169.6)..(60.3,-168)
				..controls(60.8,-166.9)and(61.4,-165.8)..(62.3,-165)
				..controls(63.1,-164.2)and(64.3,-163.4)..(65,-163.2)
				..controls(65.5,-163.1)and(72.3,-163)..(76.6,-163.1)
				..controls(77.9,-163.1)and(80.1,-163.1)..(81.4,-163)
				..controls(82.8,-162.9)and(86.3,-162.9)..(89.2,-162.9)
				..controls(93,-162.9)and(94.9,-162.8)..(96,-162.7)
				..controls(97.2,-162.5)and(101,-162.5)..(114,-162.4)
				..controls(123,-162.3)and(134.9,-162.2)..(140.2,-162.1)
				..controls(150.1,-162)and(152.2,-162.1)..(153.6,-162.4)
				..controls(154.4,-162.6)and(157.1,-163.8)..(158.1,-164.4)
				..controls(158.8,-164.8)and(159.7,-165.7)..(159.9,-166.3)
				..controls(160,-166.5)and(160.2,-167.1)..(160.2,-167.6)
				..controls(160.3,-168)and(160.6,-169.4)..(160.9,-170.6)
				..controls(161.2,-171.7)and(161.5,-173.1)..(161.6,-173.6)
				..controls(161.8,-174.4)and(161.8,-174.5)..(161.5,-174.9)
				..controls(161.2,-175.6)and(160.8,-175.9)..(159.9,-176)
				..controls(159.2,-176.1)and(154.1,-176.2)..(151.5,-176.2)
				..controls(148.1,-176.1)and(122.6,-176.1)..(121.6,-176.2)
				..controls(120.8,-176.2)and(119.8,-176.2)..(119.4,-176.1)--cycle
				(130.4,-160)..controls(129,-160)and(119.8,-160)..(110.1,-159.9)
				..controls(100.3,-159.8)and(88.5,-159.7)..(83.9,-159.6)
				..controls(77.8,-159.5)and(75.3,-159.5)..(74.5,-159.6)
				..controls(73.9,-159.7)and(72.5,-159.8)..(71.3,-159.8)
				..controls(69.6,-159.8)and(69,-159.8)..(68.5,-159.6)
				..controls(67.6,-159.3)and(66.8,-158.6)..(66.2,-157.6)
				..controls(66,-157.3)and(65.6,-156.5)..(65.3,-156.1)
				..controls(64.6,-155.1)and(64.5,-154.9)..(63.7,-152.9)
				..controls(63.2,-151.6)and(62.9,-150.9)..(61.8,-149.6)
				..controls(61.5,-149.2)and(61.1,-148.5)..(60.9,-147.9)
				..controls(60.7,-147.4)and(60.3,-146.6)..(60.1,-146.1)
				..controls(59.9,-145.5)and(59.5,-144.9)..(59,-144.3)
				..controls(58.5,-143.7)and(58.2,-143.2)..(57.9,-142.5)
				..controls(57.7,-141.9)and(57.3,-141.1)..(57,-140.7)
				..controls(56.6,-140.3)and(56.1,-139.3)..(55.8,-138.5)
				..controls(55.5,-137.8)and(55.1,-136.8)..(54.9,-136.5)
				..controls(54.7,-136.1)and(54.3,-135.2)..(53.9,-134.5)
				..controls(53.6,-133.8)and(53,-132.7)..(52.6,-132)
				..controls(51.6,-130.5)and(51.4,-130.2)..(51.1,-129)
				..controls(50.6,-127.5)and(50.6,-127.4)..(49.7,-126.3)
				..controls(49.2,-125.7)and(48.6,-124.8)..(48.3,-124.2)
				..controls(48,-123.6)and(47.5,-122.7)..(47.3,-122.2)
				..controls(47,-121.8)and(46.5,-120.7)..(46.1,-119.8)
				..controls(45.8,-118.9)and(45.2,-117.8)..(44.9,-117.2)
				..controls(44.2,-116.1)and(43.3,-114.4)..(42.3,-112.2)
				..controls(41.8,-111.3)and(41.3,-110.4)..(41.1,-110.1)
				..controls(40.5,-109.4)and(40.3,-109)..(40,-107.8)
				..controls(39.8,-107)and(39.6,-106.5)..(39.2,-106)
				..controls(38.9,-105.6)and(38.4,-104.7)..(38,-104)
				..controls(37.7,-103.4)and(37.2,-102.5)..(37,-102.1)
				..controls(36.7,-101.7)and(36.3,-100.8)..(36,-100.1)
				..controls(35.6,-98.9)and(35.3,-98.4)..(34.6,-97.6)
				..controls(34.2,-97.1)and(34.2,-96.3)..(34.6,-96)
				..controls(35,-95.7)and(35.5,-95.9)..(36.2,-96.7)
				..controls(36.5,-97)and(37.2,-97.5)..(37.6,-97.7)..controls(38,-98)and(38.6,-98.4)..(39.1,-98.9)
				..controls(39.6,-99.4)and(40.3,-100)..(40.6,-100.3)
				..controls(41.3,-100.8)and(42.4,-101.8)..(43.7,-102.9)
				..controls(44,-103.2)and(44.6,-103.7)..(45,-103.9)
				..controls(45.4,-104.2)and(45.9,-104.6)..(46.1,-104.8)
				..controls(47.1,-105.9)and(48.8,-107.1)..(50.5,-107.9)
				..controls(51.8,-108.6)and(52.3,-109)..(53.4,-110)
				..controls(54.8,-111.2)and(55.6,-111.6)..(57,-112.1)
				..controls(57.5,-112.3)and(58.5,-112.8)..(59.4,-113.3)
				..controls(60.2,-113.7)and(61.3,-114.3)..(61.9,-114.5)
				..controls(62.4,-114.7)and(63.3,-115.1)..(63.8,-115.3)
				..controls(64.4,-115.6)and(65.4,-115.9)..(66.2,-116.1)
				..controls(67,-116.3)and(68.4,-116.6)..(69.2,-116.9)
				..controls(72.3,-117.9)and(72.8,-118)..(78,-118.1)
				..controls(84,-118.3)and(86.1,-118.1)..(88.1,-117.1)
				..controls(89.4,-116.5)and(93,-115)..(93.7,-114.9)
				..controls(94.8,-114.6)and(95.4,-114.2)..(95.9,-113.4)
				..controls(96.2,-113)and(96.7,-112.5)..(97.1,-112.2)
				..controls(98,-111.6)and(101.3,-108.4)..(102.1,-107.2)
				..controls(102.4,-106.8)and(102.9,-106.2)..(103.2,-105.8)
				..controls(103.5,-105.5)and(104,-104.8)..(104.2,-104.2)
				..controls(104.5,-103.6)and(105,-102.8)..(105.3,-102.4)
				..controls(105.6,-102)and(106.1,-101.2)..(106.5,-100.6)
				..controls(107.1,-99.5)and(108.6,-97.5)..(109.5,-96.5)
				..controls(109.9,-96.1)and(110.1,-96)..(110.4,-96)
				..controls(111.1,-96)and(111.6,-96.5)..(111.8,-97.6)
				..controls(112,-98.6)and(112.3,-99.3)..(112.8,-99.9)
				..controls(113.7,-100.9)and(115.2,-103.1)..(115.7,-104.2)
				..controls(116.1,-105)and(116.5,-105.5)..(116.9,-106)
				..controls(117.3,-106.4)and(117.8,-106.9)..(118,-107.3)
				..controls(118.2,-107.6)and(118.6,-108)..(119,-108.3)
				..controls(119.6,-108.8)and(121.1,-110.3)..(121.8,-111.2)
				..controls(122.1,-111.7)and(122.5,-112)..(123.4,-112.4)
				..controls(124.8,-113.2)and(125.8,-113.8)..(127,-114.6)
				..controls(127.5,-114.9)and(128.1,-115.3)..(128.2,-115.4)
				..controls(128.4,-115.5)and(129.3,-115.8)..(130.2,-116)
				..controls(131.1,-116.2)and(132.5,-116.6)..(133.3,-116.9)
				..controls(135.7,-117.9)and(136.1,-117.9)..(142.5,-117.8)
				..controls(145.8,-117.8)and(148.4,-117.7)..(148.8,-117.6)
				..controls(149.2,-117.6)and(150.2,-117.2)..(151,-116.8)
				..controls(152,-116.4)and(152.8,-116)..(153.6,-115.9)
				..controls(154.5,-115.7)and(155.2,-115.4)..(156.3,-114.9)
				..controls(157.1,-114.5)and(158.2,-114)..(158.7,-113.8)
				..controls(159.1,-113.7)and(159.8,-113.3)..(160.2,-113)
				..controls(160.6,-112.7)and(161.3,-112.4)..(161.7,-112.2)
				..controls(162.7,-111.9)and(163.2,-111.7)..(164.1,-111.1)
				..controls(164.5,-110.8)and(165.3,-110.4)..(165.9,-110)
				..controls(166.5,-109.7)and(167.3,-109.2)..(167.6,-108.9)
				..controls(167.9,-108.6)and(168.6,-108.2)..(169,-108)
				..controls(169.4,-107.8)and(170.2,-107.3)..(170.7,-106.9)
				..controls(171.2,-106.5)and(172.1,-105.8)..(172.7,-105.4)
				..controls(173.2,-105)and(174.1,-104.4)..(174.5,-104)
				..controls(174.9,-103.7)and(175.7,-103.1)..(176.2,-102.6)
				..controls(176.7,-102.2)and(177.4,-101.5)..(177.9,-101.2)
				..controls(178.3,-100.8)and(178.8,-100.4)..(179,-100.2)
				..controls(179.2,-100)and(179.6,-99.6)..(180,-99.3)
				..controls(180.4,-99.1)and(180.9,-98.6)..(181.2,-98.4)
				..controls(181.5,-98.1)and(182.1,-97.7)..(182.5,-97.5)
				..controls(183.2,-97.2)and(183.5,-96.9)..(184.3,-95.9)
				..controls(184.6,-95.6)and(184.6,-95.6)..(184.8,-95.7)
				..controls(185,-95.8)and(185.1,-96)..(185.1,-96.2)
				..controls(185.1,-96.5)and(185,-96.7)..(184.6,-97.2)
				..controls(184.1,-97.8)and(184,-97.9)..(183.9,-98.8)
				..controls(183.7,-100.3)and(183.3,-101.3)..(182.9,-101.8)
				..controls(182.6,-102.1)and(182.3,-102.6)..(182.2,-102.9)
				..controls(182,-103.3)and(181.6,-104)..(181.2,-104.6)
				..controls(180.9,-105.2)and(180.5,-106.1)..(180.4,-106.6)
				..controls(179.9,-108)and(179.7,-108.5)..(179,-109.4)
				..controls(178.6,-109.8)and(178.2,-110.5)..(177.9,-111)
				..controls(177.7,-111.5)and(177.3,-112.3)..(177,-112.7)
				..controls(176.6,-113.2)and(176.4,-113.7)..(176.3,-114.2)
				..controls(176.2,-114.7)and(175.9,-115.2)..(175.4,-116)
				..controls(175,-116.6)and(174.4,-117.7)..(174,-118.5)
				..controls(173.7,-119.2)and(173.3,-120.1)..(173,-120.5)
				..controls(172.8,-120.9)and(172.4,-121.8)..(172.1,-122.5)
				..controls(171.8,-123.2)and(171.3,-124.1)..(171,-124.6)
				..controls(170.8,-125.1)and(170.3,-126)..(169.9,-126.6)
				..controls(169.6,-127.3)and(169.2,-128.1)..(169,-128.3)
				..controls(168.8,-128.6)and(168.5,-129.3)..(168.3,-129.8)
				..controls(168.1,-130.3)and(167.8,-131.1)..(167.5,-131.6)
				..controls(166.9,-132.5)and(164.2,-137.9)..(164.2,-138.1)
				..controls(164.2,-138.1)and(163.9,-138.8)..(163.4,-139.7)
				..controls(162.9,-140.5)and(162.4,-141.5)..(162.3,-141.9)
				..controls(161.9,-143)and(161.6,-143.4)..(160.9,-144.1)
				..controls(160.4,-144.6)and(160.3,-144.8)..(160.1,-145.4)
				..controls(160,-145.8)and(159.6,-146.8)..(159.2,-147.8)
				..controls(158.7,-148.8)and(158.1,-150)..(157.9,-150.6)
				..controls(157.6,-151.1)and(157.2,-151.9)..(156.9,-152.4)
				..controls(156.5,-152.8)and(156.1,-153.7)..(155.8,-154.4)
				..controls(155.6,-155)and(155.2,-155.9)..(154.9,-156.3)
				..controls(154.7,-156.7)and(154.4,-157.4)..(154.2,-157.9)
				..controls(154.1,-158.4)and(153.8,-159)..(153.5,-159.3)
				..controls(152.8,-160.1)and(153,-160.1)..(142.4,-160.1)
				..controls(137.2,-160.1)and(131.8,-160.1)..(130.4,-160)--cycle
				(141.5,-103.7)..controls(140.8,-103)and(139.8,-102.1)..(139.2,-101.6)
				..controls(138.5,-101.2)and(137.9,-100.6)..(137.8,-100.4)
				..controls(137,-99.5)and(134,-96.5)..(133.2,-96)
				..controls(132.7,-95.7)and(131.8,-94.9)..(131.1,-94.3)
				..controls(129.4,-92.7)and(127.9,-91.6)..(126.6,-91.1)
				..controls(125.2,-90.5)and(124.1,-90)..(122.9,-89.1)
				..controls(122.1,-88.6)and(121.3,-88.2)..(119.9,-87.7)
				..controls(118.8,-87.3)and(117.5,-86.8)..(116.9,-86.5)
				..controls(114.3,-85.5)and(113.7,-85.4)..(109.7,-85.5)
				..controls(105.3,-85.7)and(103.9,-85.9)..(101.6,-87)
				..controls(100.7,-87.4)and(99.4,-87.8)..(97.9,-88)
				..controls(97,-88.2)and(96.6,-88.3)..(95.8,-88.9)
				..controls(95.4,-89.2)and(94.7,-89.6)..(94.2,-89.8)
				..controls(92.9,-90.3)and(92.3,-90.7)..(90.9,-92.1)
				..controls(90.2,-92.7)and(89.4,-93.3)..(89,-93.5)..controls(87,-94.5)and(85.2,-96)..(83.6,-97.9)
				..controls(83.3,-98.2)and(82.5,-98.9)..(81.8,-99.4)
				..controls(80.7,-100.1)and(80.5,-100.3)..(80,-101.1)
				..controls(78.7,-103.1)and(77.2,-103.9)..(76.4,-103)
				..controls(75.7,-102.2)and(75.7,-103.2)..(75.7,-83.6)
				..controls(75.6,-73.6)and(75.6,-64.8)..(75.5,-64)
				..controls(75.5,-63.2)and(75.5,-61.4)..(75.6,-59.8)--(75.7,-57.1)--(76,-56.8)
				..controls(76.8,-56.2)and(78,-56.3)..(79.8,-57.2)
				..controls(80.4,-57.5)and(81.3,-57.8)..(81.8,-57.9)
				..controls(82.3,-58.1)and(83,-58.4)..(83.6,-58.8)..controls(84.3,-59.2)and(85,-59.5)..(86,-59.8)
				..controls(86.8,-60.1)and(88,-60.6)..(88.7,-60.9)
				..controls(89.3,-61.2)and(90.6,-61.7)..(91.4,-62)
				..controls(92.2,-62.3)and(93.7,-62.9)..(94.6,-63.3)
				..controls(95.6,-63.7)and(96.8,-64.1)..(97.4,-64.3)
				..controls(98,-64.5)and(99,-64.9)..(99.7,-65.2)
				..controls(100.4,-65.6)and(101.5,-66)..(102.3,-66.2)
				..controls(104.2,-66.8)and(105,-67.1)..(106.2,-67.8)
				..controls(107.9,-68.7)and(110.2,-69.4)..(111.3,-69.2)
				..controls(112.5,-69)and(115,-68.2)..(115.6,-67.8)
				..controls(117.3,-66.7)and(117.6,-66.6)..(119.5,-66.2)
				..controls(121.9,-65.8)and(122.1,-65.7)..(123,-65.1)
				..controls(124.1,-64.4)and(124.4,-64.3)..(126,-63.9)
				..controls(127,-63.7)and(127.7,-63.4)..(128.4,-63.1)
				..controls(129,-62.8)and(129.7,-62.5)..(130.6,-62.4)
				..controls(132.1,-62.1)and(134,-61.5)..(135.2,-61)
				..controls(135.7,-60.8)and(136.7,-60.5)..(137.3,-60.3)
				..controls(138.4,-60)and(139.3,-59.7)..(141.1,-59)
				..controls(141.6,-58.7)and(142,-58.6)..(142.5,-58.6)
				..controls(143.2,-58.6)and(143.3,-58.7)..(143.6,-59)
				..controls(144,-59.5)and(144.1,-59.9)..(144,-60.8)
				..controls(144,-61.2)and(144,-71.1)..(144,-82.8)--(143.9,-104.1)--(143.7,-104.5)
				..controls(143.5,-104.7)and(143.3,-104.8)..(143,-104.9)
				..controls(142.6,-104.9)and(142.5,-104.8)..(141.5,-103.7)--cycle
				(27.8,-91.4)..controls(27.1,-90.9)and(26.3,-90)..(26.1,-89.5)
				..controls(26,-89.2)and(25.8,-88.6)..(25.7,-87.9)
				..controls(25.4,-86.2)and(25.6,-84.8)..(26.5,-83.5)
				..controls(27.4,-82.3)and(28.3,-82)..(29.4,-82.3)
				..controls(31.1,-82.8)and(32.1,-84.9)..(32,-87.4)
				..controls(31.9,-89.3)and(31.2,-90.6)..(29.9,-91.4)
				..controls(29.1,-91.9)and(28.6,-91.9)..(27.8,-91.4)--cycle
				(190.1,-91.3)..controls(189.7,-91.2)and(189.3,-90.9)..(189.1,-90.7)
				..controls(188.3,-89.8)and(187.6,-87.7)..(187.6,-86.2)
				..controls(187.6,-85.5)and(187.7,-85.2)..(188.2,-84.4)
				..controls(189.1,-82.6)and(190.1,-81.8)..(191.3,-82.1)
				..controls(192,-82.3)and(193.5,-83.7)..(193.8,-84.4)
				..controls(194.1,-85.1)and(194.2,-87.1)..(194,-88.3)
				..controls(193.8,-89.8)and(193.3,-90.5)..(192.3,-91.3)
				..controls(191.7,-91.7)and(190.9,-91.7)..(190.1,-91.3)--cycle
				(173.7,-83.4)..controls(172.2,-82.5)and(172,-82.5)..(170.6,-82.4)
				..controls(169.7,-82.4)and(168.8,-82.4)..(168.3,-82.5)
				..controls(167.2,-82.7)and(166.9,-82.7)..(166.6,-82.2)
				..controls(166.3,-81.8)and(166.3,-81)..(166.8,-80)
				..controls(167.2,-78.9)and(167.4,-77.4)..(167.7,-74.1)
				..controls(167.8,-72.8)and(167.9,-71.8)..(167.9,-71.7)
				..controls(168,-71.7)and(168,-70.2)..(168,-68.4)..controls(168,-65.4)and(168,-64.9)..(167.7,-63)
				..controls(167.3,-60.5)and(167.1,-59.9)..(166.4,-58.4)
				..controls(166.1,-57.7)and(165.8,-56.9)..(165.7,-56.6)
				..controls(165.7,-56.2)and(165.4,-55.6)..(165.2,-55.2)
				..controls(164.6,-53.8)and(161.9,-50.7)..(161,-50.3)
				..controls(159.8,-49.7)and(157.9,-50)..(155.2,-51.2)
				..controls(154.7,-51.4)and(153.4,-51.8)..(152.2,-52.1)
				..controls(150.8,-52.4)and(149.4,-52.9)..(148.6,-53.2)
				..controls(147.8,-53.6)and(146.7,-54)..(146,-54.2)
				..controls(145.4,-54.4)and(144.5,-54.7)..(144,-54.9)
				..controls(143.6,-55.2)and(142.7,-55.5)..(141.8,-55.7)
				..controls(141,-55.9)and(139.6,-56.3)..(138.7,-56.6)
				..controls(137.8,-57)and(136.3,-57.4)..(135.3,-57.6)
				..controls(133.8,-57.9)and(133.2,-58.1)..(132.1,-58.7)
				..controls(131.1,-59.2)and(130.3,-59.4)..(129.4,-59.6)
				..controls(128.7,-59.8)and(127.8,-60)..(127.4,-60.1)
				..controls(127,-60.3)and(126,-60.5)..(125.1,-60.7)
				..controls(123.2,-61.1)and(121.7,-61.5)..(120.2,-62.1)
				..controls(119.7,-62.3)and(118.1,-62.8)..(116.7,-63.2)
				..controls(114.6,-63.9)and(114,-64.1)..(113,-64.7)
				..controls(111.9,-65.4)and(111.8,-65.4)..(111,-65.5)
				..controls(109.9,-65.5)and(109.3,-65.4)..(107.7,-64.7)
				..controls(107,-64.4)and(106,-64)..(105.3,-63.9)
				..controls(103.5,-63.5)and(102.2,-63)..(101.2,-62.5)
				..controls(100.1,-61.8)and(99.3,-61.5)..(97.8,-61.1)
				..controls(97.2,-61)and(96.2,-60.7)..(95.7,-60.4)
				..controls(95.1,-60.1)and(94.2,-59.8)..(93.5,-59.6)
				..controls(92.9,-59.4)and(91.7,-59)..(91,-58.7)..controls(90.2,-58.4)and(88.9,-58)..(88,-57.7)
				..controls(87.2,-57.5)and(86.1,-57.1)..(85.5,-56.8)
				..controls(85,-56.6)and(84.1,-56.2)..(83.5,-56)..controls(82.9,-55.7)and(81.8,-55.2)..(81,-54.9)
				..controls(80.1,-54.5)and(79.1,-54.2)..(78.7,-54.1)
				..controls(77.9,-53.8)and(77.3,-53.6)..(76.2,-52.9)
				..controls(75.7,-52.6)and(75.2,-52.4)..(74.1,-52.2)
				..controls(72.5,-51.9)and(72.4,-51.8)..(71.3,-51.1)
				..controls(69.9,-50.1)and(69.4,-49.9)..(67.9,-49.8)
				..controls(67,-49.7)and(66.6,-49.6)..(66.3,-49.4)
				..controls(66.1,-49.3)and(65.7,-49)..(65.4,-48.8)
				..controls(65.1,-48.7)and(64.2,-48.4)..(63.6,-48.2)
				..controls(62.9,-48)and(61.7,-47.6)..(60.9,-47.2)
				..controls(60.1,-46.8)and(58.9,-46.4)..(58.2,-46.2)
				..controls(57.6,-46)and(56.8,-45.7)..(56.5,-45.5)
				..controls(54.9,-44.6)and(54.6,-44.5)..(53.2,-44.2)
				..controls(51.7,-43.8)and(51,-43.6)..(49.9,-42.9)
				..controls(49.5,-42.6)and(48.9,-42.4)..(47.5,-42)
				..controls(45.7,-41.6)and(45.6,-41.5)..(44.3,-40.6)
				..controls(43.5,-40.1)and(42.8,-39.6)..(42.7,-39.4)
				..controls(42.4,-38.7)and(43,-38.1)..(44.5,-37.6)..controls(44.9,-37.5)and(46,-37)..(46.9,-36.6)
				..controls(47.9,-36.1)and(48.8,-35.8)..(49.6,-35.6)
				..controls(50.5,-35.4)and(51.3,-35.2)..(52.4,-34.6)
				..controls(53.3,-34.1)and(54.5,-33.7)..(54.9,-33.5)
				..controls(55.4,-33.4)and(56.4,-33)..(57.2,-32.6)
				..controls(58.6,-32)and(60.1,-31.5)..(61.6,-31.2)
				..controls(62.3,-31.1)and(62.7,-31)..(63.2,-30.7)
				..controls(64.1,-30)and(64.3,-29.9)..(65.3,-29.7)
				..controls(65.7,-29.6)and(66.5,-29.2)..(67.1,-28.9)
				..controls(67.6,-28.6)and(68.8,-28.2)..(69.6,-27.9)
				..controls(70.5,-27.7)and(71.6,-27.3)..(72.1,-27)
				..controls(72.6,-26.7)and(73.8,-26.2)..(74.8,-25.9)
				..controls(75.7,-25.5)and(76.7,-25.1)..(77.1,-25)
				..controls(77.4,-24.8)and(78.5,-24.4)..(79.6,-24.1)
				..controls(80.6,-23.8)and(82,-23.2)..(82.8,-22.9)
				..controls(83.5,-22.5)and(84.8,-22)..(85.8,-21.8)
				..controls(86.7,-21.5)and(87.8,-21.1)..(88.1,-20.9)
				..controls(88.4,-20.7)and(89.3,-20.3)..(90.1,-20.1)
				..controls(90.8,-19.8)and(91.9,-19.4)..(92.4,-19.1)
				..controls(93.5,-18.5)and(94.4,-18.2)..(95.4,-18.1)
				..controls(95.7,-18)and(96.4,-17.8)..(96.9,-17.5)
				..controls(97.5,-17.3)and(98.4,-16.9)..(99,-16.6)
				..controls(99.7,-16.3)and(100.9,-15.9)..(101.9,-15.6)
				..controls(102.9,-15.3)and(104.4,-14.7)..(105.3,-14.3)
				..controls(108.1,-13.2)and(108.5,-13.1)..(110.1,-13.1)
				..controls(111.4,-13.1)and(111.5,-13.1)..(112.4,-13.5)
				..controls(113.2,-14)and(115,-14.8)..(116.6,-15.4)
				..controls(116.9,-15.5)and(117.9,-15.8)..(118.8,-16.1)
				..controls(120.9,-16.6)and(121.5,-16.9)..(122.6,-17.9)
				..controls(123,-18.3)and(123.7,-18.8)..(123.9,-18.9)
				..controls(124.6,-19.4)and(126,-19.8)..(127.9,-20.2)
				..controls(129.6,-20.5)and(129.8,-20.6)..(131.1,-21.4)
				..controls(131.5,-21.7)and(132.4,-22.1)..(133.1,-22.3)
				..controls(136.1,-23.2)and(138.3,-24.1)..(140.2,-25)
				..controls(141,-25.5)and(141.5,-25.6)..(142.5,-25.8)
				..controls(145.1,-26.3)and(146,-26.6)..(149.4,-28.1)
				..controls(151.3,-28.9)and(153.5,-29.8)..(154.3,-30.1)
				..controls(155.1,-30.4)and(156.2,-30.8)..(156.6,-31)
				..controls(157.1,-31.2)and(158.2,-31.6)..(159.1,-31.8)
				..controls(160.6,-32.2)and(160.9,-32.3)..(161.7,-32.9)
				..controls(162.4,-33.3)and(162.9,-33.6)..(163.7,-33.8)
				..controls(164.3,-34)and(164.8,-34.1)..(164.9,-34.1)
				..controls(165,-34.1)and(165.6,-34.4)..(166.4,-34.8)
				..controls(167.3,-35.2)and(168.1,-35.4)..(169,-35.6)
				..controls(170,-35.8)and(170.5,-36)..(171.7,-36.6)
				..controls(172.4,-37)and(173.3,-37.4)..(173.6,-37.4)
				..controls(175.6,-37.7)and(176.2,-38)..(176.7,-39.1)
				..controls(176.9,-39.6)and(176.9,-39.7)..(176.7,-40.4)
				..controls(176.4,-41.6)and(175.8,-42.7)..(175.2,-43.4)
				..controls(174.5,-44.1)and(172.6,-45.5)..(171.7,-45.8)
				..controls(170.6,-46.2)and(169.1,-46)..(168.2,-45.4)
				..controls(167.9,-45.2)and(167.3,-44.8)..(166.9,-44.4)
				..controls(166.4,-44.1)and(165.8,-43.6)..(165.4,-43.2)
				..controls(165.1,-42.9)and(164.4,-42.3)..(163.9,-41.9)
				..controls(163.3,-41.5)and(162.7,-40.9)..(162.4,-40.5)
				..controls(161.4,-39.2)and(159.9,-38)..(158.8,-37.8)
				..controls(158.4,-37.8)and(158.3,-37.7)..(158.2,-37.2)
				..controls(158,-36.8)and(157.4,-36.1)..(157.1,-36.1)
				..controls(157,-36.1)and(156.7,-36.1)..(156.5,-36)
				..controls(156.1,-35.9)and(155.9,-35.7)..(155.7,-35.2)
				..controls(155.3,-34.6)and(155.1,-34.4)..(154,-34.1)
				..controls(153.6,-34)and(153.2,-33.7)..(152.9,-33.5)
				..controls(152.4,-32.9)and(151.6,-32.4)..(151.3,-32.4)
				..controls(150.9,-32.4)and(150.4,-32.6)..(150.4,-32.9)
				..controls(150.4,-33.2)and(150.8,-33.4)..(151.4,-33.6)
				..controls(152.1,-33.7)and(152.8,-34.1)..(153.5,-34.9)
				..controls(154.1,-35.5)and(154.3,-35.6)..(155,-35.8)
				..controls(155.7,-36)and(155.7,-36)..(156,-36.6)
				..controls(156.3,-37.3)and(156.7,-37.7)..(157.5,-37.8)
				..controls(158.1,-37.9)and(158.2,-38)..(158.4,-38.8)
				..controls(158.5,-39.4)and(158.6,-39.6)..(159.6,-40.4)
				..controls(160.2,-41)and(160.9,-41.6)..(161.3,-41.8)
				..controls(161.7,-42)and(161.9,-42.3)..(162.2,-42.8)
				..controls(162.4,-43.2)and(162.8,-43.8)..(163.1,-44.1)
				..controls(163.9,-44.8)and(165.1,-46.7)..(165.7,-48)
				..controls(166.2,-49.3)and(166.4,-49.7)..(167.2,-50.5)
				..controls(167.5,-50.8)and(167.9,-51.4)..(168.1,-51.8)
				..controls(168.3,-52.2)and(168.8,-53)..(169.1,-53.6)
				..controls(169.6,-54.3)and(169.8,-54.9)..(169.9,-55.6)
				..controls(170.4,-57.8)and(170.5,-58.1)..(171.2,-59.5)
				..controls(171.6,-60.3)and(172,-61.5)..(172.2,-62.1)
				..controls(172.3,-62.8)and(172.7,-63.7)..(172.9,-64.1)
				..controls(173.2,-64.6)and(173.4,-65.3)..(173.4,-65.6)
				..controls(173.8,-68.8)and(174.2,-70.6)..(174.9,-72.6)
				..controls(175.5,-74.5)and(175.9,-77.1)..(176.1,-80.2)
				..controls(176.3,-83.4)and(175.5,-84.4)..(173.7,-83.4)--cycle;
				\definecolor{ezs}{RGB}{247,179,26}\fill[ezs](-5.7,-156.6)--(-3.4,-157.5)--(61.9,12.7)--(59.7,13.6)--cycle;
				\definecolor{ezs}{RGB}{247,179,22}\fill[ezs](-3.8,-157.3)--(-1.5,-158.2)--(63.9,12)--(61.6,12.9)--cycle;
				\definecolor{ezs}{RGB}{248,179,19}\fill[ezs](-1.8,-158.1)--(0.4,-158.9)--(65.8,11.2)--(63.5,12.1)--cycle;
				\definecolor{ezs}{RGB}{248,178,16}\fill[ezs](0.1,-158.8)--(2.4,-159.7)--(67.7,10.5)--(65.4,11.4)--cycle;
				\definecolor{ezs}{RGB}{249,178,15}\fill[ezs](2,-159.5)--(4.3,-160.4)--(69.6,9.8)--(67.3,10.7)--cycle;
				\definecolor{ezs}{RGB}{251,177,17}\fill[ezs](3.9,-160.3)--(6.2,-161.1)--(71.5,9)--(69.2,9.9)--cycle;
				\definecolor{ezs}{RGB}{252,179,18}\fill[ezs](5.8,-161)--(8.1,-161.9)--(73.4,8.3)--(71.1,9.2)--cycle;
				\definecolor{ezs}{RGB}{252,182,17}\fill[ezs](7.7,-161.7)--(10,-162.6)--(75.3,7.6)--(73,8.5)--cycle;
				\definecolor{ezs}{RGB}{252,185,17}\fill[ezs](9.6,-162.5)--(11.9,-163.4)--(77.2,6.8)--(75,7.7)--cycle;
				\definecolor{ezs}{RGB}{251,189,18}\fill[ezs](11.5,-163.2)--(13.8,-164.1)--(79.2,6.1)--(76.9,7)--cycle;
				\definecolor{ezs}{RGB}{251,191,22}\fill[ezs](13.4,-163.9)--(15.7,-164.8)--(81.1,5.4)--(78.8,6.2)--cycle;
				\definecolor{ezs}{RGB}{252,193,25}\fill[ezs](15.4,-164.7)--(17.6,-165.6)--(83,4.6)--(80.7,5.5)--cycle;
				\definecolor{ezs}{RGB}{253,195,27}\fill[ezs](17.3,-165.4)--(19.6,-166.3)--(84.9,3.9)--(82.6,4.8)--cycle;
				\definecolor{ezs}{RGB}{253,196,31}\fill[ezs](19.2,-166.1)--(21.5,-167)--(86.8,3.2)--(84.5,4)--cycle;
				\definecolor{ezs}{RGB}{254,199,35}\fill[ezs](21.1,-166.9)--(23.4,-167.8)--(88.7,2.4)--(86.4,3.3)--cycle;
				\definecolor{ezs}{RGB}{253,201,38}\fill[ezs](23,-167.6)--(25.3,-168.5)--(90.6,1.7)--(88.3,2.6)--cycle;
				\definecolor{ezs}{RGB}{253,203,40}\fill[ezs](24.9,-168.3)--(27.2,-169.2)--(92.5,1)--(90.3,1.8)--cycle;
				\definecolor{ezs}{RGB}{253,205,42}\fill[ezs](26.8,-169.1)--(29.1,-170)--(94.4,0.2)--(92.2,1.1)--cycle;
				\definecolor{ezs}{RGB}{253,207,45}\fill[ezs](28.7,-169.8)--(31,-170.7)--(96.4,-0.5)--(94.1,0.4)--cycle;
				\definecolor{ezs}{RGB}{253,210,49}\fill[ezs](30.7,-170.5)--(32.9,-171.4)--(98.3,-1.2)--(96,-0.4)--cycle;
				\definecolor{ezs}{RGB}{253,212,54}\fill[ezs](32.6,-171.3)--(34.9,-172.2)--(100.2,-2)--(97.9,-1.1)--cycle;
				\definecolor{ezs}{RGB}{253,214,59}\fill[ezs](34.5,-172)--(36.8,-172.9)--(102.1,-2.7)--(99.8,-1.8)--cycle;
				\definecolor{ezs}{RGB}{253,216,63}\fill[ezs](36.4,-172.7)--(38.7,-173.6)--(104,-3.4)--(101.7,-2.6)--cycle;
				\definecolor{ezs}{RGB}{253,218,65}\fill[ezs](38.3,-173.5)--(40.6,-174.4)--(105.9,-4.2)--(103.6,-3.3)--cycle;
				\definecolor{ezs}{RGB}{254,220,66}\fill[ezs](40.2,-174.2)--(42.5,-175.1)--(107.8,-4.9)--(105.5,-4)--cycle;
				\definecolor{ezs}{RGB}{254,222,68}\fill[ezs](42.1,-175)--(44.4,-175.8)--(109.7,-5.6)--(107.5,-4.8)--cycle;
				\definecolor{ezs}{RGB}{254,224,71}\fill[ezs](44,-175.7)--(46.3,-176.6)--(111.7,-6.4)--(109.4,-5.5)--cycle;
				\definecolor{ezs}{RGB}{254,226,75}\fill[ezs](46,-176.4)--(48.2,-177.3)--(113.6,-7.1)--(111.3,-6.2)--cycle;
				\definecolor{ezs}{RGB}{254,228,79}\fill[ezs](47.9,-177.2)--(50.1,-178)--(115.5,-7.8)--(113.2,-7)--cycle;
				\definecolor{ezs}{RGB}{254,230,82}\fill[ezs](49.8,-177.9)--(52.1,-178.8)--(117.4,-8.6)--(115.1,-7.7)--cycle;
				\definecolor{ezs}{RGB}{254,232,85}\fill[ezs](51.7,-178.6)--(54,-179.5)--(119.3,-9.3)--(117,-8.4)--cycle;
				\definecolor{ezs}{RGB}{254,234,88}\fill[ezs](53.6,-179.4)--(55.9,-180.2)--(121.2,-10)--(118.9,-9.2)--cycle;
				\definecolor{ezs}{RGB}{254,237,91}\fill[ezs](55.5,-180.1)--(57.8,-181)--(123.1,-10.8)--(120.8,-9.9)--cycle;
				\definecolor{ezs}{RGB}{254,239,94}\fill[ezs](57.4,-180.8)--(59.7,-181.7)--(125,-11.5)--(122.8,-10.6)--cycle;
				\definecolor{ezs}{RGB}{254,241,97}\fill[ezs](59.3,-181.6)--(61.6,-182.4)--(127,-12.2)--(124.7,-11.4)--cycle;
				\definecolor{ezs}{RGB}{254,243,101}\fill[ezs](61.2,-182.3)--(63.5,-183.2)--(128.9,-13)--(126.6,-12.1)--cycle;
				\definecolor{ezs}{RGB}{254,245,103}\fill[ezs](63.2,-183)--(65.4,-183.9)--(130.8,-13.7)--(128.5,-12.8)--cycle;
				\definecolor{ezs}{RGB}{254,245,105}\fill[ezs](65.1,-183.8)--(67.4,-184.6)--(132.7,-14.4)--(130.4,-13.6)--cycle;
				\definecolor{ezs}{RGB}{254,245,105}\fill[ezs](67,-184.5)--(69.3,-185.4)--(134.6,-15.2)--(132.3,-14.3)--cycle;
				\definecolor{ezs}{RGB}{254,244,103}\fill[ezs](68.9,-185.2)--(71.2,-186.1)--(136.5,-15.9)--(134.2,-15)--cycle;
				\definecolor{ezs}{RGB}{253,243,101}\fill[ezs](70.8,-186)--(73.1,-186.8)--(138.4,-16.6)--(136.1,-15.8)--cycle;
				\definecolor{ezs}{RGB}{253,241,98}\fill[ezs](72.7,-186.7)--(75,-187.6)--(140.3,-17.4)--(138,-16.5)--cycle;
				\definecolor{ezs}{RGB}{253,239,95}\fill[ezs](74.6,-187.4)--(76.9,-188.3)--(142.2,-18.1)--(140,-17.2)--cycle;
				\definecolor{ezs}{RGB}{253,236,92}\fill[ezs](76.5,-188.2)--(78.8,-189)--(144.2,-18.8)--(141.9,-18)--cycle;
				\definecolor{ezs}{RGB}{253,234,88}\fill[ezs](78.5,-188.9)--(80.7,-189.8)--(146.1,-19.6)--(143.8,-18.7)--cycle;
				\definecolor{ezs}{RGB}{253,232,85}\fill[ezs](80.4,-189.6)--(82.7,-190.5)--(148,-20.3)--(145.7,-19.4)--cycle;
				\definecolor{ezs}{RGB}{253,230,82}\fill[ezs](82.3,-190.4)--(84.6,-191.2)--(149.9,-21)--(147.6,-20.2)--cycle;
				\definecolor{ezs}{RGB}{253,228,78}\fill[ezs](84.2,-191.1)--(86.5,-192)--(151.8,-21.8)--(149.5,-20.9)--cycle;
				\definecolor{ezs}{RGB}{253,226,74}\fill[ezs](86.1,-191.8)--(88.4,-192.7)--(153.7,-22.5)--(151.4,-21.6)--cycle;
				\definecolor{ezs}{RGB}{253,224,70}\fill[ezs](88,-192.6)--(90.3,-193.4)--(155.6,-23.3)--(153.3,-22.4)--cycle;
				\definecolor{ezs}{RGB}{253,222,67}\fill[ezs](89.9,-193.3)--(92.2,-194.2)--(157.5,-24)--(155.3,-23.1)--cycle;
				\definecolor{ezs}{RGB}{253,220,65}\fill[ezs](91.8,-194)--(94.1,-194.9)--(159.5,-24.7)--(157.2,-23.8)--cycle;
				\definecolor{ezs}{RGB}{253,217,63}\fill[ezs](93.7,-194.8)--(96,-195.6)--(161.4,-25.5)--(159.1,-24.6)--cycle;
				\definecolor{ezs}{RGB}{253,215,60}\fill[ezs](95.7,-195.5)--(97.9,-196.4)--(163.3,-26.2)--(161,-25.3)--cycle;
				\definecolor{ezs}{RGB}{253,213,56}\fill[ezs](97.6,-196.2)--(99.9,-197.1)--(165.2,-26.9)--(162.9,-26)--cycle;
				\definecolor{ezs}{RGB}{252,211,53}\fill[ezs](99.5,-197)--(101.8,-197.8)--(167.1,-27.7)--(164.8,-26.8)--cycle;
				\definecolor{ezs}{RGB}{252,209,49}\fill[ezs](101.4,-197.7)--(103.7,-198.6)--(169,-28.4)--(166.7,-27.5)--cycle;
				\definecolor{ezs}{RGB}{252,207,45}\fill[ezs](103.3,-198.4)--(105.6,-199.3)--(170.9,-29.1)--(168.6,-28.2)--cycle;
				\definecolor{ezs}{RGB}{253,205,42}\fill[ezs](105.2,-199.2)--(107.5,-200)--(172.8,-29.9)--(170.6,-29)--cycle;
				\definecolor{ezs}{RGB}{253,202,40}\fill[ezs](107.1,-199.9)--(109.4,-200.8)--(174.7,-30.6)--(172.5,-29.7)--cycle;
				\definecolor{ezs}{RGB}{253,200,38}\fill[ezs](109,-200.6)--(111.3,-201.5)--(176.7,-31.3)--(174.4,-30.4)--cycle;
				\definecolor{ezs}{RGB}{253,198,35}\fill[ezs](111,-201.4)--(113.2,-202.2)--(178.6,-32.1)--(176.3,-31.2)--cycle;
				\definecolor{ezs}{RGB}{253,196,31}\fill[ezs](112.9,-202.1)--(115.2,-203)--(180.5,-32.8)--(178.2,-31.9)--cycle;
				\definecolor{ezs}{RGB}{253,194,28}\fill[ezs](114.8,-202.8)--(117.1,-203.7)--(182.4,-33.5)--(180.1,-32.6)--cycle;
				\definecolor{ezs}{RGB}{253,192,26}\fill[ezs](116.7,-203.6)--(119,-204.5)--(184.3,-34.3)--(182,-33.4)--cycle;
				\definecolor{ezs}{RGB}{253,189,24}\fill[ezs](118.6,-204.3)--(120.9,-205.2)--(186.2,-35)--(183.9,-34.1)--cycle;
				\definecolor{ezs}{RGB}{253,186,20}\fill[ezs](120.5,-205)--(122.8,-205.9)--(188.1,-35.7)--(185.8,-34.8)--cycle;
				\definecolor{ezs}{RGB}{253,184,17}\fill[ezs](122.4,-205.8)--(124.7,-206.7)--(190,-36.5)--(187.8,-35.6)--cycle;
				\definecolor{ezs}{RGB}{252,183,15}\fill[ezs](124.3,-206.5)--(126.6,-207.4)--(192,-37.2)--(189.7,-36.3)--cycle;
				\definecolor{ezs}{RGB}{252,181,14}\fill[ezs](126.3,-207.2)--(128.5,-208.1)--(193.9,-37.9)--(191.6,-37.1)--cycle;
				\definecolor{ezs}{RGB}{252,179,13}\fill[ezs](128.2,-208)--(130.4,-208.9)--(195.8,-38.7)--(193.5,-37.8)--cycle;
				\definecolor{ezs}{RGB}{252,178,12}\fill[ezs](130.1,-208.7)--(132.4,-209.6)--(197.7,-39.4)--(195.4,-38.5)--cycle;
				\definecolor{ezs}{RGB}{251,178,10}\fill[ezs](132,-209.4)--(134.3,-210.3)--(199.6,-40.1)--(197.3,-39.3)--cycle;
				\definecolor{ezs}{RGB}{251,178,10}\fill[ezs](133.9,-210.2)--(136.2,-211.1)--(201.5,-40.9)--(199.2,-40)--cycle;
				\definecolor{ezs}{RGB}{250,177,11}\fill[ezs](135.8,-210.9)--(138.1,-211.8)--(203.4,-41.6)--(201.1,-40.7)--cycle;
				\definecolor{ezs}{RGB}{250,178,11}\fill[ezs](137.7,-211.6)--(140,-212.5)--(205.3,-42.3)--(203.1,-41.5)--cycle;
				\definecolor{ezs}{RGB}{248,180,10}\fill[ezs](139.6,-212.4)--(141.9,-213.3)--(207.3,-43.1)--(205,-42.2)--cycle;
				\definecolor{ezs}{RGB}{247,181,9}\fill[ezs](141.5,-213.1)--(143.8,-214)--(209.2,-43.8)--(206.9,-42.9)--cycle;
				\definecolor{ezs}{RGB}{246,182,9}\fill[ezs](143.5,-213.8)--(145.7,-214.7)--(211.1,-44.5)--(208.8,-43.7)--cycle;
				\definecolor{ezs}{RGB}{246,182,10}\fill[ezs](145.4,-214.6)--(147.7,-215.5)--(213,-45.3)--(210.7,-44.4)--cycle;
			\end{scope}
			% chữ EDUZY (Times đậm, làm dày nét bằng 5 lớp lệch nhẹ; căn đúng khung chữ logo gốc)
			\foreach \dx/\dy in {-0.7/0,0.7/0,0/-0.3,0/0.3,0/0}
			\node[anchor=center,inner sep=0pt,outer sep=0pt,text=ezText,align=center] at (109.6+\dx,-193.6+\dy)
			{\resizebox{1.01cm}{0.22cm}{\fontencoding{T1}\fontfamily{ptm}\fontseries{b}\selectfont EDUZY}};
		\end{tikzpicture}%
	}%
}
%====================== THIẾT LẬP THÂN ĐỀ ======================
%--- Môi trường tự luận {bt}: ex_test 3.3.4 chưa khai báo sẵn nên tạo ở đây
%    (ex_test đã tự móc \fixshowans và \listenumerate cho bt)
\makeatletter
\@ifundefined{bt}{\newtheorem{bt}{\color{gold1}\namebt}}{}
\makeatother
\renewcommand{\sectionmark}[1]{\markright{\thesection\ #1}}
%--- Khung tiêu đề (logo + tối đa 4 dòng; dòng 3, 4 để trống {} thì bỏ qua)
%  Đề thi:    \khungtieude{THI THỬ THPT QUỐC GIA NĂM 2026}{Môn Toán -- Đề số 21}
%                         {(Đề thi gồm: ...)}{Thời gian làm bài: 90 phút (không kể thời gian phát đề).}
%  Chuyên đề: \khungtieude{CHUYÊN ĐỀ 3 -- DÃY SỐ VÀ CÔNG THỨC TRUY HỒI}
%                         {Bồi dưỡng học sinh giỏi Toán THPT}{...}{}
\newcommand{\khungtieude}[4]{%
	\begin{tcolorbox}[colframe=gold1,colback=gold5!45,boxrule=0.5mm]
		\begin{minipage}[c]{3cm}
			\centerline{\eduzylogo[2.4cm]}
		\end{minipage}\hspace{0.5cm}
		\begin{minipage}[c]{12.5cm}
			\vskip 0.15 cm
			\centerline{\ind{#1}}
			\vskip 0.15 cm
			\centerline{\textbf{#2}}
			\if\relax\detokenize{#3}\relax\else
				\vskip 0.15 cm
				\centerline{\textit{#3}}
			\fi
			\if\relax\detokenize{#4}\relax\else
				\vskip 0.15 cm
				\centerline{\textit{#4}}
			\fi
			\vskip 0.15 cm
		\end{minipage}
	\end{tcolorbox}
}
%====================== BỔ SUNG CHO BÀI HỌC, CHUYÊN ĐỀ ======================
\usepackage[shortlabels]{enumitem}
\usepackage{tabularx}
%--- Danh sách gọn: \begin{enumerate}[gon] ... \end{enumerate}
\SetEnumitemKey{gon}{leftmargin=*,itemsep=1pt,parsep=0pt,topsep=3pt}
%--- Khung lý thuyết: \begin{khung}{} ... \end{khung}; tiêu đề chỉ dùng khi thật ngắn ("Mẹo", "Lỗi thường gặp")
\newtcolorbox{khung}[1]{breakable,enhanced,colframe=gold1,colback=gold5!45,
	boxrule=0.4mm,fonttitle=\bfseries\sffamily,coltitle=white,colbacktitle=gold1,
	attach boxed title to top left={xshift=6mm,yshift=-3mm},
	boxed title style={arc=1mm,boxrule=0.3mm},title={#1}}
%--- Tiêu đề mục nhỏ trong một dạng: \muc{Dấu hiệu}, \muc{Phương pháp}, \muc{Bài tự luyện}
\makeatletter
\newcommand{\muc}[1]{\par\medskip\noindent\indam{#1}\par\nobreak\smallskip\@afterindentfalse\@afterheading}
\makeatother
%--- Ô công thức cần nhớ (căn giữa, viền vàng, vừa khít nội dung)
%     \begin{center}\congthuc{$Q(n)=(n-1)!$}\end{center}
\newtcbox{\congthuc}{on line,colframe=gold1,colback=gold5!45,boxrule=0.5pt,arc=2pt,
	left=8pt,right=8pt,top=4pt,bottom=4pt}
%--- Công tắc lời giải bài tự luyện (ví dụ mẫu luôn dùng \loigiai)
%     \loigiaiTL[{đáp số}]{lời giải đầy đủ}
%     \TLfull -> in lời giải đầy đủ;  \TLdapso -> chỉ in một dòng "Đáp số";  \TLan -> ẩn hoàn toàn
%     (\CheDo là dethi thì mọi lời giải và đáp số đều ẩn)
\makeatletter
\newif\ifTLhien   \TLhientrue
\newif\ifTLdapso  \TLdapsofalse
\newcommand{\TLfull}{\TLhientrue\TLdapsofalse}
\newcommand{\TLdapso}{\TLhientrue\TLdapsotrue}
\newcommand{\TLan}{\TLhienfalse}
\newcommand{\loigiaiTL}[2][]{%
	\ifTLhien
		\ifTLdapso
			\def\@tlarg{#1}%
			\ifx\@tlarg\@empty\else
				\ifdefstring{\kindOPTN}{dethi}{}{\par\noindent{\bfseries\color{gold1}Đáp số:} #1\par}%
			\fi
		\else
			\loigiai{#2}%
		\fi
	\fi}
\makeatother
\TLdapso% mặc định: chỉ in đáp số bài tự luyện
%====================== BẢNG ĐÚNG/SAI (\choiceTFt) ======================
%  Hàng tiêu đề nền nâu đồng chữ trắng (Phát biểu | Đ | S), các hàng tô xen kẽ kem / trắng,
%  đường kẻ vàng kim, nhãn a) b) c) d) màu nâu đồng. Bản có đáp án: cột được chọn có vòng tròn
%  nâu đồng với dấu tích trắng (câu đúng tích ở cột Đúng, câu sai tích ở cột Sai).
%  Tiêu đề ghi đủ "Đúng | Sai": \OPTN{viettat=0} sau \input (cột Đúng/Sai rộng ra, phát biểu hẹp lại).
%  Dùng bảng gốc của ex_test: \def\BangDSgoc{1} trước \input{preamble-eduzy-v2}.
\providecommand{\BangDSgoc}{0}
\ifnum\BangDSgoc=0
\usepackage{colortbl,pifont}
\colorlet{ezDSke}{gold2!60!gold1}% đường kẻ bảng (nét 0,4pt như ex_test để không đổi chiều cao bảng)
\colorlet{ezDSle}{gold5!60}% hàng a), c)
\colorlet{ezDSchan}{white}% hàng b), d)
%--- dấu chọn ở cột Đúng/Sai
%    (bề rộng 0, căn giữa ô: không làm đổi bề rộng cột của ex_test, không báo tràn lề)
\newcommand{\EZdauchon}{\makebox[0pt]{\tikz[baseline=-1ex]\node[circle,fill=gold1,inner sep=0pt,
	minimum size=1.2em,text=white,font=\small]{\ding{51}};}}
\def\tickT{\EZdauchon}\def\tickF{\EZdauchon}
\AtBeginDocument{\sbox{\ccTFTab}{\EZdauchon}}
%--- nhãn a) b) c) d) màu nâu đồng (ở bản lời giải, nhãn câu đúng vẫn đóng khung như ex_test)
\renewcommand{\FalseTF}{\stepcounter{dapan}\textbf{\color{gold1}\DapAnTF\sepTF}\ignorespaces}
%--- trước khi dựng bảng: kẻ màu vàng kim (bề rộng cột giữ nguyên như ex_test); sau bảng trả lại màu đen
\newcount\EZhangTF
\newcommand{\EZbatdauTF}{%
	\global\EZhangTF=0\relax
	\arrayrulecolor{ezDSke}}
\newcommand{\EZketthucTF}{\arrayrulecolor{black}}
%--- hàng tiêu đề nền nâu đồng, chữ trắng
\makeatletter
\newcommand{\EZo@TF}[1]{\centering\bfseries\color{white}#1}
\newcommand{\EZtieudeTF}{%
	\ifcase\cotDS
		\ifnum\cotSTT=1\relax
			\def\rowHeader{\hline\rowcolor{gold1}\EZo@TF{TT}&\EZo@TF{\phatbieu}\tabularnewline}%
		\else
			\def\rowHeader{\hline\rowcolor{gold1}\EZo@TF{\phatbieu}\tabularnewline}%
		\fi
	\or
		\ifnum\cotSTT=1\relax
			\def\rowHeader{\hline\rowcolor{gold1}\EZo@TF{TT}&\EZo@TF{\phatbieu}&\EZo@TF{\Dung}\tabularnewline}%
		\else
			\def\rowHeader{\hline\rowcolor{gold1}\EZo@TF{\phatbieu}&\EZo@TF{\Dung}\tabularnewline}%
		\fi
	\or
		\ifnum\cotSTT=1\relax
			\def\rowHeader{\hline\rowcolor{gold1}\EZo@TF{TT}&\EZo@TF{\phatbieu}&\EZo@TF{\Dung}&\EZo@TF{\Sai}\tabularnewline}%
		\else
			\def\rowHeader{\hline\rowcolor{gold1}\EZo@TF{\phatbieu}&\EZo@TF{\Dung}&\EZo@TF{\Sai}\tabularnewline}%
		\fi
	\fi}
%--- tô xen kẽ từng hàng (gọi ngay sau \hline đầu mỗi hàng)
\newcommand{\EZmauhangTF}{%
	\noalign{\global\advance\EZhangTF\@ne}%
	\ifodd\EZhangTF\rowcolor{ezDSle}\else\rowcolor{ezDSchan}\fi}
%--- gắn vào \choiceTFt / \writeChoiceTFt của ex_test
\xpretocmd{\choiceTFt}{\EZbatdauTF}{}{\typeout{[preamble-eduzy-v2] Không gắn được bảng Đúng/Sai (đầu bảng)}}
\xpatchcmd{\choiceTFt}{\ifnum\dongPB=0\def\rowHeader{}\fi}{\EZtieudeTF\ifnum\dongPB=0\def\rowHeader{}\fi}
	{}{\typeout{[preamble-eduzy-v2] Không gắn được bảng Đúng/Sai (hàng tiêu đề)}}
\xpatchcmd{\writeChoiceTFt}{\hline\xdef\dotEX{\sdotEX}}{\hline\EZmauhangTF\xdef\dotEX{\sdotEX}}
	{}{\typeout{[preamble-eduzy-v2] Không gắn được bảng Đúng/Sai (tô hàng)}}
\xpatchcmd{\writeChoiceTFt}{{\end{longtable}}}{{\end{longtable}\EZketthucTF}}
	{}{\typeout{[preamble-eduzy-v2] Không gắn được bảng Đúng/Sai (cuối longtable)}}
\xpatchcmd{\writeChoiceTFt}{{\end{xtabular}\end{center}\medskip\noindent}}{{\end{xtabular}\end{center}\EZketthucTF\medskip\noindent}}
	{}{\typeout{[preamble-eduzy-v2] Không gắn được bảng Đúng/Sai (cuối xtabular)}}
\fi
%=====================================================================
\renewcommand{\parallel}{\mathrel{/\!\!/}}

%     %=====================================================================
% preamble-vatli-hinh.tex -- BỘ MACRO HÌNH VẬT LÍ (dùng chung cho mọi đề Vật lí)
% Gọi ngay sau preamble chính:
%     \input{preamble-eduzy-v2}
%     \input{preamble-vatli-hinh}
% Tên file bắt đầu bằng "preamble" nên dong-goi.sh tự chép vào zip thành phẩm;
% tools/so_hinh.py đọc file này khi dựng ảnh so hình.
% Lưu ý màu: preamble-eduzy-v2 đổi blue, cyan, teal sang tông vàng đồng (gold1, gold2).
% Cần màu xanh thật thì dùng VLxanh (hoặc khai báo \definecolor cục bộ).
% Rút từ đề mẫu VatLi2026-CaoBang-SoGD-Lan1 (thầy duyệt 06/11/2026).
%=====================================================================

%--- Mạch điện: circuitikz kiểu châu Âu (điện trở hình chữ nhật...)
\usepackage[european]{circuitikz}

%--- Màu dùng chung
\definecolor{VLdo}{RGB}{214,72,60}
\definecolor{VLxanh}{RGB}{60,100,200}
\definecolor{VLdong}{RGB}{200,128,70}
\definecolor{VLgo}{RGB}{186,132,84}
\definecolor{VLla}{RGB}{150,200,150}

%--- Mũi tên: VLmui (mũi tên ở đầu đường), VLgiua=<vị trí> (mũi tên giữa đường, mặc định 0,55)
\tikzset{
	VLmui/.style={-{Stealth[length=4.5pt,width=3.5pt]}},
	VLgiua/.style={postaction={decorate},decoration={markings,mark=at position #1 with {\arrow{Stealth[length=6.5pt,width=5pt]}}}},
	VLgiua/.default=0.55
}
% cọc nối dây của điện kế: \VLcoc{x}{y}{màu}
\newcommand{\VLcoc}[3]{%
	\begin{scope}[shift={(#1,#2)}]
		\shade[left color=#3!45!black,right color=#3!55!black,middle color=#3!80!white] (-0.17,0) rectangle (0.17,0.28);
		\shade[left color=#3!55!black,right color=#3!65!black,middle color=#3!60!white] (-0.14,0.28) rectangle (0.14,0.5);
		\fill[#3!55!black] (0,0.5) ellipse (0.14 and 0.045);
		\fill[black!85] (0,0.5) ellipse (0.055 and 0.02);
	\end{scope}}
% phích cắm có rãnh: \VLphich{x}{y}{màu}
\newcommand{\VLphich}[3]{%
	\begin{scope}[shift={(#1,#2)}]
		\shade[left color=#3!50!black,right color=#3!60!black,middle color=#3!55!white] (-0.12,0) rectangle (0.12,0.62);
		\draw[#3!45!black,very thin] foreach \y in {0.1,0.18,...,0.5}{(-0.12,\y)--(0.12,\y)};
		\shade[left color=#3!50!black,right color=#3!60!black,middle color=#3!55!white] (-0.08,0.62) rectangle (0.08,0.85);
		\fill[#3!45!black] (0,0.85) ellipse (0.08 and 0.03);
	\end{scope}}
% điện kế nhìn chính diện, gốc ở góc trái dưới, rộng 4, cao 4: \VLdienke{1: thang số | 2: thang -, +}
\newcommand{\VLdienke}[1]{%
	\fill[black!82,rounded corners=2pt] (0,0) rectangle (4,1.35);
	\fill[black!60] (0.05,1.22) rectangle (3.95,1.35);
	\fill[VLla,draw=black!60,rounded corners=1pt] (0.28,0.25) rectangle (3.72,1.02);
	\fill[gray!8,draw=gray!65,rounded corners=1.5pt] (0.08,1.35) rectangle (3.92,4.05);
	\fill[white,draw=gray!45] (0.28,1.5) rectangle (3.72,3.85);
	\fill[gray!18] (1.25,1.52) rectangle (2.75,2.05);
	\draw[gray!40,very thin] foreach \x in {1.3,1.4,...,2.7}{(\x,1.55)--(\x,2.0)};
	\draw[black!80] (2,1.35) ++(127:1.95) arc[start angle=127,end angle=53,radius=1.95];
	\foreach \g in {53,55.4667,...,127.1}{\draw[black!80,very thin] (2,1.35)++(\g:1.95)--++(\g:0.12);}
	\foreach \g in {53,65.3333,...,127.1}{\draw[black!85] (2,1.35)++(\g:1.95)--++(\g:0.22);}
	\ifnum#1=1
		\foreach \g/\s in {127/15,114.6667/10,102.3333/5,90/0,77.6667/5,65.3333/10,53/15}{\node[font=\fontsize{5}{5}\selectfont,inner sep=0] at ($(2,1.35)+(\g:2.37)$) {\s};}
		\node[font=\scriptsize\bfseries] at (2,2.75) {G};
	\else
		\node[font=\scriptsize] at ($(2,1.35)+(127:1.7)$) {$-$};
		\node[font=\scriptsize] at ($(2,1.35)+(53:1.7)$) {$+$};
	\fi
	\draw[VLdo,semithick] (2,1.8)--(2,3.42);
	\fill[black!85] (2,1.8) circle (0.13);
	\fill[gray!55] (2,1.8) circle (0.045);
	\draw[white,line width=1.2pt,opacity=0.7] (0.2,3.95)--(0.2,1.6);
}
% nam châm thẳng 3D: \VLnamcham{dài}{cao}{màu trái}{màu phải}{chữ trái}{chữ phải}
\newcommand{\VLnamcham}[6]{%
	\fill[#3!55!white] (0,#2)--(0.5*#1,#2)--(0.5*#1+0.22,#2+0.13)--(0.22,#2+0.13)--cycle;
	\fill[#4!55!white] (0.5*#1,#2)--(#1,#2)--(#1+0.22,#2+0.13)--(0.5*#1+0.22,#2+0.13)--cycle;
	\fill[#4!70!black] (#1,0)--(#1+0.22,0.13)--(#1+0.22,#2+0.13)--(#1,#2)--cycle;
	\shade[top color=#3!75!white,bottom color=#3!95!black] (0,0) rectangle (0.5*#1,#2);
	\shade[top color=#4!75!white,bottom color=#4!95!black] (0.5*#1,0) rectangle (#1,#2);
	\node[white,font=\scriptsize\bfseries] at (0.25*#1,0.5*#2) {#5};
	\node[white,font=\scriptsize\bfseries] at (0.75*#1,0.5*#2) {#6};
}
% lưới chấm (từ trường hướng ra) / lưới dấu nhân (hướng vào)
\newcommand{\VLcham}[2]{\fill (#1,#2) circle (1.6pt);}
\newcommand{\VLnhan}[2]{\draw[semithick] (#1,#2)++(-0.06,-0.06)--++(0.12,0.12) (#1,#2)++(-0.06,0.06)--++(0.12,-0.12);}
% thanh dẫn hình trụ nằm ngang / thẳng đứng (tâm, nửa dài)
\newcommand{\VLthanhngang}[3]{%
	\shade[top color=gray!60,bottom color=gray!85,middle color=gray!35,draw=black!85] (#1-#3,#2-0.1) rectangle (#1+#3,#2+0.1);
	\filldraw[fill=gray!35,draw=black!80] (#1+#3,#2) ellipse (0.04 and 0.1);}
\newcommand{\VLthanhdung}[3]{%
	\shade[left color=gray!85,right color=gray!60,middle color=gray!35,draw=black!85] (#1-0.1,#2-#3) rectangle (#1+0.1,#2+#3);
	\filldraw[fill=gray!35,draw=black!80] (#1,#2+#3) ellipse (0.1 and 0.035);}

% Tên file bắt đầu bằng "preamble" nên dong-goi.sh tự chép vào zip thành phẩm;
% tools/so_hinh.py đọc file này khi dựng ảnh so hình.
% Lưu ý màu: preamble-eduzy-v2 đổi blue, cyan, teal sang tông vàng đồng (gold1, gold2).
% Cần màu xanh thật thì dùng VLxanh (hoặc khai báo \definecolor cục bộ).
% Rút từ đề mẫu VatLi2026-CaoBang-SoGD-Lan1 (thầy duyệt 06/11/2026).
%=====================================================================

%--- Mạch điện: circuitikz kiểu châu Âu (điện trở hình chữ nhật...)
\usepackage[european]{circuitikz}

%--- Màu dùng chung
\definecolor{VLdo}{RGB}{214,72,60}
\definecolor{VLxanh}{RGB}{60,100,200}
\definecolor{VLdong}{RGB}{200,128,70}
\definecolor{VLgo}{RGB}{186,132,84}
\definecolor{VLla}{RGB}{150,200,150}

%--- Mũi tên: VLmui (mũi tên ở đầu đường), VLgiua=<vị trí> (mũi tên giữa đường, mặc định 0,55)
\tikzset{
	VLmui/.style={-{Stealth[length=4.5pt,width=3.5pt]}},
	VLgiua/.style={postaction={decorate},decoration={markings,mark=at position #1 with {\arrow{Stealth[length=6.5pt,width=5pt]}}}},
	VLgiua/.default=0.55
}
% cọc nối dây của điện kế: \VLcoc{x}{y}{màu}
\newcommand{\VLcoc}[3]{%
	\begin{scope}[shift={(#1,#2)}]
		\shade[left color=#3!45!black,right color=#3!55!black,middle color=#3!80!white] (-0.17,0) rectangle (0.17,0.28);
		\shade[left color=#3!55!black,right color=#3!65!black,middle color=#3!60!white] (-0.14,0.28) rectangle (0.14,0.5);
		\fill[#3!55!black] (0,0.5) ellipse (0.14 and 0.045);
		\fill[black!85] (0,0.5) ellipse (0.055 and 0.02);
	\end{scope}}
% phích cắm có rãnh: \VLphich{x}{y}{màu}
\newcommand{\VLphich}[3]{%
	\begin{scope}[shift={(#1,#2)}]
		\shade[left color=#3!50!black,right color=#3!60!black,middle color=#3!55!white] (-0.12,0) rectangle (0.12,0.62);
		\draw[#3!45!black,very thin] foreach \y in {0.1,0.18,...,0.5}{(-0.12,\y)--(0.12,\y)};
		\shade[left color=#3!50!black,right color=#3!60!black,middle color=#3!55!white] (-0.08,0.62) rectangle (0.08,0.85);
		\fill[#3!45!black] (0,0.85) ellipse (0.08 and 0.03);
	\end{scope}}
% điện kế nhìn chính diện, gốc ở góc trái dưới, rộng 4, cao 4: \VLdienke{1: thang số | 2: thang -, +}
\newcommand{\VLdienke}[1]{%
	\fill[black!82,rounded corners=2pt] (0,0) rectangle (4,1.35);
	\fill[black!60] (0.05,1.22) rectangle (3.95,1.35);
	\fill[VLla,draw=black!60,rounded corners=1pt] (0.28,0.25) rectangle (3.72,1.02);
	\fill[gray!8,draw=gray!65,rounded corners=1.5pt] (0.08,1.35) rectangle (3.92,4.05);
	\fill[white,draw=gray!45] (0.28,1.5) rectangle (3.72,3.85);
	\fill[gray!18] (1.25,1.52) rectangle (2.75,2.05);
	\draw[gray!40,very thin] foreach \x in {1.3,1.4,...,2.7}{(\x,1.55)--(\x,2.0)};
	\draw[black!80] (2,1.35) ++(127:1.95) arc[start angle=127,end angle=53,radius=1.95];
	\foreach \g in {53,55.4667,...,127.1}{\draw[black!80,very thin] (2,1.35)++(\g:1.95)--++(\g:0.12);}
	\foreach \g in {53,65.3333,...,127.1}{\draw[black!85] (2,1.35)++(\g:1.95)--++(\g:0.22);}
	\ifnum#1=1
		\foreach \g/\s in {127/15,114.6667/10,102.3333/5,90/0,77.6667/5,65.3333/10,53/15}{\node[font=\fontsize{5}{5}\selectfont,inner sep=0] at ($(2,1.35)+(\g:2.37)$) {\s};}
		\node[font=\scriptsize\bfseries] at (2,2.75) {G};
	\else
		\node[font=\scriptsize] at ($(2,1.35)+(127:1.7)$) {$-$};
		\node[font=\scriptsize] at ($(2,1.35)+(53:1.7)$) {$+$};
	\fi
	\draw[VLdo,semithick] (2,1.8)--(2,3.42);
	\fill[black!85] (2,1.8) circle (0.13);
	\fill[gray!55] (2,1.8) circle (0.045);
	\draw[white,line width=1.2pt,opacity=0.7] (0.2,3.95)--(0.2,1.6);
}
% nam châm thẳng 3D: \VLnamcham{dài}{cao}{màu trái}{màu phải}{chữ trái}{chữ phải}
\newcommand{\VLnamcham}[6]{%
	\fill[#3!55!white] (0,#2)--(0.5*#1,#2)--(0.5*#1+0.22,#2+0.13)--(0.22,#2+0.13)--cycle;
	\fill[#4!55!white] (0.5*#1,#2)--(#1,#2)--(#1+0.22,#2+0.13)--(0.5*#1+0.22,#2+0.13)--cycle;
	\fill[#4!70!black] (#1,0)--(#1+0.22,0.13)--(#1+0.22,#2+0.13)--(#1,#2)--cycle;
	\shade[top color=#3!75!white,bottom color=#3!95!black] (0,0) rectangle (0.5*#1,#2);
	\shade[top color=#4!75!white,bottom color=#4!95!black] (0.5*#1,0) rectangle (#1,#2);
	\node[white,font=\scriptsize\bfseries] at (0.25*#1,0.5*#2) {#5};
	\node[white,font=\scriptsize\bfseries] at (0.75*#1,0.5*#2) {#6};
}
% lưới chấm (từ trường hướng ra) / lưới dấu nhân (hướng vào)
\newcommand{\VLcham}[2]{\fill (#1,#2) circle (1.6pt);}
\newcommand{\VLnhan}[2]{\draw[semithick] (#1,#2)++(-0.06,-0.06)--++(0.12,0.12) (#1,#2)++(-0.06,0.06)--++(0.12,-0.12);}
% thanh dẫn hình trụ nằm ngang / thẳng đứng (tâm, nửa dài)
\newcommand{\VLthanhngang}[3]{%
	\shade[top color=gray!60,bottom color=gray!85,middle color=gray!35,draw=black!85] (#1-#3,#2-0.1) rectangle (#1+#3,#2+0.1);
	\filldraw[fill=gray!35,draw=black!80] (#1+#3,#2) ellipse (0.04 and 0.1);}
\newcommand{\VLthanhdung}[3]{%
	\shade[left color=gray!85,right color=gray!60,middle color=gray!35,draw=black!85] (#1-0.1,#2-#3) rectangle (#1+0.1,#2+#3);
	\filldraw[fill=gray!35,draw=black!80] (#1,#2+#3) ellipse (0.1 and 0.035);}

% Tên file bắt đầu bằng "preamble" nên dong-goi.sh tự chép vào zip thành phẩm;
% tools/so_hinh.py đọc file này khi dựng ảnh so hình.
% Lưu ý màu: preamble-eduzy-v2 đổi blue, cyan, teal sang tông vàng đồng (gold1, gold2).
% Cần màu xanh thật thì dùng VLxanh (hoặc khai báo \definecolor cục bộ).
% Rút từ đề mẫu VatLi2026-CaoBang-SoGD-Lan1 (thầy duyệt 06/11/2026).
%=====================================================================

%--- Mạch điện: circuitikz kiểu châu Âu (điện trở hình chữ nhật...)
\usepackage[european]{circuitikz}

%--- Màu dùng chung
\definecolor{VLdo}{RGB}{214,72,60}
\definecolor{VLxanh}{RGB}{60,100,200}
\definecolor{VLdong}{RGB}{200,128,70}
\definecolor{VLgo}{RGB}{186,132,84}
\definecolor{VLla}{RGB}{150,200,150}

%--- Mũi tên: VLmui (mũi tên ở đầu đường), VLgiua=<vị trí> (mũi tên giữa đường, mặc định 0,55)
\tikzset{
	VLmui/.style={-{Stealth[length=4.5pt,width=3.5pt]}},
	VLgiua/.style={postaction={decorate},decoration={markings,mark=at position #1 with {\arrow{Stealth[length=6.5pt,width=5pt]}}}},
	VLgiua/.default=0.55
}
% cọc nối dây của điện kế: \VLcoc{x}{y}{màu}
\newcommand{\VLcoc}[3]{%
	\begin{scope}[shift={(#1,#2)}]
		\shade[left color=#3!45!black,right color=#3!55!black,middle color=#3!80!white] (-0.17,0) rectangle (0.17,0.28);
		\shade[left color=#3!55!black,right color=#3!65!black,middle color=#3!60!white] (-0.14,0.28) rectangle (0.14,0.5);
		\fill[#3!55!black] (0,0.5) ellipse (0.14 and 0.045);
		\fill[black!85] (0,0.5) ellipse (0.055 and 0.02);
	\end{scope}}
% phích cắm có rãnh: \VLphich{x}{y}{màu}
\newcommand{\VLphich}[3]{%
	\begin{scope}[shift={(#1,#2)}]
		\shade[left color=#3!50!black,right color=#3!60!black,middle color=#3!55!white] (-0.12,0) rectangle (0.12,0.62);
		\draw[#3!45!black,very thin] foreach \y in {0.1,0.18,...,0.5}{(-0.12,\y)--(0.12,\y)};
		\shade[left color=#3!50!black,right color=#3!60!black,middle color=#3!55!white] (-0.08,0.62) rectangle (0.08,0.85);
		\fill[#3!45!black] (0,0.85) ellipse (0.08 and 0.03);
	\end{scope}}
% điện kế nhìn chính diện, gốc ở góc trái dưới, rộng 4, cao 4: \VLdienke{1: thang số | 2: thang -, +}
\newcommand{\VLdienke}[1]{%
	\fill[black!82,rounded corners=2pt] (0,0) rectangle (4,1.35);
	\fill[black!60] (0.05,1.22) rectangle (3.95,1.35);
	\fill[VLla,draw=black!60,rounded corners=1pt] (0.28,0.25) rectangle (3.72,1.02);
	\fill[gray!8,draw=gray!65,rounded corners=1.5pt] (0.08,1.35) rectangle (3.92,4.05);
	\fill[white,draw=gray!45] (0.28,1.5) rectangle (3.72,3.85);
	\fill[gray!18] (1.25,1.52) rectangle (2.75,2.05);
	\draw[gray!40,very thin] foreach \x in {1.3,1.4,...,2.7}{(\x,1.55)--(\x,2.0)};
	\draw[black!80] (2,1.35) ++(127:1.95) arc[start angle=127,end angle=53,radius=1.95];
	\foreach \g in {53,55.4667,...,127.1}{\draw[black!80,very thin] (2,1.35)++(\g:1.95)--++(\g:0.12);}
	\foreach \g in {53,65.3333,...,127.1}{\draw[black!85] (2,1.35)++(\g:1.95)--++(\g:0.22);}
	\ifnum#1=1
		\foreach \g/\s in {127/15,114.6667/10,102.3333/5,90/0,77.6667/5,65.3333/10,53/15}{\node[font=\fontsize{5}{5}\selectfont,inner sep=0] at ($(2,1.35)+(\g:2.37)$) {\s};}
		\node[font=\scriptsize\bfseries] at (2,2.75) {G};
	\else
		\node[font=\scriptsize] at ($(2,1.35)+(127:1.7)$) {$-$};
		\node[font=\scriptsize] at ($(2,1.35)+(53:1.7)$) {$+$};
	\fi
	\draw[VLdo,semithick] (2,1.8)--(2,3.42);
	\fill[black!85] (2,1.8) circle (0.13);
	\fill[gray!55] (2,1.8) circle (0.045);
	\draw[white,line width=1.2pt,opacity=0.7] (0.2,3.95)--(0.2,1.6);
}
% nam châm thẳng 3D: \VLnamcham{dài}{cao}{màu trái}{màu phải}{chữ trái}{chữ phải}
\newcommand{\VLnamcham}[6]{%
	\fill[#3!55!white] (0,#2)--(0.5*#1,#2)--(0.5*#1+0.22,#2+0.13)--(0.22,#2+0.13)--cycle;
	\fill[#4!55!white] (0.5*#1,#2)--(#1,#2)--(#1+0.22,#2+0.13)--(0.5*#1+0.22,#2+0.13)--cycle;
	\fill[#4!70!black] (#1,0)--(#1+0.22,0.13)--(#1+0.22,#2+0.13)--(#1,#2)--cycle;
	\shade[top color=#3!75!white,bottom color=#3!95!black] (0,0) rectangle (0.5*#1,#2);
	\shade[top color=#4!75!white,bottom color=#4!95!black] (0.5*#1,0) rectangle (#1,#2);
	\node[white,font=\scriptsize\bfseries] at (0.25*#1,0.5*#2) {#5};
	\node[white,font=\scriptsize\bfseries] at (0.75*#1,0.5*#2) {#6};
}
% lưới chấm (từ trường hướng ra) / lưới dấu nhân (hướng vào)
\newcommand{\VLcham}[2]{\fill (#1,#2) circle (1.6pt);}
\newcommand{\VLnhan}[2]{\draw[semithick] (#1,#2)++(-0.06,-0.06)--++(0.12,0.12) (#1,#2)++(-0.06,0.06)--++(0.12,-0.12);}
% thanh dẫn hình trụ nằm ngang / thẳng đứng (tâm, nửa dài)
\newcommand{\VLthanhngang}[3]{%
	\shade[top color=gray!60,bottom color=gray!85,middle color=gray!35,draw=black!85] (#1-#3,#2-0.1) rectangle (#1+#3,#2+0.1);
	\filldraw[fill=gray!35,draw=black!80] (#1+#3,#2) ellipse (0.04 and 0.1);}
\newcommand{\VLthanhdung}[3]{%
	\shade[left color=gray!85,right color=gray!60,middle color=gray!35,draw=black!85] (#1-0.1,#2-#3) rectangle (#1+0.1,#2+#3);
	\filldraw[fill=gray!35,draw=black!80] (#1,#2+#3) ellipse (0.1 and 0.035);}

\begin{document}
	\khungtieude{ĐỀ THI THỬ THPT YÊN HÒA -- LẦN 1}{Môn Vật lí}{(Đề thi gồm: 18 câu trắc nghiệm, 04 câu đúng-sai, 06 câu trả lời ngắn)}{Thời gian làm bài: 50 phút (không kể thời gian phát đề).}

	\noindent\textbf{Cho biết:} $\pi=3{,}14$; $T(\text{K})=t\,(^\circ\text{C})+273$; $R=8{,}31$~J$\cdot$mol$^{-1}\cdot$K$^{-1}$; $N_A=6{,}02\cdot 10^{23}$~hạt/mol.

	\NeedspaceCau{7.6\baselineskip}
	\noindent
	\subsection*{\indam{PHẦN I. Câu trắc nghiệm nhiều phương án lựa chọn. Thí sinh trả lời từ câu 1 đến câu 18. Mỗi câu hỏi thí sinh chỉ chọn một phương án.}}
	\setcounter{ex}{0}
	\Opensolutionfile{ans}[ans/VatLi2026-HaNoi-YenHoa-Lan1-1]

	\begin{ex}%[VL12-C1-TH]
		Nhiệt độ vào một ngày mùa hè ở Thành phố Hồ Chí Minh là $35\,^\circ\text{C}$. Nhiệt độ đó tương ứng với bao nhiêu độ Fahrenheit?
		\choice
		{$67\,^\circ\text{F}$}
		{$59\,^\circ\text{F}$}
		{$76\,^\circ\text{F}$}
		{\True $95\,^\circ\text{F}$}
		\loigiai{
			$t_{\text{F}}=1{,}8\,t_{\text{C}}+32=1{,}8\cdot 35+32=95\,^\circ\text{F}$.
		}
	\end{ex}

	\NeedspaceCau{6.8\baselineskip}
	\begin{ex}%[VL12-C1-TH]
		Một bệnh viện cần khử trùng dụng cụ y tế bằng hơi nước. Họ đã sử dụng nồi hơi để làm nóng nước từ nhiệt độ phòng $25\,^\circ\text{C}$ đến khi nước chuyển hoàn toàn thành hơi ở $100\,^\circ\text{C}$. Biết rằng mỗi lần khử trùng, bệnh viện cần đun sôi $10$~kg nước. Nhiệt hóa hơi riêng của nước ở $100\,^\circ\text{C}$ là $2{,}26\cdot 10^6$~J/kg và nhiệt dung riêng của nước là $4200$~J/(kg.K). Nhiệt lượng cần thiết để làm cho $10$~kg nước ở $25\,^\circ\text{C}$ chuyển hoàn toàn thành hơi nước ở $100\,^\circ\text{C}$ là
		\choice
		{$22{,}6$~MJ}
		{$16{,}95$~MJ}
		{$31{,}5$~MJ}
		{\True $25{,}75$~MJ}
		\loigiai{
			$Q=mc\Delta t+Lm=10\cdot 4200\cdot 75+2{,}26\cdot 10^6\cdot 10=2{,}575\cdot 10^7~\text{J}=25{,}75$~MJ.
		}
	\end{ex}

	\NeedspaceCau{7.3\baselineskip}
	\begin{ex}%[VL12-C2-NB]
		Người ta dùng tay ấn pit-tông để nén khí trong một xilanh kín. Hiện tượng nào sẽ xảy ra với khí trong xilanh?
		\choice
		{\True Thể tích bình chứa khí giảm. Áp suất khí trong bình tăng}
		{Thể tích bình chứa khí giảm. Áp suất khí trong bình giảm}
		{Thể tích bình chứa khí tăng. Áp suất khí trong bình giảm}
		{Thể tích bình chứa khí tăng. Áp suất khí trong bình tăng}
		\loigiai{
			Nén khí làm thể tích giảm, mật độ phân tử khí tăng nên áp suất tăng.
		}
	\end{ex}

	\NeedspaceCau{5.6\baselineskip}
	\begin{ex}%[VL12-C3-TH]
		Một dây dẫn thẳng, cứng, dài $20$~cm, có khối lượng $50$~g được giữ nằm yên theo phương ngang trong một từ trường có độ lớn cảm ứng từ là $0{,}49$~T và có hướng nằm ngang, vuông góc với dây. Cường độ dòng điện chạy trong dây là bao nhiêu ampe để khi dây được thả ra thì nó vẫn nằm yên? Lấy $g=9{,}8~\text{m/s}^2$.
		\choice
		{\True $5$~A}
		{$2{,}5$~A}
		{$7{,}5$~A}
		{$10$~A}
		\loigiai{
			Lực từ cân bằng trọng lực: $BIl=mg\Rightarrow I=\dfrac{mg}{Bl}=\dfrac{0{,}05\cdot 9{,}8}{0{,}49\cdot 0{,}2}=5$~A.
		}
	\end{ex}

	\NeedspaceCau{3.9\baselineskip}
	\begin{ex}%[VL12-C1-NB]
		Dụng cụ nào sau đây \textbf{không} dùng để đo nhiệt độ?
		\choice
		{Nhiệt kế thủy ngân}
		{Nhiệt kế rượu}
		{Nhiệt kế điện tử}
		{\True Tốc kế}
		\loigiai{
			Tốc kế dùng để đo tốc độ.
		}
	\end{ex}

	\NeedspaceCau{4.9\baselineskip}
	\begin{ex}%[VL12-C1-NB]
		Những quá trình chuyển thể nào của kim loại đồng được vận dụng trong việc đúc tượng bằng đồng?
		\choice
		{Bay hơi và ngưng tụ}
		{\True Nóng chảy và đông đặc}
		{Bay hơi và đông đặc}
		{Nóng chảy và bay hơi}
		\loigiai{
			Đồng được nung nóng chảy, rót vào khuôn rồi để nguội đông đặc thành tượng.
		}
	\end{ex}

	\NeedspaceCau{5.9\baselineskip}
	\begin{ex}%[VL12-C1-TH]
		% KIỂM TRA: nhiệt nóng chảy riêng của nước đá thực tế là 3,4.10^5 J/kg (đề ghi 34.10^5 J/kg); các phương án tính theo 34.10^5 (B); với 3,4.10^5 thì Q = 1,778 MJ, không có phương án
		Một gia đình đang chuẩn bị bữa ăn ngoài trời vào mùa hè. Họ muốn làm tan và đun nóng $4$~kg nước đá từ $0\,^\circ\text{C}$ để có nước $25\,^\circ\text{C}$ cho thức uống của mình. Bạn hãy tính xem cần bao nhiêu nhiệt lượng để chuyển nước đá này thành nước ở $25\,^\circ\text{C}$. Biết rằng nhiệt nóng chảy riêng của nước đá là $34\cdot 10^5$~J/kg và nhiệt dung riêng của nước là $4180$~J/(kg.K).
		\choice
		{$18{,}96$~MJ}
		{\True $14{,}018$~MJ}
		{$16{,}94$~MJ}
		{$21{,}23$~MJ}
		\loigiai{
			$Q=\lambda m+mc\Delta t=34\cdot 10^5\cdot 4+4\cdot 4180\cdot 25=1{,}4018\cdot 10^7~\text{J}=14{,}018$~MJ.
		}
	\end{ex}

	\NeedspaceCau{12.1\baselineskip}
	\begin{ex}%[VL12-C3-TH]
		Xét đoạn dây dẫn thẳng, dài mang dòng điện đặt trong từ trường đều theo phương vuông góc với đường sức từ. Hình vẽ nào dưới đây biểu diễn đúng phương, chiều của lực từ tác dụng lên dây dẫn?
		\begin{center}
			\begin{tikzpicture}[line join=round,line cap=round,font=\footnotesize]
				% Hình a
				\begin{scope}
					\VLthanhngang{1.3}{0}{1.3}
					\draw[VLmui,thick,VLdo] (1.1,0)--(1.7,0);
					\node[VLdo] at (1.4,-0.32) {$I$};
					\draw[VLmui,thick,VLluc] (0.75,0.1)--(0.75,1.45);
					\node[VLluc] at (0.48,1.35) {$\vv{F}$};
					\draw[semithick,VLxanh!85!black] (1.8,1.0) circle (0.2) ++(-0.141,-0.141)--++(0.283,0.283) (1.8,1.0)++(-0.141,0.141)--++(0.283,-0.283);
					\node[VLxanh!85!black] at (1.8,1.48) {$\vv{B}$};
					\node at (1.3,-0.95) {Hình a};
				\end{scope}
				% Hình b
				\begin{scope}[xshift=3.9cm]
					\draw[semithick,VLdo] (0.5,1.1) circle (0.22);
					\fill[VLdo] (0.5,1.1) circle (0.06);
					\node[VLdo] at (0.12,1.5) {$I$};
					\draw[VLmui,thick,VLluc] (0.5,0.88)--(0.5,-0.3);
					\node[VLluc] at (0.78,-0.15) {$\vv{F}$};
					\draw[VLmui,thick,VLxanh] (1.0,1.55)--(2.4,1.55);
					\node[VLxanh!85!black] at (2.5,1.82) {$\vv{B}$};
					\node at (1.3,-0.95) {Hình b};
				\end{scope}
				% Hình c
				\begin{scope}[xshift=7.8cm]
					\draw[VLmui,thick,VLxanh] (0.5,0.42)--(0.5,1.6);
					\node[VLxanh!85!black] at (0.22,1.5) {$\vv{B}$};
					\draw[semithick,VLdo] (0.5,0.2) circle (0.22);
					\fill[VLdo] (0.5,0.2) circle (0.06);
					\node[VLdo] at (0.08,0.2) {$I$};
					\draw[VLmui,thick,VLluc] (0.72,0.2)--(2.1,0.2);
					\node[VLluc] at (2.1,0.48) {$\vv{F}$};
					\node at (1.3,-0.95) {Hình c};
				\end{scope}
				% Hình d
				\begin{scope}[xshift=11.7cm]
					\draw[semithick,VLxanh!85!black] (1.05,1.45) circle (0.2) ++(-0.141,-0.141)--++(0.283,0.283) (1.05,1.45)++(-0.141,0.141)--++(0.283,-0.283);
					\node[VLxanh!85!black] at (0.62,1.45) {$\vv{B}$};
					\VLthanhdung{2.0}{0.75}{0.95}
					\draw[VLmui,thick,VLdo] (2.0,0.0)--(2.0,-0.45);
					\node[VLdo] at (2.3,-0.2) {$I$};
					\draw[VLmui,thick,VLluc] (1.9,0.75)--(0.5,0.75);
					\node[VLluc] at (0.6,0.42) {$\vv{F}$};
					\node at (1.3,-0.95) {Hình d};
				\end{scope}
			\end{tikzpicture}
		\end{center}
		\choice
		{\True Hình a}
		{Hình b}
		{Hình c}
		{Hình d}
		\loigiai{
			Theo quy tắc bàn tay trái: Hình a đúng ($\vv{F}$ hướng lên); Hình b lực hướng lên, Hình c lực hướng sang trái, Hình d lực hướng sang phải.
		}
	\end{ex}

	\NeedspaceCau{11.9\baselineskip}
	\begin{ex}%[VL12-C2-VD]
		Hình bên phải là đồ thị mô tả sự biến đổi trạng thái của $1$~mol khí lí tưởng trong hệ tọa độ $(V;T)$. Đồ thị của sự biến đổi trạng thái trên trong hệ tọa độ $(p;T)$ tương ứng với hình nào dưới đây?
		\begin{center}
			\begin{tikzpicture}[line join=round,line cap=round,font=\footnotesize]
				\foreach \k/\ten in {0/a,1/b,2/c,3/d}{
					\begin{scope}[xshift=\k*3.15cm]
						\draw[VLmui,thick] (0,0)--(2.35,0) node[above,inner sep=2pt] {$T$};
						\draw[VLmui,thick] (0,0)--(0,2.25) node[left,inner sep=2pt] {$p$};
						\node[below left,inner sep=2pt] at (0,0) {$O$};
						\node at (1.15,-0.4) {\ten)};
					\end{scope}}
				% a)
				\draw[very thick,VLxanh] (0.4,2.0) .. controls (0.6,0.95) and (1.0,0.6) .. (1.65,0.55)--(0.4,0.55)--cycle;
				\node[left,inner sep=2pt] at (0.4,2.0) {1};
				\node[left,inner sep=2pt] at (0.4,0.55) {3};
				\node[right,inner sep=2pt] at (1.65,0.55) {2};
				% b)
				\begin{scope}[xshift=3.15cm]
					\draw[dashed,thin] (0,0)--(0.75,0.83);
					\draw[very thick,VLxanh] (0.75,0.83)--(1.7,1.87)--(1.7,0.83)--cycle;
					\node[above left,inner sep=1pt] at (0.75,0.83) {3};
					\node[above right,inner sep=1pt] at (1.7,1.87) {2};
					\node[right,inner sep=2pt] at (1.7,0.83) {1};
				\end{scope}
				% c)
				\begin{scope}[xshift=6.3cm]
					\draw[very thick,VLxanh] (0.4,2.0) .. controls (0.6,0.95) and (1.0,0.6) .. (1.65,0.55)--(0.4,0.55)--cycle;
					\node[left,inner sep=2pt] at (0.4,2.0) {3};
					\node[left,inner sep=2pt] at (0.4,0.55) {1};
					\node[right,inner sep=2pt] at (1.65,0.55) {2};
				\end{scope}
				% d)
				\begin{scope}[xshift=9.45cm]
					\draw[dashed,thin] (0,0)--(0.75,0.83);
					\draw[very thick,VLxanh] (0.75,0.83)--(1.7,1.87)--(1.7,0.83)--cycle;
					\node[above left,inner sep=1pt] at (0.75,0.83) {1};
					\node[above right,inner sep=1pt] at (1.7,1.87) {3};
					\node[right,inner sep=2pt] at (1.7,0.83) {2};
				\end{scope}
				% đồ thị (V;T) bên phải
				\draw[gray!60] (12.35,-0.55)--(12.35,2.3);
				\begin{scope}[xshift=12.85cm]
					\draw[VLmui,thick] (0,0)--(2.35,0) node[above,inner sep=2pt] {$T$};
					\draw[VLmui,thick] (0,0)--(0,2.25) node[left,inner sep=2pt] {$V$};
					\node[below left,inner sep=2pt] at (0,0) {$O$};
					\draw[dashed,thin] (0,0)--(0.75,0.83);
					\draw[very thick,VLxanh] (0.75,0.83)--(1.7,1.87)--(1.7,0.83)--cycle;
					\node[above left,inner sep=1pt] at (0.75,0.83) {1};
					\node[above right,inner sep=1pt] at (1.7,1.87) {2};
					\node[right,inner sep=2pt] at (1.7,0.83) {3};
				\end{scope}
			\end{tikzpicture}
		\end{center}
		\choice
		{Hình a}
		{Hình c}
		{\True Hình d}
		{Hình b}
		\loigiai{
			Trên đồ thị $(V;T)$: (1) $\to$ (2) đẳng áp, $T$ tăng; (2) $\to$ (3) đẳng nhiệt, $V$ giảm nên $p$ tăng; (3) $\to$ (1) đẳng tích, $T$ giảm nên $p$ giảm, đường kéo dài qua $O$ trên đồ thị $(p;T)$. Ứng với Hình d.
		}
	\end{ex}

	\NeedspaceCau{11.8\baselineskip}
	\begin{ex}%[VL12-C1-TH]
		Nồi áp suất là dụng cụ khá phổ biến trong mỗi gia đình. Khi nấu bằng nồi áp suất đồ ăn thường chín nhanh và nhừ hơn so với nấu bằng các nồi thông thường là do khi nước trong nồi áp suất sôi thì
		% giuhinh
		\par\noindent
		\begin{minipage}[c]{0.7\linewidth}
		\choice
		{áp suất trong nồi áp suất lớn hơn, nhiệt độ nhỏ hơn so với nồi thông thường}
		{áp suất trong nồi áp suất lớn hơn, nhiệt độ như nhau so với nồi thông thường}
		{\True áp suất và nhiệt độ trong nồi áp suất đều lớn hơn so với nồi thông thường}
		{áp suất trong nồi áp suất nhỏ hơn, nhiệt độ lớn hơn so với nồi thông thường}
		\end{minipage}\hfill
		\begin{minipage}[c]{0.28\linewidth}
			\centering
			\begin{tikzpicture}[x=0.0088cm,y=-0.0088cm,line join=round,line cap=round]
				% toạ độ theo điểm ảnh của ảnh gốc (nồi áp suất inox, nhìn thẳng từ phía trước)
				\def\than{(107,170)--(107,352) .. controls (107,378) and (120,399) .. (152,409) .. controls (175,415) and (205,417) .. (235,417)--(310,417) .. controls (340,417) and (370,415) .. (393,409) .. controls (425,399) and (438,378) .. (438,352)--(438,170)--cycle}
				% bóng đổ
				\fill[gray!22] (272,417) ellipse (150 and 8);
				\fill[gray!32] (272,416) ellipse (105 and 5);
				% quai cài hai bên: móc kẹp nắp, đai kim loại, bản lề, tay cầm nhựa đen
				\foreach \s/\t in {1/0,-1/545}{
					\begin{scope}[cm={\s,0,0,1,(\t,0)}]
						\fill[black!88] (28,219) .. controls (26,214) and (30,211) .. (40,209)--(97,197)--(108,197)--(110,240)--(92,241) .. controls (70,238) and (45,229) .. (28,223)--cycle;
						\fill[gray!60] (42,212)--(90,201.5)--(90,204)--(46,214)--cycle;
						\draw[gray!70!black,line width=1.8pt] (82,148) .. controls (78,165) and (84,178) .. (101,190);
						\draw[gray!20,line width=0.6pt] (81.5,149) .. controls (77.5,165) and (83.5,177) .. (100,189);
						\shade[left color=black!85,right color=gray!55] (98,183) rectangle (118,240);
						\draw[black!85,line width=0.3pt] (98,183) rectangle (118,240);
						\fill[black!88] (76,131) .. controls (76,128) and (78,127) .. (81,127)--(104,128) .. controls (107,128) and (107,131) .. (106,139)--(90,143)--(90,151)--(77,151)--cycle;
						\fill[gray!55] (81,130) rectangle (101,132.5);
					\end{scope}}
				% thân nồi inox bóng
				\begin{scope}
					\clip \than;
					\shade[left color=gray!62,right color=gray!58,middle color=white] (107,160) rectangle (438,420);
					\fill[white] (122,160) rectangle (196,420) (402,160) rectangle (426,420);
					\shade[left color=white,right color=gray!45] (180,160) rectangle (212,420);
					% các dải phản chiếu tối trên thân
					\foreach \a/\b/\c in {212/231/88,250/262/85,270/313/96,331/346/85,366/379/88}{
						\shade[left color=gray!55,right color=black!\c] (\a-4,160) rectangle (\a,420);
						\fill[black!\c] (\a,160) rectangle (\b,420);
						\shade[left color=black!\c,right color=gray!45] (\b,160) rectangle (\b+4,420);}
					\fill[gray!55] (231,160) rectangle (250,420) (313,160) rectangle (331,420) (346,160) rectangle (366,420);
					\fill[white,opacity=0.75] (262,160) rectangle (270,420) (318,160) rectangle (324,420) (383,160) rectangle (394,420);
					\fill[gray!70] (396,160) rectangle (399,420);
					\shade[left color=black!96,right color=black!96,middle color=black!70] (272,160) rectangle (311,420);
					\fill[gray!60,opacity=0.5] (285,250) .. controls (290,300) and (290,340) .. (286,378)--(296,378) .. controls (300,330) and (298,280) .. (293,250)--cycle;
					% đáy nồi (đế nhiều lớp, sáng hơn)
					\shade[top color=white,bottom color=gray!28] (100,378) .. controls (180,384) and (365,384) .. (445,378)--(445,420)--(100,420)--cycle;
					\fill[white,opacity=0.85] (100,366) .. controls (180,374) and (365,374) .. (445,366)--(445,380) .. controls (365,386) and (180,386) .. (100,380)--cycle;
					\draw[gray!55,line width=0.5pt] (100,366) .. controls (180,374) and (365,374) .. (445,366);
					\shade[top color=black!60,bottom color=gray!30,opacity=0.55] (107,160) rectangle (438,192);
				\end{scope}
				\draw[gray!60!black,line width=0.5pt] \than;
				% nắp: mặt nắp hình nón thấp (có dải phản chiếu tối), van nhỏ, vành nắp
				\shade[top color=gray!25,bottom color=white] (100,157) .. controls (150,100) and (225,84) .. (272,84) .. controls (320,84) and (395,100) .. (445,157)--cycle;
				\shade[top color=gray!40,bottom color=gray!15] (140,157)--(254,128) .. controls (264,125) and (281,125) .. (291,128)--(416,157)--cycle;
				\fill[black!78] (140,155)--(228,123)--(256,121)--(256,133)--(196,157)--cycle;
				\fill[gray!45] (150,155)--(232,128)--(250,127)--(165,157)--cycle;
				\fill[black!78] (289,115)--(318,121)--(416,155)--(398,157)--(289,128)--cycle;
				\fill[gray!50] (300,121)--(402,153)--(392,154)--(300,126)--cycle;
				\fill[black!82] (183,117)--(196,114)--(200,136)--(180,140)--(178,134)--(185,131)--cycle;
				\shade[top color=gray!45,bottom color=gray!55,middle color=white,rounded corners=1.6pt] (91,153) rectangle (454,177);
				\fill[black!72] (104,155.5) rectangle (441,158);
				\fill[black!60] (120,171)--(425,171)--(425,174)--(120,174)--cycle;
				\fill[black!55] (205,163) rectangle (250,171) (300,163) rectangle (322,171) (345,163) rectangle (395,171);
				\fill[white,opacity=0.9] (257,162) rectangle (270,173) (326,162) rectangle (338,173);
				\draw[gray!55!black,line width=0.4pt,rounded corners=1.6pt] (91,153) rectangle (454,177);
				% tay đòn hình cung bắc ngang nắp (thanh inox dẹt)
				\draw[gray!55,line width=0.7pt] (300,101) .. controls (345,106) and (395,120) .. (432,142);
				\draw[gray!60!black,line width=4pt] (100,146) .. controls (150,97) and (225,81) .. (272,81) .. controls (320,81) and (395,97) .. (445,146);
				\draw[gray!42,line width=2.9pt] (100,146) .. controls (150,97) and (225,81) .. (272,81) .. controls (320,81) and (395,97) .. (445,146);
				\draw[gray!10,line width=1.1pt] (108,137) .. controls (156,94) and (226,79.5) .. (272,79.5) .. controls (318,79.5) and (388,94) .. (436,137);
				\draw[white,line width=0.6pt] (150,104) .. controls (190,86) and (240,79) .. (272,79) .. controls (300,79) and (330,84) .. (360,93);
				\draw[black!70,line width=0.9pt] (106,149) .. controls (156,103) and (227,87.5) .. (272,87.5) .. controls (317,87.5) and (389,103) .. (439,149);
				% trục vít giữa nắp: đế, trục mạ sáng, đai ốc
				\fill[black!85] (246,151) .. controls (252,143) and (292,143) .. (299,151)--(299,157)--(246,157)--cycle;
				\shade[left color=gray!80,right color=gray!85,middle color=white] (258,100) rectangle (286,146);
				\fill[black!78] (262,100) rectangle (267,146) (277,108) rectangle (280,142);
				\fill[gray!65!black] (255,143) rectangle (289,147);
				\fill[black!88] (259,92)--(266,86)--(279,86)--(286,92)--(283,105)--(262,105)--cycle;
				% núm vặn nhựa đen
				\fill[black!90] (245,86) .. controls (250,79) and (294,79) .. (299,86)--cycle;
				\shade[left color=gray!55,right color=black!85,middle color=black!95] (253,26) rectangle (292,82);
				\shade[left color=gray!50,right color=gray!40,middle color=gray!5,opacity=0.5] (285,34) rectangle (291,78);
				\shade[top color=black!72,bottom color=black] (233,10) .. controls (232,8) and (312,8) .. (311,10) .. controls (316,14) and (316,22) .. (311,28)--(233,28) .. controls (228,22) and (228,14) .. (233,10)--cycle;
				\fill[gray!55,rounded corners=0.6pt] (240,12.5) rectangle (304,15);
			\end{tikzpicture}
		\end{minipage}
		\loigiai{
			Nồi áp suất kín nên áp suất hơi trong nồi lớn; áp suất càng lớn thì nhiệt độ sôi của nước càng cao, thức ăn chín nhanh hơn.
		}
	\end{ex}

	\NeedspaceCau{7.5\baselineskip}
	\begin{ex}%[VL12-C1-NB]
		Quy ước về dấu nào sau đây đúng với công thức $\Delta U=A+Q$ của định luật I nhiệt động lực học?
		\choice
		{Vật nhận công $A<0$, vật nhận nhiệt $Q<0$}
		{Vật thực hiện công $A<0$, vật truyền nhiệt $Q>0$}
		{\True Vật nhận công $A>0$, vật nhận nhiệt $Q>0$}
		{Vật thực hiện công $A>0$, vật truyền nhiệt $Q<0$}
		\loigiai{
			Trong $\Delta U=A+Q$: vật nhận công thì $A>0$, vật nhận nhiệt thì $Q>0$.
		}
	\end{ex}

	\NeedspaceCau{7.7\baselineskip}
	\begin{ex}%[VL12-C2-TH]
		Một lượng khí lí tưởng biến đổi trạng thái theo đồ thị như hình bên. Quá trình biến đổi từ trạng thái 1 đến trạng thái 2 là quá trình
		\par\noindent
		\begin{minipage}[c]{0.64\linewidth}
		\choice
		{\True giãn đẳng áp}
		{nén đẳng áp}
		{giãn đẳng nhiệt}
		{nén đẳng nhiệt}
		\end{minipage}\hfill
		\begin{minipage}[c]{0.34\linewidth}
			\centering
			\begin{tikzpicture}[line join=round,line cap=round,font=\footnotesize]
				\draw[VLmui,thick] (0,0)--(2.6,0) node[above,inner sep=2pt] {$T$};
				\draw[VLmui,thick] (0,0)--(0,2.4) node[right,inner sep=2pt] {$V$};
				\node[below left,inner sep=2pt] at (0,0) {$O$};
				\draw[dashed,thin] (0,0)--(0.5,0.6);
				\draw[very thick,VLdo,VLgiua=0.45] (0.5,0.6)--(1.9,2.25);
				\filldraw[fill=black] (0.5,0.6) circle (2pt) (1.9,2.25) circle (2pt);
				\node[right,inner sep=3pt] at (0.5,0.6) {(1)};
				\node[right,inner sep=3pt] at (1.9,2.25) {(2)};
			\end{tikzpicture}
		\end{minipage}
		\loigiai{
			Đường biểu diễn trong hệ $(V;T)$ kéo dài qua $O$ nên $V\sim T$: quá trình đẳng áp; $V$ tăng nên là giãn đẳng áp.
		}
	\end{ex}

	\NeedspaceCau{10.3\baselineskip}
	\begin{ex}%[VL12-C3-TH]
		Một dây dẫn mang dòng điện được đặt giữa hai cực của một nam châm chữ U. Lực từ tác dụng lên dây dẫn có hướng như thế nào trong hình vẽ bên?
		% giuhinh
		\par\noindent
		\begin{minipage}[c]{0.6\linewidth}
		\choice
		{Hướng B}
		{\True Hướng A}
		{Hướng C}
		{Hướng D}
		\end{minipage}\hfill
		\begin{minipage}[c]{0.38\linewidth}
			\centering
			\begin{tikzpicture}[x=0.0082cm,y=-0.0082cm,line join=round,line cap=round,font=\footnotesize]
				% toạ độ theo điểm ảnh của ảnh gốc; nam châm chữ U vẽ khối (độ dày lùi về phía sau, lên trên sang trái)
				\def\Utrai{(60,177)--(60,448) arc[start angle=180,end angle=90,radius=140]--(200,531) arc[start angle=90,end angle=180,radius=83]--(117,177)--cycle}
				\def\Uphai{(340,177)--(340,448) arc[start angle=0,end angle=90,radius=140]--(200,531) arc[start angle=90,end angle=0,radius=83]--(283,177)--cycle}
				% mạch điện: cạnh sau và cạnh phải (có nguồn) nằm sau nam châm
				\coordinate (BL) at (150,47); \coordinate (BR) at (463,47); \coordinate (FR) at (548,282); \coordinate (FL) at (233,282);
				\draw[thick] (BL)--(BR)--(499,146) (503,157)--(FR);
				\draw[very thick] (464,146)--(533,146);
				\draw[very thick] (478,157)--(522,157);
				% mặt bên (độ dày): quét nửa S trước, nửa N sau
				\foreach \k in {0,...,32}{\fill[VLxanh!78!white] [shift={(-23*\k/32,-60*\k/32)}] \Uphai;}
				\foreach \k in {0,...,32}{\fill[VLdo!72!black] [shift={(-23*\k/32,-60*\k/32)}] \Utrai;}
				% vệt sáng tối trên mặt trong chỗ uốn
				\begin{scope}
					\clip (117,448) arc[start angle=180,end angle=0,radius=83]--(283,388)--(260,388) arc[start angle=0,end angle=140,radius=83]--cycle;
					\shade[left color=VLdo!60!black,right color=VLdo!85!black] (110,430) rectangle (200,540);
					\shade[left color=VLxanh!52!white,right color=VLxanh!78!white] (177,380)--(283,380)--(283,540)--(200,531)--(177,471)--cycle;
				\end{scope}
				% nét viền mặt bên
				\draw[black!70,line width=0.5pt] (37,117)--(37,388) arc[start angle=180,end angle=159,radius=140]--(69.3,498.1);
				\draw[black!70,line width=0.5pt] (260,117)--(260,388) arc[start angle=0,end angle=136.3,radius=83];
				\draw[black!70,line width=0.5pt] (177,471)--(200,531);
				% mặt trên hai cực
				\filldraw[fill=VLdo!72!orange,draw=black!70,line width=0.5pt] (60,177)--(117,177)--(94,117)--(37,117)--cycle;
				\filldraw[fill=VLxanh!72!black,draw=black!70,line width=0.5pt] (283,177)--(340,177)--(317,117)--(260,117)--cycle;
				% mặt trước
				\shade[left color=VLdo!88!black,right color=VLdo!96!white] \Utrai;
				\shade[left color=VLxanh!50!white,right color=VLxanh!64!white] \Uphai;
				\fill[VLdo!80!orange] (60,177) rectangle (117,201);
				\fill[VLxanh!80!black] (283,177) rectangle (340,205);
				\draw[black!70,line width=0.5pt] \Utrai \Uphai (60,201)--(117,201) (283,205)--(340,205);
				\node[white,font=\bfseries] at (88.5,236) {N};
				\node[white,font=\bfseries] at (311.5,238) {S};
				% cạnh trái (đoạn dây nằm giữa hai cực) và cạnh trước của mạch
				\draw[thick] (BL)--(FL)--(FR);
				\draw[VLmui,thick,VLdo] (345,47)--(305,47);
				\node[VLdo,below,inner sep=3.5pt] at (330,47) {$I$};
				\draw[VLmui,thick,VLdo] (395,282)--(440,282);
				\node[VLdo,above,inner sep=2.5pt] at (437,282) {$I$};
				% các hướng A, B, C, D
				\coordinate (O) at (190,160);
				\draw[VLmui,semithick] (O)--(190,99);
				\draw[VLmui,semithick] (O)--(190,223);
				\draw[VLmui,semithick] (O)--(252,160);
				\draw[VLmui,semithick] (O)--(126,160);
				\node[font=\scriptsize\bfseries,inner sep=0pt] at (190,82) {A};
				\node[font=\scriptsize\bfseries,inner sep=0pt] at (190,242) {C};
				\node[font=\scriptsize\bfseries,inner sep=0pt] at (241,139) {B};
				\node[font=\scriptsize\bfseries,inner sep=0pt] at (138,139) {D};
			\end{tikzpicture}
		\end{minipage}
		\loigiai{
			Đường sức từ hướng từ cực N sang cực S (sang phải); dòng điện trong đoạn dây giữa hai cực hướng ra phía trước. Theo quy tắc bàn tay trái, lực từ hướng lên: hướng A.
		}
	\end{ex}

	\NeedspaceCau{9.7\baselineskip}
	\begin{ex}%[VL12-C1-NB]
		Mô hình động học phân tử có thể dùng để giải thích cấu trúc của các chất rắn, chất lỏng và chất khí. Nhận định nào sau đây \textbf{chưa đúng}?
		\choice
		{\True Các chất rắn có hình dạng và thể tích xác định vì các phân tử cấu tạo nên nó đứng yên tại chỗ}
		{Các chất khí dễ nén hơn các chất lỏng vì lực liên kết giữa các phân tử chất khí nhỏ hơn lực liên kết giữa các phân tử chất lỏng}
		{Các chất lỏng không có hình dạng xác định vì các phân tử cấu tạo nên chất lỏng vừa dao động quanh những vị trí cân bằng, vừa có thể dời từ vị trí này sang vị trí khác}
		{Các chất rắn, lỏng, khí được cấu tạo từ các nguyên tử, phân tử}
		\loigiai{
			Các phân tử chất rắn luôn dao động quanh vị trí cân bằng xác định, không đứng yên.
		}
	\end{ex}

	\NeedspaceCau{9.3\baselineskip}
	\begin{ex}%[VL12-C3-TH]
		Cho các hình vẽ dưới đây xác định hướng của vectơ cảm ứng từ tại tâm của dây dẫn tròn mang dòng điện. Chọn đáp án đúng.
		\begin{center}
			\begin{tikzpicture}[line join=round,line cap=round,font=\footnotesize]
				\foreach \k in {0,1,2,3}{
					\draw[thick,VLdong!80!black] (\k*3.6,0) circle (0.68);
					\node at (\k*3.6,-1.1) {Hình \the\numexpr\k+1\relax};}
				% Hình 1: ngược chiều kim đồng hồ, B hướng ra
				\draw[VLmui,thick,VLdo] (5:0.68) arc[start angle=5,end angle=55,radius=0.68];
				\node[VLdo] at (50:1.0) {$I$};
				\draw[semithick,VLxanh!85!black] (0.15,0) circle (0.2); \fill[VLxanh!85!black] (0.15,0) circle (0.065);
				\node[VLxanh!85!black] at (-0.3,0) {$\vv{B}$};
				% Hình 2: cùng chiều kim đồng hồ, B hướng ra
				\begin{scope}[xshift=3.6cm]
					\draw[VLmui,thick,VLdo] (-5:0.68) arc[start angle=-5,end angle=-55,radius=0.68];
					\node[VLdo] at (-45:1.0) {$I$};
					\draw[semithick,VLxanh!85!black] (0.15,0) circle (0.2); \fill[VLxanh!85!black] (0.15,0) circle (0.065);
					\node[VLxanh!85!black] at (-0.3,0) {$\vv{B}$};
				\end{scope}
				% Hình 3: ngược chiều kim đồng hồ, B hướng vào
				\begin{scope}[xshift=7.2cm]
					\draw[VLmui,thick,VLdo] (5:0.68) arc[start angle=5,end angle=55,radius=0.68];
					\node[VLdo] at (50:1.0) {$I$};
					\draw[semithick,VLxanh!85!black] (0.15,0) circle (0.2) ++(-0.141,-0.141)--++(0.283,0.283) (0.15,0)++(-0.141,0.141)--++(0.283,-0.283);
					\node[VLxanh!85!black] at (-0.3,0) {$\vv{B}$};
				\end{scope}
				% Hình 4: cùng chiều kim đồng hồ, B nằm trong mặt phẳng hướng xuống
				\begin{scope}[xshift=10.8cm]
					\draw[VLmui,thick,VLdo] (-5:0.68) arc[start angle=-5,end angle=-55,radius=0.68];
					\node[VLdo] at (-45:1.0) {$I$};
					\fill[VLxanh!85!black] (0.1,0.1) circle (0.07);
					\draw[VLmui,thick,VLxanh] (0.1,0.1)--(0.1,-0.58);
					\node[VLxanh!85!black] at (-0.3,0.1) {$\vv{B}$};
				\end{scope}
			\end{tikzpicture}
		\end{center}
		\choice
		{Hình 4}
		{Hình 2}
		{Hình 3}
		{\True Hình 1}
		\loigiai{
			Theo quy tắc nắm tay phải: dòng điện ngược chiều kim đồng hồ thì $\vv{B}$ tại tâm hướng ra (Hình 1 đúng, Hình 3 sai); dòng điện cùng chiều kim đồng hồ thì $\vv{B}$ hướng vào (Hình 2, Hình 4 sai).
		}
	\end{ex}

	\NeedspaceCau{6.5\baselineskip}
	\begin{ex}%[VL12-C1-NB]
		Trong các phát biểu sau đây về sự bay hơi và sự sôi của chất lỏng, phát biểu nào \textbf{sai}?
		\choice
		{Sự sôi chỉ xảy ra ở nhiệt độ sôi}
		{Sự sôi là sự hóa hơi xảy ra ở cả mặt thoáng và trong lòng khối chất lỏng}
		{Sự bay hơi là sự hóa hơi chỉ xảy ra ở mặt thoáng của khối chất lỏng}
		{\True Sự bay hơi chỉ xảy ra ở một nhiệt độ nhất định}
		\loigiai{
			Sự bay hơi xảy ra ở mọi nhiệt độ.
		}
	\end{ex}

	\NeedspaceCau{6.2\baselineskip}
	\begin{ex}%[VL12-C2-NB]
		Các phân tử khí lí tưởng có các tính chất nào sau đây?
		\choice
		{Chuyển động không ngừng, tương tác hút hoặc đẩy với nhau}
		{Được coi là chất điểm, chuyển động không ngừng, tương tác hút hoặc đẩy với nhau}
		{Được coi là chất điểm, tương tác hút hoặc đẩy với nhau}
		{\True Được coi là chất điểm và chuyển động không ngừng}
		\loigiai{
			Phân tử khí lí tưởng được coi là chất điểm, chuyển động hỗn loạn không ngừng, chỉ tương tác khi va chạm.
		}
	\end{ex}

	\NeedspaceCau{8.2\baselineskip}
	\begin{ex}%[VL12-C2-TH]
		Một khối khí xác định chứa trong một xilanh kín được đậy bởi một pit-tông có thể chuyển động không ma sát với xilanh. Khi nung nóng đẳng áp lượng khí nói trên, khoảng cách trung bình giữa các phân tử khí sẽ
		\choice
		{giảm tới khoảng cách nhỏ nhất rồi tăng lên}
		{\True tăng lên}
		{giảm đi}
		{không đổi}
		\loigiai{
			Nung nóng đẳng áp: $V\sim T$, thể tích tăng nên khoảng cách trung bình giữa các phân tử tăng.
		}
	\end{ex}
	\Closesolutionfile{ans}

	\NeedspaceCau[15.4\baselineskip]{30.2\baselineskip}
	\subsection*{\indam{PHẦN II. Câu trắc nghiệm đúng sai. Thí sinh trả lời từ câu 1 đến câu 4. Trong mỗi ý a), b), c), d) ở mỗi câu, thí sinh chọn đúng hoặc sai.}}
	\setcounter{ex}{0}
	\Opensolutionfile{ans}[ans/VatLi2026-HaNoi-YenHoa-Lan1-2]

	\begin{ex}%[VL12-C3-VD]
		% tachDS
		Trong một dự án nghiên cứu tàu hỏa tốc độ cao, các kĩ sư tạo ra một mô hình tàu mini chạy trên đệm từ trường. Con tàu được mô phỏng bằng một thanh dẫn kim loại MN đặt nằm ngang trên đường ray đặc biệt (hình vẽ).
		\par\noindent
		\begin{minipage}[c]{0.62\linewidth}
		Thông số kĩ thuật như sau:\\
		+ Thanh kim loại MN có khối lượng trên mỗi đơn vị chiều dài là $10$~g/m, đặt nằm ngang.\\
		+ Hệ thống nam châm dưới đường ray tạo ra một từ trường đều hướng nằm ngang, vuông góc với thanh MN, có độ lớn cảm ứng từ $B=0{,}1$~T.\\
		+ Lấy gia tốc trọng trường $g=10~\text{m/s}^2$.
		\end{minipage}\hfill
		\begin{minipage}[c]{0.36\linewidth}
			\centering
			\begin{tikzpicture}[scale=1.09,line join=round,line cap=round,font=\footnotesize]
				\def\ux{0.669}\def\uy{0.743}\def\vx{0.743}\def\vy{-0.669}
				% đường sức từ
				\foreach \i in {0,...,8}{
					\pgfmathsetmacro{\px}{(\i-4)*0.27*\vx}\pgfmathsetmacro{\py}{(\i-4)*0.27*\vy}
					\draw[VLxanh,semithick] ({\px-1.75*\ux},{\py-1.75*\uy})--({\px+1.75*\ux},{\py+1.75*\uy});
					\draw[VLxanh,VLmui,semithick] ({\px+1.2*\ux},{\py+1.2*\uy})--({\px+1.45*\ux},{\py+1.45*\uy});}
				\node[VLxanh!85!black] at ({4*0.27*\vx+1.95*\ux},{4*0.27*\vy+1.95*\uy}) {$\vv{B}$};
				% thanh MN
				\begin{scope}[shift={({-0.05*\ux},{-0.05*\uy})},rotate=-42]
					\shade[top color=VLdong!50!white,bottom color=VLdong!85!black,middle color=VLdong!20!white] (-1.55,-0.13) arc[start angle=270,end angle=90,radius=0.13]--(1.55,0.13)--(1.55,-0.13)--cycle;
					\draw[VLdong!40!black] (-1.55,-0.13) arc[start angle=270,end angle=90,radius=0.13]--(1.55,0.13) (-1.55,-0.13)--(1.55,-0.13);
					\filldraw[fill=VLdong!35!white,draw=VLdong!40!black] (1.55,0) ellipse (0.06 and 0.13);
					\node at (-1.75,0.35) {M};
					\node at (1.85,-0.3) {N};
				\end{scope}
			\end{tikzpicture}
		\end{minipage}
		\NeedspaceBang{15.8\baselineskip}%
		\choiceTFt
		{\True Xác định thông số vận hành cơ bản: Để con tàu có thể ``bay'' lơ lửng (trạng thái cân bằng tĩnh), cường độ dòng điện $I$ chạy qua thanh phải bằng $1$~A}
		{Xác định quy trình an toàn về chiều dòng điện: Biết các đường sức từ $\vv{B}$ hướng từ trái sang phải. Để tạo ra lực từ nâng tàu lên (ngược chiều trọng lực), dòng điện phải chạy theo chiều từ N đến M}
		{\True Xử lí tình huống sự cố hệ thống: Nếu do trục trặc kĩ thuật, cảm ứng từ $B$ bất ngờ bị giảm xuống chỉ còn $0{,}05$~T trong khi dòng điện vẫn giữ nguyên như ở ý a). Khi đó con tàu sẽ rơi xuống với gia tốc $5~\text{m/s}^2$}
		{\True Bài toán tối ưu hóa tải trọng: Giả sử toa tàu được chất thêm hàng hóa khiến tổng khối lượng tăng lên thành $25$~g/m. Nếu các kĩ sư muốn giữ nguyên cường độ dòng điện $I$ ở trên để tránh làm nóng dây dẫn thì họ phải điều chỉnh độ lớn của cảm ứng từ $B$ lên đến $0{,}25$~T để tàu vẫn lơ lửng được}
		\loigiai{
			Gọi $\lambda$ là khối lượng trên mỗi đơn vị chiều dài của thanh; lơ lửng khi $BIl=\lambda lg\Rightarrow BI=\lambda g$.
			\begin{itemchoice}
				\itemch \textbf{Đúng.} $I=\dfrac{\lambda g}{B}=\dfrac{0{,}01\cdot 10}{0{,}1}=1$~A.
				\itemch \textbf{Sai.} Hình vẽ nhìn từ trên xuống; theo quy tắc bàn tay trái, lực từ hướng lên khi dòng điện chạy từ M đến N.
				\itemch \textbf{Đúng.} $a=g-\dfrac{B'I}{\lambda}=10-\dfrac{0{,}05\cdot 1}{0{,}01}=5~\text{m/s}^2$.
				\itemch \textbf{Đúng.} $B'=\dfrac{\lambda' g}{I}=\dfrac{0{,}025\cdot 10}{1}=0{,}25$~T.
			\end{itemchoice}
		}
	\end{ex}

	\NeedspaceCau[27.3\baselineskip]{34.7\baselineskip}
	\begin{ex}%[VL12-C3-TH]
		% KIỂM TRA: dòng 3 của bảng ghi I = 0,46 A không khớp B = 0,018 T (F/I = 0,370 so với 0,30 và 0,275 ở hai lần đầu; có thể là 0,6 A); không ảnh hưởng các ý vì chỉ dùng cột B
		Để đo độ lớn cảm ứng từ, người ta lần lượt thực hiện các bước như sau:\\
		% giuhinh
		\par\noindent
		\begin{minipage}[c]{0.38\linewidth}
		Bước 1: Bố trí thí nghiệm như hình vẽ.\\
		Bước 2: Điều chỉnh để cho đòn cân nằm ngang. Đọc số chỉ $F_1$ của lực kế và ghi kết quả.
		\end{minipage}\hfill
		\begin{minipage}[c]{0.6\linewidth}
			\centering
			\begin{tikzpicture}[x=0.0092cm,y=-0.0092cm,line join=round,line cap=round,font=\footnotesize]
				% toạ độ theo điểm ảnh của ảnh gốc: bộ thí nghiệm đo cảm ứng từ (hộp gỗ, cân đòn, nam châm điện) và bộ nguồn
				\colorlet{YHgo}{VLgo!72!black}
				% bóng đổ trên mặt bàn
				\fill[gray!22] (182,517) ellipse (172 and 5);
				\fill[gray!26] (645,514) ellipse (312 and 6);
				%----- bộ nguồn (13)
				\fill[black!80] (38,512) rectangle (62,521) (296,512) rectangle (322,521);
				\shade[top color=gray!14,bottom color=white] (22,373)--(72,349)--(331,345)--(325,371)--cycle;
				\shade[left color=black!78,right color=black!55] (325,371)--(331,345)--(339,348)--(339,511)--(325,515)--cycle;
				\shade[top color=white,bottom color=gray!20] (22,373)--(325,371)--(325,515)--(26,517)--cycle;
				\draw[gray!65!black,line width=0.4pt] (22,373)--(325,371)--(325,515)--(26,517)--cycle (22,373)--(72,349)--(331,345)--(325,371);
				\draw[gray!45,line width=0.3pt] (30,381)--(318,379)--(318,508)--(33,510)--cycle;
				% quai xách
				\fill[gray!55] (108,357) ellipse (14 and 2.5) (262,353.5) ellipse (14 and 2.5);
				\draw[gray!60!black,line width=1.5pt] (118,357)--(128,356) .. controls (140,352) and (142,336) .. (153,334)--(226,333) .. controls (238,334) and (240,350) .. (250,353)--(268,353);
				\draw[gray!15,line width=0.5pt] (130,354.5) .. controls (140,351) and (143,335) .. (153,333)--(226,332);
				% chữ in trên mặt máy (vệt màu), ốc
				\draw[VLxanh!75!black,line width=1.1pt,line cap=butt] (188,390)--(272,389);
				\draw[VLdo!85,line width=0.9pt,line cap=butt] (277,389)--(315,388.5);
				\draw[VLxanh!60,line width=0.5pt,line cap=butt] (60,394)--(156,393) (92,400)--(104,400);
				\draw[gray!60,line width=0.4pt,line cap=butt] (70,410)--(125,409.5) (70,418)--(125,417.5);
				\foreach \x/\y in {65/382,277/376}{\shade[ball color=gray!55] (\x,\y) circle (3.2);}
				% bốn lỗ cắm, đèn báo, nút bấm
				\foreach \x/\y/\c in {73/450/VLdo!55!black,112/447/black!80,75/485/VLdo!55!black,112/483/VLdo!90!black}{
					\fill[gray!55] (\x,\y) circle (9.5);
					\fill[\c] (\x,\y) circle (7.5);
					\fill[black!85] (\x,\y) circle (3);
					\fill[white,opacity=0.5] (\x-3,\y-3.5) circle (1.6);}
				\draw[gray!60,line width=0.4pt,line cap=butt] (48,451)--(57,451) (51,486)--(59,486) (155,420)--(165,420) (155,472)--(166,472);
				\fill[gray!50] (182,418) circle (11.5);
				\shade[ball color=VLxanh!90!black] (182,418) circle (9.5);
				\fill[gray!55] (182,445) circle (8);
				\shade[ball color=black!70] (182,445) circle (6);
				\fill[gray!50] (184,470) circle (11.5);
				\shade[ball color=VLdo] (184,470) circle (9.5);
				% núm xoay chọn điện áp
				\foreach \g in {190,215,...,350}{\draw[black!70,line width=0.4pt] ({262+34*cos(\g)},{450+34*sin(\g)})--({262+40*cos(\g)},{450+40*sin(\g)});}
				\fill[black!88] (262,450) circle (31);
				\shade[ball color=black!75] (262,450) circle (27);
				\fill[white] (238,458) circle (2);
				%----- hộp gỗ
				\shade[top color=VLgo!88!black,bottom color=VLgo!60!black] (336,478)--(945,478)--(948,509)--(334,512)--cycle;
				\shade[top color=YHgo!92!white,bottom color=YHgo!88!black] (355,322)--(905,322)--(919,478)--(350,478)--cycle;
				\shade[left color=YHgo!70!black,right color=YHgo!60!black] (905,322)--(912,322)--(928,478)--(919,478)--cycle;
				\shade[top color=VLgo!96!black,bottom color=VLgo!80!black] (340,300) rectangle (912,322);
				\shade[top color=VLgo!70!white,bottom color=VLgo!92!white] (352,280)--(900,280)--(912,300)--(340,300)--cycle;
				\draw[YHgo!55!black,line width=0.4pt] (340,300)--(912,300)--(912,322)--(340,322)--cycle (336,478)--(945,478)--(948,509)--(334,512)--cycle (355,322)--(350,478) (912,322)--(928,478) (352,280)--(340,300) (900,280)--(912,300) (352,280)--(900,280);
				\draw[YHgo!65!black,line width=0.3pt] (360,336) .. controls (450,330) and (560,340) .. (700,334) .. controls (780,331) and (850,338) .. (905,334) (356,468) .. controls (480,462) and (600,472) .. (760,466) .. controls (830,464) and (880,470) .. (916,467) (343,311) .. controls (500,308) and (700,314) .. (910,310) (340,495) .. controls (500,490) and (700,498) .. (945,493);
				\draw[YHgo!80!black,line width=0.3pt] (640,322)--(640,478);
				% công tắc đảo chiều (8), (9)
				\foreach \x/\y/\dx in {415/401/-5,855/395/7}{
					\fill[yellow!45!VLgo!55!white,draw=YHgo!60!black,line width=0.3pt,rounded corners=0.8pt] (\x-13,\y-27) rectangle (\x+13,\y+27);
					\fill[VLdo!90!black] (\x-8,\y-21) rectangle (\x+8,\y-12) (\x-8,\y+12) rectangle (\x+8,\y+21);
					\fill[gray!65!black] (\x,\y) circle (7.5);
					\shade[ball color=gray!40] (\x,\y) circle (5.5);
					\draw[gray!25,line width=1.3pt] (\x,\y)--(\x+\dx,\y-3);
					\draw[gray!70!black,line width=0.3pt] (\x+\dx,\y-3) circle (1.6);}
				% hai ampe kế (5), (6)
				\foreach \x/\y in {445/345,715/342}{
					\fill[black,opacity=0.25] (\x+3,\y+4) rectangle (\x+123,\y+117);
					\fill[gray!38,draw=gray!70!black,line width=0.4pt] (\x,\y) rectangle (\x+120,\y+114);
					\shade[top color=white,bottom color=VLxanh!7] (\x+3,\y+3) rectangle (\x+117,\y+76);
					\shade[top color=gray!30,bottom color=gray!42] (\x+3,\y+76) rectangle (\x+117,\y+111);
					\draw[gray!60!black,line width=0.3pt] (\x+3,\y+76)--(\x+117,\y+76);
					\draw[black!75,line width=0.35pt] ({\x+60+64*cos(228)},{\y+84+64*sin(228)}) arc[start angle=228,end angle=312,radius=64];
					\foreach \g in {228,232.2,...,312.1}{\draw[black!70,line width=0.25pt] ({\x+60+64*cos(\g)},{\y+84+64*sin(\g)})--({\x+60+59*cos(\g)},{\y+84+59*sin(\g)});}
					\foreach \g in {228,249,270,291,312}{\draw[black!80,line width=0.4pt] ({\x+60+64*cos(\g)},{\y+84+64*sin(\g)})--({\x+60+55*cos(\g)},{\y+84+55*sin(\g)});}
					\draw[black,line width=0.4pt] ({\x+60+12*cos(230)},{\y+84+12*sin(230)})--({\x+60+60*cos(230)},{\y+84+60*sin(230)});
					\node[font=\fontsize{5}{5}\selectfont\bfseries,inner sep=0] at (\x+60,\y+48) {A};
					\draw[black,line width=0.3pt] (\x+55,\y+56)--(\x+65,\y+56);
					\fill[gray!36,draw=gray!60!black,line width=0.25pt] (\x+43,\y+76) .. controls (\x+44,\y+61) and (\x+76,\y+61) .. (\x+77,\y+76);
					\fill[black!80] (\x+11,\y+70) circle (3.2) (\x+109,\y+68) circle (3.2);
					\shade[ball color=black!70] (\x+60,\y+94) circle (6.5);}
				% hai biến trở xoay (7)
				\foreach \x/\y/\a/\b in {603/400/628/408,678/397/703/393}{
					\fill[black,opacity=0.3] (\x+2,\y+3) circle (28);
					\fill[black!90] (\x,\y) circle (28);
					\shade[ball color=black!72] (\x,\y) circle (24);
					\fill[white] (\a,\b) circle (2);}
				%----- cân đòn: đế, trụ đỡ, đòn cân (2), gia trọng (3), khớp nối (4)
				\fill[black,opacity=0.18] (606,284) ellipse (56 and 6);
				\shade[top color=white,bottom color=gray!30] (552,277) .. controls (552,268) and (656,268) .. (656,277) .. controls (656,288) and (552,288) .. (552,277)--cycle;
				\draw[gray!60,line width=0.3pt] (552,277) .. controls (552,268) and (656,268) .. (656,277) .. controls (656,288) and (552,288) .. (552,277);
				\shade[left color=gray!42,right color=gray!52,middle color=gray!6] (590,52) rectangle (626,276);
				\shade[left color=gray!55,right color=gray!30] (603,60) rectangle (613,224);
				\draw[gray!65!black,line width=0.3pt] (590,52)--(590,276) (626,52)--(626,276);
				\fill[VLdo!88] (597,227) rectangle (622,238);
				\fill[white] (597,238) rectangle (622,249);
				\fill[VLdo!88] (602,238) rectangle (607,243) (612,238) rectangle (617,243);
				% dây và phích cắm đỏ bên trụ
				\draw[black!75,line width=0.5pt] (592,180) .. controls (572,176) and (560,200) .. (578,228);
				\draw[VLdo!55!black,line width=0.5pt] (594,186) .. controls (578,186) and (568,204) .. (580,228);
				\shade[left color=VLdo!60!black,right color=VLdo!95!black,middle color=VLdo!80!white] (573,230) rectangle (583,270);
				% đòn cân
				\shade[top color=white,bottom color=gray!40] (466,43) rectangle (738,49.5);
				\draw[gray!65!black,line width=0.3pt] (466,43) rectangle (738,49.5);
				\fill[gray!65!black] (463,22) rectangle (467.5,51);
				\shade[top color=gray!25,bottom color=gray!60] (562,50)--(652,50)--(646,58)--(568,58)--cycle;
				\draw[gray!65!black,line width=0.3pt] (562,50)--(652,50)--(646,58)--(568,58)--cycle;
				\shade[left color=black!70,right color=black!55,middle color=gray!70] (592,18) rectangle (629,50);
				\fill[gray!35] (599,24) rectangle (622,44);
				\fill[black!75] (606,27) rectangle (615,36);
				\fill[black!75] (589,46) rectangle (594,58);
				% gia trọng
				\draw[black!72,line width=0.9pt,line cap=butt] (372,37)--(463,37);
				\shade[top color=VLgo!25!gray!45,bottom color=VLgo!25!gray!80,middle color=VLgo!12!gray!18] (392,20) rectangle (418,53);
				\draw[gray!65!black,line width=0.3pt] (392,20) rectangle (418,53);
				% khớp nối với khung dây
				\draw[gray!60!black,line width=0.7pt] (727,46)--(744,42);
				\fill[gray!55] (727,46) circle (2.4);
				\shade[top color=VLgo!20!gray!30,bottom color=VLgo!20!gray!75,middle color=gray!12] (742,21) rectangle (772,46);
				\draw[gray!65!black,line width=0.3pt] (742,21) rectangle (772,46);
				\fill[gray!65!black] (748,26) rectangle (756,41) (762,30) rectangle (767,38);
				\fill[gray!60!black] (750,46) rectangle (770,53);
				%----- lực kế (11) móc vào đòn cân, neo xuống mặt hộp
				\draw[black!70,line width=0.5pt] (472,49.5)--(472,63);
				\shade[left color=gray!30,right color=gray!35,middle color=white] (465,62) rectangle (480,176);
				\shade[left color=VLdo!70!black,right color=VLdo!75!black,middle color=VLdo!95!white] (465.5,72) rectangle (479.5,103);
				\foreach \y in {108,112,...,172}{\draw[black!65,line width=0.2pt] (466,\y)--(472,\y);}
				\foreach \y in {108,128,148,168}{\draw[black!80,line width=0.3pt] (466,\y)--(475,\y);}
				\draw[gray!70!black,line width=0.3pt] (465,62) rectangle (480,176);
				\draw[black!65,line width=0.4pt] (472,176)--(472,238);
				\draw[black!65,line width=0.35pt] (472,196) circle (2.8);
				\shade[top color=white,bottom color=VLgo!35,rounded corners=0.8pt] (459,237) rectangle (487,252);
				\draw[VLgo!60!black,line width=0.25pt,rounded corners=0.8pt] (459,237) rectangle (487,252);
				\fill[black!70] (469,252) rectangle (476,269);
				\shade[top color=black!50,bottom color=black!85] (459,268)--(496,267)--(498,288)--(461,289)--cycle;
				%----- nam châm điện (1): hai tấm thép, cuộn dây, đế (12); khung dây (10)
				\fill[black,opacity=0.2] (700,289)--(845,285)--(852,294)--(706,297)--cycle;
				\shade[left color=gray!62,right color=gray!48] (751,165) rectangle (767,200);
				\shade[left color=gray!58,right color=gray!45] (777,168) rectangle (799,201);
				\shade[top color=gray!12,bottom color=gray!35] (745,205)--(799,203)--(799,230)--(745,231)--cycle;
				% cuộn dây đồng
				\fill[black!75] (746,272) rectangle (764,291);
				\shade[top color=VLdong!85!white,bottom color=VLdong!60!black,middle color=VLdong!95!black] (748,230) rectangle (823,277);
				\foreach \x in {750.5,753.5,...,822}{\draw[VLdong!45!black,line width=0.2pt] (\x,231)--(\x,276.5);}
				\fill[white,opacity=0.35] (748,236) rectangle (823,242);
				% tấm thép trái
				\shade[left color=gray!75,right color=gray!52] (697,165) rectangle (741,291);
				\shade[left color=gray!55,right color=gray!72] (741,165)--(747,168)--(747,288)--(741,291)--cycle;
				\draw[gray!72!black,line width=0.35pt] (697,165) rectangle (741,291) (741,165)--(747,168)--(747,288)--(741,291);
				\foreach \x/\y in {706/191,725/191,706/265,724/263}{\fill[gray!30,draw=black!70,line width=0.3pt] (\x,\y) circle (3.6); \draw[black!70,line width=0.3pt] (\x-2,\y-1.5)--(\x,\y+2)--(\x+2,\y-1.5);}
				% tấm thép phải
				\shade[left color=gray!68,right color=gray!50] (797,165) rectangle (827,283);
				\shade[left color=gray!30,right color=gray!45] (827,165)--(839,167)--(839,280)--(827,283)--cycle;
				\foreach \x in {803,809,815,821}{\draw[gray!70,line width=0.25pt] (\x,167)--(\x,281);}
				\draw[gray!72!black,line width=0.35pt] (797,165) rectangle (827,283) (827,165)--(839,167)--(839,280)--(827,283);
				% đế gỗ dưới cuộn dây
				\shade[top color=VLgo!45!white,bottom color=VLgo!75!white] (758,279)--(838,276)--(843,291)--(762,294)--cycle;
				\draw[VLgo!60!black,line width=0.3pt] (758,279)--(838,276)--(843,291)--(762,294)--cycle;
				% khung dây (cạnh đáy nằm trong khe từ trường)
				\draw[gray!55!black,line width=0.5pt] (757,53)--(760,66);
				\fill[VLdong!55!black] (783,49) rectangle (790,206);
				\fill[VLdong!80!white] (784.5,52) rectangle (786,204);
				\fill[VLgo!45!white,draw=VLgo!70!black,line width=0.25pt] (738,76)--(738,70)--(784,51)--(790,51)--(790,57)--(748,76)--cycle;
				\fill[VLgo!40!white,draw=VLgo!70!black,line width=0.25pt] (738,74) rectangle (748,209);
				\fill[VLdong!60!black] (748,203) rectangle (783,208);
				%----- nhãn
				\draw[thin] (377,34)--(351,39) node[left,inner sep=1pt] {(3)};
				\draw[thin] (693,42)--(686,64) node[below,inner sep=1.5pt] {(2)};
				\draw[thin] (778,27)--(818,22) node[right,inner sep=1pt] {(4)};
				\draw[thin] (790,113)--(826,104) node[right,inner sep=1pt] {(10)};
				\draw[thin] (832,182)--(866,177) node[right,inner sep=1pt] {(1)};
				\draw[thin] (808,288)--(868,269) node[right,inner sep=1pt] {(12)};
				\draw[thin] (471,139)--(440,151) node[left,inner sep=1pt] {(11)};
				\draw[thin] (125,356)--(120,302) node[above,inner sep=1.5pt] {(13)};
				\draw[thin] (527,362)--(561,351);
				\node[white] at (592,341) {(5)};
				\draw[thin] (741,363)--(713,349);
				\node[white] at (685,341) {(6)};
				\draw[thin] (622,426)--(640,449)--(657,424);
				\node[white] at (641,468) {(7)};
				\draw[thin] (410,430)--(400,450);
				\node[white] at (397,465) {(8)};
				\draw[thin] (867,421)--(883,438);
				\node[white] at (888,454) {(9)};
			\end{tikzpicture}
		\end{minipage}
		\par\noindent
		Bước 3: Bật công tắc để dòng điện chạy qua khung dây dẫn và nam châm điện. Chọn chiều của dòng điện sao cho lực từ tác dụng lên cạnh của khung dây có hướng thẳng đứng từ trên xuống.\\
		+ Đọc số chỉ của ampe kế (5) và ghi kết quả.\\
		+ Điều chỉnh để đòn cân trở lại trạng thái cân bằng nằm ngang.\\
		+ Đọc số chỉ $F_2$ của lực kế và ghi kết quả.\\
		+ Xác định độ lớn lực từ tác dụng lên cạnh của khung dây đặt trong từ trường $F=F_2-F_1$. Ghi kết quả.\\
		Bước 4: Thực hiện lại bước 3 với ít nhất hai giá trị $I$ khác nhau, thu được kết quả như bảng dưới đây:
		\begin{center}
			\renewcommand{\arraystretch}{1.25}
			\begin{tabular}{|c|c|c|c|c|c|}
				\hline
				Lần & $I$ (A) & $F_1$ (N) & $F_2$ (N) & $F=F_2-F_1$ (N) & $B=\dfrac{F}{NIL}$ (T)\\
				\hline
				$1$ & $0{,}2$ & $0{,}210$ & $0{,}270$ & $0{,}060$ & $0{,}019$\\
				\hline
				$2$ & $0{,}4$ & $0{,}210$ & $0{,}320$ & $0{,}110$ & $0{,}017$\\
				\hline
				$3$ & $0{,}46$ & $0{,}210$ & $0{,}380$ & $0{,}170$ & $0{,}018$\\
				\hline
			\end{tabular}
		\end{center}
		\NeedspaceBang{8.1\baselineskip}%
		\choiceTFt
		{\True Phép đo độ lớn lực từ $F$ trong thí nghiệm trên là phép đo gián tiếp}
		{Trọng lượng của khung dây dẫn bằng $0{,}021$~N}
		{\True Kết quả đo: $B=(0{,}018\pm 0{,}001)$~T}
		{Sai số tỉ đối của phép đo bằng $6{,}5\,\%$}
		\loigiai{
			$\overline{B}=\dfrac{0{,}019+0{,}017+0{,}018}{3}=0{,}018$~T; $\overline{\Delta B}=\dfrac{0{,}001+0{,}001+0}{3}\approx 0{,}0007$~T, làm tròn $\Delta B=0{,}001$~T.
			\begin{itemchoice}
				\itemch \textbf{Đúng.} $F$ được tính từ hai số đo trực tiếp $F_1$, $F_2$ qua $F=F_2-F_1$.
				\itemch \textbf{Sai.} Khi chưa có dòng điện, lực kế chỉ $F_1=0{,}210$~N cân bằng trọng lượng khung dây (hai cánh tay đòn bằng nhau): $P=0{,}210$~N.
				\itemch \textbf{Đúng.} $B=\overline{B}\pm\Delta B=(0{,}018\pm 0{,}001)$~T.
				\itemch \textbf{Sai.} $\delta B=\dfrac{\Delta B}{\overline{B}}=\dfrac{0{,}001}{0{,}018}\approx 5{,}6\,\%$.
			\end{itemchoice}
		}
	\end{ex}

	\NeedspaceCau[11.4\baselineskip]{20.1\baselineskip}
	\begin{ex}%[VL12-C2-TH]
		Một học sinh tiến hành thí nghiệm để kiểm chứng định luật Boyle về quá trình biến đổi đẳng nhiệt của một lượng khí xác định (xem như lí tưởng) chứa trong xilanh kín. Dụng cụ thí nghiệm được bố trí như hình vẽ, trong đó pit-tông có thể dịch chuyển, áp kế được nối với vùng không gian chứa khí. Dịch chuyển pit-tông, học sinh này thu được bảng số liệu ở bảng dưới.
		% giuhinh
		\begin{center}
			\begin{minipage}[c]{0.5\linewidth}
				\centering\footnotesize
				\setlength{\tabcolsep}{3.5pt}
				\renewcommand{\arraystretch}{1.3}
				\begin{tabular}{|c|c|c|c|}
					\hline
					Lần đo & Thể tích $V$ (ml) & Áp suất $p$ (N/cm$^2$) & $p\cdot V$\\
					\hline
					$1$ & $28{,}0$ & $10{,}5$ & $294{,}00$\\
					\hline
					$2$ & $33{,}0$ & $9{,}0$ & $297{,}00$\\
					\hline
					$3$ & $39{,}5$ & $7{,}5$ & $296{,}25$\\
					\hline
				\end{tabular}
			\end{minipage}\hfill%
			\begin{minipage}[c]{0.46\linewidth}
				\centering
				\begin{tikzpicture}[x=0.0098cm,y=-0.0098cm,line join=round,line cap=round,font=\footnotesize]
					% toạ độ theo điểm ảnh của ảnh gốc: bộ thí nghiệm định luật Boyle (xilanh nằm hơi chếch, áp kế gắn trên khối đầu phải)
					% thanh ren giằng phía sau (trên)
					\draw[gray!62,line width=1pt,line cap=butt] (232,235)--(556,203);
					\draw[white,line width=0.4pt,line cap=butt,dash pattern=on 0.8pt off 1.1pt] (232,235)--(556,203);
					% núm kéo pit-tông (nhựa đen) và cần pit-tông
					\shade[top color=gray!30,bottom color=gray!75,middle color=white] (88,295.5)--(196,279.5)--(196,286)--(88,302.5)--cycle;
					\draw[gray!60!black,line width=0.3pt] (88,295.5)--(196,279.5) (88,302.5)--(196,286);
					\fill[black!80] (80,287)--(92,286)--(93,306)--(81,308)--cycle;
					\shade[left color=black!70,right color=black!92,middle color=gray!70!black] (49,268) .. controls (48,262) and (52,261) .. (58,261)--(74,259) .. controls (80,259) and (82,262) .. (82,268)--(86,326) .. controls (86,332) and (83,334) .. (78,334)--(62,336) .. controls (56,336) and (54,333) .. (53,327)--cycle;
					\begin{scope}
						\clip (49,268) .. controls (48,262) and (52,261) .. (58,261)--(74,259) .. controls (80,259) and (82,262) .. (82,268)--(86,326) .. controls (86,332) and (83,334) .. (78,334)--(62,336) .. controls (56,336) and (54,333) .. (53,327)--cycle;
						\foreach \y in {272,282,...,322}{\draw[black,line width=0.3pt,line cap=butt] (46,\y)--(90,\y-3);}
						\fill[gray!55,opacity=0.6] (55,258) rectangle (60,338);
					\end{scope}
					% ống xilanh trong suốt
					\def\ong{(243,246)--(566,212)--(571,291)--(248,320)--cycle}
					\shade[top color=gray!28,bottom color=gray!38,middle color=gray!6] \ong;
					% cần pit-tông bên trong ống
					\fill[gray!45] (246,279)--(456,256)--(456,262)--(246,286)--cycle;
					% pit-tông (thân tối, vòng đệm đỏ nâu)
					\shade[top color=black!55,bottom color=black!70,middle color=black!88] (452,226)--(487,222.5)--(491,296)--(456,299.5)--cycle;
					\shade[top color=VLdo!55!black,bottom color=VLdo!60!black,middle color=VLdo!85!black] (487,222.5)--(499,221.3)--(503,295)--(491,296)--cycle;
					% thước chia thể tích dán dọc ống
					\fill[white] (268,262)--(561,231.5)--(562.5,256)--(269.5,285)--cycle;
					\draw[gray!55,line width=0.3pt] (268,262)--(561,231.5) (269.5,285)--(562.5,256);
					\foreach \i in {1,...,39}{
						\draw[gray!55!black,line width=0.25pt] ({268+7.32*\i},{263.5-0.762*\i})--({268.4+7.32*\i},{271.5-0.762*\i});}
					\foreach \i in {0,5,...,40}{
						\draw[black!80,line width=0.45pt] ({268+7.32*\i},{263.5-0.762*\i})--({268.7+7.32*\i},{277-0.762*\i});}
					% lớp thuỷ tinh: vệt sáng dọc ống, mép ống
					\fill[white,opacity=0.85] (246,252)--(566,218)--(566.4,225)--(246.4,259)--cycle;
					\fill[black,opacity=0.12] (247.5,306)--(570.3,278)--(571,291)--(248,320)--cycle;
					\fill[white,opacity=0.55] (247.8,311)--(570.6,283)--(570.8,286.5)--(248,314.5)--cycle;
					\draw[gray!60!black,line width=0.5pt] (243,246)--(566,212) (248,320)--(571,291);
					% thanh ren giằng phía trước (dưới)
					\draw[gray!62,line width=1pt,line cap=butt] (238,333)--(580,302);
					\draw[white,line width=0.4pt,line cap=butt,dash pattern=on 0.8pt off 1.1pt] (238,333)--(580,302);
					% tấm mica đầu trái (có đế)
					\fill[black!60] (201,320)--(250,315)--(253,335)--(205,341)--cycle;
					\fill[gray!55] (203,324)--(250,319)--(250.5,323)--(203.5,328)--cycle;
					\shade[left color=gray!25,right color=gray!45,middle color=white,opacity=0.9] (188,237)--(238,232)--(246,319)--(197,325)--cycle;
					\fill[gray!20,opacity=0.9] (188,237)--(199,229)--(248,224.5)--(238,232)--cycle;
					\fill[gray!45,opacity=0.9] (238,232)--(248,224.5)--(255,311)--(246,319)--cycle;
					\draw[gray!60!black,line width=0.4pt] (188,237)--(238,232)--(246,319)--(197,325)--cycle (188,237)--(199,229)--(248,224.5)--(238,232) (248,224.5)--(255,311)--(246,319);
					\fill[white,opacity=0.8] (193,243)--(202,242)--(209,318)--(200,319)--cycle;
					\draw[gray!55,line width=0.3pt] (212,262)--(240,259) (213,275)--(241,272) (215,300)--(243,297);
					% khối mica đầu phải (có lỗ dẫn khí, van)
					\shade[left color=gray!30,right color=gray!50] (546,200)--(574,222)--(574,300)--(546,277)--cycle;
					\shade[top color=gray!22,bottom color=gray!42,middle color=gray!10] (574,222)--(720,211)--(720,286)--(574,300)--cycle;
					\shade[left color=gray!12,right color=gray!28] (546,200)--(692,190)--(720,211)--(574,222)--cycle;
					\fill[gray!60,opacity=0.55] (574,247)--(672,239)--(672,251)--(574,260)--cycle;
					\fill[gray!60,opacity=0.55] (606,219.5)--(622,218.3)--(622,246)--(606,247.5)--cycle;
					\fill[gray!50!black,opacity=0.55] (598,283)--(624,281)--(624,295)--(598,297.5)--cycle;
					\fill[white,opacity=0.75] (578,226.5)--(716,216)--(716,220)--(578,230.5)--cycle;
					\filldraw[fill=gray!35,draw=gray!60!black,line width=0.3pt] (640,268) circle (5);
					\filldraw[fill=gray!35,draw=gray!60!black,line width=0.3pt] (589,236) circle (3.5);
					\draw[gray!60!black,line width=0.4pt] (546,200)--(574,222)--(574,300)--(546,277)--cycle (574,222)--(720,211)--(720,286)--(574,300) (546,200)--(692,190)--(720,211);
					% van khoá (nhựa đen)
					\fill[black!85] (676,241) .. controls (680,238) and (686,240) .. (692,246)--(700,236) .. controls (703,232) and (708,234) .. (708,239)--(709,250) .. controls (716,252) and (718,258) .. (717,266)--(716,283) .. controls (712,289) and (700,290) .. (694,287)--(690,262) .. controls (683,256) and (676,250) .. (676,241)--cycle;
					\fill[gray!55,opacity=0.7] (703,253)--(707,253)--(708,282)--(704,282)--cycle;
					% áp kế: cổ nối, vỏ đen, mặt số
					\shade[left color=gray!70!black,right color=gray!45,middle color=gray!20] (597,186)--(628,186)--(628,207)--(597,209)--cycle;
					\fill[black!88] (612,104) ellipse (89 and 94);
					\shade[left color=gray!55,right color=gray!35] (634,105) ellipse (87 and 93);
					\shade[inner color=white,outer color=gray!12] (634,105) ellipse (80 and 86);
					\draw[gray!55!black,line width=0.3pt] (634,105) ellipse (80 and 86);
					\foreach \i in {0,...,50}{
						\draw[black!85,line width=0.25pt] ({634+76*cos(225-5.4*\i)},{105-82*sin(225-5.4*\i)})--({634+70.5*cos(225-5.4*\i)},{105-76*sin(225-5.4*\i)});}
					\foreach \i in {0,5,...,25}{
						\draw[black,line width=0.55pt] ({634+76*cos(225-10.8*\i)},{105-82*sin(225-10.8*\i)})--({634+66*cos(225-10.8*\i)},{105-71*sin(225-10.8*\i)});
					}
					\foreach \i/\g in {0/225,5/171,10/126,15/63,20/9,25/-45}{
						\node[font=\fontsize{5.5}{5.5}\selectfont,inner sep=0] at ({634+50*cos(\g)},{105-55*sin(\g)}) {\i};}
					\node[font=\fontsize{4.5}{4.5}\selectfont,inner sep=0] at (636,160) {N/cm$^2$};
					% kim chỉ 10,5 N/cm^2 (số liệu lần đo 1), đuôi kim có đối trọng
					\fill[black!85] ({634+72*cos(111.6)},{105-78*sin(111.6)})--({634+5*cos(21.6)},{105-5*sin(21.6)})--({634+30*cos(291.6)+9*cos(21.6)},{105-32*sin(291.6)-9*sin(21.6)})--({634+30*cos(291.6)-9*cos(21.6)},{105-32*sin(291.6)+9*sin(21.6)})--({634-5*cos(21.6)},{105+5*sin(21.6)})--cycle;
					\fill[gray!40] (634,105) circle (3.2);
					\fill[white,opacity=0.5] (590,40) .. controls (610,24) and (650,20) .. (675,34) .. controls (650,28) and (615,32) .. (598,46)--cycle;
				\end{tikzpicture}
			\end{minipage}
		\end{center}
		\NeedspaceBang{9\baselineskip}%
		\choiceTFt
		{\True Áp kế là dụng cụ dùng để đo áp suất của lượng khí chứa trong xilanh}
		{Để giữ nhiệt độ khí không đổi, học sinh này phải dịch chuyển pit-tông thật nhanh}
		{\True Khi dịch chuyển pit-tông làm tăng thể tích khí thì số chỉ của áp kế sẽ giảm}
		{Theo định luật Boyle, khi thể tích khí là $56$~ml thì số chỉ của áp kế là $10{,}5$~N/cm$^2$}
		\loigiai{
			\begin{itemchoice}
				\itemch \textbf{Đúng.} Áp kế nối với vùng chứa khí, đo áp suất của khí.
				\itemch \textbf{Sai.} Phải dịch chuyển pit-tông thật chậm để khí kịp trao đổi nhiệt với môi trường, nhiệt độ khí không đổi.
				\itemch \textbf{Đúng.} Đẳng nhiệt: $pV=\text{hằng số}$, $V$ tăng thì $p$ giảm.
				\itemch \textbf{Sai.} $p=\dfrac{p_1V_1}{V}=\dfrac{10{,}5\cdot 28}{56}=5{,}25$~N/cm$^2$.
			\end{itemchoice}
		}
	\end{ex}

	\NeedspaceCau{14.9\baselineskip}
	\begin{ex}%[VL12-C1-TH]
		Bạn Linh pha chế một mẫu trà sữa bằng cách trộn các mẫu chất với nhau: nước trà đen (mẫu A), nước đá (mẫu B) và sữa tươi (mẫu C). Các mẫu chất này chỉ trao đổi nhiệt lẫn nhau mà không gây ra các phản ứng hóa học. Bỏ qua sự trao đổi nhiệt với môi trường và bình chứa. Nhiệt độ trước khi trộn của mẫu A, mẫu B và mẫu C lần lượt là $80\,^\circ\text{C}$, $-5\,^\circ\text{C}$ và $20\,^\circ\text{C}$.
		\NeedspaceBang{11\baselineskip}%
		\choiceTFt
		{\True Khi để cốc trà sữa ngoài không khí một thời gian bạn Linh thấy xung quanh thành bên ngoài cốc có một lớp nước mỏng, chứng tỏ có sự ngưng tụ hơi nước lên trên thành bên ngoài cốc}
		{\True Nhiệt độ cân bằng của hệ khi trộn cả ba mẫu với nhau nhỏ hơn $80\,^\circ\text{C}$}
		{Ban đầu các mẫu A, B, C đều ở thể lỏng}
		{Khi pha chế, bạn Linh trộn $3$ mẫu vào nhau, mẫu B có sự chuyển thể từ thể lỏng sang thể rắn}
		\loigiai{
			\begin{itemchoice}
				\itemch \textbf{Đúng.} Hơi nước trong không khí gặp thành cốc lạnh thì ngưng tụ thành lớp nước mỏng.
				\itemch \textbf{Đúng.} Mẫu A ($80\,^\circ\text{C}$) nóng nhất, tỏa nhiệt cho B, C nên nhiệt độ cân bằng nhỏ hơn $80\,^\circ\text{C}$.
				\itemch \textbf{Sai.} Mẫu B là nước đá ở $-5\,^\circ\text{C}$, ở thể rắn.
				\itemch \textbf{Sai.} Nước đá nhận nhiệt, nóng chảy: chuyển từ thể rắn sang thể lỏng.
			\end{itemchoice}
		}
	\end{ex}
	\Closesolutionfile{ans}

	\NeedspaceCau{18.8\baselineskip}
	\subsection*{\indam{PHẦN III. Câu trắc nghiệm trả lời ngắn. Thí sinh trả lời từ câu 1 đến câu 6.}}
	\setcounter{ex}{0}
	\Opensolutionfile{ans}[ans/VatLi2026-HaNoi-YenHoa-Lan1-3]

	\begin{ex}%[VL12-C2-VD]
		Để thuận tiện rút thuốc từ lọ thuốc kín, y tá thường sử dụng ống tiêm để bơm một lượng nhỏ khí vào lọ thuốc như hình vẽ. Một lọ thuốc có thể tích $0{,}9$~ml và chứa $0{,}5$~ml thuốc, áp suất của khí trong lọ là $10^5$~Pa. Một lượng khí trong ống tiêm có tiết diện $0{,}3~\text{cm}^2$, dài $0{,}4$~cm và áp suất $10^5$~Pa được y tá bơm vào lọ thuốc. Biết nhiệt độ bên trong và bên ngoài lọ thuốc bằng nhau và không thay đổi. Áp suất của lượng khí mới trong lọ thuốc là $a\cdot 10^5$~Pa. Tìm $a$ (làm tròn kết quả đến chữ số hàng phần mười).
		\begin{center}
			\begin{tikzpicture}[scale=0.7,line join=round,line cap=round]
				% lọ thuốc (úp ngược)
				\fill[VLxanh!8] (-0.55,4.85) rectangle (0.55,5.45);
				\fill[VLxanh!25] (-0.55,4.85)--(-0.55,4.35)--(-0.22,4.05)--(0.22,4.05)--(0.55,4.35)--(0.55,4.85)--cycle;
				\foreach \x/\y in {-0.4/4.7,-0.2/4.6,0/4.75,0.2/4.62,0.4/4.72,-0.3/4.42,-0.1/4.3,0.1/4.45,0.3/4.38,-0.4/4.55,0.05/4.15,0.25/4.22,-0.15/4.2,0.42/4.5,-0.02/4.53}{\fill[VLxanh!80!black] (\x,\y) circle (0.025);}
				\draw[very thick] (-0.55,5.45)--(-0.55,4.35)--(-0.22,4.05)--(0.22,4.05)--(0.55,4.35)--(0.55,5.45);
				\draw[thick] (-0.62,5.45)--(0.62,5.45);
				\draw[VLxanh!60!black,thin] (-0.55,4.85)--(0.55,4.85);
				\shade[top color=gray!40,bottom color=gray!70] (-0.28,3.8) rectangle (0.28,4.05);
				\draw[thick] (-0.28,3.8) rectangle (0.28,4.05);
				% kim tiêm
				\draw[semithick,gray!70!black] (0,4.4)--(0,3.05);
				% ống tiêm
				\fill[black] (0,2.95) ellipse (0.15 and 0.12);
				\fill[gray!12] (-0.45,0.55) rectangle (0.45,2.75);
				\fill[white,opacity=0.8] (-0.32,0.7) rectangle (-0.22,2.65);
				\draw[very thick] (-0.45,2.75)--(-0.45,0.4) (0.45,2.75)--(0.45,0.4);
				\shade[top color=white,bottom color=gray!40] (-0.7,2.75) rectangle (0.7,2.95);
				\draw[thick] (-0.7,2.75) rectangle (0.7,2.95);
				\fill[black] (-0.45,0.4) rectangle (0.45,0.6);
				\draw[thick] (-0.12,0.4)--(-0.12,0.05) (0.12,0.4)--(0.12,0.05);
				\shade[top color=white,bottom color=gray!40] (-0.75,-0.1) rectangle (0.75,0.05);
				\draw[thick] (-0.75,-0.1) rectangle (0.75,0.05);
			\end{tikzpicture}
		\end{center}
		\shortans[oly]{1{,}3}
		\loigiai{
			Khí trong lọ: $V_1=0{,}9-0{,}5=0{,}4~\text{cm}^3$; khí trong ống tiêm: $V_2=0{,}3\cdot 0{,}4=0{,}12~\text{cm}^3$.\\
			Đẳng nhiệt: $p_0(V_1+V_2)=pV_1\Rightarrow p=\dfrac{10^5\cdot(0{,}4+0{,}12)}{0{,}4}=1{,}3\cdot 10^5~\text{Pa}\Rightarrow a=1{,}3$.
		}
	\end{ex}

	\NeedspaceCau{12.5\baselineskip}
	\begin{ex}%[VL12-C1-VD]
		Quả cầu có nhiệt dung riêng $c=460$~J/(kg.K) được treo bởi sợi dây có chiều dài $l=46$~cm. Quả cầu được nâng lên đến B rồi thả rơi. Sau khi chạm tường, nó bật lên đến C ($\alpha=60^\circ$). Biết rằng $60\,\%$ độ giảm thế năng giữa hai vị trí B và C biến thành nhiệt làm nóng quả cầu. Lấy $g=10~\text{m/s}^2$. Độ tăng nhiệt độ của quả cầu là $a\cdot 10^{-3}$~K. Tính $a$ (làm tròn kết quả đến chữ số hàng đơn vị).
		\begin{center}
			\begin{tikzpicture}[scale=0.85,line join=round,line cap=round,font=\footnotesize]
				\fill[pattern=north east lines,pattern color=gray] (-0.2,-0.2) rectangle (0,2.3);
				\draw[line width=1.6pt] (0,-0.2)--(0,2.3);
				\draw[thick] (0,2)--(2,2);
				\node[above,inner sep=2pt] at (1,2) {$l$};
				\draw[dashed,thick] (0,2)--(0,0);
				\draw[dashed,thick] (0,2)--({2*sin(60)},{2-2*cos(60)});
				\draw[dashed,thick] (2,2) arc[start angle=0,end angle=-90,radius=2];
				\draw[thin] (0,1.45) arc[start angle=-90,end angle=-30,radius=0.55];
				\node at ($(0,2)+(-60:0.8)$) {$\alpha$};
				\shade[ball color=gray!70] (2,2) circle (0.1);
				\shade[ball color=gray!70] ({2*sin(60)},{2-2*cos(60)}) circle (0.1);
				\shade[ball color=gray!70] (0.1,0) circle (0.1);
				\node[above right,inner sep=2pt] at (2,2) {B};
				\node[below right,inner sep=2pt] at ({2*sin(60)},{2-2*cos(60)}) {C};
				\node[below right,inner sep=2pt] at (0.1,-0.05) {A};
			\end{tikzpicture}
		\end{center}
		\shortans[oly]{3}
		\loigiai{
			B ngang điểm treo, C thấp hơn điểm treo $l\cos\alpha$: $\Delta h=l\cos 60^\circ=0{,}23$~m.\\
			$0{,}6\,mg\Delta h=mc\Delta t\Rightarrow \Delta t=\dfrac{0{,}6\cdot 10\cdot 0{,}23}{460}=3\cdot 10^{-3}~\text{K}\Rightarrow a=3$.
		}
	\end{ex}

	\NeedspaceCau{7.6\baselineskip}
	\begin{ex}%[VL12-C2-VD]
		Một bình chứa oxygen (O$_2$) sử dụng trong y tế có thể tích $14$ lít, áp suất $15\cdot 10^6$~Pa và nhiệt độ phòng $27\,^\circ\text{C}$. Theo thông tin từ Bộ Y tế, thông thường đối với một bệnh nhân mắc bệnh hô hấp được chỉ định dùng liệu pháp oxygen thì người bệnh cần được cung cấp trung bình $6$ lít oxygen (ở điều kiện tiêu chuẩn, nhiệt độ $0\,^\circ\text{C}$, áp suất $1{,}013\cdot 10^5$~Pa) trong $15$ phút. Với tốc độ thở như vậy, người đó dùng hết bình oxygen $14$ lít sau bao nhiêu giờ (làm tròn kết quả đến chữ số hàng phần mười)?
		\par
		\shortans[oly]{78{,}6}
		\loigiai{
			Thể tích oxygen ở điều kiện tiêu chuẩn: $V_0=\dfrac{pVT_0}{p_0T}=\dfrac{15\cdot 10^6\cdot 14\cdot 273}{1{,}013\cdot 10^5\cdot 300}\approx 1886{,}5$ lít.\\
			Mỗi giờ dùng $4\cdot 6=24$ lít: $t=\dfrac{1886{,}5}{24}\approx 78{,}6$~h.
		}
	\end{ex}

	\NeedspaceCau{5.5\baselineskip}
	\begin{ex}%[VL12-C2-TH]
		Một hộp hình lập phương có cạnh $10$~cm chứa khí lí tưởng ở nhiệt độ $20\,^\circ\text{C}$ và áp suất $1{,}2\cdot 10^6$~Pa. Cho số Avogadro $N_A=6{,}02\cdot 10^{23}~\text{mol}^{-1}$, hằng số Boltzmann $k=1{,}38\cdot 10^{-23}$~J/K. Số phân tử khí trong hộp là $a\cdot 10^{23}$ phân tử. Tìm $a$ (làm tròn kết quả đến chữ số hàng phần trăm).
		\par
		\shortans[oly]{2{,}97}
		\loigiai{
			$pV=NkT\Rightarrow N=\dfrac{pV}{kT}=\dfrac{1{,}2\cdot 10^6\cdot 10^{-3}}{1{,}38\cdot 10^{-23}\cdot 293}\approx 2{,}97\cdot 10^{23}\Rightarrow a=2{,}97$.
		}
	\end{ex}

	\NeedspaceCau{20.5\baselineskip}
	\begin{ex}%[VL12-C3-VD]
		Bạn Trang đã tiến hành thí nghiệm như hình để xác định cảm ứng từ $B$ trong lòng của nam châm. Nam châm được đặt trên cân điện tử. PQ là một thanh cứng thẳng dẫn điện, đặt cố định nằm ngang, vuông góc với từ trường giữa các cực của nam châm và được nối với nguồn điện. Chiều dài của nam châm $l=15$~cm, coi từ trường trong lòng nam châm là đều, lực từ tác dụng lên phần thanh PQ ở bên ngoài nam châm là không đáng kể. Tăng dần cường độ dòng điện $I$ chạy trong dây PQ và ghi lại số chỉ $m$ của cân, bạn Trang vẽ được đồ thị $m$ theo $I$ như hình vẽ. Dùng thước đo góc, bạn xác định được $\alpha=28^\circ$. Lấy $g=9{,}8~\text{m/s}^2$. Độ lớn cảm ứng từ đo được trong lòng của nam châm là bao nhiêu mT (làm tròn kết quả đến chữ số hàng phần mười)?
		% giuhinh
		\begin{center}
			\begin{tikzpicture}[line join=round,line cap=round,font=\footnotesize]
				% bộ thí nghiệm (toạ độ theo điểm ảnh của ảnh gốc): cân điện tử, nam châm chữ U đặt trên cân, thanh PQ kẹp trên hai giá
				\begin{scope}[x=0.0115cm,y=-0.0115cm]
					\colorlet{YHcan}{VLxanh!20!black}
					% giá phải (đứng sau cân): trụ đứng, tay ngang
					\shade[left color=gray!65,right color=gray!45,middle color=gray!12] (328,33) rectangle (337,205);
					\shade[top color=gray!15,bottom color=gray!60] (333,86)--(391,121)--(389,127)--(331,92)--cycle;
					% thân cân: gờ sau, mặt trên, mặt trước trái, mặt nghiêng có màn hình
					\fill[YHcan!70!gray] (372,192)--(418,220)--(418,238)--(372,212)--cycle;
					\fill[YHcan!45!gray] (352,182)--(372,192)--(418,220)--(400,214)--cycle;
					\shade[left color=VLxanh!10!gray!75,right color=VLxanh!10!gray!95] (153,250)--(293,168)--(465,265)--(325,347)--cycle;
					\shade[top color=YHcan!88!white,bottom color=YHcan] (153,250)--(325,347)--(350,400)--(350,423)--(153,305)--cycle;
					\shade[top color=YHcan!62!gray,bottom color=YHcan!82!gray] (325,347)--(465,265)--(492,335)--(350,400)--cycle;
					\fill[YHcan!95!white] (350,400)--(492,335)--(492,346)--(350,423)--cycle;
					\draw[black!85,line width=0.3pt] (153,250)--(325,347)--(465,265) (325,347)--(350,400)--(492,335)--(465,265) (350,400)--(350,423)--(492,346)--(492,335) (153,250)--(153,305)--(350,423);
					\draw[white,opacity=0.35,line width=0.4pt] (155,251.5)--(325,348.5)--(464,267);
					% màn hình (để trống), các phím
					\filldraw[fill=VLla!18!white,draw=black!80,line width=0.3pt] (345,346)--(452,289)--(463,317)--(356,375)--cycle;
					\node[font=\fontsize{4.5}{4.5}\selectfont,inner sep=0] at (447,301) {g};
					\fill[gray!55] (439,331)--(466,319)--(468,324)--(441,336)--cycle;
					\fill[gray!8] (361,384)--(392,369)--(394,374)--(363,389)--cycle;
					% nam châm chữ U: thành sau (cực S), đáy
					\fill[gray!22] (240,268)--(300,298)--(300,236)--cycle;
					\shade[top color=gray!6,bottom color=gray!25] (240,228)--(318,186)--(318,226)--(240,268)--cycle;
					\shade[top color=VLxanh!75!white,bottom color=VLxanh!95!black] (240,180)--(318,138)--(318,186)--(240,228)--cycle;
					\fill[white] (240,227)--(318,185)--(318,190)--(240,232)--cycle;
					\draw[black!70,line width=0.3pt] (240,180)--(318,138)--(318,226) (240,180)--(240,268)--(300,298);
					\node[white,font=\bfseries] at (280,186) {S};
					% thanh PQ
					\draw[black,line width=1.9pt,line cap=butt] (166,291)--(440,108.5);
					% thành trước (cực N): mặt ngoài sáng, cạnh trên và cạnh bên màu đỏ
					\shade[left color=gray!28,right color=gray!10,middle color=white] (305,221)--(378,179)--(378,253)--(305,298)--cycle;
					\fill[VLdo!85!black] (298,216)--(372,175)--(378,179)--(305,221)--cycle;
					\fill[VLdo!70!black] (298,216)--(305,221)--(305,298)--(298,294)--cycle;
					\draw[black!60,line width=0.3pt] (305,221)--(378,179)--(378,253)--(305,298)--cycle;
					\node[font=\bfseries,black!75] at (344,223) {N};
					% kẹp P, Q
					\foreach \x/\y in {191/262,385/126}{
						\fill[YHcan!60!gray] (\x,\y)--(\x+16,\y-9)--(\x+35,\y+2)--(\x+19,\y+11)--cycle;
						\fill[YHcan] (\x,\y)--(\x+19,\y+11)--(\x+19,\y+27)--(\x,\y+16)--cycle;
						\fill[YHcan!85!gray] (\x+19,\y+11)--(\x+35,\y+2)--(\x+35,\y+18)--(\x+19,\y+27)--cycle;
						\draw[black,line width=0.25pt] (\x,\y)--(\x+16,\y-9)--(\x+35,\y+2)--(\x+35,\y+18)--(\x+19,\y+27)--(\x,\y+16)--cycle (\x,\y)--(\x+19,\y+11)--(\x+35,\y+2) (\x+19,\y+11)--(\x+19,\y+27);}
					\draw[black,line width=1.9pt,line cap=butt] (421,121)--(440,108.5);
					% giá trái: đế, trụ đứng, tay ngang
					\fill[YHcan!55!gray] (88,347)--(138,311)--(218,357)--(170,393)--cycle;
					\fill[YHcan] (88,347)--(170,393)--(170,410)--(88,363)--cycle;
					\fill[YHcan!85!gray] (170,393)--(218,357)--(218,373)--(170,410)--cycle;
					\draw[black,line width=0.25pt] (88,347)--(138,311)--(218,357)--(218,373)--(170,410)--(88,363)--cycle (88,347)--(170,393)--(218,357) (170,393)--(170,410);
					\shade[top color=gray!15,bottom color=gray!60] (135,225)--(193,259)--(191,265)--(133,231)--cycle;
					\shade[left color=gray!65,right color=gray!45,middle color=gray!12] (131,173) rectangle (138.5,334);
					\fill[black,opacity=0.25] (135,334) ellipse (6 and 2.5);
					% dây nối
					\draw[black!85,line width=0.8pt] (166,291) .. controls (158,297) and (140,302) .. (122,299) .. controls (95,294) and (78,305) .. (70,332) .. controls (63,360) and (60,388) .. (38,392);
					\draw[black!85,line width=0.8pt] (440,108.5) .. controls (462,98) and (468,130) .. (469,150) .. controls (470,178) and (492,172) .. (505,166) .. controls (522,158) and (530,178) .. (518,190) .. controls (505,202) and (478,186) .. (452,199);
					% nhãn
					\node[white] at (166,275) {P};
					\node at (438,88) {Q};
					\node at (66,286) {$I$};
					\draw[dashed,thin] (240,178)--(240,120) (312,139)--(312,74);
					\draw[<->,>=stealth,thin] (241,151)--(311,108);
					\node at (268,113) {$l$};
				\end{scope}
				% đồ thị m theo I
				\begin{scope}[shift={(7.25cm,-4.4cm)},scale=0.85]
					\draw[VLmui,thick] (0,0)--(4.2,0) node[below,inner sep=3pt] {$I$(A)};
					\draw[VLmui,thick] (0,0)--(0,3.45) node[left,inner sep=2pt] {$m$(g)};
					\node[below left,inner sep=2pt] at (0,0) {O};
					\draw[dashed] (0,0.9)--(4.0,0.9);
					\node[left,inner sep=2pt] at (0,0.9) {$80$};
					\foreach \x in {1.0,1.6,2.35,3.15,3.6}{
						\draw[dashed,gray,thin] (\x,0)--(\x,{0.9+0.5317*\x}) (0,{0.9+0.5317*\x})--(\x,{0.9+0.5317*\x});}
					\draw[very thick,VLdo] (0,0.9)--(4.0,{0.9+0.5317*4.0});
					\foreach \x in {1.0,1.6,2.35,3.15,3.6}{\filldraw[fill=black] (\x,{0.9+0.5317*\x}) circle (1.8pt);}
					\filldraw[fill=black] (0.5,1.11) circle (1.8pt);
					\draw[thin] (0.65,0.9) arc[start angle=0,end angle=28,radius=0.65];
					\node at (0.88,1.06) {$\alpha$};
				\end{scope}
			\end{tikzpicture}
		\end{center}
		\shortans[oly]{34{,}7}
		\loigiai{
			Lực từ lên thanh $F=BIl$; phản lực của lực này tác dụng lên nam châm làm số chỉ của cân tăng: $\Delta m\cdot g=BIl\Rightarrow m=m_0+\dfrac{Bl}{g}I$.\\
			Hệ số góc của đồ thị: $\dfrac{Bl}{g}=\tan 28^\circ$~g/A $=0{,}5317\cdot 10^{-3}$~kg/A.
			\[B=\dfrac{g\tan 28^\circ}{l}=\dfrac{9{,}8\cdot 0{,}5317\cdot 10^{-3}}{0{,}15}\approx 34{,}7\cdot 10^{-3}~\text{T}=34{,}7~\text{mT}.\]
		}
	\end{ex}

	\NeedspaceCau{8.6\baselineskip}
	\begin{ex}%[VL12-C1-VD]
		Hương vị của bia Hà Nội đã trở thành một thương hiệu mà nhiều người yêu thích. Mở nắp một chai bia rồi rót $200$~g bia vào cốc. Cho vào cốc $40$~g nước đá ở nhiệt độ $-2{,}8\,^\circ\text{C}$ thì ta sẽ được một cốc bia mát. Biết nhiệt dung riêng của bia là $3830$~J/(kg.K); nhiệt dung riêng của nước đá là $1800$~J/(kg.K); nhiệt nóng chảy riêng của nước đá là $3{,}4\cdot 10^5$~J/kg; nhiệt dung riêng của nước là $4200$~J/(kg.K). Ban đầu bia có nhiệt độ là $32\,^\circ\text{C}$. Bỏ qua sự trao đổi nhiệt với môi trường và sự trao đổi nhiệt với thành cốc. Sau khi nước đá tan hết, nhiệt độ của cốc bia là bao nhiêu $^\circ$C (làm tròn kết quả đến chữ số hàng phần mười)?
		\par
		\shortans[oly]{11{,}5}
		\loigiai{
			Cân bằng nhiệt: $m_1c_1(32-t)=m_2c_{\text{đ}}\cdot 2{,}8+\lambda m_2+m_2c_{\text{n}}t$
			\[0{,}2\cdot 3830\,(32-t)=0{,}04\cdot 1800\cdot 2{,}8+3{,}4\cdot 10^5\cdot 0{,}04+0{,}04\cdot 4200\,t\Rightarrow t\approx 11{,}5\,^\circ\text{C}.\]
		}
	\end{ex}

	\centerline{--- HẾT ---}
	\Closesolutionfile{ans}

	%===================== BẢNG ĐÁP ÁN =====================
	%\newpage
	%\begin{center}\indam{BẢNG ĐÁP ÁN}\end{center}
	%\noindent\indam{PHẦN I.}\\
	%\inputansbox{6}{ans/VatLi2026-HaNoi-YenHoa-Lan1-1}
	%\noindent\indam{PHẦN II.}\\
	%\inputansbox{2}{ans/VatLi2026-HaNoi-YenHoa-Lan1-2}
	%\noindent\indam{PHẦN III.}\\
	%\inputansbox{6}{ans/VatLi2026-HaNoi-YenHoa-Lan1-3}
\end{document}
